%%%%%%%%%%%%%%%%%%%%%%%%%%%%%%%%%%%%%%%%%
% Masters/Doctoral Thesis 
% LaTeX Template
% Version 2.5 (27/8/17)
%
% This template was downloaded from:
% http://www.LaTeXTemplates.com
%
% Version 2.x major modifications by:
% Vel (vel@latextemplates.com)
%
% This template is based on a template by:
% Steve Gunn (http://users.ecs.soton.ac.uk/srg/softwaretools/document/templates/)
% Sunil Patel (http://www.sunilpatel.co.uk/thesis-template/)
%
% Template license:
% CC BY-NC-SA 3.0 (http://creativecommons.org/licenses/by-nc-sa/3.0/)
%
%%%%%%%%%%%%%%%%%%%%%%%%%%%%%%%%%%%%%%%%%

%----------------------------------------------------------------------------------------
%	PACKAGES AND OTHER DOCUMENT CONFIGURATIONS
%----------------------------------------------------------------------------------------

\documentclass[
12pt, % The default document font size, options: 10pt, 11pt, 12pt
%oneside, % Two side (alternating margins) for binding by default, uncomment to switch to one side
english, % ngerman for German
singlespacing, % Single line spacing, alternatives: onehalfspacing or doublespacing
%draft, % Uncomment to enable draft mode (no pictures, no links, overfull hboxes indicated)
%nolistspacing, % If the document is onehalfspacing or doublespacing, uncomment this to set spacing in lists to single
%liststotoc, % Uncomment to add the list of figures/tables/etc to the table of contents
%toctotoc, % Uncomment to add the main table of contents to the table of contents
%parskip, % Uncomment to add space between paragraphs
%nohyperref, % Uncomment to not load the hyperref package
headsepline, % Uncomment to get a line under the header
%chapterinoneline, % Uncomment to place the chapter title next to the number on one line
consistentlayout, % Uncomment to change the layout of the declaration, abstract and acknowledgements pages to match the default layout
]{MastersDoctoralThesis} % The class file specifying the document structure


\usepackage[utf8]{inputenc} % Required for inputting international characters
\usepackage[T1]{fontenc} % Output font encoding for international characters

\usepackage{mathpazo} % Use the Palatino font by default

\usepackage[backend=biber,% style=authoryear,
natbib=true]{biblatex}
\addbibresource{refs.bib} % The filename of the bibliography

\usepackage[autostyle=true]{csquotes} % Required to generate language-dependent quotes in the bibliography
\usepackage{pifont}

\usepackage{cover} % Styles for the cover (title) page

% Character sets.

\usepackage[english]{babel}

% Item list style.

\usepackage{enumitem}

% Colors.

\usepackage{xcolor}

% To set font size.

\usepackage{moresize}
\usepackage{anyfontsize}

% Listings and algorithms.

\usepackage[noline]{algorithm2e}
\usepackage{listings,lstautogobble,lstcoq,lstvhdl}

% To strike out text.

\usepackage[normalem]{ulem}

% Maths.

\usepackage{amsfonts,amssymb,amsmath,amsthm}
\usepackage{stackengine}

%%%%%%%%%%%%%%%%%%%%%%%%%%%%%%%%%%%%%%%%%%%%%%%%%%%%%%%%%%
%       PROGRAMMING LANGUAGE SPECIFICATION PACKAGES      %
%%%%%%%%%%%%%%%%%%%%%%%%%%%%%%%%%%%%%%%%%%%%%%%%%%%%%%%%%%

% Semantics rules.

\usepackage{semantic}
\usepackage{ebproof}

% Grammar specification.

\usepackage{syntax}

%%%%%%%%%%%%%%%%%%%%%%%%%%%%%%%%%%%%
%       FIGURE-RELATED PACKAGES    %
%%%%%%%%%%%%%%%%%%%%%%%%%%%%%%%%%%%%

% Float positioning.

\usepackage{float}

% Pictures.

\usepackage{graphicx}

% Tikz libraries.

\usepackage{tikz-timing}
\usepackage{tikz}

%%%%%%%%%%%%%%%%%%%%%%%%%%%%%%%%%%%
%       TABLE-RELATED PACKAGES    %
%%%%%%%%%%%%%%%%%%%%%%%%%%%%%%%%%%%

\usepackage{hhline}
\usepackage{multirow}

%%%%%%%%%%%%%%%%%%%%%%%%%%%%%%%%%%%%%%%%%%
%       BIBLIOGRAPHY-RELATED PACKAGES    %
%%%%%%%%%%%%%%%%%%%%%%%%%%%%%%%%%%%%%%%%%%


%%% Local Variables:
%%% mode: latex
%%% TeX-master: "main"
%%% End:

% SITPN fields
\SetKwData{sitpn}{sitpn}
\SetKwData{ma}{initial\_marking}
\SetKwData{pgroups}{priority\_groups}
\SetKwData{lneigh}{lneighbors}
\SetKwData{conds}{conditions}
\SetKwData{acts}{actions}
\SetKwData{pls}{places}

\SetKwFunction{pre}{pre}
\SetKwFunction{post}{post}
\SetKwFunction{hasact}{has\_action}
\SetKwFunction{sitvals}{s\_intervals}

% SITPN state fields

\SetKwData{sstate}{s}
\SetKwData{istate}{s'}
\SetKwData{fstate}{s''}
\SetKwData{sma}{marking}
\SetKwData{sfired}{fired}
\SetKwData{sreset}{reset}
\SetKwData{sexecf}{exec\_f}
\SetKwData{sexeca}{exec\_a}
\SetKwData{sditvals}{d\_intervals}
\SetKwData{scvals}{cond\_values}

\SetKwFunction{mkstate}{make\_state}

% % Local variables

\SetKwData{redm}{residual\_marking}
\SetKwData{transf}{transitions\_2b\_fired}
\SetKwData{newm}{new\_marking}  
\SetKwData{cvalsp}{cond\_values'}
\SetKwData{execap}{exec\_a'}
\SetKwData{ditvalsp}{d\_intervals'}

% Misc functions

\SetKwFunction{gv}{getValue}
\SetKwFunction{pshb}{pushBack}
\SetKwFunction{issens}{isSensitized}
\SetKwFunction{isfirable}{isFirable}
\SetKwFunction{updateredm}{updateResidualMarking}
\SetKwFunction{umpre}{updateMarkingPre}
\SetKwFunction{umpost}{updateMarkingPost}
\SetKwFunction{getcondvals}{getConditionValues}
\SetKwFunction{getactstates}{getActionStates}
\SetKwFunction{uptimeitvals}{updateTimeIntervals}
\SetKwFunction{mapfire}{mapFire}
\SetKwFunction{env}{env}
\SetKwFunction{len}{length}
\SetKwFunction{mappre}{mapUpdateMarkingPre}
\SetKwFunction{blockedandreset}{getBlockedAndResetOrds}
\SetKwFunction{mappost}{mapUpdateMarkingPost}
\SetKwFunction{getfun}{getFunctionStates}
\SetKwFunction{consTokens}{consumeTokens}
\SetKwFunction{prodTokens}{produceTokens}
\SetKwFunction{e2bfired}{electToBeFired}

\SetKwFunction{mkstate}{make\_state}

% Highlighting sections of code.
\def\HiLi{\leavevmode\rlap{\hbox to \hsize{\color{yellow!50}\leaders\hrule height .8\baselineskip depth .5ex\hfill}}}
%%%%%%%%%%%%%%%%%%%%%%%%%%%%%%%%%%%%%%%%%%%%%%%%%
%%%%%%% TikZ. library imports and macros. %%%%%%%
%%%%%%%%%%%%%%%%%%%%%%%%%%%%%%%%%%%%%%%%%%%%%%%%%

% TikZ libraries.

\usetikzlibrary{backgrounds,decorations.pathmorphing}
\usetikzlibrary{positioning, calc, arrows, shapes, shapes.misc, automata, petri, patterns}
\usetikztiminglibrary[dual arrows]{clockarrows}

% To use tikzmark, to place and refer to marks outside the current figure

\tikzset{every picture/.style={remember picture}}
\tikzset{transition/.append style={fill=black!20, thick}}
\tikzset{timing/font=\fontsize{40}{41}\selectfont}

% TikZ Styles for Petri nets.

\tikzstyle{test}=[pre, *-]
\tikzstyle{inhib}=[pre, o-]
% New math symbols.

% Defining simple hyphen for math mode.

\mathchardef\mhyphen="2D

% Typing relation with flag.

\newcommand{\tflag}[1]{\mathop{:}\limits_{}^{#1}}

% Overriding union operator.
\newcommand{\ocup}{\mathop{\cup}\limits_{}^{\leftarrow}}

% Differentiated intersection domain.
\newcommand{\dcap}{\mathop{\cap}\limits_{}^{\neq}}

% Concat operator
\newcommand\doubleplus{+\kern-1.3ex+\kern0.8ex}
\newcommand\mdoubleplus{\ensuremath{\mathbin{+\mkern-9mu+}}}

\newcommand\steq{
  \mathrel{\stackrel{\makebox[0pt]{\mbox{\normalfont\tiny\sffamily state}}}{=}}}

\newcommand\lassign[1]{
\mathrel{\stackrel{\makebox[0pt]{\mbox{\normalfont\tiny\sffamily #1}}}{:=}}}

\newcommand\rlabelledarrow[1]{
\mathrel{\stackon[2pt]{$\rightarrow$}{\scriptsize $#1$}}}

\newcommand\srsquigarrow[2]{%
  \mathrel{%
    \begin{tikzpicture}[%
      baseline={(current bounding box.south)}
      ]
      \node[%
      ,inner sep=1ex
      ,align=center
      ] (tmp) {#2 $#1$};
      
      \path[%
      ,draw,<-
      ,decorate,decoration={%
        ,zigzag
        ,amplitude=0.7pt
        ,segment length=1.2mm,pre length=3.5pt}
      ] 
      ($(tmp.south east)-(3pt,0)$) -- ($(tmp.south west)+(3pt,0)$);
    \end{tikzpicture}
  }
}

\newcommand\xrsquigarrow[1]{%
  \mathrel{%
    \begin{tikzpicture}[%
      baseline={(current bounding box.south)}
      ]
      \node[%
      ,inner sep=1ex
      ,align=center
      ] (tmp) {\scriptsize $#1$};
      
      \path[%
      ,draw,<-
      ,decorate,decoration={%
        ,zigzag
        ,amplitude=0.7pt
        ,segment length=1.2mm,pre length=3.5pt}
      ] 
      ($(tmp.south east)-(3pt,0)$) -- ($(tmp.south west)+(3pt,0)$);
    \end{tikzpicture}
  }
}

\newcommand\xlsquigarrow[1]{%
  \mathrel{%
    \begin{tikzpicture}[%
      ,baseline={(current bounding box.south)}
      ]
      \node[%
      ,inner sep=.44ex
      ,align=center
      ] (tmp) {\scriptstyle $#1$};
      \path[%
      ,draw,<-
      ,decorate,decoration={%
        ,zigzag
        ,amplitude=0.7pt
        ,segment length=1.2mm,pre length=3.5pt
      }
      ] 
      (tmp.south west) -- (tmp.south east);
    \end{tikzpicture}
  }
}

\newcommand\srarrow[2]{%
  \mathrel{%
    \begin{tikzpicture}[%
      baseline={(current bounding box.south)}
      ]
      \node[%
      ,inner sep=1ex
      ,align=center
      ] (tmp) {#2 $#1$};
      \path[%
      ,draw,<-,] 
      ($(tmp.south east)+(0,2pt)$) -- ($(tmp.south west)+(3pt,2pt)$);
    \end{tikzpicture}
  }
}

% NEW MATH ENVIRONMENTS

\newtheorem{definition}{Definition}
\newtheorem{notation}{Notation}
\newtheorem{thm}{Theorem}
\newtheorem{remark}{Remark}
\newtheorem{hypothesis}{Hypothesis}
\newtheorem{lemma}{Lemma}
\newtheorem{proposition}{Proposition}

\makeatletter
\newcommand*{\inlineequation}[2][]{%
  \begingroup
    % Put \refstepcounter at the beginning, because
    % package `hyperref' sets the anchor here.
    \refstepcounter{equation}%
    \ifx\\#1\\%
    \else
      \label{#1}%
    \fi
    % prevent line breaks inside equation
    \relpenalty=10000 %
    \binoppenalty=10000 %
    \ensuremath{%
      % \displaystyle % larger fractions, ...
      #2%
    }%
    ~\@eqnnum
  \endgroup
}
\makeatother
%%%%%%%%%%%%%%%%%%%%%%%%%%%%%%%%%%%%%%
%%%%%%% Misc. macros and defs. %%%%%%%
%%%%%%%%%%%%%%%%%%%%%%%%%%%%%%%%%%%%%%

% Macro definitions.

\def\bbook{\textsf{B}-Book}
\def\bmth{\textsf{B}}
\def\cakeml{\textsf{CakeML}}
\def\cnrs{\textsf{CNRS}}
\def\coq{\textsf{Coq}}
\def\hilecop{\textsf{HILECOP}}
\def\inria{\textsf{Inria}}
\def\isahol{\textsf{Isabelle/HOL}}
\def\lirmm{\textsf{LIRMM}}
\def\lustre{\textsf{Lustre}}
\def\nasa{\textsf{NASA}}
\def\um{\textsf{Université de Montpellier}}
\def\neurin{\textsf{NEURINNOV}}
\def\uml{\textsf{UML}}
\def\vhdl{\textsf{VHDL}}
\def\hvhdl{$\mathcal{H}$-\textsf{VHDL}}

% Counter definitions.

\newcounter{saveenumi}
\newcommand{\tocontinue}{\setcounter{saveenumi}{\theenumi}}
\newcommand{\resume}{\setcounter{enumi}{\thesaveenumi}}

\newcommand{\backupbegin}{
  \newcounter{finalframe}
  \setcounter{finalframe}{\value{framenumber}}
}
 
\newcommand{\backupend}{
  \setcounter{framenumber}{\value{finalframe}}
}

%%%%%%%%%%%%%%%%%%%%%%%%%%%%%%%%%%%%%%%%
%% BEHAVIOR PRESERVATION PROOF MACROS %%
%%%%%%%%%%%%%%%%%%%%%%%%%%%%%%%%%%%%%%%%

\def\InCsCompT{$\mathtt{comp}(id_t,$ $\texttt{transition},$ $g_t,$ $i_t,$ $o_t)\in{}d.cs$}

\def\exT{\noindent{}By construction and by definition of $id_t$, there
  exist $g_t,i_t,o_t$ s.t. \InCsCompT.}

\def\InCsCompP{$\mathtt{comp}(id_p,$ $\texttt{place},$ $g_p,$ $i_p,$ $o_p)\in{}d.cs$}

\def\exP{\noindent{}By construction and by definition of $id_p$, there
  exist $g_p,i_p,o_p$ s.t. \InCsCompP.}

\def\upSim{$\gamma,E_c,\tau\vdash{}s\stackrel{\uparrow}{\approx}\sigma$}
\def\dwSim{$\gamma\vdash{}s\stackrel{\downarrow}{\approx}\sigma$}

\def\upSitpn{$E_c,\tau\vdash{}s\xrightarrow{\uparrow}s'$}
\def\dwSitpn{$E_c,\tau\vdash{}s\xrightarrow{\downarrow}s'$}

\newtcolorbox{myproofA}[1][Proof]{%
  fonttitle=\bfseries, 
  title={#1.},
  breakable,
  colframe=iris!5,
  colbacktitle=iris!5,
  colback=iris!5,
  coltitle=black,
  lines before break=2
}

\newenvironment{niproof}{
  \begin{myproofA}
    \setlength{\parindent}{0pt}
}
{
 \qed
\end{myproofA}
}

% \newenvironment{niproof}{%
%   \begin{proof}
%     \setlength{\parindent}{0pt}
%   }{%
%   \end{proof}}

%%% Local Variables:
%%% mode: latex
%%% TeX-master: "../main"
%%% End:


%%% Local Variables:
%%% mode: latex
%%% TeX-master: "main"
%%% End:


%----------------------------------------------------------------------------------------
%	GENERAL SETTINGS
%----------------------------------------------------------------------------------------

\setlength{\hfuzz}{20pt}
\renewcommand{\labelitemi}{--}

%----------------------------------------------------------------------------------------
%	MARGIN SETTINGS
%----------------------------------------------------------------------------------------

\geometry{
	paper=a4paper, % Change to letterpaper for US letter
	inner=1.5cm, % Inner margin
	outer=1.5cm, % Outer margin
	bindingoffset=.5cm, % Binding offset
	top=1.5cm, % Top margin
	bottom=1.5cm, % Bottom margin
	%showframe, % Uncomment to show how the type block is set on the page
}

%----------------------------------------------------------------------------------------
%	THESIS INFORMATION
%----------------------------------------------------------------------------------------

\thesistitle{Vérification formelle d'une méthodologie pour la conception et la production de systèmes numériques critiques} % Your thesis title, this is used in the title and abstract, print it elsewhere with \ttitle
\supervisor{Dr. James \textsc{Smith}} % Your supervisor's name, this is used in the title page, print it elsewhere with \supname
\examiner{} % Your examiner's name, this is not currently used anywhere in the template, print it elsewhere with \examname
\degree{Doctor of Philosophy} % Your degree name, this is used in the title page and abstract, print it elsewhere with \degreename
\author{Vincent \textsc{Iampietro}} % Your name, this is used in the title page and abstract, print it elsewhere with \authorname
\addresses{} % Your address, this is not currently used anywhere in the template, print it elsewhere with \addressname

\subject{Biological Sciences} % Your subject area, this is not currently used anywhere in the template, print it elsewhere with \subjectname
\keywords{} % Keywords for your thesis, this is not currently used anywhere in the template, print it elsewhere with \keywordnames
\university{\href{http://www.university.com}{University Name}} % Your university's name and URL, this is used in the title page and abstract, print it elsewhere with \univname
\department{\href{http://department.university.com}{Department or School Name}} % Your department's name and URL, this is used in the title page and abstract, print it elsewhere with \deptname
\group{\href{http://researchgroup.university.com}{Research Group Name}} % Your research group's name and URL, this is used in the title page, print it elsewhere with \groupname
\faculty{\href{http://faculty.university.com}{Faculty Name}} % Your faculty's name and URL, this is used in the title page and abstract, print it elsewhere with \facname

\AtBeginDocument{
\hypersetup{pdftitle=\ttitle} % Set the PDF's title to your title
\hypersetup{pdfauthor=\authorname} % Set the PDF's author to your name
\hypersetup{pdfkeywords=\keywordnames} % Set the PDF's keywords to your keywords
}

\begin{document}

% Macros to put after the beginning of the document

%%%%%%%%% MDFRAMED ENVIRONMENT DEFINITIONS.

\mdfdefinestyle{framestyle}{
  linecolor=Black}

\mdfdefinestyle{ihstyle}{
  linecolor=Green,
  backgroundcolor=green!10,
  roundcorner=5pt}

\mdfdefinestyle{qedstyle}{
  linecolor=Red,
  backgroundcolor=red!10,
  roundcorner=5pt}

\mdfdefinestyle{pcommentstyle}{%
  linewidth=1pt,
  linecolor=black,
  bottomline=false,topline=false,rightline=false,
  innerrightmargin=0pt,innertopmargin=0pt,innerbottommargin=0pt,
  innerleftmargin=1em,% Distance between vertical rule & proof content
  skipabove=.5\baselineskip}

\mdfdefinestyle{todostyle}{
  linecolor=Brown,
  backgroundcolor=BurntOrange!50,
  roundcorner=5pt}

\mdfdefinestyle{premisesstyle}{
  linecolor=Green,
  backgroundcolor=green!10,
  roundcorner=10pt}

\mdfdefinestyle{sidecondsstyle}{
  linecolor=Red,
  backgroundcolor=red!10,
  roundcorner=10pt}

\mdfdefinestyle{tablelegendstyle}{
  linecolor=Blue,
  backgroundcolor=blue!10,
  roundcorner=10pt}

\newmdenv[frametitle=Premises,style=premisesstyle]{premises}
\newmdenv[frametitle={Side conditions}, style=sidecondsstyle]{sideconds}
\newmdenv[frametitle={Table legend}, style=tablelegendstyle]{tablelegend}

\newmdenv[style=ihstyle]{ih}
\newmdenv[style=framestyle,backgroundcolor=none]{frameb}
% \newmdenv[style=qedstyle]{qedbox}
\newmdenv[style=pcommentstyle]{pcomm}
\newmdenv[style=todostyle]{todobox}

%%% Local Variables:
%%% mode: latex
%%% TeX-master: "main"
%%% End:


\frontmatter % Use roman page numbering style (i, ii, iii, iv...) for the pre-content pages

\pagestyle{plain} % Default to the plain heading style until the thesis style is called for the body content

%----------------------------------------------------------------------------------------
%	TITLE PAGE
%----------------------------------------------------------------------------------------

\title{\protect \parbox{\textwidth}{\centering Vérification formelle
    d'une méthodologie pour la conception et la production de systèmes
    numériques critiques}}

\englishtitle{\protect \parbox{\textwidth}{\centering\it Formal
    verification of a methodology for the design and production of
    safety-critical digital systems}}

\directorA{David Delahaye}
\directorB{David Andreu}
% si deux personnes pour direction these, utiliser directorB et enlever %direction-2% dans fichier couverture.sty
%\directorB{[]}

\date{16 décembre 2021}

\thispagestyle{empty}

\universite{UNIVERSITÉ DE MONTPELLIER}
% \univlogo{Figures/logo_UM.png}
% \univwallpaper{Figures/wallpaper.png}
\doctoral{: Information, Structures, Syst\`emes}
\researchunit{LIRMM}
\specialisation{Informatique}
\jury{Sandrine Blazy}{Professeure}{Université de Rennes}{Rapporteure}
\jury{Frédéric Boniol}{Maître de recherche}{ONERA, Toulouse}{Rapporteur}
\jury{David Déharbe}{Docteur, Ingénieur}{Clearsy}{Examinateur}
\jury{Marc Pouzet}{Professeur}{ENS, Paris}{Examinateur}
\jury{David Delahaye}{Professeur}{Université de Montpellier}{Directeur}
\jury{David Andreu}{MCF}{Université de Montpellier/Neurinnov}{Co-directeur}

\maketitle

%----------------------------------------------------------------------------------------
%	ACKNOWLEDGEMENTS
%----------------------------------------------------------------------------------------

% \begin{acknowledgements}
%   \addchaptertocentry{\acknowledgementname} % Add the acknowledgements to the table of contents
%   The acknowledgments and the people to thank go here, don't forget to
%   include your project advisor\ldots
% \end{acknowledgements}

%----------------------------------------------------------------------------------------
%	LIST OF CONTENTS/FIGURES/TABLES PAGES
%----------------------------------------------------------------------------------------

\setcounter{tocdepth}{2}
\tableofcontents % Prints the main table of contents

% \listoffigures % Prints the list of figures

% \listoftables % Prints the list of tables

%----------------------------------------------------------------------------------------
%	ABBREVIATIONS
%----------------------------------------------------------------------------------------

% \begin{abbreviations}{ll} % Include a list of abbreviations (a table of two columns)
%   \textbf{SITPN} & \textbf{S}ynchronously executed \textbf{I}nterpreted \textbf{T}ime \textbf{P}etri \textbf{N}et with priorities \\
%   \textbf{VHDL} & \textbf{V}ery high speed integrated circuit \textbf{H}ardware \textbf{D}escription \textbf{L}anguage \\
%   \textbf{CIS} & \textbf{C}omponent \textbf{I}nstantiation \textbf{S}tatement \\
%   \textbf{PCI} & \textbf{P}lace \textbf{C}omponent \textbf{I}nstance \\
%   \textbf{TCI} & \textbf{T}ransition \textbf{C}omponent \textbf{I}nstance \\
%   \textbf{GPL} & \textbf{G}eneric \textbf{P}rogramming \textbf{L}anguage \\
%   \textbf{HDL} & \textbf{H}ardware \textbf{D}escription \textbf{L}anguage \\
%   \textbf{LRM} & \textbf{L}anguage \textbf{R}eference \textbf{M}anual \\
%   \textbf{DSL} & \textbf{D}omain \textbf{S}pecific \textbf{L}anguage \\
%   \textbf{MDE} & \textbf{M}odel-\textbf{D}riven \textbf{E}ngineering \\
%   \textbf{FM} & \textbf{F}ormal \textbf{M}ethods \\
% \end{abbreviations}

%----------------------------------------------------------------------------------------
%	DEDICATION
%----------------------------------------------------------------------------------------

% \dedicatory{For/Dedicated to/To my\ldots} 

%----------------------------------------------------------------------------------------
%	FRENCH ABSTRACT
%----------------------------------------------------------------------------------------

\begin{frenchabstract}
  La production de circuits numériques complexes est devenue
  impossible sans l'aide des ordinateurs. La méthodologie \hilecop{}
  (HIgh LEvel hardware COmponent Programming) assiste les ingénieurs
  dans la conception et la production de circuits numériques dans le
  contexte des systèmes critiques, i.e. systèmes dont le
  malfonctionnement peut résulter en la perte de vies humaines, des
  catastrophes naturelles, des désastres économiques, etc. A titre
  d'exemple, la société
  Neurinnov\footnote{\url{https://neurinnov.com/}} applique la
  méthodologie \hilecop{} pour la production de neuroprothèses,
  considérées comme des dispositifs médicaux hautement critiques par
  la loi de régulation de
  l'UE\footnote{\url{https://eur-lex.europa.eu/eli/reg/2017/745/2020-04-24}}. Dans
  \hilecop{}, les ingénieurs produisent un modèle de circuit
  numérique. Ils utilisent un formalisme graphique qui regroupe les
  diagrammes à composant et un type particulier de réseaux de Petri
  (RdP). Ensuite, le modèle est transformé en une représentation
  textuelle intermédiaire décrite en langage \vhdl{} (Very high speed
  integrated circuit Hardware Description Language). Finalement, un
  compilateur/synthétiseur industriel génère un circuit numérique
  physique, i.e. un ASIC ou sur carte FPGA, depuis la représentation
  \vhdl{}. Ici, l'utilisation des
  RdPs % comme structure sous-jacente aux
  % modèles de circuits numériques
  est liée au contexte des systèmes critiques. Les RdPs permettent la
  vérification de propriétés sur les modèles de circuits numériques
  grâce à l'application de techniques de \textit{model-checking}.
  Cependant, une des transformations décrite dans la méthodologie
  \hilecop{} pourrait altérer le \textit{comportement} (ou
  \textit{sémantique}) des modèles initiaux, invalidant ainsi les
  précédentes étapes de vérification. Le but de cette thèse est de
  prouver que la transformation modèle-vers-texte de \hilecop{}, qui
  génère une description \vhdl{} depuis un modèle de circuit
  numérique, préserve le comportement des modèles d'entrée, i.e.: pour
  tout modèle passé en entrée de la transformation, la description
  \vhdl{} résultante se comporte de la même manière. Pour prouver
  cette propriété, nous nous inspirons des travaux menés sur la
  vérification formelle de compilateurs (notamment sur le compilateur
  C certifié \ccert{} \cite{Leroy2009}). Notre approche est celle de
  la vérification déductive interactive avec assistants de preuve.
  % Spécifiquement, nous voulons prouver cette
  % propriété dans le contexte de la vérification déductive interactive
  % avec assistants de preuve.
  Dans ce contexte, les étapes pour établir la propriété de
  préservation de comportement de la transformation sont: (1)
  formaliser la sémantique d'execution de la représentation source,
  (2) de la représentation cible, (3) décrire la transformation, et
  (4) prouver un théorème de préservation sémantique. Même en suivant
  ce processus clairement détaillé, les spécificités de la
  transformation modèle-vers-texte de \hilecop{} (comparaît notamment
  aux compilateurs) apportent de nouvelles questions recherches et des
  challenges à chaque étape. Dans cette thèse, nous utilisons
  l'assistant à la preuve \coq{} pour nous accompagner tout au long du
  processus. Finalement, nous avons prouvé que la transformation de
  \hilecop{} préserve le comportement de tous modèles initiaux. La
  mécanisation complète de la preuve avec \coq{} est un travail en
  cours.
\end{frenchabstract}

%----------------------------------------------------------------------------------------
%	ABSTRACT
%----------------------------------------------------------------------------------------

\begin{abstract}
  The complexity of digital hardware circuits makes it difficult to
  produce them without the help of computers.  The \hilecop{} (HIgh
  LEvel hardware COmponent Programming) methodology aims at the
  assistance of engineers in the design and production of such digital
  circuits. The context of production is the one of
  \textit{safety-critical} digital systems, i.e. systems which failure
  could result in direct human losses, natural catastrophes, economic
  disasters, etc. To give an example, the
  Neurinnov\footnote{\url{https://neurinnov.com/}} company leverages
  the \hilecop{} methodology to produce highly critical medical
  devices known as neuroprostheses.  In \hilecop{}, engineers rely on
  a graphical formalism, based on component diagrams and a particular
  kind of Petri nets, to produce a model of a digital circuit. Then, a
  computer program turns the model into an intermediary description
  written in \vhdl{} (Very high speed integrated circuit Hardware
  Description Language). Finally, an industrial compiler/synthesizer
  transforms the \vhdl{} description into a concrete physical circuit
  on an FPGA, or as an ASIC. The use of Petri nets permits the
  engineers to describe a \textit{formal} model of a digital
  circuit. The mathematical foundations of Petri nets enable the use
  of model-checking techniques. Thus, proofs can be brought that the
  produced models verify certain \textit{soundness} properties.
  However, even with a \textit{sound} model of a circuit, one
  transformation step could alter the behavior of the initial
  model. The goal of this thesis is to bring the formal proof that the
  model-to-text transformation from a \hilecop{} high-level model to a
  \vhdl{} description is \textit{semantic preserving} (or
  \textit{behavior preserving}); i.e. for all high-level model given
  as an input to the transformation, the resulting \vhdl{} description
  behaves similarly. To perform this task, we draw our inspiration
  from the works pertaining to the formal verification of compilers
  for programming languages (especially from the certified C compiler
  \ccert{} \cite{Leroy2009}). Specifically, we are interested in
  proving the property of semantic preservation in the context of
  \textit{deductive} verification with proof assistants. In this
  context, the steps to verify that a transformation is semantic
  preserving include: (1) the formalization of the execution semantics
  of the source representation, (2) of the target representation, (3)
  the formal description of the transformation, and (4) the proof of a
  corresponding semantic preservation theorem. In this thesis, these
  steps have been carried within the framework of the \coq{} proof
  assistant.  Even though these steps are clearly set, the
  specificities of the \hilecop{} model-to-text transformation,
  compared to compilers for programming languages, bring some
  interesting research challenges.  Finally, we have brought the
  informal \textit{paper} proof that the \hilecop{} transformation is
  semantic preserving by demonstrating a related behavior preservation
  theorem.  The full mechanization of the proof using the \coq{} proof
  assistant is an ongoing task.
\end{abstract}

%----------------------------------------------------------------------------------------
%	FRENCH EXTENDED ABSTRACT
%----------------------------------------------------------------------------------------

\chapter{Résumé étendu}
\label{chap:ext-abstract}

\section{Introduction}

Pour répondre aux contraintes liées à la conception de circuits
numériques critiques, et à l'augmentation constante de la complexité
des systèmes, le domaine de l'Ingénierie Système à Base de Modèles
(ISBM) a été développé. L'intérêt est de travailler sur des modèles de
haut niveau avec un pouvoir d'expression et des qualités de
compréhension et de lisibilité qui facilitent les interactions entre
les acteurs de la conception du circuit (i.e, les ingénieurs).
Plusieurs formalismes existent~: le langage SysML
\cite{Friedenthal2014}, des variantes du langage C \cite{Yankova2007},
ou encore les réseaux de Petri (RdPs) \cite{Yakovlev2006}, pour citer
les plus répandus.  Une fois la conception terminée, les modèles sont
physiquement synthétisés en suivant un procédé manuel ou
automatique. Il reste alors à prouver que la phase de transformation
préserve le comportement du modèle de conception.
% Dans le cadre de la production de circuits numériques critiques, les
% méthodes formelles servent notamment à des fins d'analyse des
% modèles de conception. L'analyse permet de vérifier certaines
% propriétés comportementales des circuits modélisés.  Après analyse
% et correction, les modèles sont physiquement synthétisés en suivant
% un procédé manuel ou automatique. Dans le cas d'une approche
% automatique de la synthèse, il reste à prouver que la phase de
% transformation préserve le comportement du modèle de conception.
La présente thèse s'intéresse à la vérification d'un processus d'aide
à la modélisation et à la production de circuits numériques critiques
: la méthodologie \hilecop{} (HIgh LEvel hardware COmponent
Programming).  Cette méthodologie est mise en œuvre dans le cadre de
la création de micro-contrôleurs intégrés à des dispositifs médicaux
de type neuroprothèses.  La Figure \ref{fig:hilecop-wf-french} en décrit les
principales étapes.

\begin{figure}[H]
\centering
\includegraphics[keepaspectratio=true,width=\textwidth]{Figures/Hilecop/hilecop-wf}
\caption[Principe de la méthodologie \hilecop{}.]{Principe de la
  méthodologie \hilecop{}; les double flêches horizontales
  représentent des phases de transformation; les simple flêches
  indiquent les autres types d'opérations ayant cours à une étape
  précise, ou entre étapes.}
\label{fig:hilecop-wf-french}
\end{figure}

Le concepteur de systèmes électroniques esquisse premièrement un
modèle graphique de haut niveau de son circuit (\textcircled{1}). Ce
modèle s'appuie sur le formalisme des diagrammes à composants, avec
l'addition des RdPs pour décrire le comportement interne des parties
du circuit. Dans un deuxième temps, les parties du modèle sont
assemblées et la structure des composants est effacée. Le résultat
obtenu est un réseau de Petri global décrivant le système modélisé
(\textcircled{2}).  Des outils d'analyse exploités par la méthodologie
permettent alors de vérifier certaines propriétés du modèle (caractère
borné, vivacité\dots) et présentent un compte rendu au
concepteur. Après plusieurs itérations du cycle analyse-correction, du
code VHDL est généré à partir du modèle d'implémentation
(\textcircled{3}). Dès lors, la dernière étape de la méthodologie, qui
opère la synthèse du circuit électronique depuis le code source VHDL,
est prise en charge par un compilateur/synthétiseur industriel
propriétaire (\textcircled{4}).

L'objectif de la thèse est de prouver que la transformation du modèle
d'implémentation en code VHDL (i.e, de \textcircled{2} vers
\textcircled{3} dans la Figure~\ref{fig:hilecop-wf-french})
n'introduit pas de divergences de
comportement. % Pour ce faire, la sémantique des
% modèles initiaux et des programmes cibles est considérée. 
Dans cette optique, il sera nécessaire de formaliser la sémantique des
modèles de haut niveau (RdP), du langage cible (VHDL), et de décrire
la transformation. Ensuite, la preuve de similarité comportementale
devra être établie. L'intégralité de la démarche sera mécanisée avec
l'assistant à la preuve \coq{} \cite{Coq2021}. Même si cette démarche
a été éprouvée pour la vérification de compilateurs, son application à
la conception de circuits numériques est bien moins
fréquente. L'intérêt scientifique provient de la distance qui existe
entre le modèle d'exécution du formalisme source (SITPN) et celui du
langage cible (VHDL). Cette distance devra être prise en compte lors
de la preuve de préservation de comportement.

%%% Local Variables:
%%% mode: latex
%%% TeX-master: "../../main"
%%% End:

\section{Un formalisme de haut-niveau : les réseaux de Petri}

Du fait de leur statut de modèles formels et des possibilités
d'analyse qui en résultent, les RdPs ont été retenus comme modèles de
haut niveau de la méthodologie \hilecop{}. Le but de la méthodologie
étant la conception et la production de circuits numériques
\emph{critiques}, les modèles se doivent d'être validés par analyse
formelle. Afin d'augmenter l'expressivité des modèles, les RdPs
\hilecop{} combinent plusieurs classes connues de RdPs (présentées
ci-après), mais leur particularité réside dans leur exécution
synchrone. Les RdPs \hilecop{} sont nommés SITPNs pour Synchronously
executed Interpreted Time Petri Nets with priorities.

Les SITPNs sont des RdP interprétés; des actions peuvent être
associées aux places d'un réseau, et des fonctions/conditions peuvent
être associées aux transitions. % Actions, fonctions et conditions
% décrivent l'interaction du système modélisé avec son environnement.
% Une action est exécutée tant qu'au moins une place à laquelle elle
% est associée possède un marquage non nul. Une fonction, liée à une
% ou plusieurs transitions, est exécutée lors du tir d'une de ces
% transitions.
Actions et fonctions définissent des opérations sur une ensemble de
variables, ici, des \textit{signaux} \vhdl{}.  Les conditions
associées aux transitions sont des expressions Booléennes sûr la
valeur des signaux. Dans un RdP interprété, une transition est
franchissable si elle est sensibilisée et que toutes les conditions
qui lui sont associées sont \emph{vraies}. La Figure~\ref{fig:ipn}
donne un exemple de RdP interprété.

\begin{figure}[H]
  \centering
  \includegraphics[keepaspectratio=true, width=.6\textwidth]{Figures/SITPN/interpreted-pn}
  \caption[Un exemple de réseau de Petri interprété.]{Un exemple de
    réseau de Petri interprété; sur le côté gauche, le RdP; sur le
    côté droit, les expressions Booléennes associées aux conditions et
    les opérations associées aux actions et fonctions.}
  \label{fig:ipn}
\end{figure}%

Les RdP utilisés dans \hilecop{} sont temporels; une fenêtre de tir,
i.e un intervalle de temps, peut être associée à une transition. Un
compteur de temps est lancé lorsqu'une transition devient
sensibilisée; celle-ci devient franchissable lorsque son compteur de
temps a atteint l'intervalle de tir. La Figure~\ref{fig:tpn} donne un
exemple de RdP temporel. La valeur courante des compteurs de temps est
représentée entre chevrons en dessous des intervalles temporels
associés. En résumé, une transition d'un SITPN est franchissable si
elle est sensibilisée, si toutes les conditions qui lui sont associées
sont vraies et si son compteur de temps est dans l'intervalle défini.

\begin{figure}[H]
  \centering
  \includegraphics[keepaspectratio=true, width=0.2\textwidth]{Figures/SITPN/time-pn}
  \caption[Un exemple de RdP temporel.]{Un exemple de RdP temporel. La
    valeur des compteurs de temps apparaît en rouge.}
  \label{fig:tpn}
\end{figure}

Contrairement au cas général, les SITPNs ont une politique de tir
(i.e, une sémantique) \emph{synchrone}. Fondamentalement, le tir des
transitions d'un RdP est un phéno\-mène indéterministe (si deux
transitions sont franchissables au même instant, tous les ordres de
tirs sont possibles), et asynchrone (dès qu'une transition est
franchissable, elle peut être tirée sans attente). A contrario,
l'évolution d'un SITPN est rythmée par le front montant et le front
descendant d'un signal d'horloge, comme montré dans la
Figure~\ref{fig:sync-exec}. Sur le front descendant (\textcircled{1}
de la Figure~\ref{fig:sync-exec}), toutes les transitions devant être
tirées sont déterminées, ce après mise à jour des conditions et
intervalles de temps; sur le front montant (\textcircled{2} de la
Figure~\ref{fig:sync-exec}), les précédentes transitions sont tirées,
entraînant la mis à jour du marquage du réseau et l'exécution de
fonctions. La sémantique d'évolution d'un tel réseau est synchrone et
déterministe.

% \begin{figure}[!htbp]
%   \centering
%   \includegraphics[keepaspectratio,width=.75\linewidth]{pictures/sync-exec}
%   \caption{Evolution synchrone d'un SITPN.}
%   \label{fig:sync-exec}
% \end{figure}

\begin{figure}[H]
  \centering
  \includegraphics[keepaspectratio=true, width=.9\textwidth]{Figures/SITPN/sync-exec}
  \caption{Evolution d'un SITPN synchronisée avec un signal
    d'horloge.}
  \label{fig:sync-exec}
\end{figure}

La structure et la sémantique des SITPNs ont été formalisées dans
\cite{Leroux2014, Merzoug2018}. La sémantique est exprimée comme un
système états-transitions où les transitions sont étiquetées par les
évènements d'un signal d'horloge. Il y a deux évènements possibles :
le front montant et le front descendant du signal. L'état d'un SITPN
décrit, entre autres, le marquage courant du SITPN, la valeur des
compteurs de temps et des conditions associés aux transitions, la
liste des transitions à tirer\dots La sémantique des SITPNs fixe les
règles de changement d'état en fonction des évènements d'horloge. Par
exemple, sur le front descendant d'horloge, a lliste des transitions à
tirer au prochain front montant est calculée; une règle stipule qu'une
aucune transition non franchissable au front descendant n’appartient à
l’ensemble des transitions à tirer.

La première contribution de la thèse est l'implantation en \coq{} de
la structure et de la sémantique des SITPNs.  La sémantique a été
implantée comme une relation inductive paramétrée par un SITPN, deux
états (i.e, avant et après transition), et un évènement d'horloge. La
relation présente deux cas de construction, un pour chaque évènement
d'horloge considéré.  Afin de tester notre implantation de la
sémantique des SITPNs, un interprète a été conçu, i.e un programme qui
simule les changements d'état d'un SITPN pour $n$ cycles d'horloge, en
partant de l'état initial du réseau. Cet interprète est prouvé correct
et complet vis à vis de la sémantique des SITPNs pour une évolution
sur un cycle d'horloge. L'intégralité de la formalisation et de la
mécanisation est mise à disposition du
lecteur\footnote{\url{https://github.com/viampietro/sitpns}}. Cependant,
nous avons utilisées une autre version de l'implémentation des SITPNs
en \coq{} pour effectuer la preuve de préservation sémantique. La
dernière version est plus élégante et utilisent les types
dépendants\footnote{\url{https://github.com/viampietro/ver-hilecop/tree/master/sitpn/dp}}.

%%% Local Variables:
%%% mode: latex
%%% TeX-master: "../../main"
%%% End:

\section{Un langage cible : VHDL}

% \begin{wrapfigure}{H}{.5\textwidth}
%   \centering 
%   \includegraphics[keepaspectratio,width=.5\textwidth]{pictures/design}
%   \caption{Exemple de design VHDL, un couple entité-architecture.}
%   \label{fig:design}
% \end{wrapfigure}

Il existe plusieurs techniques permettant la synthèse physique d'un
RdP. Cependant, la technique la plus étudiée est la transformation
vers la langage VHDL. Cette technique a donc été retenue par la
méthodologie \hilecop{}. Le langage VHDL permet les descriptions
structurelle et comportementale de circuits électroniques, à des fins
de simulation ou de synthèse physique. En VHDL, un \emph{design}
décrit un composant électronique en termes d'interface entrée-sortie
(l'\emph{entité}) et de comportement interne (l'\emph{architecture}).
Le comportement d'un design s'ex\-prime de deux manières : via
l'interconnexion d'in\-stances d'autres designs (des
sous-compo\-sants), ou à l'aide de \emph{processus}.  La spécificité
du langage VHDL tient à l'exécution concurrente des processus et des
sous-composants décrivant une architecture de design. Un processus
définit un bloc d'instructions séquentielles; il observe un certain
nombre de signaux qui composent sa liste de sensibilité. Le changement
d'état d'un signal de cette liste entraîne l'exécution du bloc
d'instructions du processus.  Conceptuellement, un signal VHDL
représente une connexion physique sur un circuit électronique. Les
signaux sont les principaux véhicules des valeurs dans les programmes
VHDL. % Il est possible de les
% confondre jusqu'à une certain point avec les variables de la
% programmation impérative.

La sémantique de VHDL est décrite dans une prose informelle dans le
manuel de référence du langage (MRL). De fait, interpréter un
programme VHDL, qui décrit un \emph{design} de circuit, revient à
simuler le design décrit. Dans le MRL, la sémantique de VHDL est donc
définie sous la forme d'une boucle de simulation. La boucle de
simulation spécifie la dynamique d'exécution des blocs concurrents qui
composent une architecture de design, ainsi que la propagation des
valeurs au travers des signaux.

La littérature propose de nombreuses formalisations de la sémantique
de VHDL \cite{Kloos2012}. Certaines formalisations expriment la boucle
de simulation telle qu'exhibée dans le MRL; d'au\-tres choisissent de
s'abstraire de cette boucle, et optent pour une formalisation
alternative basée sur des modèles permettant la gestion de la
concurrence et du temps (automates temporels, réseaux de Petri,
logique d'intervalles temporels\dots).

La méthodologie \hilecop{} opère la génération d'un design VHDL dans
l'optique de sa synthèse physique. Dès lors, nous ne considérons
qu'une partie \emph{synthétisable} du langage que nous définissons et
nommons $\mathcal{H}$-VHDL. De plus, les designs VHDL générés par la
méthodologie \hilecop{} décrivent des circuits synchrones, i.e, dont
l'exécution est rythmée par un signal d'horloge.  La prise en compte
d'une sous-partie synthétisable et du synchronisme nous a permis
d'exprimer la sémantique des programmes $\mathcal{H}$-VHDL en termes
d'une boucle de simulation bien plus simple en comparaison de celle
exprimée dans le MRL.  L'Algorithme~\ref{alg:sim-loop} décrit notre
boucle spécifique de simulation pour un design $\mathcal{H}$-VHDL.

\begin{figure}[H]
  \begin{algorithm}[H]
    \DontPrintSemicolon

    \SetAlFnt{\fontsize{11}{13}\selectfont}
    
    \SetAlCapFnt{\fontsize{11}{13}\selectfont}
    \SetAlCapNameFnt{\fontsize{11}{13}\selectfont}
    % \NoCaptionOfAlgo

    \caption{Simulation($\Delta$, $\sigma_{e}$, $cs$, $E_p$, $T_c$)}
    \label{alg:sim-loop}
    
    \AlFnt % overriding the new font
    
    \tcp{\textcolor{red}{Initialization phase.}}
    $\sigma_e'\leftarrow$ \texttt{RunAllOnce($\Delta$,$\sigma_{e}$,$cs$)}\;\label{line:run-once} 
    $\sigma\leftarrow$ \texttt{Stabilize($\Delta$,$\sigma_e'$,$cs$)}\;\label{line:stab-init}
    \BlankLine
    
    \tcp{\textcolor{red}{Main loop.}}
    % $T_c\leftarrow{}0$\; \label{line:init-tc}
    $\theta\leftarrow[\sigma]$\; \label{line:init-trace}
    
    \BlankLine

    \While{$T_c>0$}{
      $\sigma_i\leftarrow$ \texttt{Inject($\Delta$,$\sigma$,$E_p$,$T_c$)}\;
      $\sigma_\uparrow\leftarrow$ \texttt{RisingEdge($\Delta$,$\sigma_i$,$cs$)}\;
      $\sigma'\leftarrow$ \texttt{Stabilize($\Delta$,$\sigma_\uparrow$,$cs$)}\;
      $\sigma_\downarrow\leftarrow$ \texttt{FallingEdge($\Delta$,$\sigma'$,$cs$)}\;
      $\sigma\leftarrow$ \texttt{Stabilize($\Delta$,$\sigma_\downarrow$,$cs$)}\;

      \BlankLine
      $\theta\leftarrow\theta\mdoubleplus[\sigma',\sigma]$\; \label{line:append-states}
      $T_c\leftarrow{}T_c-1$\;
    }

    \Return{$\theta$}\;
  \end{algorithm}
\end{figure}

L'Algorithme~\ref{alg:sim-loop} est paramétré par un design élaboré
$\Delta$ et un état de design $\sigma_e$. Ces deux paramètres sont le
résultat de l'élaboration du design qui va être simulé. La paramètre
$cs$ correspond au \textit{comportement}, ou, pour être précis, à la
partie comportementale de l'architecture du design. C'est ce
comportement qui sera exécuté au cours de la simulation. La paramètre
$E_p$ est l'environnement de simulation. Il permet l'injection de
nouvelle valeur sur les ports d'entrée du design à chaque nouveau
cycle d'horloge. La paramètre $T_c$ correspond au nombre de cycles de
simulation à effectuer, i.e. le \textit{front} de simulation.

La première partie de l'Algorithme~\ref{alg:sim-loop} correspond à la
phase d'initialisation. Chaque processus et sous-composants
appartenant à $cs$ sont exécutés exactement une fois lors de cette
phase (\texttt{RunAllOnce}$(\Delta,\sigma_e,cs)$). S'en suit une phase
de stabilisation des signaux
(\texttt{Stabilize}$(\Delta,\sigma_e,cs)$) où seules les parties
combinatoires du design sont exécutées. Ensuite, vient l'exécution de
la boucle principale de simulation. La boucle principale exécute $T_c$
fois les phases d'un cycle d'horloge. Dans l'ordre, ces phases sont:
(1) injection de nouvelles valeurs dans les ports d'entrée du design
simulé, (2) exécution des processus séquentiels qui réagissent au
front montant de l'horloge, (3) stabilisation des signaux, (4)
exécution des processus séquentiels qui réagissent au front descendant
de l'horloge, (5) stabilisation des signaux. Pour un cycle d'horloge,
l'état stable obtenu au milieu du cycle et à la fin du cycle sont
ajoutés à la trace de simulation $\theta$. Cette trace de simulation
est retournée à la fin de l'Algorithme~\ref{alg:sim-loop}.

Une formalisation de la sémantique de $\mathcal{H}$-VHDL a été
effectuée sous la forme d'une sémantique opérationnelle à petit pas
pour la partie simulation, i.e. chaque état intermédiaire est
considéré dans la trace de simulation. Le corps des processus est lui
interprété avec une sémantique à grands pas.  La mécanisation en
\coq{} de la syntaxe et la sémantique de \hvhdl{} a été
réalisée. Cette sémantique s'inspire des travaux de formalisation
esquissés dans \cite{VanTassel1995, Borrione1995}. La sémantique
formalisée prend également en compte la phase d'élaboration du design,
préliminaire à la simulation. L'élaboration génère l'environnement de
simulation, i.e un couple $\Delta,\sigma_{init}$ qui se trouve en
paramètre de la boucle de simulation (voir
Algorithme~\ref{alg:sim-loop}). Durant la phase d'élaboration, une
vérification de type est effectué sur le code VHDL. La vérification de
type s'assure que la partie déclarative et la partie comportementale
du design VHDL respectent certaines règles de typage définies par le
MRL. Par exemple, pour une instruction d'affectation de valeur à un
signal, l'expression affectée doit être du même type que le signal
cible.


%%% Local Variables:
%%% mode: latex
%%% TeX-master: "../../main"
%%% End:

\section{Conclusion}
\label{sec:concl}

Le but de cette thèse a été de vérifier formellement une partie de la
méthodologie \hilecop{}, usitée dans le cadre de la conception de
circuits numériques critiques.  Spécifiquement, le travail de
vérification porte sur la phase transformant un modèle de conception,
à base de RdPs, en \emph{design} VHDL. Au final, nous avons fait la
démonstration d’un théorème de préservation de comportement pour cette
phase de transformation. La preuve a été effectuée informellement,
mais représente tout de même un volume d'une centaine de pages. La
mécanisation complète de cette preuve avec l'assistant de preuves
\coq{} est en cours de réalisation.

Comme autres perspectives de travail, les modèles d'entrée de la
transformation doivent inclure de nouveaux éléments.  Deux éléments
déjà existant dans ce formalisme restent à prendre en compte : les
macroplaces, qui permettent d'exprimer la gestion d'exceptions dans
les SITPNs, ainsi que la possibilité de spécifier des domaines
d'horloge différents au sein d'un même modèle; c'est le cas des
systèmes Globalement Asynchrones Localement Synchrones (GALS).

Le code \coq{} de la thèse, comprenant la formalisation des
sémantiques des SITPNs et du langage \hvhdl{}, l'implémentation de la
transformation, l'expression du théorème de préservation sémantique
ainsi qu'une partie de la preuve formelle, est accessible à l'adresse:

\begin{center}
  \url{https://github.com/viampietro/ver-hilecop}
\end{center}

%%% Local Variables:
%%% mode: latex
%%% TeX-master: "main"
%%% End:


%%% Local Variables:
%%% mode: latex
%%% TeX-master: "../main"
%%% End:


%----------------------------------------------------------------------------------------
%	THESIS CONTENT - CHAPTERS
%----------------------------------------------------------------------------------------

\mainmatter % Begin numeric (1,2,3...) page numbering

\pagestyle{thesis} % Return the page headers back to the "thesis" style

% Include the chapters of the thesis as separate files from the Chapters folder
% Uncomment the lines as you write the chapters

% % Chapter Template

\chapter{Introduction}
\label{chap:intro}

With the use of every human-bred technology is associated a
risk. Regarding the nature of the technology, and the broader system
in which it is involved, the consequences of a failure can be
dramatic. Thus arises the notion of the safety of systems;
\cite{Bowen1993} gives the following definition of safety:

\begin{center}
  ``Safety can [\dots] be defined as the freedom from exposure to
  danger, or the exemption from hurt, injury or loss.''
\end{center}

A \textit{safety-critical} system can be understood as a system for
which the safety aspect is the main concern, being that important
consequences, such as direct human losses, natural catastrophes, or
economic disasters, could result of the failure of the system. In this
thesis, conducted in the field of computer science, we are
particularly interested in safety-critical \textit{computer}
systems. The concept of computer system encompasses both the low-level
hardware-related and the more abstract software-related aspects
involved in computer technologies. These days, computers pervade a
considerable number of objects and technologies that pave our
every-day life, including safety-critical systems. Thus, the risk
associated with the use of computers in certain critical applications
is real. Failures of safety-critical computer systems have happened
and continue to happen; the list of critical incidents maintained by
the ACM Committee on Computers and Public Policy and Peter G. Neumann
ever since the mid 80s \cite{Neumann1994} is always
growing\footnote{The \emph{risks digest} website continues to register
  the computer-related incidents that resulted or could result in
  important damages: \url{https://catless.ncl.ac.uk/Risks/} }.

To ensure the safety of computer systems involved in critical domains
such as avionics, railway, power plants or medicine, there exists a
number of standards and norms developed by international
organizations. These standards set of a number of rules and techniques
to be followed for the design, the production and the validation of
safety-critical computer systems. To cite some well-known standards,
the \href{https://www.eurocae.net/}{EUROCAE} and
\href{https://www.rtca.org/}{RTCA} organisms has devised the
ED-12C/DO-178C and ED-80/DO-254 industry standards for the development
cycle of software and hardware computer systems involved in avionics;
the \href{https://www.cencenelec.eu/}{CENELEC} has defined the
EN-50128 standard for the development of software programs for railway
control and protection systems; the \href{https://www.iec.ch}{IEC} is
at the origin of the IEC-60880 standard for the development of
software programs involved in the control of power plants; the USA,
Canada and the EU have defined the
\href{https://www.commoncriteriaportal.org/}{Common Criteria (CC)}
referential for the evaluation of the safety of systems and software
programs. In this thesis, we are interested in verifying a
\textit{computer-related} methodology involved in the production of
safety-critical medical devices. This domain is regulated by the EU
2017/145
standard\footnote{\url{http://data.europa.eu/eli/reg/2017/745/2020-04-24}}
that sets the rules for the development of medical devices, including
how to validate the technologies involved in the production line.

The rules imposed by the standards vary with respect to the
criticality of the considered systems; for instance, in the medical
field, the regulation text 2017/145 of the EU, pertaining to the
marketing of medical devices, sets a different requirement level
whether we are considering the production of dressings (level 0), or
of neuroprotheses (highest-level, level 4). The IEC defines a SIL
(Safety Integrity Level) measure that qualifies the criticality level
of a system. The CC defines a level of Evaluation Assurance Level
(EAL), from 1 to 7, that must be met by the evaluated systems
regarding their functional requirements.

Among the mandatory procedures, prescribed by the standards, which
must be followed to validate a computer system involved in a
safety-critical system, there are tests (unit, functional or
integration tests) or simulation (especially applied to hardware
systems) are to be noted.  However, in the case of the development
safety-critical computer systems, a particular kind of methods, called
\textit{formal methods} (FM), are also applied. In formal methods, a
computer system is considered as a mathematical object
\cite{Bjorner2014}. As pointed out in \cite{Bowen1993}, ``formal
methods address \textit{correctness} issues'', that is whether a
system delivers the required service. The perks of formal methods are
to set a formal mathematical framework around a computer system.  This
framework allows us to reason about the system and prove that the
system meets some required properties. Thus, a ``formal methods''
framework for computer systems requires at least a formal requirement
language, i.e. with a formal semantics, to express the properties that
a given system must verify. The expression of these required
properties is called the \textit{specification} of the system.  We can
cite some specification formalism such as CCS \cite{Milner1980}, CSP
\cite{Hoare1978}, Petri nets \cite{Petri1962} or TLA+
\cite{Lamport1994} to describe reactive systems (i.e. systems that
continuously interact with an environment); hardware description
languages such as \vhdl{} \cite{Lipsett1986} and \textsf{Verilog}
\cite{IEEE1996} can also be considered as formal specification
languages for hardware designs if provided with a formal semantics and
embedded in a formal proof system (see for instance
\cite{Borrione1992}). We can also cite specification languages such as
VDM \cite{Bjorner1987a}, the Z notation \cite{Abrial2013} on which is
based the B language \cite{Abrial1991} (included in the broader
B-method); but also, all the theorem provers and proof assistants that
come with their own specification languages such as \isahol{}
\cite{Nipkow2002}, \coq{} \cite{Coq2021}, \textsf{PVS}
\cite{Crow1995}, etc.  A FM framework must also provide a formal proof
system to reason about the formal specification of the system. Some FM
frameworks come with means to implement the computer system or
simplified version of the system (i.e. a model) in a formal
setting. This latter kind of framework enables to check if the
implementation of a system always complies with its specification,
i.e. the \textit{soundness} of the system, and if all the aspects of
the specification are met by the implementation, i.e. the
\textit{completeness} of the system.

Even though the purpose is always to check the correctness of systems,
there exist multiple kinds of formal verification techniques. These
techniques can be separated in three groups. The first group refers to
the \textit{deductive verification} methods; the techniques aims at
establishing some proofs over a computer system in a formal proof
setting; the deduction process can either be \textit{interactive}
(i.e. conducted by a human) or automated.  The second group refers to
model-checking techniques. The third group refers to the abstract
interpretation of programs. The techniques applied to the formal
verification of hardware computer systems or software computer systems
are quite similar.  Thus, most of the techniques presented here apply
to both hardware and software developments. Note that the following
presentation of formal verification techniques is not exhaustive.

\paragraph{Abstract interpretation}

Abstract interpretation \cite{Cousot1977} aims at performing automatic
static analysis of programs by approximating the possible execution
states of the program. To do so, the concrete domains of the program
variables are related to more abstract domains notably through the use
of a lattice structure. Afterwards, invariant properties that the
program verifies can be automatically checked against the lattice
structure. For instance, abstract interpretation can help to determine
that a given program terminates. To give examples of
\textit{abstract-interpretation-based} tools, we can cite the
TERMINATOR static analyzer \cite{Cook2006} that proves termination and
other properties over C programs, or the Astrée \cite{Blanchet2002}
program analyzer for real-time embedded systems.

Similarly to abstract interpretation, symbolic execution
\cite{Boyer1975} is another method for the static analysis of programs
based on the partial execution of the considered programs. The method
consists in generating all the execution traces, also called
\textit{symbolic} traces, of a given program; the result takes the
form of a symbolic execution tree. In these traces, some of the
program inputs, i.e. the \textit{variables}, will be associated with
\textit{symbolic} expressions denoting the fact that the values of
these inputs are yet unknown. Thus, by reasoning on the definition
domains of these inputs, some properties of the program can be checked
at execution points, that is, at the nodes of the symbolic execution
tree. The property checking process is most of the time performed by a
constraint solver \cite{Amadini2020}, where constraints are expressed
over the \textit{symbolic} variables of the program. Pertaining to the
construction of the symbolic execution tree, it is obtained most of
the time by applying the rules of a structural formal semantics
\cite{Plotkin2004} associated with the language of the considered
program. Typically, a branching in the execution tree is the result of
the evaluation of a conditional statement. Each path of the execution
tree is associated with a satisfiability condition, i.e. a Boolean
formula, that determines if a given execution point is reachable or
not. Symbolic execution methods are often used to generate test suites
with an important coverage of the execution paths of programs. To give
examples of \textit{symbolic-execution-based} tools, we can cite the
DART \cite{Godefroid2005} or CUTE \cite{Sen2006} test generators.

\paragraph{Model checking}
Model-checking techniques \cite{Queille1982, Clarke1986} build a model
reflecting the execution of a computer system by enumerating all the
possible execution states of the system. Then, the \textit{execution}
model can be automatically checked against some properties that the
computer system must verify. The execution model must be a
finite-state model, i.e. the enumeration of the execution states of
the model must not be infinite.  Most of the time, the properties that
the computer system must verify are expressed through formulas of a
modal logic.  Model-checking techniques are broadly used for the
formal verification of reactive systems, especially hardware
systems. In that case, the properties that must be met by the
considered system are expressed within formulas of one of the many
temporal logics (Interval Temporal Logic \cite{Moszkowski1985}, Linear
Temporal Logic \cite{Pnueli1977}, Computation Tree Logic
\cite{Emerson1984}, etc.), which are handy to express time-related
properties. To give examples of well-known model-checkers, we can cite
the BLAST \cite{Henzinger2002}, CADP \cite{Garavel2013} or
\textsc{Uppaal} \cite{Bengtsson1996} model-checkers.

\paragraph{Deductive software verification}
``Deductive software verification aims at formally verifying that all
possible behaviors of a given program satisfy formally defined,
possibly complex properties, where the verification process is based
on some form of logical inference, i.e., `deduction' ''
\cite{Hahnle2019}.  Deductive software verification methods are
divided into two categories: interactive theorem proving methods and
automated theorem proving methods.

In the philosophy of Interactive Theorem Proving (ITP), the programmer
is responsible for the specification and the implementation of a
computer program, but he also expresses theorems and conducts the
corresponding proofs in a formal proof system. Interactive theorem
proving methods are closely tied to proof assistants (cf. \isahol{},
\coq{}, \textsf{PVS}, etc.), which offer the possibility to specify,
implement, perform proofs over a given program in the same framework.
The programmer builds the proof for a given theorem in an interactive
manner, for instance assisted by a \textit{tactic} language in the
case of the \coq{} proof assistant (see
Chapter~\ref{chap:prelim-notions} for an example). Each proof
assistant comes with its own specification language and underlying
proof system. In between the world of interactive theorem proving and
automatic theorem proving, we can also cite the
\textit{contract-based} verification methods based on the Floyd-Hoare
logic \cite{Floyd1993} and Dijkstra's weakest precondition calculus
\cite{Dijkstra1976} such as SPARK \cite{Carre1988}, ESC
\cite{Rustan1998}, the B-method \cite{Abrial1991}, Frama-C
\cite{Kirchner2015}, or the Escher Verifier \cite{Carlton2013}.

The Automated Theorem Proving (ATP) methods aim at automating the
deduction steps involve in the proof search. Multiple automated
theorem provers exist based on different proof search techniques such
as natural deduction (e.g. \isa{} \cite{Nipkow2002}), the tableaux
method (e.g. the FaCT++ reasoner \cite{Tsarkov2006}, Zenon
\cite{Bonichon2007}), or resolution algorithms (e.g. the E theorem
prover \cite{Schulz2002}, Vampire
\cite{Riazanov2001}). \\

In this thesis, we address the problem of the formal verification of a
particular program. This program transforms an input model, which is
an instance of a particular kind of Petri nets (PNs), into a program
written in a Hardware Description Language (HDL). The program, the
context in which it is involved, the specificities of the input model,
and the target HDL, will all thoroughly be presented in this
thesis. Here, we want to zoom in on the nature of the considered
program that we aim to formally verified. The transformation from an
instance of one formalism to another instance of another formalism is
analog to the case of a \textit{compiler} program. The only difference
is that here an input to the transformation is not a program of a
source \textit{language}, but rather a model of an abstract source
formalism, namely a PN model.  Thus, our formal verification task
amounts to the formal \textit{verification} of a compiler program.
The problem of compiler verification has greatly stimulated the use of
formal methods in the field of software verification. Because a
complete computer system is made out of complete chain of hardware,
firmware and software components, the ultimate goal of the
verification of such a system is to be able to prove the safety of all
the layer composing it. In this system of layers, the place of
compiler programs are mandatory as they are placed at the layer
interfaces. Indeed, one can prove that a given program and a given
hardware is safe, but what if the compilation phase from the given
program to low-level version introduces errors and behavior
divergences. With these considerations in mind, compiler verification
is an important aspect for one that needs to certify an entire
computer system.

Thus, certifying a compiler program amounts to proving that the
compiler verifies certain properties; \cite{Patrignani2019} presents
three of them. First, one can verify that a compiler is
\textit{type-preserving}, also called the \textit{subject reduction}
property. A type-preserving compiler yields a well-typed target
program given a well-typed source program. Second, one can verify that
a compiler is \textit{semantic-preserving}. A semantic-preserving
compiler yields a target program that behaves similarly to the source
program. Thirdly, one can verify that the compiler is
\textit{equivalence-preserving}. Given two source programs that verify
a certain \textit{source-level} equivalence relation, an
equivalence-preserving compiler yields two corresponding target
programs that verify a certain \textit{target-level} equivalence
relation; this target-level equivalence relation is of course somehow
related to the source-level equivalence relation. In this thesis, we
are interested in proving that a \textit{compiler-like} transformation
program is \textit{semantic-preserving}.

\cite{Leroy2009} lists several techniques that exist to establish that
a compiler is semantic-preserving.

The first method is simply called \textit{compiler
  verification}. Compiler verification aims at establishing the
semantic-preserving property of a compiler program by proving a
so-called semantic preservation theorem of the form:

\begin{center}
  For all source program $S$, and compiler $C$ from the language of
  $S$ to a target language, $S$ has the behavior $B$ (written
  $S\Downarrow{}B$) iff $C(S)$
  (i.e. the compiled version of $S$) has the behavior $B$:\\

  $\forall{}S,C,B$,
  $S\Downarrow{}B\Leftrightarrow{}C(S)\Downarrow{}B$.
  
\end{center}

Now the above form of the theorem is the strongest one, i.e. it can be
proved only for a very particular kind of source and target
languages. Other refined versions of this semantic preservation
theorem exist depending on the nature of the source and target
languages. Proving such a theorem is often performed with the help of
a proof assistant as can be witnessed in the pioneering work on the
\ccert{} compiler \cite{Leroy2006}; thus, compiler verification falls
under the hood of deductive verification methods. 

The second method is called compiler \textit{validation}. Compiler
validation does not aim at proving a theorem stating that a given
compiler program is semantic preserving for all input programs. The
strategy of compiler validation is to equip the compiler program with
a \textit{validator} program. Each time the compiler produces a target
program from a source program, the validator tries afterwards to prove
that the two programs have the same behavior. To establish such a
proof, the validator program often relies on model-checking or
abstract interpretation techniques.

The third method is called \textit{proof-carrying} compilation. In
this setting, the compiler program generates alongside the target
program a proof that this program conforms to some property. The
generated proof must be is such a format that can be verified by a
\textit{proof-checker}, built in a proof assistant for instance.

Even though it is not considered as a formal method, compiler testing
is also a way to validate the semantic preservation property for a
given compiler program. Considering the essential part that plays
compiler programs regarding the production of softwares, a lot of
efforts has been dedicated for the generation of test suites
\cite{Chen2020}.

In the thesis, we will follow the \textit{compiler verification}
technique. Consequently, our aim is to prove a semantic preservation
theorem over a transformation program and mechanize the process within
the framework of the \coq{} proof assistant.

This thesis memoir is structured as follows.

Chapter~\ref{chap:prelim-notions} introduces all the necessary
mathematical notions to understand the remainder of the memoir.
Chapter~\ref{chap:hilecop} presents the context of our work, namely
the \hilecop{} methodology, which describes a process to design and
produce safety-critical digital systems. This chapter zooms in on the
part of the methodology that we propose to formally verify; it also
exposes our research question(s) with respect to this task.
Chapter~\ref{chap:hilecop-models} presents in an informal and formal
way a specific kind of Petri net models; these models are the input to
the \hilecop{} transformation program.  Chapter~\ref{chap:hvhdl} gives
an informal presentation of the \vhdl{} language. The \vhdl{} language
is the target language in which the programs generated by the
\hilecop{} transformation are written. We also give in this chapter a
formal definition of the syntax and semantics of a subset of the
\vhdl{} language that we call \hvhdl{}.
Chapter~\ref{chap:transformation} presents the algorithm of the
\hilecop{} transformation and its implementation with the \coq{} proof
assistant.  Chapter~\ref{chap:proof} details the semantic preservation
theorem expressing that the \hilecop{} transformation is
semantic-preserving.  It also gives the high-level theorems and lemmas
involved in the proof of the semantic preservation theorem.  Finally,
Chapter~\ref{chap:concl} ends the memoir, and outlines the
perspectives regarding the full completion of the task of proving that
the \hilecop{} transformation is semantic-preserving.

The results of a literature review pertaining to the formal semantics
\vhdl{} is presented at the beginning of
Chapter~\ref{chap:hvhdl}. Similarly, the results of a literature
review pertaining to the task of compiler verification in the world of
deductive verification methods are presented at the beginning of
Chapters~\ref{chap:transformation} and \ref{chap:proof}.

% All the programming tasks of this thesis have been performed within
% the framework of the \coq{} proof assistant. The produced code, of
% which some fragments are presented all along this thesis memoir, is
% fully accessible under the following \textsf{Git} repository:

% \begin{center}
%   \url{https://github.com/viampietro/ver-hilecop}
% \end{center}



%%% Local Variables:
%%% mode: latex
%%% TeX-master: "../main"
%%% End:

% \chapter{Preliminary Notions}
\label{chap:prelim-notions}

\begin{itemize}
\item Present the \coq{} proof assistant and the concepts underlying the system
\item Present the inductive type system
\item Present what are dependent types
\item Present the proof system (including tactics)
\end{itemize}

%%% Local Variables:
%%% mode: latex
%%% TeX-master: "../main"
%%% End:

% \chapter{Implementation of the \hilecop{} high-level models}
\label{chap:hilecop-models}

\begin{todobox}
  \begin{itemize}
  \item Keep the $sitpn$ argument in the SITPN full execution
    relation, but remove it from the SITPN execution, cycle and state
    transition relations.
    
  \item Change squig arrow to straight arrow in the SITPN state
    transition relation.
  \end{itemize}
\end{todobox}

\section{SITPN Definitions}
\label{sec:sitpn-defs}

\subsection{SITPN Structure and Well-definition}
\label{sec:struct-and-wd}

\subsubsection{SITPN Structure}
\label{sec:sitpn-struct}

\begin{definition}[SITPN]
\label{def:sitpn}
A synchronously executed, extended, generalized, interpreted, and time Petri net
with priorities is a tuple\\
${<}P,T,pre,post,M_0,{\succ},\mathcal{A},\mathcal{C},\mathcal{F},
\mathbb{A},\mathbb{C},\mathbb{F},{I_s}{>}$, where we have:
%
\begin{enumerate}
\item $P=\{p_0,\ldots,p_n\}$, a set of places.
\item $T=\{p_0,\ldots,p_m\}$, a set of transitions.
\item
  $pre\in{}P\rightarrow{}T\nrightarrow(\mathbb{N}^{*}\times\{\mathtt{basic},\mathtt{inhib},\mathtt{test}\})$,
  the function associating a weight to place-transition edges.
\item $post\in{}T\rightarrow{}P\nrightarrow\mathbb{N}^{*}$, the function
  associating a weight to transition-place edges.
\item $M_0\in{}P\rightarrow\mathbb{N}$, the initial marking of the SITPN.
\item $\succ\subseteq{}(T\times{}T)$, the priority relation which is a
  strict order over the set of transitions.
\item $\mathcal{A}=\{a_0,\ldots,a_n\}$, a set of continuous actions.
\item $\mathcal{C}=\{c_0,\ldots,c_n\}$, a set of conditions.
\item $\mathcal{F}=\{f_0,\ldots,f_n\}$, a set of functions (discrete actions).
\item $\mathbb{A} \in P \rightarrow \mathcal{A} \rightarrow \mathbb{B}$, the
function associating actions to places. $\forall{}p\in{}P$,
$\forall{}a\in\mathcal{A}$, $\mathbb{A}(p,a)=1$, if $a$ is associated to $p$,
$\mathbb{A}(p,a)=0$ otherwise.
\item $\mathbb{F} \in T \rightarrow \mathcal{F} \rightarrow \mathbb{B}$, the
function associating functions to transitions. $\mathbb{F}$ is defined in the
same manner as $\mathbb{A}$.
\item $\mathbb{C} \in T \rightarrow \mathcal{C} \rightarrow\{-1,0,1\}$, the
function associating conditions to transitions.
$\forall t \in T$, $\forall c \in \mathcal{C}$,
$\mathbb{C}(t,c)=1$, if $c$ is associated to $t$,
$\mathbb{C}(t,c)=-1$, if $\bar{c}$ is associated to $t$,
$\mathbb{C}(t,c)=0$ otherwise.
\item $I_s\in{}T\nrightarrow\mathbb{I}^{+}$, the partial function
  associating static time intervals to transitions, where
  $\mathbb{I}^{+}\subseteq(\mathbb{N}^{*}\times(\mathbb{N^{*}}\sqcup\{\infty\}))$.
  $T_i$ denotes the definition domain of $I_s$, i.e. the set of time
  transitions.
\end{enumerate}
\end{definition}

\paragraph{Conflict Definition} In the definition of an SITPN, the
priority relation is a mean to solve a situation of conflict in a pair
of transitions. We will keep the definition of a conflict as simple as
possible. Informally, the transitions of a pair are in conflict if
they have an common input place, and if both are linked to this input
place by a \texttt{basic} arc. Figure~\ref{fig:basic-conflict} depicts
a situation of conflict between two transitions.

At some point of the execution of the SITPN, the marking possibly
enables the two transitions of a conflicting pair in such a manner
that the firing of one transition disables the other; then, the
conflict is said to be \emph{effective}. The behavior of PNs is
fundamentally asynchronous, and a token can only be consumed by one
transition. However, in a synchronous setting as the one of the SITPN,
all transitions are first elected to be fired, and then all fired at
the same time.  Therefore, the situation can arise where a same token
is consumed by two transitions, on behalf of them being transitions in
effective conflict that are both elected to be fired (e.g,
Figure~\ref{fig:basic-conflict}). To prevent the phenomenon of
``double spending'', the well-definition property of an $SITPN$
enforces the resolution of all conflicts, i.e, to be able to decide
which transition in a conflicting pair will be fired when the conflict
becomes effective.

\begin{figure}[H]
  \centering
  \includegraphics[keepaspectratio,width=.3\linewidth]{SITPN/struct-conflict-with-basic}
  \caption{Example of conflict between two transitions}
  \label{fig:basic-conflict}
\end{figure}

The formal definition of a conflict is as follows:

\begin{definition}[Conflict]
  For a given $sitpn\in{}SITPN$, two transitions $t,t'\in{}T$ are in
  conflict if and only if there exists a place $p\in{}P$ such that
  $p\in{}input(t)\cap{}input(t')$ and there exist
  $n,m\in\mathbb{N}^{*}$ such that $pre(p,t)=(n,\mathtt{basic})$ and
  $pre(p,t')=(m,\mathtt{basic})$.
\end{definition}

A confict group qualifies a finite set of transitions that are all in
conflict with each other through the same place. In
Figure~\ref{fig:basic-conflict}, the set $\{t_0,t_1\}$ is a conflict
group.  The formal definition of a conflict group is as follows:

\begin{definition}[Conflict Group]
  For a given $sitpn\in{}SITPN$, $T_c\subseteq{}T$ is a conflict group
  if and only if there exists a place $p$ such that
  $\forall{}t\in{}output(p),\big(\exists{}n\in\mathbb{N}^{*},~pre(p,t)=(n,\mathtt{basic})\big)\Leftrightarrow{}t\in{}T_c$.
\end{definition}

Contrary to the statement made in
\cite[p.67]{Leroux2014}, we no more consider the notion
of conflict as being transitive. To illustrate this,
Figure~\ref{fig:conflict-not-trans} shows two conflict groups:
$\{t_0,t_1\}$ and $\{t_1,t_2\}$. In a well-defined $SITPN$ (see
Section~\ref{sec:sitpn-wd}), all conflicts in a conflict group must be
dealt with, i.e, for all pair of transitions in the group the conflict
must be solved. However, we no more consider transitions $t_0$ and
$t_2$ as in conflict. We argue that even when no conflict resolution
technique is applied between transitions in the same situation as
$t_0$ and $t_2$, the execution of the $SITPN$ can neither result in
the double-spending of a token, nor in the case where a transition is
not elected to be fired even though it ought to be. Therefore, we no
more consider the construction of merged conflict group (i.e, conflict
groups must be merged into one if their intersection is not empty;
e.g, $\{t_0,t_1,t_2\}$ in Figure~\ref{fig:conflict-not-trans}) as
being necessary.

\begin{figure}[H]
  \centering
  \includegraphics[keepaspectratio,width=.4\linewidth]{SITPN/conflict-not-trans}
  \caption{Example of two separate conflict groups}
  \label{fig:conflict-not-trans}
\end{figure}

\subsubsection{Well-definition of an SITPN}
\label{sec:sitpn-wd}

A given $sitpn\in{}SITPN$ is well-defined if it enforces some
properties needed on the \hilecop{} source models before the
transformation into \vhdl{}. If the properties, layed out in
Def.~\ref{def:wd-sitpn}, are not ensured, they will lead to
compile-time errors during the transformation into \vhdl{}.

\begin{definition}[Well-defined SITPN]\label{def:wd-sitpn}
  A given $sitpn\in{}SITPN$ is well-defined if and only if:
  \begin{itemize}
  \item $T\neq\emptyset$, the set of transitions must not be empty.
  \item $P\neq\emptyset$, the set of places must not be empty.
  \item There is no isolated place, i.e, a place that has neither
    input nor output transitions:\\
    $\nexists{}p\in{}P,~input(p)=\emptyset\wedge{}output(p)=\emptyset$,
    where $input(p)$ (resp. $output(p)$) denotes the set of input
    (resp. output) transitions of $p$.
  \item There is no isolated transition, i.e, a transition that has
    neither
    input nor output places:\\
    $\nexists{}t\in{}T,~input(t)=\emptyset\wedge{}output(t)=\emptyset$,
    where $input(t)$ (resp. $output(t)$) denotes the set of input
    (resp. output) places of $t$.
  \item All conflicts must be solved by a mean of mutual exclusion:
    priority relation, mutually exclusive conditions, mutually
    exclusive time intervals, structural mutual exclusion.
  \end{itemize}
\end{definition}

\subsection{SITPN State}
\label{sec:sitpn-state}

\begin{definition}[SITPN State]
  \label{def:sitpnstate}
  For a given $sitpn\in{}SITPN$, let $S(sitpn)$ be the set of possible
  states of $sitpn$. An SITPN state $s\in{}S(sitpn)$ is a tuple
  ${<}M,I,reset_t,ex,cond{>}$, where:
  \begin{enumerate}
  \item $M\in{}P\rightarrow\mathbb{N}$ is the current marking of sitpn.
  \item $I\in{}T_i{}\rightarrow\mathbb{N}\sqcup\{\psi\}$ is the
    function mapping time transitions to their current time counter
    value or to the value \emph{locked}.
  \item $reset_t\in{}T_i\rightarrow\mathbb{B}$ is the function mapping time
    transitions to time interval reset orders (defined as Booleans).
  \item $ex\in{}\mathcal{A}\sqcup\mathcal{F}\rightarrow\mathbb{B}$ is the function
    representing the current activation (resp. execution) state of actions
    (resp. functions).
  \item $cond\in\mathcal{C}\rightarrow\mathbb{B}$ is the function representing the
    current value of conditions (defined as Booleans).
  \end{enumerate}
\end{definition}

\subsection{Fired Transitions}
\label{sec:fired-trans}

\begin{definition}[Sensitization]
  A transition $t\in{}T$ is said to be sensitized by a marking $M$,
  which is noted $t\in{}Sens(M)$, if and only if
  $\forall{}p\in{}P,\omega\in\mathbb{N}^{*},~\big(pre(p,t)=(\omega,\mathtt{basic})\vee{}pre(p,t)=(\omega,\mathtt{test})\big)\Rightarrow{}M(p)\ge{}\omega$,
  and $pre(p,t)=(\omega,\mathtt{inhib})\Rightarrow{}M(p)<{}\omega$.
\end{definition}

\begin{definition}[Sensitization by test and basic arcs]
  A transition $t\in{}T$ is said to be sensitized by its basic and
  test arcs at a marking $M$, which is noted $t\in{}Sens_{bt}(M)$, if
  and only if
  $\forall{}p\in{}P,\omega\in\mathbb{N}^{*},~\big(pre(p,t)=(\omega,\mathtt{basic})\vee{}pre(p,t)=(\omega,\mathtt{test})\big)\Rightarrow{}M(p)\ge{}\omega$.
\end{definition}

\begin{definition}[Firability]
A transition $t\in{}T$ is said to be firable at a state\break{}
$s={<}M,I,reset_t,ex,cond{>}$, which is noted $t\in{}Firable(s)$,
if and only if $t\in{}Sens(M)$, and $t\notin{}T_i$ or $I(t)\in{}I_s(t)$, and
$\forall c \in \mathcal{C}, \mathbb{C}(t, c) = 1 \Rightarrow cond(c) = 1$ and
$\mathbb{C}(t, c) = -1 \Rightarrow cond(c) = 0$.
\end{definition}

\begin{definition}[Fired]
  \label{def:fired}
  A transition $t\in{}T$ is said to be fired at the SITPN state
  $s={<}M,I,reset_t,ex,cond{>}$, which is noted $t\in{}Fired(s)$, if
  and only if $t\in{}Firable(s)$ and
  $t\in{}Sens\big(M-\sum_{t_i\in{}Pr(t)}pre(t_i)\big)$, where
  $Pr(t)=\{t_i~|~t_i\succ{}t\wedge{}t_i\in{}Fired(s)\}$.
\end{definition}

Note that $M'=M-\sum_{t_i\in{}Pr(t)}pre(t_i)$ iff
$\forall{}p\in{}P,~M'(p)=M(p)-\sum_{t_i\in{}Pr(t)}pre(p,t_i)$.

Def.~\ref{def:cons-fired} makes another definition of the set of fired
transitions at a given SITPN state. The approach of
Def.~\ref{def:cons-fired} is more constructive, therefore easing the
way to inductive reasoning.

\begin{definition}[Constructive Fired]
  \label{def:cons-fired}
  A transition $t\in{}T$ is said to be fired at the SITPN state
  $s={<}M,I,reset_t,ex,cond{>}$, iff there exists a subset
  $fset\subseteq{}T$ such that $IsFiredSet(s,fset)$ and $t\in{}fset$.
\end{definition}

\begin{lemma}[Fired Equivalence]
  For all
  $sitpn\in{}SITPN,~s\in{}S(sitpn),~t\in{}T,~t\in{}Fired(s)\Leftrightarrow{}\exists{}fset\subseteq{}T,~s.t.,~IsFiredSet(s,fset)\land{}t\in{}fset$.
\end{lemma}

\begin{proof}
  Yet to be proved.
\end{proof}

\begin{definition}[IsFiredSet]
  Given an $sitpn\in{}SITPN$, a SITPN state $s\in{}S(sitpn)$, and a
  subset $fset\in{}T$, the $IsFiredSet$ relation is defined as follows:\\
  $IsFiredSet(s,fset)\equiv{}IsFiredSetAux(s,\emptyset,T,fset)$
\end{definition}

\begin{definition}[IsFiredSetAux]
  The $IsFiredSetAux$ relation is defined by the following rules:\\
  \begin{tabular}{@{}l}
    {\fontsize{9}{11}\selectfont\textsc{IsFiredSetAuxNil}} \\
    
    {\begin{prooftree}[template={\fontsize{9}{11}\selectfont\inserttext}]        
        \infer0{$FiredAux(s, fired, \emptyset, fired)$}
      \end{prooftree}} 
  \end{tabular}
  \begin{tabular}{@{}l}
    {\fontsize{9}{11}\selectfont\textsc{IsFiredSetAuxCons}} \\
    
    {\begin{prooftree}[template={\fontsize{9}{11}\selectfont\inserttext}]

        \hypo{$IsTopPriorityList(T_i, \emptyset, \emptyset, tp)$}
        \infer[no rule]1{$ElectFired(s, fired, tp, fired')$}
        \infer[no rule]1{$FiredAux(s, fired', T_i\setminus{}tp, fset)$}
        \infer1{$FiredAux(s, fired, T_i, fset)$}
      \end{prooftree}} 
  \end{tabular}
\end{definition}

\begin{definition}[IsTopPriorityList]
  The $IsTopPriorityList$ relation is defined by the following rules:\\
  \begin{tabular}{@{}l}
    {\fontsize{9}{11}\selectfont\textsc{IsTPListNil}} \\
    
    {\begin{prooftree}[template={\fontsize{9}{11}\selectfont\inserttext}]
        \infer0{$IsTopPriorityList(\emptyset, tp, ntp, tp)$}
      \end{prooftree}} 
  \end{tabular}
  \begin{tabular}{@{}l}
    {\fontsize{9}{11}\selectfont\textsc{IsTPListTP}} \\
    
    {\begin{prooftree}[template={\fontsize{9}{11}\selectfont\inserttext}]
        \hypo{$\forall{}t'\in{}T_i\cup{}tp\cup{}ntp,t'\nsucc{}t$}
        \infer[no rule]1{$IsTopPriorityList(T_i, \{t\}\cup{}tp, ntp, tp')$}
        \infer1{$IsTopPriorityList(\{t\}\cup{}T_i, tp, ntp, tp')$}
      \end{prooftree}} 
  \end{tabular}

  \begin{tabular}{@{}l}
    {\fontsize{9}{11}\selectfont\textsc{IsTPListNTP}} \\
    
    {\begin{prooftree}[template={\fontsize{9}{11}\selectfont\inserttext}]
        \hypo{$\exists{}t'\in{}T_i\cup{}tp\cup{}ntp\wedge{}t'\succ{}t$}
        \infer[no rule]1{$IsTopPriorityList(T_i, tp, \{t\}\cup{}ntp, tp')$}
        \infer1{$IsTopPriorityList(\{t\}\cup{}T_i, tp, ntp, tp')$}
      \end{prooftree}} 
  \end{tabular}
\end{definition}

\begin{definition}[ElectFired]
  The $ElectFired$ relation is defined by the following rules:
  
  \begin{tabular}{@{}l}
    {\fontsize{8}{11}\selectfont\textsc{ElectFiredNil}} \\
    
    {\begin{prooftree}[template={\inserttext}]
        \infer0{$ElectFired(s,fired,\emptyset, fired)$}
      \end{prooftree}} 
  \end{tabular}
  \vspace{2ex}
  
  \begin{tabular}{@{}l}
    {\fontsize{8}{11}\selectfont\textsc{ElectFired$\bot$}} \\
    
    {\begin{prooftree}[template={\inserttext},right label template={\fontsize{7}{8}\selectfont\inserttext}]
        \hypo{$\neg\big(t\in{}Firable(s)\wedge{}t\in{}Sens(s.M-\sum_{t_i\in{}Pr(t,fired)}pre(t_i))\big)$}
        \infer[no rule]1{$ElectFired(s, fired, tp, fired')$}
        \infer1[$Pr(t,fired)=\{t_i~|~t_i\succ{}t\wedge{}t_i\in{}fired\}$]{$ElectFired(s, fired, \{t\}\cup{}tp, fired')$}
      \end{prooftree}} 
  \end{tabular}
  \vspace{2ex}
  
  \begin{tabular}{@{}l}
    {\fontsize{8}{11}\selectfont\textsc{ElectFired$\top$}} \\
    
    {\begin{prooftree}[template=\inserttext,right label template={\fontsize{7}{9}\selectfont\inserttext}]
        \hypo{$t\in{}Firable(s)$}
        \hypo{$t\in{}Sens(s.M-\sum_{t_i\in{}Pr(t,fired)}pre(t_i))$}
        \infer[no rule]2{$ElectFired(s, \{t\}\cup{}fired, tp, fired')$}
        \infer1[$Pr(t,fired)=\{t_i~|~t_i\succ{}t\wedge{}t_i\in{}fired\}$]{$ElectFired(s, fired, \{t\}\cup{}tp, fired')$}
      \end{prooftree}}
  \end{tabular}
\end{definition}
    
\subsection{SITPN Semantics}
\label{sec:sitpn-sem}

\begin{todobox}
  \begin{itemize}
  \item Change $E$ into $E_c$.
  \item Change typing of the state transition relation to include the
    environment.
  \end{itemize}
\end{todobox}

\begin{definition}[SITPN Semantics]
  \label{def:semantics}
  The semantics of an SITPN is the transition system
  $<S,L,E,\rightsquigarrow>$ where:
  \begin{itemize}[label=-]
  \item $S$ is the set of states of the SITPN.
  \item
    $s_0=<M_0,O_\mathbb{N},O_\mathbb{B},O_\mathbb{B},O_\mathbb{B}>$
    is the initial state of the SITPN, where $M_0$ is the initial
    marking of the SITPN, $O_\mathbb{N}$ is a function that always
    returns 0, $O_\mathbb{B}$ is a function that always returns false.
  \item $L\subseteq{}Clk\times{}\mathbb{N}$ is the set of transition
    labels, where $Clk\in\{\uparrow,\downarrow\}$. A label is a couple $(clk,\tau)$
    composed of a clock event $clk\in{}Clk$, and a time value $\tau\in\mathbb{N}$
    expressing the current count of clock cycles.

  \item

    $E\in{}\mathbb{N}\rightarrow\mathcal{C}\rightarrow\mathbb{B}$ is the
    environment function, which gives (Boolean) values to conditions ($\mathcal{C}$)
    depending on the count of clock cycles ($\mathbb{N}$).
  \item $\rightsquigarrow\subseteq{}S\times{}L\times{}S$ is the state
    transition relation, which is noted
    $E,\tau\vdash{}s\stackrel{clk}{\rightsquigarrow}s'$ where
    $s,s'\in{}S$ and $(clk,\tau)\in{}L$, and which is defined as
    follows:
    % $s\xrsquigarrow{l}s'$
    \begin{itemize}
    \item $\forall\tau\in\mathbb{N}$,
      $E,\tau\vdash{}s\stackrel{\downarrow}{\rightsquigarrow}s'$,
      where $s=<M,I,reset_t,ex,cond>$ and\\
      $s'=<M,I',reset_t,ex',cond'>$, if:
      % $s\xrsquigarrow{(\downarrow,\tau)}s'$,
      \begin{enumerate}[label=(\arabic*)]
      \item\label{item:cond-env} $cond'$ is the function giving the (Boolean)
        values of conditions that are extracted from the environment at the clock count
        $\tau$, i.e.:\\
        $\forall{}c\in{}\mathcal{C},~cond'(c)=E(\tau,c)$.
      \item\label{item:activate-actions} All the actions associated
        with at least one
        marked place in the marking $M$ are activated, i.e.:\\
        $\forall{}a\in{}\mathcal{A},~ex'(a)=\sum\limits_{p\in{}marked(M)}\mathbb{A}(p,a)$
        where $p\in{}marked(M)\equiv{}M(p)>0$.
      % \item\label{item:activate-actions} All the actions associated with at least one
      %   marked place in the marking $M$ are activated, i.e.:\\
      %   $\forall{}a\in{}\mathcal{A},~\big(\exists{}p\in{}P,~M(p)>0
      %   \land{}\mathbb{A}(p,a)=1\big)\Rightarrow{}ex'(a)=1$.
      % \item\label{item:deactivate-actions} All the actions that are only associated
      %   with unmarked places in the marking $M$ are deactivated, i.e.:\\
      %   $\forall{}a\in{}\mathcal{A},~\big(\forall{}p\in{}P,~M(p)=0
      %   \vee{}\mathbb{A}(p,a)=0\big)\Rightarrow{}ex'(a)=0$.
      \item\label{item:reset-counters} All the time transitions that
        are sensitized by the marking $M$ and received the order to
        reset their time intervals, have their time counter reset and
        incremented, i.e.:\\
        $\forall{}t\in{}T_i,~t\in{}Sens(M)\land{}reset_t(t)=1
        \Rightarrow{}I'(t)=1$.
      \item\label{item:inc-counters} All the time transitions with
        active time counters that are sensitized by the marking $M$
        and
        did not receive a reset order, have their time counters incremented, i.e.:\\
        $\forall{}t\in{}T_i,~t\in{}Sens(M)\land{}reset_t(t)=0
        \land{}(I(t)\le{}upper(I_s(t))\lor{}upper(I_s(t))=\infty)\Rightarrow{}$
        $I'(t)=I(t)+1$.
      \item\label{item:locked-counters} All the time transitions
        verifying the same
        conditions as above, but with locked counters, keep having locked counters, i.e.:\\
        $\forall{}t\in{}T_i,~t\in{}Sens(M)\land{}reset_t(t)=0
        \land{}I(t)>{}upper(I_s(t))\land{}upper(I_s(t))\neq\infty\Rightarrow{}$
        $I'(t)=I(t)$.
      \item\label{item:reset-not-sens} All the time transitions that
        are not sensitized by the marking $M$ have their time counters set to zero, i.e.:\\
        $\forall{}t\in{}T_i,~t\notin{}Sens(M)\Rightarrow{}I'(t)=0$.
      \end{enumerate}
    \item $\forall\tau\in\mathbb{N}$,
      $E,\tau\vdash{}s\stackrel{\uparrow}{\rightsquigarrow}s'$, where
      $s=<M,I,reset_t,ex,cond>$ and\\
      $s'=<M',I,reset_t',ex',cond>$, if:
      % $s'\xrsquigarrow{(\uparrow,\tau)}s''$,
      \begin{enumerate}[label=(\arabic*),resume]
      \item\label{item:new-marking} $M'$ is the new marking resulting from
        the firing of all the transitions contained in $Fired(s)$, i.e.:\\
        $M'=M-\sum_{t_i\in{}Fired(s)}\big(pre(t_i)-post(t_i)\big)$.
      \item\label{item:reset-order} To receive a reset order, a
        transition must be enabled by the
        marking at state $s$ and disabled by the transient marking, or must be fired at state $s$:\\
        $\forall{}t\in{}T_i,~t\in{}Sens(M)~\land~t\notin{}Sens_{bt}\big(M-\sum_{t_i\in{}Fired(s)}pre(t_i)\big)\lor{}t\in{}Fired(s)
        \Rightarrow{}reset'(t)=1$.
      \item\label{item:no-reset-order} A transition that is disabled
        by marking $M$ or enabled on its basic and test arcs by the
        transient marking, and that is not fired a state $s$, receives no reset order:\\
        $\forall{}t\in{}T_i,~
        \big(t\notin{}Sens(M)~\lor~t\in{}Sens_{bt}\big(M-\sum_{t_i\in{}Fired(s)}pre(t_i)\big)\big)
        \land{}t\in{}Fired(s)
        \Rightarrow{}reset'(t)=0$.
      % \item\label{item:reset-order} All the time transitions that are
      %   disabled by a basic or test arc\footnote{Inhibitor arcs are
      %     not considered here; indeed, the retrieval of tokens from an
      %     input place connected to a transition $t$ by an inhibitor
      %     arc can never disable $t$.} at the transient marking, (i.e,
      %   the marking resulting from the retrieval of tokens from the
      %   input places of fired transitions), or that belong to the set
      %   of fired
      %   transitions, receive a reset order, i.e.:\\
      %   $\forall{}t\in{}T_i,~t\notin{}Sens_{bt}\big(M-\sum_{t_i\in{}Fired(s)}pre(t_i)\big)\lor{}t\in{}Fired(s)
      %   \Rightarrow{}reset'(t)=1$.
      % \item\label{item:no-reset-order} All the time transitions that
      %   are sensitized by their basic and test arcs at
      %   the transient marking and do not belong to the set of fired transitions receive no reset order, i.e.:\\
      %   $\forall{}t\in{}T_i,~t\in{}Sens_{bt}\big(M-\sum_{t_i\in{}Fired(s)}pre(t_i)\big)\land{}t\notin{}Fired(s)
      %   \Rightarrow{}reset'(t)=0$.
      % \item\label{item:upper-reached} Time counters are locked for all
      %   time transitions that have reached the upper bound of their
      %   time intervals
      %   and were not fired at this clock cycle, i.e.:\\
      %   $\forall{}t\in{}T_i,~I(t)\ge{}upper(I_s(t))\wedge{}t\notin{}Fired(s)
      %   \Rightarrow{}I'(t)=\psi$.
      % \item\label{item:upper-not-reached} Time counters are unchanged
      %   for all the time transitions that have not yet reached the
      %   upper bound of their time
      %   intervals or were fired at this clock cycle, i.e.:\\
      %   $\forall{}t\in{}T_i,~I(t)=\psi\lor{}I(t)<{}upper(I_s(t))\lor{}t\in{}Fired(s)\Rightarrow$\\
      %   ${}I'(t)=I(t)$.
        
      \item\label{item:exec-fun} All functions associated with at least one fired transition ar executed, i.e:\\
        $\forall{}f\in{}\mathcal{F},~ex'(f)=\sum\limits_{t\in{}Fired(s)}\mathbb{F}(t,f)$.
      % \item\label{item:exec-fun} All the functions associated with
      %   fired transitions
      %   are executed, i.e.:\\
      %   $\forall{}f\in{}\mathcal{F},~\big(\exists{}t\in{}Fired(s),~\mathbb{F}(t,f)=1\big)
      %   \Rightarrow{}ex'(f)=1$.
      % \item\label{item:not-exec-fun} All the functions associated with none of the
      %   fired transitions are not executed, i.e.:\\
      %   $\forall{}f\in{}\mathcal{F},~\big(\forall{}t\in{}Fired(s),~\mathbb{F}(t,f)=0\big)
      %   \Rightarrow{}ex'(f)=0$.
      \end{enumerate}
    \end{itemize}
  \end{itemize}
\end{definition}

\subsection{SITPN Execution}
\label{sec:sitpn-exec}

\begin{definition}[SITPN Execution Cycle]
  For a given $sitpn\in{}SITPN$, two states $s,s''\in{}S(sitpn)$, a
  clock cycle count $\tau\in\mathbb{N}$, and an environment
  $E_c\in\mathbb{N}\rightarrow{}\mathcal{C}\rightarrow{}\mathbb{B}$,
  $sitpn$ passes from state $s$ to state $s''$ in one clock cycle,
  written $E,\tau\vdash{}sitpn,s\xrightarrow{\uparrow,\downarrow}s''$
  iff $\exists{}s'$
  s.t. $E_c,\tau\vdash{}sitpn,s\xrightarrow{\uparrow}s'$ and
  $E_c,\tau\vdash{}sitpn,s'\xrightarrow{\downarrow}s''$.
\end{definition}

\begin{definition}[SITPN Execution]
  For a given $sitpn\in{}SITPN$, a starting state $s\in{}S(sitpn)$, a
  clock cycle count $\tau\in\mathbb{N}$, and an environment
  $E_c\in\mathbb{N}\rightarrow{}\mathcal{C}\rightarrow{}\mathbb{B}$,
  $sitpn$ yields the execution trace $\theta$ from starting state $s$,
  written $E_c,\tau\vdash{}sitpn,s\rightarrow{}\theta$, by following
  the two rules below.
  
  \begin{tabular}{@{}l}
    {\fontsize{7}{10}\selectfont\textsc{ExecutionEnd}} \\
    
    {\begin{prooftree}[template={\fontsize{9}{12}\selectfont\inserttext}]
        \infer0 {$E_c,0\vdash$ $sitpn,s\rightarrow{}[~]$ }
      \end{prooftree}} 
  \end{tabular}  
  \begin{tabular}{@{}l}
    {\fontsize{7}{10}\selectfont\textsc{ExecutionLoop}} \\
    
    {\begin{prooftree}[template={\fontsize{8}{11}\selectfont\inserttext}]

        \hypo{$E_c,\tau\vdash{}sitpn,s\xrightarrow{\uparrow,\downarrow}s'$}

        \hypo{$E_c,\tau-1\vdash{}sitpn,s'\rightarrow{}\theta$}
        
        \infer2[$\tau>0$]{$E_c,\tau\vdash{}sitpn,s\rightarrow{}(s' :: \theta)$}
      \end{prooftree}} 
  \end{tabular}
\end{definition}

\begin{definition}[SITPN Full Execution]
  For a given $sitpn\in{}SITPN$, a clock cycle count
  $\tau\in\mathbb{N}$, and an environment
  $E_c\in\mathbb{N}\rightarrow{}\mathcal{C}\rightarrow{}\mathbb{B}$,
  $sitpn$ yields the execution trace $\theta$ starting from its
  initial state $s_0\in{}S(sitpn)$ (as defined in
  Def.~\ref{def:semantics}), written
  $E_c,\tau\vdash{}sitpn\rightarrow{}\theta$, by following the two
  rules below.
  
  \begin{tabular}{@{}l}
    {\fontsize{7}{10}\selectfont\textsc{FullExec0}} \\
    
    {\begin{prooftree}[template={\fontsize{8}{11}\selectfont\inserttext}]
        
        \infer0{$E_c,0\vdash{}sitpn\xrightarrow{full}[s_0]$}
      \end{prooftree}} 
  \end{tabular}
  
  \begin{tabular}{@{}l}
    {\fontsize{7}{10}\selectfont\textsc{FullExecCons}} \\
    
    {\begin{prooftree}[template={\fontsize{8}{11}\selectfont\inserttext}]
        \hypo{$E_c,\tau\vdash{}s_0\srarrow{\uparrow_0}{\fontsize{6}{8}\selectfont}s_0$}
        \hypo{$E_c,\tau\vdash{}s_0\srarrow{\downarrow}{\fontsize{6}{8}\selectfont}s$}
        \hypo{$E_c,\tau-1\vdash{}sitpn,s\rightarrow\theta_s$}
        \infer3[$\tau>0$]{$E_c,\tau\vdash{}sitpn\xrightarrow{full}(s_0 :: s :: \theta_s)$}
      \end{prooftree}} 
  \end{tabular}
\end{definition}

%%% Local Variables:
%%% mode: latex
%%% TeX-master: "../main"
%%% End:

% \chapter{\hvhdl{}: a target hardware description language}
\label{chap:hvhdl}

\begin{todobox}
  \begin{itemize}
  \item Change the arrow in associations of maps by the coma in abss
    (in rules about elab and sim of gen map, ipm and opm).
  \item Remark about variables that have persistent values in VHDL
    semantics. Not useful as all local variables always receive an
    inital value in the processes of the place and transition designs.
  \item Make a remark on the notation of design states,
    e.g. $\mathcal{E}''$ refers to the set of events of design state
    $\sigma''$ when there is no ambiguity.
  \item Two points of view to consider the semantics of VHDL:
    simulation or synthesis.  Simulation described in the LRM;
    synthesis, problem is that there are no standard, unlike Verilog
    (cite Verilog standard synthesis semantics).
  \item There are different kinds of component instantiaton statement
    but we only refer to design instantiation statment.
  \item Talk about error cases in the implementation part (?)
  \item Change the identifiers between double quotes into identifiers
    in type-writer font.
  \end{itemize}
\end{todobox}

The main purpose of this chapter is to present the target language of
the \hilecop{} transformation, i.e. the VHDL language. The
formalization and the implementation of the VHDL language syntax and
semantics is mandatory to reason about the programs generated by the
\hilecop{} model-to-text transformation. Thus, we want the reader to
clearly understand the structure and the semantics of language to be
able to fully grab the proof of semantic preservation presented in
Chapter~\ref{chap:proof}.  Specifically, we present here the \hvhdl{}
language, a synthesizable subset of the VHDL language. This subset
permits to encode the programs generated by the \hilecop{}
transformation. We devise a formal semantics for \hvhdl{} which is a
simplification of the simulation semantics of the VHDL language. The
formalization of the \hvhdl{} semantics and its implementation is one
contribution of this thesis to the many formalization of the VHDL
language found in the literature. The chapter structure is as follows.
In Section~\ref{sec:vhdl-lang-pres}, we give an informal presentation
of the language syntax and semantics. In
Section~\ref{sec:choosing-vhdl}, we present the state of the art
pertaining to the formalization of the VHDL language semantics. In
Section~\ref{sec:abstractSyntax}, \ref{sec:sem-rules},
\ref{sec:elab-rules} and \ref{sec:sim-rules}, we give the full
formalization of the \hvhdl{} language syntax and semantics, a subset
of the VHDL language. Section~\ref{sec:ex-full-sim} illustrates the
formal \hvhdl{} semantics with an example. Finally,
Section~\ref{sec:hvhdl-impl} outlines the implementation of the
\hvhdl{} syntax and semantics with the \coq{} proof assistant.

\section{Presentation of the VHDL language}
\label{sec:vhdl-lang-pres}
The intent here is to give an overview of the VHDL language, its
purpose, its main syntactal constructs, and an informal description of
its semantics as presented in the Language Reference Manual (LRM)
\cite{VHDL2000}. The VHDL language offers a lot of possibility in
terms of hardware (and even software) description. Here, we are not
trying to be exhaustive in our presentation of the language. We will
only maintain our description of the VHDL concepts in the scope that
is of interest to us. The readers that are interested in learning more
about the VHDL language can refer to \cite{VHDL2000},
\cite{Ashenden2010} and \cite{Pedroni2020}.

\subsection{Main concepts}
\label{sec:vhdl-main-concepts}

The VHDL acronym stands for Very high speed integrated circuit
Hardware Description Language. As its name indicates, the main purpose
of the VHDL language is to describe hardware circuits. There are two
approaches to the description of circuits. The first aims at the
simulation of the described circuits, and the second aims at the
synthesis of described circuits on physical supports. Thus, the
constructs of the VHDL language must be interpreted depending on the
purpose of the designer.  For instance, the language gives the
possibility to describe the connection of wires inside a circuit. A
wire is represented by the concept of \textit{signal}.  In the context
of circuit simulation, a \textit{signal} can be compared to a
variable; it has a given type and holds a value that fluctuates in the
course of the simulation. In the context of synthesis, a signal really
describes a physical wire and must be considered as so. From these two
approaches to circuit description arise two ways of considering the
semantics of the language (see Section~\ref{sec:choosing-vhdl}).

In VHDL, a top-level program is called a \emph{design}. A design
describes a hardware circuit. As explained in
Chapter~\ref{chap:hilecop}, the hilecop{} transformation generates a
VHDL design implementing the input SITPN model. To do so, the
transformation generates and connects the component instances of two
previously defined VHDL designs: the \emph{place} design that
implements the concept of a SITPN place, and the \emph{transition}
design that implements a SITPN transition. These designs were defined
by the INRIA CAMIN team at the creation of the \hilecop{}
methodology. In the following sections, we will be using excerpts of
the definition of the place and transition designs to illustrate the
content of VHDL programs and the rules of the VHDL language
semantics. The reader will find the source code of the place and
transition designs in concrete and abstract syntax in
Appendices~\ref{app:place-design} and \ref{app:trans-design}.

A VHDL design is composed of two descriptive parts. The first part is
called the entity and describes the interfaces of a circuit, namely:
the input and output ports, and the generic
constants. Listing~\ref{lst:trans-design-entity} is an excerpt of the
transition design's entity that defines the generic constants, the
input and output port interfaces of the
design. Figure~\ref{fig:trans-design-entity} is a visual
representation of the interfaces of the transition design.

\begin{lstlisting}[language=VHDL,label={lst:trans-design-entity},
caption={The entity part of the transition design in concrete VHDL
  syntax.},framexleftmargin=1.5em,xleftmargin=2em,numbers=left,
numberstyle=\tiny\ttfamily]
entity transition is
  generic(
    transition_type      : transition_t := NOT_TEMPORAL;
    input_arcs_number    : natural := 1; 
    conditions_number    : natural := 1;
    maximal_time_counter : natural := 1
    );        
  port(
    clock                   : in std_logic;
    reset_n                 : in std_logic;
    input_conditions        : in std_logic_vector(conditions_number-1 downto 0);
    time_A_value            : in natural range 0 to maximal_time_counter;
    time_B_value            : in natural range 0 to maximal_time_counter;
    input_arcs_valid        : in std_logic_vector(input_arcs_number-1 downto 0);
    reinit_time             : in std_logic_vector(input_arcs_number-1 downto 0);
    priority_authorizations : in std_logic_vector(input_arcs_number-1 downto 0);
    fired                   : out std_logic 
    );
end transition;
\end{lstlisting}

The generic clause of the entity holds the declaration of the generic
constants.  The purpose of generic constants is either to represent
some dimensions of the design (e.g. the size of ports, internal
signals\dots) or to represent constant values used throughout the
design. In Listing~\ref{lst:trans-design-entity}, one can see that the
\texttt{conditions_number} generic constant gives a dimension to the
type of the \texttt{input_conditions} input port, which is an array of
Boolean values with indexes ranging from 0 to
\texttt{conditions_number-1} (that is the meaning of
\texttt{std_logic_vector (conditions_number-1 downto 0)}.  The port
clause holds the declaration of input and output ports of the
design. The \texttt{in} keyword indicates the declaration of an input
port and the \texttt{out} indicates the declaration of an output port.

\begin{figure}[H]
  \centering
  \includegraphics[keepaspectratio=true, width=.6\textwidth]{Figures/H-VHDL/trans-design-entity}
  \caption[A representation of the transition design entity.]{A
    representation of the transition design entity. On the left side,
    the input port interface of the transition design; \textit{cn}
    stands for \textit{conditions\_number} and \textit{ian} stands for
    \textit{input\_arcs\_number}, i.e. two of the generic constants
    declared in the generic clause of the transition design entity;
    the numbers at the right of the input pins represent the pin
    indexes. On the right side, the output port interface of the
    transition design.}
  \label{fig:trans-design-entity}
\end{figure}

The second part of a VHDL design is called the architecture. The
architecture describes the internal behavior of the design. It
declares all the internal signals involved in the description of the
design behavior. Then, there are three ways to describe the behavior
itself: by using processes, by instantiating other designs (also
called, component instantiations), or by combining both technics (the
latter option is chosen in the VHDL designs generated by the
\hilecop{} transformation).

\paragraph{Behavior specification with processes}

The first way is to specify one or multiple processes. Processes are
concurrent statements that describes the wiring or the operations
performed on the signals of a given design.  A process declares a
sensitivy list that corresponds to the signals read in the process
statement body; also, it possibly declares internal variables. The
sensitivity list is only useful for the purpose of simulation. It
permits to resume the execution of a process when the value of one of
the signals of its sensitivity list
changes. Listing~\ref{lst:trans-design-arch} gives an excerpt of the
transition design architecture containing the declarative part of the
architecture (i.e. internal signals) and three of the processes
describing the transition design behavior, namely: the
\texttt{condition_evaluation} process, the \texttt{firable} process
and the \texttt{fired_evaluation} process.

\begin{lstlisting}[language=VHDL,label={lst:trans-design-arch},
caption={An excerpt of the architecture part of the transition design in concrete VHDL
  syntax.},framexleftmargin=1.5em,xleftmargin=2em,numbers=left,
numberstyle=\tiny\ttfamily]
architecture transition_architecture of transition is
  signal s_condition_combination : std_logic; #\label{line:scc}#
  signal s_enabled               : std_logic;
  signal s_firable               : std_logic;   
  signal s_firing_condition      : std_logic;
  signal s_priority_combination  : std_logic;
  signal s_reinit_time_counter   : std_logic; 
  signal s_time_counter          : natural range 0 to maximal_time_counter; #\label{line:stc}#
begin
  
  condition_evaluation : process(input_conditions) #\label{line:condev-ps}#
    variable v_internal_condition : std_logic;
  begin
    v_internal_condition := '1';
    
    for i in 0 to conditions_number - 1 loop #\label{line:for-begin}#
      v_internal_condition := v_internal_condition and input_conditions(i);
    end loop; #\label{line:for-end}#
    
    s_condition_combination <= v_internal_condition; #\label{line:assign-scc}#
  end process condition_evaluation;
    
  $\dots$

  firable : process(reset_n, clock)
  begin
    if (reset_n = '0') then
      s_firable <= '0';
    elsif falling_edge(clock) then
      s_firable <= s_firing_condition;
    end if;
  end process firable;

  fired_evaluation : process (s_firable, s_priority_combination)
  begin
    fired <= s_firable and s_priority_combination;
  end process fired_evaluation;      
  
end transition_architecture;
\end{lstlisting}

In Listing~\ref{lst:trans-design-arch}, from Line~\ref{line:scc} to
Line~\ref{line:stc}, the architecture declares the internal signals of
the transition design. Then, Line~\ref{line:condev-ps} begins the
declaration of the \texttt{condition_evaluation} process. The
sensitivity list of the \texttt{condition_evaluation} process holds
one signal, the \texttt{input_conditions} input port, and declares a
local variable \texttt{v_internal_condition}.

In the statement body of a process, the designer can use control flow
statements common to most of the generic programming languages (if
statement, for loops\dots), and also statements that are specific to
the VHDL language. The most representative statement, and the one of
interest to us, is the \emph{signal assignment} statement. The signal
assignment statement relate a given signal identifier to a source
expression. For instance, at Line~\ref{line:assign-scc} of
Listing~\ref{lst:trans-design-arch}, the signal assignment statement,
represented with the $\Leftarrow$ operator, assigns the value of the
internal variable \texttt{v_internal_condition} to signal
\texttt{s_condition_evaluation}; the \texttt{v_internal_variable} that
itself holds the Boolean product between the members of the
\texttt{input_conditions} input port performed in the \texttt{for}
loop of Lines~\ref{line:for-begin} to \ref{line:for-end}.

When considering a VHDL design in the point of view of hardware
synthesis, a signal assignment statement specifies a wiring between a
target signal identifier and other source
signals. Figure~\ref{fig:trans-design-arch-excerpt-1} gives a
synthesis-oriented view of the processes described in
Listing~\ref{lst:trans-design-arch}.

\begin{figure}[H]
  \centering
  \includegraphics[keepaspectratio=true, width=.8\textwidth]{Figures/H-VHDL/trans-design-arch-excerpt-1}
  \caption[Representation of a part of the transition design
  architecture.]{A representation of a part of the transition design
    architecture comprising three processes. On the left side, the
    \texttt{condition_evaluation} process connecting the
    \texttt{input_conditions} input port to the
    \texttt{s_condition_combination} internal signal; the
    \texttt{firable} process in the middle; on the right side, the
    \texttt{fired_evaluation} process connecting the
    \texttt{s_firable} and the \texttt{s_priority_combination} signals
    to the \texttt{fired} output port.}
  \label{fig:trans-design-arch-excerpt-1}
\end{figure}

In Figure~\ref{fig:trans-design-arch-excerpt-1}, the
\texttt{condition_evaluation} process is represented as an
\texttt{and} port performing the product over the elements of the
\texttt{input_conditions} input port. The \texttt{fired_evaluation}
process is a simple \texttt{and} gate connecting the \texttt{fired}
output port to the \texttt{s_firable} and
\texttt{s_priority_combination} internal signals. The \texttt{firable}
process is a \textit{synchronous} process. It responds to the event of
the \texttt{clock} signal. In its statement body, the \texttt{firable}
process assigns the value of the internal signal
\texttt{s_firing_condition} to the signal \texttt{s_firable} only at
the occurrence of the falling edge of the \texttt{clock} signal
(captured by the expression \texttt{falling_edge(clock)} where
\texttt{falling_edge} is a primitive function of the VHDL
language). In the point of view of simulation, there are no
distinction between synchronous processes and \emph{combinational}
processes, i.e. processes that are executed independently of the
occurrence of a clock signal. However, in the point of view of
synthesis, processes responding to a \texttt{clock} signal follow the
rules of the synchronous logic, whereas, combinational processes
follow the rules of combinational logic.

To complete the presentation of the statements to be found in the body
of processes, the VHDL language is also equipped with \texttt{timing}
constructs, i.e. statements that explicitly specify an amount of time
in a given time unit. The signal assignement statement possibly
specifies a time clause indicating when the assignment must be
performed. For instance, the signal assignment statement specifying
that the value of signal \texttt{b} must be assigned to signal
\texttt{a} in 3 milliseconds takes the form: \texttt{$a\Leftarrow{}b$
  in 3 ms}. When no time clause is specified for a signal assignment
statement, we talk about a $\delta$-delay signal assignment, i.e. the
application of the signal assignment is related to some $\delta$
interval corresponding the propagation time through a wire. When a
time clause is specified, we talk about an unit-delay signal
assignment. $\delta$-delay signal assignments are synthetizable,
meaning they have an equivalent implementation on a physical device,
whereas, unit-delay signal assignments can not be
synthetized. Unit-delay signal assignments do not appear neither in
the VHDL designs generated by \hilecop{} transformation nor in the
declaration of the place and transition designs. We are only
mentioning their existence here because they are the witnesses of the
two time paradigms that inhabit the simulation algorithm describe the
semantics of the VHDL language: $\delta$ time and real time.

\paragraph{Behavior specification with design instances}

The second way to specify the behavior of a design is to use other
designs, or rather instances of other designs, as components of the
circuits. In that case, the design is said to be composite as it
embeds instances of other designs in its own behavior. Also, a design
at the highest level of embedding, i.e. that is not instantiated as a
part of another design's behavior, is called a \emph{top-level}
design. The design instantiation, or component instantiation,
statement permits to instantiate a design in an embedding
architecture. When instantiating a design with a design instantiation
statement, the designer provides the component instance with an
identifier. Then, the design instance must be dimensioned; this is
performed through a generic map that associates the generic constants
of the design being instantiated to a static value. Finally, the
designer specifies how the component instance is connected to the
other elements of the architecture. A port map associates the input
ports and output ports of the component instance to expressions or to
the signals of the embedding
architecture. Listing~\ref{lst:trans-design-instantiation} shows an
example of instantiation of the \hilecop{}'s transition design. This
instance is involved in the definition of the behavior of an embedding
design called \texttt{toplevel}.

\begin{lstlisting}[language=VHDL,label={lst:trans-design-instantiation},
caption={An example of design instantiation statement in the architecture of the \texttt{toplevel} design. Here, the design being instantiated is the transition design.},framexleftmargin=1.5em,xleftmargin=2em,numbers=left,
numberstyle=\tiny\ttfamily]
architecture toplevel_architecture of toplevel is
begin
  #\dots#
  $id_t$ : entity transition 
  generic map (
    transition_type => NOT_TEMPORAL,
    input_arcs_number => 1,
    conditions_number => 1,
    maximal_time_counter => 1
  )
  port map (
    clock => clock,
    reset_n => reset_n,
    time_A_value => 0,
    time_B_value => 0,
    input_conditions(0) => $id_0$,
    input_arcs_valid(0) => $id_1$,
    priority_authorizations(0) => '1',
    reinit_time(0) => $id_2$,
    fired => $id_3$
  );
  #\dots#
end toplevel_architecture;
\end{lstlisting}

In Listing~\ref{lst:trans-design-instantiation}, the transition
component instance has the identifier $id_t$. Following the
\texttt{entity} keyword is the name of the design being instantiated;
here, the transition design. Then, the generic map associates the
generic constants of the transition design (i.e. the left side of the
arrow, also called the \emph{formal} part) to static values (i.e. the
right side of the arrow called the \emph{actual} part). This permits
the dimensiniong the component instance. For example, remember that
the \texttt{input_arcs_number} generic constant value determines the
number of elements in the composite input ports
\texttt{input_arcs_valid}, \texttt{priority_authorizations} and
\texttt{reinit_time}. The port map associates the input ports of the
transition design to expressions. For instance, the
\texttt{time_A_value} input port is connected to the constant value
$0$, and the \texttt{input_conditions} input port is connected to the
internal signal $id_0$ at index $0$. The port map also associates the
output ports with signal identifiers. Contrary to the association of
input ports, output ports can not be associated to expressions as
output port association describes a direct wiring. In the port map
described in Listing~\ref{lst:trans-design-instantiation}, the
association $\mathtt{fired}\Rightarrow{}id_3$ expresses that the
\texttt{fired} output port is connected to the signal $id_3$, where
signal $id_3$ is defined in the embedding
design. Figure~\ref{fig:trans-design-instance} illustrates the
transition design instance $id_t$ and the wiring of its input and
output port interfaces inside the toplevel design.

\begin{figure}[H]
  \centering
  \includegraphics[keepaspectratio=true, width=.8\textwidth]{Figures/H-VHDL/trans-design-instance}
  \caption[Visual representation of a design instantiation
  statement.]{Visual representation of a design instantiation
    statement. Here, the figure represents the transition design
    instance described in
    Listing~\ref{lst:trans-design-instantiation}.}
  \label{fig:trans-design-instance}
\end{figure}

% The wait statement stalls the execution of a process
% during a certain amount of time, or until some event arises on a given
% signal. An event arises on a signal when its value changes. Every
% process declared in an architecture possess a implicit wait statement
% at the end of its statement body. This statement waits on an event on
% one of the signals of the sensitivity list. This implicit wait
% statement is only useful for the purpose of simulation. Indeed, the
% wait statement is no synthetizable. 

\subsection{Informal semantics of the VHDL language}
\label{sec:vhdl-informal-sem}

Even though, in practice, there are two ways to consider a VHDL
design, i.e. a synthesis-oriented way and a simulation-oriented way,
the LRM does not define a synthesis-oriented semantics for the VHDL
language.  A synthesis-oriented semantics gives an interpretation to a
design by describing an equivalent in a lower level formalism, closer
to the physical circuit. For instance, the Verilog language gives a
synthesis-oriented semantics to its programs by defining an equivalent
RTL level description \cite{Verilog2005}.  The LRM gives an informal
semantics to VHDL designs through the definition of a simulation
algorithm. The purpose of simulation is to compute the evolution of
the values of signals during a certain time interval. Through the
simulation process, the designer is able to control the behavior of
the modeled circuits and to detect flaws in the evolution of the
signal values.

Former to the simulation, the LRM defines an elaboration phase that
permits to transform a design into a simulation-ready execution
model. The elaboration phase has several goals. First, it builds the
simulation environment and a starting design simulation state. The
simulation environment is built based on the declarative parts of the
top-level design; it maps the signals to their types. In the design
simulation state, each signal is associated with a current value and
with a driver. A driver maps time points to values and the association
between a given time point and a signal value is called a
transaction. The necessity of drivers is explained by the presence of
unit-delay signal assignments. A unit-delay signal assignment specify
a time clause indicating when a giving assignment must be performed,
e.g. \texttt{$a\Leftarrow{}b$ in 3ms} (signal \texttt{a} takes the
value of signal \texttt{b} in 3 milliseconds). Thus, when a unit-delay
signal assignment is executed in the course of a simulation, its
effect is to change the driver of the target signal by posting a new
transaction. For instance, let $T_c$ by the current simulation time,
the execution of statement \texttt{$a\Leftarrow$ true in 2ns} sets a
new transaction in the driver of signal $a$. The new transaction
associates the value \texttt{true} to the time point
$T_c+\mathtt{2ns}$. Note that without unit-delay signal assignments,
drivers are not needed as all assignments take their effects at the
current simulation time.
% Also, each signal is associated with a
% graph representing the relation to its source, i.e. the signals
% involved in the computation of the target signal value. For instance,
% the source signals of the \texttt{fired} output port represented in
% Figure~\ref{fig:trans-design-arch-excerpt-1} are the
% \texttt{s_firable} and the \texttt{s_priority_combination}
% signals. 
Second, the elaboration checks the well-formedness of the design by
performing static type-checking on the behavioral part of the
design. It also checks that the connection between signals respect
certain rules, for instance, that there are no multiply-driven
signals, i.e. signals that are written to by multiple
processes. Finally, the elaboration operates some transformations over
the VHDL code, and thus builds the \emph{execution} model. To
summarize, all concurrent statements of the behavioral part are
transformed until the top-level design behavior is only composed of
processes.

After the elaboration, the top-level design, or rather its
corresponding execution model, is ready to be simulated. Two entities
are involved in the simulation: the \emph{sea} of processes obtained
after the elaboration of the top-level design, and a \emph{kernel}
process. The kernel process orchestrates the simulation; it handles
the time of the simulation, i.e. it holds a variable describing the
current time of the simulation, and resumes the execution of
processes. Figure~\ref{fig:vhdl-sim-alg}, which is an excerpt from
\cite{Borger1995}, describes the structure of the VHDL simulation
algorithm.

\begin{figure}[H]
  \centering
  \includegraphics[keepaspectratio=true, width=.7\textwidth]{Figures/H-VHDL/vhdl-sim-alg}
  \caption[The VHDL simulation loop.]{The VHDL simulation loop. Excerpt from \cite{Borger1995}.}
  \label{fig:vhdl-sim-alg}
\end{figure}


The simulation starts with the initialization phase. During the
initialization phase all processes are run exactly once. Then, the
simulation cycles are structured as follows.  All processes execute
their statement body concurrently. New transactions are posted in the
drivers of signals for every interpreted signal assignment
statement. The execution goes on until all processes have executed
their statement body and then have reached a suspension state. When,
all processes are suspended, the kernel process takes
over. Figure~\ref{fig:kernel-ps} shows the activity diagram associated
with the kernel process.

\begin{figure}[H]
  \centering
  \includegraphics[keepaspectratio=true, width=.35\textwidth]{Figures/H-VHDL/kernel-ps}
  \caption[The activity diagram of the kernel process.]{The activity
    diagram of the kernel process. Square boxes represent activities,
    diamond nodes are decision nodes. The black circle at the top
    represents the starting point of the activities; the other black
    circle in the middle of the diagram represents the end of all
    activities.}
  \label{fig:kernel-ps}
\end{figure}

As shown in Figure~\ref{fig:kernel-ps}, the kernel process will then
determine the kind of simulation cycle that will be performed next.
There are two kinds of cycles: delta cycle or time cycle. If the value
of a signal changes at the current time point, i.e. its driver holds a
transaction at the current time point with a new value, then a delta
cycle must be performed.  Then, the simulation time does not change.
The kernel process updates the current value of signals and their
drivers, and wakes up the processes sensitive to the signals that
obtained new values. The repetition of multiple delta cycles
corresponds to the stabilization of signal values, i.e. the
propagation of values through the wires, that takes effect in a
negligeable $\delta$ time. If all signal values are stable at the
current time point, then a time cycle must be performed.  The kernel
process looks up the drivers for the next time point where the value
of a given signal will change. Then, the kernel process advances the
simulation time to this next time point before updating the signal
values and resuming the execution of processes.  The simulation goes
on like this, alternating between delta and time cycles, until the
current time value reaches the time specified for the end of the
simulation.

%%% Local Variables:
%%% mode: latex
%%% TeX-master: "../../main"
%%% End:


\section{Choosing a formal semantics for VHDL}
\label{sec:choosing-vhdl}
In the previous section, we presented the main concepts underlying the
\vhdl{} language and its informal semantics. We want to prove that the
\hilecop{} transformation that generates \vhdl{} code from SITPNs
preserves the behavior of the initial model (i.e, the SITPN model)
into the generated \vhdl{} program.  A formal semantics for the \vhdl{}
language is therefore a necessary element to be able to reason about
the generated \vhdl{} programs, and moreover to be able to compare their
behaviors with the behaviors of the source SITPN models. Keeping that
in mind, which formal semantics should we consider for \vhdl{}?

The same holds for any task: there is a trade off between finding a
tool designed by others that will fit our needs, and creating our own
tool that will mitigate the gaps between our needs and what is
available in the literature. In the present case, the tool is a formal
semantics for \vhdl{}. Adopting a fully-set semantics found in the
literature as a base ground for the implementation of a formal
semantics for \vhdl{} has multiple perks. First, it reduces the
formalization effort, which is not a lesser point considering that the
proof ahead might be long and must still be completed within the time
span of the thesis. Still, the semantics would need to be implemented
in \coq{}, if no implementation exists (or not written in
\coq{}). Second, the formal semantics of programming languages found
in the literature are often general in their approach, this to provide
a generic framework to reason about programs. However, we must not
lose sight of our goal which is to prove behavior preservation; a
generic formal semantics could turn out to be too complex, or
necessitate too much tweaking and thus hinder the fulfillment of our
task. On the other side, creating our own formal semantics for
\vhdl{}, based on the work of others, is the best way to fit our needs
in compliance with our final aim. However, the pitfalls are that the
resulting semantics might prove to be very specific, therefore
preventing others from using it. Also, a work of formalization would
be necessary which, as we already stated, would be time-consuming. In
order to determine whether we ought to use an existing semantics or
design a new one, we must first clearly specify our needs regarding
the \vhdl{} language.

\subsection{Specifying our needs: \hilecop{} and \vhdl{}}
\label{subsec:specifying-needs}

Two elements are of major influence to the specification of our needs
for a formal semantics: first, the context of \hilecop{} and the
specificities of the \vhdl{} programs that are generated; second, the
context of theorem proving. These two aspects entail the following
considerations.

\paragraph{The need for coverage} The \hilecop{} methodology generates
particular \vhdl{} programs.  Even if some transformations can be
operated on the generated programs to simplify them, the looked-for
formal semantics must be able to deal with a certain subset of the
\vhdl{} language. Especially, this subset must include:

\begin{itemize}
\item 0-delay (or $\delta$-delay) signal assignments (equivalent to
  unit-delay signal assignment with a ``0 ns'' after clause)
\item component instantiation statements with generic constant and
  port mapping
\item entity's generic constant clauses (declaration of generic
  constants in a design entity)
\end{itemize}

\hilecop{}'s \vhdl{} programs only deal with 0-delay signal
assignments because they are the only kind of signal assignments that
can be synthesized. As a matter of fact, the industrial
compiler/synthesizer used in the \hilecop{} methodology only accepts
\vhdl{} programs with no timing constructs (i.e, no delayed signal
assignments) as inputs. % However, 0-delay signal assignments do not
% specify a delay after which the assignment of a value to a signal must
% be performed. The time taken by a signal assignment to be performed
% (the time of propagation) is only relevant in relation to the clock
% period that regulates the modeled circuit. It is sufficient if the
% semantics handles only this dimension of time, even if the \vhdl{}
% language gathers both the ``dimension of delta time and the dimension
% of real time'' \cite{Reetz95}.

Regarding component instantiation statements, the \vhdl{} LRM
describes a way to transform these statements into equivalent process
statements and block constructs \cite[p.~141]{VHDL2000} during the
elaboration of the design. However, we want to preserve the
hierarchical structure provided by the component instantiation
statements arguing that it will be easier to compare the state of a
given SITPN model with a \vhdl{} design state with an explicit
hierarchical structure. Indeed, there exists a mapping between places
and transitions of an SITPN and their mirror (generated by the
transformation) place and transition component instances (PCIs and
TCIs). This one-to-one correspondence might turn out to be handy to
perform the proof of behavior preservation. Obviously, the semantics
must cover the evaluation of process statements which are the core
concurrent statements of \vhdl{} programs.

The types of signals and variables used in \hilecop{} \vhdl{} designs
must have finite ranges of values. For instance, a \vhdl{} signal that
ranges over $\mathbb{N}$ cannot be synthesized on a physical circuit.
Indeed, $\mathbb{N}$ has an infinite number of values, and would
therefore require an infinite number of latches to be physically
implemented. Moreover, as the number of latches used to implement a
digital circuit greatly impacts the power consumption of the circuit,
the types of signals and variables must be as constrained as possible
to optimize the dimensioning of the circuit. The generic constants,
declared in the entity part of a design, are involved in the
dimensioning of the circuit. The generic constants define the bound of
the array and natural range types for the different signals and
variables declared in the \texttt{place} and \texttt{transition}
designs' architecture. When a place or a transition component is
instantiated, that is during the transformation of the SITPN model
into \vhdl{} code, its generic constants receive values via a generic
map; we call it the dimensioning of the component instance. Therefore,
generic constant clauses must belong to the subset of the \vhdl{}
language covered by the semantics.

\paragraph{The need for a synchronous execution} The second property
of \hilecop{}'s generated \vhdl{} programs is their synchronous
execution. The digital circuits designed with the \hilecop{}
methodology are all synchronously executed on physical target. The
generated \vhdl{} designs declare a clock signal as an input port of
their entity port interface. Thus, the behavioral part of the designs
contains two kinds of processes: \emph{synchronous} processes, i.e.
processes that are sensitive to the clock signal, and
\emph{combinational} processes, i.e. processes that are not sensitive
to the clock signal, and that are permanently running until the
stabilization of the signal values. Synchronous processes react to the
events of the clock signal, i.e. the rising and the falling edge, and
possess blocks of sequential statements that are only executed at the
precise moment of the clock event\footnote{These blocks are guarded by
  the expressions \texttt{rising\_edge(clk)} and
  \texttt{falling\_edge(clk)}.}. Therefore, we need a semantics that
is able to deal with synchronism, and that explicitly integrates the
synchronization with a clock signal into the expression of the
simulation
cycle. % Thus, the two dimensions of time, respectively delta
% time and ``real'' time \cite{Reetz1995}, that are part of the \vhdl{}
% language, must be handled by our semantics.
% Delta time pertains to
% delta-cycles, which are triggered by the execution of delta-delay
% signal assignments, that is, the only kind of signal assignments found
% in \hilecop{} \vhdl{} programs. Real time pertains to time-steps. When
% all signal values are stable (when there are no more delta-cycles),
% the simulation advances the current time value to the next time point
% where a relevant event happens.
% A relevant event could be the
% execution of the next pending signal assignment that was associated
% with a delay, or the end of a \texttt{wait for x} statement (where
% \texttt{x} specifies a time delay). In our case, neither delayed
% signal assignments nor wait statements are part the covered \vhdl{}
% language subset.
% In our case, the only two relevant events to which the simulation
% advances during a time-step are the rising edge and the falling edge
% of the clock signal.

% \paragraph{Other considerations} Considering the kind of proof that
% needs to be established, we would rather consider an operational
% semantics for \vhdl{}. The reason is that, in the \ccert{} project
% \cite{Leroy2009}, which is one of our major inpsiration source, the
% whole C compiler toolchain is verified by reasoning over the
% operational semantics of the source and target languages.
A last
consideration pertains to whether or not the \vhdl{} semantics must
explicitly handle errors. As the SITPN semantics does not include the
production of error values, the handling of errors by the \vhdl{}
semantics is not a mandatory aspect.

\paragraph{Qualifying criterions}
\label{sec:qualifying-criterions}

We here give the list of the qualifying criterions that will help to
analyze the different \vhdl{} semantics encountered in the literature
and that are presented in the next section. The three most relevant
criterions are:

\begin{itemize}
\item \emph{Synchronism}. Regarding this criterion, there are three
  possibilities:
  \begin{itemize}
  \item Synchronism is not expressible in the considered \vhdl{}
    semantics; this completely disqualifies the adoption of the
    semantics.
  \item Synchronism is expressible in the considered \vhdl{}
    semantics. Synchronism is expressible if time-steps are handle in
    the semantics, at least to be able to represent clock events.
  \item Synchronism is explicit, i.e. the simulation loop is built
    around the occurrences of clock events.
  \end{itemize}
  We will foster the semantics that explicitly formalize a
  synchronized execution of a \vhdl{} design.
\item \emph{Component instantiation}. Either the semantics handle the
  component instantiation statement in its simulation rules, or
  component instantiation statements must be transformed in order to
  be executed. We will foster the semantics that handle component
  instantiation statements without transformation.
\item \emph{Elaboration}.  This criterion pertains to the
  formalization of the elaboration phase as integrated to the \vhdl{}
  semantics. This criterion also expresses the ability of the
  semantics to handle constrained types, i.e. arrays and natural
  ranges, and generic constant clauses that are both dealt with during
  the \emph{elaboration} phase. Either the semantics handle these
  constructs or it does not. Of course, we will foster the first kind
  of semantics.
\end{itemize}

\subsection{Looking for an existing formal semantics}
\label{subsec:looking-for-sem}

Here, we give a summary of the work found in the literature pertaining
to the formalization of the \vhdl{} language semantics. Articles are
gathered and presented depending on the type of semantics used in the
formalization (operational, denotational, axiomatic\dots). Each
semantics is analyzed regarding the needs that were previously
expressed.

\paragraph{Denotational semantics}
Some authors have been interested in giving a formal denotational
semantics to \vhdl{}. In a general manner, these authors want to reason
about \vhdl{} programs: prove properties over a \vhdl{} program, prove that
two programs are equivalent\dots

In \cite{Fuchs1995}, the authors give a denotational semantics to the
\vhdl{} language within the \textsf{Focus} \cite{Dederichs1993}
framework, a method for the development of distributed systems. Signal
values and their evolution through time are represented as streams of
values. Statements are denoted as stream-processing
functions. Processes are stream-processing functions that takes input
signal streams (signals of the sensitivity list) and yields
transaction traces (i.e, waveforms) over output signals (i.e, signal
that are written by the process). Transaction traces are merged
together as the result of the concurrent execution of
processes. % Resolution functions are used in case of multiply-driven
% signals (i.e, signals that receive a value from multiple processes).
The authors only consider 0-delay signal assignments in their
semantics, stating that it is sufficient to ``consider time at a
logical level to model both synchronous and asynchronous designs''.
However, some transformations must be applied to a design that has a
synchronous execution to express its equivalent only with 0-delay
signal assignments. Therefore, this semantics does not express
synchronism of execution in an explicit manner. Moreover, the
component instantiation statements are not dealt with, and no mention
is made of the elaboration phase.

In \cite{Breuer1995a}, the authors give a denotational, yet
relational, semantics to the \vhdl{} language. A state of a \vhdl{}
design is represented by a function binding signals to values; a
worldline is a time-ordered list of states. Statements (including
processes) are denoted in the semantics by a relation that binds an
input couple, composed of a time point and a worldline, to an output
couple of the same type. Multiple input and output couples possibly
satisfy the relation denoting a particular statement; thus, the
semantics is nondeterministic.  The semantics tries to abstract from
the formalization of the simulation cycle as it is done in the
LRM. The authors want to establish a semantics that is abstract enough
to be able to compare all other works of formalization with the
authors' semantics. The authors also give an axiomatic semantics (i.e,
in the Hoare logic style) which is proved to be sound and complete
with the first denotational semantics. A \textsf{Prolog}
\cite{Colmerauer1990} implementation of the axiomatic semantics is
given. Regarding our needs, the semantics only deals with unit-delay
signal assignments. However, this semantics enables the representation
of a $\delta$-delay signal assignment with a unit-delay signal
assignment adorned with a ``\texttt{after 0 ns}'' time clause. The
hierarchical structure of designs is not preserved, and, although
expressible, the semantics does not explicitly express a synchronous
simulation cycle.

The denotational semantics expressed in \cite{Pandey1999} uses
interval temporal logic as an underlying model. Leveraging this
underlying model, the authors are interested in proving some
properties over \vhdl{} designs to help compilers to optimize the code,
for instance, by using rewrite rules proved to be valid against the
model. Some of the proofs laid out by the authors are embedded in PVS
\cite{Owre1994}. The expression of the dynamic model uses many
concepts described in the LRM, like drivers, port association, driving
and effective values for signals. The semantics deals with both
unit-delay and $\delta$-delay signal assignments. The semantics works
on fully-elaborated designs, therefore, it does not deal with
component instantiation statements. Moreover, interval temporal logic
is useful to reason on the \vhdl{} designs in the presence of delays,
however, it looses its interest for designs presenting only 0-delay
assignments.

In \cite{Borrione1995}, the author states that ``denotational
semantics is more adequate for mathematical reasoning''. The author
formalizes the \vhdl{} semantics to prove the equivalence between
\vhdl{} programs (for instance, a specification and an
implementation). What is of major interest regarding our needs is that
the author has expressed a simulation cycle for synchronous
designs. Therefore, a distinction is made between combinational and
synchronous processes in the abstract syntax. Moreover, this work
formalizes the elaboration part of a \vhdl{} design former to the
simulation; also, the elaboration keeps the hierarchical setting of
the \vhdl{} design, that is component instantiation statements are not
replaced by processes.  Due to the time abstraction, the semantics
only deals with 0-delay signal assignments. It is explained by the
fact that the reference time-unit is the clock period (i.e, the only
known time-step), and the advancing of time, happening during the
simulation cycle as described in the LRM, is captured within the
setting of the simulation cycle. % Also, the semantics takes primary
% inputs into account (i.e, input ports of the top-level design); to
% preserve a synchronous behavior for the simulated design, the
% hypothesis is made that the values of the primary inputs are stable
% between two clock events.
% The only critic that can be made to
% this semantics regarding our needs is that it is expressed in
% denotational style.

\paragraph{Operational semantics}

Multiple works formalize an operational semantics for \vhdl{}.  These
works are interested in the formal description of the \vhdl{} simulator.
The aim is to devise a formal semantics that acts as a formal
specification for a simulator.

In \cite{Breuer1995}, a formal description of a \emph{functional}
semantics for \vhdl{} is laid out based on stream-processing
functions. The semantics is expressed with the functional programming
language \textsf{Gofer} \cite{Jones1994}, thus enabling the computation
of execution traces, that is, the computation of the streams
representing the values taken by signals over time.  As in the former
work of the same author \cite{Breuer1995a}, only unit-delay
signal assignments are dealt with, however, this time the author
describes a deterministic operational semantics. Regarding our needs,
this work is neither interested in preserving the hierarchical
structure of \vhdl{} designs, and no mention is made regarding how a
design is elaborated, nor in expressing an explicit synchronous
simulation cycle.

In \cite{Borger1995}, the authors formalize the simulation loop of the
LRM using Evolving Algebra machines (EA-machines). All important
constructs of the \vhdl{} language are represented as records; processes
are represented as concurrent agents running pseudo-codes, and the
simulation control flow is passed to and fro between the kernel
process (i.e, the simulation orchestrator) and the rest of the
processes that execute the design behavior. This semantics implements
closely the simulation loop as described in the LRM. Therefore, it is
very rich and deals with most of the \vhdl{} constructs, including the
two time paradigms of the language (i.e. $\delta$ time and unit
time). Moreover, the semantics works on fully-elaborated designs,
therefore, component instantiation statements are omitted. However, a
synchronous execution is fully expressible even if not explicitly
embedded in the expression of the simulation loop.

In \cite{VanTassel1995}, the author presents a natural semantics for
\vhdl{}. The simulation loop is expressed by inference rules, and the
execution of processes is based on the events over signals of their
corresponding sensitivity lists. The execution of statements computes
transaction traces, that is, the drivers of the signals. The semantics
deals both with unit and delta delay signal assignments. Regarding our
needs, this semantics does not entirely cover the subset of \vhdl{} we
are interested in. Component instantiation statements are not dealt
with.  A synchronous execution is expressible within the semantics,
although it would be hidden in the inference rule formalizing the
generic simulation loop.  Also, the semantics does not provide its
simulation loop with a simulation horizon (a maximum number of
simulation cycles). The simulation ends when signal values evolve no
more.  % The question of the influence of the
% environment, measured through the values of the primary inputs of a
% design, is not discussed.

In \cite{Goossens1995}, the author presents an operational semantics
for \vhdl{} in the small-step style. The semantics follows closely the
simulation cycle described in the LRM; however it is very concise and
clear. The covered \vhdl{} subset comprises both unit and delta-delay
signal assignments. There is an interesting discussion about the
non-determinism of \vhdl{}, since it is a concurrent programming
language: it entails that non-determinism is only existent at the
processes level, that is, internal sequential statement of processes
can be executed in a nondeterministic manner (referred to as A
actions, that is, \emph{internal} actions).  But at every delta or
time step (referred to as $\delta$ and T actions) of the execution,
the design state can be computed in a deterministic manner, since all
processes have reached a suspension point at the end of their inner
body. The author is interested in comparing the behaviors of two
\vhdl{} designs by proving that some relation of equivalence holds
between the two. He describes two strategies to compare \vhdl{}
programs. The first one is bisimulation; it is based on the comparison
of the sequence of actions (either A, $\delta$ or T actions) performed
by the two programs. The second one is observational equivalence; it
is based on the observation of the value of the output signals of two
\vhdl{} programs (the observees), that receive values in their input
signals from another \vhdl{} program (the observer). The observer
stimulates the entries of the observees and reaches a success state
based on its observations of the value of the outputs. Regarding our
needs, this semantics permits the description of our synchronous
simulation cycle. However, like most of the semantics presented here,
the component instantiation statement is not supported as it stands,
but it is rather transformed into the equivalent processes
statements. Small-step semantics is not needed in our case because we
are only interested in the values of signals at the delta and time
steps (for us, time steps correspond to clock events). We are not
interested in capturing the design states in the middle of the
execution of a process body. We are more interested in "weak
bisimulation", therefore forsaking the internal actions % (i.e, A
% actions, execution of a process body that does not end in a wait
% statement)
performed by a \vhdl{} design.
% A natural operational semantics in the style of Van Tassel's
% \cite{VanTassel1995} is sufficient in our case.
In \cite{Thirunarayan2001}, the authors extend the work of
\cite{Goossens1995}, especially by handling shared variables, in the
presence of which a \vhdl{} program can have a concrete
nonderterministic behavior. The authors are also interested in the
equivalence between two \vhdl{} programs, and they are interested in
determining a unique meaning property for \vhdl{} programs. The unique
meaning property states that the execution of a \vhdl{} design in the
presence of shared variables is unique. This work is interesting as it
points out the fact that the \vhdl{} language is not only subject to
``\textit{benign} nondeterminism''. By benign nonderterminism, the
authors of \cite{Thirunarayan2001} mean that the only moment where the
state of a \vhdl{} design can not be decided in a deterministic way is
when the processes are in the middle of the execution of their
statement body. However, the state of a \vhdl{} design at that moment
is of no interest; it corresponds to nothing regarding the concrete
functioning of a hardware circuit. Also, two different processes can
never be writing to the same signal at the same time. If such a design
happens, this is a case of \textit{multiply-driven} signal, which is
utterly forbidden. So, there can be no nondeterminism, regarding the
value of a signal, coming from the concurrent execution of two
processes (at least when shared variables are not involved).

% However, we are not interested in
% dealing with constructs so advanced as shared variables, therefore,
% this work is not really relevant to us.

\paragraph{Translational semantics}
Another kind of semantics, called ``translational'', formalizes the
\vhdl{} language semantics by translating a \vhdl{} design into
another formal model. Thus, the semantics of \vhdl{} is modeled by the
translation and the formal semantics of the target model. The target
model has the ability to model concurrency, which is one of the
specificity of \vhdl{}. Moreover, target models are chosen regarding
the tools they provide for analysis, and thus, a translational
semantics for \vhdl{} is often related to model checking
considerations.

In \cite{Reetz1995}, the author expresses the formal semantics of
\vhdl{} by translating a \vhdl{} design into a corresponding
\emph{flowgraph}. All \vhdl{} constructs, ranging from sequential
statements to concurrent processes, are expressed with individual
flowgraphs that are then composed together through their
interfaces. The simulation cycle of \vhdl{} is also encoded by means
of connected flow graphs: one for the ``execution part'' of the
semantics, that is, all processes run until suspension, and one for
the update part (i.e, the kernel process). Flowgraphs come with a
large amount of tools for analysis, and this translational semantics
is involved in the setting of a framework to reason about \vhdl{}
programs using multiple technics (automatic theorem proving, model
checking\dots). All these technics rely on the flowgraph formalism.

In \cite{Dohmen1995}, the author introduces a translational semantics
for \vhdl{} based on deterministic finite-state automatons. Again, the
reason for using such automatons lies in the existence of many
analysis tools. Moreover, forcing the generation of deterministic
automatons improves the time execution of model-checking technics.
The translation is performed on an elaborated \vhdl{} design; a data
space stores the values of signals and variables, and automatons
represent the control-flow of \vhdl{} statements. Each \vhdl{} statement is
associated to a specific automaton; sequence of statements are
achieved by automaton composition. The simulation kernel is also
represented by a specific automaton. Processes are composed together
with respect to synchronization states, i.e. states that permit to
pass the control from one process to another, therefore achieving
determinism in the control flow of the overall automaton.

In \cite{Olcoz1995}, the author presents a translation from \vhdl{} to
Coloured Petri Nets (CPNs) thus giving a formal semantics to the
\vhdl{} constructs. The author approach to the \vhdl{} semantics is a
strict translation of the ``event-based'' \vhdl{} simulator by means
of Petri nets.  The author translates \vhdl{} execution models (sea of
processes) into CPNs, and also translates the kernel process into a
CPN. The kernel process has previously been expressed as a \vhdl{}
process so that the translation into CPN is similar to the translation
of other processes.  Signals are not represented in the subnets,
instead, three shared variables depict the signal states: one variable
for the driving, one for the effective and one for the current value
of a given signal (see \cite[p.167]{VHDL2000} for the details on the
values associated with signals during the simulation).  Color domains
of places in the subnets represent the different types of \vhdl{}
domains.  Variables are represented by tokens.  Values in drivers are
represented by sequences of transactions (equivalent to waveforms);
the author defines a set of functions that are convenient to handle
sequences of transactions.  Sequential statements are partitioned into
two kinds: control flow (if, loop, case\dots) and notation (operations
on signals and variables) nets.  Processes subnets are made by the
fusing of each sequential statements in the process body. There is a
special \emph{Resume} place that can be set by the kernel process to
resume the activity of a process.  Concurrency is not discussed here,
as the Petri net models are inherently concurrent models.  The kernel
process is a broad CPN having some of its places interfaced with the
process subnets.  The decoloration of the Petri net enables the
analysis of the model and the detection of dead-locks.

In \cite{Deharbe1995}, the author gives a formal semantics to \vhdl{}
by transforming a \vhdl{} design into an abstract machine,
i.e. defined by a set of inputs, outputs, states and transition
function over states and outputs. The author is interested in the
verification of properties over \vhdl{} designs (temporal properties)
or to prove equivalence between designs (bisimulation).  To operate
this transformation, only a subset of \vhdl{} is considered, otherwise
a finite-state representation is not reachable.  The covered \vhdl{}
subset consists of objects with finite types, and no quantitative
timing constructs (no after clause in signal
assignments).  % The author claims that a \vhdl{}
% design is implemented by an abstract machine if they have the same
% observational behavior, i.e. given the same inputs they yield the
% same outputs.
The transformation generates a decision diagram (i.e. a control flow
graph) and a state space for each process defined in the design's
behavior. The decision diagram encodes the transition function over
states and outputs.  Process statements are composed with a special
composition operator to obtain a global abstract machine. Moreover,
the article lays out a method to transform a block statement into an
abstract machine. The initiative is to be noticed as there are only a
few papers, dealing with the formalization of the \vhdl{} semantics, that
are interested in such hierarchical constructs as block or component
instantiation statements. The article concludes with an expression of
the space of complexity entailed by the transformation of a \vhdl{}
design into an abstract machine.

Although the translational semantics described above meet most of the
qualifying criterions in relation to our needs, we are not especially
interested in implementing one of these. The main reason being that it
would require to implement the transformation from the abstract
\vhdl{} syntax to the target model, in addition to the implementation
of the semantics of the target model.

Table~\ref{tab:sum-vhdl-sem} summarizes the analysis of the \vhdl{}
semantics encountered during our literature
review. Table~\ref{tab:sum-vhdl-sem} compares the different \vhdl{}
semantics in relation to our qualifying criterions (see
Section~\ref{sec:qualifying-criterions}).

\newpage

\newcommand{\YO}{\cellcolor{YellowOrange}}
\newcommand{\Gre}{\cellcolor{Green}}
\newcommand{\LGr}{\cellcolor{lightgray}}
\newcommand{\R}{\cellcolor{Red}}

\thispagestyle{empty}

\begin{landscape}

\begin{table}[H]
  \resizebox{1.4\textwidth}{!}{%
    \renewcommand{\arraystretch}{1.5}
    \begin{tabular}{ll|c|c|c|c|c|c|c|c|c|c|c|c|c|}

      % TABLE HEADER
      \cline{3-14}
      & &  &  &  &  &  &  &  &  &  &  &  &  \\
      &
      & \begin{tabular}[t]{@{}c@{}}\rotatebox[origin=r]{270}{Fuchs and Mendler} \\ \cite{Fuchs1995} \end{tabular}  
      & \begin{tabular}[t]{@{}c@{}}\rotatebox[origin=r]{270}{Breuer et al.} \\ \cite{Breuer1995a} \end{tabular}
      & \begin{tabular}[t]{@{}c@{}}\rotatebox[origin=r]{270}{Pandey et al.} \\ \cite{Pandey1999} \\ \end{tabular}
      & \begin{tabular}[t]{@{}c@{}}\rotatebox[origin=r]{270}{Borrione and Salem } \\ \cite{Borrione1995} \\\end{tabular}
      & \begin{tabular}[t]{@{}c@{}}\rotatebox[origin=r]{270}{Breuer et al. } \\ \cite{Breuer1995} \\ \end{tabular}
      & \begin{tabular}[t]{@{}c@{}}\rotatebox[origin=r]{270}{Börger et al. } \\ \cite{Borger1995} \\ \end{tabular}
      & \begin{tabular}[t]{@{}c@{}}\rotatebox[origin=r]{270}{Van Tassel} \\ \cite{VanTassel1995} \\ \end{tabular}
      & \begin{tabular}[t]{@{}c@{}}\rotatebox[origin=r]{270}{Goossens} \\ \cite{Goossens1995} \\ \end{tabular}
      & \begin{tabular}[t]{@{}c@{}}\rotatebox[origin=r]{270}{Reetz and Kropf} \\ \cite{Reetz1995} \\ \end{tabular}
      & \begin{tabular}[t]{@{}c@{}}\rotatebox[origin=r]{270}{Döhmen and Herrmann} \\ \cite{Dohmen1995} \\ \end{tabular}
      & \begin{tabular}[t]{@{}c@{}}\rotatebox[origin=r]{270}{Olcoz} \\ \cite{Olcoz1995} \\ \end{tabular}
      & \begin{tabular}[t]{@{}c@{}}\rotatebox[origin=r]{270}{Déharbe and Borrione} \\ \cite{Deharbe1995} \\ \end{tabular}
      \\

      % SEMANTICS DESCRIPTION

      
      % kind of semantics (D, O, T).
      \hline
      \multicolumn{1}{|c|}{}
      & Kind & \YO D & \YO D, A & \YO D & \YO D & \Gre O & \Gre O & \Gre O & \Gre O & \YO T & \YO T & \YO T & \YO T \\
      
      % purpose of the formalization.
      \cline{2-14}
      \multicolumn{1}{|c|}{\multirow{-2}{*}{\begin{tabular}[c]{@{}c@{}}Semantics\\ Description\end{tabular}}}
      & Purpose & \LGr AR, ATP & \LGr AR & \LGr AR & \LGr AR & \Gre SS & \Gre SS & \Gre SS, ITP & \Gre SS, MC & \Gre ATP, MC, ITP
                                   & \Gre MC, ITP & \LGr MC & \LGr MC \\

      % QUALIFYING CRITERIONS
      
      \hline
      
      % component instantiation covered?
      \multicolumn{1}{|l|}{\multirow{3}{*}{\begin{tabular}[c]{@{}l@{}}Qualifying\\ Criterions\end{tabular}}}
      & \begin{tabular}[c]{@{}l@{}}Component\\ Instantiation\end{tabular}
      & \YO T & \YO T & \YO T & \Gre N & \YO T & \YO T & \YO T & \YO T & \YO T & \YO T & \YO T & \Gre N \\

      % Synchronism expressibility status
      \cline{2-14} 
      \multicolumn{1}{|l|}{} & Synchronism & \R NE & \R NE & \R NE & \Gre Ex & \YO E & \YO E & \YO E & \YO E & \YO E & \YO E & \YO E & \R NE \\

      % Elaboration is formalized?
      \cline{2-14} 
      \multicolumn{1}{|l|}{}
      & Elaboration &\YO $\times$ &\YO $\times$ &\YO $\times$ & \Gre \checkmark &\YO $\times$ &\YO $\times$ & \Gre \checkmark &\YO $\times$ &\YO $\times$ &\YO $\times$ &\YO $\times$ & \Gre \checkmark \\ 
      
      % EXTRA. INFORMATION.

      \hline

      \multicolumn{1}{|c|}{\multirow{2}{*}{\begin{tabular}{@{}l@{}}Extra.\\ Informations.\end{tabular}}}
      
      % implementation technology
      & \begin{tabular}[c]{@{}l@{}} Impl.\\ Technology\end{tabular}
      & \LGr Focus \cite{Dederichs1993} & \LGr Prolog \cite{Colmerauer1990} & \LGr PVS \cite{Owre1994} & \LGr ? & \LGr Gofer \cite{Jones1994}
                    & \LGr ? & \LGr HOL \cite{Hutton1994} & \LGr ? & \LGr HOL \cite{Hutton1994} & \LGr ? & \LGr ? & \LGr ? \\
      
      % particular models or data types 
      \cline{2-14} 
      \multicolumn{1}{|c|}{}
      & \begin{tabular}[!htbp]{@{}l@{}}Particular\\ Model or\\ Data Types\end{tabular}
      & \LGr \begin{tabular}[c]{@{}c@{}}Stream \\ Processing \end{tabular}
      & \LGr No & \LGr \begin{tabular}[c]{@{}c@{}}Interval\\ Temporal \\ Logic \end{tabular}
      & \LGr No & \LGr \begin{tabular}[c]{@{}c@{}}Stream \\ Processing \end{tabular}
      & \LGr \begin{tabular}[c]{@{}c@{}}Evolving \\ Algebra \\ Machines \end{tabular}
      & \LGr \begin{tabular}[c]{@{}c@{}}Natural \\ Semantics \\ (big-step) \end{tabular}
      & \LGr \begin{tabular}[c]{@{}c@{}}Structural \\ Semantics \\ (small-step) \end{tabular}
      & \LGr \begin{tabular}[c]{@{}l@{}}Flow \\ Graphs \end{tabular}
      & \LGr \begin{tabular}[c]{@{}c@{}}Finite-State \\ Automatons \end{tabular}
      & \LGr \begin{tabular}[c]{@{}c@{}}Colored \\ Petri \\ Nets \end{tabular}
      & \LGr \begin{tabular}[c]{@{}c@{}}Abstract \\ Machines \\ and \\ Decision \\ Diagrams \end{tabular} \\
      \hline
    \end{tabular}%
  }
  \caption{A comparative summary on \vhdl{} formal semantics.}
  \label{tab:sum-vhdl-sem}

  \begin{itemize}[label=-]
    \fontsize{10}{12}\selectfont
  \item Kind : D (Denotational) - A (Axiomatic) - O (Operational) - T (Translational).
  \item Purpose : AR (Abstract Reasoning) - ATP (Automatic Theorem
    Proving) - SS (Simulator Specification) - ITP (Interactive Theorem
    Proving) - MC (Model Checking).
  \item Component Instantiation : T (statement is \emph{Transformed}
    into equivalent processes) - N (statement is \emph{Natively} taken into account in the semantics).
  \item Synchronism : E (Expressible within the semantics) - NE (Not
    Expressible within the semantics) - Ex (Explicitly built in the
    semantics).
  \end{itemize}
\end{table}

\end{landscape}

\newpage

To summarize, we are interested in a semantics built for the purpose
of interactive theorem proving (ideally, with an existing
implementation in the \coq{} proof assistant). Most important, the
formal semantics must be able to deal with the expression of
synchronous designs, that is, designs synchronized with a clock
signal. Therefore, a synchronous simulation cycle must be at least
expressible within the semantics. Moreover, the semantics must handle
component instantiation statements as they are, that is, without
transforming them into equivalent processes. As a bonus, the semantics
should formalize the elaboration part of \vhdl{} semantics.

In Table~\ref{tab:sum-vhdl-sem}, cells are colored in green when the
cell's content foster the adoption of the semantics, in yellow when
the content does not go towards the adoption of the semantics but is
not disqualifying, and red when the content is a disqualifying
criterion. Regarding the semantics adoption, cells are labelled in
light grey when their content is neutral. Now comparing the entries of
Table~\ref{tab:sum-vhdl-sem} with the expression of our needs, we can
discard the semantics with a cell labelled in red, that is, most of
the denotational semantics; moreover, all translational semantics are
disqualified for the previously mentioned reasons. The candidate
semantics are the operational semantics, plus the denotational
semantics by Borrione and Salem \cite{Borrione1995}, the only
semantics that formalizes an explicitly synchronous simulation
cycle. The semantics that is the most likely to be adopted is the
Borrione and Salem's semantics. However, we prefer an operational
setting for our semantics. To lower down the complexity of proofs, we
really need a semantics that builds the synchronism into its
simulation cycle, therefore putting aside all the intricacies of the
full-blown \vhdl{} simulation cycle. Moreover, the big-step style for
an operational semantics is more relevant to us; as stated before, we
are not interested in the intermediary states of computation that a
small-step style semantics considers.  Based on these observations, we
have decided to formalize our own \vhdl{} semantics inspired from the
semantics of Borrione and Salem's \cite{Borrione1995} and Van Tassel's
\cite{VanTassel1995}. The following sections are dedicated to the
presentation of the syntax and semantics of a subset of \vhdl{} that
we baptize \hvhdl{}. \hvhdl{} embeds the subset of \vhdl{} that we are
interested in when considering the \vhdl{} designs generated by the
\hilecop{} transformation.

%%% Local Variables:
%%% mode: latex
%%% TeX-master: "../../main"
%%% End:


\section{Abstract syntax of $\mathcal{H}$-VHDL}
\label{sec:abstractSyntax}
In this section, we describe the abstract syntax of \hvhdl{}, a subset
of VHDL covering all the constructs present in the programs generated
by the \hilecop{} transformation. Terminals of the language are
written in \texttt{typewriter} font, or are enclosed in simple quotes
for symbols with no typewriter representation. The $a^{*}$ denotes a
possibly empty repetition of the element $a$; the $a^{+}$ denotes a
non-empty repetition of $a$.

\subsection{Design declaration}
\label{subsec:design-decl}

Similarly to \cite{VanTassel1995}, we define the \emph{design}
construct in the $\mathcal{H}$-VHDL's abstract syntax which has no
equivalent in the concrete syntax of VHDL.

\begin{table*}[!h]
  \begin{tabular}{lcl}
    design & ::= & \texttt{design} id$_e$ id$_a$ gens ports sigs cs \\
    gens & ::= & gdecl* \\
    ports & ::= & pdecl* \\
    sigs & ::= & sdecl* \\
  \end{tabular}
\end{table*}

% \subsection{Generic constant, port and internal signal declaration.}
% \label{subsec:ent-decl}
\begin{table*}[!h]
\begin{tabular}{lcl}
  gdecl & ::= & \texttt{(}id\texttt{,} $\tau$\texttt{,} e\texttt{)} \\
  pdecl & ::= & \texttt{(}(\vhdle|in||\vhdle|out|)\texttt{,} id\texttt{,} $\tau$\texttt{)}\\
  sdecl & ::= & \texttt{(}id, $\tau$\texttt{)} \\
\end{tabular}
\end{table*}

In the above entry, $\mathrm{id}_e$ indicates the entity identifier
and $\mathrm{id}_a$ the architecture identifier of the declared
design. The $\mathrm{gens}$ entry corresponds to the generic clause,
i.e. the declaration list for the generic constants of the design. A
generic constant is declared via the \textrm{gdecl} entry; a generic
constant declaration is a triplet composed of an identifier, a type
indication and an expression denoting the generic constant's default
value. The ports entry holds the declaration of the input and output
ports of the design. A port declaration (i.e. the pdecl entry) is a
triplet composed of a port type, i.e. \texttt{in} or \texttt{out}, an
identifier, and a type indication. The sigs entry is the list
declaring the internal signals of the design. An internal signal
declaration entry (i.e. sdecl) is a couple composed of an identifier
and a type indication. The cs entry represents the concurrent
statements composing the behavior of the design.

\subsection{Concurrent statements}
\label{subsec:conc-stmt}

\begin{table*}[!h]
\begin{tabular}{lcl}
  cs & ::= & psstmt | cistmt | cs \texttt{||} cs | \texttt{null} \\
\end{tabular}
\end{table*}
In \hvhdl{}, two kinds of concurrent statements are available to
describe the behavior of a design: process statements, represented by
the psstmt entry, and component instantiation statements, represented
by the cistmt entry. Concurrent statements are composable through the
\texttt{||} operator. We add the \texttt{null} statement to the
\hvhdl{} abstract syntax to help represent empty behaviors.

\subsubsection{Process statement}

\begin{table*}[!h]
\begin{tabular}{lcl}
  psstmt & ::= & \vhdle|process| \texttt{(}id$_p$\texttt{,} sl\texttt{,} vars\texttt{,}  ss\texttt{)} \\
  sl & ::= & id$^{*}$ \\
  vars & ::= & vdecl* \\
  vdecl & ::= & \texttt{(}id\texttt{,} $\tau$\texttt{)} \\
\end{tabular}
\end{table*}

A process statement declares a sensitivity list, i.e. the sl entry,
which is a possibly empty set of signal identifiers. In order to be
well-formed, the signals of a sensitivity list must be either internal
signals or input ports of the design, i.e. . The process possibly declares a
set internal variables, i.e. the vars entry. A variable declaration
entry is a couple composed of a variable identifier and a type
indication. The ss entry represents the sequence of statements
composing the body of the process, i.e. the part that will be executed
during the simulation.

\subsubsection{Component instantiation statement}

The VHDL LRM defines two kinds of component instantiation statement:
the instantiation of a component instance \cite[p.139]{VHDL2000} and
the instantiation of a design entity \cite[p.141]{VHDL2000}. The
component instantiation statement used in the \hvhdl{} abstract syntax
corresponds to the instantiation of a design entity.

\begin{table*}[!h]
\begin{tabular}{lcl}
  cistmt & ::= & \vhdle|comp| \texttt{(}id$_c$\texttt{,} id$_e$\texttt{,} gmap\texttt{,} ipmap\texttt{,} opmap\texttt{)} \\
  gmap & ::= & assoc$_g^{*}$ \\
  ipmap & ::= & assoc$_{ip}^{*}$ \\
  opmap & ::= & assoc$_{op}^{*}$\\
  assoc$_g$ & ::= & \texttt{(}id\texttt{,}e\texttt{)} \\
  assoc$_{ip}$ & ::= & \texttt{(}name\texttt{,}e\texttt{)} \\
  assoc$_{op}$ & ::= & \texttt{(}id\texttt{,}(name|\vhdle|open|)\texttt{)}|\texttt{(}id\texttt{(}e\texttt{)}\texttt{,}name\texttt{)}\\
\end{tabular}
\end{table*}

In the cistmt entry, the identifier $\mathrm{id}_c$ represents the
name of component instance. Identifier $\mathrm{id}_e$ points out the
name of the design, i.e. the entity identifier, being instantiated
here. The gmap entry describes the list of associations between
generic constant identifiers and expressions. The ipmap entry is the
list of associations between input port identifiers (or indexed
identifiers) and expressions. The opmap entry is the list of
associations between output port identifiers (or indexed identifiers)
and signal names, or the \texttt{open} keyword. Associating the
\texttt{open} keyword with an output port identifier indicates that
the port is not connected. The left element of an association is
called the \emph{formal} part, and the right element of an association
is called the \emph{actual} part.

\subsection{Sequential statements}

\begin{table*}[!h]
\begin{tabular}{lcl}
  ss & ::= & name $\mathtt{\Leftarrow}$ e | name \texttt{:=} e | \texttt{if} \texttt{(}e\texttt{)} ss [ss] | \texttt{for (}id\texttt{,}e\texttt{,}e\texttt{)} ss \\
     & &  | \texttt{falling} ss | \texttt{rising} ss | \texttt{rst} ss ss' | ss\texttt{;} ss | \texttt{null} \\
\end{tabular}
\end{table*}

The ss entry defines the sequential statements that compose the body
of processes. The signal assignment statement is represented with the
$\Leftarrow$ operator; the variable assignment statement with the
$\mathtt{:=}$ operator. Also, we devise three control flow statements
that have no equivalent in the VHDL syntax: the \texttt{falling} block
statement, the \texttt{rising} block statement and the \texttt{rst}
block (or reset) block statement.  The \texttt{falling} statement
(resp. \texttt{rising} ss) declares a block of sequential statements
to be executed only at the falling edge (resp. rising edge) of the
clock signal (see Section~\ref{sec:clk-phases-rules}). Also, the
\texttt{rst} statement declares two blocks, the first one must be
executed during the initialization phase of the simulation; otherwise,
the second one is executed (see Section~\ref{sec:init-rules}).  These
invented constructs are equivalent to specific if-else statements that
are commonly used in the body of a synchronous process (see
Section~\ref{sec:abs-syntax-examples} for an example of transcription
of a specific if-else statement into one of these constructs).

\subsection{Expressions, names and types}
\label{sec:expr-names}

\begin{table}[H]
  \begin{tabular}{lcl}

    e & ::= & e \texttt{and} e | e \texttt{or} e | \texttt{not} e | e \texttt{=} e | e $\neq$ e \\
      & & | e \texttt{<} e | e \texttt{<=} e | e \texttt{>} e | e \texttt{>=} e | e \texttt{+} e | e \texttt{-} e \\
      & & | name | natural | boolean | \texttt{(}e$^{+}$\texttt{)} \\
      & & \\
    name & ::= & id | id\texttt{(} e \texttt{)} \\
    boolean & ::= & \texttt{true} | \texttt{false} \\
    $\tau$ & ::= & \texttt{boolean} | \texttt{natural} \texttt{(}e\texttt{,} e\texttt{)} |
                   \texttt{array} \texttt{(}$\tau$\texttt{,} e\texttt{,} e\texttt{)} \\
  \end{tabular}
\end{table}

The expression entry, i.e. e, declares a set of operators over Boolean
expressions, and natural numbers expressions.  The natural
non-terminal represents the set of natural numbers ($\mathbb{N}$). The
id non-terminal represents the set of identifiers, comparable to the
set of strings, or any infinitly enumerable set. In the following
sections, concrete identifiers will be written in \texttt{typewriter}
font, e.g. the \texttt{place} and \texttt{transition} design
identifiers.

The $\tau$ entry corresponds to the type indication associated with
the declaration of a generic constant, a port or an internal
signal. The considered types are the \emph{Boolean} type, the
constrained natural type, and the array type. The constrained natural
type, i.e. \texttt{natural}(e,e), defines a finite interval of natural
numbers; the left-most expression of the range constraint denotes the
lower bound of the interval, and the second one denotes the upper
bound of the interval. The array type indication,
i.e. \texttt{array}($\tau$, e, e), denotes a non-empty set of elements
of type $\tau$. The elements are indexed with respect to the specified
\emph{index} constraint. The left-most expression of the index
constraint denotes the starting index (possibly different from 0) and
the right-most expression denotes the final index.

% \subsection{Examples of \hvhdl{} abstract syntax}
% \label{sec:abs-syntax-examples}

% Listing~\ref{lst:t-design-decl-part-abss} gives the translation in
% abstract \hvhdl{} syntax of the declarative part of the transition
% design presented in Listings~\ref{lst:trans-design-entity} and
% \ref{lst:trans-design-arch} in concrete VHDL syntax.

% \begin{lstlisting}[language=VHDL,label={lst:t-design-decl-part-abss},
% caption={[The transition design's declarative parts in abstract
%   \hvhdl{} syntax.]The transition design's declarative parts in
%   abstract \hvhdl{}
%   syntax.},framexleftmargin=1.5em,xleftmargin=2em,numbers=left,
% numberstyle=\tiny\ttfamily]
% design "transition" "transition_architecture"
%   -- Generic clause
%   (("transition_type", natural(0, 2), 0), 
%    ("input_arcs_number", natural(0, NATMAX), 1), 
%    ("conditions_number", natural(0, NATMAX), 1),
%    ("maximal_time_counter", natural(0, NATMAX), 1))

%   -- Port clause
%   ((in, "input_conditions", array(boolean, 0, conditions_number$-$1)),
%    (in, "time_A_value", natural(0, maximal_time_counter)),
%    (in, "time_B_value", natural(0, maximal_time_counter)),
%    (in, "input_arcs_valid", array(boolean, 0, input_arcs_number$-$1)),
%    (in, "reinit_time", array(boolean, 0, input_arcs_number$-$1)),
%    (in, "priority_authorizations", array(boolean, 0, input_arcs_number$-$1)),
%    (out, "fired", boolean))

%   -- Internal signals
%   (("s_condition_combination", boolean),
%    ("s_enabled", boolean),
%    #\dots#
%    ("s_time_counter", natural(0, maximal_time_counter)))

%   -- Concurrent statements
%   cs
% \end{lstlisting}

% In Listing~\ref{lst:t-design-decl-part-abss}, the type indication of
% the ``\texttt{transition\_type}'' generic constant has been
% transformed from the \texttt{temporal_t} enumeration type to the
% \texttt{natural(0,2)} type. The \texttt{temporal_t} enumeration type,
% defined through the three values \texttt{\{NOT\_TEMPORAL,
%   TEMPORAL\_A\_A, TEMPORAL\_A\_B\}}, is naturally transformed into a
% natural range. This transformation is only valid because the
% \texttt{temporal_t} type is only defined and used in the transition
% design which has a static behavior. Also, the \texttt{std_logic} type
% defined in the VHDL Std Logic 1164 library \cite{STDLOGIC} is
% transformed into the Boolean type, and the \texttt{std_logic_vector}
% is transformed into the array of Booleans type. This transformation is
% valid because, in the transition and place designs and also in the
% generated designs, we are only referring to the values '0' and '1'
% among all the values enumerated by the \texttt{std_logic} type (see
% \cite[p. 2]{STDLOGIC}). Also, note that the \texttt{clock} and
% \texttt{reset_n} input ports declared in the port clause of the
% transition design are removed in the hvhdl{} version. In the \hvhdl{}
% abstract syntax, the if statements that were testing the value of the
% \texttt{clock} and the \texttt{reset_n} signals have been turned into
% specific sequential statements (i.e \texttt{rst}, \texttt{falling} and
% \texttt{rising} blocks). Thus, we don't need the \texttt{clock} and
% \texttt{reset_n} signals in the port declaration list anymore.

% Listing~\ref{lst:t-design-pss-abss} presents the translation in
% \hvhdl{} abstract syntax of the two of the processes presented in
% Listing~\ref{lst:trans-design-arch}.

% \begin{lstlisting}[language=VHDL,label={lst:t-design-pss-abss},
% caption={[Two processes in \hvhdl{} abstract syntax.]The \texttt{condition_evaluation} and \texttt{firable} processes in \hvhdl{} abstract syntax. The two processes are defined in the behavior of the transition design.},framexleftmargin=1.5em,xleftmargin=2em,numbers=left,
% numberstyle=\tiny\ttfamily]
%   process ("condition_evaluation", ("input_conditions"), 
%           (("v_internal_condition", boolean)),
%     v_internal_condition := true;
%     (for ("i", 0, conditions_number $-$ 1) 
%       (v_internal_condition := v_internal_conditions and input_conditions(i)));
%     s_condition_combination <= v_internal_condition
%   ) ||
%   process ("firable", ("clk"), $\emptyset$,
%     rst (s_firable <= false)
%         (falling (s_firable <= s_firing_condition))
%   )
% \end{lstlisting}

% In Listing~\ref{lst:t-design-pss-abss}, inner blocks of sequential
% statements are enclosed between parentheses to ease the reading,
% albeit parentheses are not part of the \hvhdl{} syntax. The body of
% the \texttt{firable} process gives an example of the use of a
% \texttt{rst} block and a \texttt{falling} block. One can see, from the
% comparison of Listing~\ref{lst:trans-design-arch} and
% \ref{lst:t-design-pss-abss}, that a \texttt{rst} ss ss' statement is
% the translation of the concrete VHDL statement \texttt{if reset_n =
%   '0' then ss else ss' endif;}. Comparing the same listings, a
% \texttt{falling} ss statement is the translation of the concrete VHDL
% statement \texttt{if falling_edge(clock) then ss end if;}.

% Listing~\ref{lst:t-design-ci-abss} shows the translation in \hvhdl{}
% abstract syntax of the component instantiation statement laid out in
% Listing~\ref{lst:trans-design-instantiation}.

% \begin{lstlisting}[language=VHDL,label={lst:t-design-ci-abss},
% caption={[A component instantiation statement in \hvhdl{} abstract
%   syntax (transition component instance).] An example of instantiation
%   of the transition design in \hvhdl{} abstract syntax. The transition
%   component instance is named
%   $id_t$.},framexleftmargin=1.5em,xleftmargin=2em,numbers=left,
% numberstyle=\tiny\ttfamily]
%   comp ($id_t$, "transition", 
%     -- Generic map
%     (transition_type => 0,
%      input_arcs_number => 1,
%      conditions_number => 1,
%      maximal_time_counter => 1),
%     -- Input port map
%     (time_A_value => 0,
%      time_B_value => 0,
%      input_conditions(0) => $id_0$,
%      input_arcs_valid(0) => $id_1$,
%      priority_authorizations(0) => true,
%      reinit_time(0) => $id_2$),
%     -- Output port map
%     (fired => $id_3$)
% \end{lstlisting}

% The statement of Listing~\ref{lst:t-design-ci-abss} is almost similar
% to the one of Listing~\ref{lst:trans-design-instantiation}. The
% \texttt{NOT_TEMPORAL} value associated to the \texttt{transition_type}
% constant is turned into 0 (remember the \texttt{temporal_t}
% enumeration type is transformed into a natural range type).

%%% Local Variables:
%%% mode: latex
%%% TeX-master: "../../main"
%%% End:


\section{Preliminary definitions}
\label{sec:sem-rules}
\subsection{Semantics Domains}

Let $id$ denote the set of identifiers in the semantical domain. We
write $prefix\mhyphen{}id$ to denote arbitrary subsets of the $id$
set. The $type$ and $value$ semantical types are defined as follows:

\begin{table}[H]
  \begin{tabular}{lcl}
    $type$ & ::= & \texttt{bool} | \texttt{nat(}$nat$,$nat$\texttt{)} | \texttt{array} \texttt{(}$type$,$nat$,$nat$\texttt{)} \\
           & & \\
    $value$ & ::= & $bool$ | $nat$ | $array$ \\
    $bool$ & ::= & '$\mathbf{\top}$' | '$\mathbf{\bot}$' \\
    $nat$ & ::= & \texttt{0} | \texttt{1} | \dots | \texttt{NATMAX} \\
    $array$ & ::= & \texttt{(}$value^{+}$\texttt{)}\\

  \end{tabular}
  \caption{The $type$ and $value$ semantical types.}
  \label{tab:type-value}
\end{table}

In Table~\ref{tab:type-value}, the $type$ type is in any way similar
to the $\tau$ entry of the abstract syntax, however, all constraint
bounds in the \texttt{nat} and \texttt{array} types have been
evaluated to natural numbers. $\mathtt{NATMAX}$ denotes the maximum
value for a natural number.  The $\mathtt{NATMAX}$ value depends on
the implementation of the VHDL language; $\mathtt{NATMAX}$ must at
least be equal to $2^{31}-1$. Note that the $array$ value contains at
least one value as an array's index range contains at least one index
(that is index 0).

\begin{notation}[Partial functions]
  Here, we present our notations pertaining to partial functions:
  \begin{itemize}[label=-]
  \item The $\nrightarrow$ arrow denotes a partial function.
  \item The $\rightarrow$ denotes an application (i.e, a total
    function).
  \item For all $f\in{}A\nrightarrow{}B$, $x\in{}f$ states that x is in
    the domain of function $f$.
  \item For all $f\in{}A\nrightarrow{}B$ and $g\in{}A\nrightarrow{}C$,
    $f\subseteq{}g$ states that the domain of $f$ is a subset of the
    domain of $g$.
  \item For all $X\subset{}A$ and $f\in{}A\nrightarrow{}B$,
    $X\subseteq{}f$ states that $X$ is a subset of the domain of $f$.
  \end{itemize}
\end{notation}

Now, let us define the structure of an elaborated design which is a
structure bound to a given \hvhdl{} design and to a design store, i.e
a global environment mapping identifiers to \hvhdl{} designs. Only the
designs referenced into the global design store can be instantiated as
subcomponents of a given design. The elaborated design structure is
used in the expression of the elaboration relation presented in
Section~\nameref{sec:elab-rules}, as well as in the expression of the
simulation rules. Let $ElDesign(d,\mathcal{D})$ be the set of the
elaborated designs for a given \hvhdl{} design $d$ and a design store
$D$. An elaborated design is a composite environment built out of
multiple sub-environments.  Each sub-environment is a table,
represented as a partial function, mapping identifiers of a certain
category of constructs (e.g, input port identifiers) to their
declaration information (e.g, type indication for input ports). We
represent an elaborated design as a record where the fields are the
sub-environments. An elaborated design is defined as follows:

\begin{definition}[Elaborated Design] For a given \hvhdl{} design
  $d\in{}design$ s.t. $d=$ \vhdle|design| \textit{$id_{ent}$
    $id_{arch}$ gens ports sigs behavior} and a given design store
  $\mathcal{D}\in{}entity\mhyphen{}id\nrightarrow{}design$, an
  elaborated design $\Delta\in{}ElDesign(d,\mathcal{D})$ is a record
  ${<}Gens, Ins, Outs, Sigs, Ps, Comps{>}$ where:
  \begin{itemize}[label=$-$]
  \item
    $Gens\in{}generic\mhyphen{}id\nrightarrow{}(type\times{}value)$
    where
    $generic\mhyphen{}id=\{id~\vert~(id,\tau,e)\in{}gens\}$,
    is the partial function yielding the type and the value of generic
    constants.
  \item $Ins\in{}input\mhyphen{}id\nrightarrow{}type$ where
    $input\mhyphen{}id=\{id~\vert~(\mathtt{in},id,\tau)\in{}ports\}$,
    is the partial function yielding the type of input ports.
  \item $Outs\in{}output\mhyphen{}id\nrightarrow{}type$ where
    $output\mhyphen{}id=\{id~\vert~(\mathtt{out},id,\tau)\in{}ports\}$,
    the partial function yielding the type of output ports.
  \item
    $Sigs\in{}declared\mhyphen{}signal\mhyphen{}id\nrightarrow{}type$
    where
    $declared\mhyphen{}signal\mhyphen{}id=\{id~\vert~(id,\tau)\in{}sigs\}$,
    the partial function yelding the type of declared signals.
  \item
    $Ps\in{}process\mhyphen{}id\nrightarrow{}(variable\mhyphen{}id(id_p)\nrightarrow{}(type\times{}value))$
    where\\
    $process\mhyphen{}id=\{id_p~\vert~\mathtt{process}(id_p,sl,vars,ss)\in{}behavior\}$,
    the partial function associating processes to their local
    environment. Local environments are functions mapping local
    variable identifiers to their corresponding type and
    value. Therefore, each set of local variable identifiers
    $variable\mhyphen{}id(id_p)$ depends on the process identifier
    (represented by $id_p$) passed as the first argument of the $Ps$
    function.
  \item
    $Comps\in{}component\mhyphen{}id\nrightarrow{}ElDesign(d_e,\mathcal{D})$,
    where\\
    $component\mhyphen{}id=\{id_c~\vert~\mathtt{comp}(id_c,id_e,gm,ipm,opm)\in{}behavior\}$,
    the partial function mapping component instance ids to their
    elaborated design version. The set $ElDesign(d_e,\mathcal{D})$
    depends on the design $d_e$ from which the component identifier
    $id_c$, passed as the first argument of the $Comps$ function, is
    an instance. Design $d_e$ is retrieved from the design store
    $\mathcal{D}$ s.t. $d_e=\mathcal{D}(id_e)$.
  \end{itemize}
\end{definition}

We assume that there are no overlapping between the identifiers of the
sub-environments (i.e, an identifier belongs to at most one
sub-environment). We note $\Delta(x)$ to denote the value returned for
identifier $x$, where $x$ is looked up in the appropriate field of
$\Delta$. We note $x\in\Delta$ to state that identifier $x$ is defined
in one of $\Delta$'s fields. We note $\Delta(x)\leftarrow{}v$ the
overriding of the value associated to identifier $x$ with value $v$ in
the appropriate field of $\Delta$, $\Delta\cup{}(x,v)$ to note the
addition the mapping from identifier $x$ to value $v$ in the
appropriate field of $\Delta$, that assuming $x\notin\Delta$. We note
$x\in\mathcal{F}(\Delta)$, where $\mathcal{F}$ is a field of $\Delta$,
when more precision is needed regarding the lookup of identifier $x$
in the record $\Delta$.

Let $\Sigma(\Delta)$ be the set of design states for a given
elaborated design $\Delta$.  A design state of $\Delta$ is defined as
follows:

\begin{definition}[Design state]
  A design state $\sigma\in\Sigma(\Delta)$, for a given design
  $d\in{}design$, a given design store $\mathcal{D}$ and an elaborated
  design $\Delta\in{}ElDesign(d,\mathcal{D})$, is a record
  ${<}\mathcal{S},\mathcal{C},\mathcal{E}{>}$ where:
  \begin{itemize}[label=$-$]
  \item $\mathcal{S}\in{}signal\mhyphen{}id\rightarrow{}value$, the
    partial function yielding the current values of the design's
    signals (ports and declared signals).
  \item
    $\mathcal{C}\in{}component\mhyphen{}id\rightarrow{}\Sigma(\Delta_c)$,
    the partial function yielding the current state of design's
    component instances, where $\Delta_c=\Delta(id_c)$ and
    $id_c\in{}component\mhyphen{}id$ is the component identifier
    passed to function $\mathcal{C}$.
  \item
    $\mathcal{E}\subseteq{}signal\mhyphen{}id\cup{}component\mhyphen{}id$,
    the set of signal and component instance ids that generated an
    event at the current simulation cycle.
  \end{itemize}
\end{definition}

The $signal\mhyphen{}id$ subset is the disjoint union of
$input\mhyphen{}id$, $output\mhyphen{}id$ and
$declared\mhyphen{}signal\mhyphen{}id$. We use $\sigma(id)$ to denote
the value associated to an identifier in the signal store
$\mathcal{S}$ or in the component store $\mathcal{C}$
fields. $id\in_\mathcal{E}\sigma$ states that an identifier belongs to
the event set $\mathcal{E}$, whereas $id\in\sigma$ states that an
identifier is defined in either the signal store $\mathcal{S}$ or the
component store $\mathcal{C}$ fields. $\sigma\cup_\mathcal{E}\{id\}$
denotes a new design state, in all points similar to $\sigma$ but
enriched with the identifier $id$ in its events set.

\begin{notation}[No events design state]
  For a given \hvhdl{} design $d$, a design store $\mathcal{D}$, and
  an elaborated design $\Delta\in{}ElDesign(d,\mathcal{D})$, the
  function $NoEv\in\Sigma(\Delta)\rightarrow\Sigma(\Delta)$ returns a
  design state similar to the one passed in parameter only with an
  empty set of events. I.e, for all design state
  $\sigma\in\Sigma(\Delta)$
  s.t. $\sigma={<}\mathcal{S},\mathcal{C},\mathcal{E}{>}$,
  $NoEv(\sigma)={<}\mathcal{S},\mathcal{C},\emptyset{>}$.
\end{notation}

% \subsection{Design evaluation workflow.}
% \label{subsec:elab-sim-wflow}

% The workflow to evaluate a top-level VHDL design comprises two phases
% with some pre-requisites.

% \begin{figure}[H]
%   \centering
%   \begin{tikzpicture}

%     \node(elab) {\textcircled{1} Elaboration phase};
%     \node(sim) at ($(elab.east)+(2.5,0)$) {\textcircled{2} Simulation phase};

%     \draw (elab.east) edge[->] ($(sim.west)-(.1,0)$);
%   \end{tikzpicture}
%   \caption{Workflow for \emph{top-level} design evaluation.}
%   \label{fig:toplevel-eval}
% \end{figure}

% \begin{itemize}[label=-]
% \item Pre-requisites to the top-level design evaluation are as
%   follows:

%   \begin{enumerate}
%   \item All types and subtypes defined in VHDL standard library packages
%     are present in the environment, under the $\Gamma$ parameter.
%   \item All types and subtypes defined in user-defined packages are
%     present in the environment, under the $\Gamma$ parameter.
%   \item All globally defined constants will be considered as known
%     values in the elaboration/simulation rules.
%   \item No shared variables or signals are globally declared.
%   \end{enumerate}
  
% \item The \emph{elaboration} phase:

%   \begin{enumerate}
%   \item Builds a functor for each design of the project, including for
%     the top-level design.
%   \item Type-check designs' declarative parts.
%   \item Type-checks designs' behavioral parts.
%   \end{enumerate}

% \item The \emph{simulation} phase:
%   \begin{enumerate}
%   \item Instantiates the top-level design thanks to the $I_p$ function
%     yelding the primary input values.
%   \item Simulates the top-level design until $T_c$, the count of
%     simulation cycles, equals $0$.
%   \end{enumerate}
% \end{itemize}

% \subsection{Notation of rules.}
% \todo[inline]{Add examples of type-checking rules.}

% Semantics rules are expressed by judgments that take the following form: 

% \begin{figure}[H]

  
%   \centering
%   \begin{prooftree}
%     \hypo{env &\vdash \mathtt{a}\xrightarrow{rel}i}
%   \end{prooftree}
% \end{figure}

% Here is how to interpret the above judgment: some syntactic element
% \texttt{a} of the VHDL language is interpreted as \textit{i}
% with respect to the relation \textit{rel}, in the environment
% \textit{env}.

% Moreover, semantics rules can be with or without premises.

% \begin{figure}[H]
%   {\fontsize{8}{11}\selectfont \textsc{Axiom}}
  
%   \begin{prooftree}
%     \infer0{
%       env \vdash \mathtt{a}{}\xrightarrow{rel}i
%     }
%   \end{prooftree}
% \end{figure}

% Semantics rules without premises, as the one shown above, are called
% axioms.

% \begin{figure}[H]
%   {\fontsize{8}{11}\selectfont \textsc{WithPremises}}
  
%   \begin{prooftree}
%     \hypo{env &\vdash \mathtt{a}\xrightarrow{rel}i}
%     \hypo{env &\vdash \mathtt{b}\xrightarrow{rel}j}
%     \infer2{
%       env \vdash S\mathtt{(a,b)}\xrightarrow{rel}I(i,j)
%     }
%   \end{prooftree}
% \end{figure}

% The above example depicts a rule with premises. The premises are
% placed above the judgment bar, and the conclusion appears below.  The
% above rule expresses that if the syntactic element \texttt{a} is
% interpreted as \textit{i} and the syntactic element \texttt{b} is
% interpreted as \textit{j} both with respect to relation \textit{rel},
% then the syntactic element $S\mathtt{(a,b)}$ is interpreted as
% $I\mathtt{(i,j)}$ with respect to relation \textit{rel}. Here, $S$ is
% a relation appearing in the syntax of the VHDL language and bounding
% syntactic elements \texttt{a} and \texttt{b}, whereas $I$ is a
% relation belonging to the interpretation world, bounding elements $i$
% and $j$ which are both interpretation results.

%%% Local Variables:
%%% mode: latex
%%% TeX-master: "../../main"
%%% End:


\section{Elaboration rules}
\label{sec:elab-rules}
The goal of the elaboration phase is to build an elaborated design
$\Delta$ along with a \emph{default} state $\sigma_e$, out of a
\hvhdl{} design $d$. The elaboration relation also covers
type-checking operations for the declarative and behavioral parts of
design $d$.

\subsection{Design elaboration.}
\label{subsubsec:design-elab}

\begin{definition}[Design store]
  \label{def:design-store}
  $\mathcal{D}\in{}entity\mhyphen{}id\nrightarrow{}\mathrm{design}$ is
  the partial function mapping entity ids to their corresponding
  representation in abstract syntax. As a prerequisite to the
  elaboration of \hilecop{}-generated designs (i.e, resulting from the
  transformation of a SITPN into an \hvhdl{} design), a particular
  design store $\mathcal{D}_\mathcal{H}$ is defined. Design store
  $\mathcal{D}_\mathcal{H}$ holds the description of the
  \texttt{place} and \texttt{transition} designs which full definition
  in abstract syntax are given in appendices \ref{sec:place-abs-synt}
  and \ref{sec:trans-abs-synt}.
\end{definition}

At the beginning of the elaboration phase, a function
$\mathcal{M}_g\in{}generic\mhyphen{}id\nrightarrow{}value$ mapping the
top-level design's generic constants to values is passed as an element
of the environment. The $\mathcal{M}_g$ function is refered to as the
\emph{dimensioning} function.

\begin{table}[H]
  \centering
  \begin{tabular}{@{}l}
    {\fontsize{8}{11}\selectfont\textsc{DesignElab}} \\    
    {\begin{prooftree}

        % Elaborates gens, ports, sigs and cs.

        \hypo{\Delta_{\emptyset},\mathcal{M}_g&\vdash{}\mathrm{gens}\xrightarrow{egens}\Delta}
        \infer[no rule]1{\Delta,\sigma_\emptyset&\vdash{}
          \mathrm{ports}\xrightarrow{eports}\Delta',\sigma}
        \infer[no rule]1{\Delta',\sigma&\vdash{}\mathrm{sigs}\xrightarrow{esigs}\Delta'',\sigma'}
        \infer[no rule]1{\mathcal{D},\Delta'',\sigma'&\vdash
          \mathrm{cs}\xrightarrow{ebeh}\Delta''',\sigma''}
        
        % Conclusion
        
        \infer[template={\fontsize{10}{13}\selectfont\inserttext}]1
        {
          $\mathcal{D},\mathcal{M}_g\vdash$
          \vhdle|design| id$_e$ id$_a$ gens ports sigs cs
          $\xrightarrow{elab}$
          $\Delta''',\sigma''$
        }
      \end{prooftree}} \\
  \end{tabular}
\end{table}

where $\Delta_\emptyset$ denotes an empty elaborated design, that is
an elaborated design initialized with empty fields (empty tables). In
the same manner, $\sigma_\emptyset$ denotes an empty design state.
The effect of the $egens$, $eports$, $esigs$ and $ebeh$ that
respectively deal with the elaboration of the generic constants, the
ports, the architecture declarative part and the behavioral part of
the design, are explicited in the following sections.

\subsection{Generic clause elaboration.}
\label{subsubsec:gen-clause-elab}

\begin{premises}
  \begin{itemize}[label=-]
  \item $etype_g$ transforms a subtype indication, specifically
    attached to a generic constant declaration, into a $type$ instance
    and checks its well-formedness (see \ref{subsec:type-elab}).
  \item The $e$ relation evaluates expression $e$ to value $v$ (see
    \ref{subsec:expr-rules}).
  \item $SE_l$ states that an expression is $locally$ static (see
    \ref{subsubsec:loc-static}).
  \item $v\in_c{}T$ and $\mathcal{M}(\mathrm{id}_g)\in_c{}T$ checks
    that the default value and the value of yielded by the
    dimensioning function belongs to the type of the declared generic
    constant (see \ref{subsec:constr-satif-rel}).
  \end{itemize}  
\end{premises}

\begin{table}[H]
  \begin{tabular}{@{}l}
    {\fontsize{8}{11}\selectfont\textsc{GenElabDimen}} \\
    {\begin{prooftree}

        % Well-formed type and elaborates type.
        \hypo{\vdash\tau\xrightarrow{etype_g}T}
        
        % Evaluates expr. default.
        \hypo{\Delta_\emptyset,\sigma_\emptyset,\Lambda_\emptyset&\vdash
          \mathrm{e}\xrightarrow{e}v}

        % Checks static expr. default
        \hypo{&SE_l(\mathrm{e})}

        % Checks the type of the value yielded by the dimensioning fun.
        \hypo{v&\in_c{}T}
        \infer[no rule]1{\mathcal{M}(\mathrm{id}_g)&\in_c{}T}
        
        % Conclusion
        \infer[template={\inserttext}]4
        [{\renewcommand{\arraystretch}{1.5}
          \begin{tabular}{l}
            $\mathrm{id}_g\notin\Delta$ \\ 
            $\mathrm{id}_g\in\mathcal{M}$\\
          \end{tabular}
        }]
        {
          $\Delta,\mathcal{M}\vdash
          \mathtt{(}\mathrm{id}_g\mathtt{,}\tau\mathtt{,}\mathrm{e}\mathtt{)}
          \xrightarrow{egens}$
          $\Delta\cup{}(id_g,(T,\mathcal{M}(\mathrm{id}_g)))$
        }
      \end{prooftree}} \\
  \end{tabular}
\end{table}%

\begin{figure}[H]
  {\fontsize{8}{11}\selectfont\textsc{GenElabDefault}}
  
  \begin{prooftree}

    % Well-formed type and elaborates type.
    \hypo{\vdash\tau\xrightarrow{etype_g}T}

    % Evaluates default value.
    \hypo{
      \Delta_\emptyset,\sigma_\emptyset,\Lambda_\emptyset&
      \vdash{}\mathrm{e}\xrightarrow{e}v}

    % Checks static expr. default value.
    \hypo{&SE_l(\mathrm{e})}    

    % Well-typed default.
    \hypo{v\in_c{}T}
    
    % Conclusion
    \infer[template={\inserttext}]4
    [{\renewcommand{\arraystretch}{1.5}
      \begin{tabular}{@{}l}
        $\mathrm{id}_g\notin\Delta$ \\
        $\mathrm{id}_g\notin\mathcal{M}$\\
      \end{tabular}
    }]
    {
      $\Delta,\mathcal{M}\vdash
      \mathtt{(}\mathrm{id}_g\mathtt{,}\tau\mathtt{,}\mathrm{e}\mathtt{)}
      \xrightarrow{egens}$
      $\Delta\cup{}(id_g,(T,v))$
    }
  \end{prooftree}
\end{figure}

\begin{figure}[H]
  {\fontsize{8}{11}\selectfont\textsc{GenElabComp}}
    
  \begin{prooftree}

    % Evaluates GE1
    \hypo{\Delta,\mathcal{M}\vdash\mathrm{gdecl}\xrightarrow{egens}{}\Delta'}

    % Evaluates GE2.
    \hypo{\Delta',\mathcal{M}\vdash\mathrm{gens}\xrightarrow{egens}{}\Delta''}

    % Conclusion
    \infer[template={\inserttext}]2{
      $\Delta,\mathcal{M}\vdash
      \mathrm{gdecl}\mathtt{;}~\mathrm{gens}
      \xrightarrow{egens}\Delta''$
    }
  \end{prooftree}
\end{figure}

\subsection{Port clause elaboration.}
\label{subsubsec:port-clause-elab}

\begin{premises}
  The $defaultv$ relation yields the implicit \emph{default} value for
  a given type.
\end{premises}

% Elaborate IN port.

\begin{figure}[H]
  {\fontsize{8}{11}\selectfont\textsc{InPortElab}}
  
  \begin{prooftree}
        
    % Elaborates type.
    \hypo{\Delta\vdash\tau\xrightarrow{etype}T}

    % Retrieves leftmost value of type T.
    \hypo{\Delta\vdash{}T\xrightarrow{defaultv}v}
    
    % Conclusion
    \infer[template={\inserttext}]2
    [{
      \begin{tabular}{@{}l}
        $\mathrm{id}\notin\Delta$ \\
        $\mathrm{id}\notin\sigma$ \\
      \end{tabular}
    }]
    {
      $\Delta,\sigma\vdash
      ($\vhdle|in|$,\mathrm{id},\tau)$
      $\xrightarrow{eports}$
      $\Delta\cup{}(id,T),$
      $\sigma\cup(id,v)$
    }
  \end{prooftree}
\end{figure}

% Elaborate OUT port.

\begin{figure}[H]
  {\fontsize{8}{11}\selectfont\textsc{OutPortElab}}
  
  \begin{prooftree}

    % Elaborates type.
    \hypo{\Delta\vdash\tau\xrightarrow{etype}T}

    % Retrieves leftmost value of type T.
    \hypo{\Delta\vdash{}T\xrightarrow{defaultv}v}
    
    % Conclusion
    \infer[template={\inserttext}]2
    [{
      \begin{tabular}{@{}l}
        $\mathrm{id}\notin\Delta$ \\
        $\mathrm{id}\notin\sigma$ \\
      \end{tabular}
    }]
    {
      $\Delta,\sigma\vdash
      ($\vhdle|out|$,\mathrm{id},\tau)$
      $\xrightarrow{eports}$
      $\Delta\cup{}(id,T),$
      $\sigma\cup(id,v)$
    }
  \end{prooftree}
\end{figure}

% Inductive def. of port elab.

\begin{figure}[H]
  {\fontsize{8}{11}\selectfont\textsc{PortElabComp}}
  
  \begin{prooftree}

    % Evaluates PE1
    \hypo{\Delta,\sigma\vdash\mathrm{pdecl}\xrightarrow{eports}{}\Delta',\sigma'}

    % Evaluates PE2.
    \hypo{\Delta',\sigma'\vdash\mathrm{ports}\xrightarrow{eports}{}\Delta'',\sigma''}

    % Conclusion
    \infer[template={\inserttext}]2{
      $\Delta,\sigma\vdash
      \mathrm{pdecl}\mathtt{;}~\mathrm{ports}
      \xrightarrow{eports}\Delta'',\sigma''$
    }
  \end{prooftree}
\end{figure}

\subsection{Architecture declarative part elaboration.}
\label{subsubsec:arch-decl-part-elab}

% Signal elaboration rule.

\begin{figure}[H]
  {\fontsize{8}{11}\selectfont\textsc{SigElab}}
  
  \begin{prooftree}

    % Elaborates type.
    \hypo{\Delta\vdash\tau\xrightarrow{etype}T}
    
    % Retrieves leftmost value of type T.
    \hypo{\Delta\vdash{}T\xrightarrow{defaultv}v}
    
    % Conclusion
    \infer[template={\inserttext}]2
    [{
      \begin{tabular}{@{}l}
        $\mathrm{id}\notin\Delta$ \\
        $\mathrm{id}\notin\sigma$ \\
      \end{tabular}
    }]
    {
      $\Delta,\sigma\vdash$
      $(\mathrm{id},\tau)$
      $\xrightarrow{esigs}$
      $\Delta\cup{}(id,T),$
      $\sigma\cup(id,v)$
    }
  \end{prooftree}
\end{figure}


% Inductive elaboration rule for sequence of decl. entries.

\begin{figure}[H]
  {\fontsize{8}{11}\selectfont\textsc{SigElabComp}}
  
  \begin{prooftree}

    % Evaluates sdecl.
    \hypo{\Delta,\sigma\vdash\mathrm{sdecl}\xrightarrow{esigs}{}\Delta',\sigma'}

    % Evaluates sigs.
    \hypo{\Delta',\sigma'\vdash\mathrm{sigs}\xrightarrow{esigs}{}\Delta'',\sigma''}

    % Conclusion
    \infer[template={\inserttext}]2{
      $\Delta,\sigma\vdash
      \mathrm{sdecl}\mathtt{;}~\mathrm{sigs}
      \xrightarrow{esigs}\Delta'',\sigma''$
    }
  \end{prooftree}
\end{figure}

\subsection{Type indication elaboration.}
\label{subsec:type-elab}

The $etype$ relation checks the well-formedness of a type indication
$\tau$, and transforms it into a semantical $type$ (as defined in
Table~\ref{tab:type-value}). A subtype indication $\tau$ is well-formed
in the context $\Delta$ if $\tau$ denotes the \texttt{boolean} keyword
or the \texttt{nat} or \texttt{array} keywords with a
\emph{well-formed} constraint, and a well-formed element type in the
\texttt{array} case.

\begin{table}[H]
  \centering
  \begin{tabular}{@{}l}
    {\fontsize{8}{11}\selectfont\textsc{ETypeBool}} \\
    {\begin{prooftree}
        
        % Conclusion
        \infer[template={\inserttext}]0
        {
          $\Delta\vdash$
          $~\mathtt{boolean}$
          $\xrightarrow{etype}$
          $\mathtt{bool}$
        }
      \end{prooftree}} \\
  \end{tabular}
  \begin{tabular}{@{}l}
    {\fontsize{8}{11}\selectfont\textsc{ETypeNat}} \\
    {\begin{prooftree}
        
        % Well-formed constraint and evaluates constraint.
        \hypo{
          \Delta\vdash
          (\mathrm{e},\mathrm{e}')
          \xrightarrow{econstr}
          (v,v')
        }
        
        % Conclusion
        \infer[template={\inserttext}]1
        {
          $\Delta\vdash$
          $~\mathtt{natural}(\mathrm{e},\mathrm{e}')$
          $\xrightarrow{etype}$
          $\mathtt{nat}(v,v')$
        }
      \end{prooftree}} \\
  \end{tabular}
\end{table}

\begin{table}[H]
  \centering
  \begin{tabular}{@{}l}
    {\fontsize{8}{11}\selectfont\textsc{ETypeArray}} \\
    {\begin{prooftree}

        % Checks well-formednes and evaluates tau_ind.
        \hypo{\Delta\vdash\tau\xrightarrow{etype}T}

        % Checks well-formednes and evaluates index constraint.
        \hypo{\Delta\vdash(\mathrm{e},\mathrm{e'})\xrightarrow{econstr}(v,v')}
        
        % Conclusion
        \infer[template={\inserttext}]2
        {
          $\Delta\vdash
          ~\mathtt{array(}\tau,\mathrm{e},\mathrm{e'}\mathtt{)}$
          $\xrightarrow{etype}$
          $\mathtt{array(}T,v,v'\mathtt{)}$
        }
      \end{prooftree}} \\
  \end{tabular}
\end{table}

The $econstr$ relation checks that a constraint is well-formed and
evaluates the constraint bounds.  A constraint $c$ is well-formed in
the context $\Delta$ if:

\begin{itemize}[label=-]
\item its bounds are globally static expressions \cite[p.36]{VHDL00}
  of the \texttt{nat} type.
\item its lower bound value is inferior or equal to its upper bound
  value.
\end{itemize}

\begin{remark}[Type of constraints]
  As the VHDL language reference stays unclear about the type of range
  and index constraints \cite[p.33]{VHDL00}, we add the restriction
  that range and index constraints must have bounds of the
  \texttt{nat} type (i.e, value of type $nat$).
\end{remark}

\begin{premises}
  $SE_g$ states that an expression is \emph{globally} static (see
  \ref{subsubsec:glob-static}).
\end{premises}

\begin{table}[H]
  \centering
  \begin{tabular}{@{}l}
    {\fontsize{8}{11}\selectfont\textsc{EConstr}} \\
    {\begin{prooftree}
        
        
        % Globally static e and e'.
        \hypo{\Delta\vdash{}SE_g(\mathrm{e})}
        \infer[no rule]1{\Delta\vdash{}SE_g(\mathrm{e}')}
        
        % Evaluates e and e'
        \hypo{\Delta,\sigma_\emptyset,\Lambda_\emptyset&\vdash\mathrm{e}\xrightarrow{e}v}
        \infer[no rule]1{\Delta,\sigma_\emptyset,\Lambda_\emptyset&\vdash\mathrm{e}'\xrightarrow{e}v'}

        % Well-typed e and e'.
        \hypo{v&\in_c{}\mathtt{nat}(0,\mathtt{NATMAX})}
        \infer[no rule]1{v'&\in_c{}\mathtt{nat}(0,\mathtt{NATMAX})}
        
        % Conclusion
        \infer[template={\inserttext}]3
        [{
          \renewcommand{\arraystretch}{1.5}
          \begin{tabular}{@{}l}
            $v\leq{}v'$ \\
          \end{tabular}
        }]
        {
          $\Delta\vdash$
          $(\mathrm{e},\mathrm{e}')$
          $\xrightarrow{econstr}$
          $(v,v')$
        }
      \end{prooftree}} \\
  \end{tabular}
\end{table}

When considering a type indication in a generic constant declaration,
the definition of well-formedness differs slightly from the general
definition. A type indication $\tau$ associated to a generic constant
declaration is well-formed if $\tau$ denotes the \texttt{boolean}
keyword, or the \texttt{nat} keyword with a \emph{well-formed}
constraint.

The $etype_g$ relation is specially defined to check the
well-formedness of a subtype indication associated with a generic
constant declaration.

\begin{table}[H]
  \centering
  \begin{tabular}{@{}l}
    {\fontsize{8}{11}\selectfont\textsc{ETypeGBool}} \\
    {\begin{prooftree}
        
        % Conclusion
        \infer[template={\inserttext}]0
        {
          $\vdash$
          $~\mathtt{boolean}$
          $\xrightarrow{etype}$
          $\mathtt{bool}$
        }
      \end{prooftree}} \\
  \end{tabular}
    \begin{tabular}{@{}l}
    {\fontsize{8}{11}\selectfont\textsc{ETypeGNat}} \\
    {\begin{prooftree}
        
        % Well-formed constraint and evaluates constraint.
        \hypo{
          \Delta\vdash
          (\mathrm{e},\mathrm{e}')
          \xrightarrow{econstr_g}
          (v,v')
        }
        
        % Conclusion
        \infer[template={\inserttext}]1
        {
          $\vdash$
          $~\mathtt{natural}(\mathrm{e},\mathrm{e}')$
          $\xrightarrow{etype}$
          $\mathtt{nat}(v,v')$
        }
      \end{prooftree}} \\
  \end{tabular}
\end{table}

The $econstr_g$ relation checks that a \emph{generic} constraint (i.e,
a constraint appearing in a type indication associated with a generic
constant declaration) is well-formed and evaluates the constraint
bounds.

A \emph{generic} constraint $c$ is well-formed if:

\begin{itemize}[label=-]
\item its bounds are locally static expressions \cite[p.36]{VHDL00} of
  the \texttt{nat} type.
\item its lower bound value is inferior or equal to its upper bound
  value.
\end{itemize}

\begin{table}[H]
  \centering
  \begin{tabular}{@{}l}
    {\fontsize{8}{11}\selectfont\textsc{EConstrG}} \\
    {\begin{prooftree}
        
        % Locally static e and e'.
        \hypo{&SE_l(\mathrm{e})}
        \infer[no rule]1{&SE_l(\mathrm{e}')}
        
        % Evaluates e and e'
        \hypo{\Delta_\emptyset,\sigma_\emptyset,\Lambda_\emptyset&\vdash\mathrm{e}\xrightarrow{e}v}
        \infer[no rule]1{\Delta_\emptyset,\sigma_\emptyset,\Lambda_\emptyset&\vdash\mathrm{e}'\xrightarrow{e}v'}

        % Well-typed v and v'.
        \hypo{v&\in_c\mathtt{nat}(0,\mathtt{NATMAX})}
        \infer[no rule]1{v'&\in_c\mathtt{nat}(0,\mathtt{NATMAX})}
        
        % Conclusion
        \infer[template={\inserttext}]3
        [{
          \renewcommand{\arraystretch}{1.5}
          \begin{tabular}{@{}l}
            $v\leq{}v'$ \\
          \end{tabular}
        }]
        {
          $\vdash$
          $(\mathrm{e},\mathrm{e}')$
          $\xrightarrow{econstr_g}$
          $(v,v')$
        }
      \end{prooftree}} \\
  \end{tabular}
\end{table}

\subsection{Behavior elaboration.}
\label{subsec:beh-elab}

\subsubsection{Elaboration of concurrent statements.}
\label{subsubsec:conc-elab}

As it is, the elaboration of the composition of concurrent statements
is performed in a sequential manner.

\begin{table}[H]
  \begin{tabular}{@{}l}
  {\fontsize{8}{11}\selectfont\textsc{CsParElab}} \\
  {\begin{prooftree}

    % Elaborates first conc. stmt.
    \hypo{\mathcal{D},\Delta,\sigma\vdash\mathrm{cs}\xrightarrow{ebeh}\Delta',\sigma'}

    % Elaborates second conc. stmt.
    \hypo{\mathcal{D},\Delta',\sigma'\vdash\mathrm{cs'}\xrightarrow{ebeh}\Delta'',\sigma''}
    
    % Conclusion
    \infer[template={\inserttext}]2
    {
        $\mathcal{D},\Delta,\sigma\vdash$
        $~\mathrm{cs}~||~\mathrm{cs'}$
        $\xrightarrow{ebeh}$
        $\Delta'',\sigma''$
    }
  \end{prooftree}} \\
  \end{tabular}
  \begin{tabular}{l}
    {\fontsize{8}{11}\selectfont\textsc{CsNullElab}} \\
    {\begin{prooftree}

        % Conclusion
        \infer[template={\inserttext}]0
        {
          $\mathcal{D},\Delta,\sigma\vdash$
          $~\mathtt{null}$
          $\xrightarrow{ebeh}$
          $\Delta,\sigma$
        }
      \end{prooftree}} \\
  \end{tabular}
\end{table}

\subsubsection{Process elaboration.}
\label{subsubsec:ps-elab}

% Well-typed process with sensitivity list and vars.

\begin{premises}
  $\mathtt{valid}_{ss}$ states that a sequential statement is
  well-typed.
\end{premises}

\begin{sideconds}
  $\mathrm{sl}\subseteq{}Ins(\Delta)\cup{}Sigs(\Delta)$ indicates that
  the sensitivity list \emph{sl} must only contain signal identifiers
  that are readable, that is, \emph{input} ports and declared signals.
\end{sideconds}

\begin{figure}[H]
  {\fontsize{8}{11}\selectfont\textsc{PsElab}}
  \vspace{1em}
  
  \begin{prooftree}

    % Elaborates local context.
    \hypo{\Delta,\Lambda_\emptyset\vdash\mathrm{vars}\xrightarrow{evars}\Lambda}

    % Well-typed sequential statement.
    \hypo{\Delta,\sigma,\Lambda\vdash{}\mathtt{valid}_{ss}(\mathrm{ss})}
    
    % Conclusion
    \infer[template={\inserttext}]2
    [{\renewcommand{\arraystretch}{1.5}
      \begin{tabular}{@{}l}
        $\mathrm{id}_p\notin\Delta$ \\
        $\mathrm{sl}\subseteq{}Ins(\Delta)\cup{}Sigs(\Delta)$ \\
      \end{tabular}
    }]
    {
        $\mathcal{D},\Delta,\sigma\vdash$
        $~$\vhdle|process| \texttt{(}id$_p$\texttt{,} sl\texttt{,} vars\texttt{,} ss\texttt{)}
        $\xrightarrow{ebeh}$
        $\Delta\cup{}(id_p,\Lambda),\sigma$
    }
  \end{prooftree}
\end{figure}

\subsubsection{Process declarative part elaboration.}
\label{subsubsec:ps-decl-elab}

The $evars$ relation builds a local environment out of a process
declarative part.

% Variable declaration elaboration.

\begin{figure}[H]
  {\fontsize{8}{11}\selectfont\textsc{VarElab}}
  
  \begin{prooftree}
        
    % Elaborates type ind.
    \hypo{\Delta\vdash\tau\xrightarrow{etype}T}

    % Retrieves leftmost value of type T.
    \hypo{\vdash{}T\xrightarrow{defaultv}v}
    
    % Conclusion
    \infer[template={\inserttext}]2
    [$\mathrm{id}\notin{}\Lambda,\mathrm{id}\notin\Delta$]
    {
      $\Delta,\Lambda\vdash$
      $~(\mathrm{id},\tau)$
      $\xrightarrow{evars}$
      $\Lambda\cup{}(id,(T,v))$
    }
  \end{prooftree}
\end{figure}

% Inductive elab. rule for composition of variable declaration.

\begin{figure}[H]
  {\fontsize{8}{11}\selectfont\textsc{VarElabComp}}
  
  \begin{prooftree}

    % Evaluates VE1
    \hypo{\Delta,\Lambda\vdash\mathrm{vdecl}\xrightarrow{evars}{}\Lambda'}

    % Evaluates VE2.
    \hypo{\Delta,\Lambda'\vdash\mathrm{vars}\xrightarrow{evars}{}\Lambda''}

    % Conclusion
    \infer[template={\inserttext}]2{
      $\Delta,\Lambda\vdash
      ~\mathrm{vdecl}\mathtt{,}~\mathrm{vars}
      \xrightarrow{evars}\Lambda''$
    }
  \end{prooftree}
\end{figure}

\subsubsection{Component instantiation elaboration.}
\label{subsubsec:comp-inst-elab}

\begin{premises}
  \begin{itemize}[label=-]
  \item The $emapg$ relation binds a generic map to a function
    $\mathcal{M}\in{}id\nrightarrow{}value$ (see definition below).
  \item \texttt{valid}$_{ipm}$ (resp. \texttt{valid}$_{opm}$) states
    that a in port map (resp. out port map) is valid, i.e well-formed
    and well-typed (see \ref{subsec:valid-pm}).
  \end{itemize}
\end{premises}

\begin{sideconds}
  $\mathcal{M}\subseteq{}Gens(\Delta_c)$ checks that
  the generic map \emph{gmap} contains references to known generic
  constant identifiers only.
\end{sideconds}

\begin{table}[H]
  \begin{tabular}{@{}l}
    {\fontsize{8}{11}\selectfont\textsc{CompElab}} \\
    {\begin{prooftree}
        
        % Creates mapping and elaborates gens.
        \hypo{\mathcal{M}_\emptyset&\vdash\mathrm{gmap}\xrightarrow{emapg}\mathcal{M}}
        
        % Elaborates design.
        \hypo{
          \mathcal{D},\mathcal{M}&\vdash
          \mathcal{D}(\mathrm{id}_e)\xrightarrow{elab}\Delta_c,\sigma_c}

        % Checks in port map validity.
        \hypo{\Delta,\Delta_c,\sigma&\vdash\mathtt{valid}_{ipm}(\mathrm{ipmap})}

        % Checks out port map validity.
        \infer[no rule]1{\Delta,\Delta_c&\vdash\mathtt{valid}_{opm}(\mathrm{opmap})}
        
        % Conclusion
        \infer[template={\inserttext}]3
        [{
          \renewcommand{\arraystretch}{1.5}
          \begin{tabular}{@{}l}
            $\mathrm{id}_c\notin\Delta$, $\mathrm{id}_c\notin\sigma$ \\
            $\mathrm{id}_e\in{}\mathcal{D}$\\
            % $\mathcal{D}(\mathrm{id}_e)=$ \vhdlep{7}{9}|design| id$_{ce}$ id$_{ca}$ gens ports sigs cs \\
            $\mathcal{M}\subseteq{}Gens(\Delta_c)$ \\
          \end{tabular}
        }]
        {
          $\mathcal{D},\Delta,\sigma\vdash$
          \vhdle|comp| $(\mathrm{id}_c, \mathrm{id}_e, \mathrm{gmap}, \mathrm{ipmap}, \mathrm{opmap})\xrightarrow{ebeh}$
          $\Delta\cup{}(id_c,\Delta_c),$
          $\sigma\cup{}(id_c,\sigma_c)$
        }
      \end{prooftree}} \\
  \end{tabular}
\end{table}

A port map is a mapping between signals coming from an embedding
design ($\Delta$) and ports of an internal component instance
($\Delta_c$). The formal part of an port map entry (i.e, left of the
arrow) belongs to the internal component, whereas the actual part
(i.e, right of the arrow) refers to the embedding design. Therefore,
we need both $\Delta$ and $\Delta_c$ to verify if a port map is
well-typed leveraging the $\mathtt{valid}_{pm}$ predicate.

\begin{remark}[Valid generic map]
  Note that we are not checking the validity of the generic map.  In
  case of an ill-formed generic map, a inconsistent mapping
  $\mathcal{M}$ is generated by the \emph{emapg} that will make the
  \emph{elab} relation, taking $\mathcal{M}$ as a parameter, never
  derivable. Therefore, the \emph{elab} relation does an implicit
  validity check on the generic map.
\end{remark}

% Creates a mapping function from an assoc. list.

% Case formal part is not partial.

\begin{table}[H]
  \begin{tabular}{@{}l}
    {\fontsize{8}{11}\selectfont\textsc{AssocGElab}} \\
    
    {\begin{prooftree}

        % Locally static expression e.
        \hypo{SE_l(\mathrm{e})}
        
        % Evaluates e.
        \hypo{\Delta_\emptyset,\sigma_\emptyset,\Lambda_\emptyset\vdash\mathrm{e}\xrightarrow{e}v}

        % Conclusion
        \infer[template={\inserttext}]2
        [$\mathrm{id}_g\notin\mathcal{M}$]
        {
          $\mathcal{M}\vdash$
          $~\mathrm{id}_g\Rightarrow\mathrm{e}$
          $\xrightarrow{emapg}$
          $\mathcal{M}\cup{}(\mathrm{id}_g,v)$
        }
      \end{prooftree}} \\
  \end{tabular}
  \begin{tabular}{@{}l}
    {\fontsize{8}{11}\selectfont\textsc{GmapElabComp}} \\
    
    {\begin{prooftree}

        % Collects assoc_g.
        \hypo{\mathcal{M}\vdash\mathrm{assoc}_g\xrightarrow{emapg}\mathcal{M}'}

        % Collects lassoc_g.
        \hypo{\mathcal{M}'\vdash\mathrm{gmap}\xrightarrow{emapg}\mathcal{M}''}

        % Conclusion
        \infer[template={\inserttext}]2
        {
          $\mathcal{M}\vdash
          \mathrm{assoc}_g,~\mathrm{gmap}$
          $\xrightarrow{emapg}\mathcal{M}''$
        }
      \end{prooftree}} \\
  \end{tabular}
\end{table}

An $\mathrm{assoc}_g$ entry doesn't allow indexed id in its formal
part, due to the restriction of generic constants to scalar types.
Note that this restriction is not imposed by the VHDL language
reference; it is our stance towards a simplification of the VHDL
semantics.

\subsection{Implicit default value.}
\label{subsec:implicit-default-v}

According to the VHDL reference, when declaring a port, a signal or a
variable, these items must receive an implicit default value depending
on their types \cite[p.61, 64, 173]{VHDL00}. The $defaultv$ relation
determines the default value for a given type.

\begin{table}[H]
  \centering
  \begin{tabular}{@{}l}
    {\fontsize{10}{13}\selectfont\textsc{DefaultVBool}} \\
    {\begin{prooftree}

        % Conclusion
        \infer[template={\inserttext}]0
        {
          $\mathtt{bool}$
          $\xrightarrow{defaultv}\bot$
        }
      \end{prooftree}} \\
  \end{tabular}
  \begin{tabular}{@{}l}
    {\fontsize{10}{13}\selectfont\textsc{DefaultVCNat}} \\
    {\begin{prooftree}

        % Conclusion
        \infer[template={\inserttext}]0
        [$n\le{}m$]
        {
          $\mathtt{nat}(n,m)$
          $\xrightarrow{defaultv}n$
        }
      \end{prooftree}} \\
  \end{tabular}
\end{table}

\begin{table}[H]
  \centering
  \begin{tabular}{@{}l}
    {\fontsize{10}{13}\selectfont\textsc{DefaultVCArr}} \\
    {\begin{prooftree}

        % Retrieves default value for element subtype.
        \hypo{T\xrightarrow{defaultv}v}
        
        % Conclusion
        \infer[template={\inserttext}]1
        [
        \begin{tabular}{@{}l}
          $n\le{}m$ \\
          $size=(m-n)+1$ \\
        \end{tabular}
        ]
        {
          $\mathtt{array}(T,n,m)$
          $\xrightarrow{defaultv}$
          $\mathtt{create\_array}(size,T,v)$
        }
      \end{prooftree}} \\
  \end{tabular}
\end{table}

$\mathtt{create\_array}(size,T,v)$ creates an array of size $size$,
containing element of type $T$, where each element is initialized with
the value $v$.

%%% Local Variables:
%%% mode: latex
%%% TeX-master: "../../main"
%%% End:

\subsection{Typing relation}
\label{sec:constr-satif-rel}

The typing relation $\in_c$ checks that a given value conforms to a
given type.

\begin{table}[H]
  \centering
  \begin{tabular}{@{}l}
    {\fontsize{8}{11}\selectfont\textsc{IsBool}} \\
    {\begin{prooftree}

        % Conclusion.
        \infer0[$b\in{}\mathbb{B}$]{b\in_c{}bool}
      \end{prooftree}} \\
  \end{tabular}
  \begin{tabular}{@{}l}
    {\fontsize{8}{11}\selectfont\textsc{IsCNat}} \\
    {\begin{prooftree}

        % Conclusion.
        \infer0[$n\in{}[l,u]$]{n\in_c{}nat(l,u)}
      \end{prooftree}} \\
  \end{tabular}
  \begin{tabular}{@{}l}
    {\fontsize{8}{11}\selectfont\textsc{Array}} \\
    {\begin{prooftree}
        
        % Subelements satisfy subtype indication's constraint.
        \hypo{v_i\in_c{}T}
        
        % Conclusion.
        \infer1
        [{
          \begin{tabular}{@{}l}
            $i=1,\dots,n$ \\
            $n=(u-l)+1$ \\
          \end{tabular}
        }]
        {
          \Delta\vdash
          \mathtt{(}v_1,\dots,v_n\mathtt{)}
          \in_c
          array(T,l,u)}
      \end{prooftree}} \\
  \end{tabular}
\end{table}

\subsection{Static expressions}
\label{subsubsec:wfe}

Static expressions are either locally static or globally static; the
LRM defines locally static and globally static expressions as follows.

\subsubsection{Locally static expressions}
\label{subsubsec:loc-static}

An expression is \emph{locally} static if:

\begin{itemize}[label=-]
\item It is composed of operators and operands of a \emph{scalar} type
  (i.e, \texttt{natural} or \texttt{boolean}).
\item It is a \emph{literal} of a scalar type.
\end{itemize}

The $SE_l$ relation, defined by the following rules, states that an
expression is locally static.

\begin{table}[H]
  \begin{tabular}{@{}l}
    {\fontsize{8}{11}\selectfont\textsc{LSENat}} \\
    {\begin{prooftree}
        
        % Conclusion
        \infer[template={\inserttext}]0
        [$\mathrm{n}\in\mathbb{N}$]
        {$SE_l(\mathrm{n})$}
      \end{prooftree}} \\
  \end{tabular}
  \begin{tabular}{@{}l}
    {\fontsize{8}{11}\selectfont\textsc{LSEBool}} \\
    {\begin{prooftree}
        
        % Conclusion
        \infer[template={\inserttext}]0
        [$\mathrm{b}\in\mathbb{B}$]
        {$SE_l(\mathrm{b})$}
      \end{prooftree}} \\
  \end{tabular}
  \begin{tabular}{@{}l}
    {\fontsize{8}{11}\selectfont\textsc{LSENot}} \\
    {\begin{prooftree}

        % Locally static e.
        \hypo{SE_l(\mathrm{e})}
        
        % Conclusion
        \infer[template={\inserttext}]1
        {
          $SE_l(\mathtt{not}~\mathrm{e})$
        }
      \end{prooftree}} \\
  \end{tabular}
  \begin{tabular}{@{}l}
    {\fontsize{8}{11}\selectfont\textsc{LSEBinOp}}\\    
    {\begin{prooftree}

        % Locally static e.
        \hypo{SE_l(\mathrm{e})}

        % Locally static e'.
        \hypo{SE_l(\mathrm{e'})}
        
        % Conclusion
        \infer[template={\inserttext}]2
        [$\mathtt{op}\in\{~+,-,=,\neq,<,\le,>,\ge,\mathtt{and},\mathtt{or}~\}$]
        {
          $SE_l(\mathrm{e}~\mathtt{op}~\mathrm{e'})$
        }
      \end{prooftree}}\\
  \end{tabular}
\end{table}

\subsubsection{Globally static expressions}
\label{subsubsec:glob-static}

An expression is \emph{globally} static in the context $\Delta$ if:

\begin{itemize}[label=-]
\item It is a generic constant.
\item It is an array aggregate composed of globally static
  expressions.
\item It is a locally static expression.
\end{itemize}

The $SE_g$ relation, defined by the following rules, checks that an
expression is globally static is a given context $\Delta$.

\begin{table}[H]
  \centering
  \begin{tabular}{@{}l}
    {\fontsize{8}{11}\selectfont\textsc{GSELocal}} \\
    {\begin{prooftree}

      % Locally static e.
      \hypo{SE_l(\mathrm{e})}
      
      % Conclusion
      \infer[template={\inserttext}]1
      {
        $\Delta\vdash{}SE_g(\mathrm{e})$
      }
    \end{prooftree}} \\
  \end{tabular}
  \begin{tabular}{@{}l}
    {\fontsize{8}{11}\selectfont\textsc{GSEGen}} \\
    {\begin{prooftree}
      
      % Conclusion
      \infer[template={\inserttext}]0
      [{\renewcommand{\arraystretch}{1.5}
        \begin{tabular}{@{}l}
          $\mathrm{id}_g\in{}Gens(\Delta)$ \\
        \end{tabular}
      }]
      {
        $\Delta\vdash{}SE_g(\mathrm{id}_g)$
      }
    \end{prooftree}} \\
  \end{tabular}
  \begin{tabular}{@{}l}
    {\fontsize{8}{11}\selectfont\textsc{GSEAggregate}} \\
    {\begin{prooftree}

        % Locally static e.
        \hypo{\Delta\vdash{}SE_g(\mathrm{e}_i)}
        
        % Conclusion
        \infer[template={\inserttext}]1
        [$i=1,\dots,n$]
        {
          $\Delta\vdash{}SE_g($\texttt{(}$\mathrm{e}_1,\dots,\mathrm{e}_n$\texttt{)}$)$
        }
      \end{prooftree}} \\
  \end{tabular}
\end{table}

\subsection{Valid port map}
\label{subsec:valid-pm}

\paragraph{Valid input port map}

The $\mathtt{valid}_{ipm}$ predicate states that an input port map is
valid in the context $\Delta,\Delta_c$, where $\Delta$ is the
embedding design structure and $\Delta_c$ denotes the component
instance, owner of the input port map, if:

\begin{itemize}[label=-]
\item All ports defined in $\Delta_c$ are exactly mapped once in the
  input port map.
\item For each input port map entry, the formal and actual part are of
  the same type.
\end{itemize}

\begin{premises}
  \begin{itemize}[label=-]
  \item $list_{ipm}$ builds a set
    $\mathcal{L}\subset{}id\sqcup{}(id\times\mathbb{N})$ out of the
    input port map.
  \item $\mathtt{check}_{pm}$ checks the validity of a port map based
    on the corresponding port list (here, the input ports of
    $\Delta_c$) and the set built by the $list_{ipm}$ relation.
  \end{itemize}
\end{premises}

\begin{table}[H]
  \centering
  \begin{tabular}{l}
    {\fontsize{8}{11}\selectfont\textsc{ValidIPM}} \\
    {\begin{prooftree}

        % Lists formal parts.
        \hypo{\Delta,\Delta_c,\sigma,\mathcal{L}_\emptyset\vdash\mathrm{i}\xrightarrow{list_{ipm}}\mathcal{L}}

        % Well-formed port map.
        \hypo{\mathtt{check}_{pm}(Ins(\Delta_c),\mathcal{L})}
        
        % Conclusion
        \infer[template={\inserttext}]2
        {
          $\Delta,\Delta_c,\sigma\vdash\mathtt{valid}_{ipm}(\mathrm{i})$
        }
      \end{prooftree}} \\
  \end{tabular}
\end{table}

The $list_{ipm}$ relation builds a set composed of identifiers and/or
couples \textit{(identifier, natural number)} collected from the
identifiers and indexed identifiers found in the formal parts of an
input port map. It also checks, for each association of the input port
map, that the expression of the actual part is of the same type than
the identifier or indexed identifier of the formal part.

\begin{sideconds}
  \begin{itemize}
  \item $\mathrm{id}_f\in{}Ins(\Delta_c)$ checks that the identifier
    $\mathrm{id}_f$ is an input port identifier of $\Delta_c$.
  \item $\mathrm{id}_f\notin\mathcal{L}$ checks that the port
    identifier $\mathrm{id}_f$ is not already mapped, i.e. it is not
    already referenced in the $\mathcal{L}$ set.
  \item $\nexists{}v_i$ s.t. $(\mathrm{id}_f,v_i)\in\mathcal{L}$
    checks that a subelement of id$_f$ is not already map, that is, if
    id$_f$ denotes a signal identifier of the \texttt{array} type.
  \end{itemize}
\end{sideconds}

\begin{table}[H]
  \begin{tabular}{@{}l}
    {\fontsize{8}{11}\selectfont\textsc{ListIPMSimple}} \\
    {\begin{prooftree}

        % Evaluates actual.
        \hypo{\Delta,\sigma\vdash\mathrm{e}~\xrightarrow{e}v}

        % Well-typed actual.
        \hypo{v\in_c{}T}
        
        % Conclusion
        \infer[template={\inserttext}]
        2
        [{\renewcommand{\arraystretch}{1.5}
          \begin{tabular}{l}
            $\mathrm{id}_f\notin\mathcal{L}$, $\mathrm{id}_f\in{}Ins(\Delta_c)$\\
            $\nexists{}v_i$ s.t. $(\mathrm{id}_f,v_i)\in\mathcal{L}$\\
            $\Delta_c(\mathrm{id}_f)=T$ \\
          \end{tabular}
        }]
        {
          $\Delta,\Delta_c,\sigma,\mathcal{L}\vdash
          ~(\mathrm{id}_f,\mathrm{e})~$
          $\xrightarrow{list_{ipm}}\mathcal{L}\cup{}\{id_f\}$
        }
      \end{prooftree}} \\
  \end{tabular}
\end{table}

\begin{premises}
  $v_i\in_c{}nat(n,m)$ checks that the index value stays in the
  array bounds.
\end{premises}

\begin{sideconds}
  $\mathrm{id}_f\notin\mathcal{L}$ and
  $(\mathrm{id}_f, v_i)\notin\mathcal{L}$ checks that neither the port
  identifier $id_f$ nor the couple port identifier $id_f$ and index
  $v_i$ are already mapped.
\end{sideconds}

\begin{figure}[H]
  {\fontsize{8}{11}\selectfont\textsc{ListIPMPartial}}
  \vspace{1em}

  \begin{prooftree}

    % Static expr e_i.
    \hypo{SE_l(\mathrm{e}_i)}

    % Evaluates e_i and e.
    \hypo{&\mathrm{e}_i\xrightarrow{e}v_i}
    \infer[no rule]1{\Delta,\sigma&\vdash\mathrm{e}\xrightarrow{e}v}
    
    % Evaluates e.
    \hypo{v_i&\in_c{}nat(n,m)}
    \infer[no rule]1{v&\in_c{}T}
    
    
    % Conclusion
    \infer[template={\inserttext}]
    3
    [{\renewcommand{\arraystretch}{1.5}
      \begin{tabular}{l}
        $\mathrm{id}_f\notin\mathcal{L}$, $(\mathrm{id}_f, v_i)\notin\mathcal{L}$ \\
        $\mathrm{id}_f\in{}Ins(\Delta_c)$ \\
        $\Delta_c(\mathrm{id}_f)=array(T,n,m)$ \\
      \end{tabular}
    }]
    {
      $\Delta,\Delta_c,\sigma,\mathcal{L}\vdash
      ~(\mathrm{id}_f(\mathrm{e}_i),\mathrm{e})~$
      $\xrightarrow{list_{ipm}}
      \mathcal{L}\cup\{~(id_f,v_i)~\}$
    }
  \end{prooftree}
\end{figure}

\begin{figure}[H]
  {\fontsize{8}{11}\selectfont\textsc{ListIPMCons}}
  \vspace{1em}

  \begin{prooftree}

    % Lists assoc.
    \hypo{\Delta,\Delta_c,\sigma,\mathcal{L}\vdash\mathrm{assoc}_{ip}\xrightarrow{list_{ipm}}\mathcal{L}'}

    % Lists lassoc.
    \hypo{\Delta,\Delta_c,\sigma,\mathcal{L}'\vdash\mathrm{i}\xrightarrow{list_{ipm}}\mathcal{L}''}
    
    % Conclusion
    \infer[template={\inserttext}] 2 {
      $\Delta,\Delta_c,\sigma,\mathcal{L}\vdash
      ~\mathrm{assoc}_{ip}\mathtt{,}~\mathrm{i}~$
      $\xrightarrow{list_{ipm}}\mathcal{L}''$ }
  \end{prooftree}
\end{figure}

The $\mathtt{check}_{pm}(Ports,\mathcal{L})$ predicate states that all
port identifiers referenced in the domain of
$Ports\in{}id\nrightarrow{}type$ appear in $\mathcal{L}$ as a simple
identifier, or if the port identifier is of the $array$ type, then all
couples ($id$,$i$) must belong to $\mathcal{L}$, where $i$ denotes all
indexes of the index range and $id$ denotes the port identifier.

\begin{equation*}
  \begin{split}
    \mathtt{check}_{pm}(Ports,\mathcal{L})\equiv
    \forall{}\mathrm{id}_f\in\mathtt{dom}(Ports),~\mathrm{id}_f\in\mathcal{L}~\lor~&
    (Ports(\mathrm{id}_f)=\mathtt{array}(T,n,m)\land\\
    & \forall{}i\in[n,m],~(\mathrm{id}_f,i)\in\mathcal{L}) \\
  \end{split}
\end{equation*}

\paragraph{Valid output port map}

The $\mathtt{valid}_{opm}$ predicate states that an \emph{output} port
map is valid in the context $\Delta,\Delta_c$, where $\Delta$ is the
embedding design structure and $\Delta_c$ denotes the component
instance owner of the port map, if:

\begin{itemize}[label=-]
\item An output port identifier appears at most once in the output
  port map.
\item Two different output port identifiers cannot be connected to the
  same signal.
\item For each output port map entry, the formal and the actual part
  are of the exact same type.
\end{itemize}

We allow partially connected output port map; i.e, an output port map
where all output ports might not be present in the mapping. Such
output ports are \texttt{open} by default.

\begin{premises}
  $list_{opm}$ builds two sets
  $\mathcal{L},\mathcal{L}_{ids}\subseteq{}id\sqcup{}(id\times\mathbb{N})$
  out of the output port map $\textrm{opmap}$. $\mathcal{L}_{ids}$ is
  built incrementally to check that there are no multiply-driven
  signals resulting of the port map connection.
\end{premises}

\begin{table}[H]
  \centering
  \begin{tabular}{l}
    {\fontsize{8}{11}\selectfont\textsc{ValidOPM}} \\
    {\begin{prooftree}

        % Lists formal parts.
        \hypo{
          \Delta,\Delta_c,\mathcal{L}_\emptyset,\mathcal{L}_{ids\emptyset}\vdash
          \mathrm{o}
          \xrightarrow{list_{opm}}
          \mathcal{L},\mathcal{L}_{ids}
        }
        
        % Conclusion
        \infer[template={\inserttext}]
        1
        {
          $\Delta,\Delta_c\vdash\mathtt{valid}_{opm}(\mathrm{o})$
        }
      \end{prooftree}} \\
  \end{tabular}
\end{table}

\begin{sideconds}
  \begin{itemize}[label=-]
  \item $\mathrm{id}_f\notin\mathcal{L}$ checks that the port
    identifier $\mathrm{id}_f$ is not already mapped (i.e, is not
    already used in the formal part of a port map entry).
  \item $\mathrm{id}_a\notin\mathcal{L}_{ids}$ checks that the signal
    identifier $\mathrm{id}_a$ is not already mapped (i.e, is not
    already used in the actual part of a port map entry).
  \item $\mathrm{id}_f\in{}Outs(\Delta_c)$ checks that $id_f$ is an
    output port identifier of $\Delta_c$.
  \item $\mathrm{id}_a\in{}Sigs(\Delta)\cup{}Outs(\Delta)$ checks that
    $\mathrm{id}_a$ is either an output port or an internal signal
    identifier of $\Delta$.
  \item $\Delta_c(\mathrm{id}_f)=\Delta(\mathrm{id}_a)=T$ checks that
    $\mathrm{id}_f$ and $\mathrm{id}_a$ are exactly of the same type.
  \end{itemize}
\end{sideconds}

\begin{table}[H]
  \begin{tabular}{@{}l}
    {\fontsize{8}{11}\selectfont\textsc{ListOPMSimpleToSimple}} \\
    {\begin{prooftree}
        
        % Conclusion
        \infer[template={\inserttext}]
        0
        [{\renewcommand{\arraystretch}{1.5}
          \begin{tabular}{@{}l}
            $\mathrm{id}_f\notin\mathcal{L}$, $\mathrm{id}_a\notin\mathcal{L}_{ids}$ \\
            $\mathrm{id}_f\in{}Outs(\Delta_c)$\\
            $\mathrm{id}_a\in{}Sigs(\Delta)\cup{}Outs(\Delta)$ \\
            $\Delta_c(\mathrm{id}_f)=\Delta(\mathrm{id}_a)=T$ \\
          \end{tabular}
        }] {
          $\Delta,\Delta_c,\mathcal{L},\mathcal{L}_{ids}\vdash$
          $(\mathrm{id}_f,\mathrm{id}_a)$
          $\xrightarrow{list_{opm}}$
          $\mathcal{L}\cup{}\{id_f\},\mathcal{L}_{ids}\cup{}\{id_a\}$ }
      \end{prooftree}} \\
  \end{tabular}
\end{table}

\begin{sideconds}
  $Outs_c(\mathrm{id}_f)=T$ and
  $Sigs(\mathrm{id}_a)=\mathtt{array(}T,n,m\mathtt{)}$ checks that the
  type of $\mathrm{id}_f$ and the type of the elements of
  $\mathrm{id}_a$ are the same. Note that $\mathrm{id}_a$ be an signal
  identifier of the \texttt{array} type as $\mathrm{id}_f$ is mapped
  to one subelement of $\mathrm{id}_a$.
\end{sideconds}

\begin{table}[H]
  \begin{tabular}{@{}l}
    {\fontsize{8}{11}\selectfont\textsc{ListOPMSimpleToPartial}} \\
    {\begin{prooftree}

        % Static expr e_i and natural.
        \hypo{SE_l(\mathrm{e}_i)}

        % Evaluates e_i.
        \hypo{\mathrm{e}_i\xrightarrow{e}v_i}

        % Well-typed v_i (in array bounds).
        \hypo{v_i\in_c{}nat(n,m)}
        
        % Conclusion
        \infer[template={\inserttext}]    
        3
        [{\renewcommand{\arraystretch}{1.5}
          \begin{tabular}{@{}l}
            $\mathrm{id}_f\notin\mathcal{L}$, $\mathrm{id}_a$, $(\mathrm{id}_a, v_i)\notin\mathcal{L}_{ids}$ \\
            $\mathrm{id}_f\in{}Outs(\Delta_c)$ \\
            $\mathrm{id}_a\in{}Sigs(\Delta)\cup{}Outs(\Delta)$ \\
            $\Delta_c(\mathrm{id}_f)=T$ \\
            $\Delta(\mathrm{id}_a)=array(T,n,m)$ \\
          \end{tabular}
        }]
        {
          $\Delta,\Delta_c,\mathcal{L},\mathcal{L}_{ids}\vdash
          ~(\mathrm{id}_f,\mathrm{id}_a(\mathrm{e}_i))~$
          $\xrightarrow{list_{opm}}
          \mathcal{L}\cup{}\{id_f\},\mathcal{L}_{ids}\cup{}\{(id_a,v_i)\}$
        }
      \end{prooftree}} \\
  \end{tabular}
\end{table}

\begin{table}[H]
  \begin{tabular}{@{}l}
    {\fontsize{8}{11}\selectfont\textsc{ListOPMSimpleToOpen}} \\
    {\begin{prooftree}

        % Conclusion
        \infer[template={\inserttext}]
        0
        [{
          \begin{tabular}{l}
            $\mathrm{id}_f\notin\mathcal{L}$ \\
            $\mathrm{id}_f\in{}Outs(\Delta_c)$ \\
          \end{tabular}
        }]
        {
          $\Delta,\Delta_c,\mathcal{L},\mathcal{L}_{ids}\vdash
          ~(\mathrm{id}_f,\mathtt{open})~$
          $\xrightarrow{list_{opm}}\mathcal{L}\cup{}\{id_f\},\mathcal{L}_{ids}$
        }
      \end{prooftree}} \\
  \end{tabular}
\end{table}

\begin{remark}[Unconnected output port.]
  We forbid the case where an indexed formal part corresponding to the
  subelement of a composite output port is unconnected, i.e
  $(\mathrm{id}_f(e_i),\mathtt{open})$, as it could lead to the case
  where some subelements of a composite output port are connected
  while others are not (error case in \cite[p.7]{VHDL2000}).
\end{remark}

\begin{table}[H]
  \begin{tabular}{@{}l}
    {\fontsize{8}{11}\selectfont\textsc{ListOPMPartialToSimple}} \\
    {\begin{prooftree}

        % Static expr e_i and natural.
        \hypo{SE_l(\mathrm{e}_i)}

        % Evaluates e_i.
        \hypo{\mathrm{e}_i\xrightarrow{e}v_i}

        % Well-typed v_i (in array bounds).
        \hypo{v_i\in_c{}nat(n,m)}
        
        % Conclusion
        \infer[template={\inserttext}]
        3
        [{\renewcommand{\arraystretch}{1.5}
          \begin{tabular}{l}
            $\mathrm{id}_f,(\mathrm{id}_f,v_i)\notin\mathcal{L}$, $\mathrm{id}_a\notin\mathcal{L}_{ids}$ \\
            $\mathrm{id}_f\in{}Outs(\Delta_c)$ \\
            $\mathrm{id}_a\in{}Sigs(\Delta)\cup{}Outs(\Delta)$ \\
            $\Delta_c(\mathrm{id}_f)=array(T,n,m)$ \\
            $\Delta(\mathrm{id}_a)=T$ \\
          \end{tabular}
        }]
        {
          $\Delta,\Delta_c,\mathcal{L},\mathcal{L}_{ids}\vdash$
          $(\mathrm{id}_f(\mathrm{e}_i),\mathrm{id}_a)$
          $\xrightarrow{list_{opm}}$
          $\mathcal{L}\cup\{(id_f,v_i)\},\mathcal{L}_{ids}\cup\{id_a\}$
        }
      \end{prooftree}} \\
  \end{tabular}
\end{table}

\begin{table}[H]
  \begin{tabular}{@{}l}
    {\fontsize{8}{11}\selectfont\textsc{ListOPMPartialToPartial}} \\
    {\begin{prooftree}

        % Static expr. ei e'i.
        
        \hypo{&SE_l(\mathrm{e'}_i)}
        \infer[no rule]1{&SE_l(\mathrm{e}_i)}

        % Evaluates e_i.
        \hypo{&\mathrm{e'}_i\xrightarrow{e}v'_i}
        \infer[no rule]1{&\mathrm{e}_i\xrightarrow{e}v_i}

        % Well-typed v'_i and v_i.
        \hypo{v'_i&\in_c{}nat(n',m')}
        \infer[no rule]1{v_i&\in_c{}nat(n,m)}
        
        % Conclusion
        \infer[template={\inserttext}]
        3
        [{\renewcommand{\arraystretch}{1.5}
          \begin{tabular}{l}
            $\mathrm{id}_f,(\mathrm{id}_f,v_i)\notin\mathcal{L}$, $\mathrm{id}_a,(\mathrm{id}_a,v'_i)\notin\mathcal{L}_{ids}$ \\
            $\mathrm{id}_f\in{}Outs(\Delta_c)$ \\
            $\mathrm{id}_a\in{}Sigs(\Delta)\cup{}Outs(\Delta)$ \\
            $\Delta_c(\mathrm{id}_f)=array(T,n,m)$ \\
            $\Delta(\mathrm{id}_a)=array(T,n',m')$ \\ 
          \end{tabular}
        }]
        {
          $\Delta,\Delta_c,\mathcal{L},\mathcal{L}_{ids}\vdash$
          $(\mathrm{id}_f(\mathrm{e}_i),\mathrm{id}_a(\mathrm{e'}_i))$
          $\xrightarrow{list_{opm}}$
          \begin{tabular}{@{}l@{}}
            $\mathcal{L}\cup\{(id_f,v_i)\},$ \\
            $\mathcal{L}_{ids}\cup\{(id_a,v'_i)\}$\\
          \end{tabular}
        }
      \end{prooftree}} \\
  \end{tabular}
\end{table}

\begin{table}[H]
  \begin{tabular}{@{}l}
    {\fontsize{8}{11}\selectfont\textsc{ListOPMCons}} \\
    {\begin{prooftree}

        % Lists assoc.
        \hypo{
          \Delta,\Delta_c,\mathcal{L},\mathcal{L}_{ids}\vdash
          \mathrm{assoc}_{po}
          \xrightarrow{list_{opm}}
          \mathcal{L}',\mathcal{L}'_{ids}
        }

        % Lists lassoc.
        \hypo{
          \Delta,\Delta_c,\mathcal{L}',\mathcal{L}'_{ids}\vdash
          \mathrm{opmap}
          \xrightarrow{list_{opm}}
          \mathcal{L}'',\mathcal{L}''_{ids}
        }
        
        % Conclusion
        \infer[template={\inserttext}]
        2
        {
          $\Delta,\Delta_c,\mathcal{L},\mathcal{L}_{ids}\vdash
          ~\mathrm{assoc}_{po}\mathtt{,}~\mathrm{opmap}~$
          $\xrightarrow{list_{opm}}$
          $\mathcal{L}'',\mathcal{L}''_{ids}$
        }
      \end{prooftree}} \\
  \end{tabular}
\end{table}

\subsection{Valid sequential statements}
\label{subsubsec:valid-ss}

The $\mathtt{valid}_{ss}$ predicate states that a sequential statement
is well-typed in the context $\Delta,\sigma,\Lambda$.

\subsubsection{Well-typed signal assignment}

% Well-typed signal assignment.

\begin{premises}
  \begin{itemize}
  \item $\Delta,\sigma,\Lambda\vdash{}\mathrm{e}\xrightarrow{e}v$
    evaluates the expression assigned to signal $\mathrm{id}_s$ in the
    context $\Delta,\sigma,\Lambda$.% During the elaboration, $\sigma$
    % corresponds to the default design state, i.e. where each signal is
    % associated to its type default value.
  \item $v\in_c{}T$ checks that the value of expression $e$ conforms
    to the type of signal $\mathrm{id}_s$.
  \end{itemize}
\end{premises}

\begin{table}[H]
  \begin{tabular}{@{}l}
    {\fontsize{8}{11}\selectfont\textsc{WTSig}} \\
    {\begin{prooftree}

        % Evaluates e.
        \hypo{\Delta,\sigma,\Lambda\vdash{}\mathrm{e}\xrightarrow{e}v}
        
        % Well-typed v.
        \hypo{v\in_c{}T}
        
        % Conclusion
        \infer[template={\inserttext}]2
        [{
          \renewcommand{\arraystretch}{1.5}
          \begin{tabular}{@{}l}
            $\mathrm{id}_s\in{}Sigs(\Delta)\cup{}Outs(\Delta)$\\
            $\Delta(\mathrm{id}_s)=T$\\
          \end{tabular}
        }]
        {
          $\Delta,\sigma,\Lambda\vdash\mathtt{valid}_{ss}(\mathrm{id}_s\Leftarrow\mathrm{e})$
        }
      \end{prooftree}} \\
  \end{tabular}
\end{table}

% Well-typed indexed signal assignment.
\begin{table}[H]
  \begin{tabular}{l}
    {\fontsize{8}{11}\selectfont\textsc{WTIdxSig}} \\
    {\begin{prooftree}

        % Evaluates e and e_i.
        \hypo{\Delta,\sigma,\Lambda&\vdash\mathrm{e}\xrightarrow{e}v}
        \infer[no rule]1{\Delta,\sigma,\Lambda&\vdash\mathrm{e_i}\xrightarrow{e}v_i}
        
        % Well-typed v and v_i.
        \hypo{v&\in_c{}T}
        \infer[no rule]1{v_i&\in_c{}nat(n,m)}
        
        % Conclusion
        \infer[template={\inserttext}]2
        [{
          \renewcommand{\arraystretch}{1.5}
          \begin{tabular}{@{}l}
            $\mathrm{id}_s\in{}Sigs(\Delta)\cup{}Outs(\Delta)$ \\
            $\Delta(\mathrm{id}_s)=array(T,n,m)$ \\
          \end{tabular}
        }]
        {
          $\Delta,\sigma,\Lambda\vdash\mathtt{valid}_{ss}(\mathrm{id}_s(\mathrm{e}_i)\Leftarrow\mathrm{e})$
        }
      \end{prooftree}} \\
  \end{tabular}
\end{table}

\subsubsection{Well-typed variable assignment}

\begin{table}[H]
  \begin{tabular}{@{}l}
    {\fontsize{8}{11}\selectfont\textsc{WTVar}} \\
    {\begin{prooftree}
        
        % Evaluates e.
        \hypo{\Delta,\sigma,\Lambda\vdash\mathrm{e}\xrightarrow{e}v}

        % Well-typed v.
        \hypo{v\in_c{}T}
        
        % Conclusion
        \infer[template={\inserttext}]2
        [{
          \renewcommand{\arraystretch}{1.5}
          \begin{tabular}{@{}l}
            $\mathrm{id}_v\in\Lambda$\\
            $\Lambda(\mathrm{id}_v)=(T,val)$\\
          \end{tabular}
        }]
        {
          $\Delta,\sigma,\Lambda\vdash\mathtt{valid}_{ss}(\mathrm{id}_v := \mathrm{e})$
        }
      \end{prooftree}} \\
  \end{tabular}
\end{table}

\begin{table}[H]    
  \begin{tabular}{@{}l}
    {\fontsize{8}{11}\selectfont\textsc{WTIdxVar}} \\
    {\begin{prooftree}

        % Evaluates e and e_i.
        \hypo{\Delta,\sigma,\Lambda&\vdash\mathrm{e}\xrightarrow{e}v}
        \infer[no rule]1{\Delta,\sigma,\Lambda&\vdash\mathrm{e_i}\xrightarrow{e}v_i}
        
        % Well-typed v and v_i.
        \hypo{v&\in_c{}T}
        \infer[no rule]1{v_i&\in_c{}nat(n,m)}
        
        % Conclusion
        \infer[template={\inserttext}]2
        [{
          \renewcommand{\arraystretch}{1.5}
          \begin{tabular}{@{}l}
            $\mathrm{id}_v\in{}\Lambda$ \\
            $\Lambda(\mathrm{id}_v)=(array(T,n,m),val)$ \\
          \end{tabular}
        }]
        {
          $\Delta,\Lambda\vdash\mathtt{valid}_{ss}(\mathrm{id}_v(\mathrm{e}_i) := \mathrm{e})$
        }
      \end{prooftree}} \\
  \end{tabular}
\end{table}%

\subsubsection{Well-typed if statements}

\begin{table}[H]
  \begin{tabular}{@{}l}
    {\fontsize{8}{11}\selectfont\textsc{WTIf}} \\
    {\begin{prooftree}

        % Eval e.
        \hypo{\Delta,\sigma,\Lambda\vdash\mathrm{e}\xrightarrow{e}v}
        
        % Well-typed v.
        \hypo{v\in_c\mathtt{bool}}

        % Validss ss.
        \hypo{\Delta,\sigma,\Lambda\vdash\mathtt{valid}_{ss}(\mathrm{ss})}
        
        % Conclusion
        \infer[template={\inserttext}]3
        {
          $\Delta,\sigma,\Lambda\vdash\mathtt{valid}_{ss}($\vhdle|if| \texttt{(}e\texttt{)} ss$)$
        }
      \end{prooftree}} \\
  \end{tabular}
\end{table}%

\begin{table}[H]
  \begin{tabular}{@{}l}
    {\fontsize{8}{11}\selectfont\textsc{WTIfElse}} \\
    {\begin{prooftree}

        % Eval e.
        \hypo{\Delta,\sigma,\Lambda\vdash\mathrm{e}\xrightarrow{e}v}
        
        % Well-typed v.
        \hypo{v\in_c\mathtt{bool}}

        % Validss ss and ss'.
        \hypo{\Delta,\sigma,\Lambda\vdash\mathtt{valid}_{ss}(\mathrm{ss})}
        \infer[no rule]1{\Delta,\sigma,\Lambda\vdash\mathtt{valid}_{ss}(\mathrm{ss}')}
        
        % Conclusion
        \infer[template={\inserttext}]3
        {
          $\Delta,\sigma,\Lambda\vdash\mathtt{valid}_{ss}($\vhdle|if| \texttt{(}e\texttt{)} ss ss'$)$
        }
      \end{prooftree}} \\
  \end{tabular}
\end{table}%

\subsubsection{Well-typed loop statement}

\begin{table}[H]
  \begin{tabular}{@{}l}
    {\fontsize{8}{11}\selectfont\textsc{WTLoop}} \\
    {\begin{prooftree}

        % Eval e and e'.
        \hypo{\Delta,\sigma,\Lambda&\vdash\mathrm{e}\xrightarrow{e}v}
        \infer[no rule]1{\Delta,\sigma,\Lambda&\vdash\mathrm{e'}\xrightarrow{e}v'}

        % Well-typed v and v'.
        \hypo{v&\in_c{}{}nat(0,\mathtt{NATMAX})}
        \infer[no rule]1{v'&\in_c{}{}nat(0,\mathtt{NATMAX})}
        
        % Well-typed ss.
        \hypo{\Delta,\sigma,\Lambda'\vdash\mathtt{valid}_{ss}(\mathrm{ss})}
        
        % Conclusion
        \infer[template={\inserttext}]3
        [$\Lambda'=\Lambda\cup(\mathrm{id}_v,({}nat(v,v'), v))$]
        {
          $\Delta,\sigma,\Lambda\vdash\mathtt{valid}_{ss}($\vhdle|for| (id$_v$, e, e') ss$)$
        }
      \end{prooftree}} \\
  \end{tabular}
\end{table}%

\subsubsection{Well-typed rising and falling edge blocks}

\begin{table}[H]
  \begin{tabular}{@{}l}
    {\fontsize{8}{11}\selectfont\textsc{WTRising}} \\
    {\begin{prooftree}
        
        % Well-typed ss.
        \hypo{\Delta,\sigma,\Lambda\vdash\mathtt{valid}_{ss}(\mathrm{ss})}
        
        % Conclusion
        \infer[template={\inserttext}]1
        {
          $\Delta,\sigma,\Lambda\vdash\mathtt{valid}_{ss}($\vhdle|rising| ss$)$
        }
      \end{prooftree}} \\
  \end{tabular}
  \begin{tabular}{@{}l}
    {\fontsize{8}{11}\selectfont\textsc{WTFalling}} \\
    {\begin{prooftree}
        
        % Well-typed ss.
        \hypo{\Delta,\sigma,\Lambda\vdash\mathtt{valid}_{ss}(\mathrm{ss})}
        
        % Conclusion
        \infer[template={\inserttext}]1
        {
          $\Delta,\sigma,\Lambda\vdash\mathtt{valid}_{ss}($\vhdle|falling| ss$)$
        }
      \end{prooftree}} \\
  \end{tabular}
\end{table}%

\subsubsection{Well-typed rst blocks}

\begin{table}[H]
  \begin{tabular}{@{}l}
    {\fontsize{8}{11}\selectfont\textsc{WTRst}} \\
    {\begin{prooftree}
        
        % Well-typed ss.
        \hypo{\Delta,\sigma,\Lambda\vdash\mathtt{valid}_{ss}(\mathrm{ss})}

        % Well-typed ss'.
        \hypo{\Delta,\sigma,\Lambda\vdash\mathtt{valid}_{ss}(\mathrm{ss'})}
        
        % Conclusion
        \infer[template={\inserttext}]2
        {
          $\Delta,\sigma,\Lambda\vdash\mathtt{valid}_{ss}($\vhdle|rst| ss ss'$)$
        }
      \end{prooftree}} \\
  \end{tabular}
\end{table}%

\subsubsection{Well-typed null statement}
\label{sec:wt-null}

\begin{table}[H]
  \begin{tabular}{@{}l}
    {\fontsize{8}{11}\selectfont\textsc{WTNull}} \\
    {\begin{prooftree}
                
        % Conclusion
        \infer[template={\inserttext}]0
        {
          $\Delta,\sigma,\Lambda\vdash\mathtt{valid}_{ss}($\vhdle|null|$)$
        }
      \end{prooftree}} \\
  \end{tabular}
\end{table}


%%% Local Variables:
%%% mode: latex
%%% TeX-master: "../../main"
%%% End:


\section{Simulation rules}
\label{sec:sim-rules}
% \subsection{Well-formed environment.}
% \label{subsec:env-wf}

% To avoid the situation where an identifier met in an expression refers
% both to a variable and a signal at the same time, we need to ensure
% that identifiers referenced in the environment are unique.  We define
% the predicate $\mathcal{U}$ stating that identifier unicity is met at
% different environment levels.

% \begin{figure}[H]
%   {\fontsize{10}{13}\selectfont\textsc{WFTypedefEnv}}
  
%   \begin{prooftree}

%     % No intersection dom and val
%     \hypo{\mathtt{dom}(f)\cap\mathtt{val}(f)=\emptyset}

%     % Unique val.
%     \infer[no rule]1{
%       \forall{}id,id'\in\mathtt{dom}(f),~id\neq{}id'\wedge\mathtt{enum}(f(id))\wedge\mathtt{enum}(f(id'))
%       \Rightarrow{}f(id)\cap{}f(id')=\emptyset
%     }
    
%     % Conclusion
%     \infer1
%     [$f\in{}ident\nrightarrow{}typedef$]
%     {
%       \mathcal{U}(f)
%     }
%   \end{prooftree}
% \end{figure}

% where $\mathtt{enum}$ is a predicate stating that a $typedef$ instance
% is an enumeration of values.

% \begin{figure}[H]
%   {\fontsize{10}{13}\selectfont\textsc{WFDesign}}
  
%   \begin{prooftree}

%     % Well-formed Types.
%     \hypo{\mathcal{U}(Types)}

%     % No intersection val(Types).
%     \infer[no rule]1{
%       \mathtt{val}(Types)\cap\mathtt{dom}(\Delta)=\emptyset
%     }
    
%     % No intersection dom(Types).
%     \infer[no rule]1{
%       \mathtt{dom}(Types)\cap\big((\mathtt{dom}(\Delta)\setminus\mathtt{dom}(Types))\big)=\emptyset
%     }

%     % No intersection Gens.
%     \infer[no rule]1{
%       \mathtt{dom}(Gens)\cap\big((\mathtt{dom}(\Delta)\setminus\mathtt{dom}(Gens))\cup\mathtt{val}(Types)\big)=\emptyset
%     }

%     % No intersection Ins.
%     \infer[no rule]1{
%       \mathtt{dom}(Ins)\cap\big((\mathtt{dom}(\Delta)\setminus\mathtt{dom}(Ins))\cup\mathtt{val}(Types)\big)=\emptyset
%     }

%     % No intersection Outs.
%     \infer[no rule]1{
%       \mathtt{dom}(Outs)\cap\big((\mathtt{dom}(\Delta)\setminus\mathtt{dom}(Outs))\cup\mathtt{val}(Types)\big)=\emptyset
%     }

%     % No intersection Sigs.
%     \infer[no rule]1{
%       \mathtt{dom}(Sigs)\cap\big((\mathtt{dom}(\Delta)\setminus\mathtt{dom}(Sigs))\cup\mathtt{val}(Types)\big)=\emptyset
%     }

%     % No intersection Consts.
%     \infer[no rule]1{
%       \mathtt{dom}(Consts)\cap\big((\mathtt{dom}(\Delta)\setminus\mathtt{dom}(Consts))\cup\mathtt{val}(Types)\big)=\emptyset
%     }
    
%     % No intersection Ps.
%     \infer[no rule]1{
%       \mathtt{dom}(Ps)\cap\big((\mathtt{dom}(\Delta)\setminus\mathtt{dom}(Ps))\cup\mathtt{val}(Types)\big)=\emptyset
%     }

%     % No intersection Comps.
%     \infer[no rule]1{
%       \mathtt{dom}(Comps)\cap\big((\mathtt{dom}(\Delta)\setminus\mathtt{dom}(Comps))\cup\mathtt{val}(Types)\big)=\emptyset
%     }
    
%     % Conclusion
%     \infer1
%     [$\Delta\in{}Design$]
%     {
%       \mathcal{U}(\Delta)
%     }
%   \end{prooftree}
% \end{figure}

% where $\Delta=(Gens,Ins,Outs,Sigs,Consts,Ps,Comps,Types,Behavior)$.

% \begin{figure}[H]
%   {\fontsize{10}{13}\selectfont\textsc{WFGammaDeltaEnv}}
%   \vspace{2ex}
  
%   \begin{prooftree}

%     % Well-formed .
%     \hypo{\mathcal{U}()}

%     % Well-formed \Delta.
%     \hypo{\mathcal{U}(\Delta)}

%     % No intersection  and \Delta.
%     \hypo{(\mathtt{dom}()\cup\mathtt{val}())
%       \cap(\mathtt{dom}(\Delta)\cup\mathtt{val}(Types))=\emptyset}
    
%     % Conclusion
%     \infer3
%     [{
%       \begin{tabular}{@{}l}
%         $\in{}ident\nrightarrow{}typedef$ \\
%         $\Delta\in{}Design$ \\
%       \end{tabular}
%     }]
%     {
%       \mathcal{U}(\Delta)
%     }
%   \end{prooftree}
% \end{figure}

% \begin{figure}[H]
%   {\fontsize{10}{13}\selectfont\textsc{WFGammaDeltaLambdaEnv}}
%   \vspace{2ex}
  
%   \begin{prooftree}

%     % Well-formed  and \Delta.
%     \hypo{\mathcal{U}(\Delta)}

%     % No intersection \Lambda and .
%     \hypo{\mathtt{dom}(\Lambda)\cap(\mathtt{dom}()\cup\mathtt{val}())=\emptyset}
    
%     % No intersection \Lambda and \Delta.
%     \hypo{\mathtt{dom}(\Lambda)\cap(\mathtt{dom}(\Delta)\cup\mathtt{val}(Types))=\emptyset}
    
%     % Conclusion
%     \infer3
%     [{
%       \begin{tabular}{@{}l}
%         $\in{}ident\nrightarrow{}typedef$ \\
%         $\Delta\in{}Design$ \\
%         $\Lambda\in{}ident\nrightarrow{}(type\times{}value)$
%       \end{tabular}
%     }]
%     {
%       \mathcal{U}(\Delta,\Lambda)
%     }
%   \end{prooftree}
% \end{figure}

% To safely pursue the simulation phase for a given top-level design, we
% need to prove Theorem~\ref{thm:id-unicity} stating that if the
% elaboration of the top-level design succeeds then there is no
% overlaping of identifiers between the $$ and $\Delta$ structures
% and every local environment of processes present in the top-level
% design's behavior.

% \begin{theorem}
%   \label{thm:id-unicity}
%   $\forall{}\mathrm{design},\mathcal{M},\Delta,$
%   $\big(\mathcal{D},\mathcal{M}\vdash\mathrm{design}\xrightarrow{elab}\Delta\big)$
%   $\Rightarrow\forall\Lambda,\exists\mathrm{id}\in\mathtt{dom}(Ps)|Ps(\mathrm{id})=\Lambda,~\mathcal{U}(\Delta,\Lambda)$.
% \end{theorem}

% Note that in $$ and $\mathcal{D}$ are free variables in
% Theorem~\ref{thm:id-unicity} as they are known structures.  With
% Theorem~\ref{thm:id-unicity} proved, they can be no ambiguity due to
% identifiers while interpreting expressions during the simulation
% phase.

\subsection{Full Simulation}

The full simulation process is decomposed in two steps. The first step
is the elaboration phase that builds an elaborated version of a
\hvhdl{} design along with its default state, and type-checks the
design. Previous to the elaboration phase, the top-level design
receives a value for each of its generic constant; we refer to it as
the \emph{dimensioning} of the top-level design. The second step is
the simulation phase that executes the behavioral part of the
top-level design starting from an initial state. The simulation is
decomposed into simulation cycles. Each simulation cycle is divided in
four parts entailed by the \emph{synchronous} execution of
$\mathcal{H}$-VHDL top-level designs, i.e designs whose behavior
depend on a clock signal. The four parts are, first, the execution of
concurrent statements responding to the rising edge of the clock
signal, then, a phase of signal stabilization followed by the
execution of concurrent statements responding to the falling edge of
the clock signal, and finally another phase of signal
stabilization. At each clock event, the value of the primary inputs of
the design being currrently simulated are captured and injected in the
simulation; primary inputs receive values from the design
environment. Here, the environment is represented by a function
mapping input port identifiers to values depending on the current
count of simulation cycles and the considered clock event. This leads
to the following hypothesis:

\begin{hypothesis}[Stable primary inputs]
  \label{hyp:stable-pi}
  The values of primary inputs (i.e, input ports of the top-level
  design) are captured at each clock event, and therefore are stable
  (i.e, their values do not change) between two contiguous clock
  events.
\end{hypothesis}

Hypothesis~\nameref{hyp:stable-pi} arises from the fact that the clock
signal sample rate respects the Nyquist-Shannon sampling
theorem. Therefore, the sample rate of the design's clock is
sufficient to capture all events possibly arising in the environment.
We only need to settle the values of the primary inputs at the clock
edges.

Also, after each clock event phase follows a signal stabilization
phase in the proceedings of a simulation cycle. One more hypothesis is
needed here:

\begin{hypothesis}[Stabilization]
  \label{hyp:stabilization}
  All signals have enough time to stabilize during the signal
  stabilization phase that happens between two clock events.
\end{hypothesis}

As a \hvhdl{} design represents a physical circuit, one can assume
that the represented circuit is analyzed former to the simulation.
Therefore, one knows exactly how much time is needed to propagate
signal values through the longest physical path; as a consequence, a
proper clock frequency is set ensuring signal stabilization between
two clock events. Thus, Hypothesis~\nameref{hyp:stabilization} arises
from the previous facts.

%%% DESIGN ELABORATION AND SIMULATION.

The $full$ simulation relation takes in parameter a top-level design
d, a design store $\mathcal{D}\in{}id\nrightarrow{}design$, an
elaborated design $\Delta\in{}ElDesign(d)$, a dimensioning function
$\mathcal{M}_g\in{}Gens(\Delta)\nrightarrow{}value$, a primary input
environment
$E_p\in{}(\mathbb{N}\times{}Clk)\rightarrow{}(Ins(\Delta)\rightarrow{}value)$,
a simulation cycle count $\tau\in\mathbb{N}$, and a simulation trace
$\theta\in{}\mathtt{list}(\Sigma(\Delta))$, corresponding to the list
of states yielded by design d after $\tau$ simulation cycles. Note
that we use the pointed notation to access the behavioral part of
design d, written d.cs. It is this part of the design that is executed
during the simulation, and therefore is passed as a parameter of the
initialization and simulation relations. % The states in
% $\sigma_t$ are \emph{time}-ordered, that is, the first states of the
% list are the earliest states in the simulation history. To achieve
% conciseness, further on a singleton trace containing the only element
% $\sigma$ is written $\sigma$. The trace resulting of the concatenation
% of two traces $\sigma_t$ and $\sigma_t'$ is written
% $\sigma_t.\sigma_t'$.

\begin{table}[H]
  {\fontsize{10}{13}\selectfont\textsc{FullSim}}
  
  \begin{prooftree}[template=\inserttext]

    % Design elab.
    \hypo{$\mathcal{D},\mathcal{M}_g\vdash\mathrm{d}\xrightarrow{elab}\Delta,\sigma$}

    % Initialization.
    \hypo{$\mathcal{D},\Delta,\sigma\vdash{}\mathrm{d.cs}\xrightarrow{init}\sigma_0$}
       
    % Simulation loop.
    \hypo{$\mathcal{D},E_p,\Delta,\tau,\sigma_0\vdash{}\mathrm{d.cs}\rightarrow\theta$}
    
    \infer3 [] { $\mathcal{D},\Delta,\mathcal{M}_g,E_p,\tau\vdash$
      $\mathrm{d}\xrightarrow{full}(\sigma_0::\theta)$ }
  \end{prooftree}
\end{table}

where:

\begin{itemize}[label=-]
\item $\mathcal{M}_g\in{}Gens(\Delta)\nrightarrow{}value$, the
  function yelding the values of generic constants for a given
  top-level design, refered to as the \emph{dimensioning} function.
\item
  $E_p\in{}(\mathbb{N}\times{}Clk)\rightarrow{}(ident\nrightarrow{}value)$,
  the function yelding a mapping from primary inputs (i.e, input ports
  of the top-level design) to values at a given simulation cycle count
  (i.e, the $\mathbb{N}$ argument), and a given clock event (i.e, the
  $Clk$ argument, where $Clk=\{\uparrow,\downarrow\}$).
\item $\tau$, the number of simulation cycles to execute.  The value
  of $\tau$ is decremented at each clock cycle until it reaches zero
  (see Section~\ref{sec:sim-loop}).
\end{itemize}

\subsection{Simulation loop.}
\label{sec:sim-loop}

\begin{table}[H]

  \begin{tabular}{@{}l}
    %%% SIMULATION END.
    
    {\fontsize{10}{13}\selectfont\textsc{SimEnd}} \\
    
    {\begin{prooftree}[template=\inserttext]
        \infer0 {
          $\mathcal{D},E_p,\Delta,0,\sigma\vdash{}cs\rightarrow{}[~]$
        }
      \end{prooftree}} \\
  \end{tabular}
  \begin{tabular}{l}
    % SIMULATION LOOP
    
    {\fontsize{10}{13}\selectfont\textsc{SimLoop}} \\
    
    {\begin{prooftree}[template=\inserttext]
        
        % First column.
        \hypo{$\mathcal{D},E_p,\Delta,\tau,\sigma\vdash{}cs\xrightarrow{\uparrow,\downarrow}\sigma',\sigma''$}

        % Second column.
        \hypo{$\mathcal{D},E_p,\Delta,\tau-1,\sigma''\vdash{}cs\rightarrow\theta$}
        
        \infer2 [$\tau>0$] {
          $\mathcal{D},E_p,\Delta,\tau,\sigma\vdash{}cs\rightarrow(\sigma'
          :: \sigma'' :: \theta)$ }
      \end{prooftree}} \\
  \end{tabular}

\end{table}

\subsection{Simulation cycle.}

 To ease the reading of
forward simulation rules, we need to introduce two notations.

\begin{notation}[Overriding union]
  For all partial function $f,f'\in{}X\nrightarrow{}Y$, $f\ocup{}f'$
  denotes the overriding union of $f$ and $f'$ such that
  $f\ocup{}f'(x)=
  \begin{cases}
    f'(x) & if~x\in\mathtt{dom}(f') \\
    f(x) & otherwise \\
  \end{cases}
  $
\end{notation}

\begin{notation}[Differentiated intersection domain]
  For all partial function $f,f'\in{}X\nrightarrow{}Y$, $f\dcap{}f'$
  denotes the intersection of the domain of $f$ and $f'$ for which $f$
  and $f'$ yields different values. That is,
  $f\dcap{}f'=\{~x\in\mathtt{dom}(f)\cap\mathtt{dom}(f')~|~f(x)\neq{}f'(x)~\}$.
\end{notation}

% \begin{remark}[Initialization and signal events.]
%   In the definition of the $\dcap$ relation, adding the difference
%   between the domain of $f'$ and $f$ is mandatory, as it permits to
%   take into account elements that have a fresh value in $f'$.  This
%   helps to catch events on signals that happened to have no values in
%   a signal store $\mathcal{S}$ but have a fresh (initial) value in a
%   resulting signal store $\mathcal{S}'$.
% \end{remark}

\begin{definition}[Input port values update]
  Given an \hvhdl{} design $d\in{}design$, a design store
  $\mathcal{D}\in{}id\rightarrow{}design$, an elaborated design
  $\Delta\in{}ElDesign(d,\mathcal{D})$, a simulation environment
  $E_p\in{}(\mathbb{N}\times\{\uparrow,\downarrow\})\rightarrow(Ins(\Delta)\rightarrow{}value)$,
  let us define the relation expressing the update of the values of
  the input ports of $\Delta$ at a given design state
  $\sigma\in\Sigma(\Delta)$, clock cycle count $\tau\in{}\mathbb{N}$,
  and clock event $clk\in\{\uparrow,\downarrow\}$, and thus resulting
  in a new state $\sigma_i\in\Sigma(\Delta)$. The relation is written
  $\mathtt{Inject}_{clk}(\sigma,E_p,\tau,\sigma_i)$ and verifies that:
  $\sigma={<}\mathcal{S},\mathcal{C},\mathcal{E}{>}$ and
  $\sigma_i={<}\mathcal{S}\ocup{}E_p(\tau,clk),\mathcal{C},\mathcal{E}{>}$.
\end{definition}

The cycle relation states that the design state $\sigma''$ is the
result of the execution of concurrent statement $cs$ over one
simulation cycle, this starting from state $\sigma$.  As told in
Hypothesis~\ref{hyp:stable-pi}, we update the input port values at
each clock event. New input port values are coming from the
environment $E_p$. The updates are made in the definitions of
$\sigma_i$ and $\sigma'_i$. These definitions are expressed as side
conditions.

\begin{table}[H]
  {\fontsize{10}{13}\selectfont\textsc{SimCyc}}
  
  \begin{prooftree}[template=\inserttext]

    % Rising
    \hypo{$\mathcal{D},\Delta,\sigma_i\vdash\mathrm{cs}\xrightarrow{\uparrow}\sigma_\uparrow$}
    
    % Falling
    \infer[no rule]1{$\mathcal{D},\Delta,\sigma'_i\vdash\mathrm{cs}\xrightarrow{\downarrow}\sigma_\downarrow$}

    % Stabilize after rising.
    \hypo{$\mathcal{D},\Delta,\sigma_\uparrow\vdash\mathrm{cs}\xrightarrow{\rightsquigarrow}\sigma'$}

    % Stabilize after falling.
    \infer[no rule]1{$\mathcal{D},\Delta,\sigma_\downarrow\vdash\mathrm{cs}\xrightarrow{\rightsquigarrow}\sigma''$}
    
    \infer2
    [{
      \renewcommand{\arraystretch}{1.5}
      \begin{tabular}{@{}l}
        $\mathtt{Inject}_\uparrow(\sigma,E_p,\tau,\sigma_i)$\\
        $\mathtt{Inject}_\downarrow(\sigma',E_p,\tau,\sigma'_i)$\\
      \end{tabular}
    }] {
      $\mathcal{D},E_p,\Delta,\tau,\sigma\vdash\mathrm{cs}\xrightarrow{\uparrow,\downarrow}$
      $\sigma',\sigma''$ }
  \end{prooftree}
\end{table}

\begin{remark}[Input ports and signal store]
  For a given $\Delta\in{}Design$, $\sigma\in\Sigma(\Delta)$,
  $E_p\in{}\mathbb{N}\rightarrow{}Clk\rightarrow{}(Ins(\Delta)\rightarrow{}value)\mathbb{}$,
  $\tau\in\mathbb{N}$, $clk\in{}Clk$, we have
  $\mathtt{dom}(E_p(\tau,clk))\subseteq{}\mathtt{dom}(\mathcal{S})$,
  where $\sigma={<}\mathcal{S}, \mathcal{C}, \mathcal{E}{>}$. Indeed,
  the input ports of $\Delta$ that constitutes the domain of
  $E_p(\tau,clk)$ are a subset of the set of signals. The set of
  signals constitutes the domain of the signal store of $\sigma$ (i.e,
  $\mathcal{S}$); thus we have
  $\mathtt{dom}(E_p(\tau,clk))\subseteq{}\mathtt{dom}(\mathcal{S})$.
\end{remark}

\subsection{Initialization rules}

\begin{table}[H]
  {\fontsize{10}{13}\selectfont\textsc{Init}}
  
  \begin{prooftree}[template=\inserttext]

    % Run all processes once.
    \hypo{$\mathcal{D},\Delta,\sigma\vdash\mathrm{cs}\xrightarrow{runinit}{}\sigma'$}

    % Stabilization phase after runinit.
    \hypo{$\mathcal{D},\Delta,\sigma'\vdash\mathrm{cs}\xrightarrow{\rightsquigarrow}{}\sigma''$}

    % Conclusion.
    \infer2
    {
      $\mathcal{D},\Delta,\sigma\vdash\mathrm{cs}\xrightarrow{init}\sigma''$
    }

  \end{prooftree}
\end{table}

At the initialization phase, the block of sequential instructions of
all processes is executed exactly once ($runinit$), then a
stabilization phase follows ($stabilize$). It is during the
initialization phase that the first part of $\mathtt{rst}$ blocks is
executed. A block (\vhdle|rst| ss ss') is equivalent to (\vhdle|if|
(rst = false) \vhdle|then| ss \vhdle|else| ss') where \texttt{rst} is
a reserved signal identifier. Therefore, when considering a
\texttt{rst} block, the $runinit$ relation executes the ss block; at
every other moment of the simulation, the ss' block is executed. This
mimicks the conventional proceeding of a simulation where the
\texttt{rst} signal (for \textit{reset} signal) is set to false during
the initialization (only during the $runinit$ phase, not during the
stabilization phase), and then is set to true for the rest of the
simulation.

\subsubsection{Evaluation of a process statement}
\label{subsubsec:ps-stmt-init}

\begin{premises}
  \begin{itemize}
  \item The $i$ flag of the $ss_i$ relation indicates that all
    sequential instructions responding to the initialization phase
    (i.e, \texttt{rst} blocks) will be executed.
  \item The set of events of state $\sigma$ is emptied
    ($NoEv(\sigma)$) before the evaluation of the process statement
    body, and the resulting state is the starting state that will be
    written through the execution of the process statement body.
\end{itemize}
\end{premises}

\begin{table}[H]
  \centering
  \begin{tabular}{@{}l}
    {\fontsize{10}{13}\selectfont\textsc{PsRunInit}} \\
    {\begin{prooftree}[template=\inserttext]
        
        % Runs seq.
        \hypo{$\Delta,\sigma,NoEv(\sigma),\Lambda\vdash\mathrm{ss}\xrightarrow{ss_i}{}\sigma',\Lambda'$}

        % Conlcusion.
        \infer1
        [{
          \begin{tabular}{@{}l}
            $\Delta(\mathrm{id}_p)=\Lambda$\\
          \end{tabular}
        }]
        {
          $\mathcal{D},\Delta,\sigma\vdash$
          \vhdle|process| \texttt{(}id$_p$\texttt{,} sl\texttt{,} vars\texttt{,} ss\texttt{)}
          $\xrightarrow{runinit}\sigma'$
        }
      \end{prooftree}} \\
  \end{tabular}
\end{table}

\subsubsection{Evaluation of a component instantiation statement}
\label{subsubsec:ps-stmt-init}

\begin{sideconds}
  If $\sigma_c''$ has a non-empty set of events, then the state of
  component $\mathrm{id}_c$ must be overidden in the target state of
  the $runinit$ relation. Moreover, $\mathrm{id}_c$ must be added to
  the set of events of the embedding design. This is expressed by the
  side condition:
  $\sigma''=(\sigma'(\mathrm{id}_c)\leftarrow\sigma_c'')\cup_\mathcal{E}\{\mathrm{id}_c\}$
\end{sideconds}

\begin{table}[H]
  \centering
  \begin{tabular}{@{}l}
    {\fontsize{10}{13}\selectfont\textsc{CompRunInit}} \\
    {\begin{prooftree}
        
        % Builds mapping for in ports.
        \hypo{\Delta,\Delta_c,\sigma,\sigma_c\vdash\mathrm{i}\xrightarrow{mapip}\sigma'_c}
        
        % Executes runinit on component behavior.
        \infer[no rule]1{\mathcal{D},\Delta_c,\sigma_c'\vdash{}\mathcal{D}(\mathrm{id_e}).\mathrm{cs}\xrightarrow{runinit}\sigma_c''}
        
        % Builds mapping for out ports.
        \infer[no rule]1{
          \Delta,\Delta_c,NoEv(\sigma),\sigma_c''\vdash
          \mathrm{o}
          \xrightarrow{mapop}
          \sigma'
        }
        
        % Conclusion.
        \infer1
        [{\renewcommand{\arraystretch}{1.5}
          \begin{tabular}{@{}l}
            $id_e\in\mathcal{D}$ \\
            $\Delta(\mathrm{id}_c)=\Delta_c$, $\sigma(\mathrm{id}_c)=\sigma_c$ \\
            $\sigma''={<}\mathcal{S}',\mathcal{C}'(id_c)\leftarrow\sigma_c'',\mathcal{E}'\cup(\mathcal{C}\dcap\mathcal{C}'){>}$ \\
          \end{tabular}
        }] {
          \mathcal{D},\Delta,\sigma\vdash~
          $\vhdle|comp|$~(\mathrm{id}_c, \mathrm{id}_e, \mathrm{g},
          \mathrm{i}, \mathrm{o})
          \xrightarrow{runinit}{}\sigma''
        }
      \end{prooftree}} \\
  \end{tabular}
\end{table}

\subsubsection{Evaluation of the composition of concurrent statements}
\label{subsubsec:ps-stmt-init}

\begin{table}[H]
  \begin{tabular}{l}
    {\fontsize{10}{13}\selectfont\textsc{ParRunInit}} \\
    {\begin{prooftree}
        \hypo{\mathcal{D},\Delta,\sigma\vdash\mathrm{cs}\xrightarrow{runinit}\sigma'}
        \hypo{\mathcal{D},\Delta,\sigma\vdash\mathrm{cs'}\xrightarrow{runinit}\sigma''}
        \infer2
        [{\renewcommand{\arraystretch}{1.5}
          \begin{tabular}{@{}l@{}}
            $\mathcal{E}'\cap\mathcal{E}''=\emptyset$ \\
          \end{tabular}
        }] {
          \mathcal{D},\Delta,\sigma\vdash\mathrm{cs}~\mathtt{||}~\mathrm{cs'}\xrightarrow{runinit}\mathtt{merge}(\sigma,\sigma',\sigma'')
        }
      \end{prooftree}} \\
  \end{tabular}
  \begin{tabular}{l}
    {\fontsize{10}{13}\selectfont\textsc{NullRunInit}} \\
    {\begin{prooftree}
        \infer0
        {
          \Delta,\sigma\vdash\mathtt{null}\xrightarrow{runinit}NoEv(\sigma)
        }
      \end{prooftree}} \\
  \end{tabular}
\end{table}

The \texttt{merge} function computes a new state based on the original
state \textit{o}, and the states \textit{s} and \textit{s'} yielded by
the computation of two concurrent statements. In the resulting state,
the signal value store $\mathcal{S}_m$ is a function merging together
the signal store functions at state $o$, $s$ and $s'$. $S_m$ yields
values from the signal store $\mathcal{S}$ (resp. $\mathcal{S}'$) for
all signal that belongs to the set of events at state $s$
(resp. $s'$), and yields values from the original signal store
$\mathcal{S}_o$ for all unchanged signals. The same goes for the
resulting component instance state store $C_m$. The new set of events
$\mathcal{E}_m$ is the union between set of events at state $s$ and
$s'$. The \texttt{merge} correctly merges the state $o$, $s$ and $s'$
only if the set of events of $s$ and $s'$ are disjoint. Fortunately,
the \textsc{ParRunInit} rule that calls to the \texttt{merge} function
defines the condition of disjoint set of events as a side condition.

\begin{table}[H]
\begin{lstlisting}[language=PseudoCoq]
Definition merge(o,s,s') :=
   let o = ($\mathcal{S}_o$,$\mathcal{C}_o$,$\mathcal{E}_o$) in
   let s = ($\mathcal{S}$,$\mathcal{C}$,$\mathcal{E}$) in
   let s' = ($\mathcal{S}'$,$\mathcal{C}'$,$\mathcal{E}'$) in
   let $\mathcal{S}_m$ = $\lambda\mathrm{id}.$ if $\mathrm{id}\in\mathcal{E}$ then $\mathcal{S}(\mathrm{id})$ else if $\mathrm{id}\in\mathcal{E}'$ then $\mathcal{S}'(\mathrm{id})$ else $\mathcal{S}_o(\mathrm{id})$
   let $\mathcal{C}_m$ = $\lambda\mathrm{id}.$ if $\mathrm{id}\in\mathcal{E}$ then $\mathcal{C}(\mathrm{id})$ else if $\mathrm{id}\in\mathcal{E}'$ then $\mathcal{C}'(\mathrm{id})$ else $\mathcal{C}_o(\mathrm{id})$
   let $\mathcal{E}_m$ = $\mathcal{E}\cup\mathcal{E}'$ in ($\mathcal{S}_m$,$\mathcal{C}_m$,$\mathcal{E}_m$).
\end{lstlisting}
\end{table}

\begin{remark}[No multiply-driven signals]
  For all states $\sigma=(\mathcal{S},\mathcal{C},\mathcal{E})$ and
  $\sigma'=(\mathcal{S}',\mathcal{C}',\mathcal{E}')$ resulting from
  the execution of two concurrent statements cs and cs',
  $\mathcal{E}\cap\mathcal{E}'=\emptyset$. Otherwise, there exists
  some multiply-driven signals, which are forbidden in our semantics.
\end{remark}

\subsection{Clock phases rules}

The following rules express the evaluation of concurrent statements at
clock phases, i.e, the $\uparrow$ and $\downarrow$ phases. The clock
signal, trigerring the evaluation of synchronous process statements,
is represented by the reserved signal identifier \texttt{clk}. Thus,
synchronous processes are processes containing the \texttt{clk} in
their sensitivity list.

\subsubsection{Evaluation of a process statement}

\begin{table}[H]
  {\fontsize{10}{13}\selectfont\textsc{PsRENoClk}}
  
  \begin{prooftree}[template=\inserttext]
    \hypo{}

    \infer1
    [$\mathtt{clk}\notin\mathrm{sl}$]
    {
      $\mathcal{D},\Delta,\sigma\vdash$
      \vhdle|process| \texttt{(}id$_p$\texttt{,} sl\texttt{,} vars\texttt{,} ss\texttt{)}
      $\xrightarrow{\uparrow}$
      $\sigma$
    }
  \end{prooftree}
\end{table}

\begin{premises}
  The $\uparrow$ flag in the $ss_\uparrow$ relation indicates that
  $\mathtt{rising}$ blocks will be executed.
\end{premises}

\begin{table}[H]
  {\fontsize{10}{13}\selectfont\textsc{PsREClk}}
  
  \begin{prooftree}[template=\inserttext]

    \hypo{$\Delta,\sigma,NoEv(\sigma),\Lambda\vdash\mathrm{ss}\xrightarrow{ss_\uparrow}{}\sigma',\Lambda'$}
    
    \infer1
    [{
      \begin{tabular}{@{}l}
        $\mathtt{clk}\in\mathrm{sl}$ \\
        $\Delta(\mathrm{id}_p)=\Lambda$ \\
      \end{tabular}
    }]
    {
      $\mathcal{D},\Delta,\sigma\vdash$
      \vhdle|process| \texttt{(}id$_p$\texttt{,} sl\texttt{,} vars\texttt{,} ss\texttt{)}
      $\xrightarrow{\uparrow}$
      $\sigma'$
    }
  \end{prooftree}
\end{table}

\begin{table}[H]
  {\fontsize{10}{13}\selectfont\textsc{PsFENoClk}}
  
  \begin{prooftree}[template=\inserttext]
    \hypo{}

    \infer1
    [$\mathtt{clk}\notin\mathrm{sl}$]
    {
      $\mathcal{D},\Delta,\sigma
      \vdash$
      \vhdle|process| \texttt{(}id$_p$\texttt{,} sl\texttt{,} vars\texttt{,} ss\texttt{)}
      $\xrightarrow{\downarrow}\sigma$
    }
  \end{prooftree}
\end{table}

\begin{premises}
  The $\downarrow$ flag in the $ss_\downarrow$ relation indicates that
  $\mathtt{falling}$ blocks will be executed.
\end{premises}

\begin{table}[H]
  {\fontsize{10}{13}\selectfont\textsc{PsFEClk}}
  
  \begin{prooftree}[template=\inserttext]

    \hypo{$\Delta,\sigma,NoEv(\sigma),\Lambda\vdash\mathrm{ss}\xrightarrow{ss_\downarrow}{}\sigma',\Lambda'$}
    
    \infer1
    [{
      \begin{tabular}{@{}l}
        $\mathtt{clk}\in\mathrm{sl}$ \\
        $\Delta(\mathrm{id}_p)=\Lambda$ \\
      \end{tabular}
    }]
    {
      $\mathcal{D},\Delta,\sigma
      \vdash$
      \vhdle|process| \texttt{(}id$_p$\texttt{,} sl\texttt{,} vars\texttt{,} ss\texttt{)}
      $\xrightarrow{\downarrow}$
      $\sigma'$
    }
  \end{prooftree}
\end{table}

\subsubsection{Evaluation of a component instantiation statement}
\label{subsubsec:sync-comp-inst}

\begin{table}[H]
  \centering
  \begin{tabular}{@{}l}
    {\fontsize{10}{13}\selectfont\textsc{CompRE}} \\
    {\begin{prooftree}[template=\inserttext]
        
        % Builds mapping for in ports.
        \hypo{$\Delta,\Delta_c,\sigma,\sigma_c\vdash\mathrm{i}\xrightarrow{mapip}\sigma'_c$}
        
        % Executes rising on component behavior.
        \infer[no rule]1{$\mathcal{D},\Delta_c,\sigma_c'\vdash{}\mathcal{D}(id_e).\mathrm{cs}\xrightarrow{\uparrow}\sigma_c''$}
        
        % Builds mapping for out ports.
        \infer[no rule]1{
          $\Delta,\Delta_c,\sigma,\sigma_c''\vdash$
          $\mathrm{o}$
          $\xrightarrow{mapop}$
          $\sigma'$
        }
        
        % Conclusion.
        \infer1
        [{\renewcommand{\arraystretch}{1.5}
          \begin{tabular}{@{}l}
            $id_e\in\mathcal{D}$ \\
            $\Delta(\mathrm{id}_c)=\Delta_c$, $\sigma(\mathrm{id}_c)=\sigma_c$ \\
            $\sigma''={<}\mathcal{S}',\mathcal{C}'(id_c)\leftarrow\sigma_c'',\mathcal{E}'\cup(\mathcal{C}\dcap\mathcal{C}'){>}$ \\
          \end{tabular}
        }]
        {
          $\mathcal{D},\Delta,\sigma\vdash$
          \vhdle|comp| (id$_c$, id$_e$, g, i, o)
          $\xrightarrow{\uparrow}{}\sigma''$
        }
      \end{prooftree}} \\
  \end{tabular}
\end{table}

\begin{table}[H]
  \centering
  \begin{tabular}{@{}l}
    {\fontsize{10}{13}\selectfont\textsc{CompFE}} \\
    {\begin{prooftree}[template=\inserttext]
        
        % Builds mapping for in ports.
        \hypo{$\Delta,\Delta_c,\sigma,\sigma_c\vdash\mathrm{i}\xrightarrow{mapip}\sigma'_c$}
        
        % Executes rising on component behavior.
        \infer[no rule]1{$\mathcal{D},\Delta_c,\sigma_c'\vdash{}\mathcal{D}(id_e).\mathrm{cs}\xrightarrow{\downarrow}\sigma_c''$}
        
        % Builds mapping for out ports.
        \infer[no rule]1{
          $\Delta,\Delta_c,\sigma,\sigma_c''\vdash$
          $\mathrm{o}$
          $\xrightarrow{mapop}$
          $\sigma'$
        }
        
        % Conclusion.
        \infer1
        [{\renewcommand{\arraystretch}{1.5}
          \begin{tabular}{@{}l}
            $id_e\in\mathcal{D}$ \\
            $\Delta(\mathrm{id}_c)=\Delta_c$, $\sigma(\mathrm{id}_c)=\sigma_c$ \\
            $\sigma''={<}\mathcal{S}',\mathcal{C}'(id_c)\leftarrow\sigma_c'',\mathcal{E}'\cup(\mathcal{C}\dcap\mathcal{C}'){>}$ \\
          \end{tabular}
        }]
        {
          $\mathcal{D},\Delta,\sigma\vdash$
          \vhdle|comp| (id$_c$, id$_e$, g, i, o)
          $\xrightarrow{\downarrow}{}\sigma''$
        }
      \end{prooftree}} \\
  \end{tabular}
\end{table}

\subsubsection{Evaluation of the composition of concurrent statements}
\label{subsubsec:comp-of-cs-clk}

\begin{table}[H]
  \begin{tabular}{l}
    {\fontsize{10}{13}\selectfont\textsc{ParFE}} \\
    {\begin{prooftree}
        \hypo{\mathcal{D},\Delta,\sigma\vdash\mathrm{cs}\xrightarrow{\downarrow}\sigma'}
        \hypo{\mathcal{D},\Delta,\sigma\vdash\mathrm{cs'}\xrightarrow{\downarrow}\sigma''}
        \infer2
        [{\renewcommand{\arraystretch}{1.5}
          \begin{tabular}{@{}l@{}}
            $\mathcal{E}'\cap\mathcal{E}''=\emptyset$ \\
          \end{tabular}
        }] {
          \mathcal{D},\Delta,\sigma\vdash\mathrm{cs}~\mathtt{||}~\mathrm{cs'}\xrightarrow{\downarrow}
          \mathtt{merge}(\sigma,\sigma',\sigma'')
        }
      \end{prooftree}} \\
  \end{tabular}
  \begin{tabular}{l}
    {\fontsize{10}{13}\selectfont\textsc{NullFE}} \\
    {\begin{prooftree}
        \infer0
        {
          \Delta,\sigma\vdash\mathtt{null}\xrightarrow{\downarrow}\sigma
        }
      \end{prooftree}} \\
  \end{tabular}
\end{table}

\begin{table}[H]
  \begin{tabular}{l}
    {\fontsize{10}{13}\selectfont\textsc{ParRE}} \\
    {\begin{prooftree}
        \hypo{\mathcal{D},\Delta,\sigma\vdash\mathrm{cs}\xrightarrow{\uparrow}\sigma'}
        \hypo{\mathcal{D},\Delta,\sigma\vdash\mathrm{cs'}\xrightarrow{\uparrow}\sigma''}
        \infer2
        [{\renewcommand{\arraystretch}{1.5}
          \begin{tabular}{@{}l@{}}
            $\mathcal{E}'\cap\mathcal{E}''=\emptyset$ \\
          \end{tabular}
        }] {
          \mathcal{D},\Delta,\sigma\vdash\mathrm{cs}~\mathtt{||}~\mathrm{cs'}\xrightarrow{\uparrow}
          \mathtt{merge}(\sigma,\sigma',\sigma'')
        }
      \end{prooftree}} \\
  \end{tabular}
  \begin{tabular}{l}
    {\fontsize{10}{13}\selectfont\textsc{NullRE}} \\
    {\begin{prooftree}
        \infer0
        {\Delta,\sigma\vdash\mathtt{null}\xrightarrow{\uparrow}\sigma}
      \end{prooftree}} \\
  \end{tabular}
\end{table}

\subsection{Stabilization rules}
\label{sec:stab-rules}

\begin{table}[H]
  \begin{tabular}{@{}l}
    {\fontsize{10}{13}\selectfont\textsc{StabilizeEnd}} \\
    
    {\begin{prooftree}[template=\inserttext]

        \hypo{$\mathcal{D},\Delta,\sigma\vdash\mathrm{cs}\xrightarrow{comb}\sigma$}
        \infer1
        [$\mathcal{E}=\emptyset$]
        {
          $\mathcal{D},\Delta,\sigma\vdash\mathrm{cs}\xrightarrow{\rightsquigarrow}\sigma$
        }
      \end{prooftree}} \\
  \end{tabular}
\end{table}

\begin{table}[H]  
  \begin{tabular}{@{}l}
    {\fontsize{10}{13}\selectfont\textsc{StabilizeLoop}} \\
    {\begin{prooftree}[template=\inserttext]
        \hypo{$\mathcal{D},\Delta,\sigma\vdash\mathrm{cs}\xrightarrow{comb}\sigma'$}
        \hypo{$\mathcal{D},\Delta,\sigma'\vdash\mathrm{cs}\xrightarrow{\rightsquigarrow}\sigma''$}
        \infer2
        [{
          \begin{tabular}{@{}l}
            $\mathcal{E}\neq\emptyset$ \\
            $\mathcal{E}''=\emptyset$
          \end{tabular}
        }]
        {
          $\mathcal{D},\Delta,\sigma\vdash\mathrm{cs}\xrightarrow{\rightsquigarrow}\sigma''$
        }
      \end{prooftree}} \\
  \end{tabular}
\end{table}

\subsubsection{Evaluation of a process statement}

% \begin{table}[H]
%   {\fontsize{10}{13}\selectfont\textsc{PsCombStable}}
  
%   \begin{prooftree}[template=\inserttext]
%     \infer0
%     [{
%       \begin{tabular}{@{}l}
%         $\sigma={<}\mathcal{S},\mathcal{C},\mathcal{E}{>}$\\
%         $\mathrm{sl}\cap\mathcal{E}=\emptyset$
%       \end{tabular}
%     }]
%     {
%       $\mathcal{D},\Delta,\sigma\vdash$
%       \vhdle|process| \texttt{(}id$_p$\texttt{,} sl\texttt{,} vars\texttt{,} ss\texttt{)}
%       $\xrightarrow{comb}NoEv(\sigma)$
%     }
%   \end{prooftree}
% \end{table}

\begin{premises}
  The $c$ flag (for \textit{combinational}) on the $ss_c$ relation
  indicates that instructions responding to clock events
  (\texttt{falling} and \texttt{rising} blocks) and instructions
  executed during the initialization phase only (\texttt{rst} blocks)
  will not be considered.
\end{premises}

\begin{table}[H]
  {\fontsize{10}{13}\selectfont\textsc{PsComb}}
  
  \begin{prooftree}[template=\inserttext]
    
    % Executes ss.
    \hypo{$\Delta,\sigma,NoEv(\sigma),\Lambda\vdash\mathrm{ss}\xrightarrow{ss_c}{}\sigma',\Lambda'$}

    % Conclusion.
    \infer1
    [{
      \begin{tabular}{@{}l}
        $\Delta(\mathrm{id}_p)=\Lambda$ \\
      \end{tabular}
    }]
    {
      $\mathcal{D},\Delta,\sigma\vdash$
      \vhdle|process| \texttt{(}id$_p$\texttt{,} sl\texttt{,} vars\texttt{,} ss\texttt{)}
      $\xrightarrow{comb}\sigma'$
    }
  \end{prooftree}
\end{table}

\subsubsection{Evaluation of a component instantiation statement}

\begin{table}[H]
  \centering
  \begin{tabular}{@{}l}
    {\fontsize{10}{13}\selectfont\textsc{CompComb}} \\
    {\begin{prooftree}[template=\inserttext]
        
        % Builds mapping for in ports.
        \hypo{$\Delta,\Delta_c,\sigma,\sigma_c\vdash\mathrm{i}\xrightarrow{mapip}\sigma'_c$}
        
        % Executes comb on component behavior.
        \infer[no rule]1{$\mathcal{D},\Delta_c,\sigma_c'\vdash{}\mathcal{D}(id_e).\mathrm{cs}\xrightarrow{comb}\sigma_c''$}
        
        % Builds mapping for out ports.
        \infer[no rule]1{
          $\Delta,\Delta_c,NoEv(\sigma),\sigma_c''\vdash$
          $\mathrm{o}$
          $\xrightarrow{mapop}$
          $\sigma'$
        }
        
        % Conclusion.
        \infer1
        [{\renewcommand{\arraystretch}{1.5}
          \begin{tabular}{@{}l}
            $id_e\in\mathcal{D}$ \\
            $\Delta(\mathrm{id}_c)=\Delta_c$, $\sigma(\mathrm{id}_c)=\sigma_c$ \\
            $\sigma''={<}\mathcal{S}',\mathcal{C}'(id_c)\leftarrow\sigma_c'',\mathcal{E}'\cup(\mathcal{C}\dcap\mathcal{C}'){>}$ \\
          \end{tabular}
        }]
        {
          $\mathcal{D},\Delta,\sigma\vdash$
          \vhdle|comp| (id$_c$, id$_e$, g, i, o)
          $\xrightarrow{comb}$
          $\sigma''$
        }
      \end{prooftree}} \\
  \end{tabular}
\end{table}

\subsubsection{Evaluation of the composition of concurrent statements}

\begin{table}[H]
  \begin{tabular}{l}
    {\fontsize{10}{13}\selectfont\textsc{ParComb}} \\
    {\begin{prooftree}
        \hypo{\mathcal{D},\Delta,\sigma\vdash\mathrm{cs}\xrightarrow{comb}\sigma'}
        \hypo{\mathcal{D},\Delta,\sigma\vdash\mathrm{cs'}\xrightarrow{comb}\sigma''}
        \infer2
        [{\renewcommand{\arraystretch}{1.5}
          \begin{tabular}{@{}l@{}}
            $\mathcal{E}'\cap\mathcal{E}''=\emptyset$ \\
          \end{tabular}
        }] {
          \mathcal{D},\Delta,\sigma\vdash\mathrm{cs}~\mathtt{||}~\mathrm{cs'}\xrightarrow{comb}
          \mathtt{merge}(\sigma,\sigma',\sigma'')
        }
      \end{prooftree}} \\
  \end{tabular}
  \begin{tabular}{l}
    {\fontsize{10}{13}\selectfont\textsc{NullComb}} \\
    {\begin{prooftree}
        \infer0
        {
          \Delta,\sigma\vdash\mathtt{null}\xrightarrow{\downarrow}NoEv(\sigma)
        }
      \end{prooftree}} \\
  \end{tabular}
\end{table}

\subsection{Evaluation of input and output port maps}
\label{subsec:mapinout}

\begin{table}[H]
  \begin{tabular}{@{}l}
    {\fontsize{10}{13}\selectfont\textsc{MapipSimple}} \\
    {\begin{prooftree}

        % Evaluates e.
        \hypo{\Delta,\sigma\vdash\mathrm{e}\xrightarrow{e}v}

        % Checks that v complies with T.
        \hypo{v\in_c{}T}

        % Conclusion.
        \infer2
        [{\renewcommand{\arraystretch}{1.5}
          \begin{tabular}{@{}l}
            $\Delta_c(\mathrm{id}_s)=T$ \\
            $\sigma_c={<}\mathcal{S},\mathcal{C},\mathcal{E}{>}$ \\
            $\mathcal{S}'=\mathcal{S}(id_s)\leftarrow{}v$ \\
          \end{tabular}
        }]
        {
          \Delta,\Delta_c,\sigma,\sigma_c\vdash
          \mathrm{id}_s\Rightarrow\mathrm{e}
          \xrightarrow{mapip}
          {<}\mathcal{S}',\mathcal{C},\mathcal{E}{>}
        }
      \end{prooftree}} \\
  \end{tabular}
\end{table}

\begin{table}[H]
  \begin{tabular}{@{}l}
    {\fontsize{10}{13}\selectfont\textsc{MapipPartial}} \\
    {\begin{prooftree}

        % Evaluates e_i and e.
        \hypo{\Delta,\sigma&\vdash\mathrm{e}\xrightarrow{e}v}
        \infer[no rule]1{&\vdash\mathrm{e}_i\xrightarrow{e}v_i}        

        % Checks that v complies with T.
        \hypo{v&\in_c{}T}
        \infer[no rule]1{v_i&\in_c{}\mathtt{nat}(n,m)}

        % Conclusion.
        \infer2
        [{\renewcommand{\arraystretch}{1.5}
          \begin{tabular}{@{}l}
            $\Delta_c(\mathrm{id}_s)=\mathtt{array(}T,n,m\mathtt{)}$ \\
            $\sigma_c={<}\mathcal{S},\mathcal{C},\mathcal{E}{>}$ \\
            $\mathcal{S}'=\mathcal{S}(id_s)\leftarrow\mathtt{set\_at}(v,v_i,\mathcal{S}(id_s))$ \\
          \end{tabular}
        }]
        {
          \Delta,\Delta_c,\sigma,\sigma_c\vdash
          \mathrm{id}_s(\mathrm{e}_i)\Rightarrow\mathrm{e}
          \xrightarrow{mapip}
          {<}\mathcal{S}',\mathcal{C},\mathcal{E}{>}
        }
      \end{prooftree}} \\
  \end{tabular}
\end{table}

\begin{table}[H]
  {\fontsize{10}{13}\selectfont\textsc{MapipComp}}
  
  \begin{prooftree}    
    % Evaluates mapip on assoc_p.
    \hypo{
      \Delta,\Delta_c,\sigma,\sigma_c\vdash
      \mathrm{assoc}_{ip} 
      \xrightarrow{mapip}
      \sigma'_c
    }

    % Evaluates mapip on portmap.
    \hypo{
      \Delta,\Delta_c,\sigma,\sigma'_c\vdash
      \mathrm{ipmap}
      \xrightarrow{mapip}
      \sigma''_c
    }

    % Conclusion.
    \infer2
    {
      \Delta,\Delta_c,\sigma,\sigma_c\vdash
      \langle{}\mathrm{assoc}_{ip},~\mathrm{ipmap}\rangle
      \xrightarrow{mapip}
      \sigma''_c
    }
  \end{prooftree}
\end{table}

\begin{remark}[Out ports and $e$]
  We can not use the $e$ relation to interpret the values of output
  ports, because output ports are write-only constructs. We append the
  flag $o$ to the $e$ relation (i.e, $e_o$) to permit the
  evaluation of output port identifiers as regular signal identifier
  expressions.
\end{remark}

\begin{table}[H]
  \begin{tabular}{@{}l}
    {\fontsize{10}{13}\selectfont\textsc{MapopOpen}} \\
    {\begin{prooftree}

        % Conclusion
        \infer
        0
        {
          \Delta,\Delta_c,\sigma,\sigma_c\vdash
          \mathrm{id}_f\Rightarrow\mathtt{open}
          \xrightarrow{mapop}
          \sigma_c
        }
      \end{prooftree}} \\
  \end{tabular}
\end{table}

\begin{table}[H]
  \begin{tabular}{@{}l}
    {\fontsize{10}{13}\selectfont\textsc{MapopSimpleToSimple}} \\
    {\begin{prooftree}

        % Evaluates name.
        \hypo{\Delta_c,\sigma_c\vdash\mathrm{id}_f\xrightarrow{e_o}v}

        % Checks that v complies with T.
        \hypo{v\in_c{}T}
        
        % Conclusion.
        \infer2
        [{\renewcommand{\arraystretch}{1.5}
          \begin{tabular}{@{}l}
            $\mathrm{id}_a\in{}Sigs(\Delta)\cup{}Outs(\Delta)$ \\
            $\Delta(\mathrm{id}_a)=T$ \\
            $\sigma={<}\mathcal{S},\mathcal{C},\mathcal{E}{>}$ \\
            $\mathcal{S}'=\mathcal{S}(id_a)\leftarrow{}v$, $\mathcal{E}'=\mathcal{E}\cup(\mathcal{S}\dcap\mathcal{S}')$ \\
          \end{tabular}
        }]
        {
          \Delta,\Delta_c,\sigma,\sigma_c\vdash
          \mathrm{id}_f\Rightarrow\mathrm{id}_a
          \xrightarrow{mapop}
          {<}\mathcal{S}',\mathcal{C},\mathcal{E}'{>}
        }
      \end{prooftree}} \\
  \end{tabular}
\end{table}

\begin{table}[H]
  \begin{tabular}{@{}l}
    {\fontsize{10}{13}\selectfont\textsc{MapopSimpleToPartial}} \\
    {\begin{prooftree}

        % Evaluates id_f and e_i.
        \hypo{&\vdash\mathrm{e}_i\xrightarrow{e}v_i}
        \infer[no rule]1{&\Delta_c,\sigma_c\vdash\mathrm{id}_f\xrightarrow{e_o}v}

        % Checks that v complies with T.
        \hypo{v&\in_c{}T}
        \infer[no rule]1{v_i&\in_c{}\mathtt{nat}(n,m)}

        % Conclusion.
        \infer2
        [{\renewcommand{\arraystretch}{1.5}
          \begin{tabular}{l}
            $\mathrm{id}_a\in{}Sigs(\Delta)\cup{}Outs(\Delta)$ \\
            $\Delta(\mathrm{id}_a)=\mathtt{array(}T,n,m\mathtt{)}$ \\
            $\sigma={<}\mathcal{S},\mathcal{C},\mathcal{E}{>}$ \\
            $\mathcal{S}'=\mathcal{S}(id_a)\leftarrow\mathtt{set\_at}(v,v_i,\mathcal{S}(id_a))$ \\
            $\mathcal{E}'=\mathcal{E}\cup(\mathcal{S}\dcap\mathcal{S}')$ \\
          \end{tabular}
        }]
        {
          \Delta,\Delta_c,\sigma,\sigma_c\vdash
          \mathrm{id}_f\Rightarrow\mathrm{id}_a(\mathrm{e}_i)
          \xrightarrow{mapop}
          {<}\mathcal{S}',\mathcal{C},\mathcal{E}'{>}
        }
      \end{prooftree}} \\
  \end{tabular}
\end{table}

\begin{table}[H]
  \begin{tabular}{@{}l}
    {\fontsize{10}{13}\selectfont\textsc{MapopPartialToSimple}} \\
    {\begin{prooftree}

        % Evaluates name.
        \hypo{\Delta_c,\sigma_c\vdash\mathrm{id}_f(\mathrm{e}'_i)\xrightarrow{e_o}v}

        % Checks that v complies with T.
        \hypo{v\in_c{}T}
        
        % Conclusion.
        \infer2
        [{\renewcommand{\arraystretch}{1.5}
          \begin{tabular}{@{}l}
            $\mathrm{id}_a\in{}Sigs(\Delta)\cup{}Outs(\Delta)$ \\
            $\Delta(\mathrm{id}_a)=T$ \\
            $\sigma={<}\mathcal{S},\mathcal{C},\mathcal{E}{>}$ \\
            $\mathcal{S}'=\mathcal{S}(id_a)\leftarrow{}v$, $\mathcal{E}'=\mathcal{E}\cup(\mathcal{S}\dcap\mathcal{S}')$ \\
          \end{tabular}
        }]
        {
          \Delta,\Delta_c,\sigma,\sigma_c\vdash
          \mathrm{id}_f(\mathrm{e}'_i)\Rightarrow\mathrm{id}_a
          \xrightarrow{mapop}
          {<}\mathcal{S}',\mathcal{C},\mathcal{E}'{>}
        }
      \end{prooftree}} \\
  \end{tabular}
\end{table}

\begin{table}[H]
  \begin{tabular}{@{}l}
    {\fontsize{10}{13}\selectfont\textsc{MapopPartialToPartial}} \\
    {\begin{prooftree}

        % Evaluates id_f and e_i.
        \hypo{&\vdash\mathrm{e}_i\xrightarrow{e}v_i}
        \infer[no rule]1{&\Delta_c,\sigma_c\vdash\mathrm{id}_f(\mathrm{e}'_i)\xrightarrow{e_o}v}

        % Checks that v complies with T.
        \hypo{v&\in_c{}T}
        \infer[no rule]1{v_i&\in_c{}\mathtt{nat}(n,m)}

        % Conclusion.
        \infer2
        [{\renewcommand{\arraystretch}{1.5}
          \begin{tabular}{l}
            $\mathrm{id}_a\in{}Sigs(\Delta)\cup{}Outs(\Delta)$ \\
            $\Delta(\mathrm{id}_a)=\mathtt{array(}T,n,m\mathtt{)}$ \\
            $\sigma={<}\mathcal{S},\mathcal{C},\mathcal{E}{>}$ \\
            $\mathcal{S}'=\mathcal{S}(id_a)\leftarrow\mathtt{set\_at}(v,v_i,\mathcal{S}(id_a))$ \\
            $\mathcal{E}'=\mathcal{E}\cup(\mathcal{S}\dcap\mathcal{S}')$ \\
          \end{tabular}
        }]
        {
          \Delta,\Delta_c,\sigma,\sigma_c\vdash
          \mathrm{id}_f(\mathrm{e}'_i)\Rightarrow\mathrm{id}_a(\mathrm{e}_i)
          \xrightarrow{mapop}
          {<}\mathcal{S}',\mathcal{C},\mathcal{E}'{>}
        }
      \end{prooftree}} \\
  \end{tabular}
\end{table}


\begin{table}[H]
  {\fontsize{10}{13}\selectfont\textsc{MapopComp}}
  
  \begin{prooftree}    
    % Evaluates mapip on assoc.
    \hypo{\Delta,\Delta_c,\sigma,\sigma_c\vdash\mathrm{assoc}_{po}\xrightarrow{mapop}\sigma'}

    % Evaluates mapip on portmap.
    \hypo{\Delta,\Delta_c,\sigma',\sigma_c\vdash\mathrm{opmap}\xrightarrow{mapop}\sigma''}

    % Conclusion.
    \infer2
    {
      \Delta,\Delta_c,\sigma,\sigma_c\vdash
      \langle{}\mathrm{assoc}_{po},~\mathrm{opmap}\rangle
      \xrightarrow{mapop}\sigma''
    }
  \end{prooftree}
\end{table}

The $e_o$ relation is only defined to retrieve the value of out
ports from a store signal $\mathcal{S}$ under a design state
$\sigma={<}\mathcal{S},\mathcal{C},\mathcal{E}{>}$.

\begin{table}[H]
  \centering
  \begin{tabular}{@{}l}
    {\fontsize{10}{13}\selectfont\textsc{OutO}} \\
    {\begin{prooftree}
        \infer0
        [{
          \begin{tabular}{@{}l}
            $\mathrm{id}_s\in{}Outs(\Delta)$ \\
            $\mathrm{id}_s\in\sigma$ \\
          \end{tabular}
        }] {
          \Delta,\sigma\vdash
          \mathrm{id}_s
          \xrightarrow{e_o}
          \sigma(\mathrm{id}_s)
        }
      \end{prooftree}} \\
  \end{tabular}
\end{table}

\begin{table}[H]
  \centering
  \begin{tabular}{@{}l}
    {\fontsize{10}{13}\selectfont\textsc{IdxOutO}} \\
    {\begin{prooftree}

        % Evaluates e_i.
        \hypo{
          \vdash
          \mathrm{e}_i
          \xrightarrow{e}
          v_i
        }

        % Checks well-typed v_i.
        \hypo{v_i\in_c\mathtt{nat}(n,m)}
        
        % Conclusion.
        \infer2
        [{
          \begin{tabular}{@{}l}
            $\mathrm{id}_s\in{}Outs(\Delta)$ \\
            $\mathrm{id}_s\in\sigma$ \\
            $\Delta(\mathrm{id}_s)=\mathtt{array(}T,n,m\mathtt{)}$ \\
            $i=v_i~\mathtt{mod}~n$ \\
          \end{tabular}
        }]
        {
          \Delta,\sigma\vdash
          \mathrm{id}_s(\mathrm{e}_i)
          \xrightarrow{e_o}
          \mathtt{get\_at}(i,\sigma(\mathrm{id}_s))
        }
      \end{prooftree}} \\
  \end{tabular}
\end{table}

\subsection{Evaluation of sequential statements}
\label{subsec:seq-stmts}

The $ss$ symbol indicates that the evaluation of the considered
sequential statement does not depend on a specific flag (i.e, the $c$,
$i$, $\uparrow$ or $\downarrow$ flag). In the rules of the $ss$
relation, a $ss$ flag is tranferred from the conclusion to the
premises when an sequential statement is composed of inner sequential
blocks.

\subsubsection{Signal assigment statement} A signal assignment
generates a new design state with a modified signal store and a new
set of events. Note that there are two states on the left side of the
thesis symbol. $\sigma$ represents the state holding the current
values of signals, and $\sigma_w$ holds the new values of signals
(i.e. the \emph{written} state).

\begin{premises}
  The premise $\mathcal{S}(id_s)\in_c{}T$ checks that the value
  associated to signal $id_s$ in the signal store of $\sigma$ complies
  with type $T$, where $T$ is the type associated with signal $id_s$
  in $\Delta$.
\end{premises}

\begin{table}[H]
  \begin{tabular}{@{}l}
    {\fontsize{10}{13}\selectfont\textsc{SigAssign}} \\
    {\begin{prooftree}

        % Evaluates e.
        \hypo{\Delta,\sigma,\Lambda\vdash\mathrm{e}\xrightarrow{e}v}
        
        % Checks that v complies with T.
        \hypo{v\in_c{}T}
        
        % Conclusion.
        \infer2
        [{\renewcommand{\arraystretch}{1.5}
          \begin{tabular}{@{}l}
            $\mathrm{id}_s\in{}Sigs(\Delta)\cup{}Outs(\Delta)$ \\
            $\Delta(\mathrm{id}_s)=T$ \\
            % $\sigma={<}\mathcal{S},\mathcal{C},\mathcal{E}{>}$ \\
            $\mathcal{S}'_w=\mathcal{S}_w(id_s)\leftarrow{}v$ \\
            $\mathcal{E}'_w=\mathcal{E}_w\cup(\mathcal{S}_w\dcap\mathcal{S}'_w)$ \\
          \end{tabular}
        }]
        {
          \Delta,\sigma,\sigma_w,\Lambda\vdash
          \mathrm{id}_s\Leftarrow\mathrm{e}
          \xrightarrow{ss}
          {<}\mathcal{S}'_w,\mathcal{C}_w,\mathcal{E}'_w{>},\Lambda
        }
      \end{prooftree}} \\
  \end{tabular}
\end{table}

\begin{table}[H]
  \begin{tabular}{@{}l}
    {\fontsize{10}{13}\selectfont\textsc{IdxSigAssign}} \\
    {\begin{prooftree}

        % Evaluates e_i and e.
        \hypo{\Delta,\sigma,\Lambda&\vdash\mathrm{e}_i\xrightarrow{e}v_i}
        \infer[no rule]1{\Delta,\sigma,\Lambda&\vdash\mathrm{e}\xrightarrow{e}v}

        % Checks that v complies with T.

        \hypo{v&\in_c{}T}
        \infer[no rule]1{v_i&\in_c{}\mathtt{nat}(n,m)}
        
        % Conclusion.
        \infer2
        [{\renewcommand{\arraystretch}{1.5}
          \begin{tabular}{@{}l}
            $\mathrm{id}_s\in{}Sigs(\Delta)\cup{}Outs(\Delta)$ \\
            $\Delta(\mathrm{id}_s)=\mathtt{array(}T,n,m\mathtt{)}$ \\
            $\mathcal{S}'_w=\mathcal{S}_w(id_s)\leftarrow\mathtt{set\_at}(v,v_i,\mathcal{S}_w(id_s))$ \\
            $\mathcal{E}'_w=\mathcal{E}_w\cup(\mathcal{S}_w\dcap\mathcal{S}'_w)$ \\
          \end{tabular}
        }]
        {
          \Delta,\sigma,\sigma_w,\Lambda\vdash
          \mathrm{id}_s(\mathrm{e}_i)\Leftarrow\mathrm{e}
          \xrightarrow{ss}
          {<}\mathcal{S}'_w,\mathcal{C}_w,\mathcal{E}'_w{>},\Lambda
        }
      \end{prooftree}} \\
  \end{tabular}
\end{table}

\subsubsection{Variable assignment statement}

A variable assignment statement modifies the variable values in the
local environment.

\begin{table}[H]
  \begin{tabular}{@{}l}
    {\fontsize{10}{13}\selectfont\textsc{VarAssign}} \\    
    {\begin{prooftree}

        % Evaluates e.
        \hypo{\Delta,\sigma,\Lambda\vdash\mathrm{e}\xrightarrow{e}v}

        % Checks that v complies with T.
        \hypo{v\in_c{}T}

        % Conclusion.
        \infer2
        [{
          \begin{tabular}{@{}l}
            $\mathrm{id}_v\in\Lambda$ \\
            $\Lambda(\mathrm{id}_v)=(T,val)$ \\
          \end{tabular}
        }]
        {
          \Delta,\sigma,\sigma_w,\Lambda\vdash
          \mathrm{id}_v:=\mathrm{e}
          \xrightarrow{ss}
          \sigma_w,\Lambda(id_v)\leftarrow{}(T,v)
        }
      \end{prooftree}} \\
  \end{tabular}
\end{table}

\begin{table}[H]
  \begin{tabular}{@{}l}
    {\fontsize{10}{13}\selectfont\textsc{IdxVarAssign}} \\    
    {\begin{prooftree}

        % Evaluates e_i.
        \hypo{\Delta,\sigma,\Lambda&\vdash\mathrm{e}_i\xrightarrow{e}v_i}
        \infer[no rule]1{\Delta,\sigma,\Lambda&\vdash\mathrm{e}\xrightarrow{e}v}

        % Checks that v complies with T.
        \hypo{v_i&\in_c\mathtt{nat}(n,m)}
        \infer[no rule]1{v&\in_c{}T}

        % Conclusion.
        \infer2
        [{
          \begin{tabular}{@{}l}
            $\mathrm{id}_v\in\Lambda$ \\
            $\Lambda(\mathrm{id}_v)=(\mathtt{array(}T,n,m\mathtt{)},val)$ \\
          \end{tabular}
        }]
        {
          \Delta,\sigma,\sigma_w,\Lambda\vdash
          \mathrm{id}_v(\mathrm{e}_i):=\mathrm{e}
          \xrightarrow{ss}
          \sigma_w,\Lambda(id_v)\leftarrow{}(T,\mathtt{set\_at}(v,v_i,val))
        }
      \end{prooftree}} \\
  \end{tabular}
\end{table}

\subsubsection{If statement}
\label{subsubsec:if-stmt}

\begin{table}[H]
  \begin{tabular}{@{}l}
    {\fontsize{10}{13}\selectfont\textsc{If}$\top$} \\    
    {\begin{prooftree}

        % Evaluates condition.
        \hypo{\Delta,\sigma,\Lambda\vdash\mathrm{e}\xrightarrow{e}\top}

        % Evaluates ss.
        \hypo{\Delta,\sigma,\sigma_w,\Lambda\vdash\mathrm{ss}\xrightarrow{ss}\sigma'_w,\Lambda'}

        % Conclusion.
        \infer2
        {
          \Delta,\sigma,\sigma_w,\Lambda\vdash~
          $\vhdle|if|$~(\mathrm{e})~\mathrm{ss}
          \xrightarrow{ss}
          \sigma'_w,\Lambda'
        }
      \end{prooftree}} \\
  \end{tabular}
  \begin{tabular}{@{}l}
    {\fontsize{10}{13}\selectfont\textsc{If}$\bot$} \\
    {\begin{prooftree}

        % Evaluates condition.
        \hypo{\Delta,\sigma,\Lambda\vdash\mathrm{e}\xrightarrow{e}\bot}

        % Conclusion.
        \infer1
        {
          \Delta,\sigma,\sigma_w,\Lambda\vdash~
          $\vhdle|if|$~(\mathrm{e})~\mathrm{ss}
          \xrightarrow{ss}
          \sigma_w,\Lambda
        }
      \end{prooftree}} \\
  \end{tabular}
\end{table}

\begin{table}[H]
  \begin{tabular}{@{}l}
    {\fontsize{10}{13}\selectfont\textsc{IfElse}$\top$} \\    
    {\begin{prooftree}

        % Evaluates condition.
        \hypo{\Delta,\sigma,\Lambda\vdash\mathrm{e}\xrightarrow{expr}\top}

        % Evaluates ss.
        \hypo{\Delta,\sigma,\sigma_w,\Lambda\vdash\mathrm{ss}\xrightarrow{ss}\sigma'_w,\Lambda'}

        % Conclusion.
        \infer2
        {
          \Delta,\sigma,\sigma_w,\Lambda\vdash~
          $\vhdle|if|$~(\mathrm{e})~\mathrm{ss}~\mathrm{ss}'
          \xrightarrow{ss}
          \sigma'_w,\Lambda'
        }
      \end{prooftree}} \\
  \end{tabular}
  \begin{tabular}{@{}l}
    {\fontsize{10}{13}\selectfont\textsc{IfElse}$\bot$} \\  
    {\begin{prooftree}

        % Evaluates condition.
        \hypo{\Delta,\sigma,\Lambda\vdash\mathrm{e}\xrightarrow{e}\bot}

        % Evaluates ss.
        \hypo{\Delta,\sigma,\sigma_w,\Lambda\vdash\mathrm{ss}'\xrightarrow{ss}\sigma'_w,\Lambda'}

        % Conclusion.
        \infer2
        {
          \Delta,\sigma,\sigma_w,\Lambda\vdash~
          $\vhdle|if|$~(\mathrm{e})~\mathrm{ss}~\mathrm{ss}'
          \xrightarrow{ss}
          \sigma'_w,\Lambda'
        }
      \end{prooftree}} \\
  \end{tabular}
\end{table}

\subsubsection{Loop statement}

\begin{table}[H]
  \begin{tabular}{@{}l}
    {\fontsize{10}{13}\selectfont\textsc{Loop}$\bot$} \\
    {\begin{prooftree}

        % Evaluates upper bound check.
        \hypo{\Delta,\sigma,\Lambda_i\vdash\mathrm{id}_v=\mathrm{e}'\xrightarrow{e}\bot}

        % Evaluates ss and recursive call.
        \hypo{\Delta,\sigma,\sigma_w,\Lambda_i&\vdash\mathrm{ss}\xrightarrow{ss}\sigma'_w,\Lambda'}
        \infer[no rule]1{
          \Delta,\sigma,\sigma'_w,\Lambda'&\vdash~
          $\vhdle|for|$~(\mathrm{id}_v,\mathrm{e},\mathrm{e}')~\mathrm{ss}
          \xrightarrow{ss}\sigma''_w,\Lambda''
        }
        
        % Conclusion. 
        \infer2
        [{\renewcommand{\arraystretch}{1.5}
          \begin{tabular}{@{}l}
            $\mathrm{id}_v\in\Lambda$ \\
            $\Lambda(\mathrm{id}_v)=(T,val)$ \\
            $\Lambda_i=\Lambda(id_v)\leftarrow{}(T,val+1)$ \\
          \end{tabular}
        }]
        {
          \Delta,\sigma,\sigma_w,\Lambda\vdash~
          $\vhdle|for|$~(\mathrm{id}_v,\mathrm{e},\mathrm{e}')~\mathrm{ss}
          \xrightarrow{ss}
          \sigma''_w,\Lambda''
        }
      \end{prooftree}} \\
  \end{tabular}
\end{table}

\begin{table}[H]
  \begin{tabular}{@{}l}
    {\fontsize{10}{13}\selectfont\textsc{Loop}$\top$} \\    
    {\begin{prooftree}
                
        % Upper bound check true.
        \hypo{\Delta,\sigma,\Lambda_i\vdash\mathrm{id}_v=\mathrm{e}'\xrightarrow{e}\top}

        % Conclusion.
        \infer1
        [{\renewcommand{\arraystretch}{1.5}
          \begin{tabular}{@{}l}
            $\mathrm{id}_v\in\Lambda$ \\
            $\Lambda(\mathrm{id}_v)=(T,val)$ \\
            $\Lambda_i=\Lambda(id_v)\leftarrow(T,val+1)$ \\
          \end{tabular}
        }]
        {
          \Delta,\sigma,\sigma_w,\Lambda\vdash~
          $\vhdle|for|$~(\mathrm{id}_v,\mathrm{e},\mathrm{e}')~\mathrm{ss}
          \xrightarrow{ss}
          \sigma_w,\Lambda\setminus(id_v,\Lambda(id_v))
        }
      \end{prooftree}} \\
  \end{tabular}
\end{table}

\begin{table}[H]
  \begin{tabular}{@{}l}
    {\fontsize{10}{13}\selectfont\textsc{LoopInit}} \\
    {\begin{prooftree}
        
        % Evaluates e and e'.
        \hypo{\Delta,\sigma,\Lambda&\vdash\mathrm{e}\xrightarrow{e}v}
        \infer[no rule]1{\Delta,\sigma,\Lambda&\vdash\mathrm{e'}\xrightarrow{e}v'}
        
        % Upper bound check true.
        \hypo{\Delta,\sigma,\sigma_w,\Lambda_i\vdash~
          $\vhdle|for|$~(\mathrm{id}_v,\mathrm{e},\mathrm{e}')~\mathrm{ss}
          \xrightarrow{ss}
          \sigma'_w,\Lambda'}

        % Conclusion.
        \infer2
        [{\renewcommand{\arraystretch}{1.5}
          \begin{tabular}{@{}l}
            $\mathrm{id}_v\notin\Lambda$ \\
            $\Lambda_i=\Lambda\cup(id_v,(\mathtt{nat}(v,v'),v))$ \\
          \end{tabular}
        }]
        {
          \Delta,\sigma,\sigma_w,\Lambda\vdash~
          $\vhdle|for|$~(\mathrm{id}_v,\mathrm{e},\mathrm{e}')~\mathrm{ss}
          \xrightarrow{ss}
          \sigma'_w,\Lambda'
        }
      \end{prooftree}} \\
  \end{tabular}
\end{table}

\subsubsection{Rising and falling edge block statements}
\label{subsubsec:rise-and-fall-stmts}

\begin{table}[H]
  \begin{tabular}{l}
    {\fontsize{10}{13}\selectfont\textsc{RisingEdgeDefault}} \\    
    {\begin{prooftree}
        
        % Conclusion.
        \infer0
        [$f\neq\uparrow$]
        {
          \Delta,\sigma,\sigma_w,\Lambda\vdash~
          $\vhdle|rising|$~\mathrm{ss}
          \xrightarrow{ss_f}
          \sigma_w,\Lambda
        }
      \end{prooftree}} \\
  \end{tabular}
  \begin{tabular}{l}
    {\fontsize{10}{13}\selectfont\textsc{FallingEdgeDefault}} \\    
    {\begin{prooftree}
        
        % Conclusion.
        \infer0
        [$f\neq\downarrow$]
        {
          \Delta,\sigma,\sigma_w,\Lambda\vdash~
          $\vhdle|falling|$~\mathrm{ss}
          \xrightarrow{ss_f}
          \sigma_w,\Lambda
        }
      \end{prooftree}} \\
  \end{tabular}
\end{table}

\begin{table}[H]
  \begin{tabular}{l}
    {\fontsize{10}{13}\selectfont\textsc{RisingEdgeExec}} \\    
    {\begin{prooftree}

        % Evaluates ss.
        \hypo{\Delta,\sigma,\sigma_w,\Lambda\vdash\mathrm{ss}\xrightarrow{ss_\uparrow}\sigma'_w,\Lambda'}
        
        % Conclusion.
        \infer1
        {
          \Delta,\sigma,\sigma_w,\Lambda\vdash~
          $\vhdle|rising|$~\mathrm{ss}
          \xrightarrow{ss_\uparrow}
          \sigma'_w,\Lambda'
        }
      \end{prooftree}} \\
  \end{tabular}
  \begin{tabular}{l}
    {\fontsize{10}{13}\selectfont\textsc{FallingEdgeExec}} \\    
    {\begin{prooftree}

        % Evaluates ss.
        \hypo{\Delta,\sigma,\sigma_w,\Lambda\vdash\mathrm{ss}\xrightarrow{ss_\downarrow}\sigma'_w,\Lambda'}
        
        % Conclusion.
        \infer1
        {
          \Delta,\sigma,\sigma_w,\Lambda\vdash~
          $\vhdle|falling|$~\mathrm{ss}
          \xrightarrow{ss_\downarrow}
          \sigma'_w,\Lambda'
        }
      \end{prooftree}} \\
  \end{tabular}
\end{table}

\subsubsection{Rst block statement}

\begin{table}[H]
  \begin{tabular}{l}
    {\fontsize{10}{13}\selectfont\textsc{RstDefault}} \\    
    {\begin{prooftree}

        % Evaluates ss.
        \hypo{\Delta,\sigma,\sigma_w,\Lambda\vdash\mathrm{ss'}\xrightarrow{ss_f}\sigma'_w,\Lambda'}
        
        % Conclusion.
        \infer1
        [$f\neq{}i$]
        {
          \Delta,\sigma,\sigma_w,\Lambda\vdash~
          $\vhdle|rst|$~\mathrm{ss}~\mathrm{ss'}
          \xrightarrow{ss_f}
          \sigma'_w,\Lambda'
        }
      \end{prooftree}} \\
  \end{tabular}
  \begin{tabular}{l}
    {\fontsize{10}{13}\selectfont\textsc{RstExec}} \\    
    {\begin{prooftree}

        % Evaluates ss.
        \hypo{\Delta,\sigma,\sigma_w,\Lambda\vdash\mathrm{ss}\xrightarrow{ss_i}\sigma'_w,\Lambda'}
        
        % Conclusion.
        \infer1
        {
          \Delta,\sigma,\sigma_w,\Lambda\vdash~
          $\vhdle|rst|$~\mathrm{ss}~\mathrm{ss'}
          \xrightarrow{ss_i}
          \sigma'_w,\Lambda'
        }
      \end{prooftree}} \\
  \end{tabular}
\end{table}

\subsubsection{Composition of sequential statements and null statement}

\begin{table}[H]
  \begin{tabular}{l}
    {\fontsize{10}{13}\selectfont\textsc{SeqStmt}} \\
    {\begin{prooftree}
        
        % Evaluates ss.
        \hypo{\Delta,\sigma,\sigma_w,\Lambda\vdash\mathrm{ss}\xrightarrow{ss}\sigma'_w,\Lambda'}

        % Evaluates ss'.
        \hypo{\Delta,\sigma,\sigma'_w,\Lambda'\vdash\mathrm{ss'}\xrightarrow{ss}\sigma''_w,\Lambda''}

        % Conclusion.
        \infer2{
          \Delta,\sigma,\sigma_w,\Lambda\vdash\mathrm{ss}\mathtt{;}~\mathrm{ss}'\xrightarrow{ss}\sigma''_w,\Lambda''
        }
      \end{prooftree}} \\
  \end{tabular}
  \begin{tabular}{l}
    {\fontsize{10}{13}\selectfont\textsc{NullStmt}} \\
    {\begin{prooftree}        

        % Conclusion.
        \infer0{
          \Delta,\sigma,\sigma_w,\Lambda\vdash\mathtt{null}\xrightarrow{ss}\sigma_w,\Lambda
        }
      \end{prooftree}} \\
  \end{tabular}
\end{table}

\subsection{Evaluation of expressions}
\label{subsec:expr-rules}

% Nat and Bool literals.

\begin{table}[H]
  \begin{tabular}{l}
    {\fontsize{10}{13}\selectfont\textsc{Nat}} \\
    {\begin{prooftree}
        \infer0
        [{
          \begin{tabular}{@{}l}
            $\mathrm{n}\in\mathbb{N}$ \\
            $\mathrm{n}\le\mathtt{NATMAX}$ \\
          \end{tabular}
        }]
        {
          \Delta,\sigma,\Lambda\vdash\mathrm{n}\xrightarrow{e}n
        }
      \end{prooftree}} \\
  \end{tabular}
  \begin{tabular}{l}
    {\fontsize{10}{13}\selectfont\textsc{False}} \\
    {\begin{prooftree}
        \infer0
        {
          \Delta,\sigma,\Lambda\vdash\mathtt{false}\xrightarrow{e}\bot
        }
      \end{prooftree}} \\
  \end{tabular}
  \begin{tabular}{l}
    {\fontsize{10}{13}\selectfont\textsc{True}} \\
    {\begin{prooftree}
        \infer0
        {
          \Delta,\sigma,\Lambda\vdash\mathtt{true}\xrightarrow{e}\top
        }
      \end{prooftree}} \\
  \end{tabular}

\end{table}

% Aggregate.

\begin{table}[H]
  \centering
  \begin{tabular}{l}
    {\fontsize{10}{13}\selectfont\textsc{Aggreg}} \\
    {\begin{prooftree}

        % Evaluates e.
        \hypo{\Delta,\sigma,\Lambda\vdash\mathrm{e}_i\xrightarrow{e}v_i}

        % Conclusion.
        \infer1
        [$i=1,\dots,n$]
        {
          \Delta,\sigma,\Lambda\vdash\mathtt{(}\mathrm{e}_1,\dots,\mathrm{e}_n\mathtt{)}\xrightarrow{e}(v_1,\dots,v_n)
        }
      \end{prooftree}} \\
  \end{tabular}
\end{table}

% Name expressions.

\begin{table}[H]
  \centering
  \begin{tabular}{l}
    {\fontsize{10}{13}\selectfont\textsc{Sig}} \\
    {\begin{prooftree}
        \infer0
        [{
          \begin{tabular}{@{}l}
            $\mathrm{id}_s\in{}Sigs(\Delta)\cup{}Ins(\Delta)$ \\
          \end{tabular}
        }]
        {
          \Delta,\sigma,\Lambda\vdash\mathrm{id}_s\xrightarrow{e}\sigma(\mathrm{id}_s)
        }
      \end{prooftree}} \\
  \end{tabular}
  \begin{tabular}{@{}l}
    {\fontsize{10}{13}\selectfont\textsc{Var}} \\
    {\begin{prooftree}
        \infer0
        [{
          \begin{tabular}{@{}l}
            $\mathrm{id}_v\in\Lambda$ \\
            $\Lambda(\mathrm{id}_v)=(T,v)$ \\
          \end{tabular}
        }]
        {
          \Delta,\sigma,\Lambda\vdash\mathrm{id}_v\xrightarrow{e}v
        }
      \end{prooftree}} \\
    \end{tabular}
\end{table}

\begin{table}[H]
  \centering
  \begin{tabular}{@{}l}
    {\fontsize{10}{13}\selectfont\textsc{Gen}} \\
    {\begin{prooftree}
      \infer0
      [{
        \begin{tabular}{@{}l}
          $\mathrm{id}_g\in{}Gens(\Delta)$ \\
          $\Delta(\mathrm{id}_g)=(T,v)$ \\
        \end{tabular}
      }]
      {
        \Delta,\sigma,\Lambda\vdash\mathrm{id}_g\xrightarrow{e}v
      }
    \end{prooftree}} \\
  \end{tabular}
  % Idx name expressions.
  \begin{tabular}{@{}l}
    {\fontsize{10}{13}\selectfont\textsc{IdxSig}} \\
    {\begin{prooftree}

        % Evaluates e_i.
        \hypo{\Delta,\sigma,\Lambda\vdash\mathrm{e}_i\xrightarrow{e}v_i}

        % Well-typed v_i.
        \hypo{v_i\in_c{}\mathtt{nat}(n,m)}
        
        % Conclusion.
        \infer2
        [{
          \begin{tabular}{@{}l}
            $\mathrm{id}_s\in{}Sigs(\Delta)\cup{}Ins(\Delta)$ \\
            $\Delta(\mathrm{id}_s)=\mathtt{array(}T,n,m\mathtt{)}$ \\
            $i=v_i~\mathtt{mod}~n$ \\
          \end{tabular}
        }]
        {
          \Delta,\sigma,\Lambda\vdash
          \mathrm{id}_s(\mathrm{e}_i)
          \xrightarrow{e}
          \mathtt{get\_at}(i,\sigma(\mathrm{id}_s))
        }
      \end{prooftree}} \\
  \end{tabular}
\end{table}

\begin{table}[H]
  \centering
  \begin{tabular}{@{}l}
    {\fontsize{10}{13}\selectfont\textsc{IdxVar}} \\
    {\begin{prooftree}
        
        % Evaluates e_i.
        \hypo{\Delta,\sigma,\Lambda\vdash\mathrm{e}_i\xrightarrow{e}v_i}

        % Well-typed v_i.
        \hypo{v_i\in_c{}\mathtt{nat}(n,m)}
        
        % Conclusion.
        \infer2
        [{
          \begin{tabular}{@{}l}
            $\mathrm{id}_v\in\Lambda$ \\
            $\Lambda(\mathrm{id}_v)=(\mathtt{array(}T,n,m\mathtt{)},v)$ \\
            $i=v_i~\mathtt{mod}~n$ \\
          \end{tabular}
        }]
        {
          \Delta,\sigma,\Lambda\vdash
          \mathrm{id}_v(\mathrm{e}_i)
          \xrightarrow{e}
          \mathtt{get\_at}(i,v)
        }
      \end{prooftree}} \\
  \end{tabular}
\end{table}

where \texttt{get\_at}$(i,a)$ is a function returning the $i$-th
element of array $a$.

% Natural operators.

\begin{table}[H]
  \begin{tabular}{@{}l}
    {\fontsize{10}{13}\selectfont\textsc{NatAdd}} \\
    {\begin{prooftree}

        % Evaluates e and e'.
        \hypo{\Delta,\sigma,\Lambda&\vdash\mathrm{e}\xrightarrow{e}v}

        \hypo{\Delta,\sigma,\Lambda&\vdash\mathrm{e'}\xrightarrow{e}v'}

        % Conclusion.
        \infer2
        [$v+_{\mathbb{N}}v'\le\mathtt{NATMAX}$]
        {
          \Delta,\sigma,\Lambda\vdash\mathrm{e}~+~\mathrm{e'}\xrightarrow{e}v~+_{\mathbb{N}}~v'
        }
      \end{prooftree}} \\
  \end{tabular}
\end{table}

\begin{table}[H]
  \begin{tabular}{@{}l}
    {\fontsize{10}{13}\selectfont\textsc{NatSub}} \\
    {\begin{prooftree}

        % Evaluates e and e'.
        \hypo{\Delta,\sigma,\Lambda&\vdash\mathrm{e}\xrightarrow{e}v}
        \hypo{\Delta,\sigma,\Lambda&\vdash\mathrm{e'}\xrightarrow{e}v'}

        % Conclusion.
        \infer2
        [$v\ge{}v'$]
        {
          \Delta,\sigma,\Lambda\vdash\mathrm{e}~-~\mathrm{e'}\xrightarrow{e}v~-_{\mathbb{N}}~v'
        }
      \end{prooftree}} \\
  \end{tabular}
\end{table}

\begin{table}[H]
  \begin{tabular}{@{}l}
    {\fontsize{10}{13}\selectfont\textsc{OrdOp}} \\
    
    {\begin{prooftree}

        % Evaluates e and e'.
        \hypo{\Delta,\sigma,\Lambda&\vdash\mathrm{e}\xrightarrow{e}v}
        \hypo{\Delta,\sigma,\Lambda&\vdash\mathrm{e'}\xrightarrow{e}v'}

        % Conclusion.
        \infer2
        [$\mathrm{op}_{ordn}\in{}\{<,\le,>,\ge\}$]
        {
          \Delta,\sigma,\Lambda\vdash\mathrm{e}~\mathrm{op}_{ordn}~\mathrm{e'}\xrightarrow{e}v~op_{ord\mathbb{N}}~v'
        }
      \end{prooftree}} \\
  \end{tabular}
\end{table}

% Boolean operators.

\begin{table}[H]
  \begin{tabular}{@{}l}
    {\fontsize{10}{13}\selectfont\textsc{BoolBinOp}} \\
    {\begin{prooftree}

        % Evaluates e adn e'.
        \hypo{\Delta,\sigma,\Lambda&\vdash\mathrm{e}\xrightarrow{e}v}
        \hypo{\Delta,\sigma,\Lambda&\vdash\mathrm{e'}\xrightarrow{e}v'}

        % Conclusion.
        \infer2
        [$\mathrm{op}_{bool}\in{}\{\mathtt{and},\mathtt{or}\}$]
        {
          \Delta,\sigma,\Lambda\vdash
          \mathrm{e}~\mathrm{op}_{bool}~\mathrm{e'}
          \xrightarrow{e}
          v~op_{\mathbb{B}}~v'
        }
      \end{prooftree}} \\
  \end{tabular}
  \begin{tabular}{@{}l}
    {\fontsize{10}{13}\selectfont\textsc{NotOp}} \\    
    {\begin{prooftree}

        % Evaluates e.
        \hypo{\Delta,\sigma,\Lambda\vdash\mathrm{e}\xrightarrow{e}v}
        
        % Conclusion.
        \infer1
        {
          \Delta,\sigma,\Lambda\vdash~$\vhdle|not|$~\mathrm{e}\xrightarrow{e}\lnot{}v
        }
      \end{prooftree}} \\
  \end{tabular}
\end{table}

% Equality and difference.

\begin{table}[H]
  {\fontsize{10}{13}\selectfont\textsc{EqOp}}
  
  \begin{prooftree}

    % Evaluates e.
    \hypo{\Delta,\sigma,\Lambda\vdash\mathrm{e}\xrightarrow{e}v}

    % Evaluates e'.
    \hypo{\Delta,\sigma,\Lambda\vdash\mathrm{e'}\xrightarrow{e}v'}

    % Conclusion.
    \infer2
    {
      \Delta,\sigma,\Lambda\vdash\mathrm{e}~=~\mathrm{e'}\xrightarrow{e}eq(v,v')
    }
  \end{prooftree}
\end{table}

\begin{table}[H]
  {\fontsize{10}{13}\selectfont\textsc{DiffOp}}
  
  \begin{prooftree}

    % Evaluates e.
    \hypo{\Delta,\sigma,\Lambda\vdash\mathrm{e}~=~\mathrm{e'}\xrightarrow{e}v}

    % Conclusion.
    \infer1
    {
      \Delta,\sigma,\Lambda\vdash
      \mathrm{e}~\neq~\mathrm{e'}
      \xrightarrow{e}\lnot{}v
    }
  \end{prooftree}
\end{table}

where $eq$ is the equality relation established for all types defined
in the semantics.

\begin{table}[H]
  {\fontsize{10}{13}\selectfont\textsc{Parenth}}
  
  \begin{prooftree}

    % Evaluates e.
    \hypo{\Delta,\sigma,\Lambda\vdash\mathrm{e}\xrightarrow{e}v}

    % Conclusion.
    \infer1
    {
      \Delta,\sigma,\Lambda\vdash\mathtt{(}\mathrm{e}\mathtt{)}\xrightarrow{e}v
    }
  \end{prooftree}
\end{table}

%%% Local Variables:
%%% mode: latex
%%% TeX-master: "../../main"
%%% End:


\section{An example of full simulation}
\label{sec:ex-full-sim}
In this section, we will demonstrate the full simulation of a \hvhdl{}
top-level design on the example of Listing~\ref{lst:tl-design-ex}. The
aim here is to illustrate the formal rules of the \hvhdl{}
semantics. Listing~\ref{lst:tl-design-ex} is the result of the
transformation of the SITPN model presented in
Figure~\ref{fig:sitpn-ex-full-sim} into a \hvhdl{} design. To keep the
examples within a reasonable size, Listing~\ref{lst:tl-design-ex}, and
the other listings and derivation trees used in this section, refer to
the generic constants and ports of the transition and place designs by
their short names. See Table~\ref{tab:consts-and-sigs-ref} for a
correspondence between the short names and the full names of constants
and signals of the place and transition designs.

\begin{lstlisting}[language=VHDL,label={lst:tl-design-ex},
caption={[An example of \hvhdl{} design generated by \hilecop{}.] An example of \hvhdl{} top-level design generated by the \hilecop{} transformation.},framexleftmargin=1.5em,xleftmargin=2em,numbers=left,
numberstyle=\tiny\ttfamily]
design tl tla

-- Generic constants
$\emptyset$

-- Ports ($ports_{tl}$)
((in, $id_{c0}$, boolean), (out, $id_{f0}$, boolean), (out, $id_{a0}$, boolean))

-- Declared (internal) signals ($sigs_{tl}$)
(($id_{ft}$, boolean), ($id_{av}$, boolean), ($id_{rt}$, boolean),($id_m$, boolean))

-- Behavior ($cs_{tl}$)

-- Place component instance $id_p$
comp ($id_p$, place,
-- Generic map
((ian, 1), (oan, 1), (mm, 1)),

-- Input port map
((im, 1), (iaw(0), 1), (oat(0), 0), (oaw(0), 1), (itf(0), $id_{ft}$), (otf(0), $id_{ft}$))

-- Output port map
((oav(0), $id_{av}$), (pauths, open), (rtt(0), $id_{rt}$), (marked, $id_m$)))

||

-- Transition component instance $id_t$
comp ($id_t$, transition,
-- Generic map
((tt, 0), (ian, 1), (cn, 1)),

-- Input port map
((ic(0), $id_{c0}$), (A, 0), (B, 0), (iav(0), $id_{av}$), (rt(0), $id_{rt}$), (pauths(0), true)),

-- Output port map
((fired, $id_{ft}$)))

||

-- The marked process
process (marked, {clk}, $\emptyset$,
(rst ($id_{a0}$ <= false)
     (falling ($id_{a0}$ <= $id_m$ or false))))

||

-- The fired process
process (fired, {clk}, $\emptyset$,
(rst ($id_{f0}$ <= false)
     (rising ($id_{f0}$ <= $id_{ft}$ or false))))
\end{lstlisting}

\begin{figure}[H]
  \centering
  \includegraphics[keepaspectratio,width=.15\textwidth]{Figures/H-VHDL/sitpn-ex-full-sim}
  \caption[An SITPN model transformed into a \hvhdl{} top-level design.]{The SITPN model at the base of the generation of the top-level design presented in Listing~\ref{lst:tl-design-ex}.}
  \label{fig:sitpn-ex-full-sim}
\end{figure}

The following rule states that the full simulation of the \texttt{tl}
design (presented in Listing~\ref{lst:tl-design-ex}) over 1 clock
cycle yields the simulation trace $(\sigma_0::\sigma_1::\sigma_2)$.
The simulation over one clock cycle (the rightmost premise) yields a
trace composed of the two states: the state $\sigma_1$ at half the
clock period, and the state $\sigma_2$ at the end of the first
cycle. The full simulation happens in the context of the \hilecop's
design store $\mathcal{D}_\mathcal{H}$, the elaborated design
$\Delta$, an empty dimensioning function and an simulation environment
$E_p$. Here, $ports_{tl}$ is an alias for the list of ports of
\texttt{tl}, $sigs_{tl}$ for the list of internal signals of
\texttt{tl}, and $cs_{tl}$ for the behavior of \texttt{tl}. In what
follows, we will detail the premises of the \textsc{FullSim} rule.

\begin{figure}[H]
  \begin{prooftree}[template={\fontsize{11}{12}\selectfont\inserttext}]

    % Design elab.
    \hypo{}
    \ellipsis{}{}
    \infer1{$\mathcal{D}_\mathcal{H},\emptyset\vdash$ \vhdle|design tl| $\dots\xrightarrow{elab}\Delta,\sigma_e$}

    % Initialization.
    \hypo{}
    \ellipsis{}{}
    \infer1{$\mathcal{D}_\mathcal{H},\Delta,\sigma_e\vdash{}cs_{tl}\xrightarrow{init}\sigma_0$}
       
    % Simulation loop.
    \hypo{}
    \ellipsis{}{}
    \infer1{$\mathcal{D}_\mathcal{H},E_p,\Delta,1,\sigma_0\vdash{}cs_{tl}\rightarrow(\sigma_1::\sigma_2)$}
    
    \infer3[\fontsize{7}{10}\selectfont\textsc{FullSim}]{
        $\mathcal{D}_\mathcal{H},\Delta,\emptyset,E_p,1\vdash$      
        \vhdle|design tl tla| $\empty$ $ports_{tl}$ $sigs_{tl}$ $cs_{tl}$ $\xrightarrow{full}(\sigma_0::\sigma_1::\sigma_2)$
    }
  \end{prooftree}
  \label{fig:ex-full-sim-tl}
  \caption{The \textsc{FullSim} rule applied to the \texttt{tl} design.}
\end{figure}

\subsection{Elaboration of the \texttt{tl} design}
\label{sec:ex-elab-of-tl}

The following rule state the elaborated design $\Delta$ and the
default design state $\sigma_e$ are the result of the elaboration of
the \texttt{tl} design. From left to right in the premises of the
rule, the three premises pertain to the elaboration of the declarative
parts of the \texttt{tl} design, i.e. the generic constant declaration
list, the port declaration list and the internal signal declaration
list. The leftmost premise pertains to the elaboration of the behavior
of the \texttt{tl} design.

\begin{figure}[H]
  \begin{prooftree}[template={\fontsize{9}{12}\selectfont\inserttext}]

    % Gens elab.
    \hypo{}
    \ellipsis{}{}
    \infer1{$\emptyset,\emptyset\vdash{}gens_{tl}\xrightarrow{egens}\Delta_0$}

    % Ports elab.
    \hypo{}
    \ellipsis{}{}
    \infer1{ $\Delta_0,\emptyset\vdash$
      $ports_{tl}\xrightarrow{eports}\Delta_1,\sigma_{e1}$}

    % Sigs elab.
    \hypo{}
    \ellipsis{}{}
    \infer1{
      $\Delta_1,\sigma_{e1}\vdash$
      $sigs_{tl}\xrightarrow{esigs}\Delta_2,\sigma_{e2}$}

    % Behavior elab.
    \hypo{}
    \ellipsis{}{}
    \infer1{$\mathcal{D}_\mathcal{H},\Delta_2,\sigma_{e2}\vdash$
      $cs_{tl}\xrightarrow{ebeh}\Delta,\sigma_e$}
        
    \infer4[\fontsize{7}{10}\selectfont\textsc{DesignElab}]{
      {\fontsize{12}{12}\selectfont
      $\mathcal{D}_\mathcal{H},\emptyset\vdash$
      \vhdle|design tl tla| $gens_{tl}$ $ports_{tl}$ $sigs_{tl}$ $cs_{tl}$
      $\xrightarrow{elab}\Delta,\sigma_e$}
    }
  \end{prooftree}
  \caption{The \textsc{DesignElab} rule applied to the \texttt{tl} design.}
  \label{fig:ex-elab-tl}
\end{figure}

\subsubsection{Elaboration of the declarative parts}
\label{sec:ex-elab-decl}

The elaboration of the declarative parts populates the $Gens$, $Ins$,
$Outs$ and $Sigs$ sub-environments of the elaborated design
$\Delta$. Here is the content of the $Gens$, $Ins$, $Outs$ and $Sigs$
sub-environments of $\Delta_2$, where $\Delta_2$ is the partial
elaborated design after the elaboration of the declarative parts of
the \texttt{tl} design (passed as a parameter of third and the fourth
premises of the rule described in Figure~\ref{fig:ex-elab-tl}).

\begin{itemize}
\item $Gens(\Delta_2):=\emptyset$
\item $Ins(\Delta_2):=\{(id_{c0}, bool)\}$
\item $Outs(\Delta_2):=\{(id_{f0}, bool), (id_{a0}, bool)\}$
\item $Sigs(\Delta_2):=\{(id_{ft}, bool), (id_{av}, bool), (id_{rt}, bool), (id_m, bool)\}$
\end{itemize}

The top-level design generated by the \hilecop{} transformation all
have an empty list of generic constants (see
Chapter~\ref{chap:transformation} for more details about the
transformation). Also, all ports and internal signals are of the
Boolean type. Thus, there are no range constraint or index constraint
to solve here. The \texttt{boolean} type indication is simply
transformed into the $bool$ semantic type.

The elaboration of the declarative parts also gives a default value to
the signals in the signal store of the default design state
$\sigma_{e2}$, where $\sigma_{e2}$ is the partial default design state
at the end of the elaboration of the declarative parts of the
\texttt{tl} design (passed as a parameter of the third and the fourth
premises of the rule described in Figure~\ref{fig:ex-elab-tl}).  Here
is the content of the signal store $\mathcal{S}$ of $\sigma_{e2}$.

\begin{itemize}
\item
  $\mathcal{S}(\sigma_{e2}):=\{(id_{c0},\bot),(id_{f0},\bot),(id_{a0},\bot),(id_{ft},
  \bot), (id_{av}, \bot), (id_{rt}, \bot), (id_m, \bot)\}$
\end{itemize}

The default value associated to the $bool$ type is $\bot$, thus, all
signals of the \texttt{tl design} are initialized to $\bot$ in the
signal store of $\sigma_{e2}$.

\subsubsection{Elaboration of the behavioral part}
\label{sec:ex-elab-beh}

The behavior of the \texttt{tl} design contains two component
instantiation statements and two process statements. Each one of these
statements will be elaborated in sequence. First, we present the
elaboration of the \texttt{marked} process to illustrate the
elaboration of a process statement; then, we present the elaboration
of place component instance $id_p$ to illustrate the elaboration of a
component instantiation statement.

\paragraph{Elaboration of a process statement}

The following rule presents the elaboration of the \texttt{marked}
process defined in the behavior of the \texttt{tl} design.

\begin{figure}[H]
  \begin{prooftree}[template={\fontsize{11}{12}\selectfont\inserttext}]

    % Sigs elab.
    \infer0{
      $\Delta_2,\emptyset\vdash$
      $\emptyset\xrightarrow{evars}\emptyset$}

    % Behavior elab.
    \hypo{}
    \ellipsis{}{}
    \infer1[\fontsize{7}{10}\selectfont\textsc{WTRst}]{$\Delta_2,\sigma_{e2},\emptyset\vdash$
      $\mathtt{valid}_{ss}\Big($
      \begin{tabular}{@{}l@{}}
        $\mathtt{rst}~(id_{a0}\Leftarrow{}\mathtt{false})$ \\
        $(\mathtt{falling}(id_{a0}\Leftarrow{}id_m~\mathtt{or}~\mathtt{false}))$ \\
      \end{tabular}$\Big)$}
    
    \infer2[\fontsize{7}{10}\selectfont\textsc{PsElab}]{
      {\fontsize{12}{12}\selectfont
        $\mathcal{D}_\mathcal{H},\Delta_2,\sigma_{e2}\vdash$
        \vhdle|process(marked, {clk},| $\emptyset$, $\dots$)
        $\xrightarrow{ebeh}\Delta_2\cup(\mathtt{marked},\emptyset),\sigma_{e2}$}
    }
  \end{prooftree}
  \caption{The elaboration of the \texttt{marked} process defined in
    the behavior of the \texttt{tl} design.}
  \label{fig:ex-elab-marked-ps}
\end{figure}


The \texttt{marked} process is elaborated in the context
$\mathcal{D}_\mathcal{H},\Delta_2,\sigma_{e2}$ where $\Delta_2$ and
$\sigma_{e2}$ are the partially-built elaborated design and default
design state at a certain point of the elaboration of the behavioral
part of the \texttt{tl} design. The elaboration of a process statement
associates the process identifier to a local variable environment in
the $Ps$ sub-environment of the being-built elaborated design. The
local variable environment is built out of the variable declaration
list of the process. Here, the \texttt{marked} process has an empty
variable declaration list. Thus, the binding
(\texttt{marked},$\emptyset$) is added in the $Ps$ sub-environment of
$\Delta_2$.  % As other concurrent statements could have been
% elaborated before the \texttt{marked} process, the $\Delta_2$ and
% $\sigma_{e2}$ of the above rule are not necessarily equal the
% $\Delta_2$ and $\sigma_{e2}$ of the rightmost premise of
% Figure~\ref{fig:ex-elab-tl}.
% The elaboration of the behavior of a design only acts on the $Ps$ and
% $Comps$ sub-environments of the being-built elaborated design, and
% only of the component store of the being-built default design
% state. Thus, we have $Gens(\Delta_2)=Gens(Delta)$,
% $Ins(\Delta_2)=Ins(\Delta)$, $Outs(\Delta_2)=Outs(\Delta)$ and
% $Sigs(\Delta_2)=Sigs(\Delta)$ where $\Delta$ is the final elaborated
% design. Also, we have $\mathcal{S}(\sigma_{e2})=\mathcal{S}(\sigma_e)$
% where $\sigma_e$ is the final default design state.

The elaboration of process statement also performs static
type-checking on the process statement body leveraging the
$\mathtt{valid}_{ss}$ relation. The following rule details the static
type-checking of the statement body of the \texttt{marked} process
(rightmost premise of the rule presented in
Figure~\ref{fig:ex-elab-marked-ps}). To keep the example within a
reasonable size, we do not detail the context of all rules. We
annotate the rule names to describe the side conditions associated to
a derivation.

\begin{figure}[H]
  \begin{prooftree}[template={\fontsize{11}{12}\selectfont\inserttext}]

    % First rst block
    \infer0[\fontsize{7}{10}\selectfont\textsc{False}]{$\mathtt{false}\xrightarrow{e}\bot$}
    \infer0[\fontsize{7}{10}\selectfont\textsc{IsBool}]{$\bot\in_c{}bool$}
    \infer2[\fontsize{7}{10}\selectfont\textsc{WTSig}$^1$]{$\vdash{}\mathtt{valid}_{ss}(id_{a0}\Leftarrow{}\mathtt{false})$}

    % Second rst block
    \infer0[\fontsize{7}{10}\selectfont\textsc{Sig}$^{2}$]{$id_m\xrightarrow{e}\bot$}
    \infer0[\fontsize{7}{10}\selectfont\textsc{False}]{$\mathtt{false}\xrightarrow{e}\bot$}
    \infer2[\fontsize{7}{10}\selectfont\textsc{BoolBinOp}]{$\vdash{}id_m~\mathtt{or}~\mathtt{false}\xrightarrow{e}\bot$}
    \infer0[\fontsize{7}{10}\selectfont\textsc{IsBool}]{$\bot\in_c{}bool$}
    \infer2[\fontsize{7}{10}\selectfont\textsc{WTSig}$^{1}$]{$\vdash{}\mathtt{valid}_{ss}(id_{a0}\Leftarrow{}id_m~\mathtt{or}~\mathtt{false})$}
    \infer1[\fontsize{7}{10}\selectfont\textsc{WTFalling}]{$\vdash{}\mathtt{valid}_{ss}(\mathtt{falling}(id_{a0}\Leftarrow{}id_m~\mathtt{or}~\mathtt{false}))$}
    
    % Conclusion
    \infer2[\fontsize{7}{10}\selectfont\textsc{WTRst}]{$\Delta_2,\sigma_{e2},\emptyset\vdash$
      $\mathtt{valid}_{ss}\Big($
      \begin{tabular}{@{}l@{}}
        $\mathtt{rst}~(id_{a0}\Leftarrow{}\mathtt{false})$ \\
        $(\mathtt{falling}(id_{a0}\Leftarrow{}id_m~\mathtt{or}~\mathtt{false}))$ \\
      \end{tabular}$\Big)$}
    
  \end{prooftree}
  \caption{Static type-checking of the \texttt{marked} process
    statement body.}
  \label{fig:ex-wt-marked-ps}
\end{figure}

\noindent{}where:
\begin{enumerate}
\item $\Delta_2(id_{a0})=bool$
\item $\sigma_{e2}(id_m)=\bot$
\end{enumerate}

At the end of the elaboration of the \texttt{tl} design's behavior,
the $Ps$ sub-environment of $Delta$ is as follows:
$Ps(\Delta):=\{(\mathtt{marked},\emptyset),
(\mathtt{fired},\emptyset)\}$

\paragraph{Elaboration of a component instantiation statement}

The rule of Figure~\ref{fig:ex-elab-pci-idp} presents the elaboration
of the place component instance $id_p$ belonging to the behavior of
the \texttt{tl} design.

\begin{figure}[H]
  \begin{prooftree}[template={\fontsize{9}{10}\selectfont\inserttext}, separation=4pt]

    % Dimensioning function
    \hypo{}
    \ellipsis{}{}
    \infer1{$\emptyset\vdash{}g_p\xrightarrow{emapg}\mathcal{M}$}

    % Elaboration of the place design
    \hypo{}
    \ellipsis{}{}
    \infer1[\fontsize{7}{10}\selectfont]{
      $\mathcal{D}_\mathcal{H},\mathcal{M}\vdash\mathtt{design}~\mathtt{place}\dots\xrightarrow{elab}\Delta_p,\sigma_p$
    }

    % Valid ipmap
    \hypo{}
    \ellipsis{}{}
    \infer1{
      $\Delta_2,\Delta_p,\sigma_{e2}\vdash{}\mathtt{valid}_{ipm}(i_p)$
    }

    % Valid opmap
    \hypo{}
    \ellipsis{}{}
    \infer1{
      $\Delta_2,\Delta_p\vdash{}\mathtt{valid}_{opm}(o_p)$
    }
    
    % Conclusion
    \infer4[\fontsize{7}{10}\selectfont\textsc{CompElab}$^{1}$]{
      {\fontsize{11}{13}\selectfont
        $\mathcal{D}_\mathcal{H},\Delta_2,\sigma_{e2}\vdash$
        $\mathtt{comp}~(id_p,\mathtt{place},g_p,i_p,o_p)\xrightarrow{ebeh}$
        $\Delta_2\cup(id_p,\Delta_p),\sigma_{e2}\cup(id_p,\sigma_p)$}
    }
    
  \end{prooftree}
  \caption{The elaboration of the $id_p$ component instance defined in
    the behavior of the \texttt{tl} design.}
  \label{fig:ex-elab-pci-idp}
\end{figure}

\noindent{}where:

\begin{enumerate}
\item
  \begin{tabular}{@{}l@{}}
    $id_p\notin\Delta_2$ \\
    $id_p\notin{}\sigma_{e2}$ \\
    $\mathtt{place}\in\mathcal{D}_\mathcal{H}$ \\
    $\mathcal{M}\subseteq{}Gens(\Delta_p)$ \\
  \end{tabular}
\end{enumerate}

The elaboration of a component instantiation statement is divided in
three parts. First, a dimensioning function is built out of the
generic map of the component instance. Figure~\ref{fig:ex-emapg} shows
a part of the creation of the dimensioning functioning $\mathcal{M}$
from the generic map of the component instance $id_p$. Basically, the
elaboration of a generic map is a transformation from a syntactic
construct, i.e. the generic map, into a function, i.e. the
dimensioning function $\mathcal{M}$. For each association of the
generic map, the elaboration checks that the actual part of the
association is a locally static expression (see
Section~\ref{subsubsec:loc-static}).

\begin{figure}[H]
  \begin{prooftree}[template={\fontsize{10}{13}\selectfont\inserttext}, separation=10pt]

    % Elab of assocg
    \infer0[\fontsize{7}{10}\selectfont\textsc{LSENat}]{$SE_l(1)$}
    \infer0[\fontsize{7}{10}\selectfont\textsc{Nat}]{$1\xrightarrow{e}1$}
    \infer2[\fontsize{7}{10}\selectfont\textsc{AssocGElab}$^1$]{$\emptyset\vdash{}(\mathtt{ian},1)\xrightarrow{emapg}$
      $\{(\mathtt{ian},1)\}$}

    % Ind. call
    \hypo{}
    \ellipsis{}{}
    \infer1[\fontsize{7}{10}\selectfont\textsc{GMElab}]{$\{(\mathtt{ian},1)\}\vdash{}(\mathtt{oan},1),(\mathtt{mm},1)\xrightarrow{emapg}$
      $\{(\mathtt{ian},1),(\mathtt{oan},1),(\mathtt{mm},1)\}$}
    
    % Conclusion 
    \infer2[\fontsize{7}{10}\selectfont\textsc{GMElab}]{$\emptyset\vdash{}(\mathtt{ian},1),(\mathtt{oan},1),(\mathtt{mm},1)\xrightarrow{emapg}$
      $\{(\mathtt{ian},1),(\mathtt{oan},1),(\mathtt{mm},1)\}$}
    
  \end{prooftree}
  \caption{The elaboration of the generic map of the $id_p$ component
    instance defined in the behavior of the \texttt{tl} design.}
  \label{fig:ex-emapg}
\end{figure}

\noindent{}where:
\begin{enumerate}
\item $\mathtt{ian}\notin\emptyset$
\end{enumerate}

The second step of the elaboration of a component instance is to
retrieve from the design store the design associated with the
component instance, and to elaborate this design. Here, the design
store is the \hilecop{} design store $\mathcal{D}_\mathcal{H}$, and
the design associated with $id_p$ is the \texttt{place} design. The
dimensioning function $\mathcal{M}$ sets the value of the generic
constants declared in the \texttt{place} design.  The full code of
\texttt{place} design is available in
Appendix~\ref{app:place-design}. In
Figures~\ref{fig:ex-elab-version-of-idp} and
\ref{fig:ex-default-state-of-idp}, we give the elaborated design
$\Delta_p$ and the default design state $\sigma_p$ resulting of the
elaboration of the \texttt{place} design given the dimensioning
function $\mathcal{M}$.

\begin{figure}[H]
\begin{tabular}{@{}rl@{}}
  $\Delta_p:=\{$ & \\
               & \begin{tabular}{@{}ll@{}}
                   $Gens := $ & $\{(\mathtt{ian}, (nat(0,\mathtt{NATMAX}, 1)),$ \\
                              & $(\mathtt{oan}, (nat(0,\mathtt{NATMAX}), 1))$ \\
                              & $(\mathtt{mm}, (nat(0,\mathtt{NATMAX}), 1))\}$ \\
                 \end{tabular} \\
               & \\
               & \begin{tabular}{@{}ll@{}}
                   $Ins :=$ & $\{(\mathtt{im}, nat(0,1)),$ \\
                            & $(\mathtt{iaw}, array(nat(0,255),0,0)),$ \\
                            & $(\mathtt{oat}, array(nat(0,2),0,0)),$ \\
                            & $(\mathtt{oaw}, array(nat(0,255),0,0)),$ \\
                            & $(\mathtt{itf}, array(bool,0,0)),$ \\
                            & $(\mathtt{otf}, array(bool,0,0))\},$ \\
                 \end{tabular} \\
               & \\
               & \begin{tabular}{@{}ll@{}}
                   $Outs :=$ & $\{(\mathtt{oav}, array(bool,0,0)),$ \\
                             & $(\mathtt{pauths}, array(bool,0,0)),$ \\
                             & $(\mathtt{rtt}, array(bool,0,0))\}$ \\
                 \end{tabular} \\
                 & \\
               & \begin{tabular}{@{}ll@{}}
                   $Sigs :=$ & $\{(\mathtt{sits}, nat(0,1)),$ \\
                             & $(\mathtt{sm}, nat(0,1)),$ \\
                             & $(\mathtt{sots}, nat(0,1))\},$ \\
                 \end{tabular} \\
               & \\
               & \begin{tabular}{@{}ll@{}}
                   $Ps :=$ & $\{(\mathtt{input\_tokens\_sum}, \{("v\_internal\_its", (nat(0,1), 0))\}),$ \\
                           & $(\mathtt{output\_tokens\_sum}, \{("v\_internal\_ots", (nat(0,1), 0))\})\}$ \\
                           & $(\mathtt{priority\_evaluation}, \{("v\_saved\_ots", (nat(0,1), 0))\})\}$ \\
                 \end{tabular} \\
               & \\
               & $Comps := \emptyset$ \\
  $\}$ &  \\
\end{tabular}

\caption[An elaborated version of the place design with a given
dimensioning function.]{An elaborated version of the place design
  built with the dimensioning function deduced from the generic map of
  the component instance $id_p$.}
\label{fig:ex-elab-version-of-idp}
\end{figure}

In $\Delta_p$, all the types associated with ports and internal
signals of the \texttt{place} design have been \emph{resolved};
i.e. the expressions qualifying the bounds of the range and index
constraints in type indications have been evaluated.  For example,
\texttt{array(boolean, 0, input_arcs_number-1)} is the type indication
associated with the \texttt{input_transitions_fired} input port
(i.e. \texttt{itf}) defined in the port clause of the \texttt{place}
design. The dimensioning function $\mathcal{M}$ sets the value of the
\texttt{input_arcs_number} (i.e. \texttt{ian}) generic constant to
1. After the elaboration, the type indication \texttt{array(boolean,
  0, input_arcs_number-1)} is thus transformed into the semantic type
$array(bool, 0, 0)$. Thus, we have
$\Delta_p(\mathtt{itf})=array(bool,0,0)$ in the resulting $\Delta_p$.

Figure~\ref{fig:ex-default-state-of-idp} shows the default design
state $\sigma_p$ of $\Delta_p$.

\begin{figure}[H]
  \begin{tabular}{@{}rl@{}}
    $\sigma_p:=\{$ & \\
                   & \begin{tabular}{@{}ll@{}}
                       $\mathcal{S} :=$ & $\{(\mathtt{im}, (0)),$ $(\mathtt{iaw}, (0)),$ $(\mathtt{oat}, (0)),$\\
                                        & $(\mathtt{oaw}, (0)),$ $(\mathtt{itf}, (\bot)),$ $(\mathtt{otf}, (\bot)),$ \\
                                        & $(\mathtt{oav}, (\bot)),$ $(\mathtt{pauths},(\bot)),$ $(\mathtt{rtt},(\bot))$ \\
                                        & $(\mathtt{sits}, 0),$ $(\mathtt{sm}, 0),$ $(\mathtt{sots}, 0)\},$ \\
                     \end{tabular} \\
                   & \\
                   & $\mathcal{C}:=\emptyset$ \\
                   & $\mathcal{E}:=\emptyset$ \\
    $\}$ &  \\
  \end{tabular}

\caption[An example of default design state for the place design.]{The
  default design state $\sigma_p$ of the elaborated design $\Delta_p$.}
\label{fig:ex-default-state-of-idp}
\end{figure}

The component store of design state $\sigma_p$ is empty as there are
no component instantiation statements in the behavior of the place
design. The same stands for the $Comps$ sub-environment of $\Delta_p$.
Also, the set of events of a default design state is always empty.

The final step in the elaboration of a component instantiation
statement is to check the well-formedness and the well-typedness of
the input and output port maps. The $\mathtt{valid}_{ipm}$ and
$\mathtt{valid}_{opm}$ relations, defined in
Section~\ref{subsec:valid-pm}, state the validity of the port maps.
The rule of Figure presents a part of the construction of the
$valid_{opm}$ relation applied to the output port map of the place
component instance $id_p$. Note that $\Delta_p$ is necessary to check
the validity of the output port map of $id_p$, as it holds the
correspondence between port identifiers and port types.

\begin{figure}[H]
  \centering
  \begin{prooftree}[template={\fontsize{10}{13}\selectfont\inserttext}]

    % Elab of assocop
    \infer0[\fontsize{7}{10}\selectfont\textsc{LSENat}]{$SE_l(0)$}
    \infer0[\fontsize{7}{10}\selectfont\textsc{Nat}]{$0\xrightarrow{e}0$}
    \infer0[\fontsize{7}{10}\selectfont\textsc{IsCNat}]{$0\in_c{}nat(0,0)$}
    \infer3[$^1$]{
      $\Delta_2,\Delta_p,\emptyset,\emptyset\vdash$
      $(\mathtt{oav}(0), id_{av})$
      $\xrightarrow{list_{opm}}$
      $\{(\mathtt{oav}, 0)\},\{id_{av}\}$
    }

    % Ind. call
    \hypo{}
    \ellipsis{\fontsize{7}{10}\selectfont\textsc{ListOPMCons}$_B$}{}

    % List OPM
    \infer2[\fontsize{7}{10}\selectfont\textsc{ListOPMCons}$_A$]{
      $\Delta_2,\Delta_p,\emptyset,\emptyset\vdash$
      \begin{tabular}{@{}l@{}}
        $(\mathtt{oav}(0), id_{av}),$ $(\mathtt{pauths}, \mathtt{open})$ \\
        $(\mathtt{rtt}(0), id_{rt}),$ $(marked, id_m)$ \\
      \end{tabular}
      $\xrightarrow{list_{opm}}$
      \begin{tabular}{@{}l@{}}
        $\{(\mathtt{oav}, 0), \mathtt{pauths},$ \\
        $(\mathtt{rtt}, 0), \mathtt{marked}\}$ \\
      \end{tabular},
      \begin{tabular}{@{}l@{}}
        $\{id_{av}, id_{rt}, id_m\}$ \\
      \end{tabular}
    }
    
    % Conclusion 
    \infer1[\fontsize{7}{10}\selectfont\textsc{ValidOPM}]{
      $\Delta_2,\Delta_p\vdash$
      $\mathtt{valid}_{opm}\Big($
      \begin{tabular}{@{}l@{}}
        $(\mathtt{oav}(0), id_{av}),$ $(\mathtt{pauths}, \mathtt{open})$ \\
        $(\mathtt{rtt}(0), id_{rt}),$ $(\mathtt{marked}, id_m)$ \\
      \end{tabular}
      $\Big)$
    }
    
  \end{prooftree}
  
  \vspace{40pt}
  
  \begin{prooftree}[template={\fontsize{10}{13}\selectfont\inserttext}, separation=10pt]
    % Ind. call
    \hypo{}
    \ellipsis{}{}
    \infer1[\fontsize{7}{10}\selectfont\textsc{ListOPMCons}$_B$]{
      {$\Delta_2,\Delta_p,\{(\mathtt{oav},0)\},\{id_{av}\}\vdash$
        \begin{tabular}{@{}l@{}}
          $(\mathtt{pauths}, \mathtt{open}),$ \\
          $(\mathtt{rtt}(0), id_{rt}),$ \\
          $(marked, id_m)$ \\
        \end{tabular}
        $\xrightarrow{list_{opm}}$
        \begin{tabular}{@{}l@{}}
          $\{(\mathtt{oav}, 0),$ \\
          $\mathtt{pauths},$ \\
          $(\mathtt{rtt}, 0),$ \\
          $\mathtt{marked}\}$ \\
        \end{tabular},
        \begin{tabular}{@{}l@{}}
          $\{id_{av}, id_{rt}, id_m\}$ \\
        \end{tabular}}
    }

  \end{prooftree}
  
  \caption{An example of validity checking performed on the output
    port map of the place component instance $id_p$. The bottom proof
    tree represents the top-right premise of the top proof tree.}
  \label{fig:ex-valid-opm}
\end{figure}

\noindent{}where:
\begin{enumerate}
\item
  \begin{tabular}{@{}l@{}}
    $\Delta_p(\mathtt{oav})=array(bool,0,0)$ \\
    $\Delta_2(id_{av})=bool$ \\
    $\mathtt{oav}\notin\emptyset$ and $(\mathtt{oav},0)\notin\emptyset$ \\
    $id_{av}\notin\emptyset$ \\
  \end{tabular}
\end{enumerate}

At the end of the elaboration of the \texttt{tl} design's behavior,
the $Comps$ sub-environment of $\Delta$ is as follows:
$Comps(\Delta):=\{(id_p,\Delta_p), (id_t,\Delta_t)\}$. Here,
$\Delta_t$ represents the elaborated version of the transition design
obtained from the elaboration of the transition component instance
$id_t$.

Also, at the end of the elaboration, the component store of $\sigma_e$
is as follows:
$\mathcal{C}(\sigma_e):=\{(id_p,\sigma_p),(id_t,\sigma_t)\}$. Here,
$\sigma_t$ is the default design state of the transition component
instance $id_t$.

\subsection{Simulation of the tl design}
\label{sec:ex-tl-sim}

Let us now present the rules pertaining to the simulation of the
\texttt{tl} design, that is, pertaining to the execution of the
\texttt{tl} design's behavior with respect to our formal simulation
algorithm.

\subsubsection{Initialization}
\label{sec:ex-init-tl}

The rule of Figure~\ref{fig:ex-init-tl} presents the initialization
phase in the proceeding of the simulation of the \texttt{tl}
design. The initialization phase builds the initial state of the
simulation. The first step of the initialization, formalized by the
$runinit$ relation, runs the processes and the internal behavior of
component instances exactly once (with the execution of the first part
of reset blocks). Then, a stabilization phase follows.

\begin{figure}[H]
  \centering
  \begin{prooftree}[template={\fontsize{12}{13}\selectfont\inserttext}]

    % Runinit
    \hypo{}
    \ellipsis{}{}
    \infer1[]{
      $\mathcal{D}_\mathcal{H},\Delta,\sigma_e\vdash$
      $cs_{tl}$
      $\xrightarrow{runinit}$
      $\sigma'$
    }
    
    % Stabilization
    \hypo{}
    \ellipsis{}{}
    \infer1[]{
      $\mathcal{D}_\mathcal{H},\Delta,\sigma_e\vdash$
      $cs_{tl}$
      $\xrightarrow{\rightsquigarrow}$
      $\sigma_0$
    }
    
    % Conclusion 
    \infer2[\fontsize{7}{10}\selectfont\textsc{Init}]{
      $\mathcal{D}_\mathcal{H},\Delta,\sigma_e\vdash$
      $cs_{tl}$
      $\xrightarrow{init}$
      $\sigma_0$
    }
    
  \end{prooftree}
  
  \caption{The initialization phase, first step of the simulation of
    the \texttt{tl} design.}
  \label{fig:ex-init-tl}
\end{figure}

The rule in Figure~\ref{fig:ex-runinit-tl} presents the execution of
the \texttt{tl} design's behavior during the $runinit$ phase. The
\texttt{tl} design's behavior is defined by the composition of
concurrent statements. Here, the \texttt{marked} process is at the
head of the behavior, whereas it is not the case in
Listing~\ref{lst:tl-design-ex}. We formally proved, with the \coq{}
proof assistant, that the $\mathtt{||}$ composition operator for
concurrent statements is commutative and associative with respect to
the $runinit$ relation. In Figure, the \texttt{marked} process is
executed and yields the state $\sigma'_e$. Then, the rest of the
\texttt{tl} design's behavior is executed and yields the state
$\sigma''_e$. Finally, the starting state $\sigma_e$ and the two
states $\sigma'_e$ and $\sigma''_e$ are merged into one by the
\texttt{merge} function.

\begin{figure}[H] \centering
  \begin{prooftree}[template={\fontsize{12}{13}\selectfont\inserttext}]
    
    % Runinit marked process
    \hypo{}
    \ellipsis{}{}
    \infer1[]{
      $\mathcal{D}_\mathcal{H},\Delta,\sigma_e\vdash$
      $\mathtt{process}(\mathtt{marked},\dots)$ $\xrightarrow{runinit}$
      $\sigma'_e$ }

    % Runinit tail
    \hypo{}
    \ellipsis{}{}
    \infer1[]{
      $\mathcal{D}_\mathcal{H},\Delta,\sigma_e\vdash$ $cs'_{tl}$
      $\xrightarrow{runinit}$ $\sigma''_e$ }

    % Conclusion
    \infer2[\fontsize{7}{10}\selectfont\textsc{CompRunInit}$^1$]{
      $\mathcal{D}_\mathcal{H},\Delta,\sigma_e\vdash$
      $\mathtt{process}(\mathtt{marked},\dots)~\mathtt{||}~cs'_{tl}$
      $\xrightarrow{runinit}$
      $\mathtt{merge}(\sigma_e,\sigma'_e,\sigma''_e)$ }

  \end{prooftree}

  \caption{The $runinit$ phase applied to the concurrent statements
    composing the behavior of the \texttt{tl} design.}
  \label{fig:ex-runinit-tl}
\end{figure}

\noindent{}where:
\begin{enumerate}
\item $\mathcal{E}'_e\cap\mathcal{E}''_e$
\end{enumerate}

In what follows, we detail the execution of a process statement and of
a component instantiation statement during the first part of the
initialization, i.e. the $runinit$ phase.

\paragraph{Execution of a process statement with the $runinit$ relation}

The rule in Figure shows the execution of the \texttt{marked} process
during the $runinit$ phase. The first part of the reset block defining
the statement body of the \texttt{marked} process is executed. This
first part assigns the expression $\mathtt{false}$ to signal
$id_{a0}$.

\begin{figure}[H] \centering
  \begin{prooftree}[template={\fontsize{12}{13}\selectfont\inserttext}]

    \infer0[\fontsize{7}{10}\selectfont\textsc{False}]{$\vdash\mathtt{false}\xrightarrow{e}\bot$}
    \infer0[\fontsize{7}{10}\selectfont\textsc{IsBool}]{$\bot\in_c{}bool$}
    
    % Execute fst part of rst block
    \infer2[\fontsize{7}{10}\selectfont\textsc{SigAssign}$^2$]{
      $\mathcal{D}_\mathcal{H},\Delta,\sigma_e,\sigma_e\vdash$
      $id_{a0}\Leftarrow\mathtt{false}$
      $\xrightarrow{ss_i}$
      $\sigma'_e,\emptyset$}

    % Execute rst block
    \infer1[\fontsize{7}{10}\selectfont\textsc{RstExec}]{
      $\mathcal{D}_\mathcal{H},\Delta,\sigma_e\vdash$
      $\mathtt{rst}(id_{a0}\Leftarrow\mathtt{false})(\dots)$
      $\xrightarrow{ss_i}$
      $\sigma'_e,\emptyset$ }

    % Conclusion
    \infer1[\fontsize{7}{10}\selectfont\textsc{PsRunInit}$^1$]{
      $\mathcal{D}_\mathcal{H},\Delta,\sigma_e\vdash$
      $\mathtt{process}(\mathtt{marked},\dots)$
      $\xrightarrow{runinit}$
      $\sigma'_e$ }
  \end{prooftree}

  \caption{The $runinit$ phase applied to the concurrent statements
    composing the behavior of the \texttt{tl} design.}
  \label{fig:ex-runinit-ps}
\end{figure}

\noindent{}where:

\begin{enumerate}
\item $\Delta(\mathtt{marked})=\emptyset$
\item
  \begin{tabular}{@{}l@{}}
    $\Delta(id_{a0})=bool$ \\
    $\sigma'_e={<}\mathcal{S}'_e,\mathcal{C}'_e,\mathcal{E}'_e{>}$ \\
    $\mathcal{S}'_e=\mathcal{S}_e(id_{a0})\leftarrow\bot$ \\
    $\mathcal{E}'_e=\mathcal{E}_e\cup(\mathcal{S}_e\dcap\mathcal{S}'_e)$ \\
  \end{tabular}
\end{enumerate}

In the side conditions of the \textsc{SigAssign} rule, an new event
set $\mathcal{E}'_e$ is computed based on the event set
$\mathcal{E}_e$ joined to the expression
$\mathcal{S}_e\dcap\mathcal{S}'_e$. This expression returns the set of
signals with a different value between signal store $\mathcal{S}_e$
and singal store $\mathcal{S}'_e$. The only signal that possibly has a
different value from $\mathcal{S}_e$ to $\mathcal{S}'_e$ is the
assigned signal $id_{a0}$. Thus, this expression is a shorthand to
test if the value of signal $id_{a0}$ has changed after the execution
of the signal assignment statement. If it is the case, then the event
set receives the signal identifier $id_{a0}$; $id_{a0}$ is then an
eventful signal. In the present case, the value of signal $id_{a0}$
was $\bot$ at state $\sigma_e$ and is still $\bot$ after the execution
of the signal assignment statement. Therefore, no event is registered
on signal $id_{a0}$. When states $\sigma_e$, $\sigma'_e$ and
$\sigma''_e$ will be merged (cf. Figure~\ref{fig:ex-runinit-tl}), if
$id_{a0}$ is part of the event set of state $\sigma''_e$, then, the
merged state will return the value associated to $id_{a0}$ in state
$\sigma''_e$. We would have
$\mathtt{merge}(\sigma_e,\sigma'_e,\sigma''_e)(id_{a0})=\sigma''_e(id_{a0})$.
However, signal $id_{a0}$ would be a potentially multiply-driven
signal because both the \texttt{marked} process and the concurrent
statement $cs'_{tl}$ (cf. Figure~\ref{fig:ex-runinit-tl}) assigns the
signal value.

\paragraph{Execution of a component instantiation statement with the $runinit$ relation}

The rule of Figure presents the execution of the \texttt{place}
component instance $id_p$ during the $runinit$ phase. The execution of
a component instantiation statement is pretty much the same in all the
phases of the simulation algorithm. The difference lies in the choice
of the relation used to execute of the internal behavior of the
component instance. During the $runinit$ phase, it is the $runinit$
relation that executes the internal behavior of component instances;
during the falling edge phase, it is the $\downarrow$ relation that
executes the internal behaviors, etc.

\begin{figure}[H] \centering
  \begin{prooftree}[template={\fontsize{11}{13}\selectfont\inserttext}]
    
    % Execute mapip
    \hypo{}
    \ellipsis{}{}
    \infer1{
      $\Delta,\Delta_p,\sigma_e,\sigma_p\vdash$
      $i_p$
      $\xrightarrow{mapip}$
      $\sigma'_p$}

    % Execute internal behavior
    \hypo{}
    \ellipsis{}{}
    \infer1{
      $\mathcal{D}_\mathcal{H},\Delta_p,\sigma'_p\vdash$
      $cs_p$
      $\xrightarrow{runinit}$
      $\sigma''_p$ }

    % Execute mapop
    \hypo{}
    \ellipsis{}{}
    \infer1{
      $\Delta,\Delta_p,\sigma_e,\sigma''_p\vdash$
      $o_p$
      $\xrightarrow{mapop}$
      $\sigma'_e$}

    % Conclusion
    \infer3[\fontsize{7}{10}\selectfont\textsc{CompRunInit}$^1$]{
      $\mathcal{D}_\mathcal{H},\Delta,\sigma_e\vdash$
      $\mathtt{comp}(id_p,\mathtt{place},g_p,i_p,o_p)$
      $\xrightarrow{runinit}$
      $\sigma''_e$ }
  \end{prooftree}

  \caption[The execution of the \texttt{place} component instance
    $id_p$ during the $runinit$ phase.]{The execution of the \texttt{place} component instance
    $id_p$ during the $runinit$ phase (first part of the
    initialization).}
  \label{fig:ex-runinit-pci-idp}
\end{figure}

\noindent{}where:

\begin{enumerate}
\item
  \begin{tabular}{@{}l@{}}
    $\sigma'_e={<}\mathcal{S}'_e,\mathcal{C}'_e,\mathcal{E}'_e{>}$ \\
    $\sigma''_e={<}\mathcal{S}'_e,\mathcal{C}''_e,\mathcal{E}'_e\cup(\mathcal{C}_e\dcap\mathcal{C}''_e){>}$ \\
    $\mathcal{C}''_e=\mathcal{C}'_e(id_p)\leftarrow{}\sigma''_p$ \\
  \end{tabular}
\end{enumerate}



The execution of a component instantiation statement is decomposed in
four parts. First, the input ports of the component instance receive
new values through the evaluation of the input port map. Second, the
internal behavior of the component instance is executed. Thirdly, the
evaluation of the output port map propagates the values coming from
the output interface to the component to the signals of the embedding
design. Finally, the component instance is assigned to its new
internal state in the component store of the embedding design; here,
$\sigma''_p$ is assigned to $id_p$ in component store
$\mathcal{C}''_e$. Moreover, if the new internal state of the
component instance is different from its older internal state, then
the component instance identifier is added to the event set of the
embedding design. Here, the expression
$\mathcal{C}_e\dcap\mathcal{C}''_e$ performs the state comparison; we
have:

{\raggedleft
  \begin{tabular}{ll}
    $\mathcal{C}_e\dcap\mathcal{C}''_e$ & $=\mathcal{C}_e\dcap(\mathcal{C}'_e\leftarrow\sigma''_p)$ \\
                                      & $=\mathcal{C}_e\dcap(\mathcal{C}_e\leftarrow\sigma''_p)$ \\
                                      &$=
                                        \begin{cases}
                                          \{id_p\}~\mathtt{if}~\sigma_p\neq\sigma''_p \\
                                          \emptyset~otherwise \\
                                        \end{cases}$
  \\
\end{tabular}
}

In the second line, we have $\mathcal{C}_e=\mathcal{C}'_e$ because the
evaluation of the output port map (performed by the $mapop$ relation)
does not change the component store.

The rule of Figure~\ref{fig:ex-eval-ipm} gives a part of the
evaluation of the input port map of $id_p$.

\begin{figure}[H]
  \centering
  \begin{prooftree}[template={\fontsize{10}{13}\selectfont\inserttext}]

    % MapipSimple
    \infer0[\fontsize{7}{10}\selectfont\textsc{Nat}]{$\Delta,\sigma_e\vdash{}1\xrightarrow{e}1$}
    \infer0[\fontsize{7}{10}\selectfont\textsc{IsCNat}]{$1\in_c{}nat(0,1)$}
    \infer2[\fontsize{7}{10}\selectfont\textsc{MapipSimple}$^1$]{
      $\Delta,\Delta_p,\sigma_e,\sigma_p\vdash$
      $(\mathtt{im}, 1)\xrightarrow{mapip}\sigma'_{p0}$
    }

    % MapipComp
    \hypo{}
    \ellipsis{}{}
    \infer1{
      $\Delta,\Delta_p,\sigma_e,\sigma'_{p0}\vdash$
      \begin{tabular}{@{}l@{}}
        $(\mathtt{iaw}(0), 1)$ \\
        $(\mathtt{oat}(0), 0),$ $(\mathtt{oaw}(0), 1),$ \\
        $(\mathtt{itf}(0), id_{ft}),$ $(\mathtt{otf}(0), id_{ft})$ \\
      \end{tabular}
      $\xrightarrow{mapip}$
      $\sigma'_p$
    }
    
    % Conclusion 
    \infer2[\fontsize{7}{10}\selectfont\textsc{MapipComp}]{
      $\Delta,\Delta_p,\sigma_e,\sigma_p\vdash$
      \begin{tabular}{@{}l@{}}
        $(\mathtt{im}, 1),$ $(\mathtt{iaw}(0), 1)$ \\
        $(\mathtt{oat}(0), 0),$ $(\mathtt{oaw}(0), 1),$ \\
        $(\mathtt{itf}(0), id_{ft}),$ $(\mathtt{otf}(0), id_{ft})$ \\
      \end{tabular}
      $\xrightarrow{mapip}$
      $\sigma'_p$
    }
    
  \end{prooftree}
  
  \caption{The evaluation of the input port map of the \texttt{place}
    component instance $id_p$.}
  \label{fig:ex-eval-ipm}
\end{figure}

\noindent{}where:
\begin{enumerate}
\item
  \begin{tabular}{@{}l@{}}
    $\Delta_p(\mathtt{im})=nat(0,1)$ \\
    $\sigma_p={<}\mathcal{S},\mathcal{C},\mathcal{E}{>}$ \\
    $\mathcal{S}'=\mathcal{S}(\mathtt{im})\leftarrow{}1$ \\
    $\sigma'_{p0}={<}\mathcal{S}',\mathcal{C},\mathcal{E}{>}$ \\
  \end{tabular}
\end{enumerate}

The evaluation of the input port map of $id_p$ changes the value of
the \texttt{initial\_marking} input port (i.e. \texttt{im}). We have
$\sigma_p(\mathtt{im})=0$ and $\sigma'_p(\mathtt{im})=1$. As the value
of one of its input port has changed, the \texttt{place} component
instance $id_p$ will be registered as an eventful component instance.

The rule of Figure~\ref{fig:ex-eval-opm} gives a part of the
evaluation of the output port map of $id_p$.

\begin{figure}[H]
  \centering
  \begin{prooftree}[template={\fontsize{10}{13}\selectfont\inserttext}]

    % MapipSimple
    \infer0[]{$\vdash{}0\xrightarrow{e}0$}
    \infer0{$0\in_c{}nat(0,0)$}
    \infer2[\fontsize{7}{10}\selectfont\textsc{IdxSig}$^2$]{$\Delta_p,\sigma''_p\vdash\mathtt{oav}(0)\xrightarrow{e_o}\top$}
    \infer0[\fontsize{7}{10}\selectfont\textsc{IsBool}]{$\top\in_c{}bool$}
    \infer2[$^1$]{
      $\Delta,\Delta_p,\sigma_e,\sigma''_p\vdash$
      $(\mathtt{oav}(0), id_{av})$
      $\xrightarrow{mapop}\sigma'_{e0}$
    }

    % MapipComp
    \hypo{}
    \ellipsis{}{}
    \infer1{
      $\Delta,\Delta_p,\sigma'_{e0},\sigma''_p\vdash$
      \begin{tabular}{@{}l@{}}
        $(\mathtt{pauths}, \mathtt{open})$ \\
        $(\mathtt{rtt}(0), id_{rt}),$ $(\mathtt{marked}, id_m)$ \\
      \end{tabular}
      $\xrightarrow{mapop}$
      $\sigma'_e$
    }
    
    % Conclusion 
    \infer2[\fontsize{7}{10}\selectfont\textsc{MapopComp}]{
      $\Delta,\Delta_p,\sigma_e,\sigma''_p\vdash$
      \begin{tabular}{@{}l@{}}
        $(\mathtt{oav}(0), id_{av}),$ $(\mathtt{pauths}, \mathtt{open})$ \\
        $(\mathtt{rtt}(0), id_{rt}),$ $(\mathtt{marked}, id_m)$ \\
      \end{tabular}
      $\xrightarrow{mapop}$
      $\sigma'_e$
    }
    
  \end{prooftree}
  
  \caption{The evaluation of the output port map of the \texttt{place}
    component instance $id_p$.}
  \label{fig:ex-eval-opm}
\end{figure}

\noindent{}where:
\begin{enumerate}
\item
  \begin{tabular}{@{}l@{}}
    $\Delta(id_{av})=bool$ \\
    $\sigma_e={<}\mathcal{S},\mathcal{C},\mathcal{E}{>}$ \\
    $\mathcal{S}'=\mathcal{S}(id_{av})\leftarrow{}\top$ \\
    $\mathcal{E}'=\mathcal{E}\cup(\mathcal{S}\dcap\mathcal{S}')$ \\
    $\sigma'_{e0}={<}\mathcal{S}',\mathcal{C},\mathcal{E}'{>}$ \\
  \end{tabular} (Rule~\textsc{MapopPartialToSimple}).
\item
  \begin{tabular}{@{}l@{}}
    $\Delta_p(\mathtt{oav})=array(bool,0,0)$ \\
    $\sigma''_p(\mathtt{oav})=(\top)$ \\
    $\mathtt{get\_at}(0,(\top))=\top$ \\
  \end{tabular}
\end{enumerate}

\subsubsection{Stabilization}
\label{sec:ex-stabilize-tl}

A stabilization phase happens after the $runinit$ during the
initialization phase, but also after the rising edge phase and the
falling edge phase in the course of a simulation cycle. The
stabilization phase executes the combinational parts of the design's
behavior. The \texttt{tl} design holds no combinational processes in
its behavior. The \texttt{marked} and \texttt{fired} processes are
both synchronous. To illustrate the execution of a combinational
process during a stabilization phase, let us consider the
\texttt{fired_evaluation} process defined in the behavior of
\texttt{transition} design. The \texttt{fired_evaluation} process will
be executed with the internal behavior of the \texttt{transition}
component instance $id_t$ during the stabilization phase.  The rule of
Figure presents the execution of the internal behavior of the
\texttt{transition} component instance $id_t$. As shown, the internal
behavior $cs_t$ is executed three times before reaching a stable
state. Here, the number of execution before stabilization is
arbitrary. In Figure, $\sigma_{t0}$ corresponds to the state of $id_t$
after the $runinit$ phase and after the evaluation of its input port
map. Remember that the evaluation of the input port map of a component
instance always precedes the execution of the internal behavior of the
same component. Since $\sigma_{t0}$ and $\sigma_{t1}$ are not stable
states, it means that their event set is not empty. Thus, we have
$\mathcal{E}(\sigma_{t0})\neq\emptyset$ and
$\mathcal{E}(\sigma_{t1})\neq\emptyset$. On the contrary,
$\sigma_{t2}$ is a stable state, and thus,
$\mathcal{E}(\sigma_{t2})=\emptyset$.



%%% Local Variables:
%%% mode: latex
%%% TeX-master: "../../main"
%%% End:


\section{Implementation of the \hvhdl{} syntax and semantics}
\label{sec:hvhdl-impl}
This section presents the implementation of the \hvhdl{} abstract
syntax, and also of the elaboration and the simulation semantics of
\hvhdl{} designs with the \coq{} proof assistant. The full code is
available under the \texttt{hvhdl} folder of the following repository:
\url{https://github.com/viampietro/ver-hilecop}.

\subsection{Implementation of the \texorpdfstring{\hvhdl{}}{H-VHDL}
  abstract syntax, elaborated design and design state}
\label{sec:impl-abss-and-env}

\paragraph{\hvhdl{} abstract syntax}

The implementation of the \hvhdl{} abstract syntax is naturally done
leveraging the \texttt{Inductive} construct of the \coq{} proof
assistant. The result is strictly similar to the formal definition of
the abstract syntax given in Section~\ref{sec:abstractSyntax}. The
reader can refer to the \texttt{AbstractSyntax.v} under the
\texttt{hvhdl} folder for the details of the implementation.

\paragraph{Elaborated design}

Listing~\ref{lst:elab-design-struct} presents the implementation of
the elaborated design structure
(cf. Definition~\ref{def:elab-design}). Two definitions are involved
in the implementation of the elaborated design structure.  The first
one defines the \texttt{SemanticObject} inductive type. Each
constructor of this type corresponds to a sub-environment of the
elaborated design. For instance, the \texttt{Generic} constructor
corresponds to the couple $(type\times{}value)$ associated with a
generic constant identifier in the $Gens$ sub-environment of
Definition~\ref{def:elab-design}. The \texttt{Process} constructor
corresponds to the local variable environment associated with the
process identifiers in the $Ps$ sub-environment. A local variable
environment is implemented by the \texttt{LEnv} type. The
\texttt{LEnv} type is a map between identifiers and couples
$(type\times{}value)$. Identifiers are implemented by the
\texttt{ident} type, an alias of the \texttt{nat} type. The
\texttt{type} and \texttt{value} types are the implementation of the
semantic $type$ and $value$ presented in
Table~\ref{tab:type-value}. The \texttt{ElDesign} type implements the
elaborated design structure. It is an alias to the \texttt{IdMap
  SemanticObject} type. The \texttt{IdMap} is the type of maps from
identifiers (i.e. belonging to the \texttt{ident} type) to instances
of the type passed as an input. Here, the input is the
\texttt{SemanticObject} type. Thus, an elaborated design is
implemented as a map between identifiers and terms of the
\texttt{SemanticObject} type. We leverage the \texttt{FMaps} module
defined in the \coq{} standard library to implement the \texttt{IdMap}
type. The \texttt{IdMap} type ensures that an identifier is only
mapped once. Thus, the implementation of the elaborated design
structure verifies that there are no intersection between the domains
of sub-environments. For instance, a generic constant identifier can
not be an input port identifier, and, as it is implemented, an
identifier \texttt{id} can not be mapped to a \texttt{Generic} object
and to an \texttt{Input} object in the same instance of
\texttt{ElDesign}.
 
\begin{lstlisting}[language=coq,label={lst:elab-design-struct},
caption={[\coq{} implementation of the elaborated design structure.]The implementation of the elaborated design structure with the \coq{} proof assistant.},framexleftmargin=1.5em,xleftmargin=2em,numbers=left,
numberstyle=\tiny\ttfamily]
Inductive SemanticObject : Type :=
| Generic (t : type) (v : value)
| Input (t : type)
| Output (t : type)
| Declared (t : type)
| Process (lenv : LEnv)
| Component ($\Delta_c$ : IdMap SemanticObject).

Definition ElDesign := IdMap SemanticObject.
\end{lstlisting}

\paragraph{Design state}

Listing~\ref{lst:design-state-impl} gives the implementation of the
design state structure through the definition of the \texttt{DState}
inductive type. The constructor of the \texttt{DState} type defines
three fields: \texttt{sigstore}, implementing the signal store
$\mathcal{S}$ of the design state, \texttt{compstore}, implementing
the component store $C$, and \texttt{events}, implementing the set of
events $\mathcal{E}$ of the design state. The \texttt{sigstore} field
is a map from identifiers to values. The \texttt{compstore} field is a
map from identifiers to design states, justifying the inductive
definition of the \texttt{DState} type. The \texttt{events} field is
an instance of the \texttt{IdSet} type. The \texttt{IdSet} is the type
of sets of identifiers (i.e. sets of natural numbers). The
\texttt{IdSet} type is defined leveraging the \texttt{MSets} module of
the \coq{} standard library.

\begin{lstlisting}[language=coq,label={lst:design-state-impl},
  caption={[\coq{} implementation of the design state structure.]The implementation of the design state structure with the \coq{} proof assistant.},framexleftmargin=1.5em,xleftmargin=2em,numbers=left,
  numberstyle=\tiny\ttfamily]
Inductive DState : Type := MkDState {
  sigstore  : IdMap value;
  compstore : IdMap DState;
  events    : IdSet;
}.
\end{lstlisting}

\subsection{Implementation of the elaboration phase}
\label{sec:impl-elaboration}

The design elaboration relation, as presented in
Section~\ref{subsubsec:design-elab}, is implemented in \coq{} by the
\texttt{edesign} relation. Listing~\ref{lst:design-elab-rel-impl}
presents the definition of the \texttt{edesign} relation as an
inductive type. As usual, a n-ary relation is implemented in \coq{} by
a type defined with $n$ parameters and projecting to the \texttt{Prop}
type. The \texttt{edesign} relation as five parameters. The first
parameter is the design store $\mathcal{D}$ of type \texttt{IdMap
  design}, i.e. a map from identifiers to \hvhdl{} designs as defined
by the abstract syntax. The second parameter is the dimensioning
function $\mathcal{M}$ of type \texttt{IdMap value}, i.e. a map from
identifiers to values. The third parameter is the design being
elaborated, of type \texttt{design}. The fifth and sixth parameters
are the elaborated design (of type \texttt{ElDesign} and the default
design state (of type \texttt{DState}) resulting from the elaboration.
In Listing~\ref{lst:design-elab-rel-impl}, the \texttt{EDesign}
constructor implements the \textsc{DesignElab} rule presented in
Section~\ref{subsubsec:design-elab}. From Line~7 to Line~10, the
constructor defines the premises of Rule~\textsc{DesignElab}.  The
empty elaborated design structure, denoted $\Delta_\emptyset$, is
implemented by the \texttt{EmptyElDesign} definition, and the empty
design state structure, denoted by $\sigma_\emptyset$, is implemented
by the \texttt{EmptyDState} definition. Line~13 implements the
conclusion of Rule~\textsc{DesignElab}.

\begin{lstlisting}[language=coq,label={lst:design-elab-rel-impl},
  caption={[\coq{} implementation of the design elaboration relation.]The implementation of the design elaboration relation with the \coq{} proof assistant.},framexleftmargin=1.5em,xleftmargin=2em,numbers=left,
  numberstyle=\tiny\ttfamily]
Inductive edesign ($\mathcal{D}$ : IdMap design) : IdMap value -> design -> ElDesign -> DState -> Prop :=
| EDesign :
    forall $\mathcal{M}$ $id_e$ $id_a$ gens ports sigs behavior
           $\Delta$ $\Delta'$ $\Delta''$ $\Delta'''$ $\sigma$ $\sigma'$ $\sigma''$,

      (* Premises *)
      egens EmptyElDesign $\mathcal{M}$ gens $\Delta$ ->
      eports $\Delta$ EmptyDState ports $\Delta'$ $\sigma$ ->
      edecls $\Delta'$ $\sigma$ sigs $\Delta''$ $\sigma'$ ->
      ebeh $\mathcal{D}$ $\Delta''$ $\sigma'$ behavior $\Delta'''$ $\sigma''$ ->
      
      (* Conclusion *)
      edesign $\mathcal{D}$ $\mathcal{M}$ (design_ $id_e$ $id_a$ gens ports sigs behavior) $\Delta'''$ $\sigma''$

with ebeh ($\mathcal{D}$ : IdMap design) : ElDesign -> DState -> cs -> ElDesign -> DState -> Prop :=
$\dots$
\end{lstlisting}

The \texttt{edesign} relation requires a mutually recursive definition
with the \texttt{ebeh} relation. The mutually recursive definition is
performed leveraging the \texttt{with} clause at the end of
Listing~\ref{lst:design-elab-rel-impl}. The \texttt{ebeh} relation
needs the \texttt{edesign} relation to elaborate the component
instances found in the behavior of a
design. Listing~\ref{lst:ebeh-rel-impl} gives the details of the
\texttt{with} clause defining the \texttt{ebeh} relation.

At Line~2, the \texttt{EBehPs} constructor implements the
Rule~\textsc{PsElab} defining the elaboration of a process statement
(cf. Section~\ref{subsubsec:ps-elab}). Lines~6 and 7 implement the
premises of the rule; the \texttt{evars} relation implements the
elaboration of the local variable declaration list of the process; the
\texttt{validss} relation implements the relation that type-checks the
statement body of the process. Lines~10 to 14 implement the side
conditions of the rule. The term \coqe|~NatMap.In| $id_p$ $\Delta$
implements the side condition $id_p\notin{}\Delta$. The
\coqe|NatMap.In id m| relation states that a given identifier
\texttt{id} is a key of the \texttt{m} map. At Line~13, the
\coqe|NatSet.In $id_s$ sl| term states that $id_s$ belongs to the
identifier set \texttt{sl}. At Line~14, the term \texttt{MapsTo $id_s$
  (Input t) $\Delta$} states that the identifier $id_s$ is mapped to
\texttt{Input t} in the elaborated design $\Delta$,
i.e. $Ins(\Delta)(id_s)=\mathtt{t}$. More generally, \texttt{MapsTo}
is a ternary relation stating that a given key \texttt{k} of type
\texttt{nat}, is mapped to a value \texttt{v} of a type \texttt{A}, in
a given map \texttt{m}, i.e. \coqe|Mapsto k v m|.  Line~17 implements
the conclusion of Rule~\textsc{PsElab}. The \texttt{NatMap.add}
function binds the process identifier $id_p$ to the term
\texttt{Process $\Lambda$} in the elaborated design $\Delta$,
i.e. $\Delta\cup(id_p,\Lambda)$.

At Line~19, the \texttt{EBehComp} constructor implements the
Rule~\textsc{CompElab}
(cf. Section~\ref{subsubsec:comp-inst-elab}). This rule describes the
elaboration of a component instantiation statement. Lines~24 to 27
implement the premises of the rule. Line~25 appeals to the
\texttt{edesign} relation to elaborate the \texttt{cdesign} design
associated with the component instance $id_c$; thence, the mutually
recursive definition with the \texttt{ebeh} relation. As it is stated
at Line~32, the \texttt{cdesign} design is associated to identifier
$id_e$, i.e. the entity identifier of component instance $id_c$, in
the design store $\mathcal{D}$. Lines~30 to 33 implement the side
conditions of the rule. Line~30 checks that the identifier $id_c$ is
not already bound to a semantic object in the elaborated design
$\Delta$. Line~31 checks that the identifier $id_c$ is not already
bound in the component store of $\sigma$. Line~33 checks that all
identifiers defined in the domain of map $\mathcal{M}$, i.e. the
dimensioning function, are bound to generic constants in the
elaborated design $\Delta_c$
(i.e. $\mathcal{M}\subseteq{}Gens(\Delta_c)$). Lines~36 to 38
implement the conclusion of Rule~\textsc{CompElab}. At Line~39, the
\texttt{cstore_add} function binds $id_c$ to design state $\sigma_c$
in the component store of state $\sigma$ and returns the resulting
state.

At Line~41, the \texttt{EBehNull} constructor implements
Rule~\textsc{CsNullElab}. At Line~43, the \texttt{EBehPar} constructor
implements Rule~\textsc{CsParElab}.

\begin{lstlisting}[language=coq,label={lst:ebeh-rel-impl},
  caption={[\coq{} implementation of the behavior elaboration relation.]The implementation of the \texttt{ebeh} behavior elaboration relation with the \coq{} proof assistant.},framexleftmargin=1.5em,xleftmargin=2em,numbers=left,
  numberstyle=\tiny\ttfamily]
with ebeh ($\mathcal{D}$ : IdMap design) : ElDesign -> DState -> cs -> ElDesign -> DState -> Prop :=
| EBehPs :
    forall $id_p$ sl vars stmt $\Lambda$ $\Delta$ $\sigma$,

      (* Premises *)
      evars $\Delta$ EmptyLEnv vars $\Lambda$ ->
      validss $\Delta$ $\sigma$ $\Lambda$ stmt ->

      (* Side conditions *)
      ~NatMap.In $id_p$ $\Delta$ ->

      (forall $id_s$,
          NatSet.In $id_s$ sl ->
          exists t, MapsTo $id_s$ (Declared t) $\Delta$ \/ MapsTo $id_s$ (Input t) $\Delta$) ->

      (* Conclusion *)
      ebeh $\mathcal{D}$ $\Delta$ $\sigma$ (cs_ps $id_p$ sl vars stmt) (NatMap.add $id_p$ (Process $\Lambda$) $\Delta$) $\sigma$

| EBehComp :
    forall $\Delta$ $\sigma$ $id_c$ $id_e$ gmap ipmap opmap
           $\mathcal{M}$ $\Delta_c$ $\sigma_c$ formals actuals cdesign,

      (* Premises *)
      emapg (NatMap.empty value) gmap $\mathcal{M}$ ->
      edesign $\mathcal{D}$ $\mathcal{M}$ cdesign $\Delta_c$ $\sigma_c$ ->
      validipm $\Delta$ $\Delta_c$ $\sigma$ ipmap formals ->
      validopm $\Delta$ $\Delta_c$ opmap formals actuals ->
      
      (* Side conditions *)
      ~NatMap.In $id_c$ $\Delta$ ->
      ~NatMap.In $id_c$ (compstore $\sigma$) ->
      MapsTo $id_e$ cdesign $\mathcal{D}$ ->
      (forall g, NatMap.In g $\mathcal{M}$ -> exists t v, MapsTo g (Generic t v) $\Delta_c$) ->
      
      (* Conclusion *)
      ebeh $\mathcal{D}$ $\Delta$ $\sigma$ (cs_comp $id_c$ $id_e$ gmap ipmap opmap)
                  (NatMap.add $id_c$ (Component $\Delta_c$) $\Delta$)
                  (cstore_add $id_c$ $\sigma_c$ $\sigma$)
           
| EBehNull: forall $\Delta$ $\sigma$, ebeh $\mathcal{D}$ $\Delta$ $\sigma$ cs_null $\Delta$ $\sigma$

| EBehPar:
    forall $\Delta$ $\Delta'$ $\Delta''$ $\sigma$ $\sigma'$ $\sigma''$ cstmt cstmt',
      ebeh $\mathcal{D}$ $\Delta$ $\sigma$ cstmt $\Delta'$ $\sigma'$ ->
      ebeh $\mathcal{D}$ $\Delta'$ $\sigma'$ cstmt' $\Delta''$ $\sigma''$ ->
      ebeh $\mathcal{D}$ $\Delta$ $\sigma$ (cs_par cstmt cstmt') $\Delta''$ $\sigma''$.

\end{lstlisting}

\subsection{Implementation of the simulation algorithm}
\label{sec:impl-simulation}

The full simulation relation (cf. Section~\ref{sec:full-sim})
formalizes the \hvhdl{} simulation algorithm. The \coq{}
implementation of the full simulation relation, presented in
Listing~\ref{lst:full-sim-impl}, is a strict translation of
Rule~\textsc{FullSim}. At Lines~14 and 15, the term \texttt{(behavior
  d)} represents the concurrent statements defining the behavior of
the \hvhdl{} design \texttt{d} (i.e. d.cs in the formal rule). Line~13
corresponds to the elaboration phase, Line~14 to the initialization
phase, and Line~15 to the main simulation loop.

\begin{lstlisting}[language=coq,label={lst:full-sim-impl},
  caption={[\coq{} implementation of the full simulation relation.]The implementation of the full simulation relation with the \coq{} proof assistant.},framexleftmargin=1.5em,xleftmargin=2em,numbers=left,
  numberstyle=\tiny\ttfamily]
Inductive fullsim
          ($\mathcal{D}$ : IdMap design)
          ($\mathcal{M}$ : IdMap value)
          ($E_p$ : nat -> Clk -> IdMap value)
          ($\tau$ : nat)
          ($\Delta$ : ElDesign) 
          (d : design) : list DState -> Prop :=
  
| FullSim :
    forall $\sigma_e$ $\sigma_0$ $\theta$,
      
      (* * Premises * *)
      edesign $\mathcal{D}$ $\mathcal{M}$ d $\Delta$ $\sigma_e$ ->
      init $\mathcal{D}$ $\Delta$ $\sigma_e$ (behavior d) $\sigma_0$ ->
      simloop $\mathcal{D}$ $E_p$ $\Delta$ $\sigma_0$ (behavior d) $\tau$ $\theta$ ->
                    
      (* * Conclusion * *)
      fullsim $\mathcal{D}$ $\mathcal{M}$ $E_p$ $\tau$ $\Delta$ d ($\sigma_0$ :: $\theta$).
\end{lstlisting}

The \texttt{simloop} relation appeals to the \texttt{simcycle} that
implements the simulation cycle relation defined in
Section~\ref{sec:sim-cycle}. Listing~\ref{lst:sim-cycle-impl} presents
the implementation of the \texttt{simcycle} relation. The
\texttt{simcycle} relation is a strict transcription of the
\textsc{SimCyc} rule. At Line~13, the \texttt{vrising} relation
implements the $\uparrow$ relation, i.e. the rising edge phase of the
cycle. At Line 15, the \texttt{vfalling} relation implements the
$\downarrow$ relation, i.e. the falling edge phase of the cycle. At
Lines~14 and 16, the \texttt{stabilize} relation implements the
$\rightsquigarrow$ relation, i.e. the stabilization phases of the
simulation cycle. At Lines~18 and 19, the \texttt{IsInjectedDState}
relation implements the \texttt{Inject} relation. Line~18 states that
the $\sigma_i$ state is the result of the injection of the map
\coqe|($E_p$ $\tau$)| in the signal store of state $\sigma$.

\begin{lstlisting}[language=coq,label={lst:sim-cycle-impl},
caption={[\coq{} implementation of the simulation cycle relation.]The
  implementation of the simulation cycle relation with the \coq{}
  proof
  assistant.},framexleftmargin=1.5em,xleftmargin=2em,numbers=left,
numberstyle=\tiny\ttfamily]
Inductive simcycle ($\mathcal{D}$ : IdMap design) ($E_p$ : nat -> IdMap value) 
    ($\Delta$ : ElDesign) ($\tau$ : nat) ($\sigma$ : DState) (behavior : cs) 
    ($\sigma'$ $\sigma''$ : DState) : Prop := 
| SimCycle : forall $\sigma_i$ $\sigma_\uparrow$ $\sigma_\downarrow$,
      
      (* * Premises * *)      
      vrising $\mathcal{D}$ $\Delta$ $\sigma_i$ behavior $\sigma_\uparrow$ ->
      stabilize $\mathcal{D}$ $\Delta$ $\sigma_\uparrow$ behavior $\sigma'$ ->
      vfalling $\mathcal{D}$ $\Delta$ $\sigma'$ behavior $\sigma_\downarrow$ ->
      stabilize $\mathcal{D}$ $\Delta$ $\sigma_\downarrow$ behavior $\sigma$'' ->

      (* * Side conditions * *)
      IsInjectedDState $\sigma$ ($E_p$ $\tau$) $\sigma_i$ ->
      
      (* * Conclusion * *)
      simcycle $\mathcal{D}$ $E_p$ $\Delta$ $\tau$ $\sigma$ behavior $\sigma'$ $\sigma''$.
\end{lstlisting}


%%% Local Variables:
%%% mode: latex
%%% TeX-master: "../../main"
%%% End:


\section{Conclusion}
\label{sec:hvhdl-concl}

%%% Local Variables:
%%% mode: latex
%%% TeX-master: "../main"
%%% End:

% \chapter{The \hilecop{} model-to-text transformation}
\label{chap:transformation}

\textsc{Research question of the chapter:}\\
\fbox{\parbox{\textwidth}{\sc Is there a proper way to express our
    (i.e, \hilecop{}'s) transformation function?}}

\section{Expressing transformation functions}
\label{sec:transf-lit-rev}
In this section, we present our literature review pertaining to
transformation functions in the context of formal verification. Here,
a transformation function is understood as any kind of mapping from a
source representation to a target representation, where the source and
target representations possess a behavior of their own (i.e. they are
executable). We use the same articles to perform our literature review
in this chapter and in the following chapter,
i.e. Chapter~\ref{chap:proof}. However, our research questions,
i.e. the questions we try to give an answer to while reading the
articles, and our presentation axis differ from one chapter to the
other.  Here, the following questions guide our reading:

\begin{itemize}
\item Is there a proper way to build a transformation function? Are
  there standards depending on the application domain?
\item How can we build a modular, extensible transformation function?
\item How can we build a transformation function that will ease the
  proof of semantic preservation?
\end{itemize}

The goal is to inspire ourselves with the works of the literature, and
to see how far the correspondence holds between our specific case of
transformation, and other cases of transformations. % The results of the
% literature review are presented in two parts. The two parts have
% been prepared based on the same material. The first part will be
% focusing on the expression of the transformation functions in the
% literature, and the second part will be focusing on the proof that
% these transformations are semantic preserving ones.
The material we used for the literature review is divided in three
categories. Each category covers a specific case of transformation
function. The three categories are:

\begin{itemize}
\item Compilers for generic programming languages
\item Compilers for hardware description languages
\item Model-to-model and model-to-text transformations
\end{itemize}

Note that, in the case of compilers for programming languages, the
term \textit{translation} is preferred over transformation to talk
about the generation of a target program from a source program.

\subsection{Building transformation functions}
\label{sec:build-transf}

As the authors state in \cite{Tan2016}, ``Although theoretically
possible, verifying a compiler that is not designed for verification
would be a prohibitive amount of work in practice.'' The question is
to know how to design such a compiler? How to anticipate the fact that
we will have to prove that the compiler is semantic preserving? Now,
let us consider these questions in the more general context of
transformation functions that map a source representation to a target
one.

\paragraph{Compilers for generic programming languages} In the context
of formally verified compilers for generic programming languages, the
translation from a source program to a target program is straight
forward. While descending recursively through the AST of the input
program, each construct of the source language is mapped to one or
many constructs of the target
language. Figure~\ref{fig:java-expr-to-java-bytecode} gives an example
of the translation from \java{} program expressions to \java{}
bytecode expressions, set in the context of a compiler for \java{}
programs written within the \isahol{} theorem prover
\cite{Strecker2002}. Here, the mapping between source and target
constructs is clearly defined.
\begin{figure}[H]
  \centering
  \includegraphics[keepaspectratio,width=.6\linewidth]{Figures/Proof/java-exprs-transl}
  \caption{Translation from \java{} expressions to \java{} bytecode expressions}
  \label{fig:java-expr-to-java-bytecode}
\end{figure}
In the works pertaining to the well-known \ccert{} project
\cite{Leroy2009, Blazy2006}, the many steps that compose the compiler
from C programs to assembly programs are also clearly mapping each
construct of source program to target program constructs.  Moreover,
the pattern matching possibilities offered by languages like \coq{},
\isa{}, \hol{} and other interactive theorem provers enable a clear
and concise implementation of compilers.

The cases of optimizing compilers like \cite{Leroy2009} and \cite{Tan}
show that, to avoid writing too complex functions when passing from a
source to a target program, the compilation is decomposed into many
passes. No more than 12 passes for the CakeML compiler, and up to 7
passes for \ccert{}. This is a way to keep the translation functions
simple enough in order to ease reasoning afterwards. Indeed, the more
the gap is important between the source representation and the target
one, the more the translation function will be complex.

Another point that is noticeable while expressing a translation
function is the necessity to keep a binding between the source and the
target representations. For instance, in \ccert, when passing from
transformed C programs to an RTL representation (based on registers
and control flow graphs), a binding function $\gamma$ links the
variables of a C program to the registers generated in the RTL
representation of the program. The binding is necessary for both the
translation and the proof of semantic preservation. During the
translation, it permits to replace the variables by their
corresponding registers in the RTL code. During the proof of semantic
preservation, the link that exists between a variable and a register
indicates which elements must be compared to prove that the execution
state of the source representation is similar to the execution state
of the target representation. The generation of this binding function
must be integrated to the design of the translation function.

In \cite{Leroy2009}, and \cite{Chlipala2010}, compilers are written
within the \coq{} proof assistant. Compilers are expressed using the
state-and-error monad, thus mimicking the traits of imperative
languages into a functional programming language setting. In
Section~\ref{sec:trans-alg}, we present the \hilecop{} transformation
in the form of an imperative pseudo-code algorithm. The
state-and-error monad is well-suited to the implementation of this
kind of algorithm with a functional language like \coq{}; thus, we
chose to apply this monad to our implementation of the transformation
algorithm (see Section~\ref{sec:trans-coq-impl}).

\paragraph{Compilers for hardware description languages}

The other category of compilers that we are interested in are
compilers for hardware description languages (HDL). The \hilecop{}
methodology's goal is the design of hardware circuits. For that
reason, we are interested in studying the case of compilers for
HDLs. However, one can notice that compiling an HDL program into a
lower level representation is one level of abstraction down compared
to the transformation we propose to verify. Indeed, it corresponds to
Step~3 in the \hilecop{} methodology
(cf. Figure~\ref{fig:hilecop-wf}), i.e. the transformation of VHDL
source code into an RTL representation.

In the context of formal verification applied to HDLs compilers, only
a few works describe the specificities of their translation function.

In \cite{Braibant2013}, the authors define the FeSi language (a
refinement of the BlueSpec language, a specification language for
hardware circuit behaviors), and its implementation within the \coq{}
proof assistant.  The authors present the syntax and semantics of the
FeSi language and of the RTL language which is the target language of
the compiler.  FeSi programs are composed of simple expressions, and
actions permitting to read or write from different types of memory
(registers). Therefore, the abstract syntax is divided into the
definition of expressions and the definition of actions, i.e: control
flow instructions and operations on memory. The RTL language is
composed of expressions and write operations to registers. The authors
are more interested in proving that a FeSi specification is
well-implemented by a given \coq{} program, than giving the details of
the translation from FeSi to RTL. However, the translation seems
straight-forward, and proceeds as usual by descending through the AST
of FeSi programs.

In \cite{Bourgeat2020}, the authors present a compiler for the
language Koîka, which is also a simpler version of the BlueSpec
language. A Koîka program is composed of a list of rules; each rule
describes actions that must be performed atomically. Actions are read
and write operations on registers. A Koîka program is accompanied by a
scheduler that specify an execution order for the rules. The described
compiler transforms Koîka programs into RTL descriptions of hardware
circuits. The translation function builds an RTL circuit by
descending recursively down the AST of rules. Each action is
translated into a specific RTL representation which are afterwards
composed together to get complex circuits. The translation becomes
trickier when it comes to decide the composition of RTL circuits to
respect the execution order prescribed by the scheduler.

In \cite{Bourke}, the authors present the verification of a compiler
toolchain from \textsf{Lustre} programs to an imperative language
(Obc), and from Obc to Clight.  The Clight target is the one defined
in \ccert{} \cite{Leroy2009}.  \textsf{Lustre} permits the definition
of programs composed of nodes that are executed synchronously.  Nodes
treat input streams and yield output streams of values.  A node body
is composed of sequence of equations that determine the values of
output streams based on the input.  Obc programs are composed of class
declarations. A class declaration has a vector of memory variables, a
vector of instances of other classes, and method declarations.  The
translation turns each node of a \textsf{Lustre} program into a class
of Obc accompanied by two methods: the reset method, for the
initialization of the streams, and the step method, for the update of
values resulting of a synchronous step.

In \cite{Loow2021}, the authors describe a compiler that transforms
Verilog programs into netlists targeting given FPGA models. Verilog
programs are a lot like VHDL programs; they describe a hardware
circuit behavior in terms of processes. A netlist is composed of
registers, variables and a list of cells corresponding to
combinational components. During the translation process, the
expressions of the Verilog programs are turned into netlist cells, and
the composition of statements leads to the creation of complex
circuits by means of cell composition.

\paragraph{Model transformations}

We will now present the works pertaining to model-to-model and
model-to-text transformations in the context of formal
verification. Because of the nature of the transformation we propose
to verify, i.e a model-to-text transformation, the following works are
of particular interest to us. We will focus here on the manner to
express transformations in the case of model-to-model and
model-to-text transformations. Also, we tried to find articles related
to model transformations involving Petri nets.

In \cite{Berramla2015}, the authors observe that Model-Driven
Engineering (MDE) is all about model transformation operations. They
propose to set a formal context within the \coq{} proof assistant to
verify that model transformations preserve the structure of the source
models into the target models. To illustrate their methodology, they
choose to transform UML state machine diagrams into Petri net
models. The translation rules from source to target models are
expressed within the setting of the OMG standard QVT language
(Query/View/Transform). The QVT language offers a formal way to
express model transformations, partly based on the Object Constraint
Language (OCL). The translation rules map the different kind of
structures that can be found in UML state diagrams to specific
structures of Petri nets. Even though the two models used as source
and target of transformations are executable, the authors leverage the
formal context provided by \coq{} to prove that the expressed
transformations preserve certain structural properties.

In \cite{Calegari2011}, the authors describe a process for model
transformation where transformation rules are expressed with the Atlas
Transformation Language (ATL). Transformation rules in ATL involve
both declarative (OCL) and imperative (match rules) instructions. The
authors show how the ATL rules can easily be translated into \coq{}
relations. An example is given on the kind of model-to-model
transformations that can be implemented that way. The example is a UML
class diagram to relational database model transformation.

In \cite{Combemale2009}, the authors explore the different ways to
give a formal semantics to a Domain-Specific Language (DSL) in the
context of MDE. Here, the syntax of a given DSL is expressed with a
meta-model.  An instantiation of this meta-model (a model) yields a
DSL program. The authors specify a transformation from a DSL model to
another executable model, thus providing an \textit{translational}
semantics to the DSL model.  The authors illustrate their approach
with a source DSL named xSPEM, which is a process description
language. The target models are timed PNs. The transformation is
expressed through a structural mapping; i.e, each element of an xSPEM
model is mapped to a particular PN: an activity is mapped to a subnet,
a resource to a single place, connection from activity to resource
through parameter is mapped to a connection of transitions and places
in the resulting PN\dots 

In \cite{Dyck2019}, the authors address the problem of expressing
model transformations by using transformation graphs. Precisely, the
kind of transformation graphs that are used are called Triple Graph
Grammar (TGG). A TGG is a triplet ${<}s,c,t{>}$ where the
``correspondence model $c$ explicitly stores correspondence
relationships between source model $s$ and target model $t$''.

The work described is \cite{Fronc2011} is really close to our own
verification task. The article describes how Coloured Petri Nets
(CPNs, specifically LLVM-labelled Petri nets) are transformed into
LLVM (Low Level Virtual Machine) programs representing the state space
(the graph of reachable markings) of these PNs. LLVM is a low-level
assembly language. The aim is to enable an efficient model-checking of
the CPNs.  LLVM-labelled PNs are CPNs whose places, transitions and
arcs have LLVM constructs for color domains. Places are labelled with
data types.  Transitions are labelled with boolean expressions that
correspond to the guard of the transition. Arcs are labelled by
multisets of expressions. A marking is a function that maps each place
to a multiset of values belonging to the place's type.  The authors
define data structures (multisets, sets, markings\dots) with
interfaces, i.e. sets of operations over structures, to represent the
Petri nets in LLVM.  They define interpretation functions that draw
equivalences between Petri nets objects and LLVM data structures.  The
authors define two algorithms: $\mathtt{fire\_t}$ and
$\mathtt{succ\_t}$ to compute the graph of reachable states.  These
are the functions that transform CPNs into concrete LLVM programs.

In \cite{Meghzili2017}, the author describes a transformation from UML
state machine diagrams to Coloured Petri Nets (CPNs). The aim is to
leverage the means of analysis provided by Petri nets to certify
certain properties over UML state machine diagrams. The authors want
to verify that the transformation preserve structural properties
between source and target models. The transformation function does not
use a standard setting as QVT or ATL, or transformation graphs. It is
expressed as a specific function written in \isahol.

In \cite{Yang2014}, the author presents a transformation from
Architecture Analysis and Design Language (AADL) models to Timed
Abstract State Machines (TASMs). AADL is a language widely used in
avionics to describe both hardware and software systems. AADL doesn't
have a lot of tools to analyze and simulate the designed systems;
therefore transforming AADL models into TASM enables the use of an
important toolbox for analysis, and simulation. The transformation
from AADL to TASMs are described with ATL rules.

\paragraph{Discussions on how to build transformation functions in the
  context of semantic preservation}

Transformation functions are mappings from a source representation to
a target representation. The more the mapping from source to target is
straight-forward the easier the comparison will be when proving that
the transformation is semantic preserving. Thus, in
\cite{Leroy2009,Tan,Chlipala2010} where complex case of optimizing
compilers are presented, the compilation is split into many simple
passes to ease the verification effort coming afterwards. In the case
of the \hilecop{} transformation, we are not yet concerned with the
optimization of the generated \vhdl{} code. Thus, our transformation
algorithm performs the generation of the target \hvhdl{} design in a
single pass. We do not need to use intermediary representations
between the input SITPN model and the generated \hvhdl{} design.

Also, while transforming source programs, the compiler must often
generate fresh constructs belonging to the target language (for
instance, generating a fresh RTL register for each variable referenced
in a source C program in \cite{Leroy2009}). The compiler must keep a
binding, that is, a memory of the mapping between the elements of the
source program and their mirror in the target program. This
consideration is of interest in our case of transformation where the
elements of SITPNs are also mirrored by elements in the generated
\hvhdl{} design.

It remains hard to establish a standard way to express a
transformation function as it really depends on the form of the input
and the output representations. Compilers for programming languages
tend to be a lot more compositional than model transformations. Here,
the word \textit{compositional} means that the translation rules can
be split into simple and independent cases of translation, e.g.
translation of expressions, then translation of statements, then
translation of function bodies,\dots This is a huge advantage to
perform the proof of semantic preservation. Indeed, this decomposition
of a translation function permits to reason on simple translation
cases; yet, each of these translations cases yields a piece of target
code that can be executed or interpreted in an independent manner. In
the case of the \hilecop{}, we tried as much as possible to express
the transformation in a compositional way. First, we tried to devise
the transformation by building up transformation functions for each
element of the SITPN structure, i.e.: a transformation function for
the places, another for the transitions\dots However, due to the
interconnections that exist between the component instances of the
generated \hvhdl{} design, it is impossible to define transformation
functions that would yield stand-alone executable code.

In the world of models, there exist some standard formalisms to
express transformation rules (QVT, ATL, transformation graphs\dots).
However, the complexity of the transformation rules depends on the
richness of the elements composing the source model, and the distance
to the concepts of the target model. In our case, we were not able to
grab the perks of using such formalisms as QVT or ATL to devise our
transformation.




%%% Local Variables:
%%% mode: latex
%%% TeX-master: "../../main"
%%% End:


\section{Informal presentation of the transformation}
\label{sec:informal-transf}


\begin{itemize}
\item High-level presentation of the transformation\\
  \pnote{use the figures of the 2nd CSI}
\item point out the particularities of \hilecop{}'s transformation
  function compared to what have been presented in the lit. review.\\

  \pnote{don't know if the lit. review will give enough material to do
    so}

  \pnote{are there cases of model-to-text transformation close enough
    to our case to enable comparison?}
\item Give an example of transformation in the form of schematic figure.
\end{itemize}

\textsc{Conclusion of the section:} Hard to express our trans. fun. in
a way that facilitate compositional reasoning (if compositional
reasoning is at all possible), partly because component instances are
connected. However, the lit. review pointed out the necessity to keep
a binding between elements of the source models and elements of the
target programs (introduction of the $\gamma$ binder) to enable state
comparison, i.e on which relies the proof of semantic preservation.

\section{The transformation algorithm}
\label{sec:trans-alg}

\bnote{maybe talk about the fact that the transformation algorithm has
  been established by observing the behavior of the current \hilecop{}
  implementation, and through discussions with the \hilecop{} designer
  (a.k.a David Andreu)}

\begin{itemize}
\item Give the model-to-text transformation algorithm\\
  \pnote{give the alg. as established through our meetings, or maybe
    use a simpler alg. (without the generate infos part, maybe too
    much implementation oriented)}
\item Give the formal specification of the transformation, i.e which
  structural properties between SITPN elements and \hvhdl{} constructs\\
  \rnote{don't know if I'll have the time to do that, but it would be
    great to have it, and use it in the proof. Each time the proof
    says ``by construction'', then the reader can refer to the formal
    spec}
\item Give a concrete example of generated \hvhdl{} code from a simple
  SITPN
\end{itemize}

\section{\coq{} implementation of the \hilecop{} model-to-text transformation}
\label{sec:trans-coq-impl}

\pnote{talk about the state-and-error monad, the $\gamma$ binder, the
  use of generic list functions (fold, map, filter, iter\dots)}

\textsc{Conclusion of the chapter:} Our transformation function is not
really expressible in a compositional way as P and T instances are
bound together. It's a specific case of model-to-text
transformation. However, as much as the comparison with other works of
transformation verification holds, we have been trying to inspire
ourselves from the literature to implement the \hilecop{}
model-to-text transformation function in \coq{}.

%%% Local Variables:
%%% mode: latex
%%% TeX-master: "../main"
%%% End:

% \chapter{Proving semantic preservation in \hilecop{}}
\label{chap:proof}

\begin{todobox}
  \begin{itemize}
  \item Change $\sigma_{injr}$ and $\sigma_{injf}$ into $\sigma_i$.
  \item Define the $\mathtt{Inject_{\downarrow}}$ and
    $\mathtt{Inject_{\uparrow}}$ relations.
  \item Keep the $sitpn$ argument in the SITPN full execution
    relation, but remove it from the SITPN execution, cycle and state
    transition relations.
  \end{itemize}
\end{todobox}

\section{Preliminary Definitions}
\label{sec:prelim-defs}
% \begin{definition}[SITPN-to-\hvhdl{} Design Binder]
  Given a $sitpn\in{}SITPN$ and a \hvhdl{} design $d\in{}design$, a
  SITPN-to-\hvhdl{} design binder $\gamma\in{}WM(sitpn,d)$ is a tuple
  ${<}PMap,TMap,\mathcal{C}_{id},\mathcal{A}_{id},\mathcal{F}_{id},CMap,AMap,FMap{>}$
  where:
  \begin{itemize}
    \fontsize{9}{10}\selectfont
  \item
    $sitpn={<}P,T,pre,test,inhib,post,M_0,{\succ},\mathcal{A},\mathcal{C},\mathcal{F},
    \mathbb{A},\mathbb{C},\mathbb{F},{I_s}{>}$
  \item $d=$ \vhdle|design| \textit{$id_{ent}$ $id_{arch}$ gens ports sigs behavior}
  \item $PMap\in{}P\rightarrow{}P_{id}$ where $P_{id}=\{id~|~\mathtt{comp}(id,"place",gm,ipm,opm)\in{}behavior\}$
  \item $TMap\in{}T\rightarrow{}T_{id}$ where $T_{id}=\{id~|~\mathtt{comp}(id,"transition",gm,ipm,opm)\in{}behavior\}$
  \item $\mathcal{C}_{id}\subseteq\{id~|~(\mathtt{in}, id, t)\in{}ports\wedge{}id\notin\{"clk","rst"\}\}$
  \item $\mathcal{A}_{id}\subseteq\{id~|~(\mathtt{out}, id, t)\in{}ports\}$
  \item $\mathcal{F}_{id}\subseteq\{id~|~(\mathtt{out}, id, t)\in{}ports\}$
  \item $CMap\in\mathcal{C}\rightarrow\mathcal{C}_{id}$
  \item $AMap\in\mathcal{A}\rightarrow\mathcal{A}_{id}$
  \item $FMap\in\mathcal{F}\rightarrow\mathcal{F}_{id}$        
  \end{itemize}
\end{definition}

\begin{definition}[Similar Environments]
  For a given $sitpn\in{}SITPN$, a \hvhdl{} design $d\in{}design$, a
  design store
  $\mathcal{D}\in{}entity\mhyphen{}id\nrightarrow{}design$, an
  elaborated version $\Delta\in{}ElDesign(d,\mathcal{D})$ of design
  $d$, and a binder $\gamma\in{}WM(sitpn,d)$, the environment
  $E_p\in{}(\mathbb{N}\times{}\{\uparrow,\downarrow\})\rightarrow{}Ins(\Delta)\rightarrow{}value$,
  that yields the value of the primary input ports of $\Delta$ at a
  given simulation cycle and a given clock event, and the environment
  $E_c$, that yields the value of conditions of $sitpn$ at a given
  execution cycle, are similar, noted
  $\gamma\vdash{}E_p\stackrel{env}{=}E_c$, iff for all
  $\tau\in{}\mathbb{N}$, $clk\in\{\uparrow,\downarrow\}$,
  $c\in\mathcal{C}$, $id_c\in{}Ins(\Delta)$ s.t.  $\gamma(c)=id_c$,
  $E_p(\tau,clk)(id_c)=E_c(\tau)(c)$.
\end{definition}

\subsection{State Similarity}
\label{sec:state-sim}

\begin{definition}[General State Similarity]
  \label{def:state-sim}
  For a given $sitpn\in{}SITPN$, a \hvhdl{} design $d\in{}design$, a
  design store
  $\mathcal{D}\in{}entity\mhyphen{}id\nrightarrow{}design$, an
  elaborated design $\Delta\in{}ElDesign(d,\mathcal{D})$, and a binder
  $\gamma\in{}WM(sitpn,d)$, an SITPN state $s\in{}S(sitpn)$ and a
  design state $\sigma\in\Sigma(\Delta)$ are similar, written
  $\gamma\vdash{}s\sim\sigma$ iff
  \begin{enumerate}
  \item $\forall{}p\in{}P,id_p\in{}Comps(\Delta),\sigma_p\in{}\Sigma(\Delta(id_p))~s.t.~\gamma(p)=id_p$ and $\sigma(id_p)=\sigma_p,$\\
    $~s.M(p)=\sigma_p("s\_marking")$.
  \item
    $\forall{}t\in{}T_i,id_t\in{}Comps(\Delta),\sigma_t\in{}\Sigma(\Delta(id_t))~s.t.~\gamma(t)=id_t$ and $\sigma(id_t)=\sigma_t,$\\
    $upper(I_s(t))=\infty\land{}s.I(t)\le{}lower(I_s(t))\Rightarrow{}s.I(t)=\sigma_t("s\_tc")\land{}$\\
    $upper(I_s(t))=\infty\land{}s.I(t)>{}lower(I_s(t))\Rightarrow{}\sigma_t("s\_tc")=lower(I_s(t))\land{}$\\
    $upper(I_s(t))\neq\infty\land{}s.I(t)>{}upper(I_s(t))\Rightarrow{}\sigma_t("s\_tc")=upper(I_s(t))\land{}$\\
    $upper(I_s(t))\neq\infty\land{}s.I(t)\le{}upper(I_s(t))\Rightarrow{}s.I(t)=\sigma_t("s\_tc")$.
  \item
    $\forall{}t\in{}T_i,id_t\in{}Comps(\Delta),\sigma_t\in{}\Sigma(\Delta(id_t))~s.t.~\gamma(t)=id_t$
    and $\sigma(id_t)=\sigma_t,$\\
    $s.reset_t(t)=\sigma_t("s\_reinit\_time\_counter")$.
  \item
    $\forall{}c\in\mathcal{C},id_c\in{}Ins(\Delta)~s.t.~\gamma(c)=id_c,~s.cond(c)=\sigma(id_c)$.
  \item
    $\forall{}a\in\mathcal{A},id_a\in{}Outs(\Delta)~s.t.~\gamma(a)=id_a,~s.ex(a)=\sigma(id_a)$.
  \item
    $\forall{}f\in\mathcal{F},id_f\in{}Outs(\Delta)~s.t.~\gamma(f)=id_f,~s.ex(f)=\sigma(id_f)$.
  \end{enumerate}
\end{definition}

\begin{definition}[Post Rising Edge State Similarity]
  \label{def:post-re-state-sim}
  For a given $sitpn\in{}SITPN$, a \hvhdl{} design $d\in{}design$, an
  elaborated design $\Delta\in{}ElDesign(d)$, and a binder
  $\gamma\in{}WM(sitpn,d)$, a clock cycle count $\tau\in\mathbb{N}$,
  and an SITPN execution environment
  $E_c\in\mathbb{N}\rightarrow\mathcal{C}\rightarrow\mathbb{B}$, an
  SITPN state $s\in{}S(sitpn)$ and a design state
  $\sigma\in\Sigma(\Delta)$ are similar after a rising edge happening
  at clock cycle count $\tau$, written
  $\gamma,E_c,\tau\vdash{}s\stackrel{\uparrow}{\sim}\sigma$ iff
  \begin{enumerate}
  \item $\forall{}p\in{}P,id_p\in{}Comps(\Delta),\sigma_p\in{}\Sigma(\Delta(id_p))~s.t.~\gamma(p)=id_p$ and $\sigma(id_p)=\sigma_p,$\\
    $~s.M(p)=\sigma_p("s\_marking")$.
  \item
    $\forall{}t\in{}T_i,id_t\in{}Comps(\Delta),\sigma_t\in{}\Sigma(\Delta(id_t))~s.t.~\gamma(t)=id_t$ and $\sigma(id_t)=\sigma_t,$\\
    $upper(I_s(t))=\infty\land{}s.I(t)\le{}lower(I_s(t))\Rightarrow{}s.I(t)=\sigma_t("s\_tc")\land{}$\\
    $upper(I_s(t))=\infty\land{}s.I(t)>{}lower(I_s(t))\Rightarrow{}\sigma_t("s\_tc")=lower(I_s(t))\land{}$\\
    $upper(I_s(t))\neq\infty\land{}s.I(t)>{}upper(I_s(t))\Rightarrow{}\sigma_t("s\_tc")=upper(I_s(t))\land{}$\\
    $upper(I_s(t))\neq\infty\land{}s.I(t)\le{}upper(I_s(t))\Rightarrow{}s.I(t)=\sigma_t("s\_tc")$.
  \item
    $\forall{}t\in{}T_i,id_t\in{}Comps(\Delta),\sigma_t\in{}\Sigma(\Delta(id_t))~s.t.~\gamma(t)=id_t$
    and $\sigma(id_t)=\sigma_t,$\\
    $s.reset_t(t)=\sigma_t("s\_reinit\_time\_counter")$.
  \item
    $\forall{}a\in\mathcal{A},id_a\in{}Outs(\Delta)~s.t.~\gamma(a)=id_a,~s.ex(a)=\sigma(id_a)$.
  \item
    $\forall{}f\in\mathcal{F},id_f\in{}Outs(\Delta)~s.t.~\gamma(f)=id_f,~s.ex(f)=\sigma(id_f)$.
  \item
    $\forall{}t\in{}T,id_t\in{}Comps(\Delta),\sigma_t\in{}\Sigma(\Delta(id_t))~s.t.~\gamma(t)=id_t$
    and $\sigma(id_t)=\sigma_t,$\\
    $t\in{}Sens(s.M)\Leftrightarrow\sigma_t("s\_enabled")=\mathtt{true}$.
  \item
    $\forall{}t\in{}T,id_t\in{}Comps(\Delta),\sigma_t\in{}\Sigma(\Delta(id_t))~s.t.~\gamma(t)=id_t$
    and $\sigma(id_t)=\sigma_t,$\\
    $\sigma_t("s\_condition\_combination")=\prod_{c\in{}conds(t)}E_c(\tau,c)~.~\prod_{c\in{}\overline{conds}(t)}\mathtt{not}(E_c(\tau,c))$\\
    where $conds(t)=\{c\in\mathcal{C}~\vert~\mathbb{C}(t,c)=1\}$ and
    $\overline{conds}(t)=\{c\in\mathcal{C}~\vert~\mathbb{C}(t,c)=-1\}$.
  \end{enumerate}
\end{definition}

\begin{definition}[Post Falling Edge State Similarity]
  \label{def:post-fe-state-sim}
  For a given $sitpn\in{}SITPN$, a \hvhdl{} design $d\in{}design$, an
  elaborated design $\Delta\in{}ElDesign(d)$, and a binder
  $\gamma\in{}WM(sitpn,d)$, an SITPN state $s\in{}S(sitpn)$ and a
  design state $\sigma\in\Sigma(\Delta)$ are similar after a falling
  edge, written $\gamma\vdash{}s\stackrel{\downarrow}{\sim}\sigma$ iff
  $\gamma\vdash{}s\sim\sigma$ (Def.~\ref{def:state-sim}, general state
  similarity) and
  \begin{enumerate}
  \item $\forall{}p\in{}P,id_p\in{}Comps(\Delta),\sigma_p\in{}\Sigma(\Delta(id_p))~s.t.~\gamma(p)=id_p$ and $\sigma(id_p)=\sigma_p,$\\
    $\sum\limits_{t\in{}Fired(s)}pre(p,t)=\sigma_p("s\_output\_token\_sum")$.
  \item $\forall{}p\in{}P,id_p\in{}Comps(\Delta),\sigma_p\in{}\Sigma(\Delta(id_p))~s.t.~\gamma(p)=id_p$ and $\sigma(id_p)=\sigma_p,$\\
    $\sum\limits_{t\in{}Fired(s)}post(t,p)=\sigma_p("s\_input\_token\_sum")$.
  \item
    $\forall{}t\in{}T_i,id_t\in{}Comps(\Delta),\sigma_t\in{}\Sigma(\Delta(id_t))~s.t.~\gamma(t)=id_t$
    and $\sigma(id_t)=\sigma_t,$\\
    $t\in{}Fired(s)\Leftrightarrow\sigma_t("fired")$.
  \end{enumerate}
\end{definition}

\begin{definition}[Execution Trace Similarity]
  \label{def:exec-trace-sim}
  For a given $sitpn\in{}SITPN$, a \hvhdl{} design $d\in{}design$, a
  design store
  $\mathcal{D}\in{}entity\mhyphen{}id\nrightarrow{}design$, an
  elaborated design $\Delta\in{}ElDesign(d,\mathcal{D})$, and a binder
  $\gamma\in{}WM(sitpn,d)$, the execution trace
  $\theta_s\in{}\mathtt{list}(S(sitpn))$, and the simulation trace
  $\theta_\sigma\in\mathtt{list}(\Sigma(\Delta))$ are similar, written
  $\gamma\vdash{}\theta_s\sim\theta_\sigma$, according to the
  following rules:

  \begin{tabular}{@{}l}
    {\fontsize{9}{11}\selectfont\textsc{SimTraceNil}} \\
    
    {\begin{prooftree}[template={\fontsize{9}{11}\selectfont\inserttext}]        
        \infer0{$\gamma\vdash{}[~]\sim{}[~]$}
      \end{prooftree}} 
  \end{tabular}
  \begin{tabular}{@{}l}
    {\fontsize{9}{11}\selectfont\textsc{SimTraceCons}} \\
    
    {\begin{prooftree}[template={\fontsize{9}{11}\selectfont\inserttext}]

        \hypo{$\gamma\vdash{}s\sim\sigma$}
        \hypo{$\gamma\vdash{}\theta_s\sim{}\theta_\sigma$}
        \infer2{$\gamma\vdash{}(s :: \theta_s)\sim{}(\sigma :: \theta_\sigma)$}
      \end{prooftree}} 
  \end{tabular}
  
\end{definition}

%%% Local Variables:
%%% mode: latex
%%% TeX-master: "main"
%%% End:


\section{Behavior Preservation Theorem}
\label{sec:beh-pres-thm}
% \subsection{Proof Notations}
\label{sec:pf-notations}

\begin{itemize}
\item Frame box for pending goals:
  \framebox{$\forall{}n\in\mathbb{N},~n>0\lor{}n=0$}
\item Red frame box for completed goals: \colorbox{red!20}{$\mathtt{true}=\mathtt{true}$}
\item Green frame box for induction hypotheses: \begin{ih}$\forall{}n\in\mathbb{N},~n+1>0$\end{ih}
\item \textbf{CASE} to denote a case during a proof by case analysis.
\end{itemize}

\subsection{Behavior Preservation Theorem and Proof}
\label{sec:beh-pres-thm-pf}

%%%%%%%%%%%%%%%%%%%%%%%%%%%%%%%%%%%%%%%%%%%
%%%%%% BEHAVIOR PRESERVATION THEOREM %%%%%%
%%%%%%%%%%%%%%%%%%%%%%%%%%%%%%%%%%%%%%%%%%%

\begin{thm}[Behavior Preservation]
  \label{thm:beh-pres-hyps}
  For all $sitpn\in{}SITPN$, $d\in{}design$, $\gamma\in{}WM(sitpn,d)$,
  $\tau\in\mathbb{N}$,
  $E_c\in{}\mathbb{N}\rightarrow{}\mathcal{C}\rightarrow{}\mathbb{B}$,
  $\theta_s\in\mathtt{list}(S(sitpn))$ s.t.
  \begin{itemize}
  \item SITPN $sitpn$ translates into design $d$:
    $\lfloor{}sitpn\rfloor_\mathcal{H}=(d,\gamma)$
  \item SITPN $sitpn$ yields the execution trace $\theta_s$ after
    $\tau$ execution cycles in environment $E_c$:\\
    $E_c,\tau\vdash{}sitpn\xrightarrow{full}\theta_s$.
  \end{itemize}
  
  \noindent{}then there exists
  $\Delta\in{}ElDesign(d,\mathcal{D}_\mathcal{H})$ s.t. for all
  $E_p\in{}(\mathbb{N}\times{}\{\uparrow,\downarrow\})\rightarrow{}Ins(\Delta)\rightarrow{}value$
  verifying
  \begin{itemize}
  \item Simulation/Execution environments are similar:
    $\gamma\vdash{}E_p\stackrel{env}{=}E_c$.
  \end{itemize}
  then there exists $\theta_\sigma\in\mathtt{list}(\Sigma(\Delta))$
  s.t.
  \begin{itemize}
  \item Under the \hilecop{} design store $\mathcal{D}_\mathcal{H}$
    and with an empty generic constant dimensioning function, design d
    yields the simulation trace $\theta_\sigma$ after $\tau$
    simulation cycles, starting from its initial state:\\
    $\mathcal{D}_\mathcal{H},\Delta,\emptyset,E_p,\tau\vdash{}\mathrm{d}\xrightarrow{full}\theta_\sigma$
  \item Traces $\theta_s$ and $\theta_\sigma$ are similar: $\theta_s\sim\theta_\sigma$
  \end{itemize}
\end{thm}

\begin{proof}
  \framebox{$\exists\Delta,~\forall{}E_p,~\gamma\vdash{}E_p\stackrel{env}{=}E_c,~\exists\theta_\sigma,$
    $\mathcal{D}_\mathcal{H},\Delta,\emptyset,E_p,\tau\vdash{}\mathrm{d}\xrightarrow{full}\theta_\sigma\land\theta_s\sim\theta_\sigma$}\\
  
  \noindent{}By definition of the \hvhdl{} full simulation relation:\\
  $\mathcal{D}_\mathcal{H},\Delta,\emptyset,E_p,\tau\vdash{}\mathrm{d}\xrightarrow{full}\theta_\sigma\equiv$
  $\exists\sigma_e,\sigma_0\in\Sigma(\Delta),~\mathcal{D}_\mathcal{H},\emptyset\vdash{}d\srarrow{elab}{\fontsize{6}{8}\selectfont}(\Delta,\sigma_{e})$
  and $\mathcal{D}_\mathcal{H},\Delta,\sigma_{e}\vdash{}d.cs\srarrow{init}{\fontsize{6}{8}\selectfont}\sigma_0$\\
  and $\mathcal{D}_\mathcal{H},E_p,\Delta,\tau,\sigma_0\vdash{}\mathrm{d.cs}\rightarrow\theta_\sigma$. \\

  \noindent{}Use \nameref{thm:elab-ex}, \nameref{thm:init-ex} and
  \nameref{thm:sim-ex} theorems to show that there exists a $\Delta$,
  $\theta_\sigma$, $\sigma_e$ and $\sigma_0$ such that
  $\mathcal{D}_\mathcal{H},\emptyset\vdash{}d\srarrow{elab}{\fontsize{6}{8}\selectfont}(\Delta,\sigma_{e})$
  and
  $\mathcal{D}_\mathcal{H},\Delta,\sigma_{e}\vdash{}d.cs\srarrow{init}{\fontsize{6}{8}\selectfont}\sigma_0$
  and
  $\mathcal{D}_\mathcal{H},E_p,\Delta,\tau,\sigma_0\vdash{}\mathrm{d.cs}\rightarrow\theta_\sigma$.\\

  \noindent{}Use \nameref{thm:full-bisim} theorem to show traces similarity.
  
\end{proof}

\begin{thm}[Elaboration]
  \label{thm:elab-ex}
  For all $sitpn\in{}SITPN$, $d\in{}design$, $\gamma\in{}WM(sitpn,d)$
  s.t.
  \begin{itemize}
  \item $\lfloor{}sitpn\rfloor_\mathcal{H}=(d,\gamma)$
  \end{itemize}
  \noindent{}then there exists
  $\Delta\in{}ElDesign(d,\mathcal{D}_\mathcal{H}),\sigma_e\in\Sigma(\Delta)$ s.t. 
  \begin{itemize}
  \item
    $\mathcal{D}_\mathcal{H},\emptyset\vdash{}d\srarrow{elab}{\fontsize{6}{8}\selectfont}(\Delta,\sigma_{e})$
  \end{itemize}
\end{thm}

\begin{thm}[Initialization]
  \label{thm:init-ex}
  For all $sitpn\in{}SITPN$, $d\in{}design$, $\gamma\in{}WM(sitpn,d)$,
  $\Delta\in{}ElDesign(d,\mathcal{D}_\mathcal{H})$,
  $\sigma_e\in\Sigma(\Delta)$ s.t.
  \begin{itemize}
  \item $\lfloor{}sitpn\rfloor_\mathcal{H}=(d,\gamma)$ and
    $\mathcal{D}_\mathcal{H},\emptyset\vdash{}d\srarrow{elab}{\fontsize{6}{8}\selectfont}(\Delta,\sigma_{e})$
  \end{itemize}
  \noindent{}then there exists $\sigma_0\in\Sigma(\Delta)$ s.t. 
  \begin{itemize}
  \item $\sigma_0$ is the initial simulation state:
    $\mathcal{D}_\mathcal{H},\Delta,\sigma_{e}\vdash{}d.cs\srarrow{init}{\fontsize{6}{8}\selectfont}\sigma_0$
  \end{itemize}
\end{thm}

\begin{thm}[Simulation]
  \label{thm:sim-ex}
  For all $sitpn\in{}SITPN$, $d\in{}design$, $\gamma\in{}WM(sitpn,d)$,
  $\Delta\in{}ElDesign(d,\mathcal{D}_\mathcal{H})$,
  $\sigma_e,\sigma_0\in\Sigma(\Delta)$ s.t.
  \begin{itemize}
  \item $\lfloor{}sitpn\rfloor_\mathcal{H}=(d,\gamma)$ and
    $\mathcal{D}_\mathcal{H},\emptyset\vdash{}d\srarrow{elab}{\fontsize{6}{8}\selectfont}(\Delta,\sigma_{e})$
    and
    $\mathcal{D}_\mathcal{H},\Delta,\sigma_{e}\vdash{}d.cs\srarrow{init}{\fontsize{6}{8}\selectfont}\sigma_0$
  \end{itemize}
  
  \noindent{}then for all
  $E_p\in{}(\mathbb{N}\times{}\{\uparrow,\downarrow\})\rightarrow{}Ins(\Delta)\rightarrow{}value$,
  $\tau\in\mathbb{N}$, there exists
  $\theta_\sigma\in\mathtt{list}(\Sigma(\Delta))$ s.t.
  \begin{itemize}
  \item Design $d$ yields the simulation trace $\theta_\sigma$ after
    $\tau$ simulation cycles, starting from initial state $\sigma_0$:\\
    $\mathcal{D}_\mathcal{H},E_p,\Delta,\tau,\sigma_0\vdash{}\mathrm{d.cs}\rightarrow\theta_\sigma$
  \end{itemize}
\end{thm}

\subsection{Bisimulation Theorem and Proof}
\label{sec:bisim-thm-and-proof}

%%%%%%%%%%%%%%%%%%%%%%%%%%%%%%%%%%%%%%%%%%%
%%%%%%%%%% FULL BISIMULATION THM %%%%%%%%%%
%%%%%%%%%%%%%%%%%%%%%%%%%%%%%%%%%%%%%%%%%%%

\begin{thm}[Full Bisimulation]
  \label{thm:full-bisim}
  For all $sitpn\in{}SITPN$, $d\in{}design$, $\gamma\in{}WM(sitpn,d)$,
  $\tau\in\mathbb{N}$,
  $E_c\in{}\mathbb{N}\rightarrow{}\mathcal{C}\rightarrow{}\mathbb{B}$,
  $\theta_s\in\mathtt{list}(S(sitpn))$,
  $\Delta\in{}ElDesign(d,\mathcal{D}_\mathcal{H})$,
  $E_p\in{}(\mathbb{N}\times{}\{\uparrow,\downarrow\})\rightarrow{}Ins(\Delta)\rightarrow{}value$,
  $\theta_\sigma\in\mathtt{list}(\Sigma(\Delta))$ s.t.
  \begin{itemize}
  \item $\lfloor{}sitpn\rfloor_\mathcal{H}=(d,\gamma)$
  \item $\gamma\vdash{}E_p\stackrel{env}{=}E_c$ 
  \item $E_c,\tau\vdash{}sitpn\xrightarrow{full}\theta_s$
  \item $\mathcal{D}_\mathcal{H},\Delta,\emptyset,E_p,\tau\vdash{}\mathrm{d}\xrightarrow{full}\theta_\sigma$
  \end{itemize}
  then
 $\theta_s\sim\theta_\sigma$

\end{thm}

\begin{proof}
  Case analysis on $\tau$ (2 CASES).

  \begin{itemize}
  \item \textbf{CASE} $\tau=0$. By definition of the SITPN full execution and
    the \hvhdl{} full simulation relations:
    \begin{itemize}
    \item $\mathcal{D}_\mathcal{H},\emptyset\vdash{}d\srarrow{elab}{\fontsize{6}{8}\selectfont}(\Delta,\sigma_{e})$
    \item
      $\Delta,\sigma_{e}\vdash{}d.cs\srarrow{init}{\fontsize{6}{8}\selectfont}\sigma_0$
    \item $\theta_s=[s_0]$ and $\theta_\sigma=[\sigma_0]$
    \end{itemize}
    \framebox{$\gamma\vdash{}s_0\sim{}\sigma_0$} (by def. of similar execution trace relation).
    Solved by applying Lemma~\nameref{lem:sim-init-states}.
  \item \textbf{CASE} $\tau>0$. By definition of the SITPN full execution and
    the \hvhdl{} full execution relations:
    \begin{itemize}
    \item $E_c,\tau\vdash{}s_0\srarrow{\uparrow_0}{\fontsize{6}{8}\selectfont}s_0$
    \item $E_c,\tau\vdash{}s_0\srarrow{\downarrow}{\fontsize{6}{8}\selectfont}s$
    \item $E_c,\tau-1\vdash{}sitpn,s\rightarrow\theta_s$
    \item $\mathcal{D}_\mathcal{H},\emptyset\vdash{}d\srarrow{elab}{\fontsize{6}{8}\selectfont}(\Delta,\sigma_{e})$
    \item $\Delta,\sigma_{e}\vdash{}d.cs\srarrow{init}{\fontsize{6}{8}\selectfont}\sigma_0$
    \item $E_p,\Delta,\tau,\sigma_0\vdash{}\mathrm{d.cs}\rightarrow\theta$
    \end{itemize}
    \framebox{$\gamma\vdash{}(s_0 :: s :: \theta_s)\sim{}(\sigma_0 :: \theta)$}\\\\
    By definition of the \hvhdl{} full simulation relation, we know:
    \begin{itemize}
    \item $E_p,\Delta,\tau,\sigma_0\vdash{}d.cs\xrightarrow{\uparrow,\downarrow}\sigma$
    \item $E_p,\Delta,\tau-1,\sigma\vdash{}\mathrm{d.cs}\rightarrow\theta_\sigma$
    \end{itemize}
    where $\theta=\sigma :: \theta_\sigma$.\\
    
    Rewriting $\theta$ as $\sigma :: \theta_\sigma$, \fbox{$\gamma\vdash{}(s_0 :: s :: \theta_s)\sim{}(\sigma_0 :: \sigma :: \theta_\sigma)$}\\\\
    3 subgoals (by def. of \nameref{def:exec-trace-sim}).
    \begin{enumerate}
    \item $\gamma\vdash{}s_0\sim\sigma_0$ (solved by applying Lemma~\nameref{lem:sim-init-states}).
    \item $\gamma\vdash{}s\sim\sigma$ (solved by applying Lemma~\nameref{lem:first-cycle}).
    \item $\gamma\vdash{}\theta_s\sim\theta_\sigma$ (solved by applying Lemma~\nameref{lem:simulation}).
    \end{enumerate}
  \end{itemize}
\end{proof}

%%%%%%%%%%%%%%%%%%%%%%%%%%%%%%%
%%%%%% FIRST CYCLE LEMMA %%%%%%
%%%%%%%%%%%%%%%%%%%%%%%%%%%%%%%

\begin{lemma}[First Cycle]
  \label{lem:first-cycle}
  For all
  $sitpn\in{}SITPN,d\in{}design,\gamma\in{}WM(sitpn,d),s\in{}S(sitpn),
  \Delta\in{}ElDesign(d,\mathcal{D}_\mathcal{H}),$
  $\sigma_{e},\sigma_0,\sigma\in{}\Sigma(\Delta)$,
  $E_c\in{}\mathbb{N}\rightarrow{}\mathcal{C}\rightarrow{}\mathbb{B}$,
  $E_p\in{}(\mathbb{N}\times{}\{\uparrow,\downarrow\})\rightarrow{}Ins(\Delta)\rightarrow{}value$,
  assume that:
  \begin{itemize}
  \item $\lfloor{}sitpn\rfloor_\mathcal{H}=(d,\gamma)$ and
    $\mathcal{D}_\mathcal{H},\emptyset\vdash{}d\srarrow{elab}{\fontsize{6}{8}\selectfont}(\Delta,\sigma_{e})$
    and $\gamma\vdash{}E_p\stackrel{env}{=}E_c$
  \item $\sigma_0$ is the initial state of $\Delta$: 
    $\Delta,\sigma_{e}\vdash{}d.cs\srarrow{init}{\fontsize{6}{8}\selectfont}\sigma_0$
  \item First execution cycle for $d$: $E_p,\Delta,\tau,\sigma_0\vdash{}d.cs\xrightarrow{\uparrow,\downarrow}\sigma$
  \item
    Particular first execution cycle for $sitpn$ (first rising edge is idle):\\
    $E_c,\tau\vdash{}s_0\srarrow{\uparrow_0}{\fontsize{6}{8}\selectfont}s_0$
    and
    $E_c,\tau\vdash{}s_0\srarrow{\downarrow}{\fontsize{6}{8}\selectfont}s$
  \end{itemize}
  then $\gamma\vdash{}s\stackrel{\downarrow}{\sim}{}\sigma$.
\end{lemma}

\begin{proof}
  Let's show that the first execution cycle leads to two states
  verifying the \nameref{def:post-fe-state-sim} relation:
  \fbox{$\gamma\vdash{}s\stackrel{\downarrow}{\sim}{}\sigma$.}\\

  \noindent{}By definition of the \hvhdl{} cycle relation, we have:
  \begin{itemize}
  \item $\mathtt{Inject}_\uparrow(\sigma_0, E_p, \tau, \sigma_{injr})$
    and
    $\Delta,\sigma_{injr}\vdash\mathrm{d.cs}\xrightarrow{\uparrow}\sigma_r$
    and
    $\Delta,\sigma_r\vdash\mathrm{d.cs}\xrightarrow{\theta}\sigma'$
  \item
    $\mathtt{Inject}_\downarrow(\sigma', E_p, \tau, \sigma_{injf})$
    and
    $\Delta,\sigma_{injf}\vdash\mathrm{d.cs}\xrightarrow{\downarrow}\sigma_f$
    and
    $\Delta,\sigma_f\vdash\mathrm{d.cs}\xrightarrow{\theta'}\sigma$
  \end{itemize}
  
  \noindent{}Then, we can apply the \nameref{lem:fe} lemma to solve \fbox{$\gamma\vdash{}s\stackrel{\downarrow}{\sim}{}\sigma$.}\\

  \noindent{}One premise of the \nameref{lem:fe} lemma remains to be
  proved:
  \fbox{$\gamma,E_c,\tau\vdash{}s_0\stackrel{\uparrow}{\sim}{}\sigma'$.}\\

  \noindent{}Then, we can apply the \nameref{lem:fst-re} lemma to
  solve \fbox{$\gamma,E_c,\tau\vdash{}s_0\stackrel{\uparrow}{\sim}{}\sigma'$.}
  
\end{proof}

%%%%%%%%%%%%%%%%%%%%%%%%%%%%%%%%%%%%%%%%
%%%%%%%%%% BISIMULATION LEMMA %%%%%%%%%%
%%%%%%%%%%%%%%%%%%%%%%%%%%%%%%%%%%%%%%%%

\begin{lemma}[Bisimulation]
  \label{lem:simulation}
  For all $sitpn$, $d$, $\gamma$, $E_p$, $E_c$, $\tau$, $s$, $\theta_s$,
  $\sigma$, $\theta_\sigma$, $\Delta$, $\sigma_e$, assume that:
  \begin{itemize}
  \item $\lfloor{}sitpn\rfloor_\mathcal{H}=(d,\gamma)$ and
    $\gamma\vdash{}E_p\stackrel{env}{=}E_c$ and
    $\mathcal{D}_\mathcal{H},\emptyset\vdash{}d\srarrow{elab}{\fontsize{7}{9}\selectfont}\Delta,\sigma_e$
  \item Starting states are similar as intended after a falling edge:
    $\gamma\vdash{}s\stackrel{\downarrow}{\sim}\sigma$
  \item $E_c,\tau\vdash{}sitpn,s\rightarrow\theta_s$
  \item $E_p,\Delta,\tau,\sigma\vdash{}d.cs\rightarrow\theta_\sigma$
  \end{itemize}
  then $\gamma\vdash{}\theta_s\sim{}\theta_\sigma$.
\end{lemma}

\begin{proof}
  Induction on $\tau$.
  \begin{itemize}
  \item Base case, $\tau=0$: traces are empty, trivial.
  \item Induction case, $\tau > 0$:
    \begin{ih}
      $\forall{}s,\sigma,\theta_s,\theta_\sigma$ s.t.
      $\gamma\vdash{}s\stackrel{\downarrow}{\sim}\sigma$ and
      $E_c,\tau-1\vdash{}sitpn,s\rightarrow\theta_s$ and
      $E_p,\Delta,\tau-1,\sigma\vdash{}d.cs\rightarrow\theta_\sigma$
      then
      $\gamma\vdash{}\theta_s\sim{}\theta_\sigma$.
    \end{ih}
    By definition of the SITPN execution and the \hvhdl{} simulation
    relations for $\tau>0$: 
    \begin{itemize}
    \item $E,\tau\vdash{}sitpn,s\xrightarrow{\uparrow,\downarrow}{}s'$ and $E_c,\tau-1\vdash{}sitpn,s\rightarrow\theta_s$.
    \item $E_p,\Delta,\tau,\sigma\vdash{}\mathrm{d.cs}\xrightarrow{\uparrow,\downarrow}{\fontsize{7}{9}\selectfont}\sigma'$ and
      $E_p,\Delta,\tau-1,\sigma\vdash{}d.cs\rightarrow\theta_\sigma$.
    \end{itemize}
    \framebox{$\gamma\vdash{}(s' :: \theta_s)\sim{}(\sigma' :: \theta_\sigma)$}.\\\\
    2 subgoals (by def. of \nameref{def:exec-trace-sim}).
    \begin{enumerate}
    \item \framebox{$\gamma\vdash{}s'\sim{}\sigma'$} (solved with \nameref{lem:step}).
    \item \framebox{$\gamma\vdash{}\theta_s\sim{}\theta_\sigma$} (solved with
      \nameref{lem:step} and IH).
    \end{enumerate}
    
  \end{itemize}
\end{proof}

%%%%%%%%%%%%%%%%%%%%%%%%%%%%%%%%
%%%%%%%%%% STEP LEMMA %%%%%%%%%%
%%%%%%%%%%%%%%%%%%%%%%%%%%%%%%%%

\begin{lemma}[Step]
  \label{lem:step}
  For all $sitpn$, $d$, $\gamma$, $E_p$, $E_c$, $\tau$, $s$,
  $s''$, $\sigma$, $\sigma''$, $\Delta$, $\sigma_e$, assume that:
  \begin{itemize}
  \item $\lfloor{}sitpn\rfloor_\mathcal{H}=(d,\gamma)$ and
    $E_p\stackrel{env}{=}E_c$ and
    $\mathcal{D}_\mathcal{H},\emptyset\vdash{}d\srarrow{elab}{\fontsize{7}{9}\selectfont}\Delta,\sigma_e$
  \item $\gamma\vdash{}s\stackrel{\downarrow}{\sim}\sigma$ 
  \item From state $s$ to $s''$ in one execution cycle:
    $E_c,\tau\vdash{}sitpn,s\srarrow{\uparrow,\downarrow}{\fontsize{7}{9}\selectfont}s''$
  \item From state $\sigma$ to $\sigma''$ in one simulation cycle:
    $E_p,\Delta,\tau,\sigma\vdash{}d.cs\srarrow{\uparrow,\downarrow}{\fontsize{7}{9}\selectfont}\sigma''$
  \end{itemize}
  then $\gamma\vdash{}s''\stackrel{\downarrow}{\sim}{}\sigma''$.
\end{lemma}

\begin{proof}
  By def. of the SITPN and \hvhdl{} cycle relations:
  \begin{itemize}
  \item $E_c,\tau\vdash{}sitpn,s\xrightarrow{\uparrow}s'$ and
    $E_c,\tau\vdash{}sitpn,s'\xrightarrow{\downarrow}s''$
  \item
    $\mathtt{Inject}_\uparrow(\sigma, E_p, \tau, \sigma_{injr})$
    and
    $\Delta,\sigma_{injr}\vdash\mathrm{d.cs}\xrightarrow{\uparrow}\sigma_r$
    and
    $\Delta,\sigma_r\vdash\mathrm{d.cs}\xrightarrow{\theta}\sigma'$
  \item
    $\mathtt{Inject}_\downarrow(\sigma', E_p, \tau, \sigma_{injf})$
    and
    $\Delta,\sigma_{injf}\vdash\mathrm{d.cs}\xrightarrow{\downarrow}\sigma_f$
    and
    $\Delta,\sigma_f\vdash\mathrm{d.cs}\xrightarrow{\theta'}\sigma''$
  \end{itemize}
  Solved by applying \nameref{lem:re} and then ``Falling Edge'' lemmas.
\end{proof}

%%% Local Variables:
%%% mode: latex
%%% TeX-master: "../../main"
%%% End:


\section{Initial States}
\label{sec:init-states}
% %%%%%%%%%%%%%%%%%%%%%%%%%%%%%%%%%%%%%%%%%%%%%%%
%%%%%% SIMILAR INITIAL STATES HYPOTHESES %%%%%%
%%%%%%%%%%%%%%%%%%%%%%%%%%%%%%%%%%%%%%%%%%%%%%%

\begin{definition}[Initial State Hypotheses]
  \label{def:init-states-hyps}
  Given an $sitpn\in{}SITPN$, $d\in{}design$,
  $\gamma\in{}WM(sitpn,d)$,
  $\Delta\in{}ElDesign(d,\mathcal{D}_\mathcal{H}),\sigma_{e},\sigma_0\in{}\Sigma(\Delta)$,
  assume that:
  \begin{itemize}
  \item SITPN $sitpn$ translates into design $d$:
    $\lfloor{}sitpn\rfloor_\mathcal{H}=(d,\gamma)$
  \item $\Delta$ is the elaborated version of $d$, $\sigma_e$ is the
    default state of $\Delta$, i.e, state of $\Delta$ where all signals have their default value:\\
    $\mathcal{D}_\mathcal{H},\emptyset\vdash{}d\srarrow{elab}{\fontsize{6}{8}\selectfont}(\Delta,\sigma_{e})$
    
  \item $\sigma_0$ is the initial state of $\Delta$: 
    $\Delta,\sigma_{e}\vdash{}d.cs\srarrow{init}{\fontsize{6}{8}\selectfont}\sigma_0$
  \end{itemize}
\end{definition}

\def\inithyps{For all $sitpn\in{}SITPN$, $d\in{}design$,
  $\gamma\in{}WM(sitpn,d)$,
  $\Delta\in{}ElDesign(d,\mathcal{D}_\mathcal{H}),\sigma_{e},\sigma_0\in{}\Sigma(\Delta)$
  that verify the hypotheses of Def.~\ref{def:init-states-hyps},}

%%%%%%%%%%%%%%%%%%%%%%%%%%%%%%%%%%%%%%%%%%
%%%%%% SIMILAR INITIAL STATES LEMMA %%%%%%
%%%%%%%%%%%%%%%%%%%%%%%%%%%%%%%%%%%%%%%%%%

\begin{lemma}[Similar Initial States]
  \label{lem:sim-init-states}
  \inithyps{} then $\gamma\vdash{}s_0\sim\sigma_0$.
\end{lemma}

\begin{proof}
  By definition of \nameref{sec:state-sim}, 6 subgoals.
  \begin{frameb}
    \begin{enumerate}
    \item
      $\forall{}p\in{}P,id_p\in{}Comps(\Delta),\sigma_p^0\in\Sigma(\Delta(id_p))~s.t.~\gamma(p)=id_p$
      and $\sigma_0(id_p)=\sigma_p^0,$
      $~s_0.M(p)=\sigma_p^0("s\_marking")$.\label{it:marking-eq}
    \item
      $\forall{}t\in{}T_i,id_t\in{}Comps(\Delta),\sigma_t^0\in\Sigma(\Delta(id_t))~s.t.~\gamma(t)=id_t$ and $\sigma_0(id_t)=\sigma_t^0,$\\
      $upper(I_s(t))=\infty\land{}s_0.I(t)\le{}lower(I_s(t))\Rightarrow{}s_0.I(t)=\sigma_t^0("s\_tc")\land{}$\\
      $upper(I_s(t))=\infty\land{}s_0.I(t)>{}lower(I_s(t))\Rightarrow{}\sigma_t^0("s\_tc")=lower(I_s(t))\land{}$\\
      $upper(I_s(t))\neq\infty\land{}s_0.I(t)>{}upper(I_s(t))\Rightarrow{}\sigma_t^0("s\_tc")=upper(I_s(t))\land{}$\\
      $upper(I_s(t))\neq\infty\land{}s_0.I(t)\le{}upper(I_s(t))\Rightarrow{}s_0.I(t)=\sigma_t^0("s\_tc")$.\label{it:tc-eq}
    \item
      $\forall{}t\in{}T_i,id_t\in{}Comps(\Delta),\sigma_t^0\in\Sigma(\Delta(id_t))~s.t.~\gamma(t)=id_t$
      and
      $\sigma_0(id_t)=\sigma_t^0,~$\\
      $s_0.reset_t(t)=\sigma_t^0("s\_reinit\_time\_counter")$.\label{it:reset-eq}
    \item
      $\forall{}c\in\mathcal{C},id_c\in{}Ins(\Delta)~s.t.~\gamma(c)=id_c,~s_0.cond(c)=\sigma_0(id_c)$.\label{it:cond-eq}
    \item
      $\forall{}a\in\mathcal{A},id_a\in{}Outs(\Delta)~s.t.~\gamma(a)=id_a,~s_0.ex(a)=\sigma_0(id_a)$.\label{it:action-eq}
    \item
      $\forall{}f\in\mathcal{F},id_f\in{}Outs(\Delta)~s.t.~\gamma(f)=id_f,~s_0.ex(f)=\sigma_0(id_f)$.\label{it:fun-eq}
    \end{enumerate}
  \end{frameb}

  \begin{itemize}
  \item Apply Lemma~\nameref{lem:init-states-eq-marking} to solve \ref{it:marking-eq}.
  \item Apply Lemma~\nameref{lem:init-states-eq-tc} to solve \ref{it:tc-eq}.
  \item Apply Lemma~\nameref{lem:init-states-eq-rorders} to solve \ref{it:reset-eq}.
  \item Apply Lemma~\nameref{lem:init-states-cond-vals} to solve \ref{it:cond-eq}.
  \item Apply Lemma~\nameref{lem:init-states-act-exec} to solve \ref{it:action-eq}.
  \item Apply Lemma~\nameref{lem:init-states-fun-exec} to solve \ref{it:fun-eq}.
  \end{itemize}
\end{proof}

\subsection{Initial states and marking}
\label{sec:init-states-marking}

\begin{lemma}[Initial States Equal Marking]
  \label{lem:init-states-eq-marking}
  \inithyps{} then
  $\forall{}p\in{}P,id_p\in{}Comps(\Delta),\sigma_p^0\in\Sigma(\Delta(id_p))~s.t.~\gamma(p)=id_p$
  and $\sigma_0(id_p)=\sigma_p^0,$
  $~s_0.M(p)=\sigma_p^0("s\_marking")$.
\end{lemma}

\begin{proof}
  Given a $p\in{}P$, an $id_p\in{}Comps(\Delta)$ and a
  $\sigma_p^0\in\Sigma(\Delta(id_p))~s.t.~\gamma(p)=id_p$ and
  $\sigma_0(id_p)=\sigma_p^0,$ let's show that\\
  \framebox{$s_0.M(p)=\sigma_p^0("s\_marking")$.}\\

  \noindent{}By definition of $id_p$, there exist
  $gm_p,ipm_p,opm_p~s.t.~\mathtt{comp}(id_p,"place",gm_p,ipm_p,opm_p)\in{}d.cs$.\\

  \noindent By property of the \hvhdl{} initialization relation, the P
  design behavior (process ``\texttt{marking}''), and\\
  $\mathtt{comp}(id_p,"place",gm_p,ipm_p,opm_p)\in{}d.cs$, then
  $\sigma_p^0("s\_marking")=\sigma_p^0("initial\_marking")$.\\

  \noindent{}Rewriting $\sigma_p^0("s\_marking")$ as $\sigma_p^0("initial\_marking")$,
  \framebox{$\sigma_p^0("initial\_marking")=s_0.M(p)$.}\\
  
  \noindent By construction,
  ${<}\mathtt{id_p.initial\_marking\Rightarrow}M_0(p){>}\in{}ipm_p$. By
  property of the \hvhdl{} initialization relation, and
  $\mathtt{comp}(id_p,"place",gm_p,ipm_p,opm_p)\in{}d.cs$,
  then $\sigma_p^0("initial\_marking")=M_0(p)$.\\

  \noindent{}By definition of $s_0$, rewriting $s_0.M(p)$ as $M_0(p)$,
  \colorbox{red!20}{$\sigma_p^0("initial\_marking")=s_0.M(p)$.}
  
\end{proof}

\subsection{Initial states and time counters}
\label{sec:init-states-tc}

\begin{lemma}[Initial States Equal Time Counters]
  \label{lem:init-states-eq-tc}
  \inithyps{} then
  $\forall{}t\in{}T_i,id_t\in{}Comps(\Delta),\sigma_t^0\in\Sigma(\Delta(id_t))~s.t.~\gamma(t)=id_t$ and $\sigma_0(id_t)=\sigma_t^0,$\\
  $upper(I_s(t))=\infty\land{}s_0.I(t)\le{}lower(I_s(t))\Rightarrow{}s_0.I(t)=\sigma_t^0("s\_tc")\land{}$\\
  $upper(I_s(t))=\infty\land{}s_0.I(t)>{}lower(I_s(t))\Rightarrow{}\sigma_t^0("s\_tc")=lower(I_s(t))\land{}$\\
  $upper(I_s(t))\neq\infty\land{}s_0.I(t)>{}upper(I_s(t))\Rightarrow{}\sigma_t^0("s\_tc")=upper(I_s(t))\land{}$\\
  $upper(I_s(t))\neq\infty\land{}s_0.I(t)\le{}upper(I_s(t))\Rightarrow{}s_0.I(t)=\sigma_t^0("s\_tc")$.
\end{lemma}

\begin{proof}
  \noindent{}Given a $t\in{}T_i$, an $id_t\in{}Comps(\Delta)$ and a
  $\sigma_t^0\in\Sigma(\Delta(id_t))~s.t.~\gamma(t)=id_t$ and
  $\sigma_0(id_t)=\sigma_t^0$, let's show that:
  \begin{enumerate}
  \item \framebox{$upper(I_s(t))=\infty\land{}s_0.I(t)\le{}lower(I_s(t))\Rightarrow{}s_0.I(t)=\sigma_t^0("s\_tc")$}
  \item \framebox{$upper(I_s(t))=\infty\land{}s_0.I(t)>{}lower(I_s(t))\Rightarrow{}\sigma_t^0("s\_tc")=lower(I_s(t))$}
  \item \framebox{$upper(I_s(t))\neq\infty\land{}s_0.I(t)>{}upper(I_s(t))\Rightarrow{}\sigma_t^0("s\_tc")=upper(I_s(t))$}
  \item \framebox{$upper(I_s(t))\neq\infty\land{}s_0.I(t)\le{}upper(I_s(t))\Rightarrow{}s_0.I(t)=\sigma_t^0("s\_tc")$}
  \end{enumerate}

  \noindent{}By definition of $id_t$, there exist $gm_t,ipm_t,opm_t$
  s.t.
  $\mathtt{comp}(id_t,"transition",gm_t,ipm_t,opm_t)\in{}d.cs$.\\

  \noindent{}Then, let's show the 4 previous subgoals.
  
  \begin{enumerate}
  \item Assume $upper(I_s(t))=\infty\land{}s_0.I(t)\le{}lower(I_s(t))$, then show \framebox{${}s_0.I(t)=\sigma_t^0("s\_tc")$.}\\
    Rewriting $s_0.I(t)$ as $0$, by definition of $s_0$, \framebox{$\sigma_t^0("s\_tc")=0$.}

    \noindent By property of the \hvhdl{} initialization relation, the
    T design behavior (process ``\texttt{time\_counter}''), and
    $\mathtt{comp}(id_t,"transition",gm_t,ipm_t,opm_t)\in{}d.cs$, then
    \colorbox{red!20}{$\sigma_t^0("s\_tc")=0$.}
  \item Assume $upper(I_s(t))=\infty\land{}s_0.I(t)>{}lower(I_s(t))$,
    then show \framebox{$\sigma_t^0("s\_tc")=lower(I_s(t))$}.  By
    definition, $lower(I_s(t))\in\mathbb{N}^{*}$ and
    $s_0.I(t)=0$. Then, \colorbox{red!20}{$lower(I_s(t)){}<0$ is a
      contradiction.}
  \item Assume
    $upper(I_s(t))\neq\infty\land{}s_0.I(t)>{}upper(I_s(t))$, then
    show \framebox{$\sigma_t^0("s\_tc")=upper(I_s(t))$}.  By definition,
    $upper(I_s(t))\in\mathbb{N}^{*}$ and $s_0.I(t)=0$. Then,
    \colorbox{red!20}{$upper(I_s(t)){}<0$ is a contradiction.}
  \item Assume
    $upper(I_s(t))\neq\infty\land{}s_0.I(t)\le{}upper(I_s(t))$, then
    show \framebox{$s_0.I(t)=\sigma_t^0("s\_tc")$}.\\
 
    Rewriting $s_0.I(t)$ as $0$, by definition of $s_0$,
    \framebox{$\sigma_t^0("s\_tc")=0$.}

    \noindent By property of the \hvhdl{} initialization relation, the
    T design behavior (process ``\texttt{time\_counter}''), and
    $\mathtt{comp}(id_t,"transition",gm_t,ipm_t,opm_t)\in{}d.cs$, then
    \colorbox{red!20}{$\sigma_t^0("s\_tc")=0$.}
  \end{enumerate}
\end{proof}

\subsection{Initial states and reset orders}
\label{sec:init-states-rorders}

\begin{lemma}[Initial States Equal Reset Orders]
  \label{lem:init-states-eq-rorders}
  \inithyps{} then
  $\forall{}t\in{}T_i,id_t\in{}Comps(\Delta),\sigma_t^0\in\Sigma(\Delta(id_t))~s.t.~\gamma(t)=id_t$
  and
  $\sigma_0(id_t)=\sigma_t^0,~s_0.reset_t(t)=\sigma_t^0("s\_reinit\_time\_counter")$.
\end{lemma}

\begin{proof}
  Given a $t\in{}T_i$, an $id_t\in{}Comps(\Delta)$ and a
  $\sigma_t^0\in\Sigma(\Delta(id_t))~s.t.~\gamma(t)=id_t$, let's show
  that\\
  \framebox{$s_0.reset_t(t)=\sigma_t^0("s\_reinit\_time\_counter")$}.\\
  
  \noindent{}Rewriting $s_0.reset_t(t)$ as $false$, by definition of
  $s_0$,
  \framebox{$\sigma_t^0("s\_reinit\_time\_counter")=false$.}\\
  
  \noindent{}By definition of $id_t$, there exist
  $gm_t,ipm_t,opm_t~s.t.$
  $~\mathtt{comp}(id_t,"transition",gm_t,ipm_t,opm_t)\in{}d.cs$.\\
  
  \noindent By property of the \hvhdl{} initialization relation, the T
  design behavior (process \texttt{reinit\_time\_counter}
  \texttt{\_evaluation}), and
  $\mathtt{comp}(id_t,"transition",gm_t,ipm_t,opm_t)\in{}d.cs$,\\
  we know
  $\sigma_t^0("s\_reinit\_time\_counter")=\prod\limits_{i=0}^{\Delta(id_t)("in\_arcs\_nb")-1}\sigma_t^0("rt")(i)$,
  where $\Delta(id_t)("in\_arcs\_nb")$ is the value of the generic
  constant $"in\_arcs\_nb"$ stored in the elaborated design
  $\Delta(id_t)$ (which, by property of the \hvhdl{} elaboration
  relation, is an elaborated version of the T design).\\

  \noindent{}Rewriting
  $\sigma_t^0("s\_reinit\_time\_counter")$ as $\prod\limits_{i=0}^{\Delta(id_t)("in\_arcs\_nb")-1}\sigma_t^0("rt")(i)$,\\
  \framebox{$\prod\limits_{i=0}^{\Delta(id_t)("in\_arcs\_nb")-1}\sigma_t^0("rt")(i)=false$.}\\
  
  \noindent{}For all $t\in{}T$ (resp. $p\in{}P$), let $input(t)$
  (resp. $input(p)$) be the set of input places of $t$ (resp. input
  transitions of $p$), and let $output(t)$ (resp. $output(p)$) be the
  set of output places of $t$ (resp. output transitions of $p$).\\

  \noindent{}Case analysis on $input(t)$ (2 CASES).

  \begin{itemize}
  \item \textbf{CASE} $input(t)=\emptyset$.

    By construction,
    ${<}\mathtt{id_t.in\_arcs\_nb\Rightarrow}1{>}\in{}gm_t$, and by
    property of the elaboration relation,\\
    $\Delta(id_t)("in\_arcs\_nb")=1$.  By construction,
    $<\mathtt{id_t.rt(0)\Rightarrow}false>\in{}ipm_t$,
    and by property of the initialization relation, $\sigma_t^0("rt")(0)=false$.\\

    \noindent{}Rewriting $\Delta(id_t)("in\_arcs\_nb")$ as $1$ and
    $\sigma_t^0("rt")(0)$ as $false$,\\
    \colorbox{red!20}{$\prod\limits_{i=0}^{\Delta("in\_arcs\_nb")-1}\sigma_t^0("rt")(i)=\sigma_t^0("rt")(0)=false$.}
    
  \item \textbf{CASE} $input(t)\neq\emptyset$.

    We know
    $\prod\limits_{i=0}^{\Delta(id_t)("in\_arcs\_nb")-1}\sigma_t^0("rt")(i)=false\equiv$
    $\exists{}i\in[0,\Delta(id_t)("in\_arcs\_nb")-1]~s.t.~\sigma_t^0("rt")(i)=false$.\\
    \framebox{$\exists{}i\in[0,\Delta(id_t)("in\_arcs\_nb")-1]~s.t.~\sigma_t^0("rt")(i)=false$.}\\
    \noindent{}Since $input(t)\neq\emptyset,~\exists{}p~s.t.~p\in{}input(t)$. Let's take such a $p\in{}input(t)$.\\
    
    \noindent{}By construction, for all $p\in{}P$, there exist $id_p~s.t.~\gamma(p)=id_p$.

    \noindent{}By definition of $id_p$, there exist
    $gm_p,ipm_p,opm_p~s.t.~\mathtt{comp}(id_p,"place",gm_p,ipm_p,opm_p)\in{}d.cs$.\\

    \noindent{}By construction, for all $p\in{}P$, $t\in{}T$
    s.t. $p\in{}input(t)$ and $t\in{}output(p)$, for all $id_p,id_t$
    s.t. $\gamma(p)=id_p$ and $\gamma(t)=id_t$, for all
    $gm_p,ipm_p,opm_p$
    s.t. $\mathtt{comp}(id_p,"place",gm_p,ipm_p,opm_p)\in{}d.cs$ and
    $gm_t,ipm_t,opm_t$
    s.t. $\mathtt{comp}(id_t,"transition",gm_t,ipm_t,opm_t)\in{}d.cs$,
    there exist $i\in[0,\vert{}input(t)\vert{}-1]$,
    $j\in[0,\vert{}output(p)\vert{}-1]$, $id_{ji}$
    s.t. ${<}\mathtt{id_p.rtt(j)\Rightarrow}id_{ji}{>}\in{}opm_p$ and
    ${<}\mathtt{id_t.rt(i)\Rightarrow}id_{ji}{>}\in{}ipm_t$. Let's take such a $i$, $j$ and $id_{ji}$.\\

    \noindent{}By construction, for all $t\in{}T$
    s.t. $input(t)\neq{}\emptyset$, $id_t,gm_t,ipm_t,opm_t$
    s.t. $\gamma(t)=id_t$ and\\
    $\mathtt{comp}(id_t,"transition",gm_t,ipm_t,opm_t)\in{}d.cs$, then
    ${<}\mathtt{id_t.in\_arcs\_nb\Rightarrow}\vert{}input(t)\vert{}{>}\in{}gm_t$.\\

    \noindent{}By property of the \hvhdl{} elaboration relation and
    ${<}\mathtt{id_t.in\_arcs\_nb\Rightarrow}\vert{}input(t)\vert{}{>}\in{}gm_t$,
    we know $\Delta(id_t)("in\_arcs\_nb)=\vert{}input(t)\vert$.\\

    \noindent{}Rewriting $\Delta(id_t)("in\_arcs\_nb)$ as
    $\vert{}input(t)\vert$, we have
    $i\in[0, \Delta(id_t)("in\_arcs\_nb)-1]$. Let's take that i to
    prove the goal.\\

    \framebox{$\sigma_t^0("rt")(i)=false$.}\\

    \noindent{}By property of the \hvhdl{} initialization relation and
    ${<}\mathtt{id_t.rt(i)\Rightarrow}id_{ji}{>}\in{}ipm_t$, we know
    $\sigma_t^0("rt")(i)=\sigma_0("id_{ji}")$.\\

    \noindent{}Rewriting $\sigma_t^0("rt")(i)$ as $\sigma_0("id_{ji}")$, \framebox{$\sigma_0("id_{ji}")=false$.}\\

    \noindent{}By property of the \hvhdl{} elaboration and initialization
    relations, and
    $\mathtt{comp}(id_p,"place",gm_p,ipm_p,opm_p)\in{}d.cs$, there
    exists a $\sigma_p^0\in{}\Sigma(\Delta(id_p))$
    s.t. $\sigma_0(id_p)=\sigma_p^0$.\\

    \noindent{}By property of the \hvhdl{} initialization relation and
    $<\mathtt{id_p.rtt(j)\Rightarrow}id_{ji}>\in{}opm_p$, we know
    $\sigma_0("id_{ji}")=\sigma_p^0("rtt")(j)$.\\

    \noindent{}Rewriting $\sigma_0("id_{ji}")$ as
    $\sigma_p^0("rtt")(j)$, \framebox{$\sigma_p^0("rtt")(j)=false$.}\\

    \noindent{}By property of the \hvhdl{} initialization relation,
    the P design behavior (process
    $\mathtt{reinit\_transitions\_ti}$-\\$\mathtt{me\_evaluation}$), and
    $\mathtt{comp}(id_p,"place",gm_p,ipm_p,opm_p)\in{}d.cs$, we know
    that\\ for all $j\in{}[0,\Delta(id_p)("out\_arcs\_nb")-1]$,
    $\sigma_p^0("rtt")(j)=false$.\\

    \noindent{}By construction, for all $p\in{}P$
    s.t. $output(p)\neq{}\emptyset$,
    $id_p\in{}Comps(\Delta),gm_p,ipm_p,opm_p$ s.t. $\gamma(p)=id_p$
    and $\mathtt{comp}(id_p,"transition",gm_p,ipm_p,opm_p)\in{}d.cs$,
    then
    ${<}\mathtt{id_p.out\_arcs\_nb\Rightarrow}\vert{}output(p)\vert{}{>}\in{}gm_p$.\\

    \noindent{}By property of the \hvhdl{} elaboration relation and
    ${<}\mathtt{id_p.out\_arcs\_nb\Rightarrow}\vert{}output(p)\vert{}{>}\in{}gm_p$,
    we know $\Delta(id_p)("out\_arcs\_nb")=\vert{}output(p)\vert$.\\

    \noindent{}Rewriting $\vert{}output(p)\vert$ as
    $\Delta(id_p)("out\_arcs\_nb)$, we have
    $j\in[0, \Delta(id_p)("out\_arcs\_nb)-1]$.
    Then, we can deduce
    \colorbox{red!20}{$\sigma_p^0("rtt")(j)=false$}.
  \end{itemize}
  
\end{proof}

\subsection{Initial states and condition values}
\label{sec:init-states-cond-vals}

\begin{lemma}[Initial States Equal Condition Values]
  \label{lem:init-states-cond-vals}
  \inithyps{} 
  then
  $\forall{}c\in\mathcal{C},id_c\in{}Ins(\Delta)~s.t.~\gamma(c)=id_c,~s_0.cond(c)=\sigma_0(id_c)$.
\end{lemma}

\begin{proof}
  Given a $c\in\mathcal{C}$ and an
  $id_c\in{}Ins(\Delta)~s.t.~\gamma(c)=id_c$, let's show that
  \fbox{$s_0.cond(c)=\sigma_0(id_c)$.}\\

  \noindent{}Rewriting $s_0.cond(c)$ as $false$, by definition of
  $s_0$, \fbox{$\sigma_0(id_c)=false$.}

  \noindent{}By construction, $id_c$ is an input port identifier of
  boolean type in the \hvhdl{} design $d$.

  \noindent{}By property, of the \hvhdl{} elaboration relation,
  $\sigma_e(id_c)=false$, where $false$ is the default value
  associated to signals of the boolean type during the elaboration
  (see definition of default value in chapter \hvhdl{} semantics).

  \noindent{}By property of the \hvhdl{} initialization relation, we
  have $\sigma_e(id_c)=\sigma_0(id_c)$ (i.e, input ports are not
  assigned during the initialization phase).

  \noindent{}Rewriting $\sigma_e(id_c)$ as $false$,
  \colorbox{red!20}{$\sigma_0(id_c)=false$.}
  
\end{proof}

\subsection{Initial states and action executions}
\label{sec:init-states-act-exec}

\begin{todobox}
  Correction: $id_f$ is assigned by the reset block of the function process
\end{todobox}

\begin{lemma}[Initial States Equal Action Executions]
  \label{lem:init-states-act-exec}
  \inithyps{} 
  then
  $\forall{}a\in\mathcal{A},id_a\in{}Outs(\Delta)~s.t.~\gamma(a)=id_a,~s_0.ex(a)=\sigma_0(id_a)$.
\end{lemma}

\begin{proof}
  Given a $a\in\mathcal{A}$ and an
  $id_a\in{}Outs(\Delta)~s.t.~\gamma(a)=id_a$, let's show that
  \fbox{$s_0.ex(a)=\sigma_0(id_a)$.}\\

  \noindent{}Rewriting $s_0.ex(a)$ as $false$, by definition of
  $s_0$, \fbox{$\sigma_0(id_a)=false$.}

  \noindent{}By construction, $id_a$ is an output port identifier of
  boolean type in the \hvhdl{} design $d$.

  \noindent{}By property, of the \hvhdl{} elaboration relation,
  $\sigma_e(id_a)=false$, where $false$ is the default value
  associated to signals of the boolean type during the elaboration
  (see definition of default value in chapter \hvhdl{} semantics).

  \noindent{}By construction, we know that the output port identifier
  $id_a$ is assigned in the generated \texttt{action} process, only at
  the falling edge phase of the simulation cycle (i.e, the assignment
  takes place in a \texttt{falling} statement block).
  
  \noindent{}By property of the \hvhdl{} initialization relation, and
  we have $\sigma_e(id_a)=\sigma_0(id_a)$ (i.e, process
  \texttt{action} is idle during the initialization phase).

  \noindent{}Rewriting $\sigma_e(id_a)$ as $false$,
  \colorbox{red!20}{$\sigma_0(id_a)=false$.}
  
\end{proof}

\subsection{Initial states and function executions}
\label{sec:init-states-fun-exec}

\begin{todobox}
  Correction: $id_f$ is assigned by the reset block of the function process
\end{todobox}

\begin{lemma}[Initial States Equal Function Executions]
  \label{lem:init-states-fun-exec}
  \inithyps{} 
  then
  $\forall{}f\in\mathcal{F},id_f\in{}Outs(\Delta)~s.t.~\gamma(f)=id_f,~s_0.ex(f)=\sigma_0(id_f)$.
\end{lemma}

\begin{proof}
  Given a $f\in\mathcal{F}$ and an
  $id_f\in{}Outs(\Delta)~s.t.~\gamma(f)=id_f$, let's show that
  \fbox{$s_0.ex(f)=\sigma_0(id_f)$.}\\

  \noindent{}Rewriting $s_0.ex(f)$ as $false$, by definition of $s_0$,
  \fbox{$\sigma_0(id_f)=false$.}

  \noindent{}By construction, $id_f$ is an output port identifier of
  boolean type in the \hvhdl{} design $d$.

  \noindent{}By property, of the \hvhdl{} elaboration relation,
  $\sigma_e(id_f)=false$, where $false$ is the default value
  associated to signals of the boolean type during the elaboration
  (see definition of default value in chapter \hvhdl{} semantics).

  \noindent{}By construction, we know that the output port identifier
  $id_f$ is assigned in the generated \texttt{function} process (i.e,
  \texttt{function} is the process identifier), only at the rising
  edge phase of the simulation cycle (i.e, the assignment takes place
  in a \texttt{rising} statement block).
  
  \noindent{}By property of the \hvhdl{} initialization relation, and
  we have $\sigma_e(id_f)=\sigma_0(id_f)$ (i.e, process
  \texttt{function} is idle during the initialization phase).

  \noindent{}Rewriting $\sigma_e(id_f)$ as $false$,
  \colorbox{red!20}{$\sigma_0(id_f)=false$.}
  
\end{proof}

%%% Local Variables:
%%% mode: latex
%%% TeX-master: "../../main"
%%% End:


\section{First Rising Edge}
\label{sec:fst-re}
% \begin{definition}[First Rising Edge Hypotheses]
  \label{def:fst-re-hyps}
  Given an
  $sitpn\in{}SITPN,d\in{}design,\gamma\in{}WM(sitpn,d),
  \Delta\in{}ElDesign(d,\mathcal{D}_\mathcal{H}),$
  $\sigma_{e},\sigma_0,\sigma_i,\sigma_{\uparrow},\sigma\in{}\Sigma(\Delta)$,
  $E_c\in{}\mathbb{N}\rightarrow{}\mathcal{C}\rightarrow{}\mathbb{B}$,
  $E_p\in{}(\mathbb{N}\times{}\{\uparrow,\downarrow\})\rightarrow{}Ins(\Delta)\rightarrow{}value$,
  $\tau\in\mathbb{N}$, assume that:
  \begin{itemize}
  \item $\lfloor{}sitpn\rfloor_\mathcal{H}=(d,\gamma)$ and
    $\mathcal{D}_\mathcal{H},\emptyset\vdash{}d\srarrow{elab}{\fontsize{6}{8}\selectfont}(\Delta,\sigma_{e})$
    and $\gamma\vdash{}E_p\stackrel{env}{=}E_c$
  \item $\sigma_0$ is the initial state of $\Delta$: 
    $\Delta,\sigma_{e}\vdash{}d.cs\srarrow{init}{\fontsize{6}{8}\selectfont}\sigma_0$
  \item $E_c,\tau\vdash{}s_0\srarrow{\uparrow_0}{\fontsize{6}{8}\selectfont}s_0$
  \item $\mathtt{Inject}_\uparrow(\sigma_0, E_p, \tau, \sigma_i)$
    and
    $\Delta,\sigma_i\vdash\mathrm{d.cs}\xrightarrow{\uparrow}\sigma_{\uparrow}$
    and
    $\Delta,\sigma_{\uparrow}\vdash\mathrm{d.cs}\xrightarrow{\theta}\sigma$
  \end{itemize}
  
\end{definition}

\def\fstrehyps{For all $sitpn,d,\gamma,\Delta,$
  $\sigma_{e},\sigma_0,\sigma_i,\sigma_{\uparrow},\sigma$, $E_c$, $E_p$,
  $\tau$ that verify the hypotheses of Def.~\ref{def:fst-re-hyps},}

%%%%%%%%%%%%%%%%%%%%%%%%%%%%%%%%%%%%%%%%%%%%%
%%%%%%%%%% FIRST RISING EDGE LEMMA %%%%%%%%%%
%%%%%%%%%%%%%%%%%%%%%%%%%%%%%%%%%%%%%%%%%%%%%

\begin{lemma}[First Rising Edge]
  \label{lem:fst-re}
  \fstrehyps{} then
  $\gamma,E_c,\tau\vdash{}s_0\stackrel{\uparrow}{\sim}{}\sigma$.
\end{lemma}

\begin{proof}
  By definition of \nameref{def:post-re-state-sim}, 6 subgoals.
  \begin{frameb}
    \begin{enumerate}
    \item
      $\forall{}p\in{}P,id_p\in{}Comps(\Delta),\sigma_p\in\Sigma(\Delta(id_p))~s.t.~\gamma(p)=id_p$
      and $\sigma(id_p)=\sigma_p,$
      $~s_0.M(p)=\sigma_p("s\_marking")$.\label{it:fst-re-marking-eq}
    \item
      $\forall{}t\in{}T_i,id_t\in{}Comps(\Delta),\sigma_t\in\Sigma(\Delta(id_t))~s.t.~\gamma(t)=id_t$ and $\sigma(id_t)=\sigma_t,$\\
      $upper(I_s(t))=\infty\land{}s_0.I(t)\le{}lower(I_s(t))\Rightarrow{}s_0.I(t)=\sigma_t("s\_tc")\land{}$\\
      $upper(I_s(t))=\infty\land{}s_0.I(t)>{}lower(I_s(t))\Rightarrow{}\sigma_t("s\_tc")=lower(I_s(t))\land{}$\\
      $upper(I_s(t))\neq\infty\land{}s_0.I(t)>{}upper(I_s(t))\Rightarrow{}\sigma_t("s\_tc")=upper(I_s(t))\land{}$\\
      $upper(I_s(t))\neq\infty\land{}s_0.I(t)\le{}upper(I_s(t))\Rightarrow{}s_0.I(t)=\sigma_t("s\_tc")$.\label{it:fst-re-tc-eq}
    \item
      $\forall{}t\in{}T_i,id_t\in{}Comps(\Delta),\sigma_t\in\Sigma(\Delta(id_t))~s.t.~\gamma(t)=id_t$
      and
      $\sigma(id_t)=\sigma_t,~$\\
      $s_0.reset_t(t)=\sigma_t("s\_reinit\_time\_counter")$.\label{it:fst-re-reset-eq}
    \item
      $\forall{}a\in\mathcal{A},id_a\in{}Outs(\Delta)~s.t.~\gamma(a)=id_a,~s_0.ex(a)=\sigma(id_a)$.\label{it:fst-re-action-eq}
    \item
      $\forall{}f\in\mathcal{F},id_f\in{}Outs(\Delta)~s.t.~\gamma(f)=id_f,~s_0.ex(f)=\sigma(id_f)$.\label{it:fst-re-fun-eq}
    \item
      $\forall{}t\in{}T_i,id_t\in{}Comps(\Delta),\sigma_t\in\Sigma(\Delta(id_t))~s.t.~\gamma(t)=id_t$,
      $t\in{}Sens(s.M)\Leftrightarrow\sigma_t("s\_enabled")=\mathtt{true}$.\label{it:fst-re-sens-eq}
    \end{enumerate}
  \end{frameb}

  \begin{itemize}[label=--]
  \item Apply Lemma~\nameref{lem:fst-re-equal-marking} to solve \ref{it:fst-re-marking-eq}.
  \item Apply ``First Rising Edge Equal Time Counters'' lemma to solve \ref{it:fst-re-tc-eq}.
  \item Apply ``First Rising Edge Equal Reset Orders'' lemma to solve \ref{it:fst-re-reset-eq}.
  \item Apply ``First Rising Edge Equal Action Executions'' lemma to solve \ref{it:fst-re-action-eq}.
  \item Apply ``First Rising Edge Equal Function Executions '' lemma to solve \ref{it:fst-re-fun-eq}.
  \item Apply ``First Rising Edge Equal Sensitized'' lemma to solve \ref{it:fst-re-sens-eq}.
  \end{itemize}
  
\end{proof}

\subsection{First rising edge and marking}
\label{sec:fst-re-marking}

\begin{lemma}[First Rising Edge Equal Marking]
  \label{lem:fst-re-equal-marking}
  \fstrehyps{} then
  $\forall{}p\in{}P,id_p\in{}Comps(\Delta),\sigma_p\in\Sigma(\Delta(id_p))~s.t.~\gamma(p)=id_p$
  and $\sigma(id_p)=\sigma_p,$ $~s_0.M(p)=\sigma_p("s\_marking")$.
\end{lemma}

\begin{proof}
  Given a $p$, $id_p$, $\sigma_p$ s.t. $\gamma(p)=id_p$ and
  $\sigma(id_p)=\sigma_p$, let us show that
  \fbox{$~s_0.M(p)=\sigma_p("s\_marking")$.}
  
  \noindent{}By definition of $id_p$, there exist
  $gm_p,ipm_p,opm_p~s.t.~\mathtt{comp}(id_p,"place",gm_p,ipm_p,opm_p)\in{}d.cs$.\\
  
  \noindent By property of the \hvhdl{} elaboration relation, the
  \hvhdl{} initialization relation, the $\mathtt{Inject}_\uparrow$
  relation, the \hvhdl{} rising edge relation and
  $\mathtt{comp}(id_p,"place",gm_p,ipm_p,opm_p)\in{}d.cs$, there exist
  a
  $\sigma_p^{e},\sigma_p^{0},\sigma_p^{injr},\sigma_p^{r}\in{}\Sigma(\Delta)$
  s.t.  $\sigma_{e}(id_p)=\sigma_p^{e}$ and
  $\sigma_{0}(id_p)=\sigma_p^{0}$ and
  $\sigma_i(id_p)=\sigma_p^{injr}$ and
  $\sigma_{r}(id_p)=\sigma_p^{r}$ .

  \begin{pcomm}
    From the elaboration to the end of the first rising edge phase, an
    internal state is associated with the P component instance
    $id_p$ in the component store of the top-level design $d$.
  \end{pcomm}
  
  \noindent{} By property of the \hvhdl{} rising edge relation, the P design behavior (process ``\texttt{marking}''), and\\
  $\mathtt{comp}(id_p,"place",gm_p,ipm_p,opm_p)\in{}d.cs$, then\\
  $\sigma_p^{r}("s\_marking")=\sigma_p^{injr}("s\_marking")+\sigma_p^{injr}("s\_input\_token\_sum")-\sigma_p^{injr}("s\_output\_token\_sum")$.

  \begin{pcomm}
    Result of the execution of the process ``\texttt{marking}'' that performs the signal assignment\\
    $\mathtt{s\_marking\Leftarrow{}s\_marking+s\_input\_token\_sum-s\_output\_token\_sum}$.
  \end{pcomm}

  \noindent{} By property of the \hvhdl{} stabilize relation, the P design behavior (process ``\texttt{marking}''), and\\
  $\mathtt{comp}(id_p,"place",gm_p,ipm_p,opm_p)\in{}d.cs$, then
  $\sigma_p^{r}("s\_marking")=\sigma_p("s\_marking")$.

  \begin{pcomm}
    As it is only assigned by the process ``\texttt{marking}'', and as
    the process ``\texttt{marking}'' is never executed during the
    stabilization phase, the ``\texttt{s\_marking}'' signal has an
    invariant value during the stabilization phase.
  \end{pcomm}

  \noindent{} Rewriting $\sigma_p("s\_marking")$ as
  $\sigma_p^{r}("s\_marking")$,
  and $\sigma_p^{r}("s\_marking")$ as\\
  $\sigma_p^{injr}("s\_marking")+\sigma_p^{injr}("s\_input\_token\_sum")-\sigma_p^{injr}("s\_output\_token\_sum")$,\\
  \fbox{$s_0.M(p)=\sigma_p^{injr}("s\_marking")+\sigma_p^{injr}("s\_input\_token\_sum")-\sigma_p^{injr}("s\_output\_token\_sum")$.}\\

  \noindent{} By property of the $\mathtt{Inject}_\uparrow$ relation,
  $\sigma_p^{injr}("s\_marking")=\sigma_p^{0}("s\_marking")$ and\\
  $\sigma_p^{injr}("s\_input\_token\_sum")=\sigma_p^{0}("s\_input\_token\_sum")$
  and\\
  $\sigma_p^{injr}("s\_output\_token\_sum")=\sigma_p^{0}("s\_output\_token\_sum")$.
  Rewriting the above,\\
  \fbox{$s_0.M(p)=\sigma_p^{0}("s\_marking")+\sigma_p^{0}("s\_input\_token\_sum")-\sigma_p^{0}("s\_output\_token\_sum")$.}\\

  \begin{todobox}
    Detail the two lemmas giving this property.
  \end{todobox}
  \noindent{}By property of the \hvhdl{} initialization relation,
  $\sigma_p^{0}("s\_input\_token\_sum")=0$ and\\
  $\sigma_p^{0}("s\_output\_token\_sum")=0$. Rewriting the above,
  \fbox{$s_0.M(p)=\sigma_p^{0}("s\_marking")$.}\\

  \noindent{} Applying the \nameref{lem:init-states-eq-marking} lemma,
  \qedbox{$s_0.M(p)=\sigma_p^{0}("s\_marking")$.}
\end{proof}

\subsection{First rising edge and time counters}
\label{sec:fst-re-tc}

\begin{lemma}[First Rising Edge Equal Time Counters]
  \label{lem:fst-re-equal-tc}
  \fstrehyps{} then\\
  $\forall{}t\in{}T_i,id_t\in{}Comps(\Delta),\sigma_t\in\Sigma(\Delta(id_t))~s.t.~\gamma(t)=id_t$ and $\sigma(id_t)=\sigma_t,$\\
  $upper(I_s(t))=\infty\land{}s_0.I(t)\le{}lower(I_s(t))\Rightarrow{}s_0.I(t)=\sigma_t("s\_tc")\land{}$\\
  $upper(I_s(t))=\infty\land{}s_0.I(t)>{}lower(I_s(t))\Rightarrow{}\sigma_t("s\_tc")=lower(I_s(t))\land{}$\\
  $upper(I_s(t))\neq\infty\land{}s_0.I(t)>{}upper(I_s(t))\Rightarrow{}\sigma_t("s\_tc")=upper(I_s(t))\land{}$\\
  $upper(I_s(t))\neq\infty\land{}s_0.I(t)\le{}upper(I_s(t))\Rightarrow{}s_0.I(t)=\sigma_t("s\_tc")$.
\end{lemma}

\begin{proof}
  \noindent{}Given a $t\in{}T_i$, an $id_t\in{}Comps(\Delta)$ and a
  $\sigma_t\in\Sigma(\Delta(id_t))~s.t.~\gamma(t)=id_t$ and
  $\sigma(id_t)=\sigma_t$, let's show that:
  \begin{enumerate}
  \item \framebox{$upper(I_s(t))=\infty\land{}s_0.I(t)\le{}lower(I_s(t))\Rightarrow{}s_0.I(t)=\sigma_t("s\_tc")$}
  \item \framebox{$upper(I_s(t))=\infty\land{}s_0.I(t)>{}lower(I_s(t))\Rightarrow{}\sigma_t("s\_tc")=lower(I_s(t))$}
  \item \framebox{$upper(I_s(t))\neq\infty\land{}s_0.I(t)>{}upper(I_s(t))\Rightarrow{}\sigma_t("s\_tc")=upper(I_s(t))$}
  \item \framebox{$upper(I_s(t))\neq\infty\land{}s_0.I(t)\le{}upper(I_s(t))\Rightarrow{}s_0.I(t)=\sigma_t("s\_tc")$}
  \end{enumerate}

  \noindent{}By definition of $id_t$, there exist $gm_t,ipm_t,opm_t$
  s.t.
  $\mathtt{comp}(id_t,"transition",gm_t,ipm_t,opm_t)\in{}d.cs$.\\

  \noindent By property of the \hvhdl{} elaboration relation, the
  \hvhdl{} initialization relation, the $\mathtt{Inject}_\uparrow$
  relation, the \hvhdl{} rising edge relation and
  $\mathtt{comp}(id_t,"transition",gm_t,ipm_t,opm_t)\in{}d.cs$, there
  exist a
  $\sigma_t^{e},\sigma_t^{0},\sigma_t^{injr},\sigma_t^{r}\in{}\Sigma(\Delta)$
  s.t.  $\sigma_{e}(id_t)=\sigma_t^{e}$ and
  $\sigma_{0}(id_t)=\sigma_t^{0}$ and
  $\sigma_i(id_t)=\sigma_t^{injr}$ and
  $\sigma_{r}(id_t)=\sigma_t^{r}$ .

  \begin{pcomm}
    From the elaboration to the end of the first rising edge phase, an
    internal state is associated with the T component instance $id_t$
    in the component store of the top-level design $d$.
  \end{pcomm}
  
  \noindent{}Then, let's show the 4 previous subgoals.
  
  \begin{enumerate}
  \item Assume $upper(I_s(t))=\infty\land{}s_0.I(t)\le{}lower(I_s(t))$, then show \framebox{${}s_0.I(t)=\sigma_t("s\_tc")$.}\\
    \noindent{} By property of the $\mathtt{Inject_\uparrow}$
    relation, the \hvhdl{} rising edge and stabilize relations, and\\
    $\mathtt{comp}(id_t,"transition",gm_t,ipm_t,opm_t)\in{}d.cs$,
    $\sigma_t("s\_tc")=\sigma_t^0("s\_tc")$.
    \begin{pcomm}
      The above equality is deduced from the two following facts:
      
      \begin{itemize}
      \item The process ``\texttt{time\_counter}'' is the only process
        that assigns signal \texttt{s\_tc} in the T component
        behavior, and it is never executed during the rising edge and
        stabilization phases.
        
      \item The values of component instances' internal signals are
        invariant through the $\mathtt{Inject_\uparrow}$ relation.
      \end{itemize}
    \end{pcomm}

  \noindent{} Rewriting $\sigma_t("s\_tc")$ as $\sigma_t^0("s\_tc")$,
  \fbox{${}s_0.I(t)=\sigma_t^0("s\_tc")$.}\\

  Applying the \nameref{lem:init-states-eq-tc} lemma,
  \qedbox{${}s_0.I(t)=\sigma_t^0("s\_tc")$.}
  
  \item Assume $upper(I_s(t))=\infty\land{}s_0.I(t)>{}lower(I_s(t))$,
    then show \framebox{$\sigma_t("s\_tc")=lower(I_s(t))$}.  By
    definition, $lower(I_s(t))\in\mathbb{N}^{*}$ and
    $s_0.I(t)=0$. Then, \colorbox{red!20}{$lower(I_s(t)){}<0$ is a
      contradiction.}
  \item Assume
    $upper(I_s(t))\neq\infty\land{}s_0.I(t)>{}upper(I_s(t))$, then
    show \framebox{$\sigma_t("s\_tc")=upper(I_s(t))$}.  By definition,
    $upper(I_s(t))\in\mathbb{N}^{*}$ and $s_0.I(t)=0$. Then,
    \colorbox{red!20}{$upper(I_s(t)){}<0$ is a contradiction.}
  \item Assume
    $upper(I_s(t))\neq\infty\land{}s_0.I(t)\le{}upper(I_s(t))$, then
    show \framebox{$s_0.I(t)=\sigma_t("s\_tc")$}.\\

    \noindent{} By property of the $\mathtt{Inject_\uparrow}$
    relation, the \hvhdl{} rising edge and stabilize relations, and\\
    $\mathtt{comp}(id_t,"transition",gm_t,ipm_t,opm_t)\in{}d.cs$,
    $\sigma_t("s\_tc")=\sigma_t^0("s\_tc")$.\\

    \noindent{} Rewriting $\sigma_t("s\_tc")$ as $\sigma_t^0("s\_tc")$,
    \fbox{${}s_0.I(t)=\sigma_t^0("s\_tc")$.}\\

    Applying the \nameref{lem:init-states-eq-tc} lemma,
    \qedbox{${}s_0.I(t)=\sigma_t^0("s\_tc")$.}
  \end{enumerate}
\end{proof}

\subsection{First rising edge and condition combination}
\label{sec:fst-re-cond-comb}

\begin{lemma}[First Rising Edge Equal Condition Combination]
  \label{lem:fst-re-equal-cond-comb}
  \fstrehyps{} then\\
  $\forall{}t\in{}T,id_t\in{}Comps(\Delta)~s.t.~\gamma(t)=id_t,$\\
  $\sigma(id_t)("s\_condition\_combination")=
  \prod\limits_{c\in{}conds(t)}
  \begin{cases}
    E_c(\tau,c) & if~\mathbb{C}(t,c)=1 \\
    \mathtt{not}(E_c(\tau,c)) & if~\mathbb{C}(t,c)=-1 \\
  \end{cases}$\\
  where
  $conds(t)=\{c\in\mathcal{C}~\vert~\mathbb{C}(t,c)=1\lor\mathbb{C}(t,c)=-1\}$.
\end{lemma}

\begin{table}[h]
  \begin{tabular}{|c|c|}
    \hline
    \textbf{Full signal name} & \textbf{Alias} \\
    \hline
    $"s\_condition\_combination"$ & $"scc"$ \\
    \hline
    $"conditions\_number"$ & $"cn"$ \\
    \hline
    $"input\_conditions"$ & $"ic"$ \\
    \hline
  \end{tabular}
\end{table}

\begin{proof}
  Given a $t$, $id_t$, $\sigma_t$ s.t. $\gamma(t)=id_t$, let us show that\\
  \fbox{$\sigma(id_t)("s\_condition\_combination")=
    \prod\limits_{c\in{}conds(t)}
    \begin{cases}
      E_c(\tau,c) & if~\mathbb{C}(t,c)=1 \\
      \mathtt{not}(E_c(\tau,c)) & if~\mathbb{C}(t,c)=-1 \\
    \end{cases}$.}\\

  \noindent{}By definition of $id_t$, there exist $gm_t,ipm_t,opm_t$
  s.t. $\mathtt{comp}(id_t,"transition",gm_t,ipm_t,opm_t)\in{}d.cs$.\\

  \noindent By property of the \hvhdl{} stabilize relation, and $\mathtt{comp}(id_t,"transition",gm_t,ipm_t,opm_t)\in{}d.cs$,\\
  $\sigma(id_t)("scc")=\prod\limits_{i=0}^{\Delta(id_t)("conditions\_number")-1}\sigma(id_t)("input\_conditions")[i]$.

  \noindent{}Rewriting $\sigma(id_t)("scc")$ as
  $\prod\limits_{i=0}^{\Delta(id_t)("cn")-1}\sigma(id_t)("ic")[i]$,\\
  \fbox{$\prod\limits_{i=0}^{\Delta(id_t)("cn")-1}\sigma(id_t)("ic")[i]=
    \prod\limits_{c\in{}conds(t)}
    \begin{cases}
      E_c(\tau,c) & if~\mathbb{C}(t,c)=1 \\
      \mathtt{not}(E_c(\tau,c)) & if~\mathbb{C}(t,c)=-1 \\
    \end{cases}$.}\\

  \noindent{}Case analysis on $conds(t)$ (2 CASES):

  \begin{itemize}
  \item \textbf{CASE} $conds(t)=\emptyset$:\\
    \fbox{$\prod\limits_{i=0}^{\Delta(id_t)("cn")-1}\sigma(id_t)("ic")[i]=\mathtt{true}$.}\\
    
    \noindent{}By construction,
    ${<}\mathtt{conditions\_number\Rightarrow{}1}{>}\in{}gm_t$ and
    ${<}\mathtt{input\_conditions(0)\Rightarrow{}true}{>}\in{}ipm_t$.\\

    \noindent{}By property of the stabilize relation and
    ${<}\mathtt{conditions\_number\Rightarrow{}1}{>}\in{}gm_t$ and
    ${<}\mathtt{input\_conditions(0)\Rightarrow{}true}{>}\in{}ipm_t$,
    then $\Delta(id_t)("cn")=1$ and
    $\sigma(id_t)("ic")[0]=\mathtt{true}$.

    \noindent{}Rewriting $\Delta(id_t)("cn")$ as $1$,
    \qedbox{$\sigma(id_t)("ic")[0]=\mathtt{true}$.}
  \item \textbf{CASE} $conds(t)\neq\emptyset$:\\
    \noindent{}By construction,
    ${<}\mathtt{conditions\_number\Rightarrow{}\vert{}conds(t)\vert}{>}\in{}gm_t$,
    and by property of the stabilize relation, then
    $\Delta(id_t)("cn")=\vert{}conds(t)\vert$.
    
    Then, 2 subgoals to prove the equation:

    \begin{enumerate}
    \item\label{it:fst-re-eq-cond-comb}
      \fbox{\parbox{\linewidth}{$\forall{}c\in{}conds(t),~\exists{}i\in{}[0,\Delta(id_t)("cn")-1]~s.t.~\mathbb{C}(t,c)=1\Rightarrow{}\sigma(id_t)("ic")[i]=E_c(\tau,c)\land{}\mathbb{C}(t,c)=-1\Rightarrow{}\sigma(id_t)("ic")[i]=\mathtt{not}~E_c(\tau,c)$.}}
      Given a $c\in{}conds(t)$, let us show that\\
      \fbox{\parbox{\linewidth}{$\exists{}i\in{}[0,\Delta(id_t)("cn")-1]~s.t.~\mathbb{C}(t,c)=1\Rightarrow{}\sigma(id_t)("ic")[i]=E_c(\tau,c)\land{}\mathbb{C}(t,c)=-1\Rightarrow{}\sigma(id_t)("ic")[i]=\mathtt{not}~E_c(\tau,c)$.}}
      By definition of $c\in{}conds(t)$, there are 2 cases:
      \begin{itemize}
      \item \textbf{CASE} $\mathbb{C}(t,c)=1$:\\
        By construction, there exists $id_c\in{}Ins(\Delta)$
        s.t. $\gamma(c)=id_c$, and there
        exists $i\in{}[0,\vert{}conds(t)\vert-1]$ s.t.
        ${<}\mathtt{input\_conditions(i)\Rightarrow{}id_c}{>}\in{}ipm_t$.

        \noindent{}As $\Delta(id_t)("cn")=\vert{}conds(t)\vert$, then
        we have $i\in[0,\Delta(id_t)("cn")-1]$. Let us take this $i$
        to prove the goal,\\
        \fbox{\parbox{\linewidth}{$\mathbb{C}(t,c)=1\Rightarrow{}\sigma(id_t)("ic")[i]=E_c(\tau,c)\land{}\mathbb{C}(t,c)=-1\Rightarrow{}\sigma(id_t)("ic")[i]=\mathtt{not}~E_c(\tau,c)$.}}\\
        
        The right part of the goal is proved by contradiction, then
        what is left to prove is:
        \fbox{$\mathbb{C}(t,c)=1\Rightarrow{}\sigma(id_t)("ic")[i]=E_c(\tau,c)$.}\\

        Assuming $\mathbb{C}(t,c)=1$, let us show
        \fbox{$\sigma(id_t)("ic")[i]=E_c(\tau,c)$.}\\

        \noindent{}By property of the stabilize relation and
        ${<}\mathtt{input\_conditions(i)\Rightarrow{}id_c}{>}\in{}ipm_t$,
        then $\sigma(id_t)("ic")[i]=\sigma(id_c)$.

        \noindent{}By property of the \hvhdl{}
        $\mathtt{Inject_{\uparrow}}$, the rising edge, the stabilize
        relations, and $id_c\in{}Ins(\Delta)$, then
        $\sigma(id_c)=E_p(\tau,\uparrow)(id_c)$.

        \noindent{}By property of
        $\gamma\vdash{}E_p\stackrel{env}{=}E_c$, then
        $E_p(\tau,\uparrow)(id_c)=E_c(\tau,c)$.

        \noindent{}Rewriting the goal with the above equations,
        \qedbox{$\sigma(id_t)("ic")[i]=E_c(\tau,c)$.}

      \item \textbf{CASE} $\mathbb{C}(t,c)=-1$:\\
        By construction, there exists $id_c\in{}Ins(\Delta)$
        s.t. $\gamma(c)=id_c$, and there exists
        $i\in{}[0,\vert{}conds(t)\vert-1]$ s.t.
        ${<}\mathtt{input\_conditions(i)\Rightarrow{}not~id_c}{>}\in{}ipm_t$.

        \noindent{}As $\Delta(id_t)("cn")=\vert{}conds(t)\vert$, then
        we have $i\in[0,\Delta(id_t)("cn")-1]$. Let us take this $i$
        to prove the goal,\\
        \fbox{\parbox{\linewidth}{$\mathbb{C}(t,c)=1\Rightarrow{}\sigma(id_t)("ic")[i]=E_c(\tau,c)\land{}\mathbb{C}(t,c)=-1\Rightarrow{}\sigma(id_t)("ic")[i]=\mathtt{not}~E_c(\tau,c)$.}}\\
        
        The left part of the goal is proved by contradiction, then
        what is left to prove is:
        \fbox{$\mathbb{C}(t,c)=-1\Rightarrow{}\sigma(id_t)("ic")[i]=\mathtt{not}~E_c(\tau,c)$.}\\

        Assuming $\mathbb{C}(t,c)=-1$, let us show
        \fbox{$\sigma(id_t)("ic")[i]=\mathtt{not}~E_c(\tau,c)$.}\\

        \noindent{}By property of the stabilize relation and
        ${<}\mathtt{input\_conditions(i)\Rightarrow{}not~id_c}{>}\in{}ipm_t$,
        then $\sigma(id_t)("ic")[i]=\mathtt{not}~\sigma(id_c)$.

        \noindent{}By property of the \hvhdl{}
        $\mathtt{Inject_{\uparrow}}$, the rising edge, the stabilize
        relations, and $id_c\in{}Ins(\Delta)$, then
        $\sigma(id_c)=E_p(\tau,\uparrow)(id_c)$.

        \noindent{}By property of
        $\gamma\vdash{}E_p\stackrel{env}{=}E_c$, then
        $E_p(\tau,\uparrow)(id_c)=E_c(\tau,c)$.

        \noindent{}Rewriting the goal with the above equations,
        \qedbox{$\sigma(id_t)("ic")[i]=\mathtt{not}~E_c(\tau,c)$.}
      \end{itemize}
      
    \item
      \fbox{\parbox{\linewidth}{$\forall{}i\in{}[0,\Delta(id_t)("cn")-1],~\exists{}c\in{}conds(t),~s.t.~\mathbb{C}(t,c)=1\Rightarrow{}\sigma(id_t)("ic")[i]=E_c(\tau,c)\land{}\mathbb{C}(t,c)=-1\Rightarrow{}\sigma(id_t)("ic")[i]=\mathtt{not}~E_c(\tau,c)$.}}\\

      Given a $i\in{}[0,\Delta(id_t)("cn")-1]$, let us show\\
      \fbox{\parbox{\linewidth}{$\exists{}c\in{}conds(t),~s.t.~\mathbb{C}(t,c)=1\Rightarrow{}\sigma(id_t)("ic")[i]=E_c(\tau,c)\land{}\mathbb{C}(t,c)=-1\Rightarrow{}\sigma(id_t)("ic")[i]=\mathtt{not}~E_c(\tau,c)$.}}\\

      By construction, there exists $c\in{}conds(t)$ and
      $id_c\in{}Ins(\Delta)$ s.t.  $\gamma(c)=id_c$, and
      $\mathbb{C}(t,c)=1\Rightarrow{<}\mathtt{input\_conditions(i)\Rightarrow{}id_c}{>}\in{}ipm_t$
      and
      $\mathbb{C}(t,c)=-1\Rightarrow{<}\mathtt{input\_conditions(i)\Rightarrow{}not~id_c}{>}\in{}ipm_t$.

      Let us take such an $c\in{}conds(t)$ to prove the goal. By
      definition of $c\in{}conds(t)$, there are 2 cases: see
      \ref{it:fst-re-eq-cond-comb} for the remainder of the proof.
    \end{enumerate}
    
  \end{itemize}
  

  
\end{proof}


%%% Local Variables:
%%% mode: latex
%%% TeX-master: "../../main"
%%% End:


\section{Rising Edge}
\label{sec:rising-edge}
%%%%%%%%%%%%%%%%%%%%%%%%%%%%%%%%%%%%%%%%%%%%
%%%%%%%%%% RISING EDGE HYPOTHESES %%%%%%%%%%
%%%%%%%%%%%%%%%%%%%%%%%%%%%%%%%%%%%%%%%%%%%%

\begin{definition}[Rising Edge Hypotheses]
  \label{def:re-hyps}
  Given an $sitpn\in{}SITPN$, $d\in{}design$,
  $\gamma\in{}WM(sitpn,d)$,
  $E_c\in\mathbb{N}\rightarrow\mathcal{C}\rightarrow\mathbb{B}$,
  $\Delta\in{}ElDesign(d,\mathcal{D}_\mathcal{H})$,
  $E_p\in(\mathbb{N}\times\{\uparrow,\downarrow\})\rightarrow{}Ins(\Delta)\rightarrow{}value$,
  $\tau\in\mathbb{N}$, $s,s'\in{}S(sitpn)$,
  $\sigma_e,\sigma,\sigma_i,\sigma_\uparrow,\sigma'\in\Sigma(\Delta)$,
  $\theta\in\mathtt{list}(\Sigma(\Delta))$, assume that:
  \begin{itemize}
  \item $\lfloor{}sitpn\rfloor_\mathcal{H}=(d,\gamma)$ and
    $\gamma\vdash{}E_p\stackrel{env}{=}E_c$ and
    $\mathcal{D}_\mathcal{H},\emptyset\vdash\mathrm{d}\srarrow{elab}{\fontsize{7}{7}\selectfont}\Delta,\sigma_e$
  \item $\gamma\vdash{}s\stackrel{\downarrow}{\sim}\sigma$ 
  \item
    $E_c,\tau\vdash{}s\srarrow{\uparrow}{\fontsize{7}{7}\selectfont}s'$
  \item $\mathtt{Inject}_\uparrow(\sigma, E_p, \tau, \sigma_i)$ and
    $\Delta,\sigma_i\vdash\mathrm{d.cs}\xrightarrow{\uparrow}\sigma_\uparrow$
    and
    $\Delta,\sigma_\uparrow\vdash\mathrm{d.cs}\xrightarrow{\theta}\sigma'$
  \end{itemize}
\end{definition}

\def\rehyps{For all $sitpn$, $d$, $\gamma$, $E_c$, $E_p$, $\tau$,
  $\Delta$, $\sigma_e$, $s$, $s'$, $\sigma$, $\sigma_i$,
  $\sigma_\uparrow$, $\sigma'$, $\theta$ that verify the hypotheses of
  Def.~\ref{def:re-hyps},}

%%%%%%%%%%%%%%%%%%%%%%%%%%%%%%%%%%%%%%%
%%%%%%%%%% RISING EDGE LEMMA %%%%%%%%%%
%%%%%%%%%%%%%%%%%%%%%%%%%%%%%%%%%%%%%%%

\begin{lemma}[Rising Edge]
  \label{lem:re}
  \rehyps{} then
  $\gamma,E_c,\tau\vdash{}s'\stackrel{\uparrow}{\sim}{}\sigma'$.
\end{lemma}

\begin{proof}
  By definition of \nameref{def:post-re-state-sim}, there are 7 points
  to prove.
  \begin{frameb}
    \begin{enumerate}
    \item
      $\forall{}p\in{}P,id_p\in{}Comps(\Delta)~s'.t.~\gamma(p)=id_p,$
      $~s'.M(p)=\sigma'(id_p)("s\_marking")$.\label{it:marking-eq-re}
    \item
      $\forall{}t\in{}T_i,id_t\in{}Comps(\Delta)~s.t.~\gamma(t)=id_t$,\\
      $\big(upper(I_s(t))=\infty\land{}s'.I(t)\le{}lower(I_s(t))\Rightarrow$
      $s'.I(t)=\sigma'(id_t)("s\_time\_counter")\big)$\\
      $\land\big(upper(I_s(t))=\infty\land{}s'.I(t)>{}lower(I_s(t))\Rightarrow$
      $\sigma'(id_t)("s\_time\_counter")=lower(I_s(t))\big)$\\
      $\land\big(upper(I_s(t))\neq\infty\land{}s'.I(t)>{}upper(I_s(t))\Rightarrow$
      $\sigma'(id_t)("s\_time\_counter")=upper(I_s(t))\big)$\\
      $\land\big(upper(I_s(t))\neq\infty\land{}s'.I(t)\le{}upper(I_s(t))\Rightarrow$
      $s'.I(t)=\sigma'(id_t)("s\_time\_counter")\big)$.\label{it:time-count-eq-re}
    \item
      $\forall{}t\in{}T_i,id_t\in{}Comps(\Delta)~s.t.~\gamma(t)=id_t,$
      $s'.reset_t(t)=\sigma'(id_t)("s\_reinit\_time\_counter")$.\label{it:reset-eq-re}
    \item
      $\forall{}a\in\mathcal{A},id_a\in{}Outs(\Delta)~s.t.~\gamma(a)=id_a,~s'.ex(a)=\sigma'(id_a)$.\label{it:action-eq-re}
    \item
      $\forall{}f\in\mathcal{F},id_f\in{}Outs(\Delta)~s.t.~\gamma(f)=id_f,~s'.ex(f)=\sigma'(id_f)$.\label{it:fun-eq-re}
    \item
      $\forall{}t\in{}T,id_t\in{}Comps(\Delta)~s.t.~\gamma(t)=id_t,$
      $t\in{}Sens(s'.M)\Leftrightarrow\sigma'(id_t)("s\_enabled")=\mathtt{true}$.\label{it:sens-eq-re}
    \item
      $\forall{}t\in{}T,id_t\in{}Comps(\Delta)~s.t.~\gamma(t)=id_t,$
      $t\notin{}Sens(s'.M)\Leftrightarrow\sigma'(id_t)("s\_enabled")=\mathtt{false}$.\label{it:not-sens-eq-re}
    \item
      $\forall{}t\in{}T,id_t\in{}Comps(\Delta)~s.t.~\gamma(t)=id_t,$\\
      $\sigma'(id_t)("s\_condition\_combination")=
      \prod\limits_{c\in{}conds(t)}
      \begin{cases}
        E_c(\tau,c) & if~\mathbb{C}(t,c)=1 \\
        \mathtt{not}(E_c(\tau,c)) & if~\mathbb{C}(t,c)=-1 \\
      \end{cases}$\\
      where
      $conds(t)=\{c\in\mathcal{C}~\vert~\mathbb{C}(t,c)=1\lor\mathbb{C}(t,c)=-1\}$.\label{it:cond-comb-eq-re}
    \end{enumerate}
  \end{frameb}
  
  Each point is proved by a separate lemma:
  \begin{itemize}[label=--]
  \item Apply Lemma~\nameref{lem:re-equal-marking} to solve
    \ref{it:marking-eq-re}.
  \item Apply Lemma~\nameref{lem:re-equal-tc} lemma to solve
    \ref{it:time-count-eq-re}.
  \item Apply ``Rising Edge Equal Reset Orders'' lemma to solve
    \ref{it:reset-eq-re}.
  \item Apply ``Rising Edge Equal Action'' lemma to solve \ref{it:action-eq-re}.
  \item Apply ``Rising Edge Equal Function'' lemma to solve \ref{it:fun-eq-re}.
  \item Apply ``Rising Edge Equal Sensitized'' lemma to solve \ref{it:sens-eq-re}.
  \item Apply ``Rising Edge Equal Not Sensitized'' lemma to solve \ref{it:not-sens-eq-re}.
  \item Apply Lemma~\nameref{lem:re-equal-cond-comb} to solve
    \ref{it:cond-comb-eq-re}.
  \end{itemize}
\end{proof}

\subsection{Rising Edge and Marking}
\label{sec:re-marking}

%%%%%%%%%%%%%%%%%%%%%%%%%%%%%%%%%%%%%%%%%%%%%%%%%%%%%
%%%%%%%%%% RISING EDGE EQUAL MARKING LEMMA %%%%%%%%%%
%%%%%%%%%%%%%%%%%%%%%%%%%%%%%%%%%%%%%%%%%%%%%%%%%%%%%

\begin{lemma}[Rising Edge Equal Marking]
  \label{lem:re-equal-marking}
  \rehyps{} then $\forall{}p,id_p~s.t.~\gamma(p)=id_p$ and
  $\sigma'(id_p)=\sigma'_p,~s'.M(p)=\sigma'_p("s\_marking")$.
\end{lemma}

\begin{proof}

  Given a $p\in{}P$, let us show
  \fbox{$s'.M(p)=\sigma'(id_p)("s\_marking")$.}

  \exP
  \noindent{}By definition of the SITPN state transition relation on rising edge:
  \begin{equation}\label{eq:re-eq-marking-eqmp}
    s'.M(p)=s.M(p)-\sum\limits_{t\in{}Fired(s)}pre(p,t)+\sum\limits_{t\in{}Fired(s)}post(t,p)
  \end{equation}

  \noindent{}By property of the $\mathtt{Inject}_\uparrow$, the
  \hvhdl{} rising edge and the stabilize relations, and \InCsCompP{} :
  \begin{equation}\label{eq:re-eq-marking-eqsm}
    \begin{split}
      \sigma'(id_p)("sm")=\sigma(id_p)("sm")-\sigma(id_p)("s\_output\_token\_sum")\\
      +\sigma(id_p)("s\_input\_token\_sum")
    \end{split}
  \end{equation}
  
  \noindent{}By the definition of \nameref{def:post-fe-state-sim}
  relation:
  \begin{eqnarray}
    s.M(p)=\sigma(id_p)("sm") \label{eq:re-eq-marking-eqm}\\
    \sum\limits_{t\in{}Fired(s)}pre(p,t)=\sigma(id_p)("sots") \label{eq:re-eq-marking-eqpre}\\
    \sum\limits_{t\in{}Fired(s)}post(t,p)=\sigma(id_p)("sits") \label{eq:re-eq-marking-eqpost}
  \end{eqnarray}
  
  \noindent{}Rewriting the goal with \ref{eq:re-eq-marking-eqmp},
  \ref{eq:re-eq-marking-eqsm}, \ref{eq:re-eq-marking-eqm},
  \ref{eq:re-eq-marking-eqpre} and \ref{eq:re-eq-marking-eqpost},
  \qedbox{tautology}.

\end{proof}

%%%%%%%%%%%%%%%%%%%%%%%%%%%%%%%%%%%%%%%%%%%%%%%%%%%%%%%%%%%%%%
%%%%%%%%%% RISING EDGE EQUAL CONDITION COMBINATION  %%%%%%%%%%
%%%%%%%%%%%%%%%%%%%%%%%%%%%%%%%%%%%%%%%%%%%%%%%%%%%%%%%%%%%%%%

\subsection{Rising edge and condition combination}
\label{sec:re-cond-comb}

\begin{lemma}[Rising Edge Equal Condition Combination]
  \label{lem:re-equal-cond-comb}
  \rehyps{} then\\
  $\forall{}t\in{}T,id_t\in{}Comps(\Delta)~s.t.~\gamma(t)=id_t,$\\
  $\sigma'(id_t)("s\_condition\_combination")=
  \prod\limits_{c\in{}conds(t)}
  \begin{cases}
    E_c(\tau,c) & if~\mathbb{C}(t,c)=1 \\
    \mathtt{not}(E_c(\tau,c)) & if~\mathbb{C}(t,c)=-1 \\
  \end{cases}$\\
  where
  $conds(t)=\{c\in\mathcal{C}~\vert~\mathbb{C}(t,c)=1\lor\mathbb{C}(t,c)=-1\}$.
\end{lemma}

\begin{proof}
  Given a $t$ and an $id_t$ s.t. $\gamma(t)=id_t$, let us show\\
  \fbox{$\sigma'(id_t)("s\_condition\_combination")=
    \prod\limits_{c\in{}conds(t)}
    \begin{cases}
      E_c(\tau,c) & if~\mathbb{C}(t,c)=1 \\
      \mathtt{not}(E_c(\tau,c)) & if~\mathbb{C}(t,c)=-1 \\
    \end{cases}$.}\\

  \exT

  \noindent By property of the \hvhdl{} stabilize relation, and\\
  $\mathtt{comp}(id_t,"transition",gm_t,ipm_t,opm_t)\in{}d.cs$:
  \begin{equation}
    \sigma'(id_t)("scc")=\prod\limits_{i=0}^{\Delta(id_t)("conditions\_number")-1}\sigma'(id_t)("input\_conditions")[i]\label{eq:re-eq-cc-eqscc}
  \end{equation}

  \noindent{}Rewriting the goal with \ref{eq:re-eq-cc-eqscc},\\
  \fbox{$\prod\limits_{i=0}^{\Delta(id_t)("cn")-1}\sigma'(id_t)("ic")[i]=
    \prod\limits_{c\in{}conds(t)}
    \begin{cases}
      E_c(\tau,c) & if~\mathbb{C}(t,c)=1 \\
      \mathtt{not}(E_c(\tau,c)) & if~\mathbb{C}(t,c)=-1 \\
    \end{cases}$.}\\

  \noindent{}Case analysis on $conds(t)$ (2 CASES):

  \begin{itemize}
  \item \textbf{CASE} $conds(t)=\emptyset$:\\
    \fbox{$\prod\limits_{i=0}^{\Delta(id_t)("cn")-1}\sigma'(id_t)("ic")[i]=\mathtt{true}$.}\\
    
    \noindent{}By construction,
    ${<}\mathtt{conditions\_number\Rightarrow{}1}{>}\in{}gm_t$ and\\
    ${<}\mathtt{input\_conditions(0)\Rightarrow{}true}{>}\in{}ipm_t$.

    \noindent{}By property of the stabilize relation,
    ${<}\mathtt{conditions\_number\Rightarrow{}1}{>}\in{}gm_t$ and
    ${<}\mathtt{input\_conditions(0)\Rightarrow{}true}{>}\in{}ipm_t$:
    \begin{eqnarray}
      \Delta(id_t)("cn")=1\label{eq:re-eq-cc-eqcn1}\\
      \sigma'(id_t)("ic")[0]=\mathtt{true}\label{eq:re-eq-cc-eqic0}
    \end{eqnarray}

    \noindent{}Rewriting the goal with \ref{eq:re-eq-cc-eqcn1} and
    \ref{eq:re-eq-cc-eqic0}, \qedbox{tautology.}
    
  \item \textbf{CASE} $conds(t)\neq\emptyset$:\\
    \noindent{}By construction,
    ${<}\mathtt{conditions\_number\Rightarrow{}\vert{}conds(t)\vert}{>}\in{}gm_t$,
    and by property of the stabilize relation:
    \begin{equation}
      \Delta(id_t)("cn")=\vert{}conds(t)\vert\label{eq:re-eq-cc-eqcn}
    \end{equation}
    
    Rewriting the goal with \eqref{eq:re-eq-cc-eqcn},\\
    \fbox{$\prod\limits_{i=0}^{\vert{}conds(t)\vert-1}\sigma'(id_t)("ic")[i]=
      \prod\limits_{c\in{}conds(t)}
      \begin{cases}
        E_c(\tau,c) & if~\mathbb{C}(t,c)=1 \\
        \mathtt{not}(E_c(\tau,c)) & if~\mathbb{C}(t,c)=-1 \\
      \end{cases}$.}\\

    \noindent{}Applying Theorem~\nameref{thm:big-op-eq}, there are two points to prove:

    \begin{enumerate}
    \item \qedbox{$\vert{}conds(t)\vert=\vert{}conds(t)\vert$}
    \item $\exists$ an injection
      $\iota\in[0,\vert{}conds(t)\vert-1]\rightarrow{}conds(t)$ s.t.\\
      $\forall{}i\in[0,\vert{}conds(t)\vert-1],~
      \sigma'(id_t)("ic")[i]=
      \begin{cases}
        E_c(\tau,\iota(i)) & if~\mathbb{C}(t,\iota(i))=1 \\
        \mathtt{not}(E_c(\tau,\iota(i))) & if~\mathbb{C}(t,\iota(i))=-1 \\
      \end{cases}$
    \end{enumerate}

    \noindent{}By construction, there exists a bijection
    $\beta\in[0,\vert{}conds(t)\vert-1]\rightarrow{}conds(t)$ such
    that for all $i\in[0,\vert{}conds(t)\vert-1]$, there exists an
    $id_c\in{}Ins(\Delta)$ and:
    \begin{itemize}
    \item $\gamma(\beta(i))=id_c$
    \item $\mathbb{C}(t,\beta(i))=1$ implies ${<}\mathtt{input\_conditions(i)\Rightarrow{}id_c}{>}\in{}ipm_t$
    \item $\mathbb{C}(t,\beta(i))=-1$ implies ${<}\mathtt{input\_conditions(i)\Rightarrow{}not~id_c}{>}\in{}ipm_t$
    \end{itemize}

    Let us take such a bijection $\beta$ to prove the goal. Then,
    given an $i\in[0,\vert{}conds(t)\vert-1]$, let us show
    \fbox{$\sigma'(id_t)("ic")[i]=
      \begin{cases}
        E_c(\tau,\beta(i)) & if~\mathbb{C}(t,\beta(i))=1 \\
        \mathtt{not}(E_c(\tau,\beta(i))) & if~\mathbb{C}(t,\beta(i))=-1 \\
      \end{cases}$}

    \noindent{}By definition of $\beta(i)\in{}conds(t)$:
    \begin{equation}
      \mathbb{C}(t,\beta(i))=1\lor\mathbb{C}(t,\beta(i))=-1\label{eq:cond-or}
    \end{equation}

    Case analysis on \eqref{eq:cond-or}:
    \begin{itemize}
    \item \textbf{CASE} $\mathbb{C}(t,\beta(i))=1$: \fbox{$\sigma'(id_t)("ic")[i]=E_c(\tau,\beta(i))$}\\
      
      By property of $\beta$, there exists $id_c\in{}Ins(\Delta)$
      s.t. $\gamma(\beta(i))=id_c$ and\\
      ${<}\mathtt{input\_conditions(i)\Rightarrow{}id_c}{>}\in{}ipm_t$.

      \noindent{}By property of the stabilize relation and
      ${<}\mathtt{input\_conditions(i)\Rightarrow{}id_c}{>}\in{}ipm_t$:
      \begin{equation}
        \sigma'(id_t)("ic")[i]=\sigma'(id_c)\label{eq:eq-ic-idc}
      \end{equation}

      \noindent{}By property of the \hvhdl{} rising edge and stabilize
      relations, and $id_c\in{}Ins(\Delta)$:
      \begin{equation}
        \sigma'(id_c)=\sigma_i(id_c)\label{eq:eq-idc-idc}
      \end{equation}

      \noindent{}By property of the $\mathtt{Inject}_\uparrow$
      relation and $id_c\in{}Ins(\Delta)$:
      \begin{equation}
        \sigma_i(id_c)=E_p(\tau,\uparrow)(id_c)\label{eq:eq-idc-ep}
      \end{equation}

      \noindent{}By property of
      $\gamma\vdash{}E_p\stackrel{env}{=}E_c$:
      \begin{equation}
        E_p(\tau,\uparrow)(id_c)=E_c(\tau,c)\label{eq:eq-ep-ec}
      \end{equation}

      \noindent{}Rewriting the goal with \eqref{eq:eq-ic-idc},
      \eqref{eq:eq-idc-idc}, \eqref{eq:eq-idc-ep},
      \eqref{eq:eq-ep-ec}, \qedbox{tautology.}

    \item \textbf{CASE} $\mathbb{C}(t,c)=-1$: \fbox{$\sigma'(id_t)("ic")[i]=\mathtt{not}~E_c(\tau,\beta(i))$}\\
      By property of $\beta$, there exists $id_c\in{}Ins(\Delta)$
      s.t. $\gamma(\beta(i))=id_c$ and\\
      ${<}\mathtt{input\_conditions(i)\Rightarrow{}not~id_c}{>}\in{}ipm_t$.

      \noindent{}By property of the stabilize relation and
      ${<}\mathtt{input\_conditions(i)\Rightarrow{}\mathtt{not}~id_c}{>}\in{}ipm_t$:
      \begin{equation}
        \sigma'(id_t)("ic")[i]=\mathtt{not}~\sigma'(id_c)\label{eq:eq-ic-not-idc}
      \end{equation}

      \noindent{}Then, equations \eqref{eq:eq-idc-idc},
      \eqref{eq:eq-idc-ep} and \eqref{eq:eq-ep-ec} also hold this
      case.
      
      \noindent{}Rewriting the goal with \eqref{eq:eq-ic-not-idc},
      \eqref{eq:eq-idc-idc}, \eqref{eq:eq-idc-ep} and
      \eqref{eq:eq-ep-ec}, \qedbox{tautology.}
    \end{itemize}
  \end{itemize}
\end{proof}

\subsection{Rising edge and time counters}
\label{sec:re-tc}

\begin{lemma}[Rising Edge Equal Time Counters]
  \label{lem:re-equal-tc}
  \rehyps{} then\\
  $\forall{}t\in{}T_i,id_t\in{}Comps(\Delta)~s.t.~\gamma(t)=id_t$,\\
  $\big(upper(I_s(t))=\infty\land{}s'.I(t)\le{}lower(I_s(t))\Rightarrow$
  $s'.I(t)=\sigma'(id_t)("s\_time\_counter")\big)$\\
  $\land\big(upper(I_s(t))=\infty\land{}s'.I(t)>{}lower(I_s(t))\Rightarrow$
  $\sigma'(id_t)("s\_time\_counter")=lower(I_s(t))\big)$\\
  $\land\big(upper(I_s(t))\neq\infty\land{}s'.I(t)>{}upper(I_s(t))\Rightarrow$
  $\sigma'(id_t)("s\_time\_counter")=upper(I_s(t))\big)$\\
  $\land\big(upper(I_s(t))\neq\infty\land{}s'.I(t)\le{}upper(I_s(t))\Rightarrow$
  $s'.I(t)=\sigma'(id_t)("s\_time\_counter")\big)$.
\end{lemma}

\begin{proof}
  Given a $t\in{}T_i$ and an $id_t\in{}Comps(\Delta)$ s.t. $\gamma(t)=id_t$, let us show\\
  \noindent\fbox{\parbox{\lwidth}{$\big(upper(I_s(t))=\infty\land{}s'.I(t)\le{}lower(I_s(t))\Rightarrow$
      $s'.I(t)=\sigma'(id_t)("s\_time\_counter")\big)$\\
      $\land\big(upper(I_s(t))=\infty\land{}s'.I(t)>{}lower(I_s(t))\Rightarrow$
      $\sigma'(id_t)("s\_time\_counter")=lower(I_s(t))\big)$\\
      $\land\big(upper(I_s(t))\neq\infty\land{}s'.I(t)>{}upper(I_s(t))\Rightarrow$
      $\sigma'(id_t)("s\_time\_counter")=upper(I_s(t))\big)$\\
      $\land\big(upper(I_s(t))\neq\infty\land{}s'.I(t)\le{}upper(I_s(t))\Rightarrow$
      $s'.I(t)=\sigma'(id_t)("s\_time\_counter")\big)$}}\\

  \exT\\
  
  \noindent{}Then, there are 4 points to show:

  \begin{enumerate}
  \item\label{it:re-eq-tc-fst}
    \fbox{$upper(I_s(t))=\infty\land{}s'.I(t)>{}lower(I_s(t))\Rightarrow\sigma'(id_t)("s\_time\_counter")=lower(I_s(t)$}\\
    
    \noindent{}Assuming $upper(I_s(t))=\infty$, let us show \fbox{$\sigma'(id_t)("s\_time\_counter")=lower(I_s(t)$.}

    \noindent{}By property of the $\mathtt{Inject}_\uparrow$, \hvhdl{}
    rising edge and stabilize relations, and \\ \InCsCompT:
    \begin{equation}
      \sigma'(id_t)("s\_time\_counter")=\sigma(id_t)("s\_time\_counter")\label{eq:eq-stc-stc}
    \end{equation}

    \noindent{}By property of
    $\gamma\vdash{}s\stackrel{\downarrow}{\sim}\sigma$:
    \begin{equation}
      \sigma(id_t)("s\_time\_counter")=lower(I_s(t))\label{eq:eq-stc-low}
    \end{equation}

    \noindent{}Rewriting the goal with \eqref{eq:eq-stc-stc} and
    \eqref{eq:eq-stc-low}, \qedbox{tautology.}
    
  \item Proved in the same fashion as \ref{it:re-eq-tc-fst}.
  \item Proved in the same fashion as \ref{it:re-eq-tc-fst}.
  \item Proved in the same fashion as \ref{it:re-eq-tc-fst}.
  \end{enumerate}
\end{proof}

\subsection{Rising edge and reset orders}
\label{sec:re-reset-orders}

\begin{lemma}[Rising Edge Equal Reset Orders]
  \label{lem:re-equal-reset-orders}
  \rehyps{} then\\
  $\forall{}t\in{}T_i,id_t\in{}Comps(\Delta)~s.t.~\gamma(t)=id_t,$
  $s'.reset_t(t)=\sigma'(id_t)("s\_reinit\_time\_counter")$
\end{lemma}

\begin{proof}
  Given a $t\in{}T_i$ and an $id_t\in{}Comps(\Delta)$
  s.t. $\gamma(t)=id_t$, let us show\\
  \fbox{$s'.reset_t(t)=\sigma'(id_t)("s\_reinit\_time\_counter")$.}

  \exT

  \noindent{}By property of the \hvhdl{} stabilize relation and
  \InCsCompT{}:
  \begin{equation}
    \sigma'(id_t)("srtc")=\sum\limits_{i=0}^{\Delta(id_t)("input\_arcs\_number")-1}\sigma'(id_t)("reinit\_time")[i]\label{eq:eq-srtc-prod}
  \end{equation}

  \noindent{}Rewriting the goal with \eqref{eq:eq-srtc-prod},
  \fbox{$s'.reset_t(t)=\sum\limits_{i=0}^{\Delta(id_t)("ian")-1}\sigma'(id_t)("rt")[i]$.}\\
  
  \noindent{}Case analysis on $input(t)$ (2 CASES):

  \begin{itemize}
  \item \textbf{CASE} $input(t)=\emptyset$:\\

    By construction,
    ${<}\mathtt{input\_arcs\_number\Rightarrow{}1}{>}\in{}gm_t$, and
    by property of the elaboration relation:
    \begin{equation}
      \Delta(id_t)("ian")=1\label{eq:eq-ian-1}
    \end{equation}

    By construction, there exists an $id_{ft}\in{}Sigs(\Delta)$ s.t.
    ${<}\mathtt{reinit\_time(0)\Rightarrow{}id_{ft}}{>}\in{}ipm_t$ and
    ${<}\mathtt{fired\Rightarrow{}id_{ft}}{>}\in{}opm_t$, and by
    property of the \hvhdl{} stabilize relation and \\ \InCsCompT:
    \begin{eqnarray}
      \sigma'(id_t)("rt")[0]&=&\sigma'(id_{ft})\label{eq:eq-rt0-fired} \\
      \sigma'(id_{ft})&=&\sigma'(id_t)("fired")\label{eq:eq-fired-fired} \\
      \sigma'(id_t)("fired")&=&\sigma'(id_t)("s\_fired")\label{eq:eq-fired-sfired} \\
      \sigma'(id_t)("s\_fired")&=&\sigma'(id_t)("s\_firable").\sigma'(id_t)("s\_priority\_combination")\label{eq:eq-sfired-sfirable}
    \end{eqnarray}

    \noindent{}Rewriting the goal with \eqref{eq:eq-rt0-fired},
    \eqref{eq:eq-fired-fired}, \eqref{eq:eq-fired-sfired} and
    \eqref{eq:eq-sfired-sfirable},\\
    \fbox{$s'.reset_t(t)=\sigma'(id_t)("s\_firable").\sigma'(id_t)("s\_priority\_combination")$.}

    \noindent{}By property of the stabilize relation, and \InCsCompT:
    \begin{equation}
      \sigma'(id_t)("spc")=\prod\limits_{i=0}^{\Delta(id_t)("ian")-1}\sigma'(id_t)("priority\_authorizations")[i]\label{eq:eq-spc-sum}
    \end{equation}

    By construction,
    ${<}\mathtt{priority\_authorizations(0)\Rightarrow{}true}{>}$, and
    by property of the stabilize relation and \InCsCompT:
    \begin{equation}
      \sigma'(id_t)("priority\_authorizations")[i]=true\label{eq:eq-pauth-0}
    \end{equation}

    \noindent{}Rewriting the goal with \eqref{eq:eq-ian-1},
    \eqref{eq:eq-spc-sum} and \eqref{eq:eq-pauth-0}, and simplifying
    the equation,\\ \fbox{$s'.reset_t(t)=\sigma'(id_t)("s\_firable")$.}
    
    \noindent{}Case analysis on $t\in{}Fired(s)$ or
    $t\notin{}Fired(s)$:

    \begin{itemize}
    \item \textbf{CASE} $t\in{}Fired(s)$:
      
      \noindent{}By property of
      $E_c,\tau\vdash{}s\srarrow{\uparrow}{\fontsize{7}{7}\selectfont}s'$
      (Rule~\ref{it:reset-order}):
      \begin{equation}
        s'.reset_t(t)=\mathtt{true}\label{eq:eq-reset-fired-true}
      \end{equation}

      \noindent{}Rewriting the goal with
      \eqref{eq:eq-reset-fired-true},
      \fbox{$\sigma'(id_t)("s\_firable")=true$.}

      \noindent{}By property of the stabilize, the \hvhdl{} rising
      edge and the $\mathtt{Inject}_\uparrow$ relations, and
      \InCsCompT:
      \begin{equation}
        \sigma(id_t)("s\_firable")=\sigma'(id_t)("s\_firable")\label{eq:eq-sfirable}
      \end{equation}

      \noindent{}Rewriting the goal with \eqref{eq:eq-sfirable},
      \fbox{$\sigma(id_t)("s\_firable")=true$.}

      \noindent{}By property of
      $\gamma\vdash{}s\stackrel{\downarrow}{\sim}\sigma$:
      \begin{equation}
        t\in{}Firable(s)\Leftrightarrow\sigma(id_t)("sfa")=true\label{eq:eq-firable-sfa}
      \end{equation}

      Rewriting the goal with \eqref{eq:eq-firable-sfa},
      \fbox{$t\in{}Firable(s)$.}

      \noindent{}By property of $t\in{}Fired(s)$,
      \qedbox{$t\in{}Firable(s)$.}
      
    \item \textbf{CASE} $t\notin{}Fired(s)$:

      \noindent{}By property of $input(t)=\emptyset$:
      \begin{equation}
        t\in{}Sens(M-\sum\limits_{t_i\in{}Fired(s)}pre(t_i))\label{eq:sens-by-trans}
      \end{equation}

      \noindent{}By property of $t\notin{}Fired(s)$,
      \eqref{eq:sens-by-trans} and
      $E_c,\tau\vdash{}s\srarrow{\uparrow}{\fontsize{7}{7}\selectfont}s'$
      (Rule~\ref{it:no-reset-order}):
      \begin{equation}
        s'.reset_t(t)=false\label{eq:eq-reset-false}
      \end{equation}

      \noindent{}Rewriting the goal with \eqref{eq:eq-reset-false},
      \fbox{$\sigma'(id_t)("s\_firable")=false$.}

      \noindent{}By property of the stabilize, the \hvhdl{} rising
      edge and the $\mathtt{Inject}_\uparrow$ relations, and
      \InCsCompT, equation \eqref{eq:eq-sfirable} holds.

      \noindent{}Rewriting the goal with \eqref{eq:eq-sfirable},
      \fbox{$\sigma(id_t)("s\_firable")=false$.}
      
      
    \end{itemize}
    
  \item \textbf{CASE} $input(t)\neq{}\emptyset$:\\

    By construction,
    ${<}\mathtt{input\_arcs\_number\Rightarrow{}}\vert{}input(t)\vert{>}\in{}gm_t$, and
    by property of the \hvhdl{} elaboration relation:
    \begin{equation}
      \Delta(id_t)("ian")=\vert{}input(t)\vert\label{eq:eq-ian-inputt}
    \end{equation}

    Rewriting the goal with \eqref{eq:eq-ian-inputt},
    \fbox{$s'.reset_t(t)=\sum\limits_{i=0}^{\vert{}input(t)\vert-1}\sigma'(id_t)("rt")[i]$.}
    
    Case analysis on $t\in{}Fired(s)$ or $t\notin{}Fired(s)$:
    
    \begin{itemize}
    \item \textbf{CASE} $t\in{}Fired(s)$:
      
      \noindent{}By property of
      $E_c,\tau\vdash{}s\srarrow{\uparrow}{\fontsize{7}{7}\selectfont}s'$
      (Rule~\ref{it:reset-order}), equation
      \eqref{eq:eq-reset-fired-true} holds.

      \noindent{}Rewriting the goal with
      \eqref{eq:eq-reset-fired-true},
      \fbox{$\sum\limits_{i=0}^{\vert{}input(t)\vert-1}\sigma'(id_t)("rt")[i]=\mathtt{true}$.}
    \end{itemize}
    
    Given a $i\in[0,\vert{}input(t)\vert-1]$, let us show
    \fbox{$\sigma(id_t)("rt")[i]=false$.}
    
    By construction, and $input(t)\neq{}\emptyset$, there exist
    ${}p\in{}input(t)$ and $id_p\in{}Comps(\Delta)$
    s.t. $\gamma(p)=id_p$.

    \exP By construction for all $i\in{}[0,\vert{}input(t)\vert-1],$
    there exist $j\in[0,\vert{}output(p)\vert-1]$ and
    $id_{ji}\in{}Sigs(\Delta)$ s.t.
    ${<}\mathtt{reinit\_transition\_time(j)\Rightarrow{}id_{ji}}{>}\in{}opm_p$
    and ${<}\mathtt{reinit\_time(i)\Rightarrow{}id_{ji}}{>}\in{}ipm_t$.

    \noindent{}By property of the \hvhdl{} stabilize relation,
    ${<}\mathtt{reinit\_transition\_time(j)\Rightarrow{}id_{ji}}{>}\in{}opm_p$
    and
    ${<}\mathtt{reinit\_time(i)\Rightarrow{}id_{ji}}{>}\in{}ipm_t$,
    then $\sigma(id_t)("rt")[i]=\sigma(id_{ji})=\sigma(id_p)("rtt")[j]$.

    Rewriting $\sigma(id_t)("rt")[i]$ as $\sigma(id_{ji})$ and
    $\sigma(id_{ji})$ as $\sigma(id_p)("rtt")[j]$,
    \fbox{$\sigma(id_p)("rtt")[j]=false$.}

    \noindent{}By property of the \hvhdl{} rising edge and stabilize
    relations,\\
    \begin{equation*}
      \begin{split}
        \sigma(id_p)("rtt")[j]=& ((\sigma_0(id_p)("oat")[j]=\mathtt{BASIC}+\sigma_0(id_p)("oat")[j]=\mathtt{TEST}) \\
        & .(\sigma_0(id_p)("sm")-\sigma_0(id_p)("sots")<\sigma_0(id_p)("oaw")[j])\\
        & .(\sigma_0(id_p)("sots")>0))\\
        & +(\sigma_0(id_p)("otf")[j]) \\
      \end{split}
    \end{equation*}

    Rewriting the goal with the above equation,
    \begin{equation*}
      \fbox{$\begin{split}
          false=& ((\sigma_0(id_p)("oat")[j]=\mathtt{BASIC}+\sigma_0(id_p)("oat")[j]=\mathtt{TEST}) \\
          & .(\sigma_0(id_p)("sm")-\sigma_0(id_p)("sots")<\sigma_0(id_p)("oaw")[j])\\
          & .(\sigma_0(id_p)("sots")>0))\\
          & +(\sigma_0(id_p)("otf")[j]) \\
        \end{split}$}
    \end{equation*}
    
    \begin{todobox}
      Add a lemma + proof in section initial states for fired = false
      after initialization.
    \end{todobox}

    \noindent{}By property of the \hvhdl{} initialization and the
    $\mathtt{Inject}_{\uparrow}$ relations, then
    $\sigma_0(id_p)("otf")[j]=false$. Rewriting
    $\sigma_0(id_p)("otf")[j]$ as $false$ and simplifying the goal,
    \begin{equation*}
      \fbox{$\begin{split}
          false=& ((\sigma_0(id_p)("oat")[j]=\mathtt{BASIC}+\sigma_0(id_p)("oat")[j]=\mathtt{TEST}) \\
          & .(\sigma_0(id_p)("sm")-\sigma_0(id_p)("sots")<\sigma_0(id_p)("oaw")[j])\\
          & .(\sigma_0(id_p)("sots")>0))\\
        \end{split}$}
    \end{equation*}

    \begin{todobox}
      Add a lemma + proof in section initial states for output token
      sum = 0 after initialization.
    \end{todobox}
    
    \noindent{}By property of the \hvhdl{} initialization and the
    $\mathtt{Inject}_{\uparrow}$ relations, then
    $\sigma_0(id_p)("sots")=0$. Rewriting
    $\sigma_0(id_p)("sots")$ as $0$ and simplifying the goal,
    \qedbox{$false=false$}
    
  \end{itemize}
\end{proof}

%%% Local Variables:
%%% mode: latex
%%% TeX-master: "../../main"
%%% End:


\section{Falling Edge}
\label{sec:falling-edge}
% %%%%%%%%%%%%%%%%%%%%%%%%%%%%%%%%%%%%%%%%%%%%%
%%%%%%%%%% FALLING EDGE HYPOTHESES %%%%%%%%%%
%%%%%%%%%%%%%%%%%%%%%%%%%%%%%%%%%%%%%%%%%%%%%

\begin{definition}[Falling Edge Hypotheses]
  \label{def:fe-hyps}
  Given an $sitpn\in{}SITPN$, $d\in{}design$,
  $\gamma\in{}WM(sitpn,d)$,
  $E_c\in\mathbb{N}\rightarrow\mathcal{C}\rightarrow\mathbb{B}$,
  $\Delta\in{}ElDesign(d,\mathcal{D}_\mathcal{H})$,
  $E_p\in(\mathbb{N}\times\{\uparrow,\downarrow\})\rightarrow{}Ins(\Delta)\rightarrow{}value$,
  $\tau\in\mathbb{N}$, $s,s'\in{}S(sitpn)$,
  $\sigma_e,\sigma,\sigma_i,\sigma_\downarrow,\sigma'\in\Sigma(\Delta)$,
  assume that:
  \begin{itemize}
  \item $\lfloor{}sitpn\rfloor_\mathcal{H}=(d,\gamma)$ and
    $\gamma\vdash{}E_p\stackrel{env}{=}E_c$ and
    $\mathcal{D}_\mathcal{H},\emptyset\vdash\mathrm{d}\srarrow{elab}{\fontsize{5}{7}\selectfont}\Delta,\sigma_e$
  \item $\gamma,E_c,\tau\vdash{}s\stackrel{\uparrow}{\sim}\sigma$
  \item $E_c,\tau\vdash{}s\xrightarrow{\downarrow}s'$
  \item $\mathtt{Inject}_\downarrow(\sigma, E_p, \tau, \sigma_i)$
    and
    $\Delta,\sigma_i\vdash\mathrm{d.cs}\xrightarrow{\downarrow}\sigma_{\downarrow}$
    and
    $\Delta,\sigma_{\downarrow}\vdash\mathrm{d.cs}\xrightarrow{\rightsquigarrow}\sigma'$
  \item State $\sigma$ is a stable design state:
    $\mathcal{D}_{\mathcal{H}},\Delta,\sigma\vdash\mathrm{d.cs}\xrightarrow{comb}\sigma$
  \end{itemize}
\end{definition}

\def\fehyps{For all $sitpn$, $d$, $\gamma$, $\Delta$, $\sigma_e$,
  $E_c$, $E_p$, $\tau$, $s$, $s'$, $\sigma$, $\sigma_i$,
  $\sigma_{\downarrow}$, $\sigma'$ that verify the hypotheses of
  Def.~\ref{def:fe-hyps},}

%%%%%%%%%%%%%%%%%%%%%%%%%%%%%%%%%%%%%%% 
%%%%%%%%%% FALLING EDGE LEMMA %%%%%%%%%
%%%%%%%%%%%%%%%%%%%%%%%%%%%%%%%%%%%%%%% 

\begin{lemma}[Falling Edge]
  \label{lem:fe}
  \fehyps{} then $\gamma\vdash{}s'\stackrel{\downarrow}{\sim}{}\sigma'$.
\end{lemma}

\begin{proof}
  By definition of \nameref{def:post-fe-state-sim}, there are 12
  points to prove.
  \begin{frameb}
    \begin{enumerate}
    \item
      $\forall{}p\in{}P,id_p\in{}Comps(\Delta)~s.t.~\gamma(p)=id_p,$
      $~s'.M(p)=\sigma'(id_p)("s\_marking")$.\label{item:6}
    \item
      $\forall{}t\in{}T_i,id_t\in{}Comps(\Delta)~s.t.~\gamma(t)=id_t,$\\
      $\big(upper(I_s(t))=\infty\land{}s'.I(t)\le{}lower(I_s(t))\Rightarrow{}s'.I(t)=\sigma'(id_t)("s\_time\_counter")\big)$\\
      $\land\big(upper(I_s(t))=\infty\land{}s'.I(t)>{}lower(I_s(t))\Rightarrow{}\sigma'(id_t)("s\_time\_counter")=lower(I_s(t))\big)$\\
      $\land\big(upper(I_s(t))\neq\infty\land{}s'.I(t)>{}upper(I_s(t))\Rightarrow{}\sigma'(id_t)("s\_time\_counter")=upper(I_s(t))\big)$\\
      $\land\big(upper(I_s(t))\neq\infty\land{}s'.I(t)\le{}upper(I_s(t))\Rightarrow{}s'.I(t)=\sigma'(id_t)("s\_time\_counter")\big)$.\label{item:9}
    \item
      $\forall{}t\in{}T_i,id_t\in{}Comps(\Delta)~s.t.~\gamma(t)=id_t,$
      $s'.reset_t(t)=\sigma'(id_t)("s\_reinit\_time\_counter")$.
    \item
      $\forall{}c\in\mathcal{C},id_c\in{}Ins(\Delta)~s.t.~\gamma(c)=id_c,~s'.cond(c)=\sigma'(id_c)$.
    \item
      $\forall{}a\in\mathcal{A},id_a\in{}Outs(\Delta)~s.t.~\gamma(a)=id_a,~s'.ex(a)=\sigma'(id_a)$.
    \item
      $\forall{}f\in\mathcal{F},id_f\in{}Outs(\Delta)~s.t.~\gamma(f)=id_f,~s'.ex(f)=\sigma'(id_f)$.
    \item
      $\forall{}t\in{}T,id_t\in{}Comps(\Delta)~s.t.~\gamma(t)=id_t,$
      $t\in{}Firable(s')\Leftrightarrow\sigma'(id_t)("s\_firable")=\mathtt{true}$.
    \item
      $\forall{}t\in{}T,id_t\in{}Comps(\Delta)~s.t.~\gamma(t)=id_t,$
      $t\notin{}Firable(s')\Leftrightarrow\sigma'(id_t)("s\_firable")=\mathtt{false}$.
    \item
      $\forall{}t\in{}T,id_t\in{}Comps(\Delta)~s.t.~\gamma(t)=id_t,$
      $t\in{}Fired(s')\Leftrightarrow\sigma'(id_t)("fired")=\mathtt{true}$.
    \item
      $\forall{}t\in{}T,id_t\in{}Comps(\Delta)~s.t.~\gamma(t)=id_t,$
      $t\notin{}Fired(s')\Leftrightarrow\sigma'(id_t)("fired")=\mathtt{false}$.
    \item
      $\forall{}p\in{}P,id_p\in{}Comps(\Delta)~s.t.~\gamma(p)=id_p,$
      $\sum\limits_{t\in{}Fired(s')}pre(p,t)=\sigma'(id_p)("s\_output\_token\_sum")$.\label{item:7}
    \item
      $\forall{}p\in{}P,id_p\in{}Comps(\Delta)~s.t.~\gamma(p)=id_p,$
      $\sum\limits_{t\in{}Fired(s')}post(t,p)=\sigma'(id_p)("s\_input\_token\_sum")$.\label{item:8}
    \end{enumerate}
  \end{frameb}

  Each point is proved by a separate lemma:
  \begin{itemize}[label=--]
  \item Apply Lemma~\nameref{lem:fe-equal-marking} to solve
    \ref{item:6}.
  \item Apply Lemma~\nameref{lem:fe-equal-tc} to solve \ref{item:9}.
  \item Apply Lemma~\nameref{lem:fe-equal-ots} to solve \ref{item:7}.
  \item Apply Lemma~\nameref{lem:fe-equal-its} to solve \ref{item:8}.
  \end{itemize}
  
\end{proof}

\subsection{Falling Edge and marking}
\label{sec:fe-marking}

%%%%%%%%%%%%%%%%%%%%%%%%%%%%%%%%%%%%%%%%%%%%%%%%%%%%%%
%%%%%%%%%% FALLING EDGE EQUAL MARKING LEMMA %%%%%%%%%%
%%%%%%%%%%%%%%%%%%%%%%%%%%%%%%%%%%%%%%%%%%%%%%%%%%%%%%

\begin{lemma}[Falling Edge Equal Marking]
  \label{lem:fe-equal-marking}
  \fehyps{} then
  $\forall{}p\in{}P,id_p\in{}Comps(\Delta)~s.t.~\gamma(p)=id_p,$
  $~s'.M(p)=\sigma'(id_p)("s\_marking")$.
\end{lemma}

\begin{proof}
  Given a $p\in{}P$ and an $id\in{}Comps(\Delta)$
  s.t. $\gamma(p)=id_p$, let us show\\
  \fbox{$s'.M(p)=\sigma'(id_p)("s\_marking")$.}

  \noindent{}By definition of
  $E_c,\tau\vdash{}sitpn,s\xrightarrow{\downarrow}s'$:
  \begin{equation}
    \label{eq:42}s.M(p)=s'.M(p)
  \end{equation}

  \noindent{}By property of the $\mathtt{Inject}_\downarrow$ relation,
  the \hvhdl{} falling edge relation, the stabilize relation and
  \InCsCompP:
  \begin{equation}
    \label{eq:43}\sigma'(id_p)("s\_marking")=\sigma(id_p)("s\_marking")
  \end{equation}

  \noindent{}Rewriting the goal with \eqref{eq:42} and \eqref{eq:43}:
  \fbox{$s.M(p)=\sigma(id_p)("s\_marking")$.}

  \noindent{}By definition of
  $\gamma,E_c,\tau\vdash{}s\stackrel{\downarrow}{\sim}\sigma$:
  \qedbox{$s.M(p)=\sigma(id_p)("s\_marking")$.}
  
\end{proof}

%%%%%%%%%%%%%%%%%%%%%%%%%%%%%%%%%%%%%%%%%%%%%%%%%%%%%%%%%%%%%%%
%%%%%%%%%% FALLING EDGE EQUAL OUTPUT TOKEN SUM LEMMA %%%%%%%%%%
%%%%%%%%%%%%%%%%%%%%%%%%%%%%%%%%%%%%%%%%%%%%%%%%%%%%%%%%%%%%%%%

\begin{lemma}[Falling Edge Equal Output Token Sum]
  \label{lem:fe-equal-ots}
  \fehyps{} then $\forall{}p,id_p~s.t.~\gamma(p)=id_p$,
  $\sum\limits_{t\in{}Fired(s')}pre(p,t)=\sigma'(id_p)("s\_output\_token\_sum")$.
\end{lemma}

\begin{proof}
  Given a $p\in{}P$ and an $id_p\in{}Comps(\Delta)$, let us show\\
  \fbox{$\sum\limits_{t\in{}Fired(s')}pre(p,t)=\sigma'(id_p)("s\_output\_token\_sum")$.}

  \exP{}
  
  \noindent{}By property of the stabilize relation and \InCsCompP{}:
  \begin{equation}
    \label{eq:16}
    \sigma'(id_p)("sots")=\sum\limits_{i=0}^{\Delta(id_p)("oan")-1}
    \begin{cases}
      \sigma'(id_p)("oaw")[i]~\mathtt{if}~(\sigma'(id_p)("otf")[i]~ \\
      \hspace{19.5ex}.~\sigma'(id_p)("oat")[i]=\mathtt{BASIC}) \\
      0~otherwise \\
    \end{cases}
  \end{equation}

  \noindent{}Rewriting the goal with \eqref{eq:16}:\\
  \begin{equation*}
    \fbox{$
      \sum\limits_{t\in{}Fired(s')}pre(p,t)=\sum\limits_{i=0}^{\Delta(id_p)("oan")-1}
      \begin{cases}
        \sigma'(id_p)("oaw")[i]~\mathtt{if}~(\sigma'(id_p)("otf")[i]~ \\
        \hspace{19.5ex}.~\sigma'(id_p)("oat")[i]=\mathtt{BASIC}) \\
        0~otherwise \\
      \end{cases}$}
  \end{equation*}

  \noindent{}Let us unfold the definition of the left sum term:\\
  \fbox{
    {\begin{tabular}{c}
       $\sum\limits_{t\in{}Fired(s')}
       \begin{cases}
         \omega~\mathtt{if}~pre(p,t)=(\omega,\mathtt{basic}) \\
         0~otherwise
       \end{cases}$ \\
       $=$ \\
       $\sum\limits_{i=0}^{\Delta(id_p)("oan")-1}
       \begin{cases}
         \sigma'(id_p)("oaw")[i]~\mathtt{if}~(\sigma'(id_p)("otf")[i]~ \\
         \hspace{19.5ex}.~\sigma'(id_p)("oat")[i]=\mathtt{BASIC}) \\
         0~otherwise \\
       \end{cases}$ \\
     \end{tabular}}
 }\\

 \noindent{}To ease the reading, let us define functions
 $f\in{}Fired(s')\rightarrow\mathbb{N}$ and
 $g\in[0,\vert{}output(p)\vert-1]\rightarrow\mathbb{N}$ s.t.
 $f(t)=\begin{cases}
   \omega~\mathtt{if}~pre(p,t)=(\omega,\mathtt{basic}) \\
   0~otherwise
 \end{cases}$ and $g(i)=\begin{cases}
   \sigma'(id_p)("oaw")[i]~\mathtt{if}~(\sigma'(id_p)("otf")[i]~ \\
   \hspace{19.5ex}.~\sigma'(id_p)("oat")[i]=\mathtt{BASIC}) \\
   0~otherwise \\
 \end{cases}$

 \noindent{}Then, the goal is: \fbox{$\sum\limits_{t\in{}Fired(s')}f(t)=\sum\limits_{i=0}^{\Delta(id_p)("oan")-1}g(i)$}\\

 \noindent{}Let us perform case analysis on $output(p)$; there are two cases:

 \begin{enumerate}
 \item $output(p)=\emptyset$:\\
   
   \noindent{}By construction,
   ${<}\mathtt{output\_arcs\_number\Rightarrow{}1}{>}\in{}gm_p$,
   ${<}\mathtt{output\_arcs\_types(0)\Rightarrow{}\mathtt{BASIC}}{>}\in{}ipm_p$,
   ${<}\mathtt{output\_transitions\_fired(0)\Rightarrow{}\mathtt{true}}{>}\in{}ipm_p$,
   and
   ${<}\mathtt{output\_arcs\_weights(0)\Rightarrow{}0}{>}\in{}ipm_p$.

   \noindent{}By property of the elaboration relation and \InCsCompP:
   \begin{equation}
     \label{eq:25}\Delta(id_p)("oan")=1
   \end{equation}

   \noindent{}By property of the stabilize relation and \InCsCompP:
   \begin{eqnarray}
     \label{eq:26}\sigma'(id_p)("oat")[0]&=&\mathtt{BASIC}\\
     \label{eq:27}\sigma'(id_p)("otf")[0]&=&\mathtt{true}\\
     \label{eq:28}\sigma'(id_p)("oaw")[0]&=&0
   \end{eqnarray}

   By property of $output(p)=\emptyset$:
   \begin{equation}\label{eq:29}
     \sum\limits_{t\in{}Fired(s')}
     \begin{cases}
       \omega~\mathtt{if}~pre(p,t)=(\omega,\mathtt{basic}) \\
       0~otherwise
     \end{cases}=0
   \end{equation}

   \noindent{}Rewriting the goal with \eqref{eq:25}, \eqref{eq:26},
   \eqref{eq:27}, \eqref{eq:28} and \eqref{eq:29}, \qedbox{tautology.}
   
 \item $output(p)\neq\emptyset$:\\

   \noindent{}By construction,
   ${<}\mathtt{output\_arcs\_number\Rightarrow{}}\vert{}output(p)\vert{>}\in{}gm_p$,
   and by property of the elaboration relation:
   \begin{equation}
     \label{eq:30}\Delta(id_p)("oan")=\vert{}output(p)\vert
   \end{equation}
   
   Rewriting the goal with \eqref{eq:30}:
   \fbox{$\sum\limits_{t\in{}Fired(s')}f(t)=\sum\limits_{i=0}^{\vert{}output(p)\vert-1}g(i)$.}

   \noindent{}Let us reason by induction on the right sum term of the goal.

   \begin{itemize}
   \item \textbf{BASE CASE}:\\

     \noindent{}In that case, $0>\vert{}output\vert-1$ and
     $\sum\limits_{i=0}^{\vert{}output(p)\vert-1}g(i)=0$.

     \noindent{}As $0>\vert{}output\vert-1$, then
     $\vert{}output(p)\vert=0$, thus \qedbox{contradicting
       $output(p)\neq\emptyset$.}

   \item \textbf{INDUCTION CASE}:\\

     \noindent{}In that case, $0\le\vert{}output(p)\vert-1$.
     
     \begin{ih}
       $\forall{}F\subseteq{}Fired(s'),$
       $g(0)+\sum\limits_{t\in{}F}f(t)=g(0)+\sum\limits_{i=1}^{\vert{}output(p)\vert-1}g(i)$
     \end{ih}

     \fbox{$\sum\limits_{t\in{}Fired(s')}f(t)=g(0)+\sum\limits_{i=1}^{\vert{}output(p)\vert-1}g(i)$}

     By definition of $g$:
     \begin{equation}
       \label{eq:31}g(0)=\begin{cases}
         \sigma'(id_p)("oaw")[0]~\mathtt{if}~(\sigma'(id_p)("otf")[0] \\
         \hspace{19.5ex}.~\sigma'(id_p)("oat")[0]=\mathtt{BASIC}) \\
         0~otherwise \\
       \end{cases}
     \end{equation}

     \noindent{}Let us perform case analysis on the value of
     $\sigma'(id_p)("otf")[0]~.~\sigma'(id_p)("oat")[0]=\mathtt{BASIC}$; there are two cases:
     \begin{enumerate}
     \item $(\sigma'(id_p)("otf")[0]~.~\sigma'(id_p)("oat")[0]=\mathtt{BASIC})=\mathtt{false}$:\\
       In that case, $g(0)=0$, and then we can apply the induction
       hypothesis with $F=Fired(s')$ to solve the goal:
       \qedbox{$\sum\limits_{t\in{}Fired(s')}f(t)=\sum\limits_{i=1}^{\vert{}output(p)\vert-1}g(i)$.}

     \item $(\sigma'(id_p)("otf")[0]~.~\sigma'(id_p)("oat")[0]=\mathtt{BASIC})=\mathtt{true}$:\\
       In that case, $g(0)=\sigma'(id_p)("oaw")[0]$,
       $\sigma'(id_p)("otf")[0]=\mathtt{true}$ and\\
       $\sigma'(id_p)("oat")[0]=\mathtt{BASIC}$.

       \noindent{}By construction, there exist a $t\in{}output(t)$,
       $id_t\in{}Comps(\Delta)$ s.t. $\gamma(t)=id_t$. Let us take
       such a $t\in{}output(p)$.

       \exT{}

       \noindent{}As $t\in{}output(p)$, there exist
       $\omega\in\mathbb{N}^{*}$ and
       $a\in\{\mathtt{BASIC},\mathtt{TEST},\mathtt{INHIB}\}$
       s.t. $pre(p,t)=(\omega,a)$. Let us take an $\omega$ and $a$
       s.t. $pre(p,t)=(\omega,a)$.

       \noindent{}By construction,
       ${<}\mathtt{output\_arcs\_types(0)\Rightarrow{}}a{>}\in{}ipm_p$,\\
       ${<}\mathtt{output\_arcs\_weights(0)\Rightarrow{}}\omega{>}\in{}ipm_p$,
       and there exists $id_{ft}\in{}Sigs(\Delta)$ s.t.
       ${<}\mathtt{fired\Rightarrow}id_{ft}{>}\in{}opm_t$ and
       ${<}\mathtt{output\_transitions\_fired(0)\Rightarrow{}id_{ft}}{>}\in{}ipm_p$

       \noindent{}By property of the stabilize relation,
       $\sigma'(id_p)("oat")[0]=\mathtt{BASIC}$ and\\
       ${<}\mathtt{output\_arcs\_types(0)\Rightarrow{}a}{>}\in{}ipm_p$:
       \begin{equation}
         \label{eq:32}pre(p,t)=(\omega,\mathtt{basic})
       \end{equation}

       \noindent{}By property of the stabilize relation,
       ${<}\mathtt{fired\Rightarrow{}id_{ft}}{>}\in{}opm_t$,\\
       ${<}\mathtt{output\_transitions\_fired(0)\Rightarrow{}id_{ft}}{>}\in{}ipm_p$
       and $\sigma'(id_p)("otf")[0]=\mathtt{true}$:
       \begin{equation}
         \label{eq:33}\sigma'(id_t)("fired")=\mathtt{true}
       \end{equation}

       \noindent{}Appealing to Lemma~\ref{lem:fe-equal-fired}, we know
       $t\in{}Fired(s')$.

       \noindent{}As $t\in{}Fired(s')$, we can rewrite the left sum
       term of the goal as follows:\\
       \fbox{$f(t)+\sum\limits_{t'\in{}Fired(s')\setminus\{t\}}f(t')=g(0)+\sum\limits_{i=1}^{\vert{}output(p)\vert-1}g(i)$}

       \noindent{}We know that $g(0)=\sigma'(id_p)("oaw")[0]$, and by
       property of the stabilize relation and
       ${<}\mathtt{output\_arcs\_weights(0)\Rightarrow}\omega{>}\in{}ipm_p$:
       \begin{equation}
         \label{eq:21}\sigma'(id_p)("oaw")[0]=\omega
       \end{equation}

       Rewriting the goal with \eqref{eq:21}:\\
       \fbox{$f(t)+\sum\limits_{t'\in{}Fired(s')\setminus\{t\}}f(t')=\omega+\sum\limits_{i=1}^{\vert{}output(p)\vert-1}g(i)$}

       By definition of $f$, and as $pre(p,t)=(\omega,\mathtt{basic})$, then $f(t)=\omega$; thus, rewriting the goal:\\
       \fbox{$\omega+\sum\limits_{t'\in{}Fired(s')\setminus\{t\}}f(t')=\omega+\sum\limits_{i=1}^{\vert{}output(p)\vert-1}g(i)$}

       Then, knowing that $g(0)=\omega$, we can apply the induction
       hypothesis with $F=Fired(s')\setminus\{t\}$:
       \qedbox{$g(0)+\sum\limits_{t'\in{}Fired(s')\setminus\{t\}}f(t')=g(0)+\sum\limits_{i=1}^{\vert{}output(p)\vert-1}g(i)$.}
     \end{enumerate}
     
   \end{itemize}
   

 \end{enumerate}
 
\end{proof}


%%%%%%%%%%%%%%%%%%%%%%%%%%%%%%%%%%%%%%%%%%%%%%%%%%%%%%%%%%%%%%
%%%%%%%%%% FALLING EDGE EQUAL INPUT TOKEN SUM LEMMA %%%%%%%%%%
%%%%%%%%%%%%%%%%%%%%%%%%%%%%%%%%%%%%%%%%%%%%%%%%%%%%%%%%%%%%%%

\begin{lemma}[Falling Edge Equal Input Token Sum]
  \label{lem:fe-equal-its}
  \fehyps{} then $\forall{}p,id_p~s.t.~\gamma(p)=id_p$,
  $~\sum\limits_{t\in{}Fired(s')}post(t,p)=\sigma'_p("s\_input\_token\_sum")$.
\end{lemma}

\begin{proof}
  Given a $p\in{}P$ and an $id_p\in{}Comps(\Delta)$, let us show\\
  \fbox{$\sum\limits_{t\in{}Fired(s')}post(t,p)=\sigma'(id_p)("s\_input\_token\_sum")$.}\\

  \exP{}
  
  \noindent{}By property of the stabilize relation and \InCsCompP{}:
  \begin{equation}
    \label{eq:34}
    \sigma'(id_p)("sits")=\sum\limits_{i=0}^{\Delta(id_p)("ian")-1}
    \begin{cases}
      \sigma'(id_p)("iaw")[i]~\mathtt{if}~\sigma'(id_p)("itf")[i]\\
      0~otherwise \\
    \end{cases}
  \end{equation}

  \noindent{}Rewriting the goal with \eqref{eq:34}:\\
  \begin{equation*}
    \fbox{$
      \sum\limits_{t\in{}Fired(s')}post(t,p)=\sum\limits_{i=0}^{\Delta(id_p)("ian")-1}
      \begin{cases}
        \sigma'(id_p)("iaw")[i]~\mathtt{if}~\sigma'(id_p)("otf")[i]\\
        0~otherwise \\
      \end{cases}$}
  \end{equation*}

  \noindent{}Let us unfold the definition of the left sum term:\\
  \fbox{
    {\begin{tabular}{c}
       $\sum\limits_{t\in{}Fired(s')}
       \begin{cases}
         \omega~\mathtt{if}~post(t,p)=\omega \\
         0~otherwise
       \end{cases}$ \\
       $=$ \\
       $\sum\limits_{i=0}^{\Delta(id_p)("ian")-1}
       \begin{cases}
         \sigma'(id_p)("iaw")[i]~\mathtt{if}~\sigma'(id_p)("itf")[i]\\
         0~otherwise \\
       \end{cases}$ \\
     \end{tabular}}
 }\\

 \noindent{}Let us perform case analysis on $input(p)$; there are two cases:

 \begin{enumerate}
 \item $input(p)=\emptyset$:\\
   
   \noindent{}By construction,
   ${<}\mathtt{input\_arcs\_number\Rightarrow{}1}{>}\in{}gm_p$,
   ${<}\mathtt{input\_transitions\_fired(0)\Rightarrow{}\mathtt{true}}{>}\in{}ipm_p$,
   and
   ${<}\mathtt{input\_arcs\_weights(0)\Rightarrow{}0}{>}\in{}ipm_p$.

   \noindent{}By property of the elaboration relation and \InCsCompP:
   \begin{equation}
     \label{eq:35}\Delta(id_p)("ian")=1
   \end{equation}

   \noindent{}By property of the stabilize relation and \InCsCompP:
   \begin{eqnarray}
     \label{eq:36}\sigma'(id_p)("itf")[0]&=&\mathtt{true}\\
     \label{eq:37}\sigma'(id_p)("iaw")[0]&=&0
   \end{eqnarray}

   By property of $input(p)=\emptyset$:
   \begin{equation}\label{eq:23}
     \sum\limits_{t\in{}Fired(s')}
     \begin{cases}
       \omega~\mathtt{if}~post(t,p)=\omega \\
       0~otherwise
     \end{cases}=0
   \end{equation}

   \noindent{}Rewriting the goal with \eqref{eq:35}, \eqref{eq:36},
   \eqref{eq:37}, and \eqref{eq:23}, and simplifying the goal,
   \qedbox{tautology.}
   
 \item $input(p)\neq\emptyset$:\\

   \noindent{}By construction,
   ${<}\mathtt{input\_arcs\_number\Rightarrow{}}\vert{}input(p)\vert{>}\in{}gm_p$,
   and by property of the elaboration relation:
   \begin{equation}
     \label{eq:24}\Delta(id_p)("ian")=\vert{}input(p)\vert
   \end{equation}

   \noindent{}To ease the reading, let us define functions
   $f\in{}Fired(s')\rightarrow\mathbb{N}$ and
   $g\in[0,\vert{}input(p)\vert-1]\rightarrow\mathbb{N}$ s.t.
   $f(t)=\begin{cases}
     \omega~\mathtt{if}~post(t,p)=\omega \\
     0~otherwise
   \end{cases}$ and \\ $g(i)=\begin{cases}
     \sigma'(id_p)("iaw")[i]~\mathtt{if}~\sigma'(id_p)("itf")[i]\\
     0~otherwise \\
   \end{cases}$

   \noindent{}Then, the goal is: \fbox{$\sum\limits_{t\in{}Fired(s')}f(t)=\sum\limits_{i=0}^{\Delta(id_p)("ian")-1}g(i)$}\\
   
   Rewriting the goal with \eqref{eq:24}:
   \fbox{$\sum\limits_{t\in{}Fired(s')}f(t)=\sum\limits_{i=0}^{\vert{}input(p)\vert-1}g(i)$.}

   \noindent{}Let us reason by induction on the right sum term of the goal.

   \begin{itemize}
   \item \textbf{BASE CASE}:\\
     \noindent{}In that case, $0>\vert{}input(p)\vert-1$ and
     $\sum\limits_{i=0}^{\vert{}input(p)\vert-1}g(i)=0$.

     \noindent{}As $0>\vert{}input(p)\vert-1$, then
     $\vert{}input(p)\vert=0$, thus \qedbox{contradicting
       $input(p)\neq\emptyset$.}

   \item \textbf{INDUCTION CASE}:\\
     \noindent{}In that case, $0\le\vert{}input(p)\vert-1$.
     
     \begin{ih}
       $\forall{}F\subseteq{}Fired(s'),~$
       $g(0)+\sum\limits_{t\in{}F}f(t)=g(0)+\sum\limits_{i=1}^{\vert{}input(p)\vert-1}g(i)$
     \end{ih}

     \fbox{$\sum\limits_{t\in{}Fired(s')}f(t)=g(0)+\sum\limits_{i=1}^{\vert{}input(p)\vert-1}g(i)$}\\

     By definition of $g$:
     \begin{equation}
       \label{eq:38}g(0)=\begin{cases}
         \sigma'(id_p)("iaw")[0]~\mathtt{if}~\sigma'(id_p)("itf")[0]\\
         0~otherwise \\
       \end{cases}
     \end{equation}

     \noindent{}Let us perform case analysis on the value of
     $\sigma'(id_p)("itf")[0]$; there are two cases:
     \begin{enumerate}
     \item $\sigma'(id_p)("itf")[0]=\mathtt{false}$:\\
       In that case, $g(0)=0$, and then we can apply the induction
       hypothesis with $F=Fired(s')$ to solve the goal:
       \qedbox{$\sum\limits_{t\in{}Fired(s')}f(t)=\sum\limits_{i=1}^{\vert{}input(p)\vert-1}g(i)$.}

     \item $\sigma'(id_p)("itf")[0]=\mathtt{true}$:\\
       In that case, $g(0)=\sigma'(id_p)("iaw")[0]$ and
       $\sigma'(id_p)("itf")[0]=\mathtt{true}$ .

       \noindent{}By construction, there exist a $t\in{}input(t)$,
       $id_t\in{}Comps(\Delta)$ s.t. $\gamma(t)=id_t$. Let us take
       such a $t\in{}input(p)$.

       \exT{}

       \noindent{}As $t\in{}input(p)$, there exist
       $\omega\in\mathbb{N}^{*}$ s.t. $post(t,p)=\omega$. Let us take
       an $\omega$ s.t. $post(t,p)=\omega$.

       \noindent{}By construction,
       ${<}\mathtt{input\_arcs\_weights(0)\Rightarrow{}}\omega{>}\in{}ipm_p$,
       and there exists $id_{ft}\in{}Sigs(\Delta)$ s.t.
       ${<}\mathtt{fired\Rightarrow}id_{ft}{>}\in{}opm_t$ and
       ${<}\mathtt{input\_transitions\_fired(0)\Rightarrow{}id_{ft}}{>}\in{}ipm_p$

       \noindent{}By property of the stabilize relation and
       ${<}\mathtt{input\_arcs\_types(0)\Rightarrow{}a}{>}\in{}ipm_p$:
       \begin{equation}
         \label{eq:39}post(t,p)=\omega
       \end{equation}

       \noindent{}By property of the stabilize relation,
       ${<}\mathtt{fired\Rightarrow{}id_{ft}}{>}\in{}opm_t$,\\
       ${<}\mathtt{input\_transitions\_fired(0)\Rightarrow{}id_{ft}}{>}\in{}ipm_p$
       and $\sigma'(id_p)("itf")[0]=\mathtt{true}$:
       \begin{equation}
         \label{eq:40}\sigma'(id_t)("fired")=\mathtt{true}
       \end{equation}

       \noindent{}Appealing to Lemma~\ref{lem:fe-equal-fired} and
       \eqref{eq:40}, we know $t\in{}Fired(s')$.

       \noindent{}As $t\in{}Fired(s')$, we can rewrite the left sum
       term of the goal as follows:\\
       \fbox{$f(t)+\sum\limits_{t'\in{}Fired(s')\setminus\{t\}}f(t')=g(0)+\sum\limits_{i=1}^{\vert{}input(p)\vert-1}g(i)$}

       \noindent{}We know that $g(0)=\sigma'(id_p)("iaw")[0]$, and by
       property of the stabilize relation and
       ${<}\mathtt{input\_arcs\_weights(0)\Rightarrow}\omega{>}\in{}ipm_p$:
       \begin{equation}
         \label{eq:41}\sigma'(id_p)("iaw")[0]=\omega
       \end{equation}

       Rewriting the goal with \eqref{eq:41}:\\
       \fbox{$f(t)+\sum\limits_{t'\in{}Fired(s')\setminus\{t\}}f(t')=\omega+\sum\limits_{i=1}^{\vert{}input(p)\vert-1}g(i)$}

       By definition of $f$, and as $post(t,p)=\omega$, then $f(t)=\omega$; thus, rewriting the goal:\\
       \fbox{$\omega+\sum\limits_{t'\in{}Fired(s')\setminus\{t\}}f(t')=\omega+\sum\limits_{i=1}^{\vert{}input(p)\vert-1}g(i)$}

       Then, knowing that $g(0)=\omega$, we can apply the induction
       hypothesis with $F=Fired(s')\setminus\{t\}$:
       \qedbox{$g(0)+\sum\limits_{t'\in{}Fired(s')\setminus\{t\}}f(t')=g(0)+\sum\limits_{i=1}^{\vert{}input(p)\vert-1}g(i)$.}
     \end{enumerate}
     
   \end{itemize}
   
 \end{enumerate}
 
\end{proof}


%%%%%%%%%%%%%%%%%%%%%%%%%%%%%%%%%%%%%%%%%%%%%%%%%%%%%%
%%%%%%%%%% FALLING EDGE EQUAL TIME COUNTERS %%%%%%%%%%
%%%%%%%%%%%%%%%%%%%%%%%%%%%%%%%%%%%%%%%%%%%%%%%%%%%%%%

\subsection{Falling edge and time counters}
\label{sec:fe-equal-tc}

\begin{lemma}[Falling Edge Equal Time Counters]
  \label{lem:fe-equal-tc}
  \fehyps{} then $\forall{}t\in{}T_i,id_t\in{}Comps(\Delta)~s.t.~\gamma(t)=id_t,$\\
  $\big(upper(I_s(t))=\infty\land{}s'.I(t)\le{}lower(I_s(t))\Rightarrow{}s'.I(t)=\sigma'(id_t)("s\_time\_counter")\big)$\\
  $\land\big(upper(I_s(t))=\infty\land{}s'.I(t)>{}lower(I_s(t))\Rightarrow{}\sigma'(id_t)("s\_time\_counter")=lower(I_s(t))\big)$\\
  $\land\big(upper(I_s(t))\neq\infty\land{}s'.I(t)>{}upper(I_s(t))\Rightarrow{}\sigma'(id_t)("s\_time\_counter")=upper(I_s(t))\big)$\\
  $\land\big(upper(I_s(t))\neq\infty\land{}s'.I(t)\le{}upper(I_s(t))\Rightarrow{}s'.I(t)=\sigma'(id_t)("s\_time\_counter")\big)$.
\end{lemma}

\begin{proof}
  Given a $t\in{}T_i$ and an $id_t\in{}Comps(\Delta)$ s.t. $\gamma(t)=id_t$, let us show\\
  \noindent\fbox{\parbox{\lwidth}{$\big(upper(I_s(t))=\infty\land{}s'.I(t)\le{}lower(I_s(t))\Rightarrow$
      $s'.I(t)=\sigma'(id_t)("s\_time\_counter")\big)$\\
      $\land\big(upper(I_s(t))=\infty\land{}s'.I(t)>{}lower(I_s(t))\Rightarrow$
      $\sigma'(id_t)("s\_time\_counter")=lower(I_s(t))\big)$\\
      $\land\big(upper(I_s(t))\neq\infty\land{}s'.I(t)>{}upper(I_s(t))\Rightarrow$
      $\sigma'(id_t)("s\_time\_counter")=upper(I_s(t))\big)$\\
      $\land\big(upper(I_s(t))\neq\infty\land{}s'.I(t)\le{}upper(I_s(t))\Rightarrow$
      $s'.I(t)=\sigma'(id_t)("s\_time\_counter")\big)$}}\\

  \exT\\

  \noindent{}By property of the elaboration,
  $\mathtt{Inject}_\downarrow$, \hvhdl{} rising edge and stabilize
  relations, and \InCsCompT:
  \begin{equation}
    \begin{split}
      \sigma(id_t)("se")=\mathtt{true}\land\Delta(id_t)("tt")\neq\mathtt{NOT\_TEMPORAL}
      \land\sigma(id_t)("srtc")=\mathtt{false}\\
      \land\sigma(id_t)("stc")<\Delta(id_t)("mtc")\Rightarrow
      \sigma'(id_t)("stc")=\sigma(id_t)("stc")+1
    \end{split}
    \label{eq:44}
  \end{equation}

  \begin{equation}
    \begin{split}
      \sigma(id_t)("se")=\mathtt{true}\land\Delta(id_t)("tt")\neq\mathtt{NOT\_TEMPORAL}
      \land\sigma(id_t)("srtc")=\mathtt{false}\\
      \land\sigma(id_t)("stc")\ge\Delta(id_t)("mtc")\Rightarrow
      \sigma'(id_t)("stc")=\sigma(id_t)("stc")
    \end{split}
    \label{eq:47}
  \end{equation}

  \begin{equation}
    \begin{split}
      \sigma(id_t)("se")=\mathtt{true}\land\Delta(id_t)("tt")\neq\mathtt{NOT\_TEMPORAL}\\
      \land\sigma(id_t)("srtc")=\mathtt{true}\Rightarrow
      \sigma'(id_t)("stc")=1
    \end{split}
    \label{eq:48}
  \end{equation}

  \begin{equation}
    \begin{split}
      \sigma(id_t)("se")=\mathtt{false}\lor\Delta(id_t)("tt")=\mathtt{NOT\_TEMPORAL}\Rightarrow
      \sigma'(id_t)("stc")=0
    \end{split}
    \label{eq:49}
  \end{equation}
  
  \noindent{}Then, there are 4 points to show:

  \begin{enumerate}
  \item\label{it:re-eq-tc-fst}
    \fbox{$upper(I_s(t))=\infty\land{}s'.I(t)\le{}lower(I_s(t))\Rightarrow{}s'.I(t)=\sigma'(id_t)("s\_time\_counter")$}\\
    
    \noindent{}Assuming $upper(I_s(t))=\infty$ and
    $s'.I(t)\le{}lower(I_s(t))$, let us show\\
    \fbox{$s'.I(t)=\sigma'(id_t)("s\_time\_counter")$.}

    Case analysis on $t\in{}Sens(s.M)$; there are two cases:

    \begin{enumerate}
    \item $t\notin{}Sens(s.M)$:\\
      By definition of
      $\gamma,E_c,\tau\vdash{}s\stackrel{\downarrow}{\sim}\sigma$:
      \begin{equation}
        \label{eq:50}\sigma(id_t)("se")=\mathtt{false}
      \end{equation}

      Thanks to \eqref{eq:49} and \eqref{eq:50}:
      \begin{equation}
        \label{eq:45}\sigma'(id_t)("stc")=0
      \end{equation}

      By definition of $E_c,\tau\vdash{}s\xrightarrow{\downarrow}s'$:
      \begin{equation}
        \label{eq:46}s'.I(t)=0
      \end{equation}

      Rewriting the goal with \eqref{eq:45} and \eqref{eq:46}:
      \qedbox{tautology.}
      
    \item $t\in{}Sens(s.M)$:\\
      By definition of $\gamma,E_c,\tau\vdash{}s\stackrel{\uparrow}{\sim}\sigma$:
      \begin{equation}
        \label{eq:51}\sigma(id_t)("se")=\mathtt{true}
      \end{equation}

      By construction, and as $upper(I_s(t))=\infty$,
      ${<}\mathtt{transition\_type\Rightarrow{}TEMP\_A\_INF}{>}\in{}gm_t$. By
      property of the elaboration relation:
      \begin{equation}
        \label{eq:52}\Delta(id_t)("tt")=\mathtt{TEMP\_A\_INF}
      \end{equation}

      Case analysis on $s.reset_t(t)$; there are two cases:
      \begin{enumerate}
      \item $s.reset_t(t)=\mathtt{true}$:\\
        By definition of
        $\gamma,E_c,\tau\vdash{}s\stackrel{\uparrow}{\sim}\sigma$:
        \begin{equation}
          \label{eq:53}\sigma(id_t)("srtc")=\mathtt{true}
        \end{equation}

        Thanks to \eqref{eq:48}, \eqref{eq:51}, \eqref{eq:52} and
        \eqref{eq:53}:
        \begin{equation}
          \label{eq:54}\sigma'(id_t)("stc")=1
        \end{equation}

        By definition of $E_c,\tau\vdash{}s\xrightarrow{\downarrow}s'$:
        \begin{equation}
          \label{eq:55}s'.I(t)=1
        \end{equation}

        Rewriting the goal with \eqref{eq:54} and \eqref{eq:55}:
        \qedbox{tautology.}

      \item $s.reset_t(t)=\mathtt{false}$:
        By definition of
        $\gamma,E_c,\tau\vdash{}s\stackrel{\uparrow}{\sim}\sigma$:
        \begin{equation}\label{eq:56}
          \sigma(id_t)("srtc")=\mathtt{false}
        \end{equation}

        As $upper(I_s(t))=\infty$, there exists an
        $a\in\mathbb{N}^{*}$ s.t. $I_s(t)=[a,\infty]$. Let us take
        such an $a\in\mathbb{N}^{*}$. By construction,
        ${<}\mathtt{maximal\_time\_counter\Rightarrow}a{>}\in{}gm_t$,
        and by property of the elaboration relation:
        \begin{equation}
          \label{eq:57}\Delta(id_t)("mtc")=a
        \end{equation}

        By definition of
        $E_c,\tau\vdash{}s\xrightarrow{\downarrow}s'$, and knowing
        that $t\in{}Sens(s.M)$ and $s.reset_t(t)=\mathtt{false}$ and
        $upper(I_s(t))=\infty$:
        \begin{equation}
          \label{eq:58}s'.I(t)=s.I(t)+1
        \end{equation}

        Rewriting the goal with \eqref{eq:58}:
        \fbox{$s.I(t)+1=\sigma'(id_t)("stc")$.}
        
        We assumed that $s'.I(t)\le{}lower(I_s(t))$, and as
        $s'.I(t)=s.I(t)+1$, then $s.I(t)+1\le{}lower(I_s(t))$, then
        $s.I(t)<lower(I_s(t))$, then $s.I(t)<a$ since
        $a=lower(I_s(t))$.

        \noindent{}By definition of
        $\gamma,E_c,\tau\vdash{}s\stackrel{\uparrow}{\sim}\sigma$, and
        knowing that $s.I(t)<lower(I_s(t))$ and $upper(I_s(t))=\infty$:
        \begin{equation}
          \label{eq:59}s.I(t)=\sigma(id_t)("stc")
        \end{equation}

        Thanks to \eqref{eq:57}, \eqref{eq:59} and $s.I(t)<a$:
        \begin{equation}
          \label{eq:60}\sigma(id_t)("stc")<\Delta(id_t)("mtc")
        \end{equation}
        
        Thanks to \eqref{eq:44}, \eqref{eq:60}, \eqref{eq:56} and
        \eqref{eq:51}:
        \begin{equation}
          \label{eq:61}\sigma'(id_t)("stc")=\sigma(id_t)("stc")+1
        \end{equation}

        Rewriting the goal with \eqref{eq:61} and \eqref{eq:59}: \qedbox{tautology.}        
      \end{enumerate}
    \end{enumerate}
    
  \item
    \fbox{$upper(I_s(t))=\infty\land{}s'.I(t)>{}lower(I_s(t))\Rightarrow$
      $\sigma'(id_t)("s\_time\_counter")=lower(I_s(t)$.}

    Assuming that $upper(I_s(t))=\infty$ and
    $s'.I(t)>{}lower(I_s(t))$, let us show\\
    \fbox{$\sigma'(id_t)("s\_time\_counter")=lower(I_s(t))$.}

    As $upper(I_s(t))=\infty$, there exists an $a\in\mathbb{N}^{*}$
    s.t. $I_s(t)=[a,\infty]$. Let us take such an
    $a\in\mathbb{N}^{*}$. By construction,
    ${<}\mathtt{maximal\_time\_counter\Rightarrow}a{>}\in{}gm_t$, and
    ${<}\mathtt{transition\_type\Rightarrow{}TEMP\_A\_INF}{>}\in{}gm_t$
    by property of the elaboration relation:
    \begin{eqnarray}
      \label{eq:62}\Delta(id_t)("mtc")&=&a\\
      \label{eq:63}\Delta(id_t)("tt")&=&\mathtt{TEMP\_A\_INF}
    \end{eqnarray}
    
    Case analysis on $t\in{}Sens(s.M)$:
    \begin{enumerate}
    \item $t\notin{}Sens(s.M)$:\\
      \noindent{}By definition of
      $E_c,\tau\vdash{}s\xrightarrow{\downarrow}s'$, and knowing that
      $t\in{}Sens(s.M)$, then $s'.I(t)=0$. Since
      $lower(I_s(t))\in\mathbb{N}^{*}$, then $lower(I_s(t))>0$.
      
      \qedbox{Contradicts $s'.I(t)>lower(I_s(t))$.}
      
    \item $t\in{}Sens(s.M)$:\\
      By definition of
      $\gamma,E_c,\tau\vdash{}s\stackrel{\uparrow}{\sim}\sigma$ and
      $t\in{}Sens(s.M)$:
      \begin{equation}
        \label{eq:64}\sigma(id_t)("se")=\mathtt{true}
      \end{equation}

      Case analysis on $s.reset_t(t)$; there are two cases:
      \begin{enumerate}
      \item $s.reset_t(t)=\mathtt{true}$:\\
        By definition of
        $E_c,\tau\vdash{}s\xrightarrow{\downarrow}s'$: $s'.I(t)=1$.

        \noindent{}We assumed that $s'.I(t)>lower(I_s(t))$, then
        $1>lower(I_s(t))$.

        \qedbox{Contradicts $lower(I_s(t))>0$.}

      \item $s.reset_t(t)=\mathtt{false}$:\\
        \noindent{}By property of
        $\gamma,E_c,\tau\vdash{}s\stackrel{\uparrow}{\sim}\sigma$ and
        $s.reset_t(t)=\mathtt{false}$:
        \begin{equation}
          \sigma(id_t)("srtc")=\mathtt{false}\label{eq:68}
        \end{equation}

        By definition of
        $E_c,\tau\vdash{}s\xrightarrow{\downarrow}s'$, and knowing
        that $s'.I(t)>lower(I_s(t))$:
        \begin{equation}
          \begin{split}
            s'.I(t)=s.I(t)+1&\Rightarrow{}s.I(t)+1>lower(I_s(t))\\
            & \Rightarrow{}s.I(t)\ge{}lower(I_s(t))\\
          \end{split}
        \end{equation}

        Case analysis on $s.I(t)\ge{}lower(I_s(t))$:
        \begin{enumerate}
        \item $s.I(t)>lower(I_s(t))$: \fbox{$\sigma'(id_t)("stc")=lower(I_s(t))$.}\\
          By definition of
          $\gamma,E_c,\tau\vdash{}s\stackrel{\uparrow}{\sim}\sigma$:
          \begin{equation}
            \label{eq:69}\sigma(id_t)("stc")=lower(I_s(t))
          \end{equation}

          Thanks to \eqref{eq:47}:
          \begin{equation}
            \label{eq:70}\sigma'(id_t)("stc")=\sigma(id_t)("stc")
          \end{equation}

          Rewriting the goal with \eqref{eq:69} and \eqref{eq:70}:
          \qedbox{tautology.}
          
        \item $s.I(t)=lower(I_s(t))$: \fbox{$\sigma'(id_t)("stc")=lower(I_s(t))$.}\\
          By definition of
          $\gamma,E_c,\tau\vdash{}s\stackrel{\uparrow}{\sim}\sigma$:
          \begin{equation}
            \label{eq:65}s.I(t)=\sigma(id_t)("stc")
          \end{equation}

          Thanks to \eqref{eq:47}:
          \begin{equation}
            \label{eq:66}\sigma'(id_t)("stc")=\sigma(id_t)("stc")
          \end{equation}

          Rewriting the goal with \eqref{eq:66}, \eqref{eq:65} and
          $s.I(t)=lower(I_s(t))$: \qedbox{tautology.}
        \end{enumerate}
      \end{enumerate}
    \end{enumerate}
  \item
    \fbox{$upper(I_s(t))\neq\infty\land{}s'.I(t)>{}upper(I_s(t))\Rightarrow$
      $\sigma'(id_t)("s\_time\_counter")=upper(I_s(t))$.}

    Assuming that $upper(I_s(t))\neq\infty$ and
    $s'.I(t)>{}upper(I_s(t))$, let us show \\
    \fbox{$\sigma'(id_t)("s\_time\_counter")=upper(I_s(t))$.}

    As $upper(I_s(t))\neq\infty$, there exists an
    $a\in\mathbb{N}^{*}$, and a $b\in\mathbb{N}^{*}$
    s.t. $I_s(t)=[a,b]$. Let us take such an $a$ and $b$. By
    construction, there exists
    $tt\in\{\mathtt{TEMP\_A\_A,TEMP\_A\_B\}$ s.t.
      ${<}\mathtt{maximal\_time\_counter\Rightarrow}b{>}\in{}gm_t$,
      and ${<}\mathtt{transition\_type\Rightarrow{}}tt{>}\in{}gm_t$;
      by property of the elaboration relation:
    \begin{eqnarray}
      \label{eq:67}\Delta(id_t)("mtc")&=&b\\
      \label{eq:71}\Delta(id_t)("tt")&\neq&\mathtt{NOT\_TEMP}
    \end{eqnarray}

    Case analysis on $t\in{}Sens(s.M)$:
    \begin{enumerate}
    \item $t\notin{}Sens(s.M)$:\\
      \noindent{}By definition of
      $E_c,\tau\vdash{}s\xrightarrow{\downarrow}s'$, and knowing that
      $t\in{}Sens(s.M)$, then $s'.I(t)=0$. Since
      $upper(I_s(t))\in\mathbb{N}^{*}$, then $upper(I_s(t))>0$.
      
      \qedbox{Contradicts $s'.I(t)>upper(I_s(t))$.}
      
    \item $t\in{}Sens(s.M)$:\\
      By definition of
      $\gamma,E_c,\tau\vdash{}s\stackrel{\uparrow}{\sim}\sigma$ and
      $t\in{}Sens(s.M)$:
      \begin{equation}
        \label{eq:72}\sigma(id_t)("se")=\mathtt{true}
      \end{equation}

      Case analysis on $s.reset_t(t)$; there are two cases:
      \begin{enumerate}
      \item $s.reset_t(t)=\mathtt{true}$:\\
        By definition of
        $E_c,\tau\vdash{}s\xrightarrow{\downarrow}s'$: $s'.I(t)=1$.

        \noindent{}We assumed that $s'.I(t)>upper(I_s(t))$, then
        $1>upper(I_s(t))$.

        \qedbox{Contradicts $upper(I_s(t))>0$.}

      \item $s.reset_t(t)=\mathtt{false}$:\\
        \noindent{}By property of
        $\gamma,E_c,\tau\vdash{}s\stackrel{\uparrow}{\sim}\sigma$ and
        $s.reset_t(t)=\mathtt{false}$:
        \begin{equation}
          \sigma(id_t)("srtc")=\mathtt{false}\label{eq:73}
        \end{equation}
        
        Case analysis on $s.I(t)>{}upper(I_s(t))$ or
        $s.I(t)\le{}upper(I_s(t))$:
        \begin{enumerate}
        \item $s.I(t)>upper(I_s(t))$: \fbox{$\sigma'(id_t)("stc")=upper(I_s(t))$.}\\
          By definition of
          $E_c,\tau\vdash{}s\xrightarrow{\downarrow}s'$:
          \begin{equation}
            \label{eq:75}s'.I(t)=s.I(t)
          \end{equation}
          
          By definition of
          $\gamma,E_c,\tau\vdash{}s\stackrel{\uparrow}{\sim}\sigma$:
          \begin{equation}
            \label{eq:74}\sigma(id_t)("stc")=upper(I_s(t))
          \end{equation}

          Thanks to \eqref{eq:47}:
          \begin{equation}
            \label{eq:70}\sigma'(id_t)("stc")=\sigma(id_t)("stc")
          \end{equation}

          Rewriting the goal with \eqref{eq:70} and \eqref{eq:74}:
          \qedbox{tautology.}
          
        \item $s.I(t)\le{}upper(I_s(t))$: \fbox{$\sigma'(id_t)("stc")=upper(I_s(t))$.}\\
          Case analysis on $s.I(t)\le{}upper(I_s(t))$; there are two
          cases:
          \begin{enumerate}
          \item $s.I(t)=upper(I_s(t))$:\\
          \item $s.I(t)<upper(I_s(t))$:\\
          \end{enumerate}
        \end{enumerate}
      \end{enumerate}
    \end{enumerate}
    
    
  \item
    \fbox{$upper(I_s(t))\neq\infty\land{}s'.I(t)\le{}upper(I_s(t))\Rightarrow$
      $s'.I(t)=\sigma'(id_t)("s\_time\_counter")$.}
  \end{enumerate}
\end{proof}

\subsection{Falling edge and firable transitions}
\label{sec:fe-firable}

%%%%%%%%%%%%%%%%%%%%%%%%%%%%%%%%%%%%%%%%%%%%%%%%
%%%%%%%%%% FALLING EDGE EQUAL FIRABLE %%%%%%%%%%
%%%%%%%%%%%%%%%%%%%%%%%%%%%%%%%%%%%%%%%%%%%%%%%%

\begin{lemma}[Falling Edge Equal Firable]
  \label{lem:fe-equal-firable}
  \fehyps{} and $\forall{}t,id_t$ s.t. $\gamma(t)=id_t$ and
  $\sigma'(id_t)=\sigma'_t$, then
  $t\in{}Firable(s')\Leftrightarrow\sigma'_t("s\_firable")=\mathtt{true}$.
\end{lemma}

\begin{proof}
  
\end{proof}

%%% Local Variables:
%%% mode: latex
%%% TeX-master: "../../main"
%%% End:


%%% Local Variables:
%%% mode: latex
%%% TeX-master: "../main"
%%% End:

% \chapter{Conclusion}
\label{chap:concl}

% \section{Contributions}
% \label{sec:contribs}

% \begin{itemize}
% \item Contributions to SITPN
% \item Contributions to VHDL
% \item Contributions to transformation functions
% \item Proof
%   \begin{itemize}
%   \item Bug detections + changes in semantics
%   \item How far are we to complete the verification?
%   \end{itemize}
% \end{itemize}

% \section{Improvements and perspectives}
% \label{sec:improvs-and-perspect}

% \pnote{if proof not yet verified, first thing is not complete this job!}

% \begin{itemize}
% \item improvements on the implementation of \hvhdl{}, get closer to
%   the formal def., more use of dependent types
% \item proving the semantic preservation thm in its existential version, i.e:
%   \begin{itemize}
%   \item the generated design is always elaborable
%   \item the generated design is always simulable, i.e, it will never
%     produce errors as long as the input model verifies some properties
%     (liveness, boundedness\dots)
%   \end{itemize}
% \item consider the whole \hilecop{} high-level model
%   \begin{itemize}
%   \item macroplaces
%   \item GALS
%   \end{itemize}
% \end{itemize}

In this thesis, we were interested in the formal verification of the
\hilecop{} methodology. The \hilecop{} methodology has been devised to
assist the design and the production of safety-critical digital
systems. To summarize, the \hilecop{} methodology proposes a
high-level graphical modelling formalism; the aim of the formalism is
to provide the engineers with a framework to model safety-critical
digital systems in a way that foster the communication around the
design models (hence the graphical aspect of the formalism). The
formalism is based on component diagrams and the internal behavior of
components are described with a particular kind of Petri nets. The
mathematical foundations of the Petri net formalism provides the
possibility to formally analyze the \hilecop{} high-level models, and
thereby to bring the proof that the models verify certain
\textit{soundness} properties.  The high-level formalism of \hilecop{}
has been designed to be handy for the humans; however, the ultimate
goal of the methodology is to obtain a physical version of the
safety-critical digital system on an FPGA or ASIC. Thus, from the
state of high-level model to its concrete implementation as a hardware
circuit, the designed digital system goes through multiple
transformations.  In this thesis, we considered the formal
verification of one of these transformations: the
\textit{model-to-text} transformation from a flattened Petri net
version of the high-level model of the safety-critical digital system
(i.e. an SITPN) to its implementation into a \vhdl{} design. This
transformation is performed by a computer program. It was our purpose
to formally verify that the transformation is \textit{semantic
  preserving}; that is, for any input model of the transformation the
corresponding output model behaves in the same way as the input
model. Pragmatically, the research question that we formalized in the
introduction of this thesis was:

\begin{center}
  \textsc{Can we prove that the model-to-text transformation described
    in the \hilecop{} methodology is semantic preserving?}
\end{center}

As pointed out in the literature reviews at the beginning of
Chapters~\ref{chap:transformation} and \ref{chap:proof}, the task of
formally verifying that a transformation from a source representation
to a target one is semantic preserving has been thoroughly studied,
and in different application contexts (generic compilers,
domain-specific compilers, model transformations, etc.). From this fact
arises the will to compare the \hilecop{} model-to-text transformation
to the other kinds of transformation found in the literature. In a
research point of view, the complementary questions that gravitated
around our main research question were:

\begin{itemize}
\item What are the similarities and the differences between the
  \hilecop{} model-to-text transformation and the other kind of
  transformations that have been formally verified?
\item Are there standard technics to prove that a transformation is
  semantic preserving? Do these technics apply in the case of the
  \hilecop{} model-to-text transformation?
\end{itemize}

In other words, what elects the formal verification of the \hilecop{}
model-to-text transformation as a concrete research task, and not as
another yet interesting engineering challenge?

In this thesis, the verification of the \hilecop{} transformation has
been conducted with the help of the \coq{} proof assistant; thus, the
relation between our formalization choices and the engineering
difficulties that they brought was one topic of interest.  Especially
in the world of formal verification, and more truly in the domain of
deductive verification and interactive proving, the mechanization of
proofs with the help of a proof assistant may be very time
consuming. We believe that it is a part of the research task to
evaluate the feasibility of the mechanization of proofs within a
reasonable time-span, and also, to try to bring an understanding
regarding the
formalization choices and their impacts on the mechanization.\\

To formally verify the semantic preservation property of the
\hilecop{} transformation, we followed the usual proceeding applied in
the domain of deductive verification, which is:
\begin{enumerate}
\item Formalize the execution semantics of the source representation
  (Chapter~\ref{chap:hilecop-models}) and the target representation
  (Chapter~\ref{chap:hvhdl}), and implement them using the proof
  assistant.
\item Describe and implement the transformation within the proof
  assistant (Chapter~\ref{chap:transformation}).
\item Express the theorem of semantic preservation, prove the theorem,
  and mechanize the proof using the chosen proof assistant
  (Chapter~\ref{chap:proof}).
\end{enumerate}

%%%%%%%% ANSWERS TO RESEARCH QUESTIONS AND CONTRIBUTIONS %%%%%%%%%

Following these steps, we brought the proof that the \hilecop{}
transformation is semantic preserving. Even though the mechanization
of the proof is not completed, each step of the approach brought its
own part of contributions.

During the formalization of the SITPN semantics, i.e. the source
representation of the transformation, we helped clarify the semantics
of these models, even though it was established in two previous theses
\cite{Leroux2014,Merzoug2018}. Especially, we formalized the concept
of a well-defined SITPN, that is, an \textit{acceptable} model for the
transformation. As a matter of fact, the \hilecop{} transformation
raises errors if the input SITPN model is not well-defined. Also, all
the theorems and lemmas that we proved rely on the well-definition
condition of the input SITPN model. Through the formalization of a
well-defined SITPN, we precisely characterized the way to handle a
conflict between the transitions of an SITPN. This aspect of the SITPN
semantics is complex and has been one particularly subtle point of the
proof of semantic preservation.

The reflection around the formal semantics of the \vhdl{} language and
how it could ease the reasoning around the \hilecop{}'s \vhdl{}
programs also brought new contributions.  From the rather complex
semantics of the \vhdl{}, and all is protean expressions found in the
literature, we devised a simple simulation algorithm, and formalized
it into a simulation relation for the execution of synchronous \vhdl{}
designs. We defined the abstract syntax of a subset of \vhdl{}, called
\hvhdl{}, that suited our needs regarding the \vhdl{} programs
generated by the transformation and the ones that were previously
defined by the methodology (i.e. the \texttt{place} and
\texttt{transition} designs). However, the \hvhdl{} syntax and
semantics is well-suited to express any kind of synchronous and
synthesizable digital systems. Moreover, the implementation of the
\hvhdl{} syntax and semantics in \coq{} provides a framework to reason
about \hvhdl{} designs. To the best of our knowledge, it is the only
example of implementation of the \vhdl{} language using the \coq{}
proof assistant.

About the expression of the \hilecop{} model-to-text transformation
itself, the contribution of this thesis was the design of the
algorithm of the transformation, which had never been expressed
before, and its implementation within the \coq{} proof assistant.  For
the implementation of the transformation, we tried as much as possible
to draw our inspirations from the technics found in the
literature. Especially, we tried to produce clear, maintainable code,
through the use of functional design patterns, while anticipating the
additions of new elements in the input models. We also tried as much
as possible to implement the transformation in order to ease the
mechanization of the proof of semantic preservation.

The last part of the work was related to the expression and the proof
of the semantic preservation theorem. While expressing the theorem of
semantic preservation, we formalized the way to compare an input SITPN
model with the corresponding \vhdl{} design, result of the
transformation. This point is the angular stone of the theorem of
semantic preservation, moreover considering that the gap between the
source and the target representation is substantial.  The major
contribution of the thesis to this part of the work is, of course, to
have brought the proof of the semantic preservation theorem, more
especially the proof of Theorem~\ref{thm:full-bisim}. Although the
mechanization of the proof is far from being completed, establishing
the informal proof that the \hilecop{} model-to-text transformation is
semantic preserving has revealed a bug in the \vhdl{} implementation
of the \texttt{place} and \texttt{transition} design
(cf. Section~\ref{sec:bug-detection}).

% Answers to additional research questions

We stated above that from our very pragmatic research question arose a
lot of additional questions. These questions pertain to the position
of the \hilecop{} model-to-text transformation with respect to the
other examples of transformation for which a work of formal
verification has been realized. In other words, what makes this
transformation specific? What aspects of this transformation and of
its formal verification motivate a research interest? The very context
of the design and production of critical digital systems, in which the
\hilecop{} transformation is involved, brings out interesting research
challenges. In terms of semantics, it means that we are dealing with
\textit{reactive} systems, i.e. systems which are characterized by a
time-related execution and their interactions with an outside
environment. Considering the work done on the formal verification of
compilers for GPLs, where programs are \textit{transformational}
systems (i.e. there is a one-time end-to-end execution of the
program), this already constitutes an originality. The reactivity of
systems must be taken into account in the expression of the theorem of
semantic preservation. However, some works
\cite{Bourke,Habibi2006,Braibant2013,Yang2014,Loow2021,Hansen2007a,Bourgeat2020,Yang2016}
have already been conducted on the formal verification of hardware
description language (HDL) compilers. In that case, the source
language and its semantics permit us to describe reactive systems. As
there exist a lot of works pertaining to the formal verification of
transformations relating a source programming \textit{language} to a
target one, the first originality of the \hilecop{} transformation is
that the source representation is a \textit{graphical formalism}. This
graphical formalism is a particular brand of Petri nets with a
synchronous semantics, which is also an original point as most of the
Petri net classes have an asynchronous semantics. These SITPNs have
been designed to give as much power of expression as possible to the
engineers that are designing safety-critical digital systems. Blending
these considerations with the context of formal methods, and the
necessity to produce \textit{safe} models of critical digital systems,
the result is that the semantics of SITPNs is rather complex;
especially the handling of conflicts between transitions.

Another point of originality of the \hilecop{} transformation comes
from the distance between the SITPN models, which are yet abstract
mathematical objects, and their target representation as \vhdl{}
designs, which already deeply tied to the functioning of hardware
systems. Moreover, two designs are at the base of the representation
of SITPN models into \vhdl{} programs: the \texttt{place} and
\texttt{transition} designs. These two designs represent more or less
of a hundred lines of \vhdl{} code each. The \vhdl{} code describing
the behavior of the \texttt{place} and \texttt{transition} designs
comes with a lot of implementation-related particularities that are
sometimes difficult to relate to the semantics of SITPNs (and
sometimes impossible to relate at all, hence the bug detection,
cf. Section~\ref{sec:bug-detection}).

Finally, the \hilecop{} transformation function is itself rather
complex. It cannot be expressed by rules following the inductive
structure of the abstract syntax of a source programming language, as
it is usually done in compilers. Specifically because of the nature of
the SITPN structure, the \hilecop{} transformation has to cover a lot
of particular cases related to the form of the input models. The
specificities of the \hilecop{} transformation function relate to the
difficulties that we encountered to mechanize the proof of semantic
preservation.

%%%%%% LIMITS OF THE RESEARCH %%%%%%

Although the proof of semantic preservation has been established as a
\textit{semi-formal} paper proof, we were not able to fully mechanize
it within the \coq{} proof assistant; at least not in the time span of
the thesis. This has brought a lot of thinking about the reasons
surrounding the difficulties of the mechanization, and also about the
solutions that would allow us to go over these
difficulties. Specifically, we were wondering if the mechanization
could not be brought out entirely because of a lack of engineering
skills or because of other considerations. These considerations
included:

\begin{itemize}
\item The complexity of the \hvhdl{} semantics: during the
  mechanization of the proof, we realized that the \hvhdl{} semantics,
  and especially the rules related to the simulation loop, although
  much more simplified than the complete simulation loop of the
  \vhdl{} LRM, was not convenient to reason about the \vhdl{} designs
  generated by the transformation, nor to reason about the
  \texttt{place} and \texttt{transition} design behaviors. Therefore,
  some changes have been made in the \hvhdl{} semantics and the final
  result has been presented in Chapter~\ref{chap:hvhdl}. However, at
  the moment of the writing, we have not yet measured the impact of
  these changes on the mechanization of the proof.

\item The complexity of the source models: one solution to be able to
  complete the mechanization could have been to consider an even
  smaller subset for the source PNs models. For instance, we could
  have let aside the time and interpretation aspects in SITPNs.
  
\item The implementation of the transformation function: the current
  implementation of the transformation function corresponds to the
  implementation of a former version of the transformation
  algorithm. A new and simpler version of the transformation algorithm
  has been presented in
  Chapter~\ref{chap:transformation}. % Implementing the transformation
  % function in \coq{} as close as possible to the algorithm could
  % facilitate the proofs regarding the properties of the
  % transformation.
  The current implementation of the transformation includes some
  intermediate steps, between the input model and the final \hvhdl{}
  design, that might not be necessary and add further complexities at
  the time of proofs.
  
\item The bootstrap cost of the mechanization task: at the beginning
  of the mechanization, a lot of intermediary lemmas must be proved
  that will later be extensively used in other proofs. The consequence
  is that the overall completion of the mechanization advances very
  slowly at the beginning of the project because a lot of little
  bricks must be set. Eventually, the verification goes much faster
  when all the necessary tools are in place (notably thanks to the
  \coqeb{auto} tactic of the \coq{} proof assistant, and the hint
  databases system).
\end{itemize}

Pondering all these considerations, it remains clear that the
\hilecop{} methodology is an industrial product with all its
subtleties. Our guess is that, to complete the mechanization of the
proof, it will require one person (already acquainted with the overall
system) doing the job at full time during one year.  However, we are
confident in the fact that we cleared the way enough for the proof of
semantic preservation to be fully mechanized; now, it is only a
question of human and time resources to complete it.

Once the mechanization of the proof will be completed, we will have
the formal proof that the \hilecop{} model-to-text transformation is
semantic preserving. Then, this formal proof can help to certify the
\hilecop{} methodology as an eligible methodology for the design and
production of safety-critical digital systems.  The Neurinnov company
exploits the \hilecop{} methodology for the production of
neuroprostheses, which are considered highly critical medical
devices. To certify the neuroprostheses as eligible for market, the UE
regulation ask for the thorough testing of all the programs involved
in the production chain, thus including the \hilecop{} methodology.
For the moment, the UE regulation standards for medical devices do not
consider a formal proof as sufficient to certify a given program.  All
the standards in the domain of avionics, railways, power plants and
many others consider a formal proof as sufficient. Therefore, we are
confident in the fact that the UE regulation standards for medical
devices will soon evolve in this way.

%%%%% PERSPECTIVES %%%%%

\section{Future work and perspectives}
\label{sec:future}

In the immediate future, the first work to complete is of course the
mechanization of the proof of Theorem~\ref{thm:full-bisim}. Then, the
proofs of Theorems~\ref{thm:elab-ex}, \ref{thm:init-ex} and
\ref{thm:sim-ex}, which are actually considered as axioms, must be
completed as well. Finally, we must take into account all the aspects
of the \hilecop{} high-level models, which actually have a richer
structure than the one presented in Chapter~\ref{chap:hilecop-models}.

The first aspect to reconsider in the definition of the SITPN
structure is \textit{interpretation}. The interpretation aspect has
been simplified in the actual version of SITPNs. We could at least
consider the set of \vhdl{} signals, i.e. the variables of the
interpretation, as being a part of the SITPN structure. Depending on
the precision level we want to achieve, we can relate the conditions
to concrete Boolean expressions, and the actions/functions to concrete
operations. Moreover, we can equip the SITPN semantics with an
operational semantics to evaluate the expressions and operations
performed over \vhdl{} signals. In that case, the SITPN state would
include a mapping between the set of \vhdl{} signals and their current
value.

The two other main aspects still to be integrated are
\textit{macro\-places} and \textit{multi-clock domains}.  The
macroplace mechanism, illustrated in Figure~\ref{fig:macroplace},
enables the encapsulation of an SITPN subnet, called a refinement,
into an environment that handles exceptions.

\begin{figure}[H]
  \centering
  \includegraphics[keepaspectratio,width=.6\textwidth]{Figures/Conclusion/macroplace}
  \caption[An SITPN model with a macroplace.]{The macroplace is the
    double-lined circle that encapsulates an SITPN subnet; the subnet
    is called a \textit{refinement}. The arcs that enter and go out of
    the macroplace are particular arcs, thus with a particular
    semantics, represented by dotted arrows.}
  \label{fig:macroplace}
\end{figure}

The formal definition of the SITPN structure with macroplaces and its
formal semantics have been described in \cite{Leroux2014}. Adding
macroplaces to the actual SITPN structure will impact the
transformation function, and all the surrounding proofs. It will also
bring a new \hvhdl{} design (i.e. the one defining macroplace
components) in the loop, and will modify the code of the
\texttt{place} and \texttt{transition} designs.

In the actual semantics of SITPNs, we considered that only one clock
signal regulated the evolution of the system. However, the formalism
of the \hilecop{} high-level models includes the possibility to assign
different clock \textit{domains} to different parts of the same input
model. Thus, the modeled system is qualified as a multi-clock domain
system. It means that the different parts of the system are not
evolving at the same pace. Therefore, a mechanism of
\textit{asynchronous} message sending relates two parts with different
clock domains, and allows these parts to communicate together. The
system is said to have a Globally Asynchronous Locally Synchronous
(GALS) architecture.  The semantics of SITPNs that integrate
multi-clock domains has not been formalized yet. The multi-clock
domain aspect also implies modifying the \hvhdl{} semantics to
integrate multiple clock signals in the simulation loop.

The \coq{} proof assistant provides a way to extract \ocaml{} or
\textsf{Haskell} code from a \coq{} function. Thus, proving that our
\coq{} implementation of the \hilecop{} model-to-text transformation
is semantic preserving would allow us to extract a sound \ocaml{} or
\textsf{Haskell} program out of it. This program could then replace
the existing \textsf{Java} implementation of the \hilecop{}
methodology. At least, the engineers in charge of maintaining the
current \hilecop{} implementation are open to the idea. However, in
the certification standards of safety-critical software programs, the
executed code may not be the one over which the formal verification
has been conducted (e.g. EAL 7 in the Common Criterion standard).

All the programming tasks of this thesis have been performed within
the framework of the \coq{} proof assistant. The produced code is
fully accessible under the following \textsf{Git} repository:

\begin{center}
  \url{https://github.com/viampietro/ver-hilecop}
\end{center}

%%% Local Variables:
%%% mode: latex
%%% TeX-master: "../main"
%%% End:
 

%----------------------------------------------------------------------------------------
%	THESIS CONTENT - APPENDICES
%----------------------------------------------------------------------------------------

\appendix % Cue to tell LaTeX that the following "chapters" are Appendices

% Include the appendices of the thesis as separate files from the Appendices folder
% Uncomment the lines as you write the Appendices

% \chapter{The place design in concrete and abstract VHDL syntax}
\label{app:place-design}

\begin{lstlisting}[language=VHDL,label={lst:place-design},caption={The place design in concrete VHDL syntax.},framexleftmargin=1.5em,xleftmargin=2em,numbers=left, numberstyle=\tiny\ttfamily]
entity place is
  generic(
    input_arcs_number    : natural := 1;
    output_arcs_number   : natural := 1;
    maximal_marking      : natural := 1
  );
  port(
    clock : in std_logic;
    reset_n : in std_logic;
    initial_marking : in natural range 0 to maximal_marking; 
    input_arcs_weights : in weight_vector_t(input_arcs_number-1 downto 0);
    output_arcs_types : in arc_vector_t(output_arcs_number-1 downto 0);
    output_arcs_weights : in weight_vector_t(output_arcs_number-1 downto 0);
    input_transitions_fired : in std_logic_vector(input_arcs_number-1 downto 0);          
    output_transitions_fired : in std_logic_vector(output_arcs_number-1 downto 0);
    output_arcs_valid : out std_logic_vector(output_arcs_number-1 downto 0);
    priority_authorizations : out std_logic_vector(output_arcs_number-1 downto 0);
    reinit_transitions_time : out std_logic_vector(output_arcs_number-1 downto 0);
    marked : out std_logic 
  );
end place;

architecture place_architecture of place is

  subtype local_weight_t is natural range 0 to maximal_marking;

  signal s_input_token_sum       : local_weight_t;
  signal s_marking               : local_weight_t;
  signal s_output_token_sum      : local_weight_t; 

begin

  
  input_tokens_sum : process(input_arcs_weights, input_transitions_fired)
    variable v_internal_input_token_sum : local_weight_t;
  begin
    v_internal_input_token_sum := 0;
    
    for i in 0 to input_arcs_number - 1 loop
      if (input_transitions_fired(i) = '1') then
        
        v_internal_input_token_sum := v_internal_input_token_sum + input_arcs_weights(i);
        
      end if;
    end loop;
    
    s_input_token_sum <= v_internal_input_token_sum;
  end process input_tokens_sum;

  output_tokens_sum : process(output_arcs_types, output_arcs_weights, output_transitions_fired)
    variable v_internal_output_token_sum : local_weight_t;
  begin
    v_internal_output_token_sum := 0;
    
    for i in 0 to output_arcs_number - 1 loop
      if (output_transitions_fired(i) = '1' and output_arcs_types(i) = arc_t(BASIC)) then
        v_internal_output_token_sum := v_internal_output_token_sum + output_arcs_weights(i); 
      end if;
    end loop;
    
    s_output_token_sum <= v_internal_output_token_sum;
  end process output_tokens_sum;
    
  marking : process(clock, reset_n, initial_marking)
  begin
    if (reset_n = '0') then
      s_marking <= initial_marking;
    elsif rising_edge(clock) then
      s_marking <= s_marking + (s_input_token_sum - s_output_token_sum);
    end if;
  end process marking;

  determine_marked : process(s_marking)
  begin
    if (s_marking = 0) then
      marked <= '0';
    else
      marked <= '1';
    end if;
  end process determine_marked;        
  
  marking_validation_evaluation : process(output_arcs_types, output_arcs_weights, s_marking)
  begin
    for i in 0 to output_arcs_number - 1 loop
      if ( ( ((output_arcs_types(i) = arc_t(BASIC)) or (output_arcs_types(i) = arc_t(TEST))) and (s_marking >= output_arcs_weights(i)) )
           or ( (output_arcs_types(i) = arc_t(INHIBITOR)) and (s_marking < output_arcs_weights(i)) ) )
      then
        output_arcs_valid(i) <= '1';
      else
        output_arcs_valid(i) <= '0';
      end if;
    end loop;
  end process marking_validation_evaluation;

  priority_evaluation : process(output_arcs_types, output_arcs_weights, output_transitions_fired, s_marking)
    variable v_saved_output_token_sum : local_weight_t;
  begin
    v_saved_output_token_sum := 0;
    
    for i in 0 to output_arcs_number - 1 loop
      if (s_marking >= v_saved_output_token_sum + output_arcs_weights(i)) then
        priority_authorizations(i) <= '1';
      else
        priority_authorizations(i) <= '0';
      end if;
      
      if ((output_transitions_fired(i) = '1') and (output_arcs_types(i) = arc_t(BASIC))) then
        v_saved_output_token_sum := v_saved_output_token_sum + output_arcs_weights(i);
      end if;
      
    end loop;
  end process priority_evaluation;

  reinit_transitions_time_evaluation : process(clock, reset_n)
  begin
    if (reset_n = '0') then
      reinit_transitions_time <= (others => '0');
    elsif rising_edge(clock) then
      for i in 0 to output_arcs_number - 1 loop
        if ( (((output_arcs_types(i) = arc_t(BASIC)) or (output_arcs_types(i) = arc_t(TEST)))
                and (s_marking - s_output_token_sum < output_arcs_weights(i))
                and (s_output_token_sum > 0))
             or output_transitions_fired(i) = '1' )
        then
          reinit_transitions_time(i) <= '1';
        else
          reinit_transitions_time(i) <= '0';
        end if;
      end loop;
    end if;
  end process reinit_transitions_time_evaluation;
     
end place_architecture;

\end{lstlisting}

%%% Local Variables:
%%% mode: latex
%%% TeX-master: "../main"
%%% End:

% \chapter{The transition design in concrete and abstract VHDL syntax}
\label{app:trans-design}

\begin{lstlisting}[language=VHDL,label={lst:trans-design},caption={The transition design in concrete VHDL syntax.},framexleftmargin=1.5em,xleftmargin=2em,numbers=left, numberstyle=\tiny\ttfamily]
entity transition is
  generic(
    transition_type      : transition_t := NOT_TEMPORAL;
    input_arcs_number    : natural := 1; 
    conditions_number    : natural := 1;
    maximal_time_counter : natural := 1
    );        
  port(
    clock                   : in std_logic;
    reset_n                 : in std_logic;
    input_conditions        : in std_logic_vector(conditions_number-1 downto 0);
    time_A_value            : in natural range 0 to maximal_time_counter;
    time_B_value            : in natural range 0 to maximal_time_counter;
    input_arcs_valid        : in std_logic_vector(input_arcs_number-1 downto 0);
    reinit_time             : in std_logic_vector(input_arcs_number-1 downto 0);
    priority_authorizations : in std_logic_vector(input_arcs_number-1 downto 0);
    fired                   : out std_logic 
    );
end transition;

architecture transition_architecture of transition is

  signal s_condition_combination : std_logic; 
  signal s_enabled               : std_logic;
  signal s_firable               : std_logic;   
  signal s_firing_condition      : std_logic;
  signal s_priority_combination  : std_logic;
  signal s_reinit_time_counter   : std_logic;
  signal s_time_counter          : natural range 0 to maximal_time_counter;

begin
  
  condition_evaluation : process(input_conditions)
    variable v_internal_condition : std_logic;
  begin
    v_internal_condition := '1';
    
    for i in 0 to conditions_number $-$ 1 loop
      v_internal_condition := v_internal_condition and input_conditions(i);
    end loop;
    
    s_condition_combination <= v_internal_condition;
  end process condition_evaluation;
  
  enable_evaluation : process(input_arcs_valid)
    variable v_internal_enabled : std_logic;
  begin
    v_internal_enabled := '1';
    
    for i in 0 to input_arcs_number $-$ 1 loop
      v_internal_enabled := v_internal_enabled and input_arcs_valid(i);
    end loop;
    
    s_enabled <= v_internal_enabled;
  end process enable_evaluation;

  reinit_time_counter_evaluation : process(reinit_time, s_enabled)
    variable v_internal_reinit_time_counter : std_logic;
  begin
    v_internal_reinit_time_counter := '0';

    for i in 0 to input_arcs_number $-$ 1 loop
      v_internal_reinit_time_counter := v_internal_reinit_time_counter or reinit_time(i);
    end loop;
    
    s_reinit_time_counter <= v_internal_reinit_time_counter;
  end process reinit_time_counter_evaluation;

  time_counter : process(reset_n, clock)
  begin
    if (reset_n = '0') then
      s_time_counter <= 0;
    elsif falling_edge(clock) then
      if ((s_enabled = '1') and (transition_type /= transition_t(NOT_TEMPORAL))) then
        if (s_reinit_time_counter = '0') then
          if (s_time_counter < maximal_time_counter) then
            s_time_counter <= s_time_counter + 1;
          end if;
        else
          s_time_counter <= 1; 
        end if;
      else
        s_time_counter <= 0;
      end if;
    end if;
  end process time_counter;

  firing_condition_evaluation : process (s_enabled, s_condition_combination, s_reinit_time_counter, s_time_counter)
  begin
    if ( (s_condition_combination = '1')
         and (s_enabled = '1')
         and ( (transition_type = transition_t(NOT_TEMPORAL))
               or ( (transition_type = transition_t(TEMPORAL_A_B))
                    and (s_reinit_time_counter = '0')
                    and (s_time_counter >= (time_A_value-1))
                    and (s_time_counter < time_B_value))
               or ( (s_reinit_time_counter = '0')
                    and ( ((transition_type = transition_t(TEMPORAL_A_A))
                          and (s_time_counter = (time_A_value-1)))
                        or ( (transition_type = transition_t(TEMPORAL_A_INFINITE))
                             and (s_time_counter >= (time_A_value-1)) ))
               or ( (transition_type /= transition_t(NOT_TEMPORAL))
                    and (s_reinit_time_counter = '1')
                    and (time_A_value = 1) )
               )
    ) then
      s_firing_condition <= '1';       
    else
      s_firing_condition <= '0';
    end if;
  end process firing_condition_evaluation;
  
  priority_authorization_evaluation : process(priority_authorizations)
    variable v_priority_combination : std_logic;
  begin
    v_priority_combination := '1';
    
    for i in 0 to input_arcs_number $-$ 1 loop
      v_priority_combination := v_priority_combination and priority_authorizations(i);
    end loop;
    
    s_priority_combination <= v_priority_combination;
  end process priority_authorization_evaluation;

  firable : process(reset_n, clock)
  begin
    if (reset_n = '0') then
      s_firable <= '0';
    elsif falling_edge(clock) then
      s_firable <= s_firing_condition;
    end if;
  end process firable;

  fired_evaluation : process (s_firable, s_priority_combination)
  begin
    fired <= s_firable and s_priority_combination;
  end process fired_evaluation;      
  
end transition_architecture;
  
\end{lstlisting}

\begin{lstlisting}[language=VHDL,label={lst:trans-design-abss},caption={The \texttt{transition} design in \hvhdl{} abstract syntax.},framexleftmargin=1.5em,xleftmargin=2em,numbers=left, numberstyle=\tiny\ttfamily]
design transition transition_architecture

-- Generic clause
((transition_type, natural(0,3), 0),
 (input_arcs_number, natural(0,NATMAX), 1),
 (conditions_number, natural(0,NATMAX), 1),
 (maximal_time_counter, natural(0,NATMAX),1))

-- Port clause
((in, input_conditions, array(boolean, 0, conditions_number$-$1),
 (in, time_A_value, natural(0, maximal_time_counter)),
 (in, time_B_value, natural(0, maximal_time_counter)),
 (in, input_arcs_valid, array(boolean, 0, input_arcs_number$-$1)),
 (in, reinit_time, array(boolean, 0, input_arcs_number$-$1)),
 (in, priority_authorizations, array(boolean, 0, input_arcs_number$-$1)),
 (out, fired, boolean))

-- Architecture declarative part
((s_condition_combination, boolean),
 (s_enabled, boolean),
 (s_firable, boolean),
 (s_firing, boolean),
 (s_priority, boolean),
 (s_reinit, boolean),
 (s_time_counter, natural(0, maximal_time_counter)),

-- Behavior

process (condition_evaluation, (input_conditions), ((v_internal_condition, boolean)),
(v_internal_condition := true;
 (for (i, 0, conditions_number$-$1)
      (v_internal_condition := v_internal_condition and input_conditions(i)));
 s_condition_combination <= v_internal_condition))
  
process (enable_evaluation, (input_arcs_valid), ((v_internal_enabled, boolean)),
(v_internal_enabled := true;
 (for (i, 0, input_arcs_number$-$1)
      (v_internal_enabled := v_internal_enabled and input_arcs_valid(i)));
 s_enabled <= v_internal_enabled))

process (reinit_time_counter_evaluation, (reinit_time, s_enabled),
         ((v_internal_reinit_time_counter, boolean)),
(v_internal_reinit_time_counter := false;
 (for (i, 0, input_arcs_number$-$1)
      (v_internal_reinit_time_counter := v_internal_reinit_time_counter or reinit_time(i)));
 (s_reinit_time_counter <= v_internal_reinit_time_counter)))

process (time_counter, (clk), $\emptyset$,
(rst 
 (s_time_counter <= 0)
 (falling
  (if ((s_enabled = true) and (transition_type $\neq$ 0)) 
      (if (s_reinit_time_counter = false)
          (if (s_time_counter < maximal_time_counter)
              (s_time_counter <= s_time_counter + 1))
          (s_time_counter <= 1))
      (s_time_counter <= 0)))))

process (firing_condition_evaluation, 
         (s_enabled, s_condition_combination, s_reinit_time_counter, s_time_counter),
         $\emptyset$,
(s_firing_condition <= 
 (s_condition_combination = true)
 and (s_enabled = true)
 and ( (transition_type = 0)
       or ((s_reinit_time_counter = false) and
            (((transition_type = 1) and (s_time_counter >= (time_A_value-1)) 
                                    and (s_time_counter < time_B_value))
              or ((transition_type = 2) and (s_time_counter = (time_A_value-1)))
              or ((transition_type = 3) and (s_time_counter >= (time_A_value-1)))))                              
       or ((s_reinit_time_counter = true)
           and (transition_type $\neq$ 0)
           and (time_A_value = 1)))))
  
process (priority_authorization_evaluation, (priority_authorizations),
         ((v_priority_combination, boolean)),
(v_priority_combination := true;
 (for (i, 0, input_arcs_number$-$1)
      (v_priority_combination := v_priority_combination and priority_authorizations(i)));
 s_priority_combination <= v_priority_combination))

process (firable, (clk), $\emptyset$,
(rst (s_firable <= false)
     (falling (s_firable <= s_firing_condition))))

process (fired_evaluation, (s_firable, s_priority_combination), $\emptyset$,
(fired <= s_firable and s_priority_combination))
\end{lstlisting}

%%% Local Variables:
%%% mode: latex
%%% TeX-master: "../main"
%%% End:

% \chapter{The semantic preservation theorem and its dependencies}
\label{app:tree-beh-pres-thm}

\begin{figure}[H]
  \centering
  \resizebox{\textwidth}{!}{%
    \begin{tikzpicture}
      \tikzset{level distance=200pt}
      \Tree [.\node {
        \begin{tabular}{@{}c@{}}
          Theorem~\ref{thm:beh-pres} \\
          (\nameref{thm:beh-pres}) \\
        \end{tabular}
      };
      [.\node {
        \begin{tabular}{@{}c@{}}
          Theorem~\ref{thm:elab-ex} \\
          (\nameref{thm:elab-ex}) \\
        \end{tabular}
      }; ]
      [.\node {
        \begin{tabular}{@{}c@{}}
          Theorem~\ref{thm:init-ex} \\
          (\nameref{thm:init-ex}) \\
        \end{tabular}
      }; ]
      [.\node {
        \begin{tabular}{@{}c@{}}
          Theorem~\ref{thm:sim-ex} \\
          (\nameref{thm:sim-ex}) \\
        \end{tabular}
      }; ]
      [.\node {
        \begin{tabular}{@{}c@{}}
          Theorem~\ref{thm:full-bisim} \\
          (\nameref{thm:full-bisim}) \\
        \end{tabular}
      }; ]
      ]
    \end{tikzpicture}
  }
\end{figure}
\begin{figure}[H]
  \centering
  \resizebox{\textwidth}{!}{%
    \begin{tikzpicture}[sibling distance=40pt]
      \tikzset{level distance=200pt}
      \Tree [.\node {
        \begin{tabular}{@{}c@{}}
          Theorem~\ref{thm:full-bisim} \\
          (\nameref{thm:full-bisim}) \\
        \end{tabular}
      };
      [.\node {
        \begin{tabular}{@{}c@{}}
          Lemma~\ref{lem:sim-init-states} \\
          (\nameref{lem:sim-init-states}) \\
        \end{tabular}
      }; ]
      [.\node {
        \begin{tabular}{@{}c@{}}
          Lemma~\ref{lem:fst-re} \\
          (\nameref{lem:fst-re}) \\
        \end{tabular}
      }; ]
      [.\node(thmfe) {
        \begin{tabular}{@{}c@{}}
          Lemma~\ref{lem:fe} \\
          (\nameref{lem:fe}) \\
        \end{tabular}
      }; ]
      [.\node(thmsim) {
        {\begin{tabular}{@{}c@{}}
          Lemma~\ref{lem:simulation} \\
          (\nameref{lem:simulation}) \\
         \end{tabular}}
      };
      \node {{\begin{tabular}{@{}c@{}}
          Lemma~\ref{lem:re} \\
          (\nameref{lem:re}) \\
         \end{tabular}}}; ] ]

      \draw (thmsim.west) edge[thick, ->] (thmfe.east);
    \end{tikzpicture}
  }
\end{figure}

\begin{figure}[H]
  \centering
  \resizebox{\textwidth}{!}{%
    \begin{tikzpicture}[scale=.5, sibling distance=40pt]
      \tikzset{level distance=400pt}
      \tikzset{grow'=right}
      \tikzset{execute at begin node=\strut}
      \tikzset{every tree node/.style={anchor=base west}}
      
      \Tree [.\node {
        \begin{tabular}{@{}c@{}}
          Lemma~\ref{lem:re} \\
          {\fontsize{30}{33}\selectfont (\nameref{lem:re}) } \\
        \end{tabular}
      };
      [.\node {
        \begin{tabular}{@{}c@{}}
          Lemma~\ref{lem:re-equal-marking} \\
          \fontsize{30}{33}\selectfont  (\nameref{lem:re-equal-marking}) \\
        \end{tabular}
      }; ]
      [.\node {
        \begin{tabular}{@{}c@{}}
          Lemma~\ref{lem:re-equal-tc} \\
          \fontsize{30}{33}\selectfont (\nameref{lem:re-equal-tc}) \\
        \end{tabular}
      }; ]
      [.\node {
        \begin{tabular}{@{}c@{}}
          Lemma~\ref{lem:re-equal-reset-orders} \\
          \fontsize{30}{33}\selectfont (\nameref{lem:re-equal-reset-orders}) \\
        \end{tabular}
      }; ]
      [.\node(thmsim) {
        {\begin{tabular}{@{}c@{}}
           Lemma~\ref{lem:re-equal-action-exec} \\
           \fontsize{30}{33}\selectfont (\nameref{lem:re-equal-action-exec}) \\
         \end{tabular}}
     }; ]
     [.\node {
       \begin{tabular}{@{}c@{}}
         Lemma~\ref{lem:re-equal-fun-exec} \\
         \fontsize{30}{33}\selectfont (\nameref{lem:re-equal-fun-exec}) \\
       \end{tabular}
     }; ]
     [.\node {
       \begin{tabular}{@{}c@{}}
         Lemma~\ref{lem:re-equal-sens} \\
         \fontsize{30}{33}\selectfont (\nameref{lem:re-equal-sens}) \\
       \end{tabular}
     }; ]
     [.\node {
      \begin{tabular}{@{}c@{}}
         Lemma~\ref{lem:re-equal-not-sens} \\
         \fontsize{30}{33}\selectfont (\nameref{lem:re-equal-not-sens}) \\
      \end{tabular}
    }; ]
    [.\node {
      \begin{tabular}{@{}c@{}}
         Lemma~\ref{lem:re-equal-cond-comb} \\
         \fontsize{30}{33}\selectfont (\nameref{lem:re-equal-cond-comb}) \\
      \end{tabular}
    }; ]
    [.\node {
      \begin{tabular}{@{}c@{}}
         Lemma~\ref{lem:re-equal-cond} \\
         \fontsize{30}{33}\selectfont (\nameref{lem:re-equal-cond}) \\
      \end{tabular}
    }; ] ]

  \end{tikzpicture}
}
\end{figure}


\begin{figure}[H]
  \centering
  \resizebox{\textwidth}{!}{%
    \begin{tikzpicture}[scale=.5, sibling distance=30pt]
      \tikzset{level distance=400pt}
      \tikzset{grow'=right}
      \tikzset{execute at begin node=\strut}
      \tikzset{every tree node/.style={anchor=base west}}
      
      \Tree [.\node {
        \begin{tabular}{@{}c@{}}
          Lemma~\ref{lem:fe} \\
          {\fontsize{30}{33}\selectfont (\nameref{lem:fe}) } \\
        \end{tabular}
      };
      [.\node {
        \begin{tabular}{@{}c@{}}
          Lemma~\ref{lem:fe-equal-marking} \\
          \fontsize{30}{33}\selectfont  (\nameref{lem:fe-equal-marking}) \\
        \end{tabular}
      }; ]
      [.\node {
        \begin{tabular}{@{}c@{}}
          Lemma~\ref{lem:fe-equal-tc} \\
          \fontsize{30}{33}\selectfont (\nameref{lem:fe-equal-tc}) \\
        \end{tabular}
      }; ]
      [.\node {
        \begin{tabular}{@{}c@{}}
          Lemma~\ref{lem:fe-equal-cond-values} \\
          \fontsize{30}{33}\selectfont (\nameref{lem:fe-equal-cond-values}) \\
        \end{tabular}
      }; ]
      [.\node(thmsim) {
        {\begin{tabular}{@{}c@{}}
           Lemma~\ref{lem:fe-equal-act-exec} \\
           \fontsize{30}{33}\selectfont (\nameref{lem:fe-equal-act-exec}) \\
         \end{tabular}}
     }; ]
     [.\node {
       \begin{tabular}{@{}c@{}}
         Lemma~\ref{lem:fe-equal-fun-exec} \\
         \fontsize{30}{33}\selectfont (\nameref{lem:fe-equal-fun-exec}) \\
       \end{tabular}
     }; ]
     [.\node {
       \begin{tabular}{@{}c@{}}
         Lemma~\ref{lem:fe-equal-firable} \\
         \fontsize{30}{33}\selectfont (\nameref{lem:fe-equal-firable}) \\
       \end{tabular}
     }; ]
     [.\node {
      \begin{tabular}{@{}c@{}}
         Lemma~\ref{lem:fe-equal-not-firable} \\
         \fontsize{30}{33}\selectfont (\nameref{lem:fe-equal-not-firable}) \\
      \end{tabular}
    }; ]
    [.\node {
      \begin{tabular}{@{}c@{}}
         Lemma~\ref{lem:fe-equal-fired} \\
         \fontsize{30}{33}\selectfont (\nameref{lem:fe-equal-fired}) \\
      \end{tabular}
    }; ]
    [.\node {
      \begin{tabular}{@{}c@{}}
         Lemma~\ref{lem:fe-equal-not-fired} \\
         \fontsize{30}{33}\selectfont (\nameref{lem:fe-equal-not-fired}) \\
      \end{tabular}
    }; ]
    [.\node {
      \begin{tabular}{@{}c@{}}
         Lemma~\ref{lem:fe-equal-ots} \\
         \fontsize{30}{33}\selectfont (\nameref{lem:fe-equal-ots}) \\
      \end{tabular}
    }; ]
    [.\node {
      \begin{tabular}{@{}c@{}}
         Lemma~\ref{lem:fe-equal-its} \\
         \fontsize{30}{33}\selectfont (\nameref{lem:fe-equal-its}) \\
      \end{tabular}
    }; ] ]

  \end{tikzpicture}
}
\end{figure}

%%% Local Variables:
%%% mode: latex
%%% TeX-master: "../main"
%%% End:

% \chapter{Semantic preservation proof}
\label{app:sem-preserv-proof}

% \begin{table}[H]
%   \begin{tabular}{|c|c|c|c|}
\begin{longtable}{|c|c|c|c|}
  \hline
  \multicolumn{4}{|c|}{\textbf{Constants and signals reference}} \\
  \hline
  \textit{Full name} & \textit{Alias} & \textit{Category} & \textit{Type} \\
  \hline
  \texttt{input\_arcs\_number} & \texttt{ian} & generic constant (T) & $\mathbb{N}$ \\
  \hline
  \texttt{transition\_type} & \texttt{tt} & generic constant (T) & $\{\mathtt{not\_temp},\mathtt{temp\_a\_b},$ \\
                     & & & $\mathtt{temp\_a\_a},\mathtt{temp\_a\_inf}\}$ \\
  \hline
  \texttt{conditions\_number} & \texttt{cn} & generic constant (T) & $\mathbb{N}$ \\
  \hline
  \texttt{maximal\_time\_counter} & \texttt{mtc} & generic constant (T) & $\mathbb{N}$ \\
  \hline
  \texttt{time\_A\_value} & \texttt{A} & input port (T) & $\mathbb{N}$ \\
  \hline
  \texttt{time\_B\_value} & \texttt{B} & input port (T) & $\mathbb{N}$ \\
  \hline
  \texttt{input\_conditions} & \texttt{ic} & input port (T) & array of $\mathbb{B}$ \\
  \hline
  \texttt{reinit\_time} & \texttt{rt} & input port (T) & array of $\mathbb{B}$ \\
  \hline
  \texttt{input\_arcs\_valid} & \texttt{iav} & input port (T) & array of $\mathbb{B}$ \\
  \hline
  \texttt{priority\_authorizations} & \texttt{pauths} & input port (T) & array of $\mathbb{B}$ \\
  \hline
  \texttt{fired} & \texttt{f} & output port (T) & $\mathbb{B}$ \\
  \hline
  \texttt{s\_condition\_combination} & \texttt{scc} & internal signal (T) & $\mathbb{B}$ \\
  \hline
  \texttt{s\_reinit\_time\_counter} & \texttt{srtc} & internal signal (T) & $\mathbb{B}$ \\
  \hline
  \texttt{s\_priority\_combination} & \texttt{spc} & internal signal (T) & $\mathbb{B}$ \\
  \hline
  \texttt{s\_firable} & \texttt{sfa} & internal signal (T) & $\mathbb{B}$ \\
  \hline
  \texttt{s\_enabled} & \texttt{se} & internal signal (T) & $\mathbb{B}$ \\
  \hline
  \texttt{s\_time\_counter} & \texttt{stc} & internal signal (T) & $\mathbb{N}$ \\
  \hline
  \texttt{s\_firing\_condition} & \texttt{sfc} & internal signal (T) & $\mathbb{B}$ \\
  \hline
  \hline
  \texttt{input\_arcs\_number} & \texttt{ian} & generic constant (P) & $\mathbb{N}$ \\
  \hline
  \texttt{output\_arcs\_number} & \texttt{oan} & generic constant (P) & $\mathbb{N}$ \\
  \hline
  \texttt{maximal\_marking} & \texttt{mm} & generic constant (P) & $\mathbb{N}$ \\
  \hline
  \texttt{initial\_marking} & \texttt{im} & input port (P) & $\mathbb{N}$ \\
  \hline
  \texttt{output\_arcs\_types} & \texttt{oat} & input port (P) & array of $\{\mathtt{basic},\mathtt{test},\mathtt{inhib}\}$ \\
  \hline
  \texttt{output\_arcs\_weights} & \texttt{oaw} & input port (P) & array of $\mathbb{N}$ \\
  \hline
  \texttt{output\_transitions\_fired} & \texttt{otf} & input port (P) & array of $\mathbb{B}$ \\
  \hline
  \texttt{input\_arcs\_weights} & \texttt{iaw} & input port (P) & array of $\mathbb{N}$ \\
  \hline
  \texttt{input\_transitions\_fired} & \texttt{itf} & input port (P) & array of $\mathbb{B}$ \\
  \hline
  \texttt{output\_transitions\_fired} & \texttt{otf} & output port (P) & array of $\mathbb{B}$ \\
  \hline
  \texttt{reinit\_transitions\_time} & \texttt{rtt} & output port (P) & array of $\mathbb{B}$ \\
  \hline
  \texttt{priority\_authorizations} & \texttt{pauths} & output port (P) & array of $\mathbb{B}$ \\
  \hline
  \texttt{marked} & \texttt{m} & output port (P) & $\mathbb{B}$ \\
  \hline
  \texttt{s\_marking} & \texttt{sm} & internal signal (P) & $\mathbb{N}$ \\
  \hline
  \texttt{s\_output\_token\_sum} & \texttt{sots} & internal signal (P) & $\mathbb{N}$ \\
  \hline
  \texttt{s\_input\_token\_sum} & \texttt{sits} & internal signal (P) & $\mathbb{N}$ \\
  \hline
  % \end{tabular}
  \caption[Constants and signals reference for the \hvhdl{}
  \texttt{transition} and \texttt{place} designs.]{Constants and
  signals reference for the \hvhdl{} transition and place
  designs. In the \textit{Category} column, T (resp. P) indicates a
  generic constant, input port, output port or internal signal
  defined in the \texttt{transition} (resp. \texttt{place})
  design. }
  \label{tab:consts-and-sigs-ref}
  % \end{table}
\end{longtable}

\section{Initial States}
\label{sec:init-states}
%%%%%%%%%%%%%%%%%%%%%%%%%%%%%%%%%%%%%%%%%%%%%%%
%%%%%% SIMILAR INITIAL STATES HYPOTHESES %%%%%%
%%%%%%%%%%%%%%%%%%%%%%%%%%%%%%%%%%%%%%%%%%%%%%%

\begin{definition}[Initial state hypotheses]
  \label{def:init-states-hyps}
  Given an $sitpn\in{}SITPN$, $b\in{}P\rightarrow\mathbb{N}$,
  $d\in{}design$, $\gamma\in{}WM(sitpn,$ $d)$,
  $\Delta\in{}ElDesign,\sigma_{e},\sigma_0\in{}\Sigma$, assume that:
  \begin{itemize}
  \item SITPN $sitpn$ is transformed into the design $d$ and yields
    the binder $\gamma$: $\lfloor{}sitpn\rfloor_b=(d,\gamma)$
  \item $\Delta$ is the elaborated version of $d$, $\sigma_e$ is the
    default state of $\Delta$, i.e. the state of $\Delta$ where all signals are initialized to their default value:\\
    $\mathcal{D}_\mathcal{H},\emptyset\vdash{}d\srarrow{elab}{\fontsize{6}{8}\selectfont}(\Delta,\sigma_{e})$
    
  \item $\sigma_0$ is the initial state of $\Delta$: 
    $\Delta,\sigma_{e}\vdash{}d.cs\srarrow{init}{\fontsize{6}{8}\selectfont}\sigma_0$
  \end{itemize}
\end{definition}

\def\inithyps{For all $sitpn\in{}SITPN$,
  $b\in{}P\rightarrow\mathbb{N}$, $d\in{}design$,
  $\gamma\in{}WM(sitpn,d)$,
  $\Delta\in{}ElDesign,\sigma_{e},\sigma_0\in{}\Sigma$ that verify the
  hypotheses of Definition~\ref{def:init-states-hyps},}

%%%%%%%%%%%%%%%%%%%%%%%%%%%%%%%%%%%%%%%%%%
%%%%%% SIMILAR INITIAL STATES LEMMA %%%%%%
%%%%%%%%%%%%%%%%%%%%%%%%%%%%%%%%%%%%%%%%%%

\begin{lemma}[Similar initial states]
  \label{lem:sim-init-states}
  \inithyps{} then $\gamma\vdash{}s_0\sim\sigma_0$.
\end{lemma}

\begin{niproof}
  By definition of the \nameref{def:state-sim} relation, there are 6
  points to prove.
  \begin{frameb}
    \begin{enumerate}
    \item\label{item:init-sim-mark} $\forall{}p\in{}P,id_p\in{}Comps(\Delta)~s.t.~\gamma(p)=id_p,$
      $~s_0.M(p)=\sigma_0(id_p)(\texttt{s\_marking})$.
    \item\label{item:init-sim-tc}
      $\forall{}t\in{}T_i,id_t\in{}Comps(\Delta)~s.t.~\gamma(t)=id_t,$\\
      $\big(upper(I_s(t))=\infty\land{}s_0.I(t)\le{}lower(I_s(t))\Rightarrow{}s_0.I(t)=\sigma_0(id_t)(\texttt{s\_time\_counter})\big)$\\
      $\land\big(upper(I_s(t))=\infty\land{}s_0.I(t)>{}lower(I_s(t))\Rightarrow{}\sigma_0(id_t)(\texttt{s\_time\_counter})=lower(I_s(t))\big)$\\
      $\land\big(upper(I_s(t))\neq\infty\land{}s_0.I(t)>{}upper(I_s(t))\Rightarrow{}\sigma_0(id_t)(\texttt{s\_time\_counter})=upper(I_s(t))\big)$\\
      $\land\big(upper(I_s(t))\neq\infty\land{}s_0.I(t)\le{}upper(I_s(t))\Rightarrow{}s_0.I(t)=\sigma_0(id_t)(\texttt{s\_time\_counter})\big)$.
    \item\label{item:init-sim-reset}
      $\forall{}t\in{}T_i,id_t\in{}Comps(\Delta)~s.t.~\gamma(t)=id_t,$
      $s_0.reset_t(t)=\sigma_0(id_t)(\texttt{s\_reinit\_time\_counter})$.
    \item\label{item:init-sim-cond}
      $\forall{}c\in\mathcal{C},id_c\in{}Ins(\Delta)~s.t.~\gamma(c)=id_c,~s_0.cond(c)=\sigma_0(id_c)$.
    \item\label{item:init-sim-act}
      $\forall{}a\in\mathcal{A},id_a\in{}Outs(\Delta)~s.t.~\gamma(a)=id_a,~s_0.ex(a)=\sigma_0(id_a)$.
    \item\label{item:init-sim-fun}
      $\forall{}f\in\mathcal{F},id_f\in{}Outs(\Delta)~s.t.~\gamma(f)=id_f,~s_0.ex(f)=\sigma_0(id_f)$.
    \end{enumerate}
  \end{frameb}

  \begin{itemize}
  \item Apply the \nameref{lem:init-states-eq-marking} lemma to solve \ref{item:init-sim-mark}.
  \item Apply the \nameref{lem:init-states-eq-tc} lemma to solve \ref{item:init-sim-tc}.
  \item Apply the \nameref{lem:init-states-eq-rorders} lemma to solve \ref{item:init-sim-reset}.
  \item Apply the \nameref{lem:init-states-cond-vals} lemma to solve \ref{item:init-sim-cond}.
  \item Apply the \nameref{lem:init-states-act-exec} lemma to solve \ref{item:init-sim-act}.
  \item Apply the \nameref{lem:init-states-fun-exec} lemma to solve \ref{item:init-sim-fun}.
  \end{itemize}
\end{niproof}

\subsection{Initial states and marking}
\label{sec:init-states-marking}

\begin{lemma}[Initial states equal marking]
  \label{lem:init-states-eq-marking}
  \inithyps{} then
  $\forall{}p\in{}P,id_p\in{}Comps(\Delta)$ s.t. $\gamma(p)=id_p$,
  $~s_0.M(p)=\sigma_0(id_p)(\texttt{s\_marking})$.
\end{lemma}

\begin{niproof}
  Given a $p\in{}P$ and an $id_p\in{}Comps(\Delta)$
  s.t. $\gamma(p)=id_p$, let us show that\\
  \framebox{$s_0.M(p)=\sigma_0(id_p)(\texttt{s\_marking})$.}

  \exP{}
  
  By property of the \hvhdl{} initialization relation, \InCsCompP{},
  and through the examination of the \texttt{marking} process defined
  in the place design architecture, we can deduce
  $\sigma_0(id_p)(\texttt{s\_marking})=\sigma_0(id_p)(\texttt{initial\_marking})$.

  Rewriting $\sigma_0(id_p)(\texttt{sm})$ as
  $\sigma_0(id_p)(\texttt{initial\_marking})$,
  \framebox{$\sigma_0(id_p)(\texttt{initial\_marking})=s_0.M(p)$.}
  
  By construction,
  ${<}\mathtt{initial\_marking\Rightarrow}M_0(p){>}\in{}i_p$.

  By property of the \hvhdl{} initialization relation, and \InCsCompP,
  then $\sigma_0(id_p)(\texttt{initial\_marking})=M_0(p)$.  Rewriting
  $\sigma_0(id_p)(\texttt{initial\_marking})$ as $M_0(p)$ in the current
  goal: \framebox{$M_0(p)=s_0.M(p)$.}
  
  By definition of $s_0$, we can rewrite $s_0.M(p)$ as $M_0(p)$ in the
  current goal, \qedbox{tautology.}
  
\end{niproof}

\begin{lemma}[Null input token sum at initial state]
  \label{lem:init-states-sits-zero}
  \inithyps{} then
  $\forall{}p\in{}P,id_p\in{}Comps(\Delta)~s.t.~\gamma(p)=id_p$,
  $\sigma_0(id_p)(\texttt{s\_input\_token\_sum})=0$.
\end{lemma}

\begin{niproof}
  Given a $p$ and an $id_p$ s.t. $\gamma(p)=id_p$, let us show that
  \fbox{$\sigma_0(id_p)(\texttt{s\_input\_token\_sum})=0$.}

  \exP{}
  
  By property of the initialization relation, \InCsCompP{}, and
  through the examination of the \texttt{input\_tokens\_sum} process
  defined in the place design architecture, we can deduce:
  \begin{equation}
    \label{eq:sits-at-init-state}
    \sigma_0(id_p)(\texttt{sits})=\sum\limits_{i=0}^{\Delta(id_p)(\texttt{ian})-1}
    \begin{cases}
      \sigma_0(id_p)(\texttt{iaw})[i]~\mathtt{if}~\sigma_0(id_p)(\texttt{itf})[i]\\
      0~otherwise \\
    \end{cases}
  \end{equation}
  
  Rewriting the goal with Equation~\eqref{eq:sits-at-init-state}:\\
  \fbox{$\sum\limits_{i=0}^{\Delta(id_p)(\texttt{ian})-1}\begin{cases}
      \sigma_0(id_p)(\texttt{iaw})[i]~\mathtt{if}~\sigma_0(id_p)(\texttt{itf})[i]\\
      0~otherwise \\
    \end{cases}=0$.}

  \noindent{}Let us perform case analysis on $input(p)$; there are two cases:

  \begin{enumerate}
  \item $input(p)=\emptyset$:
    
    By construction, we have
    ${<}$\texttt{input\_arcs\_number}$\Rightarrow{}1{>}\in{}g_p$,\\
    ${<}$\texttt{input\_transitions\_fired(0)}$\Rightarrow{}\mathtt{true}{>}\in{}i_p$,\\
    and
    ${<}\texttt{input\_arcs\_weights(0)}\Rightarrow{}0{>}\in{}i_p$.

  By property of the elaboration relation, \InCsCompP{}, and
  ${<}$\texttt{input\_arcs\_number}$\Rightarrow{}1{>}\in{}g_p$, we
  can deduce $\Delta(id_p)(\texttt{ian})=1$.

  By property of the initialization relation, \InCsCompP,
  ${<}$\texttt{input\_transitions\_fired(0)}$\Rightarrow{}\mathtt{true}{>}\in{}i_p$
  and
  ${<}\mathtt{input\_arcs\_weights(0)\Rightarrow{}0}{>}\in{}i_p$, we
  can deduce $\sigma_0(id_p)(\texttt{itf})[0]=\mathtt{true}$ and
  $\sigma_0(id_p)(\texttt{iaw})[0]=0$.

  Rewriting the goal with $\Delta(id_p)(\texttt{ian})=1$,
  $\sigma_0(id_p)(\texttt{itf})[0]=\mathtt{true}$,
  $\sigma_0(id_p)(\texttt{iaw})[0]=0$ and simplifying the goal,
  \qedbox{tautology.}
  
\item $input(p)\neq\emptyset$:

  By construction,
  ${<}\texttt{input\_arcs\_number}\Rightarrow{}\vert{}input(p)\vert{>}\in{}g_p$,
  and by property of the elaboration relation, and \InCsCompP{}, we
  can deduce $\Delta(id_p)(\texttt{ian})=\vert{}input(p)\vert$.
  
  Let us reason by induction on the sum term of the goal.

  \begin{itemize}
  \item \textbf{BASE CASE}: The sum term equals 0, then \qedbox{tautology.}

  \item \textbf{INDUCTION CASE}:
    \begin{ih}
      $\sum\limits_{i=1}^{\Delta(id_p)(\texttt{ian})-1}\begin{cases}
        \sigma_0(id_p)(\texttt{iaw})[i]~\mathtt{if}~\sigma_0(id_p)(\texttt{itf})[i]\\
        0~otherwise \\
      \end{cases}=0$
    \end{ih}

    \begin{frameb}
      \begin{tabular}{c}
        $\begin{cases}
          \sigma_0(id_p)(\texttt{iaw})[0]~\mathtt{if}~\sigma_0(id_p)(\texttt{itf})[0]\\
          0~otherwise \\
        \end{cases}$ \\
        $+$ \\
        $\sum\limits_{i=1}^{\Delta(id_p)(\texttt{ian})-1}\begin{cases}
        \sigma_0(id_p)(\texttt{iaw})[i]~\mathtt{if}~\sigma_0(id_p)(\texttt{itf})[i]\\
        0~otherwise \\
      \end{cases}=0$
      \end{tabular}
    \end{frameb}

    Using the induction hypothesis to rewrite the goal:\\
    \fbox{$\begin{cases}
        \sigma_0(id_p)(\texttt{iaw})[0]~\mathtt{if}~\sigma_0(id_p)(\texttt{itf})[0]\\
        0~otherwise \\
      \end{cases}=0$}
    
    Since $input(p)\neq\emptyset$, by construction, there exist an
    $id_t\in{}Comps(\Delta),g_t,i_t,o_t$ s.t. \InCsCompT{},
    $id_{ft}\in{}Sigs(\Delta)$ s.t.
    ${<}\mathtt{fired\Rightarrow}id_{ft}{>}\in{}o_t$ and\\
    ${<}\mathtt{input\_transitions\_fired(0)\Rightarrow{}id_{ft}}{>}\in{}i_p$.

    By property of the initialization relation, \InCsCompP{},
    \InCsCompT{}, ${<}\mathtt{fired\Rightarrow}id_{ft}{>}\in{}o_t$
    and
    ${<}$\texttt{input\_transitions\_fired(0)}$\Rightarrow{}\mathtt{id_{ft}}{>}\in{}i_p$,
    we can deduce $\sigma_0(id_p)(\texttt{itf})[0]=\sigma_0(id_t)(\texttt{fired})$.

    Rewriting the goal with $\sigma_0(id_p)(\texttt{itf})[0]=\sigma_0(id_t)(\texttt{fired})$:\\
    \fbox{$\begin{cases}
        \sigma_0(id_p)(\texttt{iaw})[0]~\mathtt{if}~\sigma_0(id_t)(\texttt{fired})\\
        0~otherwise \\
      \end{cases}=0$}

    Appealing to Lemma~\ref{lem:init-states-fired-false}, we can
    deduce $\sigma_0(id_t)(\texttt{fired})=\mathtt{false}$.

    Rewriting the goal with $\sigma_0(id_t)(\texttt{fired})=\mathtt{false}$,
    and simplifying the goal, \qedbox{tautology.}
  \end{itemize}
\end{enumerate}
  
\end{niproof}

\begin{lemma}[Null output token sum at initial state]
  \label{lem:init-states-sots-zero}
  \inithyps{} then
  $\forall{}p\in{}P,id_p\in{}Comps(\Delta)~s.t.~\gamma(p)=id_p$,
  $\sigma_0(id_p)(\texttt{s\_output\_token\_sum})=0$.
\end{lemma}

\begin{niproof}
  The proof is similar to the proof of
  Lemma~\ref{lem:init-states-sits-zero}.
\end{niproof}

\subsection{Initial states and time counters}
\label{sec:init-states-tc}

\begin{lemma}[Initial states equal time counters]
  \label{lem:init-states-eq-tc}
  \inithyps{} then
  $\forall{}t\in{}T_i,id_t\in{}Comps(\Delta)$ s.t. $\gamma(t)=id_t$,\\
  $upper(I_s(t))=\infty\land{}s_0.I(t)\le{}lower(I_s(t))\Rightarrow{}s_0.I(t)=\sigma_0(id_t)(\texttt{s\_time\_counter})\land{}$\\
  $upper(I_s(t))=\infty\land{}s_0.I(t)>{}lower(I_s(t))\Rightarrow{}\sigma_0(id_t)(\texttt{s\_time\_counter})=lower(I_s(t))\land{}$\\
  $upper(I_s(t))\neq\infty\land{}s_0.I(t)>{}upper(I_s(t))\Rightarrow{}\sigma_0(id_t)(\texttt{s\_time\_counter})=upper(I_s(t))\land{}$\\
  $upper(I_s(t))\neq\infty\land{}s_0.I(t)\le{}upper(I_s(t))\Rightarrow{}s_0.I(t)=\sigma_0(id_t)(\texttt{s\_time\_counter})$.
\end{lemma}

\begin{niproof}
  Given a $t\in{}T_i$ and an $id_t\in{}Comps(\Delta)$
  s.t. $\gamma(t)=id_t$, let us show that:
  \begin{enumerate}
  \item \framebox{$upper(I_s(t))=\infty\land{}s_0.I(t)\le{}lower(I_s(t))\Rightarrow{}s_0.I(t)=\sigma_0(id_t)(\texttt{s\_time\_counter})$}
  \item \framebox{$upper(I_s(t))=\infty\land{}s_0.I(t)>{}lower(I_s(t))\Rightarrow{}\sigma_0(id_t)(\texttt{s\_time\_counter})=lower(I_s(t))$}
  \item \framebox{$upper(I_s(t))\neq\infty\land{}s_0.I(t)>{}upper(I_s(t))\Rightarrow{}\sigma_0(id_t)(\texttt{s\_time\_counter})=upper(I_s(t))$}
  \item \framebox{$upper(I_s(t))\neq\infty\land{}s_0.I(t)\le{}upper(I_s(t))\Rightarrow{}s_0.I(t)=\sigma_0(id_t)(\texttt{s\_time\_counter})$}
  \end{enumerate}

  \exP

  Then, let us show the 4 previous points.
  
  \begin{enumerate}
  \item Assuming that $upper(I_s(t))=\infty\land{}s_0.I(t)\le{}lower(I_s(t))$, then let us show\\
    \framebox{${}s_0.I(t)=\sigma_0(id_t)(\texttt{s\_time\_counter})$.}
    
    Rewriting $s_0.I(t)$ as $0$, by definition of $s_0$,
    \framebox{$\sigma_0(id_t)(\texttt{s\_time\_counter})=0$.}

    \noindent By property of the \hvhdl{} initialization relation,
    \InCsCompT, and through the examination of the
    \texttt{time\_counter} process defined in the transition design
    architecture, we can deduce
    \qedbox{$\sigma_0(id_t)(\texttt{s\_time\_counter})=0$.}
    
  \item Assuming that $upper(I_s(t))=\infty$ and
    $s_0.I(t)>{}lower(I_s(t))$, let us show\\
    \framebox{$\sigma_0(id_t)(\texttt{s\_time\_counter})=lower(I_s(t))$}.

    By definition, $lower(I_s(t))\in\mathbb{N}^{*}$ and
    $s_0.I(t)=0$. Then, \qedbox{$lower(I_s(t)){}<0$ is a
      contradiction.}
    
  \item Assuming that $upper(I_s(t))\neq\infty$ and
    $s_0.I(t)>{}upper(I_s(t))$, let us show\\
    \framebox{$\sigma_0(id_t)(\texttt{s\_time\_counter})=upper(I_s(t))$}.

    By definition, $upper(I_s(t))\in\mathbb{N}^{*}$ and
    $s_0.I(t)=0$. Then, \qedbox{$upper(I_s(t)){}<0$ is a
      contradiction.}
    
  \item Assuming that $upper(I_s(t))\neq\infty$ and
    $s_0.I(t)\le{}upper(I_s(t))$, let us
    show\\ \framebox{$s_0.I(t)=\sigma_0(id_t)(\texttt{s\_time\_counter})$}.
 
    Rewriting $s_0.I(t)$ as $0$, by definition of $s_0$,
    \framebox{$\sigma_0(id_t)(\texttt{s\_time\_counter})=0$.}

    By property of the \hvhdl{} initialization relation, \InCsCompT,
    and through the examination of the \texttt{time\_counter} process
    defined in the transition design architecture, we can deduce
    \qedbox{$\sigma_0(id_t)(\texttt{s\_time\_counter})=0$.}
  \end{enumerate}
\end{niproof}

\subsection{Initial states and reset orders}
\label{sec:init-states-rorders}

\begin{lemma}[Initial states equal reset orders]
  \label{lem:init-states-eq-rorders}
  \inithyps{} then $\forall{}t\in{}T_i,id_t\in{}Comps(\Delta)$
  s.t. $\gamma(t)=id_t$,
  $s_0.reset_t(t)=\sigma_0(id_t)(\texttt{s\_reinit\_time\_counter})$.
\end{lemma}

\begin{niproof}
  Given a $t\in{}T_i$ and an $id_t\in{}Comps(\Delta)$ s.t.
  $\gamma(t)=id_t$, let us show
  that\\
  \framebox{$s_0.reset_t(t)=\sigma_0(id_t)(\texttt{s\_reinit\_time\_counter})$.}
  
  Rewriting $s_0.reset_t(t)$ as $\mathtt{false}$, by definition of
  $s_0$,
  \fbox{$\sigma_0(id_t)(\texttt{s\_reinit\_time\_counter})=\mathtt{false}$.}
  
  \exT{}
  
  \noindent By property of the \hvhdl{} initialization relation,
  \InCsCompT, and through the examination of the \texttt{reinit\_time\_counter\_evaluation} process defined in the \texttt{transition} design architecture\\
  we can deduce
  $\sigma_0(id_t)(\texttt{s\_reinit\_time\_counter})=\prod\limits_{i=0}^{\Delta(id_t)(\texttt{ian})-1}\sigma_0(id_t)(\texttt{rt})[i]$.

  Rewriting
  $\sigma_0(id_t)(\texttt{s\_reinit\_time\_counter})$ as $\prod\limits_{i=0}^{\Delta(id_t)(\texttt{ian})-1}\sigma_0(id_t)(\texttt{rt})[i]$,\\
  \fbox{$\prod\limits_{i=0}^{\Delta(id_t)(\texttt{ian})-1}\sigma_0(id_t)(\texttt{rt})[i]=\mathtt{false}$.}
  
  For all $t\in{}T$ (resp. $p\in{}P$), let $input(t)$
  (resp. $input(p)$) be the set of input places of $t$ (resp. input
  transitions of $p$), and let $output(t)$ (resp. $output(p)$) be the
  set of output places of $t$ (resp. output transitions of $p$).

  Let us perform case analysis on $input(t)$; there are 2 cases:

  \begin{itemize}
  \item \textbf{CASE} $input(t)=\emptyset$.

    By construction,
    ${<}\mathtt{input\_arcs\_number\Rightarrow}1{>}\in{}g_t$, and by
    property of the elaboration relation, and \InCsCompT{}, we can
    deduce $\Delta(id_t)(\texttt{ian})=1$.

    By construction, $<\mathtt{reinit\_time(0)\Rightarrow{}false}>\in{}i_t$, and
    by property of the initialization relation and \InCsCompT, we can
    deduce $\sigma_0(id_t)(\texttt{rt})[0]=\mathtt{false}$.

    Rewriting $\Delta(id_t)(\texttt{ian})$ as $1$ and
    $\sigma_0(id_t)(\texttt{rt})[0]$ as $\mathtt{false}$, \qedbox{tautology.}
    
  \item \textbf{CASE} $input(t)\neq\emptyset$.

    To prove the current goal, we can equivalently prove that\\
    \fbox{$\exists{}i\in[0,\Delta(id_t)(\texttt{ian})-1]~s.t.~\sigma_0(id_t)(\texttt{rt})[i]=\mathtt{false}$.}

    Since $input(t)\neq\emptyset,~\exists{}p~s.t.~p\in{}input(t)$. Let
    us take such a $p\in{}input(t)$.
    
    By construction, for all $p\in{}P$, there exist
    $id_p~s.t.~\gamma(p)=id_p$.

    \exP{}

    \noindent{}By construction, there exist
    $i\in[0,\vert{}input(t)\vert{}-1]$,
    $j\in[0,\vert{}output(p)\vert{}-1]$, $id_{ji}\in{}Sigs(\Delta)$
    s.t.
    ${<}\mathtt{reinit\_transitions\_time(j)\Rightarrow}id_{ji}{>}\in{}o_p$
    and
    ${<}\mathtt{reinit\_time(i)\Rightarrow}id_{ji}{>}\in{}i_t$. Let us take such a $i$, $j$ and $id_{ji}$.

    By construction and $input(t)\neq\emptyset$,
    ${<}\mathtt{input\_arcs\_number\Rightarrow}\vert{}input(t)\vert{}{>}\in{}g_t$.

    By property of the \hvhdl{} elaboration relation and
    ${<}\mathtt{input\_arcs\_number\Rightarrow}\vert{}input(t)\vert{}{>}\in{}g_t$,
    we can deduce $\Delta(id_t)(\texttt{ian})=\vert{}input(t)\vert$.

    Since $\Delta(id_t)(\texttt{ian})=\vert{}input(t)\vert$ and we have an
    $i\in[0,\vert{}input(t)\vert-1]$, then, we have an
    $i\in[0, \Delta(id_t)(\texttt{ian})-1]$. Let us take that i to
    prove the goal.

    Then, we must show \framebox{$\sigma_0(id_t)(\texttt{rt})[i]=\mathtt{false}$.}

    By property of the \hvhdl{} initialization relation and
    ${<}\mathtt{reinit\_time(i)\Rightarrow}id_{ji}{>}\in{}i_t$, we can
    deduce $\sigma_0(id_t)(\texttt{rt})[i]=\sigma_0(id_{ji})$.

    Rewriting $\sigma_0(id_t)(\texttt{rt})[i]$ as $\sigma_0(id_{ji})$,
    \framebox{$\sigma_0(id_{ji})=\mathtt{false}$.}

    By property of the \hvhdl{} initialization relation and\\
    ${<}\mathtt{reinit\_transitions\_time(j)\Rightarrow}id_{ji}{>}\in{}o_p$,
    we can deduce
    $\sigma_0(id_{ji})=\sigma_0(id_p)(\texttt{rtt})[j]$.

    Rewriting $\sigma_0(id_{ji})$ as
    $\sigma_0(id_p)(\texttt{rtt})[j]$,
    \framebox{$\sigma_0(id_p)(\texttt{rtt})[j]=\mathtt{false}$.}

    Since $t\in{}output(p)$, then we know that
    $output(p)\neq\emptyset$.

    Then, by construction,
    ${<}\mathtt{output\_arcs\_number\Rightarrow}\vert{}output(p)\vert{>}\in{}g_p$.

    By property of the elaboration relation and
    ${<}\mathtt{output\_arcs\_number}\Rightarrow\vert{}output(p)\vert{>}\in{}g_p$,
    we can deduce that $\Delta(id_p)(\texttt{oan})=\vert{}output(p)\vert$.

    Since $\Delta(id_p)(\texttt{oan})=\vert{}output(p)\vert$ and
    $j\in[0,\vert{}output(p)\vert-1]$, then
    $j\in[0,\Delta(id_p)(\texttt{oan})-1]$.
    
    By property of the \hvhdl{} initialization relation, \InCsCompP,
    through the examination of the
    \texttt{reinit\_transitions\_time\_evaluation} process defined in
    the \texttt{place} design architecture, and since
    $j\in[0,\Delta(id_p)(\texttt{oan})-1]$,
    \qedbox{$\sigma_0(id_p)(\texttt{rtt})[j]=\mathtt{false}$.}

  \end{itemize}
  
\end{niproof}

\subsection{Initial states and condition values}
\label{sec:init-states-cond-vals}

\begin{lemma}[Initial states equal condition values]
  \label{lem:init-states-cond-vals}
  \inithyps{} 
  then
  $\forall{}c\in\mathcal{C},id_c\in{}Ins(\Delta)~s.t.~\gamma(c)=id_c,~s_0.cond(c)=\sigma_0(id_c)$.
\end{lemma}

\begin{niproof}
  Given a $c\in\mathcal{C}$ and an
  $id_c\in{}Ins(\Delta)~s.t.~\gamma(c)=id_c$, let us show that
  \fbox{$s_0.cond(c)=\sigma_0(id_c)$.}\\

  Rewriting $s_0.cond(c)$ as $\mathtt{false}$, by definition of $s_0$,
  \fbox{$\sigma_0(id_c)=\mathtt{false}$.}

  By construction, $id_c$ is an input port identifier of Boolean type
  in the \hvhdl{} design $d$, and thus, by property of the \hvhdl{}
  elaboration relation, we can deduce $\sigma_e(id_c)=\mathtt{false}$.

  By property of the \hvhdl{} initialization relation and
  $id_c\in{}Ins(\Delta)$, we can deduce
  $\sigma_e(id_c)=\sigma_0(id_c)$.

  Rewriting $\sigma_0(id_c)$ as $\sigma_e(id_c)$ and $\sigma_e(id_c)$
  as $\mathtt{false}$, \qedbox{tautology.}
  
\end{niproof}

\subsection{Initial states and action executions}
\label{sec:init-states-act-exec}

\begin{lemma}[Initial states equal action executions]
  \label{lem:init-states-act-exec}
  \inithyps{} 
  then
  $\forall{}a\in\mathcal{A},id_a\in{}Outs(\Delta)~s.t.~\gamma(a)=id_a,~s_0.ex(a)=\sigma_0(id_a)$.
\end{lemma}

\begin{niproof}
  Given a $a\in\mathcal{A}$ and an
  $id_a\in{}Outs(\Delta)~s.t.~\gamma(a)=id_a$, let us show that
  \fbox{$s_0.ex(a)=\sigma_0(id_a)$.}\\

  Rewriting $s_0.ex(a)$ as $\mathtt{false}$, by definition of $s_0$,
  \fbox{$\sigma_0(id_a)=\mathtt{false}$.}

  By construction, $id_a$ is an output port identifier of Boolean type
  in the \hvhdl{} design $d$. Moreover, we know that the output port
  identifier $id_a$ is assigned to $\mathtt{false}$ in the generated
  \texttt{action} process during the initialization phase (i.e. the
  assignment is a part of a \emph{reset} block). Thus, we can deduce
  that $\sigma_0(id_a)=\mathtt{false}$.
  
  Rewriting $\sigma_0(id_a)$ as $\mathtt{false}$, \qedbox{tautology.}
  
\end{niproof}

\subsection{Initial states and function executions}
\label{sec:init-states-fun-exec}

\begin{lemma}[Initial states equal function executions]
  \label{lem:init-states-fun-exec}
  \inithyps{} 
  then
  $\forall{}f\in\mathcal{F},id_f\in{}Outs(\Delta)~s.t.~\gamma(f)=id_f,~s_0.ex(f)=\sigma_0(id_f)$.
\end{lemma}

\begin{niproof}
  Given a $f\in\mathcal{F}$ and an
  $id_f\in{}Outs(\Delta)~s.t.~\gamma(f)=id_f$, let us show that
  \fbox{$s_0.ex(f)=\sigma_0(id_f)$.}\\

  Rewriting $s_0.ex(f)$ as $\mathtt{false}$, by definition of $s_0$,
  \fbox{$\sigma_0(id_f)=\mathtt{false}$.}

  By construction, $id_f$ is an output port identifier of Boolean type
  in the \hvhdl{} design $d$, and thus, by property of the \hvhdl{}
  elaboration relation, we can deduce $\sigma_e(id_f)=\mathtt{false}$.

  By construction, and by property of the initialization relation, we
  know that the output port identifier $id_f$ is assigned to
  \texttt{false} in the generated \texttt{function} process during the
  initialization phase (i.e. the assignment is a part of a
  \emph{reset} block). Thus, we can deduce
  $\sigma_0(id_f)=\mathtt{false}$.
  
  \noindent{}Rewriting $\sigma_0(id_f)$ as $\mathtt{false}$,
  \qedbox{tautology.}
  
\end{niproof}

\subsection{Initial states and fired transitions}
\label{sec:init-states-fired-false}

\begin{lemma}[No fired at initial state]
  \label{lem:init-states-fired-false}
  $\forall{}d\in{}design,\Delta\in{}ElDesign,\sigma_e,\sigma_0\in\Sigma,id_t\in{}Comps(\Delta),$
  $g_t,$ $i_t,$ $o_t$ s.t. :
  \begin{itemize}
  \item $\mathcal{D}_\mathcal{H},\emptyset\vdash{}d.cs\xrightarrow{elab}\sigma_0$
  \item $\Delta,\sigma_e\vdash{}d.cs\xrightarrow{init}\sigma_0$
  \item $\mathtt{comp}(id_t$, $\texttt{transition}$, $g_t$, $i_t$,
    $o_t)\in{}d.cs$
  \end{itemize}
  then $\sigma_0(id_t)(\texttt{fired})=\mathtt{false}$.
\end{lemma}

\begin{niproof}
  Assuming all the above hypotheses, let us show
  \fbox{$\sigma_0(id_t)(\texttt{fired})=\mathtt{false}$.}

  By property of the initialization relation, \InCsCompT, and through
  the examination of the \texttt{fired\_evaluation} process defined in
  the transition design architecture, we can deduce:
  \begin{equation}
    \label{eq:fired-at-init-state}
    \sigma_0(id_t)(\texttt{fired})=\sigma_0(id_t)(\texttt{s\_firable})~.~\sigma_0(id_t)(\texttt{s\_priority\_combination})
  \end{equation}

  Rewriting the goal with Equation~\eqref{eq:fired-at-init-state}:
  \fbox{$\sigma_0(id_t)(\texttt{sfa})~.~\sigma_0(id_t)(\texttt{spc})=\mathtt{false}$.}
  
  By property of the initialization relation, \InCsCompT, and through
  the examination of the \texttt{firable} process defined in the
  \texttt{transition} design architecture, we can deduce
  $\sigma_0(id_t)(\texttt{sfa})=\mathtt{false}$.

  Rewriting the goal with $\sigma_0(id_t)(\texttt{sfa})=\mathtt{false}$ and
  simplifying the goal, \qedbox{tautology.}
\end{niproof}


%%% Local Variables:
%%% mode: latex
%%% TeX-master: "../../main"
%%% End:


\section{First Rising Edge}
\label{sec:fst-re}
\begin{definition}[First rising edge hypotheses]
  \label{def:fst-re-hyps}
  Given an $sitpn\in{}SITPN,$ $d\in{}design,$\\
  $\gamma\in{}WM(sitpn,d),$ $\Delta\in{}ElDesign,$
  $\sigma_{e},\sigma_0,\sigma_i,\sigma_{\uparrow},\sigma\in{}\Sigma$,
  $E_c\in{}\mathbb{N}\rightarrow{}\mathcal{C}\rightarrow{}\mathbb{B}$,\\
  $E_p\in\mathbb{N}\times{}\{\uparrow,\downarrow\})\rightarrow{}Ins(\Delta)\rightarrow{}value$,
  $\tau\in\mathbb{N}$, assume that:
  \begin{itemize}
  \item $\lfloor{}sitpn\rfloor_\mathcal{H}=(d,\gamma)$ and
    $\mathcal{D}_\mathcal{H},\emptyset\vdash{}d\srarrow{elab}{\fontsize{6}{8}\selectfont}(\Delta,\sigma_{e})$
    and $\gamma\vdash{}E_p\stackrel{env}{=}E_c$
  \item $\sigma_0$ is the initial state of $\Delta$: 
    $\Delta,\sigma_{e}\vdash{}d.cs\srarrow{init}{\fontsize{6}{8}\selectfont}\sigma_0$
  \item $E_c,\tau\vdash{}s_0\srarrow{\uparrow_0}{\fontsize{6}{8}\selectfont}s_0$
  \item $\mathtt{Inject}(\sigma_0, E_p, \tau, \sigma_i)$
    and
    $\Delta,\sigma_i\vdash\mathrm{d.cs}\xrightarrow{\uparrow}\sigma_{\uparrow}$
    and
    $\Delta,\sigma_{\uparrow}\vdash\mathrm{d.cs}\xrightarrow{\theta}\sigma$
  \end{itemize}
  
\end{definition}

\def\fstrehyps{For all $sitpn,d,\gamma,\Delta,$
  $\sigma_{e},\sigma_0,\sigma_i,\sigma_{\uparrow},\sigma$, $E_c$,
  $E_p$, $\tau$ that verify the hypotheses of
  Definition~\ref{def:fst-re-hyps},}

%%%%%%%%%%%%%%%%%%%%%%%%%%%%%%%%%%%%%%%%%%%%%
%%%%%%%%%% FIRST RISING EDGE LEMMA %%%%%%%%%%
%%%%%%%%%%%%%%%%%%%%%%%%%%%%%%%%%%%%%%%%%%%%%

\begin{lemma}[First rising edge]
  \label{lem:fst-re}
  \fstrehyps{} then
  $\gamma,E_c,\tau\vdash{}s_0\stackrel{\uparrow}{\approx}{}\sigma$.
\end{lemma}

\begin{proof}
  By definition of the \nameref{def:full-post-re-state-sim} relation,
  there are 8 points to prove.
  \begin{frameb}
    \begin{enumerate}
    \item
      $\forall{}p\in{}P,id_p\in{}Comps(\Delta)~s.t.~\gamma(p)=id_p,~s_0.M(p)=\sigma(id_p)(\texttt{s\_marking})$.
      \label{item:fst-re-marking-eq}
    \item
      $\forall{}t\in{}T_i,id_t\in{}Comps(\Delta)~s.t.~\gamma(t)=id_t,$\\
      $\big(upper(I_s(t))=\infty\land{}s_0.I(t)\le{}lower(I_s(t))\Rightarrow{}s_0.I(t)=\sigma(id_t)(\texttt{s\_time\_counter})\big)$\\
      $\land\big(upper(I_s(t))=\infty\land{}s_0.I(t)>{}lower(I_s(t))\Rightarrow{}\sigma(id_t)(\texttt{s\_time\_counter})=lower(I_s(t))\big)$\\
      $\land\big(upper(I_s(t))\neq\infty\land{}s_0.I(t)>{}upper(I_s(t))\Rightarrow{}\sigma(id_t)(\texttt{s\_time\_counter})=upper(I_s(t))\big)$\\
      $\land\big(upper(I_s(t))\neq\infty\land{}s_0.I(t)\le{}upper(I_s(t))\Rightarrow{}s_0.I(t)=\sigma(id_t)(\texttt{s\_time\_counter})\big)$.
      \label{item:fst-re-tc-eq}
    \item
      $\forall{}t\in{}T_i,id_t\in{}Comps(\Delta)~s.t.~\gamma(t)=id_t,$
      $s_0.reset_t(t)=\sigma(id_t)(\texttt{s\_reinit\_time\_counter})$.
      \label{item:fst-re-reset-eq}
    \item
      $\forall{}a\in\mathcal{A},id_a\in{}Outs(\Delta)~s.t.~\gamma(a)=id_a,~s_0.ex(a)=\sigma(id_a)$.
      \label{item:fst-re-action-eq}
    \item
      $\forall{}f\in\mathcal{F},id_f\in{}Outs(\Delta)~s.t.~\gamma(f)=id_f,~s_0.ex(f)=\sigma(id_f)$.
      \label{item:fst-re-fun-eq}
    \item $\forall{}t\in{}T,id_t\in{}Comps(\Delta)~s.t.~\gamma(t)=id_t,$
      $t\in{}Sens(s_0.M)\Leftrightarrow\sigma(id_t)(\texttt{s\_enabled})=\mathtt{true}$.
      \label{item:fst-re-sens-eq}
    \item $\forall{}t\in{}T,id_t\in{}Comps(\Delta)~s.t.~\gamma(t)=id_t,$
      $t\notin{}Sens(s_0.M)\Leftrightarrow\sigma(id_t)(\texttt{s\_enabled})=\mathtt{false}$.
      \label{item:fst-re-sens-neq}
    \item
      $\forall{}t\in{}T,id_t\in{}Comps(\Delta)~s.t.~\gamma(t)=id_t,$\\
      $\sigma(id_t)(\texttt{s\_condition\_combination})=
      \prod\limits_{c\in{}conds(t)}
      \begin{cases}
        E_c(\tau,c) & if~\mathbb{C}(t,c)=1 \\
        \mathtt{not}(E_c(\tau,c)) & if~\mathbb{C}(t,c)=-1 \\
      \end{cases}$\\
      where
      $conds(t)=\{c\in\mathcal{C}~\vert~\mathbb{C}(t,c)=1\lor\mathbb{C}(t,c)=-1\}$.
      \label{item:fst-re-cond-comb-eq}
    \item $\forall{}c\in\mathcal{C},id_c\in{}Ins(\Delta)$
      s.t. $\gamma(c)=id_c$, $\sigma(id_c)=E_c(\tau,c)$.\label{item:fst-re-cond-eq}
    \end{enumerate}
  \end{frameb}

  \begin{itemize}
  \item Apply the \nameref{lem:fst-re-equal-marking} lemma to solve \ref{item:fst-re-marking-eq}.
  \item Apply the \nameref{lem:fst-re-equal-tc} lemma to solve \ref{item:fst-re-tc-eq}.
  \item Apply the \nameref{lem:fst-re-equal-reset-orders} lemma to solve \ref{item:fst-re-reset-eq}.
  \item Apply the \nameref{lem:fst-re-equal-action-ex} lemma to solve
    \ref{item:fst-re-action-eq}.
  \item Apply the \nameref{lem:fst-re-equal-fun-ex} lemma to solve
    \ref{item:fst-re-fun-eq}.
  \item Apply the \nameref{lem:fst-re-equal-sens} lemma to solve
    \ref{item:fst-re-sens-eq}.
  \item Apply the \nameref{lem:fst-re-neq-sens} lemma to solve
    \ref{item:fst-re-sens-neq}.
  \item Apply the \nameref{lem:fst-re-equal-cond-comb} lemma to solve
    \ref{item:fst-re-cond-comb-eq}.
  \item Apply the \nameref{lem:fst-re-equal-cond} lemma to solve
    \ref{item:fst-re-cond-eq}.
  \end{itemize}
  
\end{proof}

\subsection{First rising edge and marking}
\label{sec:fst-re-marking}

\begin{lemma}[First rising edge equal marking]
  \label{lem:fst-re-equal-marking}
  \fstrehyps{} then
  $\forall{}p\in{}P,id_p\in{}Comps(\Delta)~s.t.~\gamma(p)=id_p$,
  $~s_0.M(p)=\sigma(id_p)(\texttt{s\_marking})$.
\end{lemma}

\begin{niproof}
  Given a $p$ and an $id_p$ s.t. $\gamma(p)=id_p$, let us show that
  \fbox{$~s_0.M(p)=\sigma(id_p)(\texttt{s\_marking})$.}
  
  \exP{}
    
  By property of the \texttt{Inject} relation, the \hvhdl{} rising
  edge relation, the stabilize relation, \InCsCompP, and through the
  examination of the \texttt{marking} process defined in the place
  design architecture, we can deduce:
  \begin{equation}
    \sigma(id_p)(\texttt{sm})=\sigma_0(id_p)(\texttt{sm})+\sigma_0(id_p)(\texttt{sits})-\sigma_0(id_p)(\texttt{sots})\label{eq:eq-sm-after-fst-re}
  \end{equation}

  Rewriting the goal with Equation~\eqref{eq:eq-sm-after-fst-re}:\\
  \fbox{$s_0.M(p)=\sigma_0(id_p)(\texttt{sm})+\sigma_0(id_p)(\texttt{sits})-\sigma_0(id_p)(\texttt{sots})$.}

  Appealing to Lemmas~\ref{lem:init-states-sits-zero} and
  \ref{lem:init-states-sots-zero}, we can deduce
  $\sigma_0(id_p)(\texttt{sits})=0$ and $\sigma_0(id_p)(\texttt{sots})=0$. Rewriting
  the goal with $\sigma_0(id_p)(\texttt{sits})=0$ and
  $\sigma_0(id_p)(\texttt{sots})=0$,
  \fbox{$s_0.M(p)=\sigma_0(id_p)(\texttt{sm})$.}\\

  Appealing to Lemma~\ref{lem:init-states-eq-marking},
  \qedbox{$s_0.M(p)=\sigma_0(id_p)(\texttt{sm})$.}
\end{niproof}

\subsection{First rising edge and time counters}
\label{sec:fst-re-tc}

\begin{lemma}[First rising edge equal time counters]
  \label{lem:fst-re-equal-tc}
  \fstrehyps{} then\\
  $\forall{}t\in{}T_i,id_t\in{}Comps(\Delta)~s.t.~\gamma(t)=id_t$,\\
  $upper(I_s(t))=\infty\land{}s_0.I(t)\le{}lower(I_s(t))\Rightarrow{}s_0.I(t)=\sigma(id_t)(\texttt{s\_time\_counter})\land{}$\\
  $upper(I_s(t))=\infty\land{}s_0.I(t)>{}lower(I_s(t))\Rightarrow{}\sigma(id_t)(\texttt{s\_time\_counter})=lower(I_s(t))\land{}$\\
  $upper(I_s(t))\neq\infty\land{}s_0.I(t)>{}upper(I_s(t))\Rightarrow{}\sigma(id_t)(\texttt{s\_time\_counter})=upper(I_s(t))\land{}$\\
  $upper(I_s(t))\neq\infty\land{}s_0.I(t)\le{}upper(I_s(t))\Rightarrow{}s_0.I(t)=\sigma(id_t)(\texttt{s\_time\_counter})$.
\end{lemma}

\begin{niproof}
  Given a $t\in{}T_i$ and an $id_t\in{}Comps(\Delta)$
  s.t. $\gamma(t)=id_t$, let us show that:
  \begin{enumerate}
  \item \framebox{$upper(I_s(t))=\infty\land{}s_0.I(t)\le{}lower(I_s(t))\Rightarrow{}s_0.I(t)=\sigma(id_t)(\texttt{s\_time\_counter})$}
  \item \framebox{$upper(I_s(t))=\infty\land{}s_0.I(t)>{}lower(I_s(t))\Rightarrow{}\sigma(id_t)(\texttt{s\_time\_counter})=lower(I_s(t))$}
  \item \framebox{$upper(I_s(t))\neq\infty\land{}s_0.I(t)>{}upper(I_s(t))\Rightarrow{}\sigma(id_t)(\texttt{s\_time\_counter})=upper(I_s(t))$}
  \item \framebox{$upper(I_s(t))\neq\infty\land{}s_0.I(t)\le{}upper(I_s(t))\Rightarrow{}s_0.I(t)=\sigma(id_t)(\texttt{s\_time\_counter})$}
  \end{enumerate}

  \exT{}

  Then, let us show the 4 previous points:
  
  \begin{enumerate}
  \item Assuming that $upper(I_s(t))=\infty$ and
    $s_0.I(t)\le{}lower(I_s(t))$, let us show\\
    \framebox{${}s_0.I(t)=\sigma(id_t)(\texttt{stc})$.}
    
    By property of the $\mathtt{Inject}$ relation, the
    \hvhdl{} rising edge and stabilize relations, and \InCsCompT{}, we
    can deduce $\sigma(id_t)(\texttt{stc})=\sigma_0(id_t)(\texttt{stc})$.
    
    Rewriting $\sigma(id_t)(\texttt{stc})$ as $\sigma_0(id_t)(\texttt{stc})$,
    \fbox{${}s_0.I(t)=\sigma_0(id_t)(\texttt{stc})$.}

    Appealing to Lemma~\ref{lem:init-states-eq-tc},
    \qedbox{${}s_0.I(t)=\sigma_0(id_t)(\texttt{stc})$.}
  
  \item Assuming that $upper(I_s(t))=\infty$ and
    ${}s_0.I(t)>{}lower(I_s(t))$, let us show\\
    \framebox{$\sigma(id_t)(\texttt{stc})=lower(I_s(t))$.}

    By definition, $lower(I_s(t))\in\mathbb{N}^{*}$ and
    $s_0.I(t)=0$. Then, \qedbox{$lower(I_s(t)){}<0$ is a
      contradiction.}
  \item Assuming that $upper(I_s(t))\neq\infty$ and
    $s_0.I(t)>{}upper(I_s(t))$, let us show\\
    \framebox{$\sigma(id_t)(\texttt{stc})=upper(I_s(t))$}.

    By definition, $upper(I_s(t))\in\mathbb{N}^{*}$ and
    $s_0.I(t)=0$. Then, \qedbox{$upper(I_s(t)){}<0$ is a
      contradiction.}
    
  \item Assuming that $upper(I_s(t))\neq\infty$ and
    $s_0.I(t)\le{}upper(I_s(t))$, let us show\\
    \framebox{$s_0.I(t)=\sigma(id_t)(\texttt{stc})$.}

    By property of the $\mathtt{Inject}$ relation, the
    \hvhdl{} rising edge and stabilize relations, and \InCsCompT{}, we
    can deduce $\sigma(id_t)(\texttt{stc})=\sigma_0(id_t)(\texttt{stc})$.

    Rewriting $\sigma(id_t)(\texttt{stc})$ as $\sigma_0(id_t)(\texttt{stc})$,
    \fbox{${}s_0.I(t)=\sigma_0(id_t)(\texttt{stc})$.}

    Appealing to Lemma~\ref{lem:init-states-eq-tc},
    \qedbox{${}s_0.I(t)=\sigma_0(id_t)(\texttt{stc})$.}
  \end{enumerate}
\end{niproof}

\subsection{First rising edge and reset orders}
\label{sec:fst-re-reset-orders}

\begin{lemma}[First rising edge equal reset orders]
  \label{lem:fst-re-equal-reset-orders}
  \fstrehyps{} then\\
  $\forall{}t\in{}T,id_t\in{}Comps(\Delta)~s.t.~\gamma(t)=id_t,$
  $s_0.reset_t(t)=\sigma(id_t)(\texttt{s\_reinit\_time\_counter})$.
\end{lemma}

\begin{niproof}
  Given a $t\in{}T$ and an $id_t\in{}Comps(\Delta)$
  s.t. $\gamma(t)=id_t$, let us show that\\
  \fbox{$s_0.reset_t(t)=\sigma(id_t)(\texttt{srtc})$.}

  \exT{}

  \noindent{}By property of the stabilize relation, \InCsCompT{}, and
  through the examination of the
  \texttt{reinit\_time\_counter\_evaluation} process defined in the
  transition design architecture, we can deduce:  
  \begin{equation}
    \sigma(id_t)(\texttt{srtc})=\sum\limits_{i=0}^{\Delta(id_t)(\texttt{input\_arcs\_number})-1}\sigma(id_t)(\texttt{reinit\_time})[i]\label{eq:srtc-at-init-states}
  \end{equation}

  Rewriting the goal with Equation~\eqref{eq:srtc-at-init-states}:
  \fbox{$s_0.reset_t(t)=\sum\limits_{i=0}^{\Delta(id_t)(\texttt{ian})-1}\sigma(id_t)(\texttt{rt})[i]$.}
  
  Let us perform case analysis on $input(t)$; there are two cases:

  \begin{itemize}
  \item \textbf{CASE} $input(t)=\emptyset$:

    By construction,
    ${<}\mathtt{input\_arcs\_number\Rightarrow{}1}{>}\in{}g_t$, and
    by property of the \hvhdl{} elaboration relation, we can deduce
    $\Delta(id_t)(\texttt{ian})=1$.

    By construction,
    $<\mathtt{reinit\_time(0)\Rightarrow{}false}>\in{}i_t$, and by
    property of the \hvhdl{} stabilize relation,
    $\sigma(id_t)(\texttt{rt})[0]=\mathtt{false}$.

    Rewriting the goal with $\Delta(id_t)(\texttt{ian})=1$ and
    $\sigma(id_t)(\texttt{rt})[0]=\mathtt{false}$,
    \fbox{$s_0.reset_t(t)=\mathtt{false}$.}

    By definition of $s_0$, \qedbox{$s_0.reset_t(t)=\mathtt{false}$.}
    
  \item \textbf{CASE} $input(t)\neq{}\emptyset$:

    By construction,
    ${<}\mathtt{input\_arcs\_number\Rightarrow{}}\vert{}input(t)\vert{>}\in{}g_t$,
    and by property of the \hvhdl{} elaboration relation, we can
    deduce $\Delta(id_t)(\texttt{ian})=\vert{}input(t)\vert$.

    Rewriting $\Delta(id_t)(\texttt{ian})$ as $\vert{}input(t)\vert$,
    \fbox{$s_0.reset_t(t)=\sum\limits_{i=0}^{\vert{}input(t)\vert-1}\sigma(id_t)(\texttt{rt})[i]$.}

    By definition of $s_0$, $s_0.reset_t(t)=\mathtt{false}$. Rewriting
    $s_0.reset_t(t)$ as $\mathtt{false}$,\\
    \fbox{$\sum\limits_{i=0}^{\vert{}input(t)\vert-1}\sigma(id_t)(\texttt{rt})[i]=\mathtt{false}$.}

    Given a $i\in[0,\vert{}input(t)\vert-1]$, let us show
    \fbox{$\sigma(id_t)(\texttt{rt})[i]=\mathtt{false}$.}
    
    By construction, and since $input(t)\neq{}\emptyset$, there exist
    a $p\in{}input(t)$, an $id_p\in{}Comps(\Delta)$
    s.t. $\gamma(p)=id_p$, a $g_p,$ an $i_p,$ an $o_p$
    s.t. \InCsCompP{}, and there exist a
    $j\in[0,\vert{}output(p)\vert-1]$ and an
    $id_{ji}\in{}Sigs(\Delta)$ s.t.
    ${<}\mathtt{reinit\_transition\_time(j)\Rightarrow{}id_{ji}}{>}\in{}o_p$
    and
    ${<}\mathtt{reinit\_time(i)\Rightarrow{}id_{ji}}{>}\in{}i_t$.

    By property of the stabilize relation,
    ${<}$\texttt{reinit\_transition\_time(j)$\Rightarrow{}id_{ji}$}${>}\in{}o_p$
    and\\
    ${<}$\texttt{reinit\_time(i)$\Rightarrow{}id_{ji}$}${>}\in{}i_t$,
    we can deduce
    $\sigma(id_t)(\texttt{rt})[i]=\sigma(id_{ji})=\sigma(id_p)(\texttt{rtt})[j]$.

    Rewriting $\sigma(id_t)(\texttt{rt})[i]$ as $\sigma(id_{ji})$ and
    $\sigma(id_{ji})$ as $\sigma(id_p)(\texttt{rtt})[j]$,
    \fbox{$\sigma(id_p)(\texttt{rtt})[j]=\mathtt{false}$.}

    By property of the \hvhdl{} rising edge and stabilize relations,
    \InCsCompP{}, and through the examination of the process defined
    in the place design architecture, we can deduce:
    \begin{equation}
      \label{eq:rtt-at-init}
      \begin{split}
        \sigma(id_p)(\texttt{rtt})[j]=& ((\sigma_0(id_p)(\texttt{oat})[j]=\mathtt{basic}+\sigma_0(id_p)(\texttt{oat})[j]=\mathtt{test}) \\
        & .(\sigma_0(id_p)(\texttt{sm})-\sigma_0(id_p)(\texttt{sots})<\sigma_0(id_p)(\texttt{oaw})[j])\\
        & .(\sigma_0(id_p)(\texttt{sots})>0))\\
        & +(\sigma_0(id_p)(\texttt{otf})[j]) \\
      \end{split}
    \end{equation}

    Rewriting the goal with Equation~\eqref{eq:rtt-at-init},
    \begin{equation*}
      \fbox{$\begin{split}
          \mathtt{false}=& ((\sigma_0(id_p)(\texttt{oat})[j]=\mathtt{basic}+\sigma_0(id_p)(\texttt{oat})[j]=\mathtt{test}) \\
          & .(\sigma_0(id_p)(\texttt{sm})-\sigma_0(id_p)(\texttt{sots})<\sigma_0(id_p)(\texttt{oaw})[j])\\
          & .(\sigma_0(id_p)(\texttt{sots})>0))\\
          & +(\sigma_0(id_p)(\texttt{otf})[j]) \\
        \end{split}$}
    \end{equation*}

    By construction, there exists an $id_{fj}\in{}Sigs(\Delta)$ s.t.
    ${<}$\texttt{fired}$\Rightarrow{}\mathtt{id_{fj}}{>}\in{}o_t$
    and
    ${<}$\texttt{output\_transitions\_fired(j)}$\Rightarrow{}\mathtt{id_{fj}}{>}\in{}i_p$.

    By property of the initialization relation,
    ${<}$\texttt{fired}$\Rightarrow{}\mathtt{id_{fj}}{>}\in{}o_t$
    and
    ${<}$\texttt{output\_transitions\_fired(j)}$\Rightarrow{}\mathtt{id_{fj}}{>}\in{}i_p$,
    we can deduce
    $\sigma_0(id_p)(\texttt{otf})[j]=\sigma_0(id_{fj})=\sigma_0(id_t)(\texttt{fired})$.
    
    Appealing to Lemma~\ref{lem:init-states-fired-false}, we can
    deduce $\sigma_0(id_t)(\texttt{fired})=\mathtt{false}$ and consequently\\
    $\sigma_0(id_p)(\texttt{otf})[j]=\mathtt{false}$.

    Rewriting $\sigma_0(id_p)(\texttt{otf})[j]$ as $\mathtt{false}$ and
    simplifying the goal,
    \begin{equation*}
      \fbox{$\begin{split}
          false=& ((\sigma_0(id_p)(\texttt{oat})[j]=\mathtt{BASIC}+\sigma_0(id_p)(\texttt{oat})[j]=\mathtt{TEST}) \\
          & .(\sigma_0(id_p)(\texttt{sm})-\sigma_0(id_p)(\texttt{sots})<\sigma_0(id_p)(\texttt{oaw})[j])\\
          & .(\sigma_0(id_p)(\texttt{sots})>0))\\
        \end{split}$}
    \end{equation*}

    Appealing to Lemma~\ref{lem:init-states-sots-zero}, we can deduce
    $\sigma_0(id_p)(\texttt{sots})=0$.

    Rewriting $\sigma_0(id_p)(\texttt{sots})$ as $0$ and simplifying the
    goal, \qedbox{tautology.}
    
  \end{itemize}
\end{niproof}

\subsection{First rising edge and action executions}
\label{sec:fst-re-actions-ex}

\begin{lemma}[First rising edge equal action executions]
  \label{lem:fst-re-equal-action-ex}
  \fstrehyps{} then\\
  $\forall{}a\in\mathcal{A},id_a\in{}Outs(\Delta)~s.t.~\gamma(a)=id_a,~s_0.ex(a)=\sigma(id_a)$.\\
\end{lemma}

\begin{niproof}
  Given an $a\in\mathcal{A}$ and an
  $id_a\in{}Outs(\Delta)~s.t.~\gamma(a)=id_a$, let us show that
  \fbox{$s_0.ex(a)=\sigma(id_a)$.}

  By construction, $id_a$ is an output port identifier of Boolean type
  in the \hvhdl{} design $d$. The generated \texttt{action} process
  assigns a value to the output port $id_a$ only during the
  initialization phase or a falling edge phase.

  \noindent{}By property of the $\mathtt{Inject}$, \hvhdl{}
  rising edge and stabilize relations, we can deduce
  $\sigma(id_a)=\sigma_0(id_a)$.
  
  Rewriting $\sigma(id_a)$ as $\sigma_0(id_a)$,
  \fbox{$s_0.ex(a)=\sigma_0(id_a)$.}  Appealing to
  Lemma~\ref{lem:init-states-act-exec},
  \qedbox{$s_0.ex(a)=\sigma_0(id_a)$.}
  
\end{niproof}

\subsection{First rising edge and function executions}
\label{sec:fst-re-fun-ex}

\begin{lemma}[First rising edge equal function executions]
  \label{lem:fst-re-equal-fun-ex}
  \fstrehyps{} then\\
  $\forall{}f\in\mathcal{F},id_f\in{}Outs(\Delta)~s.t.~\gamma(f)=id_f,~s_0.ex(f)=\sigma(id_f)$.
\end{lemma}

\begin{niproof}
  Given an $f\in\mathcal{F}$ and an $id_f\in{}Outs(\Delta)$
  s.t. $\gamma(f)=id_f$, let us show that
  \fbox{$s_0.ex(f)=\sigma(id_f)$.}

  Rewriting $s_0.ex(f)$ as $\mathtt{false}$, by definition of $s_0$,
  \fbox{$\sigma(id_f)=\mathtt{false}$.}

  By construction, $id_f$ is an output port identifier of Boolean type
  in the \hvhdl{} design $d$. The generated \texttt{function} process
  assigns a value to the output port $id_f$ only during the
  initialization phase or during a rising edge phase.
  
  By construction, the \texttt{function} process is defined in the
  behavior of design $d$, i.e.\\
  $\mathtt{ps}(\texttt{function}, \emptyset, sl, ss)\in{}d.cs$.
  
  Let $trs(f)$ be the set of transitions associated to function $f$,
  i.e $trs(f)=\{t\in{}T~\vert~\mathbb{F}(t,f)=true\}$.

  Let us perform case analysis on $trs(f)$; there are two cases:
  
  \begin{itemize}
  \item \textbf{CASE} $trs(f)=\emptyset$:
    
    By construction,
    $\mathtt{id_f\Leftarrow{}false}\in{}ss_{\uparrow}$ where
    $ss_\uparrow$ is the part of the \texttt{``function''} process
    body executed during a rising edge phase (i.e. a rising edge block
    statement).

    By property of the \hvhdl{} rising edge and the stabilize
    relation, \qedbox{$\sigma(id_f)=\mathtt{false}$.}
    
  \item \textbf{CASE} $trs(f)\neq\emptyset$:
    
    By construction,
    $\mathtt{id_f\Leftarrow{}id_{ft_0}+\dots+id_{ft_n}}\in{}ss_\uparrow$
    where $ss_\uparrow$ is the part of the \texttt{``function''}
    process body executed during the rising edge phase, and
    $n=\vert{}trs(f)\vert-1$, and for all $i\in[0,n-1]$, $id_{ft_i}$
    is a internal signal of design $d$.

    By property of the $\mathtt{Inject}$, the \hvhdl{} rising
    edge and stabilize relations, we can deduce
    $\sigma(id_f)=\sigma_0(id_{ft_0})+\dots+\sigma_0(id_{ft_n})$.

    Rewriting $\sigma(id_f)$ as
    $\sigma_0(id_{ft_0})+\dots+\sigma_0(id_{ft_n})$,
    \fbox{$\sigma_0(id_{ft_0})+\dots+\sigma_0(id_{ft_n})=\mathtt{false}$.}

    By construction, for all $id_{ft_i}$, there exist a
    $t_i\in{}trs(f)$ and an $id_{t_i}$ s.t. $\gamma(t_i)=id_{t_i}$.

    By construction and by definition of $id_{t_i}$, there exist
    $g_{t_i}$, $i_{t_i}$ and $o_{t_i}$ s.t.
    $\mathtt{comp}($$id_{t_i},$ $\texttt{transition},$
    $g_{t_i},$ $i_{t_i},$ $o_{t_i})\in{}d.cs$.

    By construction, we have
    ${<}\mathtt{fired\Rightarrow{id_{ft_i}}}{>}\in{}o_{t_i}$, and by
    property of the initialization relation, we have
    $\sigma_0(id_{ft_i})=\sigma_0(id_{t_i})(\texttt{fired})$.

    \noindent{}Rewriting $\sigma_0(id_{ft_i})$ as
    $\sigma_0(id_{t_i})(\texttt{fired})$,
    \fbox{$\sigma_0(id_{t_0})(\texttt{fired})+\dots+\sigma_0(id_{t_n})(\texttt{fired})=\mathtt{false}$.}

    Appealing to Lemma~\ref{lem:init-states-fired-false}, we can
    deduce $\sigma_0(id_{t_i})(\texttt{fired})=\mathtt{false}$.

    Rewriting all $\sigma_0(id_{t_i})(\texttt{fired})$ as $\mathtt{false}$
    and simplifying the goal, \qedbox{tautology.}
    
  \end{itemize}

\end{niproof}

\subsection{First rising edge and sensitization}
\label{sec:fst-re-sens}

\begin{lemma}[First rising edge equal sensitized]
  \label{lem:fst-re-equal-sens}
  \fstrehyps{} then\\
  $\forall{}t\in{}T,id_t\in{}Comps(\Delta)~s.t.~\gamma(t)=id_t,$
  $t\in{}Sens(s_0.M)\Leftrightarrow\sigma(id_t)(\texttt{s\_enabled})=\mathtt{true}$.
\end{lemma}

\begin{niproof}
  See the proof of Lemma~\ref{lem:re-equal-sens}.
\end{niproof}

\begin{lemma}[First rising edge not equal sensitized]
  \label{lem:fst-re-neq-sens}
  \fstrehyps{} then\\
  $\forall{}t\in{}T,id_t\in{}Comps(\Delta)~s.t.~\gamma(t)=id_t,$
  $t\notin{}Sens(s_0.M)\Leftrightarrow\sigma(id_t)(\texttt{s\_enabled})=\mathtt{false}$.
\end{lemma}

\begin{niproof}
  See the proof of Lemma~\ref{lem:re-equal-not-sens}.
\end{niproof}

\subsection{First rising edge and conditions}
\label{sec:fst-re-cond-comb}

\begin{lemma}[First rising edge equal condition combination]
  \label{lem:fst-re-equal-cond-comb}
  \fstrehyps{} then\\
  $\forall{}t\in{}T,id_t\in{}Comps(\Delta)~s.t.~\gamma(t)=id_t,$\\
  $\sigma(id_t)(\texttt{s\_condition\_combination})=
  \prod\limits_{c\in{}conds(t)}
  \begin{cases}
    E_c(\tau,c) & if~\mathbb{C}(t,c)=1 \\
    \mathtt{not}(E_c(\tau,c)) & if~\mathbb{C}(t,c)=-1 \\
  \end{cases}$\\
  where
  $conds(t)=\{c\in\mathcal{C}~\vert~\mathbb{C}(t,c)=1\lor\mathbb{C}(t,c)=-1\}$.
\end{lemma}

\begin{niproof}
  See the proof of Lemma~\ref{lem:re-equal-cond-comb}.
\end{niproof}

\begin{lemma}[First rising edge equal conditions]
  \label{lem:fst-re-equal-cond}
  \fstrehyps{} then\\
  $\forall{}c\in\mathcal{C},id_c\in{}Ins(\Delta)$
  s.t. $\gamma(c)=id_c$, $\sigma(id_c)=E_c(\tau,c)$.
\end{lemma}

\begin{niproof}
  See the proof of Lemma~\ref{lem:re-equal-cond}.
\end{niproof}


%%% Local Variables:
%%% mode: latex
%%% TeX-master: "../../main"
%%% End:


\section{Rising Edge}
\label{sec:rising-edge}
%%%%%%%%%%%%%%%%%%%%%%%%%%%%%%%%%%%%%%%%%%%%
%%%%%%%%%% RISING EDGE HYPOTHESES %%%%%%%%%%
%%%%%%%%%%%%%%%%%%%%%%%%%%%%%%%%%%%%%%%%%%%%

\begin{definition}[Rising edge hypotheses]
  \label{def:re-hyps}
  Given an $sitpn\in{}SITPN$, $b\in{}P\rightarrow\mathbb{N}$,
  $d\in{}design$, $\gamma\in{}WM(sitpn,d)$,
  $E_c\in\mathbb{N}\rightarrow\mathcal{C}\rightarrow\mathbb{B}$,
  $\Delta\in{}ElDesign$,
  $E_p\in\mathbb{N}\rightarrow{}Ins(\Delta)\rightarrow{}value$,
  $\tau\in\mathbb{N}$, $s,s'\in{}S(sitpn)$,
  $\sigma_e,\sigma,\sigma_i,\sigma_\uparrow,\sigma'\in\Sigma$, assume
  that:
  \begin{itemize}
  \item $\lfloor{}sitpn\rfloor_b=(d,\gamma)$ and
    $\gamma\vdash{}E_p\stackrel{env}{=}E_c$ and
    $\mathcal{D}_\mathcal{H},\emptyset\vdash\mathrm{d}\srarrow{elab}{\fontsize{7}{7}\selectfont}\Delta,\sigma_e$
  \item $\gamma\vdash{}s\stackrel{\downarrow}{\approx}\sigma$
  \item
    $E_c,\tau\vdash{}s\srarrow{\uparrow}{\fontsize{7}{7}\selectfont}s'$
  \item $\mathtt{Inject}(\sigma, E_p, \tau, \sigma_i)$ and
    $\mathcal{D}_{\mathcal{H}},\Delta,\sigma_i\vdash\mathrm{d.cs}\xrightarrow{\uparrow}\sigma_\uparrow$
    and
    $\mathcal{D}_{\mathcal{H}},\Delta,\sigma_\uparrow\vdash\mathrm{d.cs}\xrightarrow{\rightsquigarrow}\sigma'$
  \item State $\sigma$ is a stable design state:
    $\mathcal{D}_{\mathcal{H}},\Delta,\sigma\vdash\mathrm{d.cs}\xrightarrow{comb}\sigma$
  \end{itemize}
\end{definition}

\def\rehyps{For all $sitpn$, $b$, $d$, $\gamma$, $E_c$, $E_p$, $\tau$,
  $\Delta$, $\sigma_e$, $s$, $s'$, $\sigma$, $\sigma_i$,
  $\sigma_\uparrow$, $\sigma'$ that verify the hypotheses of
  Definition~\ref{def:re-hyps},}

\subsection{Rising edge and Marking}
\label{sec:re-marking}

%%%%%%%%%%%%%%%%%%%%%%%%%%%%%%%%%%%%%%%%%%%%%%%%%%%%%
%%%%%%%%%% RISING EDGE EQUAL MARKING LEMMA %%%%%%%%%%
%%%%%%%%%%%%%%%%%%%%%%%%%%%%%%%%%%%%%%%%%%%%%%%%%%%%%

\begin{lemma}[Rising edge equal marking]
  \label{lem:re-equal-marking}
  \rehyps{} then $\forall{}p,id_p~s.t.~\gamma(p)=id_p$,
  $s'.M(p)=\sigma'(id_p)(\texttt{s\_marking})$.
\end{lemma}

\begin{niproof}

  Given a $p\in{}P$, let us show
  \fbox{$s'.M(p)=\sigma'(id_p)(\texttt{s\_marking})$.}

  \exP{}
  
  By definition of the SITPN state transition relation on rising edge:
  \begin{equation}\label{eq:re-eq-marking-eqmp}
    s'.M(p)=s.M(p)-\sum\limits_{t\in{}Fired(s)}pre(p,t)+\sum\limits_{t\in{}Fired(s)}post(t,p)
  \end{equation}

  By property of the $\mathtt{Inject}$, the \hvhdl{} rising
  edge and the stabilize relations, \InCsCompP{}, and through the
  examination of the \texttt{marking} process defined in the place
  design architecture, we can deduce:
  \begin{equation}\label{eq:re-eq-marking-eqsm}
    \begin{split}
      \sigma'(id_p)(\texttt{sm})=\sigma(id_p)(\texttt{sm})-\sigma(id_p)(\texttt{s\_output\_token\_sum})\\
      +\sigma(id_p)(\texttt{s\_input\_token\_sum})
    \end{split}
  \end{equation}

  Rewriting the goal with \ref{eq:re-eq-marking-eqmp} and
  \ref{eq:re-eq-marking-eqsm},
  
  \fbox{
    \begin{tabular}{@{}c@{}}
      $s.M(p)-\sum\limits_{t\in{}Fired(s)}pre(p,t)+\sum\limits_{t\in{}Fired(s)}post(t,p)$\\
      $=$ \\
      $\sigma(id_p)(\texttt{sm})-\sigma(id_p)(\texttt{sots})
      +\sigma(id_p)(\texttt{sits})$ \\
    \end{tabular}
  }\\
  
  By definition of the \nameref{def:full-post-fe-state-sim} relation,
  we can deduce $s.M(p)=\sigma(id_p)(\texttt{sm})$,
  $\sum\limits_{t\in{}Fired(s)}pre(p,t)=\sigma(id_p)(\texttt{sots})$ and
  $\sum\limits_{t\in{}Fired(s)}post(t,p)=\sigma(id_p)(\texttt{sits})$, and thus, 

  \qedbox{
    \begin{tabular}{c}
      $s.M(p)-\sum\limits_{t\in{}Fired(s)}pre(p,t)+\sum\limits_{t\in{}Fired(s)}post(t,p)$ \\
      $=$ \\
      $\sigma(id_p)(\texttt{sm})-\sigma(id_p)(\texttt{sots})+\sigma(id_p)(\texttt{sits})$ \\
    \end{tabular}
  }
\end{niproof}

%%%%%%%%%%%%%%%%%%%%%%%%%%%%%%%%%%%%%%%%%%%%%%%%%%%%%%%%%%%%%%
%%%%%%%%%% RISING EDGE EQUAL CONDITION COMBINATION  %%%%%%%%%%
%%%%%%%%%%%%%%%%%%%%%%%%%%%%%%%%%%%%%%%%%%%%%%%%%%%%%%%%%%%%%%

\subsection{Rising edge and conditions}
\label{sec:re-cond-comb}

\begin{lemma}[Rising edge equal condition combination]
  \label{lem:re-equal-cond-comb}
  \rehyps{} then\\
  $\forall{}t\in{}T,id_t\in{}Comps(\Delta)~s.t.~\gamma(t)=id_t,$\\
  $\sigma'(id_t)(\texttt{s\_condition\_combination})=
  \prod\limits_{c\in{}conds(t)}
  \begin{cases}
    E_c(\tau,c) & if~\mathbb{C}(t,c)=1 \\
    \mathtt{not}(E_c(\tau,c)) & if~\mathbb{C}(t,c)=-1 \\
  \end{cases}$\\
  where
  $conds(t)=\{c\in\mathcal{C}~\vert~\mathbb{C}(t,c)=1\lor\mathbb{C}(t,c)=-1\}$.
\end{lemma}

\begin{niproof}
  Given a $t$ and an $id_t$ s.t. $\gamma(t)=id_t$, let us show\\
  \fbox{$\sigma'(id_t)(\texttt{s\_condition\_combination})=
    \prod\limits_{c\in{}conds(t)}
    \begin{cases}
      E_c(\tau,c) & if~\mathbb{C}(t,c)=1 \\
      \mathtt{not}(E_c(\tau,c)) & if~\mathbb{C}(t,c)=-1 \\
    \end{cases}$.}\\

  \exT

  \noindent By property of the \hvhdl{} stabilize relation,
  \InCsCompT{}, and through the examination of the
  \texttt{condition_evaluation} process defined in the transition
  design architecture, we can deduce:
  \begin{equation}
    \sigma'(id_t)(\texttt{scc})=\prod\limits_{i=0}^{\Delta(id_t)(\texttt{conditions\_number})-1}\sigma'(id_t)(\texttt{input\_conditions})[i]\label{eq:re-eq-cc-eqscc}
  \end{equation}

  \noindent{}Rewriting the goal with \ref{eq:re-eq-cc-eqscc},\\
  \fbox{$\prod\limits_{i=0}^{\Delta(id_t)(\texttt{cn})-1}\sigma'(id_t)(\texttt{ic})[i]=
    \prod\limits_{c\in{}conds(t)}
    \begin{cases}
      E_c(\tau,c) & if~\mathbb{C}(t,c)=1 \\
      \mathtt{not}(E_c(\tau,c)) & if~\mathbb{C}(t,c)=-1 \\
    \end{cases}$.}\\

  \noindent{}Let us perform case analysis on $conds(t)$; there are two cases:

  \begin{itemize}
  \item \textbf{CASE} $conds(t)=\emptyset$:
    \fbox{$\prod\limits_{i=0}^{\Delta(id_t)(\texttt{cn})-1}\sigma'(id_t)(\texttt{ic})[i]=\mathtt{true}$.}\\
    
    By construction, ${<}\mathtt{cn\Rightarrow{}1}{>}\in{}g_t$ and
    ${<}\mathtt{ic(0)\Rightarrow{}true}{>}\in{}i_t$.

    By property of the stabilize relation,
    ${<}\mathtt{cn\Rightarrow{}1}{>}\in{}g_t$ and
    ${<}$\texttt{ic(0)}$\mathtt{\Rightarrow{}true}{>}\in{}i_t$, we can
    deduce $\Delta(id_t)(\texttt{cn})=1$ and
    $\sigma'(id_t)(\texttt{ic})[0]=\mathtt{true}$.

    Rewriting the goal with $\Delta(id_t)(\texttt{cn})=1$ and
    $\sigma'(id_t)(\texttt{ic})[0]=\mathtt{true}$, \qedbox{tautology.}
    
  \item \textbf{CASE} $conds(t)\neq\emptyset$:\\
    By construction,
    ${<}\mathtt{cn\Rightarrow{}\vert{}conds(t)\vert}{>}\in{}g_t$, and
    by property of the stabilize relation, we can deduce
    $\Delta(id_t)(\texttt{cn})=\vert{}conds(t)\vert$.
    
    Rewriting the goal with $\Delta(id_t)(\texttt{cn})=\vert{}conds(t)\vert$:\\
    \fbox{$\prod\limits_{i=0}^{\vert{}conds(t)\vert-1}\sigma'(id_t)(\texttt{ic})[i]=
      \prod\limits_{c\in{}conds(t)}
      \begin{cases}
        E_c(\tau,c) & if~\mathbb{C}(t,c)=1 \\
        \mathtt{not}(E_c(\tau,c)) & if~\mathbb{C}(t,c)=-1 \\
      \end{cases}$}\\

    There exists a mapping, given by the transformation function,
    between the set $conds(t)$ and the indexes of
    $[0,\vert{}conds(t){}\vert-1]$.

    Let $\beta\in{}conds(t)\rightarrow[0,\vert{}conds(t){}\vert-1]$ be
    this mapping.

    To prove the current goal, it suffices to prove that for all
    condition $c\in{}conds(t)$, we have

    \begin{frameb}
      $\bigg(\begin{cases}
        E_c(\tau,c) & if~\mathbb{C}(t,c)=1 \\
        \mathtt{not}(E_c(\tau,c)) & if~\mathbb{C}(t,c)=-1 \\
      \end{cases}\bigg)=\sigma'(id_t)(\texttt{ic})[\beta(c)]$
    \end{frameb}

    Given a $c\in{}conds(t)$, let us show the above goal.
    
    By construction, for all $c\in{}conds(t)$, there exists an
    $id_c\in{}Ins(\Delta)$ such that
    \begin{itemize}
    \item $\gamma(c)=id_c$
    \item $\mathbb{C}(t,c)=1$ implies
      ${<}\mathtt{ic(}\beta(c)\mathtt{)\Rightarrow{}id_c}{>}\in{}i_t$
    \item $\mathbb{C}(t,c)=-1$ implies
      ${<}\mathtt{ic(}\beta(c)\mathtt{)\Rightarrow{}not~id_c}{>}\in{}i_t$
    \end{itemize}

    Let us take such an $id_c$ with the above properties.

    By definition of $c\in{}conds(t)$, we have
    $\mathbb{C}(t,c)=1\lor\mathbb{C}(t,c)=-1$. Let us perform case
    analysis on $\mathbb{C}(t,c)=1\lor\mathbb{C}(t,c)=-1$:

    \begin{itemize}
    \item \textbf{CASE} $\mathbb{C}(t,c)=1$:

      In that case, we must show:
      \fbox{$E_c(\tau,c)=\sigma'(id_t)(\texttt{ic})[\beta(c)]$}

      By assumption, we have
      ${<}\mathtt{ic(}\beta(c)\mathtt{)\Rightarrow{}id_c}{>}\in{}i_t$
      and by property of the stabilize relation, we can deduce
      $\sigma(id_t)(\texttt{ic})[\beta(c)]=\sigma'(id_c)$.

      Rewriting the goal with
      $\sigma(id_t)(\texttt{ic})[\beta(c)]=\sigma'(id_c)$:

      \fbox{$E_c(\tau,c)=\sigma'(id_c)$}

      By property of the $\mathtt{Inject}$ relation and
      $id_c\in{}Ins(\Delta)$, we can deduce
      $\sigma'(id_c)=E_p(\tau)(id_c)$.

      By property of $\gamma\vdash{}E_p\stackrel{env}{=}E_c$, we can
      deduce $E_p(\tau)(id_c)=E_c(\tau,c)$.

      Rewriting the goal with $\sigma'(id_c)=E_p(\tau)(id_c)$ and
      $E_p(\tau)(id_c)=E_c(\tau,c)$:
      \fbox{$E_c(\tau,c)=E_c(\tau,c)$}, then \qedbox{tautology.}

    \item \textbf{CASE} $\mathbb{C}(t,c)=-1$:

      In that case, we must show:
      \fbox{$\mathtt{not}~E_c(\tau,c)=\sigma'(id_t)(\texttt{ic})[\beta(c)]$}

      By assumption, we have
      ${<}\mathtt{ic(}\beta(c)\mathtt{)\Rightarrow{}not~id_c}{>}\in{}i_t$
      and by property of the stabilize relation, we can deduce
      $\sigma(id_t)(\texttt{ic})[\beta(c)]=\mathtt{not}~\sigma'(id_c)$.

      Rewriting the goal with
      $\sigma(id_t)(\texttt{ic})[\beta(c)]=\mathtt{not}~\sigma'(id_c)$:

      \fbox{$\mathtt{not}~E_c(\tau,c)=\mathtt{not}~\sigma'(id_c)$}

      By property of the $\mathtt{Inject}$ relation and
      $id_c\in{}Ins(\Delta)$, we can deduce
      $\sigma'(id_c)=E_p(\tau)(id_c)$.

      By property of $\gamma\vdash{}E_p\stackrel{env}{=}E_c$, we can
      deduce $E_p(\tau)(id_c)=E_c(\tau,c)$.

      Rewriting the goal with $\sigma'(id_c)=E_p(\tau)(id_c)$ and
      $E_p(\tau)(id_c)=E_c(\tau,c)$:
      \fbox{$\mathtt{not}~E_c(\tau,c)=\mathtt{not}~E_c(\tau,c)$}, then
      \qedbox{tautology.}
      
    \end{itemize}
    
  \end{itemize}
  
\end{niproof}

\begin{lemma}[Rising edge equal conditions]
  \label{lem:re-equal-cond}
  \rehyps{} then\\
  $\forall{}c\in\mathcal{C},id_c\in{}Ins(\Delta)$
  s.t. $\gamma(c)=id_c$, $\sigma'(id_c)=E_c(\tau,c)$.
\end{lemma}

\begin{niproof}
  Given a $c\in\mathcal{C}$ and an $id_c\in{}Ins(\Delta)$ such that
  $\gamma(c)=id_c$, let us show
  \begin{frameb}
    $\sigma'(id_c)=E_c(\tau,c)$
  \end{frameb}

  By property of the $\mathtt{Inject}$ relation and
  $id_c\in{}Ins(\Delta)$, we can deduce
  $\sigma'(id_c)=E_p(\tau)(id_c)$.

  By property of $\gamma\vdash{}E_p\stackrel{env}{=}E_c$, we can
  deduce $E_p(\tau)(id_c)=E_c(\tau,c)$.

  Rewriting the goal with $\sigma'(id_c)=E_p(\tau)(id_c)$ and
  $E_p(\tau)(id_c)=E_c(\tau,c)$, \qedbox{tautology}.
\end{niproof}

\subsection{Rising edge and time counters}
\label{sec:re-tc}

\begin{lemma}[Rising edge equal time counters]
  \label{lem:re-equal-tc}
  \rehyps{} then\\
  $\forall{}t\in{}T_i,id_t\in{}Comps(\Delta)~s.t.~\gamma(t)=id_t$,\\
  $\big(u(I_s(t))=\infty\land{}s'.I(t)\le{}l(I_s(t))\Rightarrow$
  $s'.I(t)=\sigma'(id_t)(\texttt{s\_time\_counter})\big)$\\
  $\land\big(u(I_s(t))=\infty\land{}s'.I(t)>{}l(I_s(t))\Rightarrow$
  $\sigma'(id_t)(\texttt{s\_time\_counter})=l(I_s(t))\big)$\\
  $\land\big(u(I_s(t))\neq\infty\land{}s'.I(t)>{}u(I_s(t))\Rightarrow$
  $\sigma'(id_t)(\texttt{s\_time\_counter})=u(I_s(t))\big)$\\
  $\land\big(u(I_s(t))\neq\infty\land{}s'.I(t)\le{}u(I_s(t))\Rightarrow$
  $s'.I(t)=\sigma'(id_t)(\texttt{s\_time\_counter})\big)$.
\end{lemma}

\begin{niproof}
  Given a $t\in{}T_i$ and an $id_t\in{}Comps(\Delta)$ s.t. $\gamma(t)=id_t$, let us show\\
  \noindent\fbox{\parbox{\lwidth}{$\big(u(I_s(t))=\infty\land{}s'.I(t)\le{}l(I_s(t))\Rightarrow$
      $s'.I(t)=\sigma'(id_t)(\texttt{s\_time\_counter})\big)$\\
      $\land\big(u(I_s(t))=\infty\land{}s'.I(t)>{}l(I_s(t))\Rightarrow$
      $\sigma'(id_t)(\texttt{s\_time\_counter})=l(I_s(t))\big)$\\
      $\land\big(u(I_s(t))\neq\infty\land{}s'.I(t)>{}u(I_s(t))\Rightarrow$
      $\sigma'(id_t)(\texttt{s\_time\_counter})=u(I_s(t))\big)$\\
      $\land\big(u(I_s(t))\neq\infty\land{}s'.I(t)\le{}u(I_s(t))\Rightarrow$
      $s'.I(t)=\sigma'(id_t)(\texttt{s\_time\_counter})\big)$}}\\

  \exT{}

  Then, there are 4 points to show:

  \begin{enumerate}
  \item\label{it:re-eq-tc-fst}
    \fbox{$u(I_s(t))=\infty\land{}s'.I(t)\le{}l(I_s(t))\Rightarrow{}s'.I(t)=\sigma'(id_t)(\texttt{s\_time\_counter})$}\\
    
    Assuming that $u(I_s(t))=\infty$ and
    $s'.I(t)\le{}l(I_s(t))$, let us show\\
    \fbox{$s'.I(t)=\sigma'(id_t)(\texttt{s\_time\_counter})$.}

    By property of the $\mathtt{Inject}$, \hvhdl{} rising edge and
    stabilize relations, \InCsCompT, and through the examination of
    the \texttt{time_counter} process defined in the transition design
    architecture, we can deduce
    $\sigma'(id_t)(\texttt{stc})=\sigma(id_t)(\texttt{stc})$.

    By property of
    $\gamma\vdash{}s\stackrel{\downarrow}{\approx}\sigma$, we can
    deduce $s.I(t)=\sigma(id_t)(\texttt{stc})$.

    Rewriting the goal with $\sigma'(id_t)(\texttt{stc})=\sigma(id_t)(\texttt{stc})$
    and $s.I(t)=\sigma(id_t)(\texttt{stc})$, \qedbox{tautology.}
    
  \item
    \fbox{$u(I_s(t))=\infty\land{}s'.I(t)>{}l(I_s(t))\Rightarrow
      \sigma'(id_t)(\texttt{s\_time\_counter})=l(I_s(t)$.}

    \noindent{}Proved in the same fashion as \ref{it:re-eq-tc-fst}.
  \item
    \fbox{$u(I_s(t))\neq\infty\land{}s'.I(t)>{}u(I_s(t))\Rightarrow
      \sigma'(id_t)(\texttt{s\_time\_counter})=u(I_s(t)$.}

    \noindent{}Proved in the same fashion as \ref{it:re-eq-tc-fst}.
    
  \item
    \fbox{$u(I_s(t))\neq\infty\land{}s'.I(t)\le{}u(I_s(t))\Rightarrow{}s'.I(t)=\sigma'(id_t)(\texttt{s\_time\_counter})$}
    
    \noindent{}Proved in the same fashion as \ref{it:re-eq-tc-fst}.
  \end{enumerate}
  
\end{niproof}

\subsection{Rising edge and reset orders}
\label{sec:re-reset-orders}

\begin{lemma}[Rising edge equal reset orders]
  \label{lem:re-equal-reset-orders}
  \rehyps{} then\\
  $\forall{}t\in{}T_i,id_t\in{}Comps(\Delta)~s.t.~\gamma(t)=id_t,$
  $s'.reset_t(t)=\sigma'(id_t)(\texttt{s\_reinit\_time\_counter})$
\end{lemma}

\begin{niproof}
  Given a $t\in{}T_i$ and an $id_t\in{}Comps(\Delta)$
  s.t. $\gamma(t)=id_t$, let us show\\
  \fbox{$s'.reset_t(t)=\sigma'(id_t)(\texttt{s\_reinit\_time\_counter})$.}

  \exT

  By property of the \hvhdl{} stabilize relation, \InCsCompT, and
  through the examination of the
  \texttt{reinit_time_counter_evaluation} process defined in the
  transition design architecture, we can deduce:
  \begin{equation}
    \sigma'(id_t)(\texttt{srtc})=\sum\limits_{i=0}^{\Delta(id_t)(\texttt{input\_arcs\_number})-1}\sigma'(id_t)(\texttt{reinit\_time})[i]\label{eq:eq-srtc-prod}
  \end{equation}

  Rewriting the goal with \eqref{eq:eq-srtc-prod},
  \fbox{$s'.reset_t(t)=\sum\limits_{i=0}^{\Delta(id_t)(\texttt{ian})-1}\sigma'(id_t)(\texttt{rt})[i]$.}
  
  Let us perform case analysis on $input(t)$; there are two cases:

  \begin{itemize}
  \item \textbf{CASE} $input(t)=\emptyset$:

    By construction,
    ${<}\mathtt{input\_arcs\_number\Rightarrow{}1}{>}\in{}g_t$, and
    by property of the elaboration relation, we can deduce
    $\Delta(id_t)(\texttt{ian})=1$.

    By construction, there exists an $id_{ft}\in{}Sigs(\Delta)$ s.t.
    ${<}\mathtt{reinit\_time(0)\Rightarrow{}id_{ft}}{>}\in{}i_t$ and
    ${<}\mathtt{fired\Rightarrow{}id_{ft}}{>}\in{}o_t$, and by
    property of the stabilize relation and \InCsCompT, we can deduce
    $\sigma'(id_t)(\texttt{rt})[0]=\sigma'(id_{ft})=\sigma'(id_t)(\texttt{fired})$.

    Rewriting the goal with $\Delta(id_t)(\texttt{ian})=1$ and
    $\sigma'(id_t)(\texttt{rt})[0]=\sigma'(id_{ft})=\sigma'(id_t)(\texttt{fired})$:
    \fbox{$s'.reset_t(t)=\sigma'(id_t)(\texttt{fired})$.}
    
    By property of the stabilize relation, \InCsCompT{}, and through
    the examination of the \texttt{fired_evaluation} process, we can deduce:
    \begin{equation}
      \sigma'(id_t)(\texttt{fired})=\sigma'(id_t)(\texttt{s\_firable})~.~\sigma'(id_t)(\texttt{s\_priority\_combination})\label{eq:fired-at-re}
    \end{equation}

    Rewriting the goal with \eqref{eq:fired-at-re}:\\
    \fbox{$s'.reset_t(t)=\sigma'(id_t)(\texttt{s\_firable})~.~\sigma'(id_t)(\texttt{s\_priority\_combination})$.}

    By property of the stabilize relation, \InCsCompT, and through the
    examination of the \texttt{priority_authorization_evaluation}
    process defined in the transition design architecture, we can
    deduce:
    \begin{equation}
      \sigma'(id_t)(\texttt{spc})=\prod\limits_{i=0}^{\Delta(id_t)(\texttt{ian})-1}\sigma'(id_t)(\texttt{priority\_authorizations})[i]\label{eq:eq-spc-sum}
    \end{equation}

    As $\Delta(id_t)(\texttt{ian})=1$, we can deduce
    $\prod\limits_{i=0}^{\Delta(id_t)(\texttt{ian})-1}\sigma'(id_t)(\texttt{pauths})[i]=\sigma'(id_t)(\texttt{pauths})[0]$.

    Rewriting the goal with \eqref{eq:eq-spc-sum} and
    $\prod\limits_{i=0}^{\Delta(id_t)(\texttt{ian})-1}\sigma'(id_t)(\texttt{pauths})[i]=\sigma'(id_t)(\texttt{pauths})[0]$:\\
    \fbox{$s'.reset_t(t)=\sigma'(id_t)(\texttt{s\_firable})~.~\sigma'(id_t)(\texttt{pauths})[0]$.}
    
    By construction,
    ${<}\mathtt{priority\_authorizations(0)\Rightarrow{}true}{>}\in{}i_t$,
    and by property of the stabilize relation and \InCsCompT, we can
    deduce\\ $\sigma'(id_t)(\texttt{pauths})[0]=\mathtt{true}$.
    
    Rewriting the goal with $\sigma'(id_t)(\texttt{pauths})[0]=\mathtt{true}$
    , and simplifying the equation:\\
    \fbox{$s'.reset_t(t)=\sigma'(id_t)(\texttt{s\_firable})$.}
    
    Let us perform case analysis on $t\in{}Fired(s)$ or
    $t\notin{}Fired(s)$:

    \begin{itemize}
    \item \textbf{CASE} $t\in{}Fired(s)$:
      
      By property of
      $E_c,\tau\vdash{}s\srarrow{\uparrow}{\fontsize{7}{7}\selectfont}s'$
      (Rule~\ref{it:reset-order}), we can deduce
      $s'.reset_t(t)=\mathtt{true}$.

      Rewriting the goal with $s'.reset_t(t)=\mathtt{true}$:
      \fbox{$\sigma'(id_t)(\texttt{s\_firable})=\mathtt{true}$.}

      By property of the stabilize, the \hvhdl{} rising edge and the
      $\mathtt{Inject}$ relations, \InCsCompT, and through the
      examination of the \texttt{firable} process defined in the
      transition design architecture, we can deduce\\
      $\sigma(id_t)(\texttt{s\_firable})=\sigma'(id_t)(\texttt{s\_firable})$.

      Rewriting the goal with
      $\sigma(id_t)(\texttt{s\_firable})=\sigma'(id_t)(\texttt{s\_firable})$,
      we have \\
      \fbox{$\sigma(id_t)(\texttt{s\_firable})=\mathtt{true}$.}

      By property of
      $\gamma\vdash{}s\stackrel{\downarrow}{\approx}\sigma$, we can
      deduce
      $t\in{}Firable(s)\Leftrightarrow\sigma(id_t)(\texttt{sfa})=\mathtt{true}$.

      Rewriting the goal with
      $t\in{}Firable(s)\Leftrightarrow\sigma(id_t)(\texttt{sfa})=\mathtt{true}$,
      \fbox{$t\in{}Firable(s)$.}

      By property of $t\in{}Fired(s)$, \qedbox{$t\in{}Firable(s)$.}
      
    \item \textbf{CASE} $t\notin{}Fired(s)$:

      \noindent{}By property of $input(t)=\emptyset$, there does not
      exist any input place connected to $t$ by a $\mathtt{basic}$ or
      $\mathtt{test}$ arc. Thus, by property of
      $E_c,\tau\vdash{}s\srarrow{\uparrow}{\fontsize{7}{7}\selectfont}s'$
      (Rule~\ref{it:reset-order}), we can deduce
      $s'.reset_t(t)=\mathtt{false}$.

      Rewriting the goal with $s'.reset_t(t)=\mathtt{false}$:
      \fbox{$\sigma'(id_t)(\texttt{s\_firable})=\mathtt{false}$.}

      By property of the stabilize, the \hvhdl{} rising edge and the
      $\mathtt{Inject}$ relations, \InCsCompT, and through
      the examination of the \texttt{firable} process defined in the
      transition design architecture, we can deduce
      $\sigma(id_t)(\texttt{sfa})=\sigma'(id_t)(\texttt{sfa})$.

      Rewriting the goal with
      $\sigma(id_t)(\texttt{sfa})=\sigma'(id_t)(\texttt{sfa})$,
      \fbox{$\sigma(id_t)(\texttt{sfa})=\mathtt{false}$.}

      By property of
      $\gamma\vdash{}s\stackrel{\downarrow}{\approx}\sigma$, we can
      deduce
      $t\notin{}Firable(s)\Leftrightarrow\sigma(id_t)(\texttt{sfa})=\mathtt{false}$.
      
      By property of $t\notin{}Fired(s)$ and $input(t)=\emptyset$,
      \qedbox{$t\notin{}Firable(s)$}.
    \end{itemize}
    
  \item \textbf{CASE} $input(t)\neq{}\emptyset$:

    By construction,
    ${<}\mathtt{input\_arcs\_number\Rightarrow{}}\vert{}input(t)\vert{>}\in{}g_t$,
    and by property of the elaboration relation, we can deduce
    $\Delta(id_t)(\texttt{ian})=\vert{}input(t)\vert$.

    Rewriting the goal with
    $\Delta(id_t)(\texttt{ian})=\vert{}input(t)\vert$,
    \fbox{$s'.reset_t(t)=\sum\limits_{i=0}^{\vert{}input(t)\vert-1}\sigma'(id_t)(\texttt{rt})[i]$.}
    
    Let us perform case analysis on $t\in{}Fired(s)$ or
    $t\notin{}Fired(s)$:
    
    \begin{itemize}
    \item \textbf{CASE} $t\in{}Fired(s)$:
      
      By property of
      $E_c,\tau\vdash{}s\srarrow{\uparrow}{\fontsize{7}{7}\selectfont}s'$
      (Rule~\ref{it:reset-order}), we can deduce
      $s'.reset_t(t)=\mathtt{true}$.

      Rewriting the goal with $s'.reset_t(t)=\mathtt{true}$,
      \fbox{$\sum\limits_{i=0}^{\vert{}input(t)\vert-1}\sigma'(id_t)(\texttt{rt})[i]=\mathtt{true}$.}

      To prove the goal, let us show
      \fbox{$\exists{}i\in[0,\vert{}input(t)\vert-1]$
        s.t. $\sigma'(id_t)(\texttt{rt})[i]=\mathtt{true}$.}

      By construction, and $input(t)\neq{}\emptyset$, there exist
      ${}p\in{}input(t)$ and $id_p\in{}Comps(\Delta)$
      s.t. $\gamma(p)=id_p$.

      \exP{}

      By construction, there exist an
      $i\in{}[0,\vert{}input(t)\vert-1],$ a
      $j\in[0,\vert{}output(p)\vert-1]$ and $id_{ji}\in{}Sigs(\Delta)$
      s.t.
      ${<}\mathtt{reinit\_transition\_time(j)\Rightarrow{}id_{ji}}{>}\in{}o_p$
      and\\
      ${<}\mathtt{reinit\_time(i)\Rightarrow{}id_{ji}}{>}\in{}i_t$. Let
      us take such an $i$, $j$ and $id_{ji}$, and let us use $i$ to
      prove the goal: \fbox{$\sigma'(id_t)(\texttt{rt})[i]=\mathtt{true}$.}

      \noindent{}By property of the stabilize relation,
      ${<}\mathtt{reinit\_transition\_time(j)\Rightarrow{}id_{ji}}{>}\in{}o_p$
      and
      ${<}\mathtt{reinit\_time(i)\Rightarrow{}id_{ji}}{>}\in{}i_t$,
      we can deduce
      $\sigma'(id_t)(\texttt{rt})[i]=\sigma'(id_{ji})=\sigma'(id_p)(\texttt{rtt})[j]$.

      Rewriting the goal with
      $\sigma'(id_t)(\texttt{rt})[i]=\sigma'(id_{ji})=\sigma'(id_p)(\texttt{rtt})[j]$,
      \fbox{$\sigma'(id_p)(\texttt{rtt})[j]=\mathtt{true}$.}

      By property of the $\mathtt{Inject}$, the \hvhdl{}
      rising edge and the stabilize relations, \InCsCompP, and through
      the examination of the
      \texttt{reinit_transitions_time_evaluation} process defined in
      the place design architecture, we can deduce:
      \begin{equation}
        \label{eq:eq-rtt-j}
        \begin{split}
          \sigma'(id_p)(\texttt{rtt})[j]=& \big((\sigma(id_p)(\texttt{oat})[j]=\mathtt{basic}+\sigma(id_p)(\texttt{oat})[j]=\mathtt{test}) \\
          & .(\sigma(id_p)(\texttt{sm})-\sigma(id_p)(\texttt{sots})<\sigma(id_p)(\texttt{oaw})[j])\\
          & .(\sigma(id_p)(\texttt{sots})>0)\big)\\
          & +\sigma(id_p)(\texttt{otf})[j] \\
        \end{split}
      \end{equation}

      Rewriting the goal with \eqref{eq:eq-rtt-j},
      \begin{equation*}
        \fbox{$\begin{split}
            \mathtt{true}=& ((\sigma(id_p)(\texttt{oat})[j]=\mathtt{basic}+\sigma(id_p)(\texttt{oat})[j]=\mathtt{test}) \\
            & .(\sigma(id_p)(\texttt{sm})-\sigma(id_p)(\texttt{sots})<\sigma(id_p)(\texttt{oaw})[j])\\
            & .(\sigma(id_p)(\texttt{sots})>0))\\
            & +(\sigma(id_p)(\texttt{otf})[j]) \\
          \end{split}$}
      \end{equation*}

      By construction, there exists $id_{ft}\in{}Sigs(\Delta)$ such that\\
      ${<}\mathtt{output\_transitions\_fired(j)\Rightarrow{}id_{ft}}{>}\in{}i_p$
      and ${<}\mathtt{fired\Rightarrow{}id_{ft}}{>}\in{}o_t$. By
      property of state $\sigma$, which is a stable state, we have
      $\sigma(id_t)(\texttt{fired})=\sigma(id_{ft})=\sigma(id_p)(\texttt{otf})[j]$.

      Rewriting the goal with
      $\sigma(id_t)(\texttt{fired})=\sigma(id_{ft})=\sigma(id_p)(\texttt{otf})[j]$,
      \begin{equation*}
        \fbox{$\begin{split}
            \mathtt{true}=& ((\sigma(id_p)(\texttt{oat})[j]=\mathtt{basic}+\sigma(id_p)(\texttt{oat})[j]=\mathtt{test}) \\
            & .(\sigma(id_p)(\texttt{sm})-\sigma(id_p)(\texttt{sots})<\sigma(id_p)(\texttt{oaw})[j])\\
            & .(\sigma(id_p)(\texttt{sots})>0))\\
            & +\sigma(id_t)(\texttt{fired}) \\
          \end{split}$}
      \end{equation*}

      By property of
      $\gamma\vdash{}s\stackrel{\downarrow}{\approx}\sigma$, we can
      deduce
      $t\in{}Fired(s)\Leftrightarrow\sigma(id_t)(\texttt{fired})=\mathtt{true}$.

      Rewriting the goal with
      $t\in{}Fired(s)\Leftrightarrow\sigma(id_t)(\texttt{fired})=\mathtt{true}$
      and simplify the goal, then \qedbox{tautology}.

    \item \textbf{CASE} $t\notin{}Fired(s)$: Then, there are two cases
      that will determine the value of $s'.reset_t(t)$. Either there
      exists a place $p$ with an output token sum greater than zero,
      that is connected to $t$ by an $\mathtt{basic}$ or
      $\mathtt{test}$ arc, and such that the transient marking of $p$
      disables $t$; or such a place does not exist (the predicate is
      decidable).

      \begin{itemize}
      \item \textbf{CASE} there exists such a place $p$ as described
        above:

        Then, let us take such a place $p$ and $\omega\in\mathbb{N}^{*}$ s.t.:
        \begin{enumerate}
        \item $\sum\limits_{t_i\in{}Fired(s)}pre(p,t_i)>0$\label{item:1}
        \item
          $pre(p,t)=(\omega,\mathtt{basic})\lor{}pre(p,t)=(\omega,\mathtt{test})$\label{item:2}
        \item
          $s.M(p)-\sum\limits_{t_i\in{}Fired(s)}pre(p,t_i)<\omega$\label{item:3}
        \end{enumerate}

        We will only consider the case where
        $pre(p,t)=(\omega,\mathtt{basic})$; the proof is the similar
        when $pre(p,t)=(\omega,\mathtt{test})$.

        Assuming that $p$ exists, and by property of
        $E_c,\tau\vdash{}s\srarrow{\uparrow}{\fontsize{7}{7}\selectfont}s'$
        (Rule~\ref{it:reset-order}), we can deduce
        $s'.reset_t(t)=\mathtt{true}$.

        Rewriting the goal with $s'.reset_t(t)=\mathtt{true}$,
        \fbox{$\sum\limits_{i=0}^{\vert{}input(t)\vert-1}\sigma'(id_t)(\texttt{rt})[i]=\mathtt{true}$.}

        \noindent{}To prove the goal, let us show
        \fbox{$\exists{}i\in[0,\vert{}input(t)\vert-1]$
          s.t. $\sigma'(id_t)(\texttt{rt})[i]=\mathtt{true}$.}

        By construction, there exists $id_p\in{}Comps(\Delta)$
        s.t. $\gamma(p)=id_p$.

        \exP{}

        By construction, there exist an
        $i\in{}[0,\vert{}input(t)\vert-1],$ a
        $j\in[0,\vert{}output(p)\vert-1]$ and
        $id_{ji}\in{}Sigs(\Delta)$ s.t.
        ${<}\mathtt{reinit\_transition\_time(j)\Rightarrow{}id_{ji}}{>}\in{}o_p$
        and\\
        ${<}\mathtt{reinit\_time(i)\Rightarrow{}id_{ji}}{>}\in{}i_t$. Let
        us take such an $i$, $j$ and $id_{ji}$, and let us use $i$ to
        prove the goal: \fbox{$\sigma'(id_t)(\texttt{rt})[i]=\mathtt{true}$.}

        By property of the stabilize relation,
        ${<}\mathtt{reinit\_transition\_time(j)\Rightarrow{}id_{ji}}{>}\in{}o_p$
        and
        ${<}\mathtt{reinit\_time(i)\Rightarrow{}id_{ji}}{>}\in{}i_t$,
        we have
        $\sigma'(id_t)(\texttt{rt})[i]=\sigma'(id_{ji})=\sigma'(id_p)(\texttt{rtt})[j]$.

        Rewriting the goal with
        $\sigma'(id_t)(\texttt{rt})[i]=\sigma'(id_{ji})=\sigma'(id_p)(\texttt{rtt})[j]$, we have\\
        \fbox{$\sigma'(id_p)(\texttt{rtt})[j]=\mathtt{true}$.}

        By property of the $\mathtt{Inject}$, the \hvhdl{}
        rising edge and the stabilize relation, and through the
        examination of the \texttt{reinit_transitions_time_evaluation}
        process defined in the place design architecture, we can
        deduce:
        \begin{equation}
          \label{eq:eq-rtt-j-2}
          \begin{split}
            \sigma'(id_p)(\texttt{rtt})[j]=& \big((\sigma(id_p)(\texttt{oat})[j]=\mathtt{basic}+\sigma(id_p)(\texttt{oat})[j]=\mathtt{test}) \\
            & .(\sigma(id_p)(\texttt{sm})-\sigma(id_p)(\texttt{sots})<\sigma(id_p)(\texttt{oaw})[j])\\
            & .(\sigma(id_p)(\texttt{sots})>0)\big)\\
            & +\sigma(id_p)(\texttt{otf})[j] \\
          \end{split}
        \end{equation}

        Rewriting the goal with \eqref{eq:eq-rtt-j-2},
        \begin{equation*}
          \fbox{$\begin{split}
              \mathtt{true}=& ((\sigma(id_p)(\texttt{oat})[j]=\mathtt{basic}+\sigma(id_p)(\texttt{oat})[j]=\mathtt{test}) \\
              & .(\sigma(id_p)(\texttt{sm})-\sigma(id_p)(\texttt{sots})<\sigma(id_p)(\texttt{oaw})[j])\\
              & .(\sigma(id_p)(\texttt{sots})>0))\\
              & +\sigma(id_p)(\texttt{otf})[j] \\
            \end{split}$}
        \end{equation*}

        By construction,
        ${<}\mathtt{output\_arcs\_types(j)\Rightarrow{}basic}{>}\in{}i_p$
        and\\
        ${<}\mathtt{output\_arcs\_weights(j)\Rightarrow{}}\omega{>}\in{}i_p$.
        
        By property of the stabilize relation and \InCsCompP, we can
        deduce $\sigma'(id_p)(\texttt{oat})[j]=\mathtt{basic}$ and
        $\sigma'(id_p)(\texttt{oaw})[j]=\omega$.

        By property of
        $\gamma\vdash{}s\stackrel{\downarrow}{\approx}\sigma$, we can
        deduce $\sigma(id_p)(\texttt{sm})=s.M(p)$ and
        $\sigma(id_p)(\texttt{sots})=\sum\limits_{t_i\in{}Fired(s)}pre(p,t_i)$.

        \noindent{}Rewriting the goal with
        $\sigma'(id_p)(\texttt{oat})[j]=\mathtt{basic}$,
        $\sigma'(id_p)(\texttt{oaw})[j]=\omega$, $\sigma(id_p)(\texttt{sm})=s.M(p)$
        and
        $\sigma(id_p)(\texttt{sots})=\sum\limits_{t_i\in{}Fired(s)}pre(p,t_i)$,
        and simplifying the goal:
        \begin{equation*}
          \fbox{
            \begin{tabular}{@{}c@{}}
              $\big((s.M(p)-\sum\limits_{t_i\in{}Fired(s)}pre(p,t_i)<\omega)~.~(\sum\limits_{t_i\in{}Fired(s)}pre(p,t_i)>0)\big)+\sigma(id_t)(\texttt{fired})$ \\
              $=$\\
              $\mathtt{true}$\\
            \end{tabular}
          }
        \end{equation*}

        We assumed that
        $s.M(p)-\sum\limits_{t_i\in{}Fired(s)}pre(p,t_i)<\omega$ and
        $\sum\limits_{t_i\in{}Fired(s)}pre(p,t_i)>0$.

        Thus, by assumption: \\
        \qedbox{
          \begin{tabular}{@{}c@{}}
            $\big((s.M(p)-\sum\limits_{t_i\in{}Fired(s)}pre(p,t_i)<\omega)~.~(\sum\limits_{t_i\in{}Fired(s)}pre(p,t_i)>0)\big)+\sigma(id_t)(\texttt{fired})$ \\
            $=$ \\
            $\mathtt{true}$\\
          \end{tabular}
        }
        
      \item \textbf{CASE} such a place does not exist:\\
        Then, let us assume that, for all place $p\in{}P$
        \begin{enumerate}
        \item $\sum\limits_{t_i\in{}Fired(s)}pre(p,t_i)=0$\label{item:4}
        \item or
          $\forall{}\omega\in\mathbb{N}^{*},~pre(p,t)=(\omega,\mathtt{basic})\lor{}pre(p,t)=(\omega,\mathtt{test})
          \Rightarrow{}$
          
          \hspace{3ex}$s.M(p)-\sum\limits_{t_i\in{}Fired(s)}pre(p,t_i)\ge\omega$.\label{item:5}
        \end{enumerate}
        
        In that case, by property of
        $E_c,\tau\vdash{}s\srarrow{\uparrow}{\fontsize{7}{7}\selectfont}s'$
        (Rule~\ref{it:reset-order}), we can deduce
        $s'.reset_t(t)=\mathtt{false}$.
        
        Rewriting the goal with $s'.reset_t(t)=\mathtt{false}$:
        \fbox{$\sum\limits_{i=0}^{\vert{}input(t)\vert-1}\sigma'(id_t)(\texttt{rt})[i]=\mathtt{false}$.}

        To prove the goal, let us show
        \fbox{$\forall{}i\in[0,\vert{}input(t)\vert-1],~\sigma'(id_t)(\texttt{rt})[i]=\mathtt{false}$.}\\

        Given an $i\in[0,\vert{}input(t)\vert-1]$, let us show
        \fbox{$\sigma'(id_t)(\texttt{rt})[i]=\mathtt{false}$.}\\

        By construction, there exist a $p\in{}input(t)$, an
        $id_p\in{}Comps(\Delta)$, $g_p$, $i_p$, $o_p$, \\ a
        $j\in[0,\vert{}output(p)\vert-1]$, an
        $id_{ji}\in{}Sigs(\Delta)$ s.t. $\gamma(p)=id_p$ and
        \InCsCompP{} and
        ${<}\mathtt{reinit\_transition\_time(j)\Rightarrow{}id_{ji}}{>}\in{}o_p$
        and
        ${<}$\texttt{reinit\_time(i)}$\Rightarrow{}\mathtt{id_{ji}}{>}\in{}i_t$. Let
        us take such a $p$, $id_p$, $g_p$, $i_p$, $o_p$, $j$ and
        $id_{ji}$.\\

        By property of the stabilize relation,
        ${<}\mathtt{reinit\_transition\_time(j)\Rightarrow{}id_{ji}}{>}\in{}o_p$
        and
        ${<}\mathtt{reinit\_time(i)\Rightarrow{}id_{ji}}{>}\in{}i_t$,
        we have
        $\sigma'(id_t)(\texttt{rt})[i]=\sigma'(id_{ji})=\sigma'(id_p)(\texttt{rtt})[j]$.

        Rewriting the goal with
        $\sigma'(id_t)(\texttt{rt})[i]=\sigma'(id_{ji})=\sigma'(id_p)(\texttt{rtt})[j]$:\\
        \fbox{$\sigma'(id_p)(\texttt{rtt})[j]=\mathtt{false}$.}

        By property of the $\mathtt{Inject}$, the \hvhdl{}
        rising edge and the stabilize relations, \InCsCompP, and
        through the examination of the
        \texttt{reinit_transitions_time_evaluation} process defined in
        the place design architecture, we can deduce:
        \begin{equation}
          \label{eq:10}
          \begin{split}
            \sigma'(id_p)(\texttt{rtt})[j]=& \big((\sigma(id_p)(\texttt{oat})[j]=\mathtt{basic}+\sigma(id_p)(\texttt{oat})[j]=\mathtt{test}) \\
            & .(\sigma(id_p)(\texttt{sm})-\sigma(id_p)(\texttt{sots})<\sigma(id_p)(\texttt{oaw})[j])\\
            & .(\sigma(id_p)(\texttt{sots})>0)\big)\\
            & +\sigma(id_p)(\texttt{otf})[j] \\
          \end{split}
        \end{equation}

        Rewriting the goal with \eqref{eq:10},
        \begin{equation*}
          \fbox{$\begin{split}
              \mathtt{false}=& ((\sigma(id_p)(\texttt{oat})[j]=\mathtt{basic}+\sigma(id_p)(\texttt{oat})[j]=\mathtt{test}) \\
              & .(\sigma(id_p)(\texttt{sm})-\sigma(id_p)(\texttt{sots})<\sigma(id_p)(\texttt{oaw})[j])\\
              & .(\sigma(id_p)(\texttt{sots})>0))\\
              & +\sigma(id_p)(\texttt{otf})[j]) \\
            \end{split}$}
        \end{equation*}
        
        By construction, there exists $id_{ft}\in{}Sigs(\Delta)$ such that\\
        ${<}\mathtt{output\_transitions\_fired(j)\Rightarrow{}id_{ft}}{>}\in{}i_p$
        and ${<}\mathtt{fired\Rightarrow{}id_{ft}}{>}\in{}o_t$. By
        property of state $\sigma$ as being a stable state, we have
        $\sigma(id_t)(\texttt{fired})=\sigma(id_{ft})=\sigma(id_p)(\texttt{otf})[j]$.

        Rewriting the goal with
        $\sigma(id_t)(\texttt{fired})=\sigma(id_{ft})=\sigma(id_p)(\texttt{otf})[j]$:
        \begin{equation*}
          \fbox{$\begin{split}
              \mathtt{false}=& ((\sigma(id_p)(\texttt{oat})[j]=\mathtt{basic}+\sigma(id_p)(\texttt{oat})[j]=\mathtt{test}) \\
              & .(\sigma(id_p)(\texttt{sm})-\sigma(id_p)(\texttt{sots})<\sigma(id_p)(\texttt{oaw})[j])\\
              & .(\sigma(id_p)(\texttt{sots})>0))\\
              & +\sigma(id_t)(\texttt{fired}) \\
            \end{split}$}
        \end{equation*}

        By property of
        $\gamma\vdash{}s\stackrel{\downarrow}{\approx}\sigma$, we can
        deduce
        $t\notin{}Fired(s)\Leftrightarrow\sigma(id_t)(\texttt{fired})=\mathtt{false}$

        Rewriting the goal with
        $t\notin{}Fired(s)\Leftrightarrow\sigma(id_t)(\texttt{fired})=\mathtt{false}$
        and simplifying the goal:
        \begin{equation*}
          \fbox{$\begin{split}
              \mathtt{false}=& ((\sigma(id_p)(\texttt{oat})[j]=\mathtt{basic}+\sigma(id_p)(\texttt{oat})[j]=\mathtt{test}) \\
              & .(\sigma(id_p)(\texttt{sm})-\sigma(id_p)(\texttt{sots})<\sigma(id_p)(\texttt{oaw})[j])\\
              & .(\sigma(id_p)(\texttt{sots})>0))\\
            \end{split}$}
        \end{equation*}

        Then, based on the assumptions made at the beginning of case,
        there are two cases:
        \begin{enumerate}
        \item \textbf{CASE}
          $\sum\limits_{t_i\in{}Fired(s)}pre(p,t_i)=0$:\\
          
          By property of
          $\gamma\vdash{}s\stackrel{\downarrow}{\approx}\sigma$, we
          can deduce
          $\sum\limits_{t_i\in{}Fired(s)}pre(p,t_i)=\sigma(id_p)(\texttt{sots})$.

          Rewriting the goal with
          $\sum\limits_{t_i\in{}Fired(s)}pre(p,t_i)=\sigma(id_p)(\texttt{sots})$
          and $\sum\limits_{t_i\in{}Fired(s)}pre(p,t_i)=0$, and
          simplifying the goal: \qedbox{tautology.}
          
        \item \textbf{CASE}
          $\forall{}\omega\in\mathbb{N}^{*},~pre(p,t)=(\omega,\mathtt{basic})\lor{}pre(p,t)=(\omega,\mathtt{test})
          \Rightarrow{}$

          \hspace{7ex}$s.M(p)-\sum\limits_{t_i\in{}Fired(s)}pre(p,t_i)\ge\omega$:

          Let us perform case analysis on $pre(p,t)$; there are two
          cases:

          \begin{enumerate}
          \item \textbf{CASE} $pre(p,t)=(\omega,\mathtt{basic})$ or $pre(p,t)=(\omega,\mathtt{basic})$:\\
            \noindent{}By construction,
            ${<}\mathtt{output\_arcs\_weights(j)\Rightarrow{}}\omega{>}\in{}i_p$.
            
            \noindent{}By property of stable state $\sigma$ and
            \InCsCompP, we can deduce $\sigma(id_p)(\texttt{oaw})[j]=\omega$.

            By property of
            $\gamma\vdash{}s\stackrel{\downarrow}{\approx}\sigma$, we
            can deduce $\sigma(id_p)(\texttt{sm})=s.M(p)$ and
            $\sigma(id_p)(\texttt{sots})=\sum\limits_{t_i\in{}Fired(s)}pre(p,t_i)$.

            Rewriting the goal with
            $\sigma(id_p)(\texttt{oaw})[j]=\omega$,
            $\sigma(id_p)(\texttt{sm})=s.M(p)$ and\\
            $\sigma(id_p)(\texttt{sots})=\sum\limits_{t_i\in{}Fired(s)}pre(p,t_i)$:
            \begin{equation*}
              \fbox{$\begin{split}
                  \mathtt{false}=& ((\sigma(id_p)(\texttt{oat})[j]=\mathtt{basic}+\sigma(id_p)(\texttt{oat})[j]=\mathtt{test}) \\
                  & .(s.M(p)-\sum\limits_{t_i\in{}Fired(s)}pre(p,t_i)<\omega)\\
                  & .(\sum\limits_{t_i\in{}Fired(s)}pre(p,t_i)>0))\\
                \end{split}$}
            \end{equation*}
            
            We assumed that
            $s.M(p)-\sum\limits_{t_i\in{}Fired(s)}pre(p,t_i)\ge\omega$,
            and then we can deduce
            $s.M(p)-\sum\limits_{t_i\in{}Fired(s)}pre(p,t_i)<\omega=\mathtt{false}$.
            

            Rewriting the goal with
            $s.M(p)-\sum\limits_{t_i\in{}Fired(s)}pre(p,t_i)<\omega=\mathtt{false}$,
            and simplifying the goal, \qedbox{tautology.}

          \item \textbf{CASE} $pre(p,t)=(\omega,\mathtt{inhib})$:
            
            By construction,
            ${<}\mathtt{output\_arcs\_types(j)\Rightarrow{}inhib}{>}\in{}i_p$.
            
            By property of stable state $\sigma$ and \InCsCompP, we
            can deduce $\sigma(id_p)(\texttt{oat})[j]=\mathtt{inhib}$.

            Rewriting the goal with
            $\sigma(id_p)(\texttt{oat})[j]=\mathtt{inhib}$, and simplifying
            the goal, we have a \qedbox{tautology.}
          \end{enumerate}
        \end{enumerate}
      \end{itemize}
    \end{itemize}
  \end{itemize}
\end{niproof}

\subsection{Rising edge and action executions}
\label{sec:re-action-exec}

\begin{lemma}[Rising edge equal action executions]
  \label{lem:re-equal-action-exec}
  \rehyps{} then\\
  $\forall{}a\in\mathcal{A},id_a\in{}Outs(\Delta)~s.t.~\gamma(a)=id_a,~s'.ex(a)=\sigma'(id_a)$.
\end{lemma}

\begin{niproof}
  Given an $a\in\mathcal{A}$ and an
  $id_a\in{}Outs(\Delta)~s.t.~\gamma(a)=id_a$, let us show
  \fbox{$s'.ex(a)=\sigma'(id_a)$.}

  By property of
  $E_c,\tau\vdash{}s\srarrow{\uparrow}{\fontsize{7}{7}\selectfont}s'$,
  we can deduce $s.ex(a)=s'.ex(a)$.
  
  By construction, $id_a$ is an output port identifier of Boolean type
  in the \hvhdl{} design $d$. The generated \texttt{``action''}
  process is responsible for the assignment of the $id a$ only during
  the initialization phase or during a falling edge phase.

  By property of the \hvhdl{} $\mathtt{Inject}$, rising
  edge, stabilize relations, and the \texttt{``action''} process, we
  can deduce $\sigma(id_a)=\sigma'(id_a)$.

  Rewriting the goal with $s.ex(a)=s'.ex(a)$ and
  $\sigma(id_a)=\sigma'(id_a)$, \fbox{$s.ex(a)=\sigma(id_a)$.}

  By property of
  $\gamma\vdash{}s\stackrel{\downarrow}{\approx}\sigma$,
  \qedbox{$s.ex(a)=\sigma(id_a)$.}
\end{niproof}

\subsection{Rising edge and function executions}
\label{sec:re-fun-exec}

\begin{lemma}[Rising edge equal function executions]
  \label{lem:re-equal-fun-exec}
  \rehyps{} then\\
  $\forall{}f\in\mathcal{F},id_f\in{}Outs(\Delta)~s.t.~\gamma(f)=id_f,~s'.ex(f)=\sigma'(id_f)$.
\end{lemma}

\begin{niproof}
  Given an $f\in\mathcal{F}$ and an $id_f\in{}Outs(\Delta)$
  s.t. $\gamma(f)=id_f$, let us show \fbox{$s'.ex(f)=\sigma'(id_f)$.}

  By property of
  $E_c,\tau\vdash{}s\srarrow{\uparrow}{\fontsize{7}{7}\selectfont}s'$
  (Rule~\ref{it:exec-fun}):
  \begin{equation}
    s'.ex(f)=\sum\limits_{t\in{}Fired(s)}\mathbb{F}(t,f)\label{eq:eq-exf}
  \end{equation}

  By construction, $id_f$ is an output port identifier of Boolean type
  in the \hvhdl{} design $d$. The generated \texttt{function} process
  assigns a value to the output port $id_f$ only during the
  initialization phase or during a rising edge phase.
  
  By construction, the \texttt{function} process is defined in the
  behavior of design $d$, i.e.\\
  $\mathtt{ps}(\texttt{function}, \emptyset, sl, ss)\in{}d.cs$.
  
  Let $trs(f)$ be the set of transitions associated to function $f$,
  i.e. $trs(f)=\{t\in{}T~\vert~\mathbb{F}(t,f)=\mathtt{true}\}$.

  Let us perform case analysis on $trs(f)$; there are two cases:
  
  \begin{itemize}
  \item \textbf{CASE} $trs(f)=\emptyset$:
    
    By construction,
    $\mathtt{id_f\Leftarrow{}false}\in{}ss_{\uparrow}$ where
    $ss_\uparrow$ is the part of the \texttt{function} process body
    executed during a rising edge phase.

    By property of the \hvhdl{} rising edge, the stabilize relations
    and $\mathtt{ps}(\texttt{function}, \emptyset, sl, ss)\in{}d.cs$, we can
    deduce $\sigma'(id_f)=\mathtt{false}$.
    
    \noindent{}By property of
    $\sum\limits_{t\in{}Fired(s)}\mathbb{F}(t,f)$ and
    $trs(f)=\emptyset$, we can deduce
    $\sum\limits_{t\in{}Fired(s)}\mathbb{F}(t,f)=\mathtt{false}$.

    \noindent{}Rewriting the goal with \eqref{eq:eq-exf},
    $\sigma'(id_f)=\mathtt{false}$ and
    $\sum\limits_{t\in{}Fired(s)}\mathbb{F}(t,f)=\mathtt{false}$:
    \qedbox{tautology.}
    
  \item \textbf{CASE} $trs(f)\neq\emptyset$:
    
    By construction,
    $\mathtt{id_f\Leftarrow{}id_{ft_0}+\dots+id_{ft_n}}\in{}ss_\uparrow$,
    where $id_{ft_i}\in{}Sigs(\Delta)$, $ss_\uparrow$ is the part of
    the \texttt{function} process body executed during a rising edge
    phase, and $n=\vert{}trs(f)\vert-1$.

    By property of the $\mathtt{Inject}$, the
    \hvhdl{} rising edge, the stabilize relations, and\\
    $\mathtt{ps}(\texttt{function}, \emptyset, sl, ss)\in{}d.cs$, we can
    deduce:
    \begin{equation}
      \sigma'(id_f)=\sigma(id_{ft_0})+\dots+\sigma(id_{ft_n})\label{eq:eq-idf-prod}
    \end{equation}

    Rewriting the goal with \eqref{eq:eq-exf} and
    \eqref{eq:eq-idf-prod},
    \fbox{$\sum\limits_{t\in{}Fired(s)}\mathbb{F}(t,f)=\sigma(id_{ft_0})+\dots+\sigma(id_{ft_n})$.}

    Let us reason on the value of
    $\sigma(id_{ft_0})+\dots+\sigma(id_{ft_n})$; there are two cases:

    \begin{itemize}
    \item \textbf{CASE} $\sigma(id_{ft_0})+\dots+\sigma(id_{ft_n})=\mathtt{true}$:
      
      Then, we can rewrite the goal as follows:
      \fbox{$\sum\limits_{t\in{}Fired(s)}\mathbb{F}(t,f)=\mathtt{true}$.}

      To prove the above goal, let us show
      \fbox{$\exists{}t\in{}Fired(s)~s.t.~\mathbb{F}(t,f)=\mathtt{true}$.}

      From $\sigma(id_{ft_0})+\dots+\sigma(id_{ft_n})=\mathtt{true}$,
      we can deduce
      $\exists{}id_{ft_i}~s.t.~\sigma(id_{ft_i})=\mathtt{true}$. Let
      us take such an $id_{ft_i}$.
      
      By construction, there exist a $t\in{}trs(f)$, an
      $id_t\in{}Comps(\Delta)$, $g_t$, $i_t$, $o_t$ such that:
      \begin{itemize}
      \item $\gamma(t)=id_t$
      \item \InCsCompT{}
      \item ${<}\mathtt{fired\Rightarrow{id_{ft_i}}}{>}\in{}o_{t}$
      \end{itemize}
      
      \noindent{}By property of $\sigma$ as being a stable design
      state, and \InCsCompT, we can deduce
      $\sigma(id_{t})(\texttt{fired})=\sigma(id_{ft_i})$, and thus that
      $\sigma(id_{t})(\texttt{fired})=\mathtt{true}$.

      By property of
      $\gamma\vdash{}s\stackrel{\downarrow}{\approx}\sigma$, we can
      deduce $t\in{}Fired(s)$.
      
      Let us use $t$ to prove the goal:
      \fbox{$\mathbb{F}(t,f)=\mathtt{true}$.}

      By definition of $t\in{}trs(f)$,
      \qedbox{$\mathbb{F}(t,f)=\mathtt{true}$.}

    \item \textbf{CASE} $\sigma(id_{ft_0})+\dots+\sigma(id_{ft_n})=\mathtt{false}$:\\
      \noindent{}Then, we can rewrite the goal as follows:
      \fbox{$\sum\limits_{t\in{}Fired(s)}\mathbb{F}(t,f)=\mathtt{false}$.}

      \noindent{}To prove the above goal, let us show
      \fbox{$\forall{}t\in{}Fired(s)~s.t.~\mathbb{F}(t,f)=\mathtt{false}$.}

      \noindent{}Given a $t\in{}Fired(s)$, let us show
      \fbox{$\mathbb{F}(t,f)=\mathtt{false}$.}

      \noindent{}Let us perform case analysis on $\mathbb{F}(t,f)$; there are 2 cases:

      \begin{itemize}
      \item \textbf{CASE} \qedbox{$\mathbb{F}(t,f)=\mathtt{false}$.}
      \item \textbf{CASE} $\mathbb{F}(t,f)=\mathtt{true}$:\\
        
        By construction, there exist an $id_{t}\in{}Comps(\Delta)$,
        $g_{t}$, $i_{t}$, $o_{t}$ and
        $id_{ft_i}\in{}Sigs(\Delta)$ such that:
        \begin{itemize}
        \item $\gamma(t)=id_{t}$
        \item \InCsCompT
        \item ${<}\mathtt{fired\Rightarrow{id_{ft_i}}}{>}\in{}o_{t}$
        \end{itemize}

        By property of stable design state $\sigma$ and \InCsCompT, we
        can deduce $\sigma(id_t)(\texttt{fired})=\sigma(id_{ft_i})$.

        By property of
        $\gamma\vdash{}s\stackrel{\downarrow}{\approx}\sigma$, we can
        deduce
        $t\in{}Fired(s)\Leftrightarrow{}\sigma(id_t)(\texttt{fired})=\mathtt{true}$.

        \noindent{}Since $t\in{}Fired(s)$, we can deduce
        $\sigma(id_t)(\texttt{fired})=\mathtt{true}$, and from
        $\sigma(id_t)(\texttt{fired})=\sigma(id_{ft_i})$, we can deduce
        $\sigma(id_{ft_i})=\mathtt{true}$.

        Then, \qedbox{$\sigma(id_{ft_i})=\mathtt{true}$ contradicts
          $\sigma(id_{ft_0})+\dots+\sigma(id_{ft_n})=\mathtt{false}$.}
      \end{itemize}
    \end{itemize}
  \end{itemize}
\end{niproof}

\subsection{Rising edge and sensitization}
\label{sec:re-sens}

\begin{lemma}[Rising edge equal sensitized]
  \label{lem:re-equal-sens}
  \rehyps{} then\\
  $\forall{}t\in{}T,id_t\in{}Comps(\Delta)~s.t.~\gamma(t)=id_t,$
  $t\in{}Sens(s'.M)\Leftrightarrow\sigma'(id_t)(\texttt{s\_enabled})=\mathtt{true}$.
\end{lemma}

\begin{niproof}
  Given a $t\in{}T$ and an $id_t\in{}Comps(\Delta)$
  s.t. $\gamma(t)=id_t$, let us show\\
  \fbox{$t\in{}Sens(s'.M)\Leftrightarrow\sigma'(id_t)(\texttt{s\_enabled})=\mathtt{true}$.}\\
  
  \exT{}. Then, the proof is in two parts:

  \begin{enumerate}
  \item Assuming that $t\in{}Sens(s'.M)$, let us show
    \fbox{$\sigma'(id_t)(\texttt{s\_enabled})=\mathtt{true}$.}

    \noindent{}By property of the stabilize relation, \InCsCompT{},
    and through the examination of the \texttt{enable_evaluation}
    process defined in the transition design architecture:
    \begin{equation}
      \sigma'(id_t)(\texttt{se})=\prod\limits_{i=0}^{\Delta(id_t)(\texttt{ian})-1}\sigma'(id_t)(\texttt{input\_arcs\_valid})[i]\label{eq:eq-senabled-prod}
    \end{equation}

    Rewriting the goal with \eqref{eq:eq-senabled-prod},
    \fbox{$\prod\limits_{i=0}^{\Delta(id_t)(\texttt{ian})-1}\sigma'(id_t)(\texttt{iav})[i]=\mathtt{true}$.}
    
    \noindent{}To prove the goal, let us show that
    \fbox{$\forall{}i\in[0,\Delta(id_t)(\texttt{ian})-1],~\sigma'(id_t)(\texttt{iav})[i]=\mathtt{true}$.}

    \noindent{}Given an $i\in[0,\Delta(id_t)(\texttt{ian})-1]$, let us show
    \fbox{$\sigma'(id_t)(\texttt{iav})[i]=\mathtt{true}$.}
    
    \noindent{}Let us perform case analysis on $input(t)$.

    \begin{itemize}
    \item \textbf{CASE} $input(t)=\emptyset$:

      By construction,
      ${<}\mathtt{input\_arcs\_number\Rightarrow{}1}{>}\in{}g_t$ and\\
      ${<}\mathtt{input\_arcs\_valid(0)\Rightarrow{}true}{>}\in{}i_t$.

      By property of the elaboration and stabilize relations and
      \InCsCompT{}, we can deduce $\Delta(id_t)(\texttt{ian})=1$ and
      $\sigma'(id_t)(\texttt{iav})[0]=\mathtt{true}$.

      Thanks to $\Delta(id_t)(\texttt{ian})=1$, we can deduce that $i=0$.

      Rewriting the goal with $\sigma'(id_t)(\texttt{iav})[0]=\mathtt{true}$,
      \qedbox{tautology.}
      
    \item \textbf{CASE} $input(t)\neq\emptyset$:\\
      \noindent{}By construction,
      ${<}\mathtt{input\_arcs\_number\Rightarrow{}}\vert{}input(t)\vert{>}\in{}g_t$.

      By property of the elaboration relation and \InCsCompT, we can
      deduce $\Delta(id_t)(\texttt{ian})=\vert{}input(t)\vert$.
      
      Thanks to $\Delta(id_t)(\texttt{ian})=\vert{}input(t)\vert$, we know
      that $i\in[0,\vert{}input(t)\vert-1]$.

      By construction, there exist a $p\in{}input(t)$,
      $id_p\in{}Comps(\Delta)$, $g_p$, $i_p$, $o_p$,
      $j\in{}[0,\vert{}output(p)\vert-1]$ and
      $id_{ji}\in{}Sigs(\Delta)$ s.t. $\gamma(p)=id_p$ and\\
      \InCsCompP{} and
      ${<}\mathtt{output\_arcs\_valid(j)\Rightarrow{}id_{ji}}{>}\in{}o_p$
      and
      ${<}\mathtt{input\_arcs\_valid(i)\Rightarrow{}id_{ji}}{>}\in{}i_t$.

      By property of the stabilize relation, \InCsCompT{} and
      \InCsCompP, we can deduce
      $\sigma'(id_t)(\texttt{iav})[i]=\sigma'(id_{ji})=\sigma'(id_p)(\texttt{oav})[j]$.

      Rewriting the goal with
      $\sigma'(id_t)(\texttt{iav})[i]=\sigma'(id_{ji})=\sigma'(id_p)(\texttt{oav})[j]$:\\
      \fbox{$\sigma'(id_p)(\texttt{oav})[j]=\mathtt{true}$.}

      By property of the stabilize relation, \InCsCompP, and through
      the examination of the \texttt{marking_validation_evaluation}
      process defined in the place design architecture, we can deduce:
      \begin{equation}
        \label{eq:eq-oav-sens}
        \begin{split}
          \sigma'(id_p)(\texttt{oav})[j]=& \big((\sigma'(id_p)(\texttt{oat})[j]=\mathtt{basic}+\sigma'(id_p)(\texttt{oat})[j]=\mathtt{test}) \\
          & \quad.~\sigma'(id_p)(\texttt{sm})\ge\sigma'(id_p)(\texttt{oaw})[j]\big)\\
          & +\big(\sigma'(id_p)(\texttt{oat})[j]=\mathtt{inhib}~.~\sigma'(id_p)(\texttt{sm})<\sigma'(id_p)(\texttt{oaw})[j]\big)\\
        \end{split}
      \end{equation}

      Rewriting the goal with \eqref{eq:eq-oav-sens},
      \begin{equation*}
        \fbox{$
          \begin{split}
            \mathtt{true}=& \big((\sigma'(id_p)(\texttt{oat})[j]=\mathtt{basic}+\sigma'(id_p)(\texttt{oat})[j]=\mathtt{test}) \\
            & \quad.~\sigma'(id_p)(\texttt{sm})\ge\sigma'(id_p)(\texttt{oaw})[j]\big)\\
            & +\big(\sigma'(id_p)(\texttt{oat})[j]=\mathtt{inhib}~.~\sigma'(id_p)(\texttt{sm})<\sigma'(id_p)(\texttt{oaw})[j]\big)\\
          \end{split}
          $}
      \end{equation*}

      Let us perform case analysis on $pre(p,t)$; there are 3 cases:
      \begin{itemize}
      \item \textbf{CASE} $pre(p,t)=(\omega,\mathtt{basic})$:\\

        \noindent{}By construction,
        ${<}\mathtt{output\_arcs\_types(j)\Rightarrow{}basic}{>}\in{}i_p$
        and\\
        ${<}\mathtt{output\_arcs\_weights(j)\Rightarrow{}}\omega{>}\in{}i_p$.

        By property of the stabilize relation and \InCsCompP,

        we can deduce $\sigma'(id_p)(\texttt{oat})[j]=\mathtt{basic}$ and
        $\sigma'(id_p)(\texttt{oaw})[j]=\omega$.

        Rewriting the goal with
        $\sigma'(id_p)(\texttt{oat})[j]=\mathtt{basic}$
        and $\sigma'(id_p)(\texttt{oaw})[j]=\omega$, and simplifying the goal:\\
        \fbox{$\sigma'(id_p)(\texttt{sm})\ge\omega=\mathtt{true}$.}

        \noindent{}Appealing to Lemma~\ref{lem:re-equal-marking}, we
        can deduce $s'.M(p)=\sigma'(id_p)(\texttt{sm})$.

        \noindent{}Rewriting the goal with
        $s'.M(p)=\sigma'(id_p)(\texttt{sm})$:
        \fbox{$s'.M(p)\ge\omega=\mathtt{true}$.}

        \noindent{}By definition of $t\in{}Sens(s'.M)$,
        \qedbox{$s'.M(p)\ge\omega=\mathtt{true}$.}
        
      \item \textbf{CASE} $pre(p,t)=(\omega,\mathtt{test})$: same as
        above.
      \item \textbf{CASE} $pre(p,t)=(\omega,\mathtt{inhib})$:

        \noindent{}By construction,
        ${<}\mathtt{output\_arcs\_types(j)\Rightarrow{}inhib}{>}\in{}i_p$
        and\\
        ${<}\mathtt{output\_arcs\_weights(j)\Rightarrow{}}\omega{>}\in{}i_p$.

        \noindent{}By property of the stabilize relation and
        \InCsCompP, we can deduce
        $\sigma'(id_p)(\texttt{oat})[j]=\mathtt{inhib}$ and
        $\sigma'(id_p)(\texttt{oaw})[j]=\omega$.

        \noindent{}Rewriting the goal with
        $\sigma'(id_p)(\texttt{oat})[j]=\mathtt{inhib}$ and
        $\sigma'(id_p)(\texttt{oaw})[j]=\omega$, and simplifying the goal:
        \fbox{$\sigma'(id_p)(\texttt{sm})<\omega=\mathtt{true}$.}

        \noindent{}Appealing to Lemma~\ref{lem:re-equal-marking}, we
        can deduce $s'.M(p)=\sigma'(id_p)(\texttt{sm})$.

        \noindent{}Rewriting the goal with
        $s'.M(p)=\sigma'(id_p)(\texttt{sm})$:
        \fbox{$s'.M(p)<\omega=\mathtt{true}$.}

        \noindent{}By definition of $t\in{}Sens(s'.M)$,
        \qedbox{$s'.M(p)<\omega=\mathtt{true}$.}
        
      \end{itemize}
    \end{itemize}
    
  \item Assuming that $\sigma'(id_t)(\texttt{s\_enabled})=\mathtt{true}$, let
    us show \fbox{$t\in{}Sens(s'.M)$.}

    \noindent{}By definition of $t\in{}Sens(s'.M)$, let us show
    \begin{equation*}
      \fbox{\parbox{\lwidth}{$
          \forall{}p\in{}P,\omega\in\mathbb{N}^{*},
          ~\big(pre(p,t)=(\omega,\mathtt{basic})\lor{}pre(p,t)=(\omega,\mathtt{test})\Rightarrow{}s'.M(p)\ge\omega\big)\land
          \big(pre(p,t)=(\omega,\mathtt{inhib})\Rightarrow{}s'.M(p)<\omega\big)
          $}}
    \end{equation*}

    Given a $p\in{}P$ and an $\omega\in\mathbb{N}^{*}$, let us show\\
    \fbox{$pre(p,t)=(\omega,\mathtt{basic})\lor{}pre(p,t)=(\omega,\mathtt{test})\Rightarrow{}s'.M(p)\ge\omega$}
    and \\
    \fbox{$pre(p,t)=(\omega,\mathtt{inhib})\Rightarrow{}s'.M(p)<\omega$.}

    \begin{enumerate}
    \item Assuming
      $pre(p,t)=(\omega,\mathtt{basic})\lor{}pre(p,t)=(\omega,\mathtt{test})$,
      let us show \fbox{$s'.M(p)\ge\omega$.}

      \noindent{}The proceeding is the same for
      $pre(p,t)=(\omega,\mathtt{basic})$ and
      $pre(p,t)=(\omega,\mathtt{test})$. Therefore, we will only cover
      the case where $pre(p,t)=(\omega,\mathtt{basic})$.
      
      \noindent{}By property of the stabilize relation and \InCsCompT,
      equation \eqref{eq:eq-senabled-prod} holds.
      
      Rewriting $\sigma'(id_t)(\texttt{se})=\mathtt{true}$ with
      \eqref{eq:eq-senabled-prod}, we can deduce:\\
      $\prod\limits_{i=0}^{\Delta(id_t)(\texttt{ian})-1}\sigma'(id_t)(\texttt{iav})[i]=\mathtt{true}$.

      \noindent{}Then, we can deduce that
      $\forall{}i\in[0,\Delta(id_t)(\texttt{ian})-1],~\sigma'(id_t)(\texttt{iav})[i]=\mathtt{true}$.

      By construction, there exist an $id_p\in{}Comps(\Delta)$,
      $g_p$, $i_p$, $o_p$, $i\in[0,\vert{}input(t)\vert-1]$,
      $j\in{}[0,\vert{}output(p)\vert-1]$ and
      $id_{ji}\in{}Sigs(\Delta)$ s.t.  $\gamma(p)=id_p$ and\\
      \InCsCompP{} and
      ${<}\mathtt{output\_arcs\_valid(j)\Rightarrow{}id_{ji}}{>}\in{}o_p$
      and
      ${<}\mathtt{input\_arcs\_valid(i)\Rightarrow{}id_{ji}}{>}\in{}i_t$. Let
      us take such an $id_p\in{}Comps(\Delta)$, $g_p$, $i_p$,
      $o_p$, $i\in[0,\vert{}input(t)\vert-1]$,
      $j\in{}[0,\vert{}output(p)\vert-1]$ and
      $id_{ji}\in{}Sigs(\Delta)$.

      \noindent{}By construction,
      ${<}\mathtt{input\_arcs\_number\Rightarrow{}}\vert{}input(t)\vert{>}\in{}g_t$.

      \noindent{}By property of the elaboration relation and
      \InCsCompT, we can deduce
      $\Delta(id_t)(\texttt{ian})=\vert{}input(t)\vert$.

      \noindent{}Thanks to $\Delta(id_t)(\texttt{ian})=\vert{}input(t)\vert$,
      we can deduce that
      $\forall{}i\in[0,\vert{}input(t)\vert-1],$\\
      $\sigma'(id_t)(\texttt{iav})[i]=\mathtt{true}$.

      \noindent{}Having such an $i\in[0,\vert{}input(t)\vert-1]$, we
      can deduce that $\sigma'(id_t)(\texttt{iav})[i]=\mathtt{true}$.

      \noindent{}By property of the stabilize relation, \InCsCompT{}
      and \InCsCompP, we can deduce
      $\sigma'(id_t)(\texttt{iav})[i]=\sigma'(id_{ji})=\sigma'(id_p)(\texttt{oav})[j]$.

      \noindent{}Thanks to
      $\sigma'(id_t)(\texttt{iav})[i]=\sigma'(id_{ji})=\sigma'(id_p)(\texttt{oav})[j]$,
      we have $\sigma'(id_p)(\texttt{oav})[j]=\mathtt{true}$.

      By property of the stabilize relation and \InCsCompP, equation
      \eqref{eq:eq-oav-sens} holds. Thanks to \eqref{eq:eq-oav-sens},
      we can deduce that:
      \begin{equation}
        \label{eq:sens-true}
        \begin{split}
          \mathtt{true}=& \big((\sigma'(id_p)(\texttt{oat})[j]=\mathtt{basic}+\sigma'(id_p)(\texttt{oat})[j]=\mathtt{test}) \\
          & \quad.~\sigma'(id_p)(\texttt{sm})\ge\sigma'(id_p)(\texttt{oaw})[j]\big)\\
          & +\big(\sigma'(id_p)(\texttt{oat})[j]=\mathtt{inhib}~.~\sigma'(id_p)(\texttt{sm})<\sigma'(id_p)(\texttt{oaw})[j]\big)\\
        \end{split}
      \end{equation}

      \noindent{}By construction,
      ${<}\mathtt{output\_arcs\_types(j)\Rightarrow{}basic}{>}\in{}i_p$
      and\\
      ${<}\mathtt{output\_arcs\_weights(j)\Rightarrow{}}\omega{>}\in{}i_p$.

      \noindent{}By property of the stabilize relation and \InCsCompP,
      we can deduce $\sigma'(id_p)(\texttt{oat})[j]=\mathtt{basic}$ and
      $\sigma'(id_p)(\texttt{oaw})[j]=\omega$.

      \noindent{}Thanks to $\sigma'(id_p)(\texttt{oat})[j]=\mathtt{basic}$,
      $\sigma'(id_p)(\texttt{oaw})[j]=\omega$, and simplifying
      Equation~\eqref{eq:sens-true}, we can deduce
      $\sigma'(id_p)(\texttt{sm})\ge\omega=\mathtt{true}$.

      \noindent{}Appealing to Lemma~\ref{lem:re-equal-marking},
      \qedbox{$s'.M(p)\ge\omega$.}

    \item Assuming $pre(p,t)=(\omega,\mathtt{inhib})$, let us show
      \fbox{$s'.M(p)<\omega$.}

      The proceeding is the same as in the preceding case. Here, we
      will start the proof where the two cases are diverging, i.e:

      \noindent{}By construction,
      ${<}\mathtt{output\_arcs\_types(j)\Rightarrow{}inhib}{>}\in{}i_p$
      and\\
      ${<}\mathtt{output\_arcs\_weights(j)\Rightarrow{}}\omega{>}\in{}i_p$.

      \noindent{}By property of the stabilize relation and \InCsCompP,
      we can deduce $\sigma'(id_p)(\texttt{oat})[j]=\mathtt{inhib}$ and
      $\sigma'(id_p)(\texttt{oaw})[j]=\omega$.

      \noindent{}Thanks to $\sigma'(id_p)(\texttt{oat})[j]=\mathtt{inhib}$
      and $\sigma'(id_p)(\texttt{oaw})[j]=\omega$, and simplifying
      Equation~\eqref{eq:sens-true}, we can deduce
      $\sigma'(id_p)(\texttt{sm})<\omega=\mathtt{true}$.

      \noindent{}Appealing to Lemma~\ref{lem:re-equal-marking},
      \qedbox{$s'.M(p)<\omega$.}
    \end{enumerate}

  \end{enumerate}
\end{niproof}

\begin{lemma}[Rising edge equal not sensitized]
  \label{lem:re-equal-not-sens}
  \rehyps{} then\\
  $\forall{}t\in{}T,id_t\in{}Comps(\Delta)~s.t.~\gamma(t)=id_t,$
  $t\notin{}Sens(s'.M)\Leftrightarrow\sigma'(id_t)(\texttt{s\_enabled})=\mathtt{false}$.
\end{lemma}

\begin{niproof}
  Proving the above lemma is trivial by appealing to
  Lemma~\ref{lem:re-equal-sens} and by reasoning on
  contrapositives.
\end{niproof}

%%% Local Variables:
%%% mode: latex
%%% TeX-master: "../../main"
%%% End:


\section{Falling Edge}
\label{sec:falling-edge}
\subsection{Falling Edge and marking}
\label{sec:fe-marking}

%%%%%%%%%%%%%%%%%%%%%%%%%%%%%%%%%%%%%%%%%%%%%%%%%%%%%%
%%%%%%%%%% FALLING EDGE EQUAL MARKING LEMMA %%%%%%%%%%
%%%%%%%%%%%%%%%%%%%%%%%%%%%%%%%%%%%%%%%%%%%%%%%%%%%%%%

\begin{lemma}[Falling edge equal marking]
  \label{lem:fe-equal-marking}
  \fehyps{} then
  $\forall{}p\in{}P,id_p\in{}Comps(\Delta)~s.t.~\gamma(p)=id_p,$
  $~s'.M(p)=\sigma'(id_p)("s\_marking")$.
\end{lemma}

\begin{niproof}
  Given a $p\in{}P$ and an $id\in{}Comps(\Delta)$
  s.t. $\gamma(p)=id_p$, let us show\\
  \fbox{$s'.M(p)=\sigma'(id_p)("s\_marking")$.}

  By definition of
  $E_c,\tau\vdash{}sitpn,s\xrightarrow{\downarrow}s'$, we can deduce
  $s.M(p)=s'.M(p)$.

  By property of the $\mathtt{Inject}_\downarrow$ relation, the
  \hvhdl{} falling edge relation, the stabilize relation and
  \InCsCompP, and through the examination of the \texttt{marking}
  process defined in the place design architecture, we can deduce
  $\sigma'(id_p)("s\_marking")=\sigma(id_p)("s\_marking")$.

  Rewriting the goal with $s.M(p)=s'.M(p)$ and
  $\sigma'(id_p)("sm")=\sigma(id_p)("sm")$:
  \fbox{$s.M(p)=\sigma(id_p)("s\_marking")$.}

  \noindent{}By definition of
  $\gamma,E_c,\tau\vdash{}s\stackrel{\downarrow}{\approx}\sigma$:
  \qedbox{$s.M(p)=\sigma(id_p)("s\_marking")$.}
  
\end{niproof}

%%%%%%%%%%%%%%%%%%%%%%%%%%%%%%%%%%%%%%%%%%%%%%%%%%%%%%%%%%%%%%%
%%%%%%%%%% FALLING EDGE EQUAL OUTPUT TOKEN SUM LEMMA %%%%%%%%%%
%%%%%%%%%%%%%%%%%%%%%%%%%%%%%%%%%%%%%%%%%%%%%%%%%%%%%%%%%%%%%%%

\begin{lemma}[Falling Edge Equal Output Token Sum]
  \label{lem:fe-equal-ots}
  \fehyps{} then $\forall{}p,id_p~s.t.~\gamma(p)=id_p$,
  $\sum\limits_{t\in{}Fired(s')}pre(p,t)=\sigma'(id_p)("s\_output\_token\_sum")$.
\end{lemma}

\begin{niproof}
  Given a $p\in{}P$ and an $id_p\in{}Comps(\Delta)$, let us show\\
  \fbox{$\sum\limits_{t\in{}Fired(s')}pre(p,t)=\sigma'(id_p)("s\_output\_token\_sum")$.}

  \exP{}
  
  By property of the stabilize relation, \InCsCompP{}, and through the
  examination of the \texttt{output_tokens_sum} process defined in the
  place design architecture:
  \begin{equation}
    \label{eq:sots-at-fe}
    \sigma'(id_p)("sots")=\sum\limits_{i=0}^{\Delta(id_p)("oan")-1}
    \begin{cases}
      \sigma'(id_p)("oaw")[i]~\mathtt{if}~(\sigma'(id_p)("otf")[i]~ \\
      \hspace{19.5ex}.~\sigma'(id_p)("oat")[i]=\mathtt{basic}) \\
      0~otherwise \\
    \end{cases}
  \end{equation}

  Rewriting the goal with \eqref{eq:sots-at-fe}:\\
  \begin{equation*}
    \fbox{$
      \sum\limits_{t\in{}Fired(s')}pre(p,t)=\sum\limits_{i=0}^{\Delta(id_p)("oan")-1}
      \begin{cases}
        \sigma'(id_p)("oaw")[i]~\mathtt{if}~(\sigma'(id_p)("otf")[i]~ \\
        \hspace{19.5ex}.~\sigma'(id_p)("oat")[i]=\mathtt{basic}) \\
        0~otherwise \\
      \end{cases}$}
  \end{equation*}

  \noindent{}Let us unfold the definition of the left sum term:\\
  \fbox{
    {\begin{tabular}{c}
       $\sum\limits_{t\in{}Fired(s')}
       \begin{cases}
         \omega~\mathtt{if}~pre(p,t)=(\omega,\mathtt{basic}) \\
         0~otherwise
       \end{cases}$ \\
       $=$ \\
       $\sum\limits_{i=0}^{\Delta(id_p)("oan")-1}
       \begin{cases}
         \sigma'(id_p)("oaw")[i]~\mathtt{if}~(\sigma'(id_p)("otf")[i]~ \\
         \hspace{19.5ex}.~\sigma'(id_p)("oat")[i]=\mathtt{basic}) \\
         0~otherwise \\
       \end{cases}$ \\
     \end{tabular}}
 }\\

 \noindent{}To ease the reading, let us define functions
 $f\in{}Fired(s')\rightarrow\mathbb{N}$ and
 $g\in[0,\vert{}output(p)\vert-1]\rightarrow\mathbb{N}$ s.t.
 $f(t)=\begin{cases}
   \omega~\mathtt{if}~pre(p,t)=(\omega,\mathtt{basic}) \\
   0~otherwise
 \end{cases}$ and $g(i)=\begin{cases}
   \sigma'(id_p)("oaw")[i]~\mathtt{if}~(\sigma'(id_p)("otf")[i]~ \\
   \hspace{19.5ex}.~\sigma'(id_p)("oat")[i]=\mathtt{basic}) \\
   0~otherwise \\
 \end{cases}$

 \noindent{}Then, the goal is: \fbox{$\sum\limits_{t\in{}Fired(s')}f(t)=\sum\limits_{i=0}^{\Delta(id_p)("oan")-1}g(i)$}\\

 \noindent{}Let us perform case analysis on $output(p)$; there are two cases:

 \begin{itemize}
 \item \textbf{CASE} $output(p)=\emptyset$:
   
   By construction,
   ${<}$\texttt{output\_arcs\_number}$\mathtt{\Rightarrow{}1}{>}\in{}gm_p$,
   ${<}$\texttt{output\_arcs\_types(0)}$\Rightarrow{}\mathtt{basic}{>}\in{}ipm_p$,
   ${<}$\texttt{output\_transitions\_fired(0)}$\Rightarrow{}\mathtt{true}{>}\in{}ipm_p$,
   and
   ${<}$\texttt{output\_arcs\_weights(0)}$\mathtt{\Rightarrow{}0}{>}\in{}ipm_p$.

   By property of the elaboration relation and \InCsCompP, we can
   deduce $\Delta(id_p)("oan")=1$.

   By property of the stabilize relation and \InCsCompP, we can deduce
   $\sigma'(id_p)("oat")[0]=\mathtt{basic}$,
   $\sigma'(id_p)("otf")[0]=\mathtt{true}$ and
   $\sigma'(id_p)("oaw")[0]=0$.

   By property of $output(p)=\emptyset$, we can deduce\\
   $\sum\limits_{t\in{}Fired(s')}
   \begin{cases}
     \omega~\mathtt{if}~pre(p,t)=(\omega,\mathtt{basic}) \\
     0~otherwise
   \end{cases}=0$

   \noindent{}Rewriting the goal with $\Delta(id_p)("oan")=1$,
   $\sigma'(id_p)("oat")[0]=\mathtt{basic}$,
   $\sigma'(id_p)("otf")[0]=\mathtt{true}$,
   $\sigma'(id_p)("oaw")[0]=0$ and $\sum\limits_{t\in{}Fired(s')}
   \begin{cases}
     \omega~\mathtt{if}~pre(p,t)=(\omega,\mathtt{basic}) \\
     0~otherwise
   \end{cases}=0$, \qedbox{tautology.}
   
 \item \textbf{CASE} $output(p)\neq\emptyset$:\\

   \noindent{}By construction,
   ${<}\mathtt{output\_arcs\_number\Rightarrow{}}\vert{}output(p)\vert{>}\in{}gm_p$,
   and by property of the elaboration relation, we can deduce
   $\Delta(id_p)("oan")=\vert{}output(p)\vert$.
   
   Rewriting the goal with
   $\Delta(id_p)("oan")=\vert{}output(p)\vert$:
   \fbox{$\sum\limits_{t\in{}Fired(s')}f(t)=\sum\limits_{i=0}^{\vert{}output(p)\vert-1}g(i)$.}

   \noindent{}Let us reason by induction on the right sum term of the goal.

   \begin{itemize}
   \item \textbf{BASE CASE}:

     In that case, $0>\vert{}output\vert-1$ and
     $\sum\limits_{i=0}^{\vert{}output(p)\vert-1}g(i)=0$.

     \noindent{}As $0>\vert{}output\vert-1$, then
     $\vert{}output(p)\vert=0$, thus \qedbox{contradicting
       $output(p)\neq\emptyset$.}

   \item \textbf{INDUCTION CASE}:

     In that case, $0\le\vert{}output(p)\vert-1$.
     
     \begin{ih}
       $\forall{}F\subseteq{}Fired(s'),$
       $g(0)+\sum\limits_{t\in{}F}f(t)=g(0)+\sum\limits_{i=1}^{\vert{}output(p)\vert-1}g(i)$
     \end{ih}

     \fbox{$\sum\limits_{t\in{}Fired(s')}f(t)=g(0)+\sum\limits_{i=1}^{\vert{}output(p)\vert-1}g(i)$}

     By definition of $g$:
     \begin{equation}
       \label{eq:31}g(0)=\begin{cases}
         \sigma'(id_p)("oaw")[0]~\mathtt{if}~(\sigma'(id_p)("otf")[0] \\
         \hspace{19.5ex}.~\sigma'(id_p)("oat")[0]=\mathtt{basic}) \\
         0~otherwise \\
       \end{cases}
     \end{equation}

     \noindent{}Let us perform case analysis on the value of
     $\sigma'(id_p)("otf")[0]~.~\sigma'(id_p)("oat")[0]=\mathtt{basic}$; there are two cases:
     \begin{enumerate}
     \item $(\sigma'(id_p)("otf")[0]~.~\sigma'(id_p)("oat")[0]=\mathtt{basic})=\mathtt{false}$:\\
       In that case, $g(0)=0$, and then we can apply the induction
       hypothesis with $F=Fired(s')$ to solve the goal:
       \qedbox{$\sum\limits_{t\in{}Fired(s')}f(t)=\sum\limits_{i=1}^{\vert{}output(p)\vert-1}g(i)$.}

     \item $(\sigma'(id_p)("otf")[0]~.~\sigma'(id_p)("oat")[0]=\mathtt{basic})=\mathtt{true}$:\\
       In that case, $g(0)=\sigma'(id_p)("oaw")[0]$,
       $\sigma'(id_p)("otf")[0]=\mathtt{true}$ and
       $\sigma'(id_p)("oat")[0]=\mathtt{basic}$.

       \noindent{}By construction, there exist a $t\in{}output(t)$,
       $id_t\in{}Comps(\Delta)$ s.t. $\gamma(t)=id_t$, and there exist
       $gm_t$, $ipm_t$, $opm_t$ s.t. $\mathtt{comp}(id_t,$
       $"transition",$ $gm_t,$ $ipm_t,$ $opm_t)\in{}d.cs$, and there
       exist an $\omega\in\mathbb{N}^{*}$, an
       $a\in\{\mathtt{basic},\mathtt{test},\mathtt{inhib}\}$ and an
       $id_{ft}\in{}Sigs(\Delta)$ such that:
       \begin{itemize}
       \item $pre(p,t)=(\omega,a)$
       \item ${<}\mathtt{output\_arcs\_types(0)\Rightarrow{}}a{>}\in{}ipm_p$
       \item ${<}\mathtt{output\_arcs\_weights(0)\Rightarrow{}}\omega{>}\in{}ipm_p$
       \item ${<}\mathtt{fired\Rightarrow}id_{ft}{>}\in{}opm_t$
       \item
         ${<}\mathtt{output\_transitions\_fired(0)\Rightarrow{}id_{ft}}{>}\in{}ipm_p$
       \end{itemize}
       
       \noindent{}By property of the stabilize relation,
       $\sigma'(id_p)("oat")[0]=\mathtt{basic}$ and
       ${<}$\texttt{output\_arcs\_types(0)}$\mathtt{\Rightarrow{}a}{>}\in{}ipm_p$,
       we can deduce $pre(p,t)=(\omega,\mathtt{basic})$.

       \noindent{}By property of the stabilize relation,
       ${<}\mathtt{fired\Rightarrow{}id_{ft}}{>}\in{}opm_t$,\\
       ${<}\mathtt{output\_transitions\_fired(0)\Rightarrow{}id_{ft}}{>}\in{}ipm_p$
       and $\sigma'(id_p)("otf")[0]=\mathtt{true}$, we can deduce
       $\sigma'(id_t)("fired")=\mathtt{true}$.

       Appealing to Lemma~\ref{lem:fe-equal-fired}, and thanks to
       $\sigma'(id_t)("fired")=\mathtt{true}$, we can deduce
       $t\in{}Fired(s')$.

       With $t\in{}Fired(s')$, we can rewrite the left sum
       term of the goal as follows:\\
       \fbox{$f(t)+\sum\limits_{t'\in{}Fired(s')\setminus\{t\}}f(t')=g(0)+\sum\limits_{i=1}^{\vert{}output(p)\vert-1}g(i)$}

       \noindent{}We know that $g(0)=\sigma'(id_p)("oaw")[0]$, and by
       property of the stabilize relation and
       ${<}$\texttt{output\_arcs\_weights(0)}$\Rightarrow\omega{>}\in{}ipm_p$,
       we can deduce $\sigma'(id_p)("oaw")[0]=\omega$.

       Rewriting the goal with $\sigma'(id_p)("oaw")[0]=\omega$:\\
       \fbox{$f(t)+\sum\limits_{t'\in{}Fired(s')\setminus\{t\}}f(t')=\omega+\sum\limits_{i=1}^{\vert{}output(p)\vert-1}g(i)$}

       By definition of $f$, and as $pre(p,t)=(\omega,\mathtt{basic})$, then $f(t)=\omega$; thus, rewriting the goal:\\
       \fbox{$\omega+\sum\limits_{t'\in{}Fired(s')\setminus\{t\}}f(t')=\omega+\sum\limits_{i=1}^{\vert{}output(p)\vert-1}g(i)$}

       Then, knowing that $g(0)=\omega$, we can apply the induction
       hypothesis with $F=Fired(s')\setminus\{t\}$:
       \qedbox{$g(0)+\sum\limits_{t'\in{}Fired(s')\setminus\{t\}}f(t')=g(0)+\sum\limits_{i=1}^{\vert{}output(p)\vert-1}g(i)$.}
     \end{enumerate}
     
   \end{itemize}

 \end{itemize}
 
\end{niproof}


%%%%%%%%%%%%%%%%%%%%%%%%%%%%%%%%%%%%%%%%%%%%%%%%%%%%%%%%%%%%%%
%%%%%%%%%% FALLING EDGE EQUAL INPUT TOKEN SUM LEMMA %%%%%%%%%%
%%%%%%%%%%%%%%%%%%%%%%%%%%%%%%%%%%%%%%%%%%%%%%%%%%%%%%%%%%%%%%

\begin{lemma}[Falling Edge Equal Input Token Sum]
  \label{lem:fe-equal-its}
  \fehyps{} then $\forall{}p,id_p~s.t.~\gamma(p)=id_p$,
  $~\sum\limits_{t\in{}Fired(s')}post(t,p)=\sigma'_p("s\_input\_token\_sum")$.
\end{lemma}

\begin{niproof}
  Given a $p\in{}P$ and an $id_p\in{}Comps(\Delta)$, let us show\\
  \fbox{$\sum\limits_{t\in{}Fired(s')}post(t,p)=\sigma'(id_p)("s\_input\_token\_sum")$.}\\

  \exP{}
  
  By property of the stabilize relation, \InCsCompP, and through the
  examination of the \texttt{input_tokens_sum} process defined in the
  place design architecture:
  \begin{equation}
    \label{eq:sits-at-fe}
    \sigma'(id_p)("sits")=\sum\limits_{i=0}^{\Delta(id_p)("ian")-1}
    \begin{cases}
      \sigma'(id_p)("iaw")[i]~\mathtt{if}~\sigma'(id_p)("itf")[i]\\
      0~otherwise \\
    \end{cases}
  \end{equation}

  Rewriting the goal with \eqref{eq:sits-at-fe}:\\
  \begin{equation*}
    \fbox{$
      \sum\limits_{t\in{}Fired(s')}post(t,p)=\sum\limits_{i=0}^{\Delta(id_p)("ian")-1}
      \begin{cases}
        \sigma'(id_p)("iaw")[i]~\mathtt{if}~\sigma'(id_p)("otf")[i]\\
        0~otherwise \\
      \end{cases}$}
  \end{equation*}

  Let us unfold the definition of the left sum term:\\
  \begin{frameb}
    \begin{tabular}{c} $\sum\limits_{t\in{}Fired(s')}
      \begin{cases}
        \omega~\mathtt{if}~post(t,p)=\omega \\
        0~otherwise
      \end{cases}$ \\
      $=$ \\
      $\sum\limits_{i=0}^{\Delta(id_p)("ian")-1}
      \begin{cases}
        \sigma'(id_p)("iaw")[i]~\mathtt{if}~\sigma'(id_p)("itf")[i]\\
        0~otherwise \\
      \end{cases}$ \\
    \end{tabular}
  \end{frameb}
  
  Let us perform case analysis on $input(p)$; there are two cases:
  
  \begin{itemize}
  \item \textbf{CASE} $input(p)=\emptyset$:\\
    
    By construction,
    ${<}$\texttt{input\_arcs\_number}$\mathtt{\Rightarrow{}1}{>}\in{}gm_p$,
    ${<}$\texttt{input\_transitions\_fired(0)}$\Rightarrow{}\mathtt{true}{>}\in{}ipm_p$,
    and
    ${<}$\texttt{input\_arcs\_weights(0)}$\mathtt{\Rightarrow{}0}{>}\in{}ipm_p$.

    By property of the elaboration relation and \InCsCompP, we can
    deduce $\Delta(id_p)("ian")=1$.
    
    By property of the stabilize relation and \InCsCompP, we can deduce
    $\sigma'(id_p)("itf")[0]=\mathtt{true}$ and
    $\sigma'(id_p)("iaw")[0]=0$.

    By property of $input(p)=\emptyset$, we can deduce
    $\sum\limits_{t\in{}Fired(s')}
    \begin{cases}
      \omega~\mathtt{if}~post(t,p)=\omega \\
      0~otherwise
    \end{cases}=0$.

    \noindent{}Rewriting the goal with $\Delta(id_p)("ian")=1$,
    $\sigma'(id_p)("itf")[0]=\mathtt{true}$,
    $\sigma'(id_p)("iaw")[0]=0$, and $\sum\limits_{t\in{}Fired(s')}
    \begin{cases}
      \omega~\mathtt{if}~post(t,p)=\omega \\
      0~otherwise
    \end{cases}=0$, and simplifying the goal: \qedbox{tautology.}
    
  \item \textbf{CASE} $input(p)\neq\emptyset$:

    \noindent{}By construction,
    ${<}\mathtt{input\_arcs\_number\Rightarrow{}}\vert{}input(p)\vert{>}\in{}gm_p$,
    and by property of the elaboration relation, we can deduce
    \label{eq:24}$\Delta(id_p)("ian")=\vert{}input(p)\vert$.

    \noindent{}To ease the reading, let us define functions
    $f\in{}Fired(s')\rightarrow\mathbb{N}$ and
    $g\in[0,\vert{}input(p)\vert-1]\rightarrow\mathbb{N}$ s.t.
    $f(t)=\begin{cases}
      \omega~\mathtt{if}~post(t,p)=\omega \\
      0~otherwise
    \end{cases}$ and $g(i)=\begin{cases}
      \sigma'(id_p)("iaw")[i]~\mathtt{if}~\sigma'(id_p)("itf")[i]\\
      0~otherwise \\
    \end{cases}$

    \noindent{}Then, the goal is: \fbox{$\sum\limits_{t\in{}Fired(s')}f(t)=\sum\limits_{i=0}^{\Delta(id_p)("ian")-1}g(i)$}\\
    
    Rewriting the goal with $\Delta(id_p)("ian")=\vert{}input(p)\vert$:
    \fbox{$\sum\limits_{t\in{}Fired(s')}f(t)=\sum\limits_{i=0}^{\vert{}input(p)\vert-1}g(i)$.}

    \noindent{}Let us reason by induction on the right sum term of
    the goal.

    \begin{itemize}
    \item \textbf{BASE CASE}: In that case,
      $0>\vert{}input(p)\vert-1$ and
      $\sum\limits_{i=0}^{\vert{}input(p)\vert-1}g(i)=0$.

      As $0>\vert{}input(p)\vert-1$, then $\vert{}input(p)\vert=0$,
      thus \qedbox{contradicting $input(p)\neq\emptyset$.}
      
    \item \textbf{INDUCTION CASE}: In that case,
      $0\le\vert{}input(p)\vert-1$.
      
      \begin{ih}
        $\forall{}F\subseteq{}Fired(s'),~$
        $g(0)+\sum\limits_{t\in{}F}f(t)=g(0)+\sum\limits_{i=1}^{\vert{}input(p)\vert-1}g(i)$
      \end{ih}

      \fbox{$\sum\limits_{t\in{}Fired(s')}f(t)=g(0)+\sum\limits_{i=1}^{\vert{}input(p)\vert-1}g(i)$}\\

      By definition of $g$, we can deduce 
      $g(0)=\begin{cases}
        \sigma'(id_p)("iaw")[0]~\mathtt{if}~\sigma'(id_p)("itf")[0]\\
        0~otherwise \\
      \end{cases}$

      \noindent{}Let us perform case analysis on the value of
      $\sigma'(id_p)("itf")[0]$; there are two cases:
      \begin{enumerate}
      \item $\sigma'(id_p)("itf")[0]=\mathtt{false}$:\\
        In that case, $g(0)=0$, and then we can apply the induction
        hypothesis with $F=Fired(s')$ to solve the goal:
        \qedbox{$\sum\limits_{t\in{}Fired(s')}f(t)=\sum\limits_{i=1}^{\vert{}input(p)\vert-1}g(i)$.}

      \item $\sigma'(id_p)("itf")[0]=\mathtt{true}$:\\
        In that case, $g(0)=\sigma'(id_p)("iaw")[0]$ and
        $\sigma'(id_p)("itf")[0]=\mathtt{true}$ .

        \noindent{}By construction, there exist a $t\in{}input(t)$,
        an $id_t\in{}Comps(\Delta)$ s.t. $\gamma(t)=id_t$, $gm_t$,
        $ipm_t$, $opm_t$ s.t. $\mathtt{comp}(id_t,$ $"transition",$
        $gm_t,$ $ipm_t,$ $opm)\in{}d.cs$, an
        $\omega\in\mathbb{N}^{*}$ and an $id_{ft}\in{}Sigs(\Delta)$
        such that:
        \begin{itemize}
        \item $post(t,p)=\omega$
        \item
          ${<}\mathtt{input\_arcs\_weights(0)\Rightarrow{}}\omega{>}\in{}ipm_p$
        \item ${<}\mathtt{fired\Rightarrow}id_{ft}{>}\in{}opm_t$
        \item
          ${<}\mathtt{input\_transitions\_fired(0)\Rightarrow{}id_{ft}}{>}\in{}ipm_p$
        \end{itemize}
        
        By property of the stabilize relation,
        ${<}$\texttt{fired}$\Rightarrow{}\mathtt{id_{ft}}{>}\in{}opm_t$,
        ${<}$\texttt{input\_transitions\_fired(0)}$\mathtt{\Rightarrow{}id_{ft}}{>}\in{}ipm_p$
        and $\sigma'(id_p)("itf")[0]=\mathtt{true}$, we can deduce
        $\sigma'(id_t)("fired")=\mathtt{true}$.

        Appealing to Lemma~\ref{lem:fe-equal-fired} and
        $\sigma'(id_t)("fired")=\mathtt{true}$, we can deduce
        $t\in{}Fired(s')$.

        As $t\in{}Fired(s')$, we can rewrite the left sum
        term of the goal as follows:\\
        \fbox{$f(t)+\sum\limits_{t'\in{}Fired(s')\setminus\{t\}}f(t')=g(0)+\sum\limits_{i=1}^{\vert{}input(p)\vert-1}g(i)$}

        We know that $g(0)=\sigma'(id_p)("iaw")[0]$, and by property
        of the stabilize relation and
        ${<}$\texttt{input\_arcs\_weights(0)}$\Rightarrow\omega{>}\in{}ipm_p$,
        we can deduce $\sigma'(id_p)("iaw")[0]=\omega$.

        Rewriting the goal with $\sigma'(id_p)("iaw")[0]=\omega$:\\
        \fbox{$f(t)+\sum\limits_{t'\in{}Fired(s')\setminus\{t\}}f(t')=\omega+\sum\limits_{i=1}^{\vert{}input(p)\vert-1}g(i)$}

        By definition of $f$, and as $post(t,p)=\omega$, then $f(t)=\omega$; thus, rewriting the goal:\\
        \fbox{$\omega+\sum\limits_{t'\in{}Fired(s')\setminus\{t\}}f(t')=\omega+\sum\limits_{i=1}^{\vert{}input(p)\vert-1}g(i)$}

        Then, knowing that $g(0)=\omega$, we can apply the induction
        hypothesis with $F=Fired(s')\setminus\{t\}$:
        \qedbox{$g(0)+\sum\limits_{t'\in{}Fired(s')\setminus\{t\}}f(t')=g(0)+\sum\limits_{i=1}^{\vert{}input(p)\vert-1}g(i)$.}
      \end{enumerate}
      
    \end{itemize}
    
  \end{itemize}
  
\end{niproof}

%%%%%%%%%%%%%%%%%%%%%%%%%%%%%%%%%%%%%%%%%%%%%%%%%%%%%%
%%%%%%%%%% FALLING EDGE EQUAL TIME COUNTERS %%%%%%%%%%
%%%%%%%%%%%%%%%%%%%%%%%%%%%%%%%%%%%%%%%%%%%%%%%%%%%%%%

\subsection{Falling edge and time counters}
\label{sec:fe-equal-tc}

\begin{lemma}[Falling edge equal time counters]
  \label{lem:fe-equal-tc}
  \fehyps{} then $\forall{}t\in{}T_i,id_t\in{}Comps(\Delta)~s.t.~\gamma(t)=id_t,$\\
  $\big(upper(I_s(t))=\infty\land{}s'.I(t)\le{}lower(I_s(t))\Rightarrow{}s'.I(t)=\sigma'(id_t)("s\_time\_counter")\big)$\\
  $\land\big(upper(I_s(t))=\infty\land{}s'.I(t)>{}lower(I_s(t))\Rightarrow{}\sigma'(id_t)("s\_time\_counter")=lower(I_s(t))\big)$\\
  $\land\big(upper(I_s(t))\neq\infty\land{}s'.I(t)>{}upper(I_s(t))\Rightarrow{}\sigma'(id_t)("s\_time\_counter")=upper(I_s(t))\big)$\\
  $\land\big(upper(I_s(t))\neq\infty\land{}s'.I(t)\le{}upper(I_s(t))\Rightarrow{}s'.I(t)=\sigma'(id_t)("s\_time\_counter")\big)$.
\end{lemma}

\begin{niproof}
  Given a $t\in{}T_i$ and an $id_t\in{}Comps(\Delta)$ s.t. $\gamma(t)=id_t$, let us show\\
  \fbox{\parbox{\lwidth}{$\big(upper(I_s(t))=\infty\land{}s'.I(t)\le{}lower(I_s(t))\Rightarrow$
      $s'.I(t)=\sigma'(id_t)("s\_time\_counter")\big)$\\
      $\land\big(upper(I_s(t))=\infty\land{}s'.I(t)>{}lower(I_s(t))\Rightarrow$
      $\sigma'(id_t)("s\_time\_counter")=lower(I_s(t))\big)$\\
      $\land\big(upper(I_s(t))\neq\infty\land{}s'.I(t)>{}upper(I_s(t))\Rightarrow$
      $\sigma'(id_t)("s\_time\_counter")=upper(I_s(t))\big)$\\
      $\land\big(upper(I_s(t))\neq\infty\land{}s'.I(t)\le{}upper(I_s(t))\Rightarrow$
      $s'.I(t)=\sigma'(id_t)("s\_time\_counter")\big)$}}\\

  \exT{}

  \noindent{}By property of the elaboration,
  $\mathtt{Inject}_\downarrow$, \hvhdl{} rising edge and stabilize
  relations, \InCsCompT{}, and through the examination of the
  \texttt{time_counter} process defined in the transition design
  architecture, we can deduce:
  \begin{equation}
    \begin{split}
      \sigma(id_t)("se")=\mathtt{true}\land\Delta(id_t)("tt")\neq\mathtt{NOT\_TEMPORAL}
      \land\sigma(id_t)("srtc")=\mathtt{false}\\
      \land\sigma(id_t)("stc")<\Delta(id_t)("mtc")\Rightarrow
      \sigma'(id_t)("stc")=\sigma(id_t)("stc")+1
    \end{split}
    \label{eq:etnrlt}
  \end{equation}

  \begin{equation}
    \begin{split}
      \sigma(id_t)("se")=\mathtt{true}\land\Delta(id_t)("tt")\neq\mathtt{NOT\_TEMPORAL}
      \land\sigma(id_t)("srtc")=\mathtt{false}\\
      \land\sigma(id_t)("stc")\ge\Delta(id_t)("mtc")\Rightarrow
      \sigma'(id_t)("stc")=\sigma(id_t)("stc")
    \end{split}
    \label{eq:etnrge}
  \end{equation}

  \begin{equation}
    \begin{split}
      \sigma(id_t)("se")=\mathtt{true}\land\Delta(id_t)("tt")\neq\mathtt{NOT\_TEMPORAL}\\
      \land\sigma(id_t)("srtc")=\mathtt{true}\Rightarrow
      \sigma'(id_t)("stc")=1
    \end{split}
    \label{eq:etr}
  \end{equation}

  \begin{equation}
    \begin{split}
      \sigma(id_t)("se")=\mathtt{false}\lor\Delta(id_t)("tt")=\mathtt{NOT\_TEMPORAL}\Rightarrow
      \sigma'(id_t)("stc")=0
    \end{split}
    \label{eq:ne-or-nt}
  \end{equation}
  
  \noindent{}Then, there are 4 points to show:

  \begin{enumerate}
  \item\label{it:fe-eq-tc-fst}
    \fbox{$upper(I_s(t))=\infty\land{}s'.I(t)\le{}lower(I_s(t))\Rightarrow{}s'.I(t)=\sigma'(id_t)("s\_time\_counter")$}\\
    
    Assuming $upper(I_s(t))=\infty$ and
    $s'.I(t)\le{}lower(I_s(t))$, let us show\\
    \fbox{$s'.I(t)=\sigma'(id_t)("s\_time\_counter")$.}

    Let us perform case analysis on $t\in{}Sens(s.M)$; there are two
    cases:

    \begin{enumerate}
    \item \textbf{CASE} $t\notin{}Sens(s.M)$:\\
      By definition of \upSim, we can deduce
      $\sigma(id_t)("se")=\mathtt{false}$.

      Appealing to \eqref{eq:ne-or-nt} and
      $\sigma(id_t)("se")=\mathtt{false}$, we can deduce
      $\sigma'(id_t)("stc")=0$.

      By definition of \dwSitpn{} (Rule~\ref{it:reset-counters}), we
      can deduce $s'.I(t)=0$.

      Rewriting the goal with $\sigma'(id_t)("stc")=0$ and
      $s'.I(t)=0$: \qedbox{tautology.}
      
    \item \textbf{CASE} $t\in{}Sens(s.M)$:
      
      By definition of \upSim, we can deduce
      $\sigma(id_t)("se")=\mathtt{true}$.

      By construction, and as $upper(I_s(t))=\infty$,
      ${<}\mathtt{transition\_type\Rightarrow{}TEMP\_A\_INF}{>}\in{}gm_t$. By
      property of the elaboration relation, we have
      $\Delta(id_t)("tt")=\mathtt{TEMP\_A\_INF}$.

      Let us perform case analysis on $s.reset_t(t)$; there are two
      cases:
      \begin{enumerate}
      \item \textbf{CASE} $s.reset_t(t)=\mathtt{true}$:

        By definition of \upSim, $\sigma(id_t)("srtc")=\mathtt{true}$.

        Appealing to \eqref{eq:etr},
        $\sigma(id_t)("se")=\mathtt{true}$,
        $\Delta(id_t)("tt")=\mathtt{TEMP\_A\_INF}$ and
        $\sigma(id_t)("srtc")=\mathtt{true}$, we can deduce
        $\sigma'(id_t)("stc")=1$.

        By definition of \dwSitpn (Rule~\ref{it:reset-counters}), we
        can deduce $s'.I(t)=1$.

        Rewriting the goal with $\sigma'(id_t)("stc")=1$ and
        $s'.I(t)=1$: \qedbox{tautology.}

      \item \textbf{CASE} $s.reset_t(t)=\mathtt{false}$:
        
        By definition of \upSim, we have
        $\sigma(id_t)("srtc")=\mathtt{false}$.

        As $upper(I_s(t))=\infty$, there exists an
        $a\in\mathbb{N}^{*}$ s.t. $I_s(t)=[a,\infty]$. Let us take
        such an $a\in\mathbb{N}^{*}$. By construction,
        ${<}\mathtt{maximal\_time\_counter\Rightarrow}~a{>}\in{}gm_t$,
        and by property of the elaboration relation, we have
        $\Delta(id_t)("mtc")=a$.

        By definition of \dwSitpn (Rule~\ref{it:inc-counters}), and
        knowing that $t\in{}Sens(s.M)$, $s.reset_t(t)=\mathtt{false}$
        and $upper(I_s(t))=\infty$, we can deduce $s'.I(t)=s.I(t)+1$.

        Rewriting the goal with $s'.I(t)=s.I(t)+1$:
        \fbox{$s.I(t)+1=\sigma'(id_t)("stc")$.}
        
        We assumed that $s'.I(t)\le{}lower(I_s(t))$, and as
        $s'.I(t)=s.I(t)+1$, then $s.I(t)+1\le{}lower(I_s(t))$, then
        $s.I(t)<lower(I_s(t))$, then $s.I(t)<a$ since
        $a=lower(I_s(t))$.

        \noindent{}By definition of
        $\gamma,E_c,\tau\vdash{}s\stackrel{\uparrow}{\approx}\sigma$,
        and knowing that $s.I(t)<lower(I_s(t))$ and
        $upper(I_s(t))=\infty$, we can deduce
        $s.I(t)=\sigma(id_t)("stc")$.

        Appealing to $\Delta(id_t)("mtc")=a$,
        $s.I(t)=\sigma(id_t)("stc")$ and $s.I(t)<a$, we can deduce
        $\sigma(id_t)("stc")<\Delta(id_t)("mtc")$.
        
        Appealing to \eqref{eq:etnrlt},
        $\sigma(id_t)("stc")<\Delta(id_t)("mtc")$,
        $\sigma(id_t)("srtc")=\mathtt{false}$ and\\
        $\sigma(id_t)("se")=\mathtt{true}$, we can deduce:
        $\sigma'(id_t)("stc")=\sigma(id_t)("stc")+1$.

        Rewriting the goal with
        $\sigma'(id_t)("stc")=\sigma(id_t)("stc")+1$ and
        $s.I(t)=\sigma(id_t)("stc")$: \qedbox{tautology.}
      \end{enumerate}
    \end{enumerate}
    
  \item
    \fbox{$upper(I_s(t))=\infty\land{}s'.I(t)>{}lower(I_s(t))\Rightarrow$
      $\sigma'(id_t)("s\_time\_counter")=lower(I_s(t)$.}

    Assuming that $upper(I_s(t))=\infty$ and
    $s'.I(t)>{}lower(I_s(t))$, let us show\\
    \fbox{$\sigma'(id_t)("s\_time\_counter")=lower(I_s(t))$.}

    As $upper(I_s(t))=\infty$, there exists an $a\in\mathbb{N}^{*}$
    s.t. $I_s(t)=[a,\infty]$. Let us take such an
    $a\in\mathbb{N}^{*}$.

    By construction,
    ${<}$\texttt{maximal\_time\_counter}$\Rightarrow{}a{>}\in{}gm_t$,
    and
    ${<}$\texttt{transition\_type}$\Rightarrow{}$\texttt{TEMP\_A\_INF}${>}\in{}gm_t$
    by property of the elaboration relation, we can deduce
    $\Delta(id_t)("mtc")=a$ and
    $\Delta(id_t)("tt")=\mathtt{TEMP\_A\_INF}$.
    
    Let us perform case analysis on $t\in{}Sens(s.M)$:
    \begin{enumerate}
    \item \textbf{CASE} $t\notin{}Sens(s.M)$:
      
      By definition of $E_c,\tau\vdash{}s\xrightarrow{\downarrow}s'$
      (Rule~\ref{it:reset-not-sens}), and knowing that
      $t\in{}Sens(s.M)$, we can deduce $s'.I(t)=0$. Since
      $lower(I_s(t))\in\mathbb{N}^{*}$, then $lower(I_s(t))>0$.
      
      \qedbox{Contradicts $s'.I(t)>lower(I_s(t))$.}
      
    \item \textbf{CASE} $t\in{}Sens(s.M)$:
      
      By definition of
      $\gamma,E_c,\tau\vdash{}s\stackrel{\uparrow}{\sim}\sigma$ and
      $t\in{}Sens(s.M)$, we can deduce
      $\sigma(id_t)("se")=\mathtt{true}$.

      Let us perform case analysis on $s.reset_t(t)$; there are two
      cases:
      \begin{enumerate}
      \item \textbf{CASE} $s.reset_t(t)=\mathtt{true}$:
        
        By definition of
        $E_c,\tau\vdash{}s\xrightarrow{\downarrow}s'$: $s'.I(t)=1$.

        We assumed that $s'.I(t)>lower(I_s(t))$, then
        $1>lower(I_s(t))$.

        \qedbox{Contradicts $lower(I_s(t))>0$.}

      \item \textbf{CASE} $s.reset_t(t)=\mathtt{false}$:
        
        By property of
        $\gamma,E_c,\tau\vdash{}s\stackrel{\uparrow}{\approx}\sigma$
        and $s.reset_t(t)=\mathtt{false}$, we can deduce
        $\sigma(id_t)("srtc")=\mathtt{false}$.

        By definition of $E_c,\tau\vdash{}s\xrightarrow{\downarrow}s'$
        (Rule~\ref{it:inc-counters}), and knowing that
        $s'.I(t)>lower(I_s(t))$, we can deduce
        \begin{equation*}
          \begin{split}
            s'.I(t)=s.I(t)+1&\Rightarrow{}s.I(t)+1>lower(I_s(t))\\
            & \Rightarrow{}s.I(t)\ge{}lower(I_s(t))\\
          \end{split}
        \end{equation*}

        Let us perform case analysis on $s.I(t)\ge{}lower(I_s(t))$:
        \begin{enumerate}
        \item \textbf{CASE} $s.I(t)>lower(I_s(t))$: \fbox{$\sigma'(id_t)("stc")=lower(I_s(t))$.}\\
          By definition of
          $\gamma,E_c,\tau\vdash{}s\stackrel{\uparrow}{\approx}\sigma$,
          we can deduce $\sigma(id_t)("stc")=lower(I_s(t))$.

          Appealing to \eqref{eq:etnrge}, we can deduce
          $\sigma'(id_t)("stc")=\sigma(id_t)("stc")$.

          Rewriting the goal with
          $\sigma'(id_t)("stc")=\sigma(id_t)("stc")$ and
          $\sigma(id_t)("stc")=lower(I_s(t))$: \qedbox{tautology.}
          
        \item \textbf{CASE} $s.I(t)=lower(I_s(t))$: \fbox{$\sigma'(id_t)("stc")=lower(I_s(t))$.}\\
          By definition of
          $\gamma,E_c,\tau\vdash{}s\stackrel{\uparrow}{\approx}\sigma$,
          we can deduce $s.I(t)=\sigma(id_t)("stc")$.

          Appealing to \eqref{eq:etnrge}, we can deduce
          $\sigma'(id_t)("stc")=\sigma(id_t)("stc")$.

          Rewriting the goal with
          $\sigma'(id_t)("stc")=\sigma(id_t)("stc")$,
          $s.I(t)=\sigma(id_t)("stc")$ and $s.I(t)=lower(I_s(t))$:
          \qedbox{tautology.}
        \end{enumerate}
      \end{enumerate}
    \end{enumerate}
  \item
    \fbox{$upper(I_s(t))\neq\infty\land{}s'.I(t)>{}upper(I_s(t))\Rightarrow$
      $\sigma'(id_t)("s\_time\_counter")=upper(I_s(t))$.}

    Assuming that $upper(I_s(t))\neq\infty$ and
    $s'.I(t)>{}upper(I_s(t))$, let us show \\
    \fbox{$\sigma'(id_t)("s\_time\_counter")=upper(I_s(t))$.}

    As $upper(I_s(t))\neq\infty$, there exists an
    $a\in\mathbb{N}^{*}$, and a $b\in\mathbb{N}^{*}$
    s.t. $I_s(t)=[a,b]$. Let us take such an $a$ and $b$.

    By construction,
    ${<}\mathtt{maximal\_time\_counter\Rightarrow}b{>}\in{}gm_t$ and
    there exists $tt\in\{\mathtt{TEMP\_A\_A,TEMP\_A\_B\}}$ s.t.
    ${<}\mathtt{transition\_type\Rightarrow}tt{>}\in{}gm_t$.

    By property of the elaboration relation and \InCsCompT, we can
    deduce $\Delta(id_t)("mtc")=b=upper(I_s(t))$ and
    $\Delta(id_t)("tt")\neq\mathtt{NOT\_TEMP}$.

    Let us perform case analysis on $t\in{}Sens(s.M)$:
      \begin{enumerate}
      \item \textbf{CASE} $t\notin{}Sens(s.M)$:\\
        By definition of $E_c,\tau\vdash{}s\xrightarrow{\downarrow}s'$
        (Rule~\ref{it:reset-not-sens}), and knowing that
        $t\in{}Sens(s.M)$, then $s'.I(t)=0$. Since
        $upper(I_s(t))\in\mathbb{N}^{*}$, then $upper(I_s(t))>0$.
        
        \qedbox{Contradicts $s'.I(t)>upper(I_s(t))$.}
        
      \item \textbf{CASE} $t\in{}Sens(s.M)$:\\
        By definition of
        $\gamma,E_c,\tau\vdash{}s\stackrel{\uparrow}{\approx}\sigma$
        and $t\in{}Sens(s.M)$, we can deduce
        $\sigma(id_t)("se")=\mathtt{true}$.

        Let us perform case analysis on $s.reset_t(t)$; there are two
        cases:
        \begin{enumerate}
        \item \textbf{CASE} $s.reset_t(t)=\mathtt{true}$:\\
          By definition of
          $E_c,\tau\vdash{}s\xrightarrow{\downarrow}s'$
          (Rule~\ref{it:reset-counters}), we can deduce $s'.I(t)=1$.

          \noindent{}We assumed that $s'.I(t)>upper(I_s(t))$, then we
          can deduce $1>upper(I_s(t))$.

          \qedbox{Contradicts $upper(I_s(t))>0$.}

        \item \textbf{CASE} $s.reset_t(t)=\mathtt{false}$:\\
          By property of
          $\gamma,E_c,\tau\vdash{}s\stackrel{\uparrow}{\approx}\sigma$
          and $s.reset_t(t)=\mathtt{false}$, we can deduce
          $\sigma(id_t)("srtc")=\mathtt{false}$.
          
          Let us perform case analysis on $s.I(t)>{}upper(I_s(t))$ or
          $s.I(t)\le{}upper(I_s(t))$:
          \begin{enumerate}
          \item \textbf{CASE} $s.I(t)>upper(I_s(t))$: \fbox{$\sigma'(id_t)("stc")=upper(I_s(t))$.}\\
            By definition of \dwSitpn (Rule~\ref{it:locked-counters}),
            we can deduce $s'.I(t)=s.I(t)$.
            
            By definition of
            $\gamma,E_c,\tau\vdash{}s\stackrel{\uparrow}{\approx}\sigma$,
            we can deduce $\sigma(id_t)("stc")=upper(I_s(t))$.

            Appealing to \eqref{eq:etnrge}, we have
            $\sigma'(id_t)("stc")=\sigma(id_t)("stc")$.

            Rewriting the goal with
            $\sigma'(id_t)("stc")=\sigma(id_t)("stc")$ and
            $\sigma(id_t)("stc")=upper(I_s(t))$: \qedbox{tautology.}
            
          \item \textbf{CASE} $s.I(t)\le{}upper(I_s(t))$: \fbox{$\sigma'(id_t)("stc")=upper(I_s(t))$.}\\

            By definition of \upSim, we can deduce
            $s.I(t)=\sigma(id_t)("stc")$.
            
            Let us perform case analysis on
            $s.I(t)\le{}upper(I_s(t))$; there are two cases:
            \begin{itemize}
            \item \textbf{CASE} $s.I(t)=upper(I_s(t))$:\\

              Appealing to $\Delta(id_t)("mtc")=b=upper(I_s(t))$,
              $s.I(t)=\sigma(id_t)("stc")$ and $s.I(t)=upper(I_s(t))$,
              we can deduce
              $\Delta(id_t)("mtc")\le\sigma(id_t)("stc")$.

              Appealing to $\Delta(id_t)("mtc")\le\sigma(id_t)("stc")$
              and \eqref{eq:etnrge}, we can deduce
              $\sigma'(id_t)("stc")=\sigma(id_t)("stc")$.

              Rewriting the goal with
              $\sigma'(id_t)("stc")=\sigma(id_t)("stc")$,
              $s.I(t)=\sigma(id_t)("stc")$ and $s.I(t)=upper(I_s(t))$:
              \qedbox{tautology.}
              
            \item \textbf{CASE} $s.I(t)<upper(I_s(t))$:\\
              
              By definition of \dwSitpn (Rule~\ref{it:inc-counters}),
              we can deduce $s'.I(t)=s.I(t)+1$.

              From $s'.I(t)=s.I(t)+1$ and $s.I(t)<upper(I_s(t))$, we
              can deduce $s'.I(t)\le{}upper(I_s(t))$;
              \qedbox{contradicts $s'.I(t)>upper(I_s(t))$.}
            \end{itemize}
          \end{enumerate}
        \end{enumerate}
      \end{enumerate}
      
    \item
      \fbox{$upper(I_s(t))\neq\infty\land{}s'.I(t)\le{}upper(I_s(t))\Rightarrow$
        $s'.I(t)=\sigma'(id_t)("s\_time\_counter")$.}

      Assuming that $upper(I_s(t))\neq\infty$ and
      $s'.I(t)\le{}upper(I_s(t))$, let us show\\
      \fbox{$s'.I(t)=\sigma'(id_t)("s\_time\_counter")$.}

      As $upper(I_s(t))\neq\infty$, there exists an
      $a\in\mathbb{N}^{*}$, and a $b\in\mathbb{N}^{*}$
      s.t. $I_s(t)=[a,b]$. Let us take such an $a$ and $b$.

      By construction,
      ${<}\mathtt{maximal\_time\_counter\Rightarrow}b{>}\in{}gm_t$ and
      there exists $tt\in\{$\texttt{TEMP\_A\_A,TEMP\_A\_B}$\}$ s.t.
      ${<}\mathtt{transition\_type\Rightarrow}tt{>}\in{}gm_t$; by
      property of the elaboration relation, we can deduce
      $\Delta(id_t)("mtc")=b=upper(I_s(t))$ and
      $\Delta(id_t)("tt")\neq\mathtt{NOT\_TEMP}$.
      
      Let us perform case analysis on $t\in{}Sens(s.M)$:
      \begin{enumerate}
      \item \textbf{CASE} $t\notin{}Sens(s.M)$:\\
        
        By definition of \upSim, we have
        $\sigma(id_t)("se")=\mathtt{false}$.

        \noindent{}Appealing \eqref{eq:ne-or-nt} and
        $\sigma(id_t)("se")=\mathtt{false}$, we have
        $\sigma'(id_t)("stc")=0$.

        By definition of \dwSitpn (Rule~\ref{it:reset-not-sens}), we
        have $s'.I(t)=0$.

        Rewriting the goal with $\sigma'(id_t)("stc")=0$ and
        $s'.I(t)=0$: \qedbox{tautology.}
        
      \item \textbf{CASE} $t\in{}Sens(s.M)$:\\
        
        By definition of \upSim, we have
        $\sigma(id_t)("se")=\mathtt{true}$.

        Let us perform case analysis on $s.reset_t(t)$:
        \begin{enumerate}
        \item \textbf{CASE} $s.reset_t(t)=\mathtt{true}$:\\
          
          By definition of \upSim, we have
          $\sigma(id_t)("srtc")=\mathtt{true}$.

          Appealing to \eqref{eq:etr},
          $\Delta(id_t)("tt")\neq\mathtt{NOT\_TEMP}$,
          $\sigma(id_t)("se")=\mathtt{true}$ and
          $\sigma(id_t)("srtc")=\mathtt{true}$, we have
          $\sigma'(id_t)("stc")=1$.

          By definition of \dwSitpn (Rule~\ref{it:reset-counters}), we
          have $s'.I(t)=1$.

          Rewriting the goal with $\sigma'(id_t)("stc")=1$ and
          $s'.I(t)=1$, \qedbox{tautology.}
          
        \item \textbf{CASE} $s.reset_t(t)=\mathtt{false}$:\\
          
          By definition of \upSim, we have
          $\sigma(id_t)("srtc")=\mathtt{false}$.

          Let us perform case analysis on $s.I(t)>upper(I_s(t))$ or
          $s.I(t)\le{}upper(I_s(t))$:
          \begin{enumerate}
          \item \textbf{CASE} $s.I(t)>upper(I_s(t))$:\\
            
            By definition of \dwSitpn, we have $s.I(t)=s'.I(t)$, and
            thus, $s'.I(t)>upper(I_s(t))$. \qedbox{Contradicts
              $s'.I(t)\le{}upper(I_s(t))$.}
            
          \item \textbf{CASE} $s.I(t)\le{}upper(I_s(t))$:\\
            
            By definition of \upSim, we have
            $s.I(t)=\sigma(id_t)("stc")$.

            \begin{itemize}
            \item \textbf{CASE} $s.I(t)<upper(I_s(t))$:\\
              From $s.I(t)<upper(I_s(t))$,
              $s.I(t)=\sigma(id_t)("stc")$ and
              $\Delta(id_t)("mtc")=b=upper(I_s(t))$, we can deduce
              $\sigma(id_t)("stc")<\Delta(id_t)("mtc")$.

              From \eqref{eq:etnrlt},
              $\sigma(id_t)("se")=\mathtt{true}$,
              $\Delta(id_t)("tt")\neq\mathtt{NOT\_TEMP}$,
              $\sigma(id_t)("srtc")=\mathtt{false}$ and
              $\sigma(id_t)("stc")<\Delta(id_t)("mtc")$, we can deduce
              $\sigma'(id_t)("stc")=\sigma(id_t)("stc")+1$.

              By definition of \dwSitpn (Rule~\ref{it:inc-counters}),
              we can deduce $s'.I(t)=s.I(t)+1$.

              Rewriting the goal with
              $\sigma'(id_t)("stc")=\sigma(id_t)("stc")+1$ and
              $s'.I(t)=s.I(t)+1$, \qedbox{tautology.}
              
            \item \textbf{CASE} $s.I(t)=upper(I_s(t))$:\\
              By definition of \dwSitpn (Rule~\ref{it:inc-counters}),
              we know that $s'.I(t)=s.I(t)+1$. We assumed that
              $s'.I(t)\le{}upper(I_s(t))$; thus,
              $s.I(t)+1\le{}upper(I_s(t))$.

              \qedbox{Contradicts $s.I(t)=upper(I_s(t))$.}
              
            \end{itemize}
            
          \end{enumerate}
          
        \end{enumerate}
        
      \end{enumerate}
    \end{enumerate}
\end{niproof}

\subsection{Falling edge and condition values}
\label{sec:fe-cond-values}

\begin{lemma}[Falling edge equal condition values]
  \label{lem:fe-equal-cond-values}
  \fehyps{} then
  $\forall{}c\in\mathcal{C},id_c\in{}Ins(\Delta)~s.t.~\gamma(c)=id_c,~s'.cond(c)=\sigma'(id_c)$.
\end{lemma}

\begin{niproof}
  Given a $c\in\mathcal{C}$ and an $id_c\in{}Ins(\Delta)$
  s.t. $\gamma(c)=id_c$, let us show
  \fbox{$s'.cond(c)=\sigma'(id_c)$.}

  By definition of \dwSitpn (Rule~\ref{it:cond-env}), we have
  $s'.cond(c)=E_c(\tau,c)$.

  By property of the $\mathtt{Inject_\downarrow}$, the \hvhdl{}
  falling edge, the stabilize relations and $id_c\in{}Ins(\Delta)$, we
  have $\sigma'(id_c)=E_p(\tau,\downarrow)(id_c)$.

  \noindent{}Rewriting the goal with $s'.cond(c)=E_c(\tau,c)$ and
  $\sigma'(id_c)=E_p(\tau,\downarrow)(id_c)$:
  \fbox{$E_c(\tau,c)=E_p(\tau,\downarrow)(id_c)$}
  
  \noindent{}By definition of $\gamma\vdash{}E_p\stackrel{env}{=}E_c$:
  \qedbox{$E_c(\tau,c)=E_p(\tau,\downarrow)(id_c)$.}
  
\end{niproof}

\subsection{Falling and action executions}
\label{sec:fe-equal-act-exec}

\begin{lemma}[Falling edge equal action executions]
  \label{lem:fe-equal-act-exec}
  \fehyps{} then
  $\forall{}a\in\mathcal{A},id_a\in{}Outs(\Delta)~s.t.~\gamma(a)=id_a,~s'.ex(a)=\sigma'(id_a)$.
\end{lemma}

\begin{niproof}
  Given an $a\in\mathcal{A}$ and an $id_a\in{}Outs(\Delta)$
  s.t. $\gamma(a)=id_a$, let us show \fbox{$s'.ex(a)=\sigma'(id_a)$.}

  \noindent{}By property of \dwSitpn (Rule~\ref{it:activate-actions}):
  \begin{equation}
    s'.ex(a)=\sum\limits_{p\in{}marked(s.M)}\mathbb{A}(p,a)\label{eq:fe-eq-exa}
  \end{equation}
  
  By construction, the generated \texttt{action} process is a part of
  design $d$'s behavior, i.e there exist an
  $sl\subseteq{}Sigs(\Delta)$ and an $ss_a\in{}ss$ s.t.
  $\mathtt{ps}("action", \emptyset, sl, ss)\in{}d.cs$.
  
  By construction $id_a$ is only assigned in the body of the
  \texttt{action} process during the initialization or a falling edge
  phase.

  Let $pls(a)$ be the set of actions associated to action $a$, i.e
  $pls(a)=\{p\in{}P~\vert~\mathbb{A}(p,a)=true\}$. Then, depending on
  $pls(a)$, there are two cases of assignment of output port $id_a$:
  
  \begin{itemize}
  \item \textbf{CASE} $pls(a)=\emptyset$:\\
    \noindent{}By construction,
    $\mathtt{id_a\Leftarrow{}false}\in{}ss_{a\downarrow}$ where
    $ss_{a\downarrow}$ is the part of the \texttt{``action''} process
    body executed during a falling edge phase.

    \noindent{}By property of the \hvhdl{} falling edge relation, the
    stabilize relation and
    $\mathtt{ps}("action", \emptyset, sl, ss_a)\in{}d.cs$, we can
    deduce $\sigma'(id_a)=false$.
    
    \noindent{}By property of
    $\sum\limits_{p\in{}marked(s.M)}\mathbb{A}(p,a)$ and
    $pls(a)=\emptyset$, we can deduce
    $\sum\limits_{p\in{}marked(s.M)}\mathbb{A}(p,a)=\mathtt{false}$.

    \noindent{}Rewriting the goal with \eqref{eq:fe-eq-exa},
    $\sigma'(id_a)=false$ and
    $\sum\limits_{p\in{}marked(s.M)}\mathbb{A}(p,a)=\mathtt{false}$,
    \qedbox{tautology.}
    
  \item \textbf{CASE} $pls(a)\neq\emptyset$:\\
    \noindent{}By construction,
    $\mathtt{id_a\Leftarrow{}id_{mp_0}+\dots+id_{mp_n}}\in{}ss_{a\downarrow}$,
    where $id_{mp_i}\in{}Sigs(\Delta)$, $ss_{a\downarrow}$ is the part
    of the \texttt{action} process body executed during the falling
    edge phase, and $n=\vert{}pls(a)\vert-1$.

    \noindent{}By property of the $\mathtt{Inject}_\downarrow$, the
    \hvhdl{} falling edge relation, the stabilize relation, and\\
    $\mathtt{ps}("action",$ $\emptyset,$ $sl,$ $ss)\in{}d.cs$:
    \begin{equation}
      \sigma'(id_a)=\sigma(id_{mp_0})+\dots+\sigma(id_{mp_n})\label{eq:fe-eq-ida-sum}
    \end{equation}

    Rewriting the goal with \eqref{eq:fe-eq-exa} and
    \eqref{eq:fe-eq-ida-sum}:
    \fbox{$\sum\limits_{p\in{}marked(s.M)}\mathbb{A}(p,a)=\sigma(id_{mp_0})+\dots+\sigma(id_{mp_n})$.}

    Let us reason on the value of $\sigma(id_{mp_0})+\dots+\sigma(id_{mp_n})$; there are two cases:

    \begin{itemize}
    \item \textbf{CASE} $\sigma(id_{mp_0})+\dots+\sigma(id_{mp_n})=\mathtt{true}$:\\
      \noindent{}Then, we can rewrite the goal as follows:
      \fbox{$\sum\limits_{p\in{}marked(s.M)}\mathbb{A}(p,a)=\mathtt{true}$.}

      \noindent{}To prove the above goal, let us show
      \fbox{$\exists{}p\in{}marked(s.M)~s.t.~\mathbb{A}(p,a)=\mathtt{true}$.}

      \noindent{}From
      $\sigma(id_{mp_0})+\dots+\sigma(id_{mp_n})=\mathtt{true}$, we
      can deduce that
      $\exists{}id_{mp_i}~s.t.~\sigma(id_{mp_i})=\mathtt{true}$. Let
      us take an $id_{mp_i}$ s.t. $\sigma(id_{mp_i})=\mathtt{true}$.

      \noindent{}By construction, there exist a $p\in{}pls(a)$, an
      $id_{p}\in{}Comps(\Delta)$, $gm_{p}$, $ipm_{p}$ and
      $opm_{p}$ such that:
      \begin{itemize}
      \item $\gamma(p)=id_{p}$
      \item \InCsCompP
      \item ${<}\mathtt{marked\Rightarrow{}id_{mp_i}}{>}\in{}opm_{p}$
      \end{itemize}
      Let us take such a $p$, $id_{p}$, $gm_{p}$, $ipm_{p}$ and
      $opm_{p}$.

      By property of stable $\sigma$ and \InCsCompP, we can deduce
      $\sigma(id_{mp_i})=\sigma(id_{p})("marked")$.

      By property of stable $\sigma$, \InCsCompP, and through the
      examination of the \texttt{determine_marked} process defined in
      the place design architecture, we can deduce:
      \begin{equation}
        \sigma(id_{p})("marked")=\sigma(id_{p})("sm")>0\label{eq:fe-gt-sm-zero}
      \end{equation}

      \noindent{}From $\sigma(id_{mp_i})=\sigma(id_{p})("marked")$,
      \eqref{eq:fe-gt-sm-zero} and $\sigma(id_{mp_i})=\mathtt{true}$,
      we can deduce that $\sigma(id_{p})("mar\-ked")=\mathtt{true}$
      and $(\sigma(id_{p})("sm")>0)=\mathtt{true}$.
      
      \noindent{}By property of \upSim, we have
      $s.M(p)=\sigma(id_{p})("sm")$.
      
      \noindent{}From $s.M(p)=\sigma(id_{p})("sm")$ and
      $(\sigma(id_{p})("sm")>0)=\mathtt{true}$, we can deduce
      $p\in{}marked(s.M)$, i.e $s.M(p)>0$.

      Let us use $p$ to prove the goal:
      \fbox{$\mathbb{A}(p,a)=\mathtt{true}$.}

      By definition of $p\in{}pls(a)$,
      \qedbox{$\mathbb{A}(p,a)=\mathtt{true}$.}

    \item \textbf{CASE} $\sigma(id_{mp_0})+\dots+\sigma(id_{mp_n})=\mathtt{false}$:\\
      
      Then, we can rewrite the goal as follows:
      \fbox{$\sum\limits_{p\in{}marked(s.M)}\mathbb{A}(p,a)=\mathtt{false}$.}

      \noindent{}To prove the above goal, let us show
      \fbox{$\forall{}p\in{}marked(s.M)~s.t.~\mathbb{A}(p,a)=\mathtt{false}$.}

      \noindent{}Given a $p\in{}marked(s.M)$, let us show
      \fbox{$\mathbb{A}(p,a)=\mathtt{false}$.}

      \noindent{}Let us perform case analysis on $\mathbb{A}(p,a)$;
      there are 2 cases:

      \begin{itemize}
      \item \textbf{CASE} \qedbox{$\mathbb{A}(p,a)=\mathtt{false}$.}
      \item \textbf{CASE} $\mathbb{A}(p,a)=\mathtt{true}$:\\
        By construction, there exist an $id_{p}\in{}Comps(\Delta)$,
        $gm_{tp}$, $ipm_{p}$, $opm_{p}$ and
        $id_{mp_i}\in{}Sigs(\Delta)$ such that:
        \begin{itemize}
        \item $\gamma(p)=id_{p}$
        \item \InCsCompP
        \item ${<}\mathtt{marked\Rightarrow{}id_{mp_i}}{>}\in{}opm_{p}$
        \end{itemize}

        Let us take such a $id_{p}$, $gm_{p}$, $ipm_{p}$, $opm_{p}$
        and $id_{mp_i}$.

        By property of stable $\sigma$, \InCsCompP, and
        ${<}\mathtt{marked\Rightarrow{}id_{mp_i}}{>}\in{}opm_{p}$, we
        can deduce $\sigma(id_{mp_i})=\sigma(id_{p})("marked")$.

        By property of stable $\sigma$, \InCsCompP, and through the
        examination of the \texttt{determine_marked} process defined
        in the place design architecture, we can deduce:
        \begin{equation}
          \sigma(id_{p})("marked")=(\sigma(id_{p})("sm")>0)\label{eq:fe-eq-marked-gt-sm-zero}
        \end{equation}

        From
        $\sigma(id_{mp_0})+\dots+\sigma(id_{mp_n})=\mathtt{false}$, we
        can deduce $\sigma(id_{mp_i})=\mathtt{false}$.

        From $\sigma(id_{p})("marked")=\mathtt{false}$, we can deduce
        $(\sigma(id_{p})("sm")>0)=\mathtt{false}$.

        By definition of \upSim, we have $s.M(p)=\sigma(id_p)("sm")$,
        and thus, we can deduce that $s.M(p)=0$ (equivalent to
        $(s.M(p)>0)=\mathtt{false}$).

        \noindent{}Contradicts \qedbox{$p\in{}marked(s.M)$} (i.e,
        $s.M(p)>0$).
      \end{itemize}
    \end{itemize}
  \end{itemize}
\end{niproof}

\subsection{Falling edge and function executions}
\label{sec:fe-fun-exec}

\begin{lemma}[Falling edge equal function executions]
  \label{lem:fe-equal-fun-exec}
  \fehyps{} then
  $\forall{}f\in\mathcal{F},id_f\in{}Outs(\Delta)~s.t.~\gamma(f)=id_f,~s'.ex(f)=\sigma'(id_f)$.
\end{lemma}

\begin{niproof}
  Given an $f\in\mathcal{F}$ and an
  $id_f\in{}Outs(\Delta)~s.t.~\gamma(f)=id_f$, let us show
  \fbox{$s'.ex(f)=\sigma'(id_f)$.}\\

  \noindent{}By property of \dwSitpn, we can deduce
  $s.ex(f)=s'.ex(f)$.
  
  \noindent{}By construction, $id_f$ is an output port identifier of
  boolean type in the \hvhdl{} design $d$ assigned by the
  \texttt{function} process only during the initialization or during a
  rising edge phase.

  \noindent{}By property of the \hvhdl{} $\mathtt{Inject_{\uparrow}}$,
  rising edge, stabilize relations, and the \texttt{function} process,
  we can deduce $\sigma(id_f)=\sigma'(id_f)$.

  \noindent{}Rewriting the goal with $s.ex(f)=s'.ex(f)$ and
  $\sigma(id_f)=\sigma'(id_f)$, \fbox{$s.ex(f)=\sigma(id_f)$.}

  \noindent{}By definition of \upSim, \qedbox{$s.ex(f)=\sigma(id_f)$.}
\end{niproof}

\subsection{Falling edge and firable transitions}
\label{sec:fe-firable}

%%%%%%%%%%%%%%%%%%%%%%%%%%%%%%%%%%%%%%%%%%%%%%%%
%%%%%%%%%% FALLING EDGE EQUAL FIRABLE %%%%%%%%%%
%%%%%%%%%%%%%%%%%%%%%%%%%%%%%%%%%%%%%%%%%%%%%%%%

\begin{lemma}[Falling edge equal firable]
  \label{lem:fe-equal-firable}
  \fehyps{} then
  $\forall{}t\in{}T,id_t\in{}Comps(\Delta)~s.t.~\gamma(t)=id_t,$
  $t\in{}Firable(s')\Leftrightarrow\sigma'(id_t)("s\_firable")=\mathtt{true}$.
\end{lemma}

\begin{niproof}
  Given a $t\in{}T$ and $id_t\in{}Comps(\Delta)$
  s.t. $\gamma(t)=id_t$, let us show that\\
  \fbox{$t\in{}Firable(s')\Leftrightarrow\sigma'(id_t)("s\_firable")=\mathtt{true}$.}\\

  The proof is in two parts:
  \begin{enumerate}
  \item Assuming that $t\in{}Firable(s')$, let us show
    \fbox{$\sigma'(id_t)("s\_firable")=\mathtt{true}$.}
    
    Appealing to Lemma~\ref{lem:fe-equal-firable-1}:
    \qedbox{$\sigma'(id_t)("s\_firable")=\mathtt{true}$.}
    
  \item Assuming that $\sigma'(id_t)("s\_firable")=\mathtt{true}$, let
    us show \fbox{$t\in{}Firable(s')$.}
    
    Appealing to Lemma~\ref{lem:fe-equal-firable-2}:
    \qedbox{$t\in{}Firable(s')$.}
  \end{enumerate}
  
\end{niproof}

\begin{lemma}[Falling edge equal firable 1]
  \label{lem:fe-equal-firable-1}
  \fehyps{} then
  $\forall{}t\in{}T,id_t\in{}Comps(\Delta)~s.t.~\gamma(t)=id_t,$
  $t\in{}Firable(s')\Rightarrow\sigma'(id_t)("s\_firable")=\mathtt{true}$.
\end{lemma}

\begin{niproof}
  Given a $t\in{}T$ and $id_t\in{}Comps(\Delta)$
  s.t. $\gamma(t)=id_t$, and assuming that $t\in{}Firable(s')$, let us
  show \fbox{$\sigma'(id_t)("s\_firable")=\mathtt{true}$.}\\
  
  \exT{}

  By property of the $\mathtt{Inject}_\downarrow$ relation, the
  \hvhdl{} falling edge relation, the stabilize relation, \InCsCompT,
  and through the examination of the \texttt{firable} process defined
  in the transition design architecture, we can deduce:
  \begin{equation}
    \label{eq:fe-eq-sfa}
    \sigma'(id_t)("sfa")=\sigma(id_t)("se")~.~\sigma(id_t)("scc")~.~\mathtt{checktc}(\Delta(id_t),\sigma(id_t))
  \end{equation}

  Term $\mathtt{checktc}(\Delta(id_t),\sigma(id_t))$ is defined as
  follows:
  \begin{equation}
    \label{eq:checktc}
    \begin{split}
      \mathtt{checktc}(\Delta(id_t),\sigma(id_t))=&\bigg(\mathtt{not}~\sigma(id_t)("srtc")~.~\\
      & \quad\begin{split}
        \big[\big(\Delta(id_t)("tt")=\mathtt{TEMP\_A\_B}~&.~(\sigma(id_t)("stc")\ge{}\sigma(id_t)("A")-1)~\\
        &.~(\sigma(id_t)("stc")\le{}\sigma(id_t)("B")-1)\big)\\
      \end{split} \\
      & \quad+(\Delta(id_t)("tt")=\mathtt{TEMP\_A\_A}~.~(\sigma(id_t)("stc")={}\sigma(id_t)("A")-1))\\
      &
      \quad+(\Delta(id_t)("tt")=\mathtt{TEMP\_A\_INF}~.~(\sigma(id_t)("stc")\ge{}\sigma(id_t)("A")-1))\big]\bigg) \\
      & +\big(\sigma(id_t)("srtc")~.~\Delta(id_t)("tt")\neq\mathtt{NOT\_TEMP}~.~\sigma(id_t)("A")=1\big)\\
      & +\Delta(id_t)("tt")=\mathtt{NOT\_TEMP}\\
    \end{split}
  \end{equation}

  Rewriting the goal with \eqref{eq:fe-eq-sfa}:
  \fbox{$\sigma(id_t)("se")~.~\sigma(id_t)("scc")~.~\mathtt{checktc}(\Delta(id_t),\sigma(id_t))=\mathtt{true}$.}

  Then, there are three points to prove:
  \begin{enumerate}
  \item \fbox{$\sigma(id_t)("se")=\mathtt{true}$}:\\

    From $t\in{}Firable(s')$, we can deduce
    $t\in{}Sens(s'.M)$. By definition of \dwSitpn, we have $s.M=s'.M$,
    and thus, we can deduce $t\in{}Sens(s.M)$.

    By definition of \upSim, we know that $t\in{}Sens(s.M)$
    implies \qedbox{$\sigma(id_t)("se")=\mathtt{true}$.}
    
  \item \fbox{$\sigma(id_t)("scc")=\mathtt{true}$}:\\

    By definition of \upSim:
    \begin{equation}
      \sigma(id_t)("scc")=\prod\limits_{c\in{}conds(t)}
      \begin{cases}
        E_c(\tau,c) & if~\mathbb{C}(t,c)=1 \\
        \mathtt{not}(E_c(\tau,c)) & if~\mathbb{C}(t,c)=-1 \\
      \end{cases}
      \label{eq:fe-eq-scc-prod}
    \end{equation}
    where
    $conds(t)=\{c\in\mathcal{C}~\vert~\mathbb{C}(t,c)=1\lor\mathbb{C}(t,c)=-1\}$.

    Rewriting the goal with \eqref{eq:fe-eq-scc-prod}:
    \fbox{$\prod\limits_{c\in{}conds(t)}
      \begin{cases}
        E_c(\tau,c) & if~\mathbb{C}(t,c)=1 \\
        \mathtt{not}(E_c(\tau,c)) & if~\mathbb{C}(t,c)=-1 \\
      \end{cases}=\mathtt{true}$.}

    To ease the reading, let us define $f(c)=\begin{cases}
          E_c(\tau,c) & if~\mathbb{C}(t,c)=1 \\
          \mathtt{not}(E_c(\tau,c)) & if~\mathbb{C}(t,c)=-1 \\
        \end{cases}$.
    
    Let us reason by induction on the left term of the
    goal:

    \begin{itemize}
    \item \textbf{BASE CASE}: \qedbox{$\mathtt{true}=\mathtt{true}$.}
    \item \textbf{INDUCTION CASE}:
      \begin{ih}
        $\prod\limits_{c'\in{}conds(t)\setminus\{c\}}f(c')=\mathtt{true}$
      \end{ih}

      \fbox{$f(c)~.~\prod\limits_{c'\in{}conds(t)\setminus\{c\}}f(c')=\mathtt{true}$.}

      Rewriting the goal with the induction hypothesis, simplifying
      the goal, and unfolding the definition of $f(c)$:
      \fbox{$\begin{cases}
          E_c(\tau,c) & if~\mathbb{C}(t,c)=1 \\
          \mathtt{not}(E_c(\tau,c)) & if~\mathbb{C}(t,c)=-1 \\
        \end{cases}=\mathtt{true}$.}

      As $c\in{}conds(t)$, let us perform case analysis on
      $\mathbb{C}(t,c)=1\lor\mathbb{C}(t,c)=-1$:
      \begin{enumerate}
      \item \textbf{CASE} $\mathbb{C}(t,c)=1$: \fbox{$E_c(\tau,c)=\mathtt{true}$.}\\
        
        By definition of $t\in{}Firable(s')$, we can deduce that
        $s'.cond(c)=\mathtt{true}$. By definition of \dwSitpn{}
        (Rule~\ref{it:cond-env}), we have
        $s'.cond(c)=E_c(\tau,c)$. Thus,
        \qedbox{$E_c(\tau,c)=\mathtt{true}$.}
        
      \item $\mathbb{C}(t,c)=-1$: \fbox{$\mathtt{not}~E_c(\tau,c)=\mathtt{true}$.}\\

        By definition of $t\in{}Firable(s')$, we can deduce that
        $s'.cond(c)=\mathtt{false}$. By definition of \dwSitpn{}
        (Rule~\ref{it:cond-env}), we have
        $s'.cond(c)=E_c(\tau,c)$. Thus,
        \qedbox{$\mathtt{not}~E_c(\tau,c)=\mathtt{true}$.}
      \end{enumerate}
    \end{itemize}

  \item \fbox{$\mathtt{checktc}(\Delta(id_t),\sigma(id_t))=\mathtt{true}$}:\\

    By definition of $t\in{}Firable(s')$, we have
    $t\notin{}T_i\lor{}s'.I(t)\in{}I_s(t)$. Let us perform case
    analysis on $t\notin{}T_i\lor{}s'.I(t)\in{}I_s(t)$:\\

    \begin{enumerate}
    \item \textbf{CASE} $t\notin{}T_i$: \fbox{$\mathtt{checktc}(\Delta(id_t),\sigma(id_t))=\mathtt{true}$}\\

      By construction,
      ${<}\mathtt{transition\_type\Rightarrow{}NOT\_TEMP}{>}\in{}gm_t$,
      and by property of the elaboration relation, we have
      $\Delta(id_t)("tt")=\mathtt{NOT\_TEMP}$.

      From $\Delta(id_t)("tt")=\mathtt{NOT\_TEMP}$, and by definition
      of $\mathtt{checktc}(\Delta(id_t),\sigma(id_t))$, we can deduce
      \qedbox{$\mathtt{checktc}(\Delta(id_t),\sigma(id_t))=\mathtt{true}$.}

    \item \textbf{CASE} $s'.I(t)\in{}I_s(t)$: \fbox{$\mathtt{checktc}(\Delta(id_t),\sigma(id_t))=\mathtt{true}$}\\

      From $s'.I(t)\in{}I_s(t)$, we can deduce that $t\in{}T_i$. Thus,
      by construction, there exists
      $tt\in\{\mathtt{TEMP\_A\_B},\mathtt{TEMP\_A\_A},\mathtt{TEMP\_A\_INF}\}$
      s.t. ${<}\mathtt{transition\_type\Rightarrow}tt{>}\in{}gm_t$. By
      property of the elaboration relation, we have
      $\Delta(id_t)("tt")=tt$, and thus, we know
      $\Delta(id_t)("tt")\neq{}\mathtt{NOT\_TEMP}$. Therefore, we can
      simplfy the term $\mathtt{checktc}(\Delta(id_t),\sigma(id_t))$
      as follows:
      \begin{equation}        
        \label{eq:eq-checktc-minus-not-temp}
        \begin{split}
          \mathtt{checktc}(\Delta(id_t),\sigma(id_t))=&\bigg(\mathtt{not}~\sigma(id_t)("srtc")~.~\\
          & \quad\begin{split}
            \big[\big(\Delta(id_t)("tt")=\mathtt{TEMP\_A\_B}~&.~(\sigma(id_t)("stc")\ge{}\sigma(id_t)("A")-1)~\\
            &.~(\sigma(id_t)("stc")\le{}\sigma(id_t)("B")-1)\big)\\
          \end{split} \\
          & \quad{\begin{array}{l}
                    +(\Delta(id_t)("tt")=\mathtt{TEMP\_A\_A}~.\\
                    \quad~(\sigma(id_t)("stc")={}\sigma(id_t)("A")-1)) \\
                  \end{array}}\\
          &
          \quad{\begin{array}{l}
                  +(\Delta(id_t)("tt")=\mathtt{TEMP\_A\_INF}~. \\
                  \quad~(\sigma(id_t)("stc")\ge{}\sigma(id_t)("A")-1))\big]\bigg) \\
                \end{array}}\\
          & +\big(\sigma(id_t)("srtc")~.~\sigma(id_t)("A")=1\big)\\
        \end{split}
      \end{equation}
      
      By definition of \upSim, we have
      $s.reset_t(t)=\sigma(id_t)("srtc")$.

      Let us perform case analysis on the value
      $s.reset_t(t)$:

      \begin{enumerate}
      \item \textbf{CASE} $s.reset_t(t)=\mathtt{true}$: \fbox{$\mathtt{checktc}(\Delta(id_t),\sigma(id_t))=\mathtt{true}$}\\

        From $s.reset_t(t)=\sigma(id_t)("srtc")$, we can deduce that
        $\sigma(id_t)("srtc")=\mathtt{true}$.

        From $\sigma(id_t)("srtc")=\mathtt{true}$, we can
        simplify the term
        $\mathtt{checktc}(\Delta(id_t),\sigma(id_t))$ as follows:
        \begin{equation}
          \label{eq:eq-checktc-a-eq-1}
            \mathtt{checktc}(\Delta(id_t),\sigma(id_t))=\big(\sigma(id_t)("A")=1\big)
        \end{equation}

        Rewriting the goal with \eqref{eq:eq-checktc-a-eq-1}, and
        simplifying the goal: \fbox{$\sigma(id_t)("A")=1$.}

        By definition of \dwSitpn{} (Rule~\ref{it:reset-counters}), from
        $t\in{}Sens(s.M)$ and $s.reset_t(t)=\mathtt{true}$, we can
        deduce $s'.I(t)=1$. We know that $s'.I(t)\in{}I_s(t)$, and
        thus, we have $1\in{}I_s(t)$.

        By definition of $1\in{}I_s(t)$, there exist an
        $a\in\mathbb{N}^{*}$ and a
        $ni\in{}\mathbb{N}^{*}\sqcup\{\infty\}$ s.t. $I_s(t)=[a,ni]$
        and $1\in[a,ni]$.

        By definition of $1\in[a,ni]$, we have $a\le{}1$, and
        since $a\in\mathbb{N}^{*}$, we can deduce $a=1$.

        By construction,
        ${<}\mathtt{time\_A\_value\Rightarrow}{}a{>}\in{}ipm_t$, and
        by property of stable $\sigma$, we have
        \qedbox{$\sigma(id_t)("A")=a=1$.}

      \item \textbf{CASE} $s.reset_t(t)=\mathtt{false}$: \fbox{$\mathtt{checktc}(\Delta(id_t),\sigma(id_t))=\mathtt{true}$}\\

        From $s.reset_t(t)=\sigma(id_t)("srtc")$, we can deduce
        $\sigma(id_t)("srtc")=\mathtt{false}$.

        From $\sigma(id_t)("srtc")=\mathtt{false}$, we can simplify
        the term $\mathtt{checktc}(\Delta(id_t),\sigma(id_t))$ as
        follows:
        \begin{equation}
          \label{eq:eq-checktc-srtc-false}
          \begin{array}{l}
            \multicolumn{1}{c}{\mathtt{checktc}(\Delta(id_t),\sigma(id_t))}\\
            \multicolumn{1}{c}{=}\\
            \begin{array}{ll}
              \big(\Delta(id_t)("tt")=\mathtt{TEMP\_A\_B}~&.~(\sigma(id_t)("stc")\ge{}\sigma(id_t)("A")-1)~\\
                                                          &.~(\sigma(id_t)("stc")\le{}\sigma(id_t)("B")-1)\big)\\
            \end{array}\\
            +(\Delta(id_t)("tt")=\mathtt{TEMP\_A\_A}~.~(\sigma(id_t)("stc")={}\sigma(id_t)("A")-1))\\
            +(\Delta(id_t)("tt")=\mathtt{TEMP\_A\_INF}~.~(\sigma(id_t)("stc")\ge{}\sigma(id_t)("A")-1)) \\
          \end{array}
        \end{equation}

        Let us perform case analysis on $I_s(t)$; there are two cases:
        \begin{itemize}
        \item \textbf{CASE} $I_s(t)=[a,b]$ where
          $a,b\in\mathbb{N}^{*}$; then, either $a=b$ or $a\neq{}b$:
          
          \begin{itemize}
          \item \textbf{CASE} $a=b$:\\
            Then, we have $I_s(t)=[a,a]$, and by construction
            ${<}$\texttt{transition\_type}$\Rightarrow{}$
            \texttt{TEMP\_A\_A}${>}\in{}gm_t$. By property of the
            elaboration relation, we have
            $\Delta(id_t)$$("tt")=$\texttt{TEMP\_A\_A}; thus we can
            simplify the $\mathtt{checktc}$ term as follows:
            \begin{equation}
              \label{eq:eq-checktc-temp-a-a}
              \mathtt{checktc}(\Delta(id_t),\sigma(id_t))=(\sigma(id_t)("stc")={}\sigma(id_t)("A")-1)
            \end{equation}

            Rewriting the goal with \eqref{eq:eq-checktc-temp-a-a},
            and simplifying the goal:\\
            \fbox{$\sigma(id_t)("stc")={}\sigma(id_t)("A")-1$.}\\
            
            From $s'.I(t)\in[a,a]$, we can deduce that
            $s'.I(t)=a$. Let us perform case analysis on
            $s.I(t)<upper(I_s(t))$ or $s.I(t)\ge{}upper(I_s(t))$:
            
            \begin{itemize}
            \item \textbf{CASE} $s.I(t)<upper(I_s(t))$:\\
              By definition of \upSim, we have
              $s.I(t)=\sigma(id_t)("stc")$. By definition of
              \dwSitpn{} (Rule~\ref{it:inc-counters}), we have
              $s'.I(t)=s.I(t)+1$. From $s'.I(t)=a$ and
              $s'.I(t)=s.I(t)+1$, we can deduce $a-1=s.I(t)$.

              By construction,
              ${<}\mathtt{time\_A\_value\Rightarrow}{}a{>}\in{}ipm_t$,
              and by property of stable $\sigma$, we have
              $\sigma(id_t)("A")=a$.

              Rewriting the goal with $\sigma(id_t)("A")=a$,
              $s.I(t)=\sigma(id_t)("stc")$, and $a-1=s.I(t)$:
              \qedbox{tautology.}
              
            \item \textbf{CASE} $s.I(t)\ge{}upper(I_s(t))$:\\
              In the case where $s.I(t)>upper(I_s(t))$, then
              $s.I(t)>a$.  By definition of \dwSitpn{}
              (Rule~\ref{it:locked-counters}), we have
              $s.I(t)=s'.I(t)=a$. Then, \qedbox{$a>a$ is a
                contradiction.}\\

              In the case where $s.I(t)=upper(I_s(t))$, then
              $s.I(t)=a$. By definition of \dwSitpn{}
              (Rule~\ref{it:inc-counters}), we have
              $s'.I(t)=s.I(t)+1$. Then, we have $s'.I(t)=a$ and
              $s'.I(t)=a+1$. Then, \qedbox{$a=a+1$ is a
                contradiction.}
            \end{itemize}
            
          \item \textbf{CASE} $a\neq{}b$: \fbox{$\mathtt{checktc}(\Delta(id_t),\sigma(id_t))=\mathtt{true}$}\\
            Then, we have $I_s(t)=[a,b]$, and by construction
            ${<}$\texttt{transition\_type}$\Rightarrow{}$
            \texttt{TEMP\_A\_B}${>}\in{}gm_t$. By property of the
            elaboration relation, we have
            $\Delta(id_t)$$("tt")=$\texttt{TEMP\_A\_B}; thus we can
            simplify the term $\mathtt{checktc}$ as follows:
            \begin{equation}
              \label{eq:eq-checktc-temp-a-b}
              \begin{array}{c}
                \multicolumn{1}{c}{\mathtt{checktc}(\Delta(id_t),\sigma(id_t))}\\
                \multicolumn{1}{c}{=}\\
                (\sigma(id_t)("stc")\ge{}\sigma(id_t)("A")-1)~.~(\sigma(id_t)("stc")\le{}\sigma(id_t)("B")-1)\\
              \end{array}\\
            \end{equation}

            Rewriting the goal with \eqref{eq:eq-checktc-temp-a-b},
            and simplifying the goal:\\
            \fbox{$(\sigma(id_t)("stc")\ge{}\sigma(id_t)("A")-1)\land(\sigma(id_t)("stc")\le{}\sigma(id_t)("B")-1)$.}\\
            
            Let us perform case analysis on $s.I(t)<upper(I_s(t))$ or
            $s.I(t)\ge{}upper(I_s(t))$:
            \begin{itemize}
            \item \textbf{CASE} $s.I(t)<upper(I_s(t))$:\\
              By definition of \upSim, we have
              $s.I(t)=\sigma(id_t)("stc")$. By definition of \dwSitpn
              (Rule~\ref{it:inc-counters}), we have
              $s'.I(t)=s.I(t)+1$. By definition of $s'.I(t)\in[a,b]$:
              
              $\Rightarrow{}a\le{}s'.I(t)\le{}b$.
              
              $\Rightarrow{}a\le{}s'.I(t)\land{}s'.I(t)\le{}b$

              $\Rightarrow{}a\le{}s.I(t)+1\land{}s.I(t)+1\le{}b$

              $\Rightarrow{}a-1\le{}s.I(t)\land{}s.I(t)\le{}b-1$

              By construction,
              ${<}\mathtt{time\_A\_value\Rightarrow}{}a{>}\in{}ipm_t$
              and
              ${<}\mathtt{time\_B\_value\Rightarrow}{}b{>}\in{}ipm_t$,
              and by property of stable $\sigma$, we have
              $\sigma(id_t)("A")=a$ and
              $\sigma(id_t)("B")=b$.

              Rewriting the goal with
              $\sigma(id_t)("A")=a$, $\sigma(id_t)("B")=b$ and
              $s.I(t)=\sigma(id_t)("stc")$:
              \qedbox{$a-1\le{}s.I(t)\land{}s.I(t)\le{}b-1$.}
              
            \item \textbf{CASE} $s.I(t)\ge{}upper(I_s(t))$:\\
              
              In the case where $s.I(t)>upper(I_s(t))$, then
              $s.I(t)>b$. By definition of \dwSitpn{}
              (Rule~\ref{it:locked-counters}), we have
              $s.I(t)=s'.I(t)=b$. Then, \qedbox{$b>b$ is a
                contradiction.}\\

              In the case where $s.I(t)=upper(I_s(t))$, then
              $s.I(t)=b$. By definition of \dwSitpn{}
              (Rule~\ref{it:inc-counters}), we have
              $s'.I(t)=s.I(t)+1$.

              By definition of $s'.I(t)\in[a,b]$, we have
              $s'.I(t)\le{}b$:

              $\Rightarrow{}s.I(t)+1\le{}b$

              $\Rightarrow{}$ \qedbox{$b+1\le{}b$ is contradiction.}
            \end{itemize}
          \end{itemize}
          
        \item \textbf{CASE} $I_s(t)=[a,\infty]$ where
          $a\in{}\mathbb{N}^{*}$: \fbox{$\mathtt{checktc}(\Delta(id_t),\sigma(id_t))=\mathtt{true}$}
 
          By construction ${<}\mathtt{transition\_type\Rightarrow{}}$
          $\mathtt{TEMP\_A\_INF}{>}\in{}gm_t$. By property of the
          elaboration relation, we have
          $\Delta(id_t)$$("tt")=\mathtt{TEMP\_A\_INF}$; thus we can
          simplify the term $\mathtt{checktc}$ as follows:
          \begin{equation}
            \label{eq:eq-checktc-temp-a-inf}
            \mathtt{checktc}(\Delta(id_t),\sigma(id_t))=(\sigma(id_t)("stc")\ge{}\sigma(id_t)("A")-1))
          \end{equation}

          Rewriting the goal with \eqref{eq:eq-checktc-temp-a-inf},
          and simplifying the goal:\\
          \fbox{$\sigma(id_t)("stc")\ge{}\sigma(id_t)("A")-1$.}\\

          From $s'.I(t)\in[a,\infty]$, we can deduce
          $a\le{}s'.I(t)$. Then, let us perform case analysis on
          $s.I(t)\le{}lower(I_s(t))$ or $s.I(t)>lower(I_s(t))$:
          \begin{itemize}
          \item \textbf{CASE} $s.I(t)\le{}lower(I_s(t))$:

            By definition of \upSim, we have
            $s.I(t)=\sigma(id_t)("stc")$.

            By definition of \dwSitpn{} (Rule~\ref{it:inc-counters}),
            we have $s'.I(t)=s.I(t)+1$:

            $\Rightarrow{}s'.I(t)\ge{}a$

            $\Rightarrow{}s.I(t)+1\ge{}a$

            $\Rightarrow{}s.I(t)\ge{}a-1$

            By construction,
            ${<}\mathtt{time\_A\_value\Rightarrow}{}a{>}\in{}ipm_t$,
            and by property of stable $\sigma$, we have
            $\sigma(id_t)("A")=a$.

            Rewriting the goal with $\sigma(id_t)("A")=a$ and
            $s.I(t)=\sigma(id_t)("stc")$:\\ \qedbox{$s.I(t)\ge{}a-1$.}
            
          \item \textbf{CASE} $s.I(t)>lower(I_s(t))$:

            By definition of \upSim, we have
            $\sigma(id_t)("stc")=lower(I_s(t))=a$.

            By construction,
            ${<}\mathtt{time\_A\_value\Rightarrow}{}a{>}\in{}ipm_t$,
            and by property of stable $\sigma$, we have
            $\sigma(id_t)("A")=a$.

            Rewriting the goal with $\sigma(id_t)("stc")=a$ and
            $\sigma(id_t)("A")=a$: \qedbox{$a\ge{}a-1$.}
          \end{itemize}
          
        \end{itemize}
      \end{enumerate}
    \end{enumerate}
  \end{enumerate}
\end{niproof}

\begin{lemma}[Falling Edge Equal Firable 2]
  \label{lem:fe-equal-firable-2}
  \fehyps{} then
  $\forall{}t\in{}T,id_t\in{}Comps(\Delta)~s.t.~\gamma(t)=id_t,$
  $\sigma'(id_t)("s\_firable")=\mathtt{true}\Rightarrow{}t\in{}Firable(s')$.
\end{lemma}

\begin{niproof}

  Given a $t\in{}T$ and $id_t\in{}Comps(\Delta)$
  s.t. $\gamma(t)=id_t$, and assuming that
  $\sigma'(id_t)("s\_firable")=\mathtt{true}$, let us
  show \fbox{$t\in{}Firable(s')$.}\\
  
  \exT{}

  By property of the $\mathtt{Inject}_\downarrow$ relation, the
  \hvhdl{} falling edge relation, the stabilize relation, \InCsCompT,
  and through the examination of the \texttt{firable} process defined
  in the transition design architecture, we can deduce:
  \begin{equation}
    \label{eq:fe-eq-sfa-true}
    \sigma'(id_t)("sfa")=\sigma(id_t)("se")~.~\sigma(id_t)("scc")~.~\mathtt{checktc}(\Delta(id_t),\sigma(id_t))=\mathtt{true}
  \end{equation}

  From \eqref{eq:fe-eq-sfa-true}, we can deduce:
  \begin{eqnarray}
    \label{eq:fe-eq-se-true}\sigma(id_t)("se")=\mathtt{true}\\
    \label{eq:fe-eq-scc-true}\sigma(id_t)("scc")=\mathtt{true}\\
    \label{eq:fe-eq-checktc-true}\mathtt{checktc}(\Delta(id_t),\sigma(id_t))=\mathtt{true}
  \end{eqnarray}

  Term $\mathtt{checktc}(\Delta(id_t),\sigma(id_t))$ as the same
  definition as in Lemma~\nameref{lem:fe-equal-firable-1}.

  By definition of $t\in{}Firable(s')$, there are three points to prove:
  \begin{enumerate}
  \item \fbox{$t\in{}Sens(s'.M)$}
  \item
    \fbox{\parbox{\lwidth}{$\forall{}c\in{}\mathcal{C},~\mathbb{C}(t,c)=1\Rightarrow{}s'.cond(c)=\mathtt{true}$
        and $\mathbb{C}(t,c)=-1\Rightarrow{}s'.cond(c)=\mathtt{false}$}}
  \item \fbox{$t\notin{}T_i\lor{}s'.I(t)\in{}I_s(t)$}
  \end{enumerate}

  Let us prove these three points:
  \begin{enumerate}
  \item \fbox{$t\in{}Sens(s'.M)$}:

    By definition of \dwSitpn, we have $s.M=s'.M$. Rewriting the goal
    with $s.M=s'.M$: \fbox{$t\in{}Sens(s.M)$.}

    By definition of \upSim, we have
    $\sigma(id_t)("se")=\mathtt{true}\Leftrightarrow{}t\in{}Sens(s.M)$.

    From $\sigma(id_t)("se")=\mathtt{true}$, we can deduce: \qedbox{$t\in{}Sens(s.M)$.}
    
  \item
    \fbox{\parbox{\lwidth}{$\forall{}c\in{}\mathcal{C},~\mathbb{C}(t,c)=1\Rightarrow{}s'.cond(c)=\mathtt{true}$
        and
        $\mathbb{C}(t,c)=-1\Rightarrow{}s'.cond(c)=\mathtt{false}$}}

    Given a $c\in\mathcal{C}$, there are two points to prove:
    \begin{enumerate}
    \item
      \fbox{$\mathbb{C}(t,c)=1\Rightarrow{}s'.cond(c)=\mathtt{true}$.}
    \item
      \fbox{$\mathbb{C}(t,c)=-1\Rightarrow{}s'.cond(c)=\mathtt{false}$.}
    \end{enumerate}

    Let us prove these two points:
    \begin{enumerate}
    \item Assuming that $\mathbb{C}(t,c)=1$, let us show
      \fbox{$s'.cond(c)=\mathtt{true}$.}

      By definition of \upSim, we have:
      \begin{equation}
        \sigma(id_t)("scc")=\prod\limits_{c'\in{}conds(t)}
        \begin{cases}
          E_c(\tau,c') & if~\mathbb{C}(t,c')=1 \\
          \mathtt{not}(E_c(\tau,c')) & if~\mathbb{C}(t,c')=-1 \\
        \end{cases}=\mathtt{true}
        \label{eq:fe-eq-scc-prod-true}
      \end{equation}
      where
      $conds(t)=\{c_i\in\mathcal{C}~\vert~\mathbb{C}(t,c_i)=1\lor\mathbb{C}(t,c_i)=-1\}$.

      From $\mathbb{C}(t,c)=1$, we can deduce $c\in{}conds(t)$. By
      definition of the product expression, we have:
      \begin{equation}
        \label{eq:fe-eq-scc-prod-true-decomp}
        E_c(\tau,c)~.~\prod\limits_{c'\in{}conds(t)\setminus\{c\}}
        \begin{cases}
          E_c(\tau,c') & if~\mathbb{C}(t,c')=1 \\
          \mathtt{not}(E_c(\tau,c')) & if~\mathbb{C}(t,c')=-1 \\
        \end{cases}=\mathtt{true}
      \end{equation}

      From \eqref{eq:fe-eq-scc-prod-true-decomp}, we can deduce that
      $E_c(\tau,c)=\mathtt{true}$.

      By definition of \dwSitpn{} (Rule~\ref{it:cond-env}), we have
      $s'.cond(c)=E_c(\tau,c)$.
      
      Rewriting the goal with $s'.cond(c)=E_c(\tau,c)$ and
      $E_c(\tau,c)=\mathtt{true}$: \qedbox{tautology.}
      
    \item Assuming that $\mathbb{C}(t,c)=-1$, let us show
      \fbox{$s'.cond(c)=\mathtt{false}$.}

      By definition of \upSim, we have:
      \begin{equation}
        \sigma(id_t)("scc")=\prod\limits_{c'\in{}conds(t)}
        \begin{cases}
          E_c(\tau,c') & if~\mathbb{C}(t,c')=1 \\
          \mathtt{not}(E_c(\tau,c')) & if~\mathbb{C}(t,c')=-1 \\
        \end{cases}=\mathtt{true}
        \label{eq:fe-eq-scc-prod-false}
      \end{equation}
      where
      $conds(t)=\{c'\in\mathcal{C}~\vert~\mathbb{C}(t,c')=1\lor\mathbb{C}(t,c')=-1\}$.

      From $\mathbb{C}(t,c)=-1$, we can deduce $c\in{}conds(t)$. By
      definition of the product expression, we have:
      \begin{equation}
        \label{eq:fe-eq-scc-prod-false-decomp}
        \mathtt{not}~E_c(\tau,c)~.~\prod\limits_{c'\in{}conds(t)\setminus\{c\}}
        \begin{cases}
          E_c(\tau,c') & if~\mathbb{C}(t,c')=1 \\
          \mathtt{not}(E_c(\tau,c')) & if~\mathbb{C}(t,c')=-1 \\
        \end{cases}=\mathtt{true}
      \end{equation}

      From \eqref{eq:fe-eq-scc-prod-false-decomp}, we can deduce that
      $E_c(\tau,c)=\mathtt{false}$.

      By definition of \dwSitpn{} (Rule~\ref{it:cond-env}), we have
      $s'.cond(c)=E_c(\tau,c)$.
      
      Rewriting the goal with $s'.cond(c)=E_c(\tau,c)$ and
      $E_c(\tau,c)=\mathtt{false}$: \qedbox{tautology.}
    \end{enumerate}
    
  \item \fbox{$t\notin{}T_i\lor{}s'.I(t)\in{}I_s(t)$}

    Reasoning on
    $\mathtt{checktc}(\Delta(id_t),\sigma(id_t))=\mathtt{true}$, there
    are 3 cases:
    \begin{enumerate}
    \item
      $\big(\mathtt{not}~\sigma(id_t)("srtc")~.~[\dots]\big)=\mathtt{true}$\footnote{See
        equation~\eqref{eq:checktc} for the full definition.}
    \item
      $\big(\sigma(id_t)("srtc")~.~\Delta(id_t)("tt")\neq\mathtt{NOT\_TEMP}~.~\sigma(id_t)("A")=1\big)=\mathtt{true}$
    \item
      $\big(\Delta(id_t)("tt")=\mathtt{NOT\_TEMP}\big)=\mathtt{true}$
    \end{enumerate}
    
    \begin{enumerate}
    \item \textbf{CASE}
      $\big(\mathtt{not}~\sigma(id_t)("srtc")~.~[\dots]\big)=\mathtt{true}$:

      Then, we can deduce
      $\mathtt{not}~\sigma(id_t)("srtc")=\mathtt{true}$ and
      $[\dots]=\mathtt{true}$. From
      $\mathtt{not}~\sigma(id_t)("srtc")=\mathtt{true}$, we can deduce
      $\sigma(id_t)("srtc")=\mathtt{false}$, and from
      $[\dots]=\mathtt{true}$, we have three other cases:

      \begin{enumerate}
      \item \textbf{CASE}
        $\big(\Delta(id_t)("tt")=\mathtt{TEMP\_A\_B}~.~(\sigma(id_t)("stc")\ge{}\sigma(id_t)("A")-1)~.~(\sigma(id_t)("stc")\le{}\sigma(id_t)("B")-1)\big)=\mathtt{true}$
      \item \textbf{CASE}
        $(\Delta(id_t)("tt")=\mathtt{TEMP\_A\_A}~.~(\sigma(id_t)("stc")={}\sigma(id_t)("A")-1))=\mathtt{true}$
      \item \textbf{CASE}
        $(\Delta(id_t)("tt")=\mathtt{TEMP\_A\_INF}~.~(\sigma(id_t)("stc")\ge{}\sigma(id_t)("A")-1))=\mathtt{true}$
      \end{enumerate}

      Let us prove the goal is these three contexts:
      \begin{enumerate}
      \item \textbf{CASE}
        $\big(\Delta(id_t)("tt")=\mathtt{TEMP\_A\_B}~.~(\sigma(id_t)("stc")\ge{}\sigma(id_t)("A")-1)~.~(\sigma(id_t)("stc")\le{}\sigma(id_t)("B")-1)\big)=\mathtt{true}$:

        Then, converting boolean equalities into intuitionistic
        predicates, we have:
        \begin{itemize}
        \item $\Delta(id_t)("tt")=\mathtt{TEMP\_A\_B}$
        \item $\sigma(id_t)("stc")\ge{}\sigma(id_t)("A")-1$
        \item $\sigma(id_t)("stc")\le{}\sigma(id_t)("B")-1$
        \end{itemize}

        By property of the elaboration relation, and
        $\Delta(id_t)("tt")=\mathtt{TEMP\_A\_B}$, there exist
        $a,b\in\mathbb{N}^{*}$ s.t. $I_s(t)=[a,b]$. Let us take such
        an $a$ and $b$. Then, let us show \fbox{$s'.I(t)\in{}I_s(t)$.}

        Rewriting the goal with $I_s(t)=[a,b]$:
        \fbox{$s'.I(t)\in[a,b]$.}
        
        By construction, ${<}\mathtt{time\_A\_value\Rightarrow}{}a{>}$
        and ${<}\mathtt{time\_B\_value\Rightarrow}{}b{>}$, and by
        property of stable $\sigma$, we have $\sigma(id_t)("A")=a$ and
        $\sigma(id_t)("B")=b$.

        Rewriting the goal with $\sigma(id_t)("A")=a$ and
        $\sigma(id_t)("B")=b$, and by definition of $\in$:
        \fbox{$\sigma(id_t)("A")\le{}s'.I(t)\le{}\sigma(id_t)("B")$.}

        Now, let us perform case analysis on
        $s.I(t)\le{}upper(I_s(t))$ or $s.I(t)>upper(I_s(t))$:
        \begin{itemize}
        \item \textbf{CASE} $s.I(t)\le{}upper(I_s(t))$:
          
          By definition of \upSim, we have $s.I(t)=\sigma(id_t)("stc")$.

          From $\sigma(id_t)("se")=\mathtt{true}$, we can deduce
          $t\in{}Sens(s.M)$, and from
          $\sigma(id_t)("srtc")=\mathtt{false}$, we can deduce
          $s.reset_t(t)=\mathtt{false}$. Then, by definition of
          \dwSitpn{} (Rule~\ref{it:inc-counters}), we have
          $s'.I(t)=s.I(t)+1$.

          $\Rightarrow{}$
          \fbox{$\sigma(id_t)("A")\le{}s.I(t)+1\le{}\sigma(id_t)("B")$}
          (by $s'.I(t)=s.I(t)+1$)

          $\Rightarrow{}$
          \fbox{$\sigma(id_t)("A")\le{}\sigma(id_t)("stc")+1\le{}\sigma(id_t)("B")$}
          (by $s.I(t)=\sigma(id_t)("stc")$)
          
          $\Rightarrow{}$
          \fbox{$\sigma(id_t)("A")-1\le{}\sigma(id_t)("stc")\le{}\sigma(id_t)("B")-1$}
          
          We assumed $\sigma(id_t)("stc")\ge{}\sigma(id_t)("A")-1$ and
          $\sigma(id_t)("stc")\le{}\sigma(id_t)("B")-1$, and thus we can deduce:
          \qedbox{$\sigma(id_t)("A")-1\le{}\sigma(id_t)("stc")\le{}\sigma(id_t)("B")-1$}
          
        \item \textbf{CASE} $s.I(t)>upper(I_s(t))$:

          By definition of \upSim, we have
          $\sigma(id_t)("stc")=upper(I_s(t))=b$.

          Then, from $\sigma(id_t)("stc")\le{}\sigma(id_t)("B")-1$,
          $\sigma(id_t)("stc")=upper(I_s(t))=b$ and
          $\sigma(id_t)("B")=b$, we can deduce the following
          contradiction:\\
          \qedbox{$\sigma(id_t)("B")\le{}\sigma(id_t)("B")-1$.}
        \end{itemize}
      \item
        $(\Delta(id_t)("tt")=\mathtt{TEMP\_A\_A}~.~(\sigma(id_t)("stc")={}\sigma(id_t)("A")-1))=\mathtt{true}$:

        Then, converting boolean equalities into intuitionistic
        predicates, we have:
        \begin{itemize}
        \item $\Delta(id_t)("tt")=\mathtt{TEMP\_A\_A}$
        \item $\sigma(id_t)("stc")=\sigma(id_t)("A")-1$
        \end{itemize}

        By property of the elaboration relation, and
        $\Delta(id_t)("tt")=\mathtt{TEMP\_A\_A}$, there exist
        $a\in\mathbb{N}^{*}$ s.t. $I_s(t)=[a,a]$. Let us take such an
        $a$. Then, let us show \fbox{$s'.I(t)\in{}I_s(t)$.}

        Rewriting the goal with $I_s(t)=[a,a]$:
        \fbox{$s'.I(t)\in[a,a]$.}
        
        By construction,
        ${<}\mathtt{time\_A\_value\Rightarrow}{}a{>}$, and by property
        of stable $\sigma$, we have $\sigma(id_t)("A")=a$.

        Rewriting the goal with $\sigma(id_t)("A")=a$, unfolding the
        definition of $\in$, and simplifying the goal:
        \fbox{$s'.I(t)=\sigma(id_t)("A")$.}

        Now, let us perform case analysis on
        $s.I(t)\le{}upper(I_s(t))$ or $s.I(t)>upper(I_s(t))$:
        \begin{itemize}
        \item \textbf{CASE} $s.I(t)\le{}upper(I_s(t))$:
          
          By definition of \upSim, we have $s.I(t)=\sigma(id_t)("stc")$.

          From $\sigma(id_t)("se")=\mathtt{true}$, we can deduce
          $t\in{}Sens(s.M)$, and from
          $\sigma(id_t)("srtc")=\mathtt{false}$, we can deduce
          $s.reset_t(t)=\mathtt{false}$. Then, by definition of
          \dwSitpn{} (Rule~\ref{it:inc-counters}), we have
          $s'.I(t)=s.I(t)+1$.

          $\Rightarrow{}$
          \fbox{$s.I(t)+1=\sigma(id_t)("A")$} (by $s'.I(t)=s.I(t)+1$)

          $\Rightarrow{}$
          \fbox{$\sigma(id_t)("stc")+1=\sigma(id_t)("A")$}
          (by $s.I(t)=\sigma(id_t)("stc")$)

          $\Rightarrow{}$
          \qedbox{$\sigma(id_t)("stc")=\sigma(id_t)("A")-1$} (assumption)
          
        \item \textbf{CASE} $s.I(t)>upper(I_s(t))$:

          By definition of \upSim, we have
          $\sigma(id_t)("stc")=upper(I_s(t))=a$.

          Then, from $\sigma(id_t)("stc")={}\sigma(id_t)("A")-1$,
          $\sigma(id_t)("stc")=upper(I_s(t))=a$,
          $\sigma(id_t)("A")=a$, and $a\in\mathbb{N}^{*}$, we can
          derive the following
          contradiction:\\
          \qedbox{$\sigma(id_t)("A")=\sigma(id_t)("A")-1$.}
        \end{itemize}
        
      \item
        $(\Delta(id_t)("tt")=\mathtt{TEMP\_A\_INF}~.~(\sigma(id_t)("stc")\ge\sigma(id_t)("A")-1))=\mathtt{true}$:

        Then, converting boolean equalities into intuitionistic
        predicates, we have:
        \begin{itemize}
        \item $\Delta(id_t)("tt")=\mathtt{TEMP\_A\_INF}$
        \item $\sigma(id_t)("stc")\ge\sigma(id_t)("A")-1$
        \end{itemize}

        By property of the elaboration relation, and
        $\Delta(id_t)("tt")=\mathtt{TEMP\_A\_INF}$, there exist
        $a\in\mathbb{N}^{*}$ s.t. $I_s(t)=[a,\infty]$. Let us take
        such an $a$. Then, let us show \fbox{$s'.I(t)\in{}I_s(t)$.}

        Rewriting the goal with $I_s(t)=[a,\infty]$:
        \fbox{$s'.I(t)\in[a,\infty]$.}
        
        By construction,
        ${<}\mathtt{time\_A\_value\Rightarrow}{}a{>}$, and by property
        of stable $\sigma$, we have $\sigma(id_t)("A")=a$.

        Rewriting the goal with $\sigma(id_t)("A")=a$, unfolding the
        definition of $\in$, and simplifying the goal:
        \fbox{$\sigma(id_t)("A")\le{}s'.I(t)$.}

        Now, let us perform case analysis on
        $s.I(t)\le{}lower(I_s(t))$ or $s.I(t)>lower(I_s(t))$:
        \begin{itemize}
        \item \textbf{CASE} $s.I(t)\le{}lower(I_s(t))$:
          
          By definition of \upSim, we have $s.I(t)=\sigma(id_t)("stc")$.

          From $\sigma(id_t)("se")=\mathtt{true}$, we can deduce
          $t\in{}Sens(s.M)$, and from
          $\sigma(id_t)("srtc")=\mathtt{false}$, we can deduce
          $s.reset_t(t)=\mathtt{false}$. Then, by definition of
          \dwSitpn{} (Rule~\ref{it:inc-counters}), we have
          $s'.I(t)=s.I(t)+1$.

          $\Rightarrow{}$
          \fbox{$\sigma(id_t)("A")\le{}s.I(t)+1$} (by $s'.I(t)=s.I(t)+1$)

          $\Rightarrow{}$
          \fbox{$\sigma(id_t)("A")\le\sigma(id_t)("stc")+1$} (by
          $s.I(t)=\sigma(id_t)("stc")$)

          $\Rightarrow{}$
          \qedbox{$\sigma(id_t)("A")-1\le\sigma(id_t)("stc")$}
          (assumption)
          
        \item \textbf{CASE} $s.I(t)>lower(I_s(t))$:

          By definition of \upSim, we have
          $\sigma(id_t)("stc")=lower(I_s(t))=a$.

          From $\sigma(id_t)("se")=\mathtt{true}$, we can deduce
          $t\in{}Sens(s.M)$, and from
          $\sigma(id_t)("srtc")=\mathtt{false}$, we can deduce
          $s.reset_t(t)=\mathtt{false}$. Then, by definition of
          \dwSitpn{} (Rule~\ref{it:inc-counters}), we have
          $s'.I(t)=s.I(t)+1$.

          $\Rightarrow{}$ \fbox{$\sigma(id_t)("A")\le{}s.I(t)+1$} (by
          $s'.I(t)=s.I(t)+1$)

          $\Rightarrow{}$ \fbox{$a\le{}s.I(t)+1$} (by
          $\sigma(id_t)("A")=a$)
          
          $\Rightarrow{}$ \fbox{$a<s.I(t)$}

          $\Rightarrow{}$ \qedbox{$lower(I_s(t))<s.I(t)$} (assumption)
        \end{itemize}
      \end{enumerate}

    \item
      $\big(\sigma(id_t)("srtc")~.~\Delta(id_t)("tt")\neq\mathtt{NOT\_TEMP}~.~\sigma(id_t)("A")=1\big)=\mathtt{true}$

      Then, converting boolean equalities into intuitionistic predicates, we have:
      \begin{itemize}
      \item $\sigma(id_t)("srtc")=\mathtt{true}$
      \item $\Delta(id_t)("tt")\neq\mathtt{NOT\_TEMP}$
      \item $\sigma(id_t)("A")=1$
      \end{itemize}

      By property of the elaboration relation, and
      $\Delta(id_t)("tt")\neq\mathtt{NOT\_TEMP}$, there exist an
      $a\in\mathbb{N}^{*}$ and a $ni\in\mathbb{N}^{*}\sqcup\{\infty\}$
      s.t. $I_s(t)=[a,ni]$. Let us take such an $a$ and $ni$.

      By construction,
      ${<}\mathtt{time\_A\_value\Rightarrow}{}a{>}\in{}ipm_t$, and by
      property of stable $\sigma$, we have
      $\sigma(id_t)("A")=a$. Thus, we can deduce $a=1$ and
      $I_s(t)=[1,ni]$.

      By definition of \upSim, from
      $\sigma(id_t)("se")=\mathtt{true}$, we can deduce
      $t\in{}Sens(s.M)$, and from
      $\sigma(id_t)("srtc")=\mathtt{true}$, we can deduce
      $s.reset_t(t)=\mathtt{true}$.

      By definition of \dwSitpn{} (Rule~\ref{it:reset-counters}),
      $t\in{}Sens(s.M)$ and $s.reset_t(t)=\mathtt{true}$, we have
      $s'.I(t)=1$.

      Now, let us show \fbox{$s'.I(t)\in{}I_s(t)$}.

      Rewriting the goal with $s'.I(t)=1$ and $I_s(t)=[1,ni]$:
      \qedbox{$1\in[1,ni]$.}
    \item
      $\big(\Delta(id_t)("tt")=\mathtt{NOT\_TEMP}\big)=\mathtt{true}$

      Let us show \fbox{$t\notin{}T_i$.}
      
      By property of the elaboration relation and
      $\Delta(id_t)("tt")=\mathtt{NOT\_TEMP}$, we have
      \qedbox{$t\notin{}T_i$.}
    \end{enumerate}
  \end{enumerate}
\end{niproof}

\begin{lemma}[Falling edge equal not firable]
  \label{lem:fe-equal-not-firable}
  \fehyps{} then
  $\forall{}t\in{}T,id_t\in{}Comps(\Delta)~s.t.~\gamma(t)=id_t,$
  $t\notin{}Firable(s')\Leftrightarrow\sigma'(id_t)("s\_firable")=\mathtt{false}$.
\end{lemma}

\begin{proof}
  Proving the above lemma is trivial by appealing to
  Lemma~\ref{lem:fe-equal-firable} and by reasoning on
  contrapositives.
\end{proof}

\subsection{Falling edge and fired transitions}
\label{sec:fe-fired}

%%%%%%%%%%%%%%%%%%%%%%%%%%%%%%%%%%%%%%%%%%%%%%%%%%%%%%%%
%%%%%%%%%% FALLING EDGE EQUAL FIRED SET LEMMA %%%%%%%%%%
%%%%%%%%%%%%%%%%%%%%%%%%%%%%%%%%%%%%%%%%%%%%%%%%%%%%%%%%

\begin{lemma}[Falling Edge Equal Fired Set]
  \label{lem:fe-equal-fset}
  \fehyps{} then
  $\forall{}t\in{}T,~id_t\in{}Comps(\Delta)~s.t.~\gamma(t)=id_t$,
  $\forall{}Fset\subseteq{}T,~s.t.~IsFiredSet(s',Fset),~t\in{}Fset\Leftrightarrow\sigma'(id_t)(\texttt{fired})=\mathtt{true}$.
\end{lemma}

\begin{niproof}
  Given a $t\in{}T$, and $id_t\in{}Comps(\Delta)$, and a
  $Fset\subseteq{}T$ s.t. $IsFiredSet(s',Fset)$, let us show
  \fbox{$t\in{}Fset\Leftrightarrow\sigma'(id_t)(\texttt{fired})=\mathtt{true}$.}\\

  By definition of $IsFiredSet(s',Fset)$, we have
  $IsFiredSetAux(s',T,\emptyset,Fset)$.

  Then, we can appeal to Lemma~\ref{lem:fe-equal-fset-aux} to solve
  the goal, but first we must prove the following \emph{extra
    hypothesis} (i.e, one of the premise of
  Lemma~\nameref{lem:fe-equal-fset-aux}):

  \fbox{\parbox{\lwidth}{$\forall{}t'\in{}T,id_{t'}\in{}Comps(\Delta)$
      s.t. $\gamma(t')=id_{t'},$\\
      $(t'\in{}\emptyset\Rightarrow\sigma'(id_{t'})(\texttt{fired})=\mathtt{true})$
      $\land$
      $(\sigma'(id_{t'})(\texttt{fired})=\mathtt{true}\Rightarrow{}t'\in{}\emptyset~\lor~{}t'\in{}T)$.}}\\

  Given a $t'\in{}T$ and an $id_{t'}\in{}Comps(\Delta)$
  s.t. $\gamma(t')=id_{t'}$, there are two points to prove:
  \begin{enumerate}
  \item
    \fbox{$t'\in{}\emptyset\Rightarrow\sigma'(id_{t'})(\texttt{fired})=\mathtt{true}$}
  \item
    \fbox{$\sigma'(id_{t'})(\texttt{fired})=\mathtt{true}\Rightarrow{}t'\in{}\emptyset~\lor~{}t'\in{}T$}
  \end{enumerate}

  Let us show these two points:
  \begin{enumerate}
  \item Assuming $t'\in{}\emptyset$, let us show
    \fbox{$\sigma'(id_{t'})(\texttt{fired})=\mathtt{true}$.}

    \qedbox{$t'\in{}\emptyset$ is a contradiction.}
  \item Assuming $\sigma'(id_{t'})(\texttt{fired})=\mathtt{true}$, let us
    show \fbox{$t'\in{}\emptyset~\lor~{}t'\in{}T$.}

    By definition, \qedbox{$t'\in{}T$.}
  \end{enumerate}
  
\end{niproof}

%%%%%%%%%%%%%%%%%%%%%%%%%%%%%%%%%%%%%%%%%%%%%%%%%%%%%%%%%%%%
%%%%%%%%%% FALLING EDGE EQUAL FIRED SET AUX LEMMA %%%%%%%%%%
%%%%%%%%%%%%%%%%%%%%%%%%%%%%%%%%%%%%%%%%%%%%%%%%%%%%%%%%%%%%

\begin{lemma}[Falling edge equal fired set aux]
  \label{lem:fe-equal-fset-aux}
  \fehyps{} then
  $\forall{}t\in{}T,id_t\in{}Comps(\Delta)~s.t.~\gamma(t)=id_t$,
  $\forall{}F\subseteq{}T,~T_s\subseteq{}T,~Fset\subseteq{}T,$
  assume that:
  \begin{itemize}
  \item $IsFiredSetAux(s',T_s,F,Fset)$
  \item EH (Extra. Hypothesis):\\
    $\forall{}t'\in{}T,~id_{t'}\in{}Comps(\Delta)$
    s.t. $\gamma(t')=id_{t'}$,\\
    $(t'\in{}F\Rightarrow\sigma'(id_{t'})(\texttt{fired})=\mathtt{true})$
    $\land$
    $(\sigma'(id_{t'})(\texttt{fired})=\mathtt{true}\Rightarrow{}t'\in{}F~\lor~{}t'\in{}T_s)$.
  \end{itemize}
  then
  $t\in{}Fset\Leftrightarrow\sigma'(id_t)(\texttt{fired})=\mathtt{true}$.
\end{lemma}

\begin{niproof}
  Given a $t\in{}T$, an $id_t\in{}Comps(\Delta)$, a
  $T_s,F,Fset\subseteq{}T$, and assuming\\
  $IsFiredSetAux(s',T_s,F,Fset)$, let us show
  \begin{frameb}
    $\big(\forall{}t'\in{}T,~id_{t'}\in{}Comps(\Delta)$
    s.t. $\gamma(t')=id_{t'}$,\\
    $(t'\in{}F\Rightarrow\sigma'(id_{t'})(\texttt{fired})=\mathtt{true})$
    $\land$
    $(\sigma'(id_{t'})(\texttt{fired})=\mathtt{true}\Rightarrow{}t'\in{}F~\lor~{}t'\in{}T_s)\big)\Rightarrow$
    $t\in{}Fset\Leftrightarrow\sigma'(id_t)(\texttt{fired})=\mathtt{true}$.
  \end{frameb}

  Let us use rule induction on $IsFiredSetAux(s',T_s,F,Fset)$. Let us
  define the property $P$ taken into account in the induction scheme
  as follows

  \begin{center}
    \begin{tabular}{c}
      $P(s',T_s,F,Fset)$ \\
      $\equiv$ \\
      $(t'\in{}F\Rightarrow\sigma'(id_{t'})(\texttt{fired})=\mathtt{true})$
      $\land$
      $\big(\sigma'(id_{t'})(\texttt{fired})=\mathtt{true}\Rightarrow{}t'\in{}F~\lor~{}t'\in{}T_s)\big)\Rightarrow$ \\
      $t\in{}Fset\Leftrightarrow\sigma'(id_t)(\texttt{fired})=\mathtt{true}$ \\
    \end{tabular}
  \end{center}

  \begin{itemize}
  \item \textbf{CASE} \textsc{FSetEmp}: we must show $P(s',\emptyset,F,F)$, i.e.
    \begin{frameb}
      $\big(\forall{}t'\in{}T,~id_{t'}\in{}Comps(\Delta)$
      s.t. $\gamma(t')=id_{t'}$,\\
      $(t'\in{}F\Rightarrow\sigma'(id_{t'})(\texttt{fired})=\mathtt{true})$
      $\land$
      $(\sigma'(id_{t'})(\texttt{fired})=\mathtt{true}\Rightarrow{}t'\in{}F~\lor~{}t'\in{}\emptyset)\big)\Rightarrow$
      $t\in{}F\Leftrightarrow\sigma'(id_t)(\texttt{fired})=\mathtt{true}$.
    \end{frameb}
    
    Assuming
    \begin{center}
      \begin{tabular}{l}
        $\forall{}t'\in{}T,~id_{t'}\in{}Comps(\Delta)$
        s.t. $\gamma(t')=id_{t'}$,\\
        $(t'\in{}F\Rightarrow\sigma'(id_{t'})(\texttt{fired})=\mathtt{true})$
        $\land$
        $(\sigma'(id_{t'})(\texttt{fired})=\mathtt{true}\Rightarrow{}t'\in{}F~\lor~{}t'\in{}\emptyset)$ \\
      \end{tabular}
    \end{center}
    we can easily show
    \qedbox{$t\in{}F\Leftrightarrow\sigma'(id_t)(\texttt{fired})=\mathtt{true}$.}
    
  \item \textbf{CASE} \textsc{FSetFired}:

    Assuming
    \begin{itemize}
    \item $t\in{}Firable(s')$
    \item $t\in{}Sens(s'.M-\sum\limits_{t_i\in{}Pr(t,F)}pre(t_i))$
    \item $IsFiredSetAux(s',T_s,F\cup\{t\},Fset)$
    \item $\nexists{}t'\in{}T_s~s.t.~t'\succ{}t$
    \item $Pr(t,F)=\{t'~\vert~t'\succ{}t\land{}t'\in{}F\}$
    \end{itemize}
    and the induction hypothesis (i.e. $P(s',T_s,F\cup\{t\},Fset)$)
    \begin{ih}
      \begin{tabular}{l}
        $\big(\forall{}t'\in{}T,~id_{t'}\in{}Comps(\Delta)$
        s.t. $\gamma(t')=id_{t'}$,\\
        $(t'\in{}F\cup\{t\}\Rightarrow\sigma'(id_{t'})(\texttt{fired})=\mathtt{true})$\\
        $\land$
        $(\sigma'(id_{t'})(\texttt{fired})=\mathtt{true}\Rightarrow{}t'\in{}F\cup\{t\}~\lor~{}t'\in{}T_s)\big)\Rightarrow$ \\
        $t\in{}Fset\Leftrightarrow\sigma'(id_t)(\texttt{fired})=\mathtt{true}$
      \end{tabular}
    \end{ih}

    we must show

    \begin{frameb}
      \begin{tabular}{l}
        $\big(\forall{}t'\in{}T,~id_{t'}\in{}Comps(\Delta)$
        s.t. $\gamma(t')=id_{t'}$,\\
        $(t'\in{}F\Rightarrow\sigma'(id_{t'})(\texttt{fired})=\mathtt{true})$\\
        $\land$
        $(\sigma'(id_{t'})(\texttt{fired})=\mathtt{true}\Rightarrow{}t'\in{}F~\lor~{}t'\in{}T_s\cup\{t\})\big)\Rightarrow$ \\
        $t\in{}Fset\Leftrightarrow\sigma'(id_t)(\texttt{fired})=\mathtt{true}$
      \end{tabular}
    \end{frameb}

    Assuming the following hypothesis that we will call EH (for Extra
    Hypothesis)

    \begin{center}
      \begin{tabular}{l}
        $\forall{}t'\in{}T,~id_{t'}\in{}Comps(\Delta)$
        s.t. $\gamma(t')=id_{t'}$,\\
        $(t'\in{}F\Rightarrow\sigma'(id_{t'})(\texttt{fired})=\mathtt{true})$
        $\land$
        $(\sigma'(id_{t'})(\texttt{fired})=\mathtt{true}\Rightarrow{}t'\in{}F~\lor~{}t'\in{}T_s\cup\{t\})$ \\
      \end{tabular}
    \end{center}

    we must show
    \begin{frameb}
      $t\in{}Fset\Leftrightarrow\sigma'(id_t)(\texttt{fired})=\mathtt{true}$
    \end{frameb}
    
    Appealing to the induction hypothesis, to prove the current goal,
    it is sufficient to prove that

    \begin{frameb}
      \begin{tabular}{l}
        $\forall{}t'\in{}T,~id_{t'}\in{}Comps(\Delta)$
        s.t. $\gamma(t')=id_{t'}$,\\
        $(t'\in{}F\cup\{t\}\Rightarrow\sigma'(id_{t'})(\texttt{fired})=\mathtt{true})$\\
        $\land$
        $(\sigma'(id_{t'})(\texttt{fired})=\mathtt{true}\Rightarrow{}t'\in{}F\cup\{t\}~\lor~{}t'\in{}T_s)$ \\
      \end{tabular}
    \end{frameb}

    Given a $t'\in{}T$, an $id_{t'}\in{}Comps(\Delta)$ s.t. $\gamma(t')=id_{t'}$, we must show that
    \begin{frameb}
      \begin{tabular}{l}
        $(t'\in{}F\cup\{t\}\Rightarrow\sigma'(id_{t'})(\texttt{fired})=\mathtt{true})$\\
        $\land$
        $(\sigma'(id_{t'})(\texttt{fired})=\mathtt{true}\Rightarrow{}t'\in{}F\cup\{t\}~\lor~{}t'\in{}T_s)$ \\
      \end{tabular}
    \end{frameb}

    There are two points to prove
    \begin{enumerate}
    \item Assuming $t'\in{}F\cup\{t\}$, then
      $\sigma'(id_{t'})(\texttt{fired})=\mathtt{true}$
    \item Assuming $\sigma'(id_{t'})(\texttt{fired})=\mathtt{true}$,
      then $t'\in{}F\cup\{t\}~\lor~{}t'\in{}T_s$
    \end{enumerate}

    \begin{enumerate}
    \item Assuming $t'\in{}F\cup\{t\}$, let us show \fbox{$\sigma'(id_{t'})(\texttt{fired})=\mathtt{true}$}.
      Let us perform case analysis on $t'\in{}F\cup\{t\}$; there are 2 cases:
      \begin{itemize}
      \item \textbf{CASE} $t'\in{}F$: Appealing to EH, the goal is trivially proved.
      \item \textbf{CASE} $t'=t$: Then, $id_t=id_{t'}$, and we must
        show \fbox{$\sigma'(id_t)(\texttt{fired})=\mathtt{true}$}.

        By definition of $id_{t}$, there exist a $g_{t}$, $i_{t}$,
        $o_{t}$ s.t.  $\mathtt{comp}(id_{t},\texttt{transition},$
        $g_{t},$ $i_{t}$, $o_{t})\in{}d.cs$.

        By property of the stabilize relation and
        $\mathtt{comp}(id_{t},\texttt{transition},$ $g_{t},$ $i_{t}$,
        $o_{t})\in{}d.cs$, and through the examination of the
        \texttt{fired_evaluation} process defined in the
        \texttt{transition} design architecture:
        \begin{equation*}
          \sigma(id_{t})(\texttt{fired})=\sigma(id_{t})(\texttt{sfa})~.~\sigma(id_{t})(\texttt{spc})
        \end{equation*}

        Rewriting the goal with the above equation:
        \fbox{$\sigma(id_{t})(\texttt{sfa})~.~\sigma(id_{t})(\texttt{spc})=\mathtt{true}$.}
        
        Then, there are two points to prove: 
        \begin{enumerate}
        \item \fbox{$\sigma(id_{t})(\texttt{sfa})=\mathtt{true}$.}
          
          Appealing to Lemma~\ref{lem:fe-equal-firable}, and since
          $t\in{}Firable(s')$, we can deduce
          \qedbox{$\sigma(id_{t})(\texttt{sfa})=\mathtt{true}$.}
        \item \fbox{$\sigma(id_{t})(\texttt{spc})=\mathtt{true}$.}
          
          Appealing to Lemma~\ref{lem:stab-compute-pcomb}, and since
          $t\in{}Sens(s'M-\sum\limits_{t_i\in{}Pr(t,F)}pre(t_i))$, we
          can deduce
          \qedbox{$\sigma(id_{t})(\texttt{spc})=\mathtt{true}$.}
        \end{enumerate}
      \end{itemize}
      
    \item Assuming $\sigma'(id_{t'})(\texttt{fired})=\mathtt{true}$,
      let us show
      \fbox{$t'\in{}F\cup\{t\}~\lor~{}t'\in{}T_s$}. Appealing to EH,
      we can deduce that $t'\in{}F~\lor~{}t'\in{}T_s\cup\{t\}$. Then,
      the goal is trivially shown.
    \end{enumerate}
    
  \item \textbf{CASE} \textsc{FSetNotFirable}:
    Assuming
    \begin{itemize}
    \item $t\notin{}Firable(s')$
    \item $IsFiredSetAux(s',T_s,F,Fset)$
    \item $\nexists{}t'\in{}T_s~s.t.~t'\succ{}t$
    \end{itemize}
    and the induction hypothesis (i.e. $P(s',T_s,F,Fset)$)
    \begin{ih}
      \begin{tabular}{l}
        $\big(\forall{}t'\in{}T,~id_{t'}\in{}Comps(\Delta)$
        s.t. $\gamma(t')=id_{t'}$,\\
        $(t'\in{}F\Rightarrow\sigma'(id_{t'})(\texttt{fired})=\mathtt{true})$\\
        $\land$
        $(\sigma'(id_{t'})(\texttt{fired})=\mathtt{true}\Rightarrow{}t'\in{}F~\lor~{}t'\in{}T_s)\big)\Rightarrow$ \\
        $t\in{}Fset\Leftrightarrow\sigma'(id_t)(\texttt{fired})=\mathtt{true}$
      \end{tabular}
    \end{ih}

    we must show

    \begin{frameb}
      \begin{tabular}{l}
        $\big(\forall{}t'\in{}T,~id_{t'}\in{}Comps(\Delta)$
        s.t. $\gamma(t')=id_{t'}$,\\
        $(t'\in{}F\Rightarrow\sigma'(id_{t'})(\texttt{fired})=\mathtt{true})$\\
        $\land$
        $(\sigma'(id_{t'})(\texttt{fired})=\mathtt{true}\Rightarrow{}t'\in{}F~\lor~{}t'\in{}T_s\cup\{t\})\big)\Rightarrow$ \\
        $t\in{}Fset\Leftrightarrow\sigma'(id_t)(\texttt{fired})=\mathtt{true}$
      \end{tabular}
    \end{frameb}

    Assuming the following hypothesis that we will call EH (for Extra
    Hypothesis)

    \begin{center}
      \begin{tabular}{l}
        $\forall{}t'\in{}T,~id_{t'}\in{}Comps(\Delta)$
        s.t. $\gamma(t')=id_{t'}$,\\
        $(t'\in{}F\Rightarrow\sigma'(id_{t'})(\texttt{fired})=\mathtt{true})$
        $\land$
        $(\sigma'(id_{t'})(\texttt{fired})=\mathtt{true}\Rightarrow{}t'\in{}F~\lor~{}t'\in{}T_s\cup\{t\})$ \\
      \end{tabular}
    \end{center}

    we must show
    \begin{frameb}
      $t\in{}Fset\Leftrightarrow\sigma'(id_t)(\texttt{fired})=\mathtt{true}$
    \end{frameb}
    
    Appealing to the induction hypothesis, to prove the current goal,
    it is sufficient to prove that

    \begin{frameb}
      \begin{tabular}{l}
        $\forall{}t'\in{}T,~id_{t'}\in{}Comps(\Delta)$
        s.t. $\gamma(t')=id_{t'}$,\\
        $(t'\in{}F\Rightarrow\sigma'(id_{t'})(\texttt{fired})=\mathtt{true})$
        $\land$
        $(\sigma'(id_{t'})(\texttt{fired})=\mathtt{true}\Rightarrow{}t'\in{}F~\lor~{}t'\in{}T_s)$ \\
      \end{tabular}
    \end{frameb}

    Given a $t'\in{}T$, an $id_{t'}\in{}Comps(\Delta)$ s.t. $\gamma(t')=id_{t'}$, we must show that
    \begin{frameb}
      \begin{tabular}{l}
        $(t'\in{}F\Rightarrow\sigma'(id_{t'})(\texttt{fired})=\mathtt{true})$
        $\land$
        $(\sigma'(id_{t'})(\texttt{fired})=\mathtt{true}\Rightarrow{}t'\in{}F~\lor~{}t'\in{}T_s)$ \\
      \end{tabular}
    \end{frameb}

    There are two points to prove
    \begin{enumerate}
    \item Assuming $t'\in{}F$, then
      $\sigma'(id_{t'})(\texttt{fired})=\mathtt{true}$
    \item Assuming $\sigma'(id_{t'})(\texttt{fired})=\mathtt{true}$,
      then $t'\in{}F~\lor~{}t'\in{}T_s$
    \end{enumerate}

    \begin{enumerate}
    \item Assuming $t'\in{}F$, let us show
      \fbox{$\sigma'(id_{t'})(\texttt{fired})=\mathtt{true}$}.

      Appealing to EH, the goal is trivially shown.
      
    \item Assuming $\sigma'(id_{t'})(\texttt{fired})=\mathtt{true}$,
      let us show \fbox{$t'\in{}F~\lor~{}t'\in{}T_s$}.

      Appealing to EH, we can deduce $t'\in{}F~\lor~{}t'\in{}T_s\cup\{t\}$.
      Let us perform case analysis on $t'\in{}F~\lor~{}t'\in{}T_s\cup\{t\}$; there are 2 cases:
      \begin{itemize}
      \item \textbf{CASE} $t'\in{}F$: trivially shown, as it is an assumption.
      \item \textbf{CASE} $t'\in{}T_s\cup\{t\}$: In the case where
        $t'\in{}T_s$, the goal is trivially shown. In the case where
        $t'=t$, we can prove a contradiction based on
        $t\notin{}Firable(s')$ and
        $\sigma'(id_{t'})(\texttt{fired})=\mathtt{true}$.

        Since $t=t'$, then $id_t=id_{t'}$, and we know that
        $\sigma'(id_t)(\texttt{fired})=\mathtt{true}$.

        By definition of $id_{t}$, there exist a $g_{t}$, $i_{t}$,
        $o_{t}$ s.t.  $\mathtt{comp}(id_{t},\texttt{transition},$
        $g_{t},$ $i_{t}$, $o_{t})\in{}d.cs$.

        By property of the stabilize relation and
        $\mathtt{comp}(id_{t},\texttt{transition},$ $g_{t},$ $i_{t}$,
        $o_{t})\in{}d.cs$, and through the examination of the
        \texttt{fired_evaluation} process defined in the
        \texttt{transition} design architecture, we can deduce
        \begin{equation*}
          \sigma(id_{t})(\texttt{fired})=\sigma(id_{t})(\texttt{sfa})~.~\sigma(id_{t})(\texttt{spc})=\mathtt{true}
        \end{equation*}
        
        Thus, we have
        \begin{equation*}
          \sigma(id_{t})(\texttt{sfa})=\mathtt{true}
        \end{equation*}
        and, appealing to Lemma~\ref{lem:fe-equal-firable}, we can
        deduce $t\in{}Firable(s')$, which directly contradicts
        \qedbox{ $t\notin{}Firable(s')$.}
      \end{itemize}
    \end{enumerate}
    
  \item \textbf{CASE} \textsc{FSetNotSens}:
    Assuming
    \begin{itemize}
    \item $t\notin{}Sens(s'.M-\sum\limits_{t_i\in{}Pr(t,F)}pre(t_i))$
    \item $IsFiredSetAux(s',T_s,F,Fset)$
    \item $\nexists{}t'\in{}T_s~s.t.~t'\succ{}t$
    \item $Pr(t,F)=\{t'~\vert~t'\succ{}t\land{}t'\in{}F\}$
    \end{itemize}
    and the induction hypothesis (i.e. $P(s',T_s,F,Fset)$)
    \begin{ih}
      \begin{tabular}{l}
        $\big(\forall{}t'\in{}T,~id_{t'}\in{}Comps(\Delta)$
        s.t. $\gamma(t')=id_{t'}$,\\
        $(t'\in{}F\Rightarrow\sigma'(id_{t'})(\texttt{fired})=\mathtt{true})$\\
        $\land$
        $(\sigma'(id_{t'})(\texttt{fired})=\mathtt{true}\Rightarrow{}t'\in{}F~\lor~{}t'\in{}T_s)\big)\Rightarrow$ \\
        $t\in{}Fset\Leftrightarrow\sigma'(id_t)(\texttt{fired})=\mathtt{true}$
      \end{tabular}
    \end{ih}

    we must show

    \begin{frameb}
      \begin{tabular}{l}
        $\big(\forall{}t'\in{}T,~id_{t'}\in{}Comps(\Delta)$
        s.t. $\gamma(t')=id_{t'}$,\\
        $(t'\in{}F\Rightarrow\sigma'(id_{t'})(\texttt{fired})=\mathtt{true})$\\
        $\land$
        $(\sigma'(id_{t'})(\texttt{fired})=\mathtt{true}\Rightarrow{}t'\in{}F~\lor~{}t'\in{}T_s\cup\{t\})\big)\Rightarrow$ \\
        $t\in{}Fset\Leftrightarrow\sigma'(id_t)(\texttt{fired})=\mathtt{true}$
      \end{tabular}
    \end{frameb}

    Assuming the following hypothesis, which we will call EH (for
    Extra Hypothesis)

    \begin{center}
      \begin{tabular}{l}
        $\forall{}t'\in{}T,~id_{t'}\in{}Comps(\Delta)$
        s.t. $\gamma(t')=id_{t'}$,\\
        $(t'\in{}F\Rightarrow\sigma'(id_{t'})(\texttt{fired})=\mathtt{true})$
        $\land$
        $(\sigma'(id_{t'})(\texttt{fired})=\mathtt{true}\Rightarrow{}t'\in{}F~\lor~{}t'\in{}T_s\cup\{t\})$ \\
      \end{tabular}
    \end{center}

    we must show
    \begin{frameb}
      $t\in{}Fset\Leftrightarrow\sigma'(id_t)(\texttt{fired})=\mathtt{true}$
    \end{frameb}
    
    Appealing to the induction hypothesis, to prove the current goal,
    it is sufficient to prove that

    \begin{frameb}
      \begin{tabular}{l}
        $\forall{}t'\in{}T,~id_{t'}\in{}Comps(\Delta)$
        s.t. $\gamma(t')=id_{t'}$,\\
        $(t'\in{}F\Rightarrow\sigma'(id_{t'})(\texttt{fired})=\mathtt{true})$
        $\land$
        $(\sigma'(id_{t'})(\texttt{fired})=\mathtt{true}\Rightarrow{}t'\in{}F~\lor~{}t'\in{}T_s)$ \\
      \end{tabular}
    \end{frameb}

    Given a $t'\in{}T$, an $id_{t'}\in{}Comps(\Delta)$ s.t. $\gamma(t')=id_{t'}$, we must show that
    \begin{frameb}
      \begin{tabular}{l}
        $(t'\in{}F\Rightarrow\sigma'(id_{t'})(\texttt{fired})=\mathtt{true})$
        $\land$
        $(\sigma'(id_{t'})(\texttt{fired})=\mathtt{true}\Rightarrow{}t'\in{}F~\lor~{}t'\in{}T_s)$ \\
      \end{tabular}
    \end{frameb}

    There are two points to prove
    \begin{enumerate}
    \item Assuming $t'\in{}F$, then
      $\sigma'(id_{t'})(\texttt{fired})=\mathtt{true}$
    \item Assuming $\sigma'(id_{t'})(\texttt{fired})=\mathtt{true}$,
      then $t'\in{}F~\lor~{}t'\in{}T_s$
    \end{enumerate}

    \begin{enumerate}
    \item Assuming $t'\in{}F$, let us show
      \fbox{$\sigma'(id_{t'})(\texttt{fired})=\mathtt{true}$}.

      Appealing to EH, the goal is trivially shown.
      
    \item Assuming $\sigma'(id_{t'})(\texttt{fired})=\mathtt{true}$,
      let us show \fbox{$t'\in{}F~\lor~{}t'\in{}T_s$}.

      Appealing to EH, we can deduce $t'\in{}F~\lor~{}t'\in{}T_s\cup\{t\}$.
      Let us perform case analysis on $t'\in{}F~\lor~{}t'\in{}T_s\cup\{t\}$; there are 2 cases:
      \begin{itemize}
      \item \textbf{CASE} $t'\in{}F$: trivially shown, as it is an assumption.
      \item \textbf{CASE} $t'\in{}T_s\cup\{t\}$: In the case where
        $t'\in{}T_s$, the goal is trivially shown. In the case where
        $t'=t$, we can prove a contradiction based on
        $t\notin{}Sens(s'.M-\sum\limits_{t_i\in{}Pr(t,F)}pre(t_i))$
        and $\sigma'(id_{t'})(\texttt{fired})=\mathtt{true}$.

        Since $t=t'$, then $id_t=id_{t'}$, and we know that
        $\sigma'(id_t)(\texttt{fired})=\mathtt{true}$.

        By definition of $id_{t}$, there exist a $g_{t}$, $i_{t}$,
        $o_{t}$ s.t.  $\mathtt{comp}(id_{t},\texttt{transition},$
        $g_{t},$ $i_{t}$, $o_{t})\in{}d.cs$.

        By property of the stabilize relation and
        $\mathtt{comp}(id_{t},\texttt{transition},$ $g_{t},$ $i_{t}$,
        $o_{t})\in{}d.cs$, and through the examination of the
        \texttt{fired_evaluation} process defined in the
        \texttt{transition} design architecture, we can deduce
        \begin{equation*}
          \sigma(id_{t})(\texttt{fired})=\sigma(id_{t})(\texttt{sfa})~.~\sigma(id_{t})(\texttt{spc})=\mathtt{true}
        \end{equation*}
        
        Thus, we have
        \begin{equation*}
          \sigma(id_{t})(\texttt{spc})=\mathtt{true}
        \end{equation*}
        and, appealing to Lemma~\ref{lem:stab-compute-pcomb}, we can
        deduce
        $t\in{}Sens(s'.M-\sum\limits_{t_i\in{}Pr(t,F)}pre(t_i))$,
        which directly contradicts \qedbox{
          $t\notin{}Sens(s'.M-\sum\limits_{t_i\in{}Pr(t,F)}pre(t_i))$.}
      \end{itemize}
    \end{enumerate}
  \end{itemize}
  
\end{niproof}

%%%%%%%%%%%%%%%%%%%%%%%%%%%%%%%%%%%%%%%%%%%%%%%%%%%%%%%%%%%%%%%%%%%%%%%%%%%%%%%
%%%%%%%%%% STABILIZE COMPUTE PRIORITY COMBINATION AFTER FALLING EDGE %%%%%%%%%%
%%%%%%%%%%%%%%%%%%%%%%%%%%%%%%%%%%%%%%%%%%%%%%%%%%%%%%%%%%%%%%%%%%%%%%%%%%%%%%%

\begin{lemma}[Stabilize compute priority combination after falling edge]
  \label{lem:stab-compute-pcomb}
  \fehyps{} then
  $\forall{}t\in{}T,id_t\in{}Comps(\Delta)$ s.t. $\gamma(t)=id_t$,
  $\forall{}T_s,F,~Fset\subseteq{}T$ assume that:
  \begin{itemize}
  \item $t\in{}Firable(s')$
  \item $\nexists{}t'\in{}T_s~s.t.~t'\succ{}t$
  \item EH: $\forall{}t'\in{}T,id_{t'}\in{}Comps(\Delta)$
    s.t. $\gamma(t')=id_{t'}$,\\
    $(t'\in{}F\Rightarrow\sigma'(id_{t'})(\texttt{fired})=\mathtt{true})$
    $\land$
    $(\sigma'(id_{t'})(\texttt{fired})=\mathtt{true}\Rightarrow{}t'\in{}F~\lor~{}t'\in{}T_s)$.

  \end{itemize}
  then
  $t\in{}Sens(s'.M-\sum\limits_{t_i\in{}Pr(t,F)}pre(t_i))\Leftrightarrow\sigma'(id_t)(\texttt{spc})=\mathtt{true}$
\end{lemma}

\begin{niproof}

  Given a $t\in{}T$ and an $id_t\in{}Comps(\Delta)$
  s.t. $\gamma(t)=id_t$, a $T_s$, $F$, $Fset\subseteq{}T$ and assuming
  \begin{itemize}
  \item $t\in{}Firable(s')$
  \item $\nexists{}t'\in{}T_s~s.t.~t'\succ{}t$
  \item EH: $\forall{}t'\in{}T,id_{t'}\in{}Comps(\Delta)$
    s.t. $\gamma(t')=id_{t'}$,\\
    $(t'\in{}F\Rightarrow\sigma'(id_{t'})(\texttt{fired})=\mathtt{true})$
    $\land$
    $(\sigma'(id_{t'})(\texttt{fired})=\mathtt{true}\Rightarrow{}t'\in{}F~\lor~{}t'\in{}T_s)$.
  \end{itemize}
  
  let us show
  \begin{frameb}
    $t\in{}Sens(s'.M-\sum\limits_{t_i\in{}Pr(t,F)}pre(t_i))\Leftrightarrow\sigma'(id_t)(\texttt{spc})=\mathtt{true}$.
  \end{frameb}

  \exT{}
  
  By property of the stabilize relation, \InCsCompT, and through the
  examination of the \texttt{priority_authorization_evaluation}
  process defined in the \texttt{transition} design architecture, we
  can deduce:
  \begin{equation*}
    \sigma'(id_t)("spc")=\prod\limits_{i=0}^{\Delta(id_t)(\texttt{ian})-1}\sigma'(id_t)(\texttt{pauths})[i]\label{eq:frd-eq-spc-prod-pauths}
  \end{equation*}

  Rewriting the goal with the above equation:  
  \begin{frameb}
    $t\in{}Sens(s'.M-\sum\limits_{t_i\in{}Pr(t,F)}pre(t_i))\Leftrightarrow\prod\limits_{i=0}^{\Delta(id_t)(\texttt{ian})-1}\sigma'(id_t)(\texttt{pauths})[i]=\mathtt{true}$.
  \end{frameb}
  
  Then, the proof is in two parts:
  \begin{enumerate}
  \item
    $t\in{}Sens(s'.M-\sum\limits_{t_i\in{}Pr(t,F)}pre(t_i))\Rightarrow\prod\limits_{i=0}^{\Delta(id_t)(\texttt{ian})-1}\sigma'(id_t)(\texttt{pauths})[i]=\mathtt{true}$
  \item
    $\prod\limits_{i=0}^{\Delta(id_t)(\texttt{ian})-1}\sigma'(id_t)(\texttt{pauths})[i]=\mathtt{true}\Rightarrow{}t\in{}Sens(s'.M-\sum\limits_{t_i\in{}Pr(t,F)}pre(t_i))$
  \end{enumerate}

  Let us prove both sides of the equivalence:
  \begin{enumerate}
  \item\label{item:stab-comp-spc-fst-case} Assuming that
    $t\in{}Sens(s'.M-\sum\limits_{t_i\in{}Pr(t,F)}pre(t_i))$, let us
    show\\
    \fbox{$\prod\limits_{i=0}^{\Delta(id_t)(\texttt{ian})-1}\sigma'(id_t)(\texttt{pauths})[i]=\mathtt{true}$.}

    Let us perform case analysis on $input(t)$; there are 2 cases:
    \begin{itemize}
    \item \textbf{CASE} $input(t)=\emptyset$:

      By construction,
      ${<}\mathtt{input\_arcs\_number\Rightarrow}{}1{>}\in{}g_t$ and
      ${<}$\texttt{priority\_authorizations(0)}$\Rightarrow{}\mathtt{true}{>}\in{}i_t$.

      By property of the elaboration relation, we have
      $\Delta(id_t)(\texttt{ian})=1$, and by property of the stabilize
      relation, we have
      $\sigma'(id_t)(\texttt{pauths})[0]=\mathtt{true}$.
      
      Rewriting the goal with $\Delta(id_t)(\texttt{ian})=1$ and
      $\sigma'(id_t)(\texttt{pauths})[0]=\mathtt{true}$, and
      simplifying the goal: \qedbox{tautology.}
      
    \item \textbf{CASE} $input(t)\neq{}\emptyset$:

      Then, let us show an equivalent goal:\\
      \fbox{$\forall{}i\in[0,\Delta(id_t)(\texttt{ian})-1],~\sigma'(id_t)(\texttt{pauths})[i]=\mathtt{true}$.}

      Given an $i\in{}[0,\Delta(id_t)(\texttt{ian})-1]$, let us show
      \fbox{$\sigma'(id_t)(\texttt{pauths})[i]=\mathtt{true}$.}

      By construction,
      ${<}\mathtt{input\_arcs\_number\Rightarrow}{}\vert{}input(t)\vert{>}\in{}g_t$.

      By property of the elaboration relation, we have
      $\Delta(id_t)(\texttt{ian})=\vert{}input(t)\vert$. Then, we can deduce
      $i\in{}[0,\vert{}input(t)\vert-1]$.
      
      By construction, for all $i\in{}[0,\vert{}input(t)\vert-1]$,
      there exist a $p\in{}input(t)$ and an $id_p\in{}Comps(\Delta)$
      s.t. $\gamma(p)=id_p$, there exist a $g_p$, $i_p$, $o_p$
      s.t. $\mathtt{comp}(id_p,$ $\texttt{place},$ $g_p,$ $i_p,$
      $o_p)\in{}d.cs$, and there exist a
      $j\in{}[0,\vert{}output(p)\vert]$ and an
      $id_{ji}\in{}Sigs(\Delta)$ s.t.\\
      ${<}\mathtt{input\_arcs\_valid(i)\Rightarrow{}id_{ji}}{>}\in{}i_t$
      and
      ${<}\mathtt{output\_arcs\_valid(j)\Rightarrow{}id_{ji}}{>}\in{}o_t$.
      Let us take such a $p\in{}input(t)$, $id_p\in{}Comps(\Delta)$,
      $g_p$, $i_p$, $o_p$, $j\in{}[0,$ $\vert{}output(p)\vert]$ and
      $id_{ji}\in{}Sigs(\Delta)$.\\

      Now, let us perform case analysis on the nature of the arc
      connecting $p$ and $t$; there are 2 cases:
      
      \begin{itemize}
      \item \textbf{CASE} $pre(p,t)=(\omega,\mathtt{test})$ or
        $pre(p,t)=(\omega,\mathtt{inhib})$:

        By construction,
        ${<}\mathtt{priority\_authorizations(i)\Rightarrow{}true}{>}\in{}i_t$,
        and by property of the stabilize relation:
        \qedbox{$\sigma'(id_t)(\texttt{pauths})[i]=\mathtt{true}$.}

      \item \textbf{CASE} $pre(p,t)=(\omega,\mathtt{basic})$:

        Let us define
        $output_c(p)=\{t\in{}T~\vert~\exists{}\omega,~pre(p,t)=(\omega,\mathtt{basic})\}$,
        the set of output transitions of $p$ that are in
        conflict. Then, there are two cases, one for each way to solve
        the conflicts between the output transitions of $p$:

        \begin{itemize}
        \item \textbf{CASE} For all pair of transitions in
          $output_c(p)$, all conflicts are solved by mutual exclusion:

          By construction,
          ${<}\mathtt{priority\_authorizations(i)\Rightarrow{}true}{>}\in{}i_t$,
          and by property of the stabilize relation:
          \qedbox{$\sigma'(id_t)(\texttt{pauths})[i]=\mathtt{true}$.}
        \item \textbf{CASE} The priority relation is a strict total
          order over the set $output_c(p)$:
        \end{itemize}
        By construction, there exists an $id'_{ji}\in{}Sigs(\Delta)$
        s.t.\\
        ${<}\mathtt{priority\_authorizations(i)\Rightarrow{}id'_{ji}}{>}\in{}i_t$
        and\\
        ${<}\mathtt{priority\_authorizations(j)\Rightarrow{}id'_{ji}}{>}\in{}o_p$.

        By property of the stabilize relation, \InCsCompT{} and
        \InCsCompP, we can deduce:
        \begin{equation*}
          \label{eq:frd-eq-tpauthsi-ppauthsj}\sigma'(id_t)(\texttt{pauths})[i]=\sigma'(id'_{ji})=\sigma'(id_p)(\texttt{pauths})[j]\\
        \end{equation*}

        Rewriting the goal with the above equation:
        \fbox{$\sigma'(id_p)(\texttt{pauths})[j]=\mathtt{true}$.}

        By property of the stabilize relation, \InCsCompP, and
        through the examination of the \texttt{priority_evaluation}
        process defined in the \texttt{place} design behavior, we
        can deduce:
        \begin{equation}
          \label{eq:frd-eq-pauthsj}
          \sigma'(id_p)(\texttt{pauths})[j]=(\sigma'(id_p)(\texttt{sm})\ge{}\mathtt{vsots}+\sigma'(id_p)(\texttt{oaw})[j])
        \end{equation}

        Let us define the $\mathtt{vsots}$ term as follows:
        \begin{equation}
          \label{eq:frd-vsots}
          \mathtt{vsots}=\sum\limits_{i=0}^{j-1}
          \begin{cases}
            \sigma'(id_p)(\texttt{oaw})[i]~\mathtt{if}~\sigma'(id_p)(\texttt{otf})[i].\\
            \hspace{19ex}\sigma'(id_p)(\texttt{oat})[i]=\mathtt{basic}\\
            0~otherwise\\
          \end{cases}
        \end{equation}

        Rewriting the goal with \eqref{eq:frd-eq-pauthsj}:
        \fbox{$\sigma'(id_p)(\texttt{sm})\ge{}\mathtt{vsots}+\sigma'(id_p)(\texttt{oaw})[j]$}

        By definition of
        $t\in{}Sens(s'.M-\sum\limits_{t_i\in{}Pr(t,F)}pre(t_i))$, we
        can deduce:\\
        $s'.M(p)\ge{}\sum\limits_{t_i\in{}Pr(t,F)}pre(p,t_i)+\omega$.
        
        Then, there are three points to prove:
        \begin{enumerate}
        \item \fbox{$s'.M(p)=\sigma'(id_p)(\texttt{sm})$}
        \item \fbox{$\omega=\sigma'(id_p)(\texttt{oaw})[j]$}
        \item \fbox{$\sum\limits_{t_i\in{}Pr(t,F)}pre(p,t_i)=\mathtt{vsots}$}
        \end{enumerate}

        Let us prove these three points:
        \begin{enumerate}
        \item \fbox{$s'.M(p)=\sigma'(id_p)(\texttt{sm})$}

          Appealing to Lemma~\ref{lem:fe-equal-marking},
          \qedbox{$s'.M(p)=\sigma'(id_p)(\texttt{sm})$.}
        \item \fbox{$\omega=\sigma'(id_p)(\texttt{oaw})[j]$}

          By construction, and as
          $pre(p,t)=(\omega,\mathtt{basic})$, we know that
          ${<}$\texttt{output\_arcs\_weights(j)}$\Rightarrow\omega{>}\in{}i_p$.

          By property of the stabilize relation and \InCsCompP:\\
          \qedbox{$\omega=\sigma'(id_p)(\texttt{oaw})[j]$.}
          
        \item
          \fbox{$\sum\limits_{t_i\in{}Pr(t,F)}pre(p,t_i)=\mathtt{vsots}$}
        \end{enumerate}
        
        Let us replace the left and right term of the equality by
        their full definition:

        \begin{frameb}
          \begin{tabular}{c}
            $\sum\limits_{t_i\in{}Pr(t,F)}
            \begin{cases}
              \omega~\mathtt{if}~pre(p,t_i)=(\omega,\mathtt{basic})\\
              0~otherwise
            \end{cases}$ \\
            $=$ \\
            $\sum\limits_{i=0}^{j-1}
            \begin{cases}
              \sigma'(id_p)(\texttt{oaw})[i]~\mathtt{if}~\sigma'(id_p)(\texttt{otf})[i].\\
              \hspace{19ex}\sigma'(id_p)(\texttt{oat})[i]=\mathtt{basic}\\
              0~otherwise\\
            \end{cases}$ \\
          \end{tabular}
        \end{frameb}

        Now, we must reason on the priority status of transition
        $t$ regarding the group of conflicting output transitions
        of $p$. There 2 cases:
          
        \begin{itemize}
        \item \textbf{CASE} $t$ is the top-priority transition in
          the group of conflicting output transitions of $p$:

          In that case, the set $Pr(t,F)$ is empty and, by
          construction, $j=0$. Thus, the goal is a tautology
          \qedbox{$0=0$.}
          
        \item \textbf{CASE} $t$ is not the top-priority transition
          in the group of conflicting output transitions of $p$:
        \end{itemize}
        In that case, we know that there is a least one element
        in $Pr(t,F)$ and the index $j>0$.

        Let us replace the sum terms in the goal by equivalent terms:
        \begin{frameb}
          \begin{tabular}{c}
            $\sum\limits_{t_i\in{}Pr_p}
            \begin{cases}
              \omega~\mathtt{if}~pre(p,t_i)=(\omega,\mathtt{basic})~\mathtt{and}~t_i\in{}F\\
              0~otherwise
            \end{cases}$ \\
            $=$ \\
            $\sum\limits_{i\in{}IPr_p}
            \begin{cases}
              \sigma'(id_p)(\texttt{oaw})[i]~\mathtt{if}~\sigma'(id_p)(\texttt{otf})[i]\\
              0~otherwise\\
            \end{cases}$ \\
          \end{tabular}
        \end{frameb}

        Let us define the set $Pr_p$ as
        \begin{center}
          $Pr_p=\{t_i~\vert~t_i\succ{}t\land\exists{}\omega~s.t.~pre(p,t_i)=(\omega,\mathtt{basic})\}$
        \end{center}
        and set $IPr_p$ as
        \begin{center}
          $IPr_p=\{i~\vert~i\in[0,j-1]\land\sigma'(id_p)(\texttt{oat})[i]=\texttt{basic}\}$
        \end{center}

        Let us define $f(t_i)$ as
        \begin{center}
          $f(t_i)=\begin{cases}
            \omega~\mathtt{if}~pre(p,t_i)=(\omega,\mathtt{basic})~\mathtt{and}~t_i\in{}F\\
            0~otherwise
          \end{cases}$
        \end{center}
        and $g(i)$ as              
        \begin{center}
          $g(i)=\begin{cases}
            \sigma'(id_p)(\texttt{oaw})[i]~\mathtt{if}~\sigma'(id_p)(\texttt{otf})[i]\\
            0~otherwise\\
          \end{cases}$
        \end{center}

        then, we must prove
        \fbox{$\sum\limits_{t_i\in{}Pr_p}f(t_i)=\sum\limits_{i\in{}IPr_p}g(i)$.}

        To prove the above equality, it is sufficient to prove that
        there exists a bijection $\beta$ from $Pr_p$ to $IPr_p$ such
        that for all $t_i\in{}Pr_p$, $f(t_i)=g(\beta(t_i))$. Let us
        use the function $\beta$ that takes a $t_i\in{}Pr_p$ and
        yields the index denoting the position of $t_i$ in the
        priority-ordered version of set $Pr_p$. Since we assumed that
        a total order existed over the conflicting output transitions
        of place $p$, then there exists a total ordering of the
        transitions of set $Pr_p$. By property of the \hilecop{}
        transformation function, we know that the index returned by
        the function $\beta$ belongs to the interval $[0,j-1]$ and
        verifies $\sigma'(id_p)(\texttt{oat})[i]=\texttt{basic}$.
        Given a $t_i\in{}Pr_p$, we must show
        \fbox{$f(t_i)=g(\beta(t_i))$.}

        Let us unfold terms $f(t_i)$ and $g(\beta(t_i))$ to their full definition:
        \begin{frameb}
          \begin{tabular}{c}
            $\begin{cases}
              \omega~\mathtt{if}~pre(p,t_i)=(\omega,\mathtt{basic})~\mathtt{and}~t_i\in{}F\\
              0~otherwise
            \end{cases}$ \\
            $=$ \\
            $\begin{cases}
              \sigma'(id_p)(\texttt{oaw})[\beta(t_i)]~\mathtt{if}~\sigma'(id_p)(\texttt{otf})[\beta(t_i)]\\
              0~otherwise\\
            \end{cases}$ \\
          \end{tabular}

          By construction, there exists an
          $id_{t_i}\in{}Comps(\Delta)$ such that
          $\gamma(t_i)=id_{t_i}$, and there exist $g_{t_i}$, $i_{t_i}$
          and $o_{t_i}$ such that
          $\mathtt{comp}(id_{t_i},\texttt{transition},g_{t_i},i_{t_i},o_{t_i})\in{}d.cs$.
          
          By property of the function $\beta$ and by construction, we
          can deduce that the element of index $\beta(t_i)$ of the
          \texttt{otf} output port of PCI $id_p$ is connected the
          \texttt{fired} output port of TCI $id_{t_i}$.
        \end{frameb}
      \end{itemize}
    \end{itemize}
    
  \item Assuming that
    $\prod\limits_{i=0}^{\Delta(id_t)("ian")-1}\sigma'(id_t)("pauths")[i]=\mathtt{true}$,
    let us show\\
    \fbox{$t\in{}Sens(s'.M-\sum\limits_{t_i\in{}Pr(t,fired)}pre(t_i))$.}
    
    By definition of $t\in{}Sens(s'.M-\sum\limits_{t_i\in{}Pr(t,fired)}pre(t_i))$:

    \begin{frameb}
      $\forall{}p\in{}P,\omega\in\mathbb{N}^{*},~$\\
      $\big((pre(p,t)=(\omega,\mathtt{basic})\lor{}pre(p,t)=(\omega,\mathtt{test}))\Rightarrow$
      ${}s'.M(p)-\sum\limits_{t_i\in{}Pr(t,fired)}pre(p,t_i)\ge\omega\big)$\\
      $\land{}$ $\big(pre(p,t)=(\omega,\mathtt{inhib})\Rightarrow$
      ${}s'.M(p)-\sum\limits_{t_i\in{}Pr(t,fired)}pre(p,t_i)<\omega\big)$
    \end{frameb}

    Given a $p\in{}P$ and an $\omega\in{}\mathbb{N}^{*}$, let us show 
    \begin{frameb}
      $\big((pre(p,t)=(\omega,\mathtt{basic})\lor{}pre(p,t)=(\omega,\mathtt{test}))\Rightarrow$
      ${}s'.M(p)-\sum\limits_{t_i\in{}Pr(t,fired)}pre(p,t_i)\ge\omega\big)$\\
      $\land{}$ $\big(pre(p,t)=(\omega,\mathtt{inhib})\Rightarrow$
      ${}s'.M(p)-\sum\limits_{t_i\in{}Pr(t,fired)}pre(p,t_i)<\omega\big)$
    \end{frameb}

    By construction, there exists an $id_p\in{}Comps(\Delta)$
    s.t. $\gamma(p)=id_p$. \exP{}
    
    There are three different cases:
    \begin{enumerate}
    \item Assuming that $pre(p,t)=(\omega,\mathtt{test})$, let us show
      \fbox{$s'.M(p)-\sum\limits_{t_i\in{}Pr(t,fired)}pre(p,t_i)\ge\omega$.}

      Then, assuming that the priority relation is well-defined, there
      exists no transition $t_i$ connected by a $\mathtt{basic}$ arc
      to $p$ that verifies $t_i\succ{}t$. This is because $t$ is
      connected to $p$ by a $\mathtt{test}$ arc; thus, $t$ is not in
      conflict with the other output transitions of $p$; thus, there
      is no relation of priority between $t$ and the output of $p$.

      Then, we can deduce that
      $\sum\limits_{t_i\in{}Pr(t,fired)}pre(p,t_i)=0$.

      Then, the new goal is $s'.M(p)\ge\omega$.

      Knowing that $t\in{}Firable(s')$, thus, $t\in{}Sens(s'.M)$,
      thus, we have \qedbox{$s'.M(p)\ge\omega$.}
      
    \item Assuming that $pre(p,t)=(\omega,\mathtt{inhib})$, let us
      show
      \fbox{$s'.M(p)-\sum\limits_{t_i\in{}Pr(t,fired)}pre(p,t_i)<\omega$.}

      Use the same strategy as above.
      
    \item Assuming that $pre(p,t)=(\omega,\mathtt{basic})$, let us
      show
      \fbox{$s'.M(p)-\sum\limits_{t_i\in{}Pr(t,fired)}pre(p,t_i)\ge\omega$.}

      Then, there are two cases:

      \begin{enumerate}
      \item \textbf{CASE} For all pair of transitions in
        $output_c(p)$, all conflicts are solved by mutual exclusion.

        Then, assuming that the priority relation is well-defined, it
        must not be defined over the set $output_c(t)$, and we know
        that $t\in{}output_c(p)$ since
        $pre(p,t)=(\omega,\mathtt{basic})$.

        Then, there exists no transition $t_i$ connected to $p$ by a
        $\mathtt{basic}$ arc that verifies $t_i\succ{}t$.

        Then, we can deduce $\sum\limits_{t_i\in{}Pr(t,fired)}pre(p,t_i)=0$.
        
        Then, the new goal is $s'.M(p)\ge\omega$.

        We know $t\in{}Firable(s')$, thus, $t\in{}Sens(s'.M)$, thus,
        \qedbox{$s'.M(p)\ge\omega$.}
        
      \item \textbf{CASE} The priority relation is a strict total
        order over the set $output_c(p)$.
        
        By construction, there exists $id_t\in{}Comps(\Delta)$
        s.t. $\gamma(t)=id_t$. \exT{}

        By construction, there exist
        $j\in{}[0,\vert{}input(t)\vert-1]$,
        $k\in[0,\vert{}output(t)\vert-1]$, and
        $id_{kj}\in{}Sigs(\Delta)$ s.t.
        ${<}\mathtt{priority\_authorizations(j)\Rightarrow{}id_{kj}}{>}\in{}ipm_t$
        and\\
        ${<}\mathtt{priority\_authorizations(k)\Rightarrow{}id_{kj}}{>}\in{}opm_p$.
        Let us take such an $j$, $k$ and $id_{kj}$.

        From
        $\prod\limits_{i=0}^{\Delta(id_t)("ian")-1}\sigma'(id_t)("pauths")[i]=\mathtt{true}$,
        we can deduce that for all $i\in{}[0,$ $\Delta(id_t)$
        $("ian")-1]$, $\sigma'(id_t)("pauths")[i]=\mathtt{true}$.

        By construction,
        ${<}\mathtt{input\_arcs\_number\Rightarrow}{}\vert{}input(t)\vert{>}\in{}gm_t$,
        and by property of the elaboration relation, we have
        $\Delta(id_t)("ian")=\vert{}input(t)\vert$. Then, from
        $j\in{}[0,\vert{}input(t)\vert-1]$, we can deduce
        $j\in[0,\Delta(id_t)("ian")-1]$. And, from
        $\forall{}i\in{}[0,\Delta(id_t)("ian")-1],~\sigma'(id_t)$ $("pauths")[i]=$ $\mathtt{true}$,
        we can deduce $\sigma'(id_t)("pauths")[j]=\mathtt{true}$.
        
        By property of the stabilize relation, \InCsCompP{} and
        \InCsCompT{}:
        \begin{equation}
          \label{eq:frd-eq-pauthsk-pauthsj}
          \sigma'(id_p)("pauths")[k]=\sigma'(id_{kj})=\sigma'(id_t)("pauths")[j]=\mathtt{true}
        \end{equation}

        By property of the stabilize relation and \InCsCompP:
        \begin{equation}
          \label{eq:frd-eq-pauthsk}
          \sigma'(id_p)("pauths")[k]=(\sigma'(id_p)("sm")\ge{}\mathtt{vsots}+\sigma'(id_p)("oaw")[k])
        \end{equation}

        Let us define the $\mathtt{vsots}$ term as follows:
        \begin{equation}
          \label{eq:frd-vsots-snd-case}
          \mathtt{vsots}=\sum\limits_{i=0}^{k-1}
          \begin{cases}
            \sigma'(id_p)("oaw")[i]~\mathtt{if}~\sigma'(id_p)("otf")[i].\\
            \hspace{19ex}\sigma'(id_p)("oat")[i]=\mathtt{basic}\\
            0~otherwise\\
          \end{cases}
        \end{equation}
        
        From \eqref{eq:frd-eq-pauthsk-pauthsj} and
        \eqref{eq:frd-eq-pauthsk}, we can deduce that
        $\sigma'(id_p)("sm")\ge{}\mathtt{vsots}+\sigma'(id_p)("oaw")[k]$.

        Then, there are three points to prove:
        \begin{enumerate}
        \item \fbox{$s'.M(p)=\sigma'(id_p)("sm")$}
        \item \fbox{$\omega=\sigma'(id_p)("oaw")[k]$}
        \item \fbox{$\sum\limits_{t_i\in{}Pr(t,fired)}pre(p,t_i)=\mathtt{vsots}$}
        \end{enumerate}
        
        See \ref{item:stab-comp-spc-fst-case} for the remainder of the
        proof.
        
      \end{enumerate}
      
    \end{enumerate}
  \end{enumerate}
  
\end{niproof}

%%%%%%%%%%%%%%%%%%%%%%%%%%%%%%%%%%%%%%%%%%%%%%%%%%%%%%%%
%%%%%%%%%% FALLING EDGE NOT EQUAL FIRED LEMMA %%%%%%%%%%
%%%%%%%%%%%%%%%%%%%%%%%%%%%%%%%%%%%%%%%%%%%%%%%%%%%%%%%%

\begin{lemma}[Falling Edge Equal Not Fired]
  \label{lem:fe-equal-not-fired}
  \fehyps{} then $\forall{}t,id_t~s.t.~\gamma(t)=id_t,$
  $~t\notin{}Fired(s')\Leftrightarrow\sigma'_t("fired")=\mathtt{false}$.
\end{lemma}

\begin{proof}
  Proving the above lemma is trivial by appealing to
  Lemma~\nameref{lem:fe-equal-fired} and by reasoning on
  contrapositives.
\end{proof}

%%% Local Variables:
%%% mode: latex
%%% TeX-master: "../../main"
%%% End:


%%% Local Variables:
%%% mode: latex
%%% TeX-master: "../main"
%%% End:



%----------------------------------------------------------------------------------------
%	BIBLIOGRAPHY
%----------------------------------------------------------------------------------------

\printbibliography[heading=bibintoc]

%----------------------------------------------------------------------------------------

\end{document}  

